\input{preamble}

% OK, start here.
%
\begin{document}
\setlength{\emergencystretch}{4em}
\title{Compléments sur les morphismes}


\maketitle

\phantomsection
\label{section-phantom}

\tableofcontents

\section{Introduction}
\label{section-introduction}

\noindent
Dans ce chapitre, nous poursuivons l'étude des propriétés des morphismes de schémas.
Une référence fondamentale est \cite{EGA}.






\section{Épaississements}
\label{section-thickenings}

\noindent
La terminologie suivante n'est peut-être pas entièrement usuelle, mais elle est commode.

\begin{definition}
\label{definition-thickening}
Épaississements.
\begin{enumerate}
\item On dit qu'un schéma $X'$ est un {\it épaississement} d'un schéma $X$ si
$X$ est un sous-schéma fermé de $X'$ et si les espaces topologiques sous-jacents
sont égaux.
\item On dit qu'un schéma $X'$ est un {\it épaississement du premier ordre} d'un schéma $X$ si
$X$ est un sous-schéma fermé de $X'$ et si le faisceau quasi-cohérent d'idéaux
$\mathcal{I} \subset \mathcal{O}_{X'}$ définissant $X$ est de carré nul.
\item On dit qu'un schéma $X'$ est un {\it épaississement d'ordre fini} d'un schéma $X$
si $X$ est un sous-schéma fermé de $X'$ et si le faisceau quasi-cohérent d'idéaux
$\mathcal{I} \subset \mathcal{O}_{X'}$ définissant $X$ est nilpotent, c'est-à-dire
s'il existe un entier $n \geq 0$ tel que $\mathcal{I}^{n + 1} = 0$.
\item On dit qu'un schéma $X'$ est un {\it épaississement d'ordre $n$} d'un schéma $X$
si $X$ est un sous-schéma fermé de $X$ et si $\mathcal{I}^{n + 1} = 0$,
où $\mathcal{I} \subset \mathcal{O}_{X'}$ est le
faisceau quasi-cohérent d'idéaux définissant $X$.
\item Étant donnés deux épaississements $X \subset X'$ et $Y \subset Y'$, un
{\it morphisme d'épaississements} est un morphisme $f' : X' \to Y'$ tel que
$f'(X) \subset Y$, c'est-à-dire tel que $f'|_X$ se factorise par le sous-schéma
fermé $Y$. Dans cette situation, on pose $f = f'|_X : X \to Y$ et on dit
que $(f, f') : (X \subset X') \to (Y \subset Y')$ est un morphisme
d'épaississements.
\item Soit $S$ un schéma. On définit de même les {\it épaississements sur $S$} et
les {\it morphismes d'épaississements sur $S$}. Cela signifie que les schémas
$X, X', Y, Y'$ ci-dessus sont des schémas sur $S$ et que les morphismes
$X \to X'$, $Y \to Y'$ et $f' : X' \to Y'$ sont des morphismes sur $S$.
\end{enumerate}
\end{definition}

\noindent
Soit $i_X : X \to X'$ un épaississement. Toute section locale du noyau
$\mathcal{I} = \Ker(i_X^\sharp)$ est localement nilpotente. Cela donne
un certain contrôle sur le faisceau structural $\mathcal{O}_{X'}$, mais,
techniquement, la classe des épaississements d'ordre fini $X \subset X'$
est beaucoup plus facile à manier. En effet, si $X'$ est un épaississement d'ordre $n$,
on dispose d'une filtration
$$
0 = \mathcal{I}^{n + 1} \subset
\mathcal{I}^n \subset
\mathcal{I}^{n - 1} \subset \ldots \subset
\mathcal{I} \subset \mathcal{O}_{X'}
$$
et l'on voit que $X'$ est filtré par des sous-espaces fermés
$$
X = X_1 \subset X_2 \subset \ldots \subset X_n \subset X_{n + 1} = X'
$$
tels que chaque paire $X_i \subset X_{i + 1}$ soit un épaississement du premier ordre
sur $S$. Par de simples raisonnements par récurrence, de nombreux résultats établis pour
les épaississements du premier ordre se reformulent en résultats sur les épaississements d'ordre fini.

\medskip\noindent
Les épaississements du premier ordre se décrivent comme suit (voir
Modules, lemme \ref{modules-lemma-double-structure-gives-derivation}).

\begin{lemma}
\label{lemma-first-order-thickening}
Soit $X$ un schéma sur une base $S$. Considérons une suite exacte courte
$$
0 \to \mathcal{I} \to \mathcal{A} \to \mathcal{O}_X \to 0
$$
de faisceaux sur $X$, où $\mathcal{A}$ est un faisceau de
$f^{-1}\mathcal{O}_S$-algèbres,
$\mathcal{A} \to \mathcal{O}_X$ est une surjection
de faisceaux de $f^{-1}\mathcal{O}_S$-algèbres et $\mathcal{I}$ en est le noyau.
Si
\begin{enumerate}
\item $\mathcal{I}$ est un idéal de carré nul dans $\mathcal{A}$ et
\item $\mathcal{I}$ est quasi-cohérent en tant que $\mathcal{O}_X$-module,
\end{enumerate}
alors $X' = (X, \mathcal{A})$ est un schéma et $X \to X'$ est un épaississement
du premier ordre sur $S$. De plus, tout épaississement du premier ordre sur
$S$ est de cette forme.
\end{lemma}

\begin{proof}
Il est clair que $X'$ est un espace localement annelé. Soit $U = \Spec(B)$
un ouvert affine de $X$. Posons $A = \Gamma(U, \mathcal{A})$. Remarquons que,
puisque $H^1(U, \mathcal{I}) = 0$ (voir Cohomologie des schémas, lemme
\ref{coherent-lemma-quasi-coherent-affine-cohomology-zero}),
l'application $A \to B$ est surjective. Par hypothèse, le noyau
$I = \mathcal{I}(U)$ est un idéal de carré nul de l'anneau $A$.
D'après
Schémas, lemme \ref{schemes-lemma-morphism-into-affine},
il existe un morphisme canonique d'espaces localement annelés
$$
(U, \mathcal{A}|_U) \longrightarrow \Spec(A)
$$
provenant de l'application $B \to \Gamma(U, \mathcal{A})$. Ce morphisme
s'inscrit dans le diagramme commutatif
$$
\xymatrix{
(U, \mathcal{O}_X|_U) \ar[d] \ar[r] & \Spec(B) \ar[d] \\
(U, \mathcal{A}|_U) \ar[r] & \Spec(A)
}
$$
et l'on voit donc qu'il est un homéomorphisme sur les espaces topologiques sous-jacents.
Pour montrer que c'est un isomorphisme, il suffit donc de vérifier qu'il induit
un isomorphisme sur les anneaux locaux.
Pour $u \in U$ correspondant à l'idéal premier $\mathfrak p \subset A$,
on obtient un diagramme commutatif de suites exactes courtes
$$
\xymatrix{
0 \ar[r] &
I_{\mathfrak p} \ar[r] \ar[d] &
A_{\mathfrak p} \ar[r] \ar[d] &
B_{\mathfrak p} \ar[r] \ar[d] & 0 \\
0 \ar[r] &
\mathcal{I}_u \ar[r] &
\mathcal{A}_u \ar[r] &
\mathcal{O}_{X, u} \ar[r] & 0.
}
$$
Les flèches verticales de gauche et de droite sont des isomorphismes, car
$\mathcal{I}$ et $\mathcal{O}_X$ sont des faisceaux quasi-cohérents.
La flèche du milieu est donc elle aussi un isomorphisme. Ainsi, tout point
de $X' = (X, \mathcal{A})$ possède un voisinage affine et $X'$ est un
schéma, comme voulu.
\end{proof}

\begin{lemma}
\label{lemma-thickening-affine-scheme}
\begin{slogan}
Le caractère affine est invariant par épaississement
\end{slogan}
\begin{reference}
Le cas d'un épaississement d'ordre fini est \cite[Proposition 5.1.9]{EGA1}.
\end{reference}
Tout épaississement d'un schéma affine est affine.
\end{lemma}

\begin{proof}
C'est un cas particulier de
Limites, proposition \ref{limits-proposition-affine}.
\end{proof}

\begin{proof}[Démonstration pour un épaississement d'ordre fini]
Supposons que $X \subset X'$ soit un épaississement d'ordre fini avec $X$ affine.
On peut alors appliquer le critère de Serre pour prouver que $X'$ est affine. Plus précisément,
on utilisera Cohomologie des schémas, lemme
\ref{coherent-lemma-quasi-compact-h1-zero-covering}. Soit $\mathcal{F}$ un
$\mathcal{O}_{X'}$-module quasi-cohérent. Il suffit de montrer que
$H^1(X', \mathcal{F}) = 0$. Notons $i : X \to X'$ l'immersion fermée donnée
et posons
$\mathcal{I} = \Ker(i^\sharp : \mathcal{O}_{X'} \to i_*\mathcal{O}_X)$.
D'après la discussion des épaississements d'ordre fini (qui suit
la définition \ref{definition-thickening}), il existe un entier $n \geq 0$
et une filtration
$$
0 = \mathcal{F}_{n + 1} \subset \mathcal{F}_n \subset
\mathcal{F}_{n - 1} \subset \ldots \subset
\mathcal{F}_0 = \mathcal{F}
$$
par des sous-modules quasi-cohérents tels que $\mathcal{F}_a/\mathcal{F}_{a + 1}$ soit
annulé par $\mathcal{I}$. On peut en effet prendre
$\mathcal{F}_a = \mathcal{I}^a\mathcal{F}$. Alors
$\mathcal{F}_a/\mathcal{F}_{a + 1} = i_*\mathcal{G}_a$ pour un certain
$\mathcal{O}_X$-module quasi-cohérent $\mathcal{G}_a$, voir Morphismes, lemme
\ref{morphisms-lemma-i-star-equivalence}. On obtient
$$
H^1(X', \mathcal{F}_a/\mathcal{F}_{a + 1}) =
H^1(X', i_*\mathcal{G}_a) = H^1(X, \mathcal{G}_a) = 0
$$
La deuxième égalité résulte de Cohomologie des schémas, lemme
\ref{coherent-lemma-relative-affine-cohomology},
et la dernière de Cohomologie des schémas, lemme
\ref{coherent-lemma-quasi-coherent-affine-cohomology-zero}.
Ainsi, $\mathcal{F}$ possède une filtration finie dont les quotients successifs
ont leur premier groupe de cohomologie nul ; un simple
raisonnement par récurrence montre donc que $H^1(X', \mathcal{F}) = 0$.
\end{proof}

\begin{lemma}
\label{lemma-base-change-thickening}
Soit $S \subset S'$ un épaississement de schémas. Soit $X' \to S'$ un morphisme
et posons $X = S \times_{S'} X'$. Alors $(X \subset X') \to (S \subset S')$
est un morphisme d'épaississements. Si $S \subset S'$ est un épaississement du premier
ordre (resp.\ d'ordre fini), alors $X \subset X'$ est un épaississement du premier
ordre (resp.\ d'ordre fini).
\end{lemma}

\begin{proof}
Omise.
\end{proof}

\begin{lemma}
\label{lemma-composition-thickening}
\begin{slogan}
Les composés d'épaississements sont des épaississements
\end{slogan}
Si $S \subset S'$ et $S' \subset S''$ sont des épaississements, alors
$S \subset S''$ en est un également.
\end{lemma}

\begin{proof}
Omise.
\end{proof}

\begin{lemma}
\label{lemma-descending-property-thickening}
La propriété d'être un épaississement est locale pour la topologie fpqc.
Il en va de même pour les épaississements du premier ordre.
\end{lemma}

\begin{proof}
L'énoncé signifie ceci : soit $X \to X'$ un morphisme
de schémas et soit $\{g_i : X'_i \to X'\}$
un recouvrement fpqc tel que le changement de base $X_i \to X'_i$
soit un épaississement pour tout $i$. Alors $X \to X'$ est un épaississement.
Puisque les morphismes $g_i$ sont conjointement surjectifs, on en déduit
que $X \to X'$ est surjectif. D'après
Descente, lemme \ref{descent-lemma-descending-property-closed-immersion},
on en déduit que $X \to X'$ est une immersion fermée.
Ainsi, $X \to X'$ est un épaississement. On omet la démonstration dans le
cas des épaississements du premier ordre.
\end{proof}









\section{Morphismes d'épaississements}
\label{section-morphisms-thickenings}

\noindent
Si $(f, f') : (X \subset X') \to (Y \subset Y')$ est un morphisme
d'épaississements de schémas, alors les propriétés du morphisme
$f$ se transmettent souvent à $f'$. Il existe plusieurs variantes.

\begin{lemma}
\label{lemma-thicken-property-morphisms}
Soit $(f, f') : (X \subset X') \to (S \subset S')$ un morphisme
d'épaississements. Alors
\begin{enumerate}
\item $f$ est un morphisme affine si et seulement si $f'$ est un morphisme affine,
\item $f$ est un morphisme surjectif si et seulement si $f'$ est un morphisme surjectif,
\item $f$ est quasi-compact si et seulement si $f'$ est quasi-compact,
\item $f$ est universellement fermé si et seulement si $f'$ est universellement fermé,
\item $f$ est entier si et seulement si $f'$ est entier,
\item $f$ est (quasi-)séparé si et seulement si $f'$ est (quasi-)séparé,
\item $f$ est universellement injectif si et seulement si $f'$ est universellement injectif,
\item $f$ est universellement ouvert si et seulement si $f'$ est universellement ouvert,
\item $f$ est quasi-affine si et seulement si $f'$ est quasi-affine, et
\item ajouter d'autres propriétés ici.
\end{enumerate}
\end{lemma}

\begin{proof}
Remarquons que $S \to S'$ et $X \to X'$ sont des homéomorphismes universels
(voir par exemple
Morphismes, lemme \ref{morphisms-lemma-reduction-universal-homeomorphism}).
Cela entraîne immédiatement (2), (3), (4), (7) et (8).
Le point (1) résulte du lemme \ref{lemma-thickening-affine-scheme},
qui établit une correspondance biunivoque entre
les ouverts affines de $S$ et de $S'$, ainsi qu'entre ceux de $X$ et de $X'$.
Le point (9) résulte de
Limites, lemme \ref{limits-lemma-thickening-quasi-affine},
et de la remarque précédente sur les ouverts affines de $S$ et de $S'$.
Le point (5) résulte de (1) et (4), d'après
Morphismes, lemme \ref{morphisms-lemma-integral-universally-closed}.
Enfin, remarquons que

$$
S \times_X S = S \times_{X'} S \to S \times_{X'} S' \to S' \times_{X'} S'
$$
est un épaississement (les deux flèches sont des épaississements d'après
le lemme \ref{lemma-base-change-thickening}).
En appliquant (3) et (4) au morphisme
$(S \subset S') \to (S \times_X S \to S' \times_{X'} S')$,
on obtient donc (6).
\end{proof}

\begin{lemma}
\label{lemma-thicken-property-relatively-ample}
Soit $(f, f') : (X \subset X') \to (S \subset S')$ un morphisme
d'épaississements. Soit $\mathcal{L}'$ un faisceau inversible sur $X'$
et notons $\mathcal{L}$ sa restriction à $X$.
Alors $\mathcal{L}'$ est $f'$-ample si et seulement si
$\mathcal{L}$ est $f$-ample.
\end{lemma}

\begin{proof}
Rappelons que l'amplitude relative se vérifie sur chaque
ouvert affine de la base, voir
Morphismes, définition \ref{morphisms-definition-relatively-ample}.
D'après le lemme \ref{lemma-thickening-affine-scheme},
il existe une correspondance biunivoque entre
les ouverts affines de $S$ et de $S'$.
On peut donc supposer que $S$ et $S'$ sont affines,
et l'on est ramené à prouver que
$\mathcal{L}'$ est ample si et seulement si
$\mathcal{L}$ est ample.
C'est Limites, lemme \ref{limits-lemma-ample-on-reduction}.
\end{proof}

\begin{lemma}
\label{lemma-thicken-property-morphisms-cartesian}
Soit $(f, f') : (X \subset X') \to (S \subset S')$ un morphisme
d'épaississements tel que $X = S \times_{S'} X'$. Si $S \subset S'$
est un épaississement d'ordre fini, alors
\begin{enumerate}
\item $f$ est une immersion fermée si et seulement si $f'$ est une immersion fermée,
\item $f$ est localement de type fini si et seulement si $f'$ est
localement de type fini,
\item $f$ est localement quasi-fini si et seulement si $f'$ est localement
quasi-fini,
\item $f$ est localement de type fini et de dimension relative $d$ si et
seulement si $f'$ est localement de type fini et de dimension relative $d$,
\item $\Omega_{X/S} = 0$ si et seulement si $\Omega_{X'/S'} = 0$,
\item $f$ est non ramifié si et seulement si $f'$ est non ramifié,
\item $f$ est propre si et seulement si $f'$ est propre,
\item $f$ est fini si et seulement si $f'$ est fini,
\item $f$ est un monomorphisme si et seulement si $f'$ est un monomorphisme,
\item $f$ est une immersion si et seulement si $f'$ est une immersion, et
\item ajouter d'autres propriétés ici.
\end{enumerate}
\end{lemma}

\begin{proof}
Les propriétés $\mathcal{P}$ énumérées dans le lemme sont toutes stables
par changement de base ; si $f'$ possède la propriété $\mathcal{P}$,
il en va donc de même de $f$. Voir
Schémas, lemmes \ref{schemes-lemma-base-change-immersion} et
\ref{schemes-lemma-base-change-monomorphism},
et
Morphismes, lemmes
\ref{morphisms-lemma-base-change-finite-type},
\ref{morphisms-lemma-base-change-quasi-finite},
\ref{morphisms-lemma-base-change-relative-dimension-d},
\ref{morphisms-lemma-base-change-differentials},
\ref{morphisms-lemma-base-change-unramified},
\ref{morphisms-lemma-base-change-proper} et
\ref{morphisms-lemma-base-change-finite}.

\medskip\noindent
Dans chaque cas, le sens intéressant consiste donc à supposer
que $f$ possède la propriété et à en déduire que $f'$ la possède aussi.
Par récurrence sur l'ordre de l'épaississement, on peut
supposer que $S \subset S'$ est un épaississement du premier ordre, voir
la discussion qui suit immédiatement
la définition \ref{definition-thickening}.

\medskip\noindent
La plupart des démonstrations utiliseront une réduction au cas affine. Soient
$U' \subset S'$ un ouvert affine et $V' \subset X'$ un ouvert affine
au-dessus de $U'$. Écrivons $U' = \Spec(A')$ et notons $I \subset A'$ l'idéal
définissant le sous-schéma fermé $U' \cap S$. Écrivons $V' = \Spec(B')$.
Alors $V' \cap X = \Spec(B'/IB')$. En posant $A = A'/I$ et
$B = B'/IB'$, on obtient un diagramme commutatif
$$
\xymatrix{
0 \ar[r] &
IB' \ar[r] &
B' \ar[r] &
B \ar[r] & 0 \\
0 \ar[r] &
IA' \ar[r] \ar[u] &
A' \ar[r] \ar[u] &
A \ar[r] \ar[u] & 0
}
$$
dont les lignes sont exactes et où $I^2 = 0$.

\medskip\noindent
Traduction algébrique de (1) : si $A \to B$ est surjectif,
alors $A' \to B'$ est surjectif. Cela résulte du
lemme de Nakayama (Algèbre, lemme \ref{algebra-lemma-NAK}).

\medskip\noindent
Traduction algébrique de (2) : si $A \to B$ est un homomorphisme d'anneaux
de type fini, alors $A' \to B'$ est de type fini. Cela résulte du
lemme de Nakayama (Algèbre, lemme \ref{algebra-lemma-NAK})
appliqué à un homomorphisme $A'[x_1, \ldots, x_n] \to B'$ tel que
$A[x_1, \ldots, x_n] \to B$ soit surjectif.

\medskip\noindent
Démonstration de (3). Elle résulte de (2) et du fait que la quasi-finitude d'un morphisme
localement de type fini se vérifie sur les fibres, voir
Morphismes, lemme \ref{morphisms-lemma-quasi-finite-at-point-characterize}.

\medskip\noindent
Démonstration de (4). Elle résulte de (2) et du fait que la propriété supplémentaire d'« être de
dimension relative $d$ » se vérifie sur les fibres (par définition, voir
Morphismes, définition \ref{morphisms-definition-relative-dimension-d}).

\medskip\noindent
Traduction algébrique de (5) : si $\Omega_{B/A} = 0$, alors
$\Omega_{B'/A'} = 0$. D'après
Algèbre, lemme \ref{algebra-lemma-differentials-base-change},
on a $0 = \Omega_{B/A} = \Omega_{B'/A'}/I\Omega_{B'/A'}$.
Ainsi, $\Omega_{B'/A'} = 0$ par le
lemme de Nakayama (Algèbre, lemme \ref{algebra-lemma-NAK}).

\medskip\noindent
Traduction algébrique de (6) : si $A \to B$ est un homomorphisme non ramifié,
alors $A' \to B'$ est non ramifié. Puisque $A \to B$ est de type
fini, on voit que $A' \to B'$ est de type fini d'après (2).
Comme $A \to B$ est non ramifié, on a $\Omega_{B/A} = 0$. D'après
(5), on a $\Omega_{B'/A'} = 0$. Ainsi, $A' \to B'$ est non ramifié.

\medskip\noindent
Démonstration de (7). Elle résulte de la combinaison de (2) avec
les résultats du lemme \ref{lemma-thicken-property-morphisms}
et du fait que propre signifie quasi-compact $+$
séparé $+$ localement de type fini $+$ universellement fermé.

\medskip\noindent
Démonstration de (8). Elle résulte de la combinaison de (2) avec
les résultats du lemme \ref{lemma-thicken-property-morphisms}
et du fait que fini signifie entier $+$ localement
de type fini (Morphismes, lemme \ref{morphisms-lemma-finite-integral}).

\medskip\noindent
Démonstration de (9). Puisque $f$ est un monomorphisme, on a $X = X \times_S X$.
On peut appliquer les résultats déjà démontrés au morphisme d'épaississements
$(X \subset X') \to (X \times_S X \subset X' \times_{S'} X')$.
On en déduit que $X' \to X' \times_{S'} X'$ est une immersion fermée d'après (1).
En fait, c'est un épaississement du premier ordre, car l'idéal définissant
l'immersion fermée
$X' \to X' \times_{S'} X'$ est contenu dans l'image inverse de l'idéal
$\mathcal{I} \subset \mathcal{O}_{S'}$ définissant $S$ dans $S'$.
En effet, $X = X \times_S X = (X' \times_{S'} X') \times_{S'} S$
est contenu dans $X'$. D'après
Morphismes, lemme \ref{morphisms-lemma-differentials-diagonal},
il suffit donc de montrer que
$\Omega_{X'/S'} = 0$, ce qui résulte de (5)
et de l'énoncé correspondant pour $X/S$.

\medskip\noindent
Démonstration de (10). Si $f : X \to S$ est une immersion, il se factorise en
$X \to U \to S$, où $U \to S$ est une immersion ouverte et $X \to U$ une
immersion fermée. Soit $U' \subset S'$ le sous-schéma ouvert dont
l'espace topologique sous-jacent est celui de $U$. Alors $X' \to S'$
se factorise par $U'$ et l'on en déduit que $X' \to U'$ est une immersion
fermée d'après (1). Cela achève la démonstration.
\end{proof}

\noindent
Le lemme suivant est une variante du précédent. Au lieu de supposer
que les épaississements considérés sont d'ordre fini (ce qui permet de transmettre
la propriété d'être localement de type fini de $f$ à $f'$),
on suppose directement que $f$ et $f'$ sont tous deux localement de
type fini.

\begin{lemma}
\label{lemma-properties-that-extend-over-thickenings}
Soit $(f, f') : (X \subset X') \to (Y \to Y')$ un morphisme
d'épaississements. Supposons que $f$ et $f'$ soient localement de type fini
et que $X = Y \times_{Y'} X'$. Alors
\begin{enumerate}
\item $f$ est localement quasi-fini si et seulement si $f'$ est localement quasi-fini,
\item $f$ est fini si et seulement si $f'$ est fini,
\item $f$ est une immersion fermée si et seulement si $f'$ est une immersion fermée,
\item $\Omega_{X/Y} = 0$ si et seulement si $\Omega_{X'/Y'} = 0$,
\item $f$ est non ramifié si et seulement si $f'$ est non ramifié,
\item $f$ est un monomorphisme si et seulement si $f'$ est un monomorphisme,
\item $f$ est une immersion si et seulement si $f'$ est une immersion,
\item $f$ est propre si et seulement si $f'$ est propre, et
\item ajouter d'autres propriétés ici.
\end{enumerate}
\end{lemma}

\begin{proof}
Les propriétés $\mathcal{P}$ énumérées dans le lemme sont toutes stables
par changement de base ; si $f'$ possède la propriété $\mathcal{P}$,
il en va donc de même de $f$. Voir
Schémas, lemmes \ref{schemes-lemma-base-change-immersion} et
\ref{schemes-lemma-base-change-monomorphism},
et
Morphismes, lemmes
\ref{morphisms-lemma-base-change-quasi-finite},
\ref{morphisms-lemma-base-change-relative-dimension-d},
\ref{morphisms-lemma-base-change-differentials},
\ref{morphisms-lemma-base-change-unramified},
\ref{morphisms-lemma-base-change-proper} et
\ref{morphisms-lemma-base-change-finite}.
Dans chaque cas, il suffit donc de prouver que, si $f$ possède
la propriété voulue, $f'$ la possède aussi.

\medskip\noindent
Un morphisme est localement quasi-fini si et seulement s'il est localement
de type fini et si ses fibres schématiques sont des espaces discrets, voir
Morphismes, lemme \ref{morphisms-lemma-locally-quasi-finite-fibres}.
Puisque le passage à un épaississement ne modifie pas
l'espace topologique sous-jacent, on voit que $f'$ est localement quasi-fini si
(et seulement si) $f$ l'est. Cela prouve (1).

\medskip\noindent
Le cas (2) résulte du cas (5) du lemme \ref{lemma-thicken-property-morphisms}
et du fait que les morphismes finis sont précisément
les morphismes entiers localement de type fini
(Morphismes, lemme \ref{morphisms-lemma-finite-integral}).

\medskip\noindent
Cas (3). Il résulte immédiatement de
Morphismes, lemme \ref{morphisms-lemma-check-closed-infinitesimally}.

\medskip\noindent
Le cas (4) résulte de l'énoncé algébrique suivant : soient $A$ un anneau et
$I \subset A$ un idéal localement nilpotent. Soit $B$ une
$A$-algèbre de type fini. Si $\Omega_{(B/IB)/(A/I)} = 0$, alors $\Omega_{B/A} = 0$.
En effet, l'hypothèse signifie que $I\Omega_{B/A} = 0$, voir
Algèbre, lemme \ref{algebra-lemma-differentials-base-change}.
D'autre part, $\Omega_{B/A}$ est un $B$-module de type fini, voir
Algèbre, lemme \ref{algebra-lemma-differentials-finitely-generated}.
L'annulation de $\Omega_{B/A}$ résulte donc du lemme de
Nakayama (Algèbre, lemme \ref{algebra-lemma-NAK}) et du fait
que $IB$ est contenu dans le radical de Jacobson de $B$.

\medskip\noindent
Le cas (5) résulte immédiatement de (4) et de
Morphismes, lemme \ref{morphisms-lemma-unramified-omega-zero}.

\medskip\noindent
Démonstration de (6). Puisque $f$ est un monomorphisme, on a $X = X \times_Y X$.
On peut appliquer les résultats déjà démontrés au morphisme d'épaississements
$(X \subset X') \to (X \times_Y X \subset X' \times_{Y'} X')$.
On en déduit que $\Delta_{X'/Y'} : X' \to X' \times_{Y'} X'$
est une immersion fermée d'après (3). En fait, $\Delta_{X'/Y'}$ est une bijection sur
les ensembles sous-jacents ; $\Delta_{X'/Y'}$ est donc un épaississement. D'autre part,
$\Delta_{X'/Y'}$ est localement de présentation finie d'après
Morphismes, lemme \ref{morphisms-lemma-diagonal-morphism-finite-type}.
Autrement dit, $\Delta_{X'/Y'}(X')$ est défini par
un faisceau quasi-cohérent d'idéaux
$\mathcal{J} \subset \mathcal{O}_{X' \times_{Y'} X'}$ de type fini.
Puisque $\Omega_{X'/Y'} = 0$ d'après (5), on voit que
le faisceau conormal de $X' \to X' \times_{Y'} X'$ est nul d'après
Morphismes, lemme \ref{morphisms-lemma-differentials-diagonal}.
Autrement dit, $\mathcal{J}/\mathcal{J}^2 = 0$.
Cela entraîne que $\Delta_{X'/Y'}$ est un isomorphisme, par exemple
d'après Algèbre, lemme \ref{algebra-lemma-ideal-is-squared-union-connected}.

\medskip\noindent
Démonstration de (7). Si $f : X \to Y$ est une immersion, il se factorise en
$X \to V \to Y$, où $V \to Y$ est une immersion ouverte et $X \to V$ une
immersion fermée. Soit $V' \subset Y'$ le sous-schéma ouvert dont
l'espace topologique sous-jacent est celui de $V$. Alors $X' \to V'$
se factorise par $V'$ et l'on en déduit que $X' \to V'$ est une immersion
fermée d'après (3).

\medskip\noindent
Le cas (8) résulte du lemme \ref{lemma-thicken-property-morphisms}
et de la définition des morphismes propres comme morphismes quasi-compacts,
universellement fermés, séparés et localement de type fini.
\end{proof}








\section{Groupes de Picard des épaississements}
\label{section-picard-group-thickening}

\noindent
Quelques résultats sur les groupes de Picard des épaississements.

\begin{lemma}
\label{lemma-picard-group-first-order-thickening}
Soit $X \subset X'$ un épaississement du premier ordre
de faisceau d'idéaux $\mathcal{I}$. Il existe alors une suite
exacte canonique
$$
\xymatrix{
0 \ar[r] &
H^0(X, \mathcal{I}) \ar[r] &
H^0(X', \mathcal{O}_{X'}^*) \ar[r] &
H^0(X, \mathcal{O}^*_X) \ar `r[d] `d[l] `l[llld] `d[dll] [dll] \\
& H^1(X, \mathcal{I}) \ar[r] &
\Pic(X') \ar[r] &
\Pic(X) \ar `r[d] `d[l] `l[llld] `d[dll] [dll] \\
& H^2(X, \mathcal{I}) \ar[r] & \ldots \ar[r] & \ldots
}
$$
de groupes abéliens.
\end{lemma}

\begin{proof}
C'est la suite exacte longue de cohomologie associée à la
suite exacte courte de faisceaux de groupes abéliens
$$
0 \to \mathcal{I} \to \mathcal{O}_{X'}^* \to \mathcal{O}_X^* \to 0
$$
dans laquelle la première application envoie une section locale $f$ de $\mathcal{I}$
sur la section inversible $1 + f$ de $\mathcal{O}_{X'}$.
On utilise aussi l'identification du groupe de Picard d'un
espace annelé au premier groupe de cohomologie du faisceau
des fonctions inversibles, voir
Cohomologie, lemme \ref{cohomology-lemma-h1-invertible}.
\end{proof}

\begin{lemma}
\label{lemma-torsion-pic-thickening}
Soit $X \subset X'$ un épaississement. Soit $n$ un entier
inversible dans $\mathcal{O}_X$. Alors l'application
$\Pic(X')[n] \to \Pic(X)[n]$ est bijective.
\end{lemma}

\begin{proof}[Démonstration pour un épaississement d'ordre fini]
D'après le principe général expliqué à la suite de
la définition \ref{definition-thickening},
on est ramené au cas d'un épaississement du premier ordre.
On peut alors utiliser le lemme \ref{lemma-picard-group-first-order-thickening}
pour voir qu'il suffit de montrer que
$H^1(X, \mathcal{I})[n]$, $H^1(X, \mathcal{I})/n$ et
$H^2(X, \mathcal{I})[n]$ sont nuls.
Cela résulte du fait que la multiplication par $n$ sur $\mathcal{I}$
est un isomorphisme, puisque c'est un $\mathcal{O}_X$-module.
\end{proof}

\begin{proof}[Démonstration dans le cas général]
Soit $\mathcal{I} \subset \mathcal{O}_{X'}$ le faisceau quasi-cohérent d'idéaux
définissant $X$. On dispose alors d'une suite exacte courte de
groupes abéliens
$$
0 \to (1 + \mathcal{I})^* \to \mathcal{O}_{X'}^* \to \mathcal{O}_X^* \to 0
$$
On obtient une suite exacte longue de cohomologie comme dans l'énoncé du
lemme \ref{lemma-picard-group-first-order-thickening},
où $H^i(X, \mathcal{I})$ est remplacé par $H^i(X, (1 + \mathcal{I})^*)$.
Il suffit donc de montrer que l'élévation à la puissance $n$ est un
isomorphisme $(1 + \mathcal{I})^* \to (1 + \mathcal{I})^*$.
En prenant les sections sur les ouverts affines, cela résulte de
Algèbre, lemme \ref{algebra-lemma-lift-nth-roots}.
\end{proof}





\section{Voisinages infinitésimaux}
\label{section-first-order-infinitesimal-neighbourhood}

\noindent
Voici une construction naturelle d'épaississements d'ordre fini.
Soit $i : Z \to X$ une immersion de schémas. Choisissons un
sous-schéma ouvert $U \subset X$ tel que $i$ identifie $Z$ à un sous-schéma
fermé $Z \subset U$. Soit $\mathcal{I} \subset \mathcal{O}_U$ le
faisceau quasi-cohérent d'idéaux définissant $Z$ dans $U$. Pour $n \geq 1$,
on peut considérer le sous-schéma fermé $Z_n \subset U$
défini par le faisceau quasi-cohérent d'idéaux $\mathcal{I}^{n + 1}$.

\begin{definition}
\label{definition-first-order-infinitesimal-neighbourhood}
Soit $i : Z \to X$ une immersion de schémas.
\begin{enumerate}
\item Le {\it voisinage infinitésimal du premier ordre} de $Z$ dans $X$ est
l'épaississement du premier ordre $Z \subset Z_1$ sur $X$ décrit ci-dessus.
\item Le {\it voisinage infinitésimal d'ordre $n$} de $Z$ dans $X$ est
l'épaississement d'ordre $n$ $Z \subset Z_n$ sur $X$ décrit ci-dessus.
\end{enumerate}
\end{definition}

\noindent
Ces épaississements possèdent la propriété universelle suivante (qui dissipe
toute crainte que la construction ci-dessus ne dépende du choix de l'ouvert
$U$).

\begin{lemma}
\label{lemma-first-order-infinitesimal-neighbourhood}
Soit $i : Z \to X$ une immersion de schémas.
\begin{enumerate}
\item Le voisinage infinitésimal du premier ordre $Z'$ de $Z$ dans $X$
possède la propriété universelle suivante : pour tout diagramme commutatif
$$
\xymatrix{
Z \ar[d]_i & T \ar[l]^a \ar[d] \\
X & T' \ar[l]_b
}
$$
où $T \subset T'$ est un épaississement du premier ordre sur $X$, il existe
un unique morphisme $(a', a) : (T \subset T') \to (Z \subset Z')$
d'épaississements sur $X$.
\item Pour $n \geq 1$, le voisinage infinitésimal d'ordre $n$, $Z_n$,
de $Z$ dans $X$ possède la propriété universelle suivante : pour tout diagramme
commutatif
$$
\xymatrix{
Z \ar[d]_i & T \ar[l]^a \ar[d] \\
X & T' \ar[l]_b
}
$$
où $T \subset T'$ est un épaississement d'ordre $n$ sur $X$, il existe
un unique morphisme $(a', a) : (T \subset T') \to (Z \subset Z_n)$
d'épaississements sur $X$.
\end{enumerate}
\end{lemma}

\begin{proof}
On démontrera seulement (1).
Soit $U \subset X$ l'ouvert utilisé dans la construction de $Z'$, c'est-à-dire un
ouvert tel que $Z$ soit identifié à un sous-schéma fermé de $U$ défini par
le faisceau quasi-cohérent d'idéaux $\mathcal{I}$.
Puisque $|T| = |T'|$, on voit que $b(T') \subset U$. On peut donc
considérer $b$ comme un morphisme à valeurs dans $U$. Soit $\mathcal{J} \subset \mathcal{O}_{T'}$
l'idéal définissant $T$. Puisque $b(T) \subset Z$ d'après le diagramme ci-dessus,
on voit que $b^\sharp(b^{-1}\mathcal{I}) \subset \mathcal{J}$. Comme
$T'$ est un épaississement du premier ordre de $T$, on a $\mathcal{J}^2 = 0$,
d'où $b^\sharp(b^{-1}(\mathcal{I}^2)) = 0$. D'après
Schémas, lemme \ref{schemes-lemma-characterize-closed-subspace},
cela entraîne que $b$ se factorise par $Z'$. Notons $a' : T' \to Z'$
cette factorisation ; tout est alors clair.
\end{proof}

\begin{lemma}
\label{lemma-infinitesimal-neighbourhood-conormal}
Soit $i : Z \to X$ une immersion de schémas. Soit $Z \subset Z'$
le voisinage infinitésimal du premier ordre de $Z$ dans $X$.
Alors le diagramme
$$
\xymatrix{
Z \ar[r] \ar[d] & Z' \ar[d] \\
Z \ar[r] & X
}
$$
induit un homomorphisme de faisceaux conormaux $\mathcal{C}_{Z/X} \to \mathcal{C}_{Z/Z'}$ d'après
Morphismes, lemme \ref{morphisms-lemma-conormal-functorial}.
Cet homomorphisme est un isomorphisme.
\end{lemma}

\begin{proof}
Cela résulte immédiatement de la construction de $Z'$ ci-dessus.
\end{proof}

\section{Morphismes formellement non ramifiés}
\label{section-formally-unramified}

\noindent
Rappelons qu'un homomorphisme d'anneaux $R \to A$ est dit {\it formellement non ramifié}
(voir Algèbre, définition \ref{algebra-definition-formally-unramified})
si, pour tout diagramme commutatif de flèches pleines
$$
\xymatrix{
A \ar[r] \ar@{-->}[rd] & B/I \\
R \ar[r] \ar[u] & B \ar[u]
}
$$
où $I \subset B$ est un idéal de carré nul, il existe au plus une flèche
pointillée qui rende le diagramme commutatif. Cela conduit
à l'analogue suivant pour les morphismes de schémas.

\begin{definition}
\label{definition-formally-unramified}
Soit $f : X \to S$ un morphisme de schémas.
On dit que $f$ est {\it formellement non ramifié} si, pour tout diagramme commutatif de flèches pleines
$$
\xymatrix{
X \ar[d]_f & T \ar[d]^i \ar[l] \\
S & T' \ar[l] \ar@{-->}[lu]
}
$$
où $T \subset T'$ est un épaississement du premier ordre de schémas affines sur $S$,
il existe au plus une flèche pointillée rendant le diagramme commutatif.
\end{definition}

\noindent
On prouve d'abord quelques lemmes formels, c'est-à-dire des lemmes dont la démonstration
consiste à tracer les diagrammes correspondants.

\begin{lemma}
\label{lemma-formally-unramified-not-affine}
Si $f : X \to S$ est un morphisme formellement non ramifié, alors, pour
tout diagramme commutatif de flèches pleines
$$
\xymatrix{
X \ar[d]_f & T \ar[d]^i \ar[l] \\
S & T' \ar[l] \ar@{-->}[lu]
}
$$
où $T \subset T'$ est un épaississement du premier ordre de schémas sur $S$,
il existe au plus une flèche pointillée rendant le diagramme commutatif.
Autrement dit, dans
la définition \ref{definition-formally-unramified},
on peut omettre la condition que $T$ soit affine.
\end{lemma}

\begin{proof}
Cela vient de ce qu'un morphisme est déterminé par ses restrictions
aux ouverts affines.
\end{proof}

\begin{lemma}
\label{lemma-composition-formally-unramified}
Un composé de morphismes formellement non ramifiés est formellement non ramifié.
\end{lemma}

\begin{proof}
C'est formel.
\end{proof}

\begin{lemma}
\label{lemma-base-change-formally-unramified}
Tout changement de base d'un morphisme formellement non ramifié est formellement non ramifié.
\end{lemma}

\begin{proof}
C'est formel.
\end{proof}

\begin{lemma}
\label{lemma-formally-unramified-on-opens}
Soit $f : X \to S$ un morphisme de schémas.
Soient $U \subset X$ et $V \subset S$ des ouverts tels que
$f(U) \subset V$. Si $f$ est formellement non ramifié, il en va de même de $f|_U : U \to V$.
\end{lemma}

\begin{proof}
Considérons un diagramme de flèches pleines
$$
\xymatrix{
U \ar[d]_{f|_U} & T \ar[d]^i \ar[l]^a \\
V & T' \ar[l] \ar@{-->}[lu]
}
$$
comme dans la définition \ref{definition-formally-unramified}. Si $f$ est formellement
non ramifié, il existe au plus un
$S$-morphisme $a' : T' \to X$ tel que $a'|_T = a$.
Il existe donc, a fortiori, au plus un tel morphisme à valeurs dans $U$.
\end{proof}

\begin{lemma}
\label{lemma-affine-formally-unramified}
Soit $f : X \to S$ un morphisme de schémas.
Supposons que $X$ et $S$ soient affines.
Alors $f$ est formellement non ramifié si et seulement si
$\mathcal{O}_S(S) \to \mathcal{O}_X(X)$ est un homomorphisme d'anneaux
formellement non ramifié.
\end{lemma}

\begin{proof}
Cela résulte immédiatement des définitions
(définition \ref{definition-formally-unramified} et
Algèbre, définition \ref{algebra-definition-formally-unramified}),
par l'équivalence entre les catégories des anneaux et des schémas affines,
voir
Schémas, lemme \ref{schemes-lemma-category-affine-schemes}.
\end{proof}

\noindent
Voici une caractérisation en termes du faisceau des différentielles.

\begin{lemma}
\label{lemma-formally-unramified-differentials}
Soit $f : X \to S$ un morphisme de schémas.
Alors $f$ est formellement non ramifié si et seulement si $\Omega_{X/S} = 0$.
\end{lemma}

\begin{proof}
Rappelons quelques arguments de la démonstration de
Morphismes, lemme \ref{morphisms-lemma-differentials-affine}.
Soit $W \subset X \times_S X$ un ouvert tel que
$\Delta : X \to X \times_S X$ induise une immersion fermée dans $W$.
Soit $\mathcal{J} \subset \mathcal{O}_W$ le faisceau d'idéaux de cette
immersion fermée. Soit $X' \subset W$ le sous-schéma fermé
défini par le faisceau quasi-cohérent d'idéaux $\mathcal{J}^2$.
Considérons les deux morphismes $p_1, p_2 : X' \to X$ induits par
les deux projections $X \times_S X \to X$.
Remarquons que $p_1$ et $p_2$ coïncident après composition avec $\Delta : X \to X'$
et que $X \to X'$ est une immersion fermée définie par un idéal
de carré nul. De plus, on dispose d'une suite exacte courte
$$
0 \to \mathcal{J}/\mathcal{J}^2 \to \mathcal{O}_{X'} \to \mathcal{O}_X \to 0
$$
et $\Omega_{X/S} = \mathcal{J}/\mathcal{J}^2$. De plus,
$\mathcal{J}/\mathcal{J}^2$ est engendré par les sections
locales $p_1^\sharp(f) - p_2^\sharp(f)$, où $f$ est une section locale de
$\mathcal{O}_X$.

\medskip\noindent
Supposons que $f : X \to S$ soit formellement non ramifié.
Par hypothèse, cela signifie que $p_1 = p_2$ après restriction à tout
ouvert affine $T' \subset X'$. Donc $p_1 = p_2$. D'après ce qui précède,
on en déduit que $\Omega_{X/S} = \mathcal{J}/\mathcal{J}^2 = 0$.

\medskip\noindent
Réciproquement, supposons que $\Omega_{X/S} = 0$. Alors $X' = X$. Prenons deux
morphismes $f'_1, f'_2 : T' \to X$ pouvant servir de flèches pointillées dans
le diagramme de
la définition \ref{definition-formally-unramified}.
On obtient ainsi un morphisme $(f'_1, f'_2) : T' \to X \times_S X$.
Puisque $f'_1|_T = f'_2|_T$ et $|T| =|T'|$, on voit que l'image de $T'$
par $(f'_1, f'_2)$ est contenue dans l'ouvert $W$ choisi ci-dessus. Comme
$(f'_1, f'_2)(T) \subset \Delta(X)$ et que $T$ est défini par un idéal
de carré nul dans $T'$, on voit que $(f'_1, f'_2)$ se factorise par $X'$.
Puisque $X' = X$, on en déduit que $f_1' = f'_2$, comme voulu.
\end{proof}

\begin{lemma}
\label{lemma-formally-unramified}
Soit $f : X \to Y$ un morphisme de schémas. Les conditions suivantes sont équivalentes :
\begin{enumerate}
\item $f$ est formellement non ramifié,
\item pour tout $x \in X$, il existe des ouverts $x \in U \subset X$ et
$f(x) \in V \subset Y$ avec $f(U) \subset V$ tels que
$f|_U : U \to V$ soit formellement non ramifié,
\item pour tout couple d'ouverts affines $U \subset X$ et $V \subset Y$
avec $f(U) \subset V$, l'homomorphisme d'anneaux $\mathcal{O}_Y(V) \to \mathcal{O}_X(U)$
est formellement non ramifié, et
\item il existe un recouvrement ouvert affine $Y = \bigcup V_j$ et,
pour chaque $j$, un recouvrement ouvert affine $f^{-1}(V_j) = \bigcup U_{ji}$
tels que $\mathcal{O}_Y(V) \to \mathcal{O}_X(U)$ soit un homomorphisme d'anneaux
formellement non ramifié pour tous $j$ et $i$.
\end{enumerate}
\end{lemma}

\begin{proof}
En combinant les lemmes \ref{lemma-formally-unramified-on-opens} et
\ref{lemma-affine-formally-unramified}, on voit que (1) $\Rightarrow$ (3).
L'implication (3) $\Rightarrow$ (4) est immédiate.
On a (4) $\Rightarrow$ (2) d'après le lemme \ref{lemma-affine-formally-unramified}.
Si (2) est vérifiée, alors $\Omega_{X/S}$ s'annule au voisinage
de chaque point d'après le lemme \ref{lemma-formally-unramified-differentials}
(on utilise aussi Morphismes, lemme
\ref{morphisms-lemma-differentials-restrict-open}) ; le
lemme \ref{lemma-formally-unramified-differentials} montre alors que $f$
est formellement non ramifié.
\end{proof}

\begin{lemma}
\label{lemma-unramified-formally-unramified}
\begin{slogan}
Les morphismes non ramifiés sont exactement les morphismes formellement non ramifiés
qui sont localement de type fini.
\end{slogan}
Soit $f : X \to S$ un morphisme de schémas.
Les conditions suivantes sont équivalentes :
\begin{enumerate}
\item Le morphisme $f$ est non ramifié (resp.\ G-non ramifié), et
\item le morphisme $f$ est localement de type fini (resp.\ localement de
présentation finie) et formellement non ramifié.
\end{enumerate}
\end{lemma}

\begin{proof}
Utiliser le lemme \ref{lemma-formally-unramified-differentials} et
Morphismes, lemme \ref{morphisms-lemma-unramified-omega-zero}.
\end{proof}










\section{Épaississements universels du premier ordre}
\label{section-universal-thickening}

\noindent
Soit $h : Z \to X$ un morphisme de schémas. Un {\it épaississement universel du premier
ordre} de $Z$ sur $X$ est un épaississement du premier ordre $Z \subset Z'$
sur $X$ tel que, pour tout épaississement du premier ordre $T \subset T'$
sur $X$ et tout diagramme commutatif de flèches pleines
$$
\xymatrix{
& Z \ar[ld] & & T \ar[rd] \ar[ll]^a \\
Z' \ar[rrd] & & & & T' \ar@{..>}[llll]_{a'} \ar[lld]^b \\
 & & X
}
$$
il existe une unique flèche pointillée rendant le diagramme commutatif.
Remarquons que, dans cette situation, $(a, a') : (T \subset T') \to (Z \subset Z')$
est un morphisme d'épaississements sur $X$. Ainsi, s'il existe un épaississement universel
du premier ordre, il est unique à isomorphisme unique près.
En général, un épaississement universel du premier ordre
n'existe pas, mais il en existe un si $h$ est formellement non ramifié.

\begin{lemma}
\label{lemma-universal-thickening}
Soit $h : Z \to X$ un morphisme de schémas formellement non ramifié.
Il existe un épaississement universel du premier ordre $Z \subset Z'$ de
$Z$ sur $X$.
\end{lemma}

\begin{proof}
Au cours de cette démonstration, on dira que $Z \subset Z'$ est un épaississement universel du premier
ordre de $Z$ sur $X$ s'il satisfait à la condition du lemme.
On construira l'épaississement universel du premier ordre $Z \subset Z'$ sur $X$
par recollement, en partant du cas affine, donné par
Algèbre, lemme \ref{algebra-lemma-universal-thickening}.
Commençons par quelques remarques générales.

\medskip\noindent
S'il existe un épaississement universel du premier ordre de $Z$ sur $X$, il est unique
à isomorphisme unique près. De plus, supposons que $V \subset Z$ et
$U \subset X$ soient des sous-schémas ouverts tels que $h(V) \subset U$. Soit
$Z \subset Z'$ un épaississement universel du premier ordre de $Z$ sur $X$.
Soit $V' \subset Z'$ le sous-schéma ouvert tel que $V = Z \cap V'$.
On affirme alors que $V \subset V'$ est l'épaississement universel du premier ordre de
$V$ sur $U$. En effet, supposons donné un diagramme
$$
\xymatrix{
V \ar[d]_h & T \ar[l]^a \ar[d] \\
U & T' \ar[l]_b
}
$$
où $T \subset T'$ est un épaississement du premier ordre sur $U$. Par la propriété
universelle de $Z'$, on obtient $(a, a') : (T \subset T') \to (Z \subset Z')$.
Mais, puisqu'on a l'égalité $|T| = |T'|$ des espaces topologiques sous-jacents,
on voit que $a'(T') \subset V'$. On peut donc considérer $(a, a')$
comme un morphisme d'épaississements $(a, a') : (T \subset T') \to (V \subset V')$
sur $U$. L'unicité est également claire. De façon tout à fait analogue, on prouve
que, si $h(Z) \subset U$ et si $Z \subset Z'$ est un épaississement universel du premier
ordre sur $U$, alors $Z \subset Z'$ est un épaississement universel du premier ordre
sur $X$.

\medskip\noindent
Avant de recoller les morceaux affines, montrons que le lemme est vrai si
$Z$ et $X$ sont affines. Écrivons $X = \Spec(R)$ et $Z = \Spec(S)$. D'après
Algèbre, lemme \ref{algebra-lemma-universal-thickening},
il existe un épaississement du premier ordre $Z \subset Z'$ sur $X$
qui possède la propriété universelle du lemme pour les diagrammes
$$
\xymatrix{
Z \ar[d]_h & T \ar[l]^a \ar[d] \\
X & T' \ar[l]_b
}
$$
où $T, T'$ sont affines. Pour un diagramme quelconque, on peut choisir un recouvrement
ouvert affine $T' = \bigcup T'_i$ et l'on obtient des morphismes
$a'_i : T'_i \to Z'$ sur $X$ tels que $a'_i|_{T_i} = a|_{T_i}$.
Par unicité, on voit que $a'_i$ et $a'_j$ coïncident sur tout ouvert affine
de $T'_i \cap T'_j$. Les morphismes $a'_i$ se recollent donc en un morphisme global
$a' : T' \to Z'$ sur $X$, comme voulu. Le lemme est ainsi prouvé si $X$ et $Z$
sont affines.

\medskip\noindent
Choisissons un recouvrement ouvert affine $Z = \bigcup Z_i$ tel que chaque $Z_i$
s'envoie dans un ouvert affine $U_i$ de $X$. D'après
le lemme \ref{lemma-formally-unramified-on-opens},
les morphismes $Z_i \to U_i$ sont formellement non ramifiés.
Le cas affine fournit donc des épaississements universels du premier ordre
$Z_i \subset Z_i'$ sur $U_i$. D'après les remarques générales ci-dessus,
$Z_i \subset Z_i'$ est aussi un épaississement universel du premier ordre de
$Z_i$ sur $X$. Soit $Z'_{i, j} \subset Z'_i$ le sous-schéma ouvert
tel que $Z_i \cap Z_j = Z'_{i, j} \cap Z_i$. D'après les remarques générales,
on voit que $Z'_{i, j}$ et $Z'_{j, i}$ sont tous deux des épaississements universels du premier
ordre de $Z_i \cap Z_j$ sur $X$. Ainsi, d'après
la première de ces remarques, il existe un isomorphisme canonique
$\varphi_{ij} : Z'_{i, j} \to Z'_{j, i}$ induisant l'identité sur
$Z_i \cap Z_j$. On affirme que ces morphismes vérifient la condition de cocycle de
Schémas, section \ref{schemes-section-glueing-schemes}.
(Vérification omise. Indication : utiliser le fait que $Z'_{i, j} \cap Z'_{i, k}$ est
l'épaississement universel du premier ordre de $Z_i \cap Z_j \cap Z_k$, ce qui le détermine
à isomorphisme unique près d'après ce qui précède.)
On peut donc utiliser les résultats de
Schémas, section \ref{schemes-section-glueing-schemes},
pour obtenir un épaississement du premier ordre $Z \subset Z'$ sur $X$ tel que
le sous-schéma ouvert $Z'_i \subset Z'$ vérifiant $Z_i = Z'_i \cap Z$
soit un épaississement universel du premier ordre de $Z_i$ sur $X$.

\medskip\noindent
Cela entraîne formellement que $Z'$ est un épaississement universel du premier
ordre de $Z$ sur $X$. En effet, on dispose de la propriété universelle pour tout
diagramme
$$
\xymatrix{
Z \ar[d]_h & T \ar[l]^a \ar[d] \\
X & T' \ar[l]_b
}
$$
où $a(T)$ est contenu dans l'un des $Z_i$. Pour un diagramme quelconque, on peut
choisir un recouvrement ouvert $T' = \bigcup T'_i$ tel que $a(T_i) \subset Z_i$.
On obtient des morphismes $a'_i : T'_i \to Z'$ sur $X$ tels que
$a'_i|_{T_i} = a|_{T_i}$. On voit que $a'_i$ et $a'_j$ coïncident nécessairement
sur $T'_i \cap T'_j$, car $a'_i|_{T'_i \cap T'_j}$ et
$a'_j|_{T'_i \cap T'_j}$ résolvent tous deux le problème de morphisme à valeurs dans
l'épaississement universel du premier ordre $Z'_i \cap Z'_j$ de $Z_i \cap Z_j$ sur $X$.
Les morphismes $a'_i$ se recollent donc en un morphisme global
$a' : T' \to Z'$ sur $X$, comme voulu. Cela achève la démonstration.
\end{proof}

\begin{definition}
\label{definition-universal-thickening}
Soit $h : Z \to X$ un morphisme de schémas formellement non ramifié.
\begin{enumerate}
\item L'{\it épaississement universel du premier ordre} de $Z$ sur $X$
est l'épaississement $Z \subset Z'$ construit dans
le lemme \ref{lemma-universal-thickening}.
\item Le {\it faisceau conormal de $Z$ sur $X$} est le faisceau conormal
de $Z$ dans son épaississement universel du premier ordre $Z'$ sur $X$.
\end{enumerate}
Dans cette situation, on note souvent le faisceau conormal $\mathcal{C}_{Z/X}$.
\end{definition}

\noindent
On voit ainsi qu'il existe une suite exacte courte de faisceaux
$$
0 \to \mathcal{C}_{Z/X} \to \mathcal{O}_{Z'} \to \mathcal{O}_Z \to 0
$$
sur $Z$.
Le lemme suivant prouve que cette définition est compatible
avec celle donnée dans le cas où $Z \to X$ est une immersion.

\begin{lemma}
\label{lemma-immersion-universal-thickening}
Soit $i : Z \to X$ une immersion de schémas. Alors
\begin{enumerate}
\item $i$ est formellement non ramifié,
\item l'épaississement universel du premier ordre de $Z$ sur $X$ est le voisinage
infinitésimal du premier ordre de $Z$ dans $X$ de
la définition \ref{definition-first-order-infinitesimal-neighbourhood}, et
\item le faisceau conormal de $i$ au sens de
Morphismes, définition \ref{morphisms-definition-conormal-sheaf},
coïncide avec le faisceau conormal de $i$ au sens de
la définition \ref{definition-universal-thickening}.
\end{enumerate}
\end{lemma}

\begin{proof}
D'après
Morphismes, lemmes \ref{morphisms-lemma-open-immersion-unramified} et
\ref{morphisms-lemma-closed-immersion-unramified},
une immersion est non ramifiée, donc formellement non ramifiée d'après
le lemme \ref{lemma-unramified-formally-unramified}.
Les autres assertions résultent de la combinaison
des lemmes \ref{lemma-first-order-infinitesimal-neighbourhood} et
\ref{lemma-infinitesimal-neighbourhood-conormal}
avec les définitions.
\end{proof}

\begin{lemma}
\label{lemma-universal-thickening-unramified}
Soit $Z \to X$ un morphisme de schémas formellement non ramifié.
Alors l'épaississement universel du premier ordre $Z'$ est formellement
non ramifié sur $X$.
\end{lemma}

\begin{proof}
Il y a deux démonstrations. La première consiste à montrer que $\Omega_{Z'/X} = 0$
en travaillant localement sur des ouverts affines et en appliquant
Algèbre, lemme \ref{algebra-lemma-differentials-universal-thickening}.
Alors
le lemme \ref{lemma-formally-unramified-differentials}
donne le résultat voulu.
La seconde est l'argument direct suivant.

\medskip\noindent
Soit $T \subset T'$ un épaississement du premier ordre. Soit
$$
\xymatrix{
Z' \ar[d] & T \ar[l]^c \ar[d] \\
X & T' \ar[l] \ar[lu]^{a, b}
}
$$
un diagramme commutatif. Considérons deux morphismes $a, b : T' \to Z'$
s'inscrivant dans ce diagramme. Posons $T_0 = c^{-1}(Z) \subset T$ et
$T'_a = a^{-1}(Z)$ (au sens schématique).
Puisque $Z'$ est un épaississement du premier ordre de $Z$, on voit que $T'$
est un épaississement du premier ordre de $T'_a$. De plus, puisque $c = a|_T$, on voit que
$T_0 = T \cap T'_a$ (au sens schématique). Comme $T'$ est un épaississement du premier
ordre de $T$, il en résulte que $T'_a$
est un épaississement du premier ordre de $T_0$. Or $a|_{T'_a}$ et $b|_{T'_a}$
sont des morphismes de $T'_a$ dans $Z'$ sur $X$ qui coïncident sur $T_0$ en tant que
morphismes à valeurs dans $Z$. Par la propriété universelle de $Z'$, on en déduit donc que
$a|_{T'_a} = b|_{T'_a}$. Ainsi, $a$ et $b$ sont des morphismes définis sur
l'épaississement du premier ordre $T'$ de $T'_a$ dont les restrictions à
$T'_a$ coïncident comme morphismes à valeurs dans $Z$. En utilisant encore la propriété universelle de
$Z'$, on en déduit donc que $a = b$. Autrement dit, la propriété définissant
un morphisme formellement non ramifié est vérifiée par $Z' \to X$, comme voulu.
\end{proof}

\begin{lemma}
\label{lemma-universal-thickening-functorial}
Considérons un diagramme commutatif de schémas
$$
\xymatrix{
Z \ar[r]_h \ar[d]_f & X \ar[d]^g \\
W \ar[r]^{h'} & Y
}
$$
où $h$ et $h'$ sont formellement non ramifiés. Soit $Z \subset Z'$ l'épaississement universel
du premier ordre de $Z$ sur $X$. Soit $W \subset W'$ l'épaississement universel
du premier ordre de $W$ sur $Y$. Il existe un morphisme canonique
$(f, f') : (Z, Z') \to (W, W')$ d'épaississements sur $Y$ qui s'inscrit dans
le diagramme commutatif suivant
$$
\xymatrix{
& & & Z' \ar[ld] \ar[d]^{f'} \\
Z \ar[rr] \ar[d]_f \ar[rrru] & & X \ar[d] & W' \ar[ld] \\
W \ar[rrru]|(.667)\hole \ar[rr] & & Y
}
$$
En particulier, le morphisme $(f, f')$ d'épaississements induit un homomorphisme
de faisceaux conormaux $f^*\mathcal{C}_{W/Y} \to \mathcal{C}_{Z/X}$.
\end{lemma}

\begin{proof}
La première assertion résulte immédiatement de la propriété universelle de $W'$.
L'homomorphisme induit sur les faisceaux conormaux est celui de
Morphismes, lemme \ref{morphisms-lemma-conormal-functorial},
appliqué à $(Z \subset Z') \to (W \subset W')$.
\end{proof}

\begin{lemma}
\label{lemma-universal-thickening-fibre-product}
Soit
$$
\xymatrix{
Z \ar[r]_h \ar[d]_f & X \ar[d]^g \\
W \ar[r]^{h'} & Y
}
$$
un diagramme cartésien dans la catégorie des schémas, où
$h'$ est formellement non ramifié. Alors $h$ est formellement non ramifié et, si
$W \subset W'$ est l'épaississement universel du premier ordre de $W$ sur $Y$,
alors $Z = X \times_Y W \subset X \times_Y W'$ est l'épaississement universel
du premier ordre de $Z$ sur $X$. En particulier, l'homomorphisme canonique
$f^*\mathcal{C}_{W/Y} \to \mathcal{C}_{Z/X}$ du
lemme \ref{lemma-universal-thickening-functorial}
est surjectif.
\end{lemma}

\begin{proof}
Le morphisme $h$ est formellement non ramifié d'après
le lemme \ref{lemma-base-change-formally-unramified}.
Il est clair que $X \times_Y W'$ est un épaississement du premier ordre.
On vérifie directement qu'il possède la propriété universelle,
car $W'$ la possède (par la propriété universelle des
produits fibrés). Voir
Morphismes, lemme \ref{morphisms-lemma-conormal-functorial-flat},
pour en déduire la surjectivité de l'homomorphisme de faisceaux conormaux.
\end{proof}

\begin{lemma}
\label{lemma-universal-thickening-fibre-product-flat}
Soit
$$
\xymatrix{
Z \ar[r]_h \ar[d]_f & X \ar[d]^g \\
W \ar[r]^{h'} & Y
}
$$
un diagramme cartésien dans la catégorie des schémas, où
$h'$ est formellement non ramifié et $g$ plat. Dans ce cas, le morphisme
correspondant $Z' \to W'$ des épaississements universels du premier ordre est plat, et
$f^*\mathcal{C}_{W/Y} \to \mathcal{C}_{Z/X}$ est un isomorphisme.
\end{lemma}

\begin{proof}
La platitude est préservée par changement de base, voir
Morphismes, lemme \ref{morphisms-lemma-base-change-flat}.
La première assertion résulte donc de la description de
$W'$ dans le lemme \ref{lemma-universal-thickening-fibre-product}.
Il est clair que $X \times_Y W'$ est un épaississement du premier ordre.
On vérifie directement qu'il possède la propriété universelle,
car $W'$ la possède (par la propriété universelle des
produits fibrés). Voir
Morphismes, lemme \ref{morphisms-lemma-conormal-functorial-flat},
pour en déduire que l'homomorphisme de faisceaux conormaux est un isomorphisme.
\end{proof}

\begin{lemma}
\label{lemma-universal-thickening-localize}
La formation des épaississements universels du premier ordre commute à la restriction aux ouverts.
Plus précisément, soit $h : Z \to X$ un morphisme de schémas formellement non ramifié.
Soient $V \subset Z$, $U \subset X$ des ouverts tels que $h(V) \subset U$.
Soit $Z'$ l'épaississement universel du premier ordre de $Z$ sur $X$.
Alors $h|_V : V \to U$ est formellement non ramifié et l'épaississement universel du premier
ordre de $V$ sur $U$ est le sous-schéma ouvert $V' \subset Z'$
tel que $V = Z \cap V'$. En particulier,
$\mathcal{C}_{Z/X}|_V = \mathcal{C}_{V/U}$.
\end{lemma}

\begin{proof}
La première assertion est
le lemme \ref{lemma-formally-unramified-on-opens}.
La compatibilité des épaississements universels peut se déduire de la démonstration du
lemme \ref{lemma-universal-thickening},
ou d'Algèbre,
lemme \ref{algebra-lemma-universal-thickening-localize},
ou encore du
lemme \ref{lemma-universal-thickening-fibre-product-flat}.
\end{proof}

\begin{lemma}
\label{lemma-differentials-universally-unramified}
Soit $h : Z \to X$ un morphisme formellement non ramifié de schémas sur $S$.
Soit $Z \subset Z'$ l'épaississement universel du premier ordre de $Z$
sur $X$, de morphisme structural $h' : Z' \to X$. L'homomorphisme canonique
$$
c_{h'} : (h')^*\Omega_{X/S} \longrightarrow \Omega_{Z'/S}
$$
induit un isomorphisme
$h^*\Omega_{X/S} \to \Omega_{Z'/S} \otimes \mathcal{O}_Z$.
\end{lemma}

\begin{proof}
L'homomorphisme $c_{h'}$ est celui défini dans
Morphismes, lemme \ref{morphisms-lemma-functoriality-differentials}.
Si $i : Z \to Z'$ est l'immersion fermée donnée, alors
$i^*c_{h'}$ est un homomorphisme
$h^*\Omega_{X/S} \to \Omega_{Z'/S} \otimes \mathcal{O}_Z$.
Pour vérifier que c'est un isomorphisme, on se ramène au cas affine
par localisation, voir
le lemme \ref{lemma-universal-thickening-localize}
et
Morphismes, lemme \ref{morphisms-lemma-differentials-restrict-open}.
Dans ce cas, le résultat est
Algèbre, lemme \ref{algebra-lemma-differentials-universal-thickening}.
\end{proof}

\begin{lemma}
\label{lemma-universally-unramified-differentials-sequence}
Soit $h : Z \to X$ un morphisme formellement non ramifié de schémas sur $S$.
Il existe une suite exacte canonique
$$
\mathcal{C}_{Z/X} \to h^*\Omega_{X/S} \to \Omega_{Z/S} \to 0.
$$
La première flèche est induite par $\text{d}_{Z'/S}$, où
$Z'$ est le voisinage universel du premier ordre de $Z$ sur $X$.
\end{lemma}

\begin{proof}
On sait qu'il existe une suite exacte canonique
$$
\mathcal{C}_{Z/Z'} \to
\Omega_{Z'/S} \otimes \mathcal{O}_Z \to
\Omega_{Z/S} \to 0.
$$
voir
Morphismes, lemme \ref{morphisms-lemma-differentials-relative-immersion}.
Le résultat s'en déduit donc en appliquant
le lemme \ref{lemma-differentials-universally-unramified}.
\end{proof}

\begin{lemma}
\label{lemma-two-unramified-morphisms}
Soit
$$
\xymatrix{
Z \ar[r]_i \ar[rd]_j & X \ar[d] \\
& Y
}
$$
un diagramme commutatif de schémas où $i$ et $j$ sont formellement
non ramifiés. Il existe alors une suite exacte canonique
$$
\mathcal{C}_{Z/Y} \to
\mathcal{C}_{Z/X} \to
i^*\Omega_{X/Y} \to 0
$$
dont la première flèche provient du
lemme \ref{lemma-universal-thickening-functorial}
et la seconde du
lemme \ref{lemma-universally-unramified-differentials-sequence}.
\end{lemma}

\begin{proof}
Notons $Z \to Z'$ l'épaississement universel du premier ordre de $Z$ sur $X$.
Notons $Z \to Z''$ l'épaississement universel du premier ordre de $Z$ sur $Y$.
D'après
le lemme \ref{lemma-universally-unramified-differentials-sequence},
il existe un morphisme canonique $Z' \to Z''$ donnant un diagramme
commutatif
$$
\xymatrix{
Z \ar[r]_{i'} \ar[rd]_{j'} & Z' \ar[r] \ar[d] & X \ar[d] \\
& Z'' \ar[r] & Y
}
$$
Appliquons
Morphismes, lemme \ref{morphisms-lemma-two-immersions},
au triangle de gauche pour obtenir une suite exacte
$$
\mathcal{C}_{Z/Z''} \to
\mathcal{C}_{Z/Z'} \to
(i')^*\Omega_{Z'/Z''} \to 0
$$
Comme $Z''$ est formellement non ramifié sur $Y$ (voir
le lemme \ref{lemma-universal-thickening-unramified}),
on a
$\Omega_{Z'/Z''} = \Omega_{Z/Y}$ (en combinant
le lemme \ref{lemma-formally-unramified-differentials}
et
Morphismes, lemme \ref{morphisms-lemma-triangle-differentials}).
On a alors $(i')^*\Omega_{Z'/Y} = i^*\Omega_{X/Y}$ d'après
le lemme \ref{lemma-differentials-universally-unramified}.
\end{proof}

\begin{lemma}
\label{lemma-transitivity-conormal}
Soient $Z \to Y \to X$ des morphismes de schémas formellement non ramifiés.
\begin{enumerate}
\item Si $Z \subset Z'$ est l'épaississement universel du premier ordre de $Z$
sur $X$ et si $Y \subset Y'$ est l'épaississement universel du premier ordre de $Y$
sur $X$, alors il existe un morphisme $Z' \to Y'$ et $Y \times_{Y'} Z'$ est
l'épaississement universel du premier ordre de $Z$ sur $Y$.
\item Il existe une suite exacte canonique
$$
i^*\mathcal{C}_{Y/X} \to
\mathcal{C}_{Z/X} \to
\mathcal{C}_{Z/Y} \to 0
$$
dont les homomorphismes proviennent du
lemme \ref{lemma-universal-thickening-functorial}
et où $i : Z \to Y$ est le premier morphisme.
\end{enumerate}
\end{lemma}

\begin{proof}
Le morphisme $h : Z' \to Y'$ de (1) provient du
lemme \ref{lemma-universal-thickening-functorial}.
Le fait que $Y \times_{Y'} Z'$ soit l'épaississement universel du premier
ordre de $Z$ sur $Y$ résulte immédiatement des propriétés universelles
de $Z'$ et de $Y'$. D'après
Morphismes, lemme \ref{morphisms-lemma-transitivity-conormal},
on dispose d'une suite exacte
$$
(i')^*\mathcal{C}_{Y \times_{Y'} Z'/Z'} \to
\mathcal{C}_{Z/Z'} \to
\mathcal{C}_{Z/Y \times_{Y'} Z'} \to 0
$$
où $i' : Z \to Y \times_{Y'} Z'$ est le morphisme donné. D'après
Morphismes, lemme \ref{morphisms-lemma-conormal-functorial-flat},
il existe une surjection
$h^*\mathcal{C}_{Y/Y'} \to \mathcal{C}_{Y \times_{Y'} Z'/Z'}$.
En combinant cela avec les égalités
$\mathcal{C}_{Y/Y'} = \mathcal{C}_{Y/X}$,
$\mathcal{C}_{Z/Z'} = \mathcal{C}_{Z/X}$ et
$\mathcal{C}_{Z/Y \times_{Y'} Z'} = \mathcal{C}_{Z/Y}$,
on obtient le lemme.
\end{proof}












\section{Morphismes formellement étales}
\label{section-formally-etale}

\noindent
Rappelons qu'un homomorphisme d'anneaux $R \to A$ est dit {\it formellement étale}
(voir Algèbre, définition \ref{algebra-definition-formally-etale})
si, pour tout diagramme commutatif de flèches pleines
$$
\xymatrix{
A \ar[r] \ar@{-->}[rd] & B/I \\
R \ar[r] \ar[u] & B \ar[u]
}
$$
où $I \subset B$ est un idéal de carré nul, il existe
exactement une flèche pointillée qui rende le diagramme commutatif. Cela conduit
à l'analogue suivant pour les morphismes de schémas.

\begin{definition}
\label{definition-formally-etale}
Soit $f : X \to S$ un morphisme de schémas.
On dit que $f$ est {\it formellement étale} si, pour tout diagramme commutatif de flèches pleines
$$
\xymatrix{
X \ar[d]_f & T \ar[d]^i \ar[l] \\
S & T' \ar[l] \ar@{-->}[lu]
}
$$
où $T \subset T'$ est un épaississement du premier ordre de schémas affines sur $S$,
il existe exactement une flèche pointillée rendant le diagramme commutatif.
\end{definition}

\noindent
Il est clair qu'un morphisme formellement étale est formellement non ramifié.
Ainsi, si $f : X \to S$ est formellement étale, alors $\Omega_{X/S}$ est nul, voir
le lemme \ref{lemma-formally-unramified-differentials}.

\begin{lemma}
\label{lemma-formally-etale-not-affine}
Si $f : X \to S$ est un morphisme formellement étale, alors, pour
tout diagramme commutatif de flèches pleines
$$
\xymatrix{
X \ar[d]_f & T \ar[d]^i \ar[l] \\
S & T' \ar[l] \ar@{-->}[lu]
}
$$
où $T \subset T'$ est un épaississement du premier ordre de schémas sur $S$,
il existe exactement une flèche pointillée rendant le diagramme commutatif.
Autrement dit, dans
la définition \ref{definition-formally-etale},
on peut omettre la condition que $T$ soit affine.
\end{lemma}

\begin{proof}
Soit $T' = \bigcup T'_i$ un recouvrement ouvert affine et posons
$T_i = T \cap T'_i$. On obtient alors des morphismes $a'_i : T'_i \to X$ s'inscrivant
dans le diagramme. Par unicité, on voit que $a'_i$ et $a'_j$ coïncident sur
tout sous-schéma ouvert affine de $T'_i \cap T'_j$. Ainsi, $a'_i$ et
$a'_j$ coïncident sur $T'_i \cap T'_j$. On voit donc que les morphismes $a'_i$
se recollent en un morphisme global $a' : T' \to X$. L'unicité de
$a'$ a été établie dans
le lemme \ref{lemma-formally-unramified-not-affine}.
\end{proof}

\begin{lemma}
\label{lemma-composition-formally-etale}
Un composé de morphismes formellement étales est formellement étale.
\end{lemma}

\begin{proof}
C'est formel.
\end{proof}

\begin{lemma}
\label{lemma-base-change-formally-etale}
Tout changement de base d'un morphisme formellement étale est formellement étale.
\end{lemma}

\begin{proof}
C'est formel.
\end{proof}

\begin{lemma}
\label{lemma-formally-etale-on-opens}
Soit $f : X \to S$ un morphisme de schémas.
Soient $U \subset X$ et $V \subset S$ des sous-schémas ouverts tels que
$f(U) \subset V$. Si $f$ est formellement étale, il en va de même de $f|_U : U \to V$.
\end{lemma}

\begin{proof}
Considérons un diagramme de flèches pleines
$$
\xymatrix{
U \ar[d]_{f|_U} & T \ar[d]^i \ar[l]^a \\
V & T' \ar[l] \ar@{-->}[lu]
}
$$
comme dans la définition \ref{definition-formally-etale}. Si $f$ est formellement
étale, il existe exactement un $S$-morphisme $a' : T' \to X$
tel que $a'|_T = a$. Puisque $|T'| = |T|$, on en déduit que $a'(T') \subset U$,
ce qui donne l'unique morphisme de $T'$ dans $U$ voulu.
\end{proof}

\begin{lemma}
\label{lemma-characterize-formally-etale}
Soit $f : X \to S$ un morphisme de schémas.
Les conditions suivantes sont équivalentes :
\begin{enumerate}
\item $f$ est formellement étale,
\item $f$ est formellement non ramifié et l'épaississement universel du premier ordre
de $X$ sur $S$ est égal à $X$,
\item $f$ est formellement non ramifié et $\mathcal{C}_{X/S} = 0$, et
\item $\Omega_{X/S} = 0$ et $\mathcal{C}_{X/S} = 0$.
\end{enumerate}
\end{lemma}

\begin{proof}
En fait, la dernière assertion n'a de sens que parce que $\Omega_{X/S} = 0$
entraîne que $\mathcal{C}_{X/S}$ est défini en vertu du
lemme \ref{lemma-formally-unramified-differentials}
et de
la définition \ref{definition-universal-thickening}.
Cela montre aussi clairement que (3) et (4) sont équivalentes.

\medskip\noindent
Chacune des hypothèses (1), (2) et (3) entraîne que $f$ est formellement
non ramifié. On peut donc supposer $f$ formellement non ramifié. L'équivalence
de (1), (2) et (3) résulte de la propriété universelle de l'épaississement universel
du premier ordre $X'$ de $X$ sur $S$ et du fait que
$X = X' \Leftrightarrow \mathcal{C}_{X/S} = 0$, puisque
$\mathcal{C}_{X/S} = \mathcal{C}_{X/X'}$ est, par définition,
le faisceau d'idéaux de $X$ dans $X'$.
\end{proof}

\begin{lemma}
\label{lemma-unramified-flat-formally-etale}
Un morphisme non ramifié et plat est formellement étale.
\end{lemma}

\begin{proof}
Supposons $X \to S$ non ramifié et plat. Alors $\Delta : X \to X \times_S X$
est une immersion ouverte, voir
Morphismes, lemme \ref{morphisms-lemma-diagonal-unramified-morphism}.
Il faut montrer que $\mathcal{C}_{X/S}$ est nul.
Considérons les deux projections $p, q : X \times_S X \to X$.
Comme $f$ est formellement non ramifié (voir
le lemme \ref{lemma-unramified-formally-unramified}),
$q$ est formellement non ramifié (voir
le lemme \ref{lemma-base-change-formally-unramified}).
Comme $f$ est plat, $p$ est plat, voir
Morphismes, lemme \ref{morphisms-lemma-base-change-flat}.
Ainsi, $p^*\mathcal{C}_{X/S} = \mathcal{C}_q$ d'après
le lemme \ref{lemma-universal-thickening-fibre-product-flat},
où $\mathcal{C}_q$ désigne le faisceau conormal du morphisme formellement
non ramifié $q : X \times_S X \to X$.
Mais $\Delta(X) \subset X \times_S X$ est un sous-schéma ouvert
qui s'envoie isomorphiquement sur $X$ par $q$. D'après
le lemme \ref{lemma-universal-thickening-localize},
on voit donc que $\mathcal{C}_q|_{\Delta(X)} = \mathcal{C}_{X/X} = 0$.
Autrement dit, l'image inverse de $\mathcal{C}_{X/S}$ sur $X$ par
le morphisme identité est nulle, c'est-à-dire $\mathcal{C}_{X/S} = 0$.
\end{proof}

\begin{lemma}
\label{lemma-affine-formally-etale}
Soit $f : X \to S$ un morphisme de schémas.
Supposons que $X$ et $S$ soient affines.
Alors $f$ est formellement étale si et seulement si
$\mathcal{O}_S(S) \to \mathcal{O}_X(X)$ est un homomorphisme d'anneaux
formellement étale.
\end{lemma}

\begin{proof}
Cela résulte immédiatement des définitions
(définition \ref{definition-formally-etale} et
Algèbre, définition \ref{algebra-definition-formally-etale}),
par l'équivalence entre les catégories des anneaux et des schémas affines,
voir
Schémas, lemme \ref{schemes-lemma-category-affine-schemes}.
\end{proof}

\begin{lemma}
\label{lemma-formally-etale}
Soit $f : X \to Y$ un morphisme de schémas. Les conditions suivantes sont équivalentes :
\begin{enumerate}
\item $f$ est formellement étale,
\item pour tout $x \in X$, il existe des ouverts $x \in U \subset X$ et
$f(x) \in V \subset Y$ avec $f(U) \subset V$ tels que
$f|_U : U \to V$ soit formellement étale,
\item pour tout couple d'ouverts affines $U \subset X$ et $V \subset Y$
avec $f(U) \subset V$, l'homomorphisme d'anneaux $\mathcal{O}_Y(V) \to \mathcal{O}_X(U)$
est formellement étale, et
\item il existe un recouvrement ouvert affine $Y = \bigcup V_j$ et,
pour chaque $j$, un recouvrement ouvert affine $f^{-1}(V_j) = \bigcup U_{ji}$
tels que $\mathcal{O}_Y(V) \to \mathcal{O}_X(U)$ soit un homomorphisme d'anneaux
formellement étale pour tous $j$ et $i$.
\end{enumerate}
\end{lemma}

\begin{proof}
En combinant les lemmes \ref{lemma-formally-etale-on-opens} et
\ref{lemma-affine-formally-etale}, on voit que (1) $\Rightarrow$ (3).
L'implication (3) $\Rightarrow$ (4) est immédiate.
On a (4) $\Rightarrow$ (2) d'après le lemme \ref{lemma-affine-formally-etale}.
Si (2) est vérifiée, alors $\Omega_{X/S}$ et $\mathcal{C}_{X/S}$ s'annulent
au voisinage de chaque point d'après
le lemme \ref{lemma-characterize-formally-etale}
(on utilise aussi Morphismes, lemme
\ref{morphisms-lemma-differentials-restrict-open}, et
le lemme \ref{lemma-universal-thickening-localize}) ; le
lemme \ref{lemma-characterize-formally-etale} montre alors que $f$
est formellement étale.
\end{proof}

\begin{lemma}
\label{lemma-etale-formally-etale}
Soit $f : X \to S$ un morphisme de schémas.
Les conditions suivantes sont équivalentes :
\begin{enumerate}
\item Le morphisme $f$ est étale, et
\item le morphisme $f$ est localement de présentation finie et
formellement étale.
\end{enumerate}
\end{lemma}

\begin{proof}
Supposons $f$ étale.
Un morphisme étale est localement de présentation finie, plat et non ramifié, voir
Morphismes, section \ref{morphisms-section-etale}.
Ainsi, $f$ est localement de présentation finie et formellement étale, voir
le lemme \ref{lemma-unramified-flat-formally-etale}.

\medskip\noindent
Réciproquement, supposons que $f$ soit localement de présentation finie et
formellement étale. La propriété d'être étale est locale pour la topologie de Zariski sur
$X$ et $S$, voir
Morphismes, lemme \ref{morphisms-lemma-etale-characterize}.
D'après
le lemme \ref{lemma-formally-etale-on-opens},
on peut recouvrir $X$ par des ouverts affines $U$ s'envoyant dans des ouverts affines
$V$ tels que $U \to V$ soit formellement étale (et de présentation finie, voir
Morphismes,
lemme \ref{morphisms-lemma-locally-finite-presentation-characterize}).
D'après
le lemme \ref{lemma-affine-formally-etale},
on voit que les homomorphismes d'anneaux $\mathcal{O}(V) \to \mathcal{O}(U)$ sont
formellement étales (et de présentation finie).
On conclut par
Algèbre, lemme \ref{algebra-lemma-formally-etale-etale}.
(On donnera une autre démonstration de cette implication lors de l'étude
des morphismes formellement lisses.)
\end{proof}














\section{Déformations infinitésimales des morphismes}
\label{section-action-by-derivations}

\noindent
Dans cette section, on explique comment utiliser une dérivation pour
déformer infinitésimalement un morphisme. Dans toute cette section, on utilise le fait qu'un
faisceau sur un épaississement $X'$ de $X$ peut être considéré comme un faisceau sur $X$.

\begin{lemma}
\label{lemma-difference-derivation}
Soit $S$ un schéma.
Soient $X \subset X'$ et $Y \subset Y'$ deux épaississements du premier ordre
sur $S$. Soient $(a, a'), (b, b') : (X \subset X') \to (Y \subset Y')$
deux morphismes d'épaississements sur $S$. Supposons que
\begin{enumerate}
\item $a = b$, et
\item les deux homomorphismes $a^*\mathcal{C}_{Y/Y'} \to \mathcal{C}_{X/X'}$
(Morphismes, lemme \ref{morphisms-lemma-conormal-functorial})
soient égaux.
\end{enumerate}
Alors l'application $(a')^\sharp - (b')^\sharp$ se factorise en
$$
\mathcal{O}_{Y'} \to \mathcal{O}_Y \xrightarrow{D}
a_*\mathcal{C}_{X/X'} \to a_*\mathcal{O}_{X'}
$$
où $D$ est une $\mathcal{O}_S$-dérivation.
\end{lemma}

\begin{proof}
Au lieu de travailler sur $Y$, on travaille sur $X$. L'avantage est que le foncteur image inverse
$a^{-1}$ est exact. Grâce à (1) et (2), on obtient un diagramme commutatif
à lignes exactes
$$
\xymatrix{
0 \ar[r] &
\mathcal{C}_{X/X'} \ar[r] &
\mathcal{O}_{X'} \ar[r] &
\mathcal{O}_X \ar[r] & 0 \\
0 \ar[r] &
a^{-1}\mathcal{C}_{Y/Y'} \ar[r] \ar[u] &
a^{-1}\mathcal{O}_{Y'}
\ar[r] \ar@<1ex>[u]^{(a')^\sharp} \ar@<-1ex>[u]_{(b')^\sharp} &
a^{-1}\mathcal{O}_Y \ar[r] \ar[u] & 0
}
$$
Or, dans une telle situation, c'est un fait général que la différence des homomorphismes de
$\mathcal{O}_S$-algèbres $(a')^\sharp$ et $(b')^\sharp$ est une
$\mathcal{O}_S$-dérivation de $a^{-1}\mathcal{O}_Y$ dans $\mathcal{C}_{X/X'}$.
Par adjonction des foncteurs $a^{-1}$ et $a_*$, cela revient
à se donner une $\mathcal{O}_S$-dérivation de
$\mathcal{O}_Y$ dans $a_*\mathcal{C}_{X/X'}$. Certains détails sont omis.
\end{proof}

\noindent
Remarquons que, dans la situation du lemme précédent, on peut écrire
$D$ sous la forme
\begin{equation}
\label{equation-D}
D = \text{d}_{Y/S} \circ \theta
\end{equation}
où $\theta$ est un homomorphisme $\mathcal{O}_Y$-linéaire
$\theta : \Omega_{Y/S} \to a_*\mathcal{C}_{X/X'}$.
Bien entendu, par une nouvelle adjonction, on peut alors considérer $\theta$ comme un
homomorphisme $\mathcal{O}_X$-linéaire
$\theta : a^*\Omega_{Y/S} \to \mathcal{C}_{X/X'}$.

\begin{lemma}
\label{lemma-action-by-derivations}
Soit $S$ un schéma.
Soit $(a, a') : (X \subset X') \to (Y \subset Y')$
un morphisme d'épaississements du premier ordre sur $S$.
Soit
$$
\theta : a^*\Omega_{Y/S} \to \mathcal{C}_{X/X'}
$$
un homomorphisme $\mathcal{O}_X$-linéaire. Il existe alors un unique morphisme de couples
$(b, b') : (X \subset X') \to (Y \subset Y')$ tel que
(1) et (2) du
lemme \ref{lemma-difference-derivation}
soient vérifiées et que la dérivation $D$ et $\theta$ soient liées par
l'équation (\ref{equation-D}).
\end{lemma}

\begin{proof}
On pose simplement $b = a$ et on définit $(b')^\sharp$ comme l'application
$$
(a')^\sharp + D : a^{-1}\mathcal{O}_{Y'} \to \mathcal{O}_{X'}
$$
où $D$ est donné par l'équation (\ref{equation-D}). On omet de vérifier
que $(b')^\sharp$ est un homomorphisme de faisceaux de $\mathcal{O}_S$-algèbres et
que (1) et (2) du
lemme \ref{lemma-difference-derivation}
sont vérifiées. L'équation (\ref{equation-D}) est satisfaite par construction.
\end{proof}

\begin{remark}
\label{remark-action-by-derivations}
Les hypothèses et notations sont celles du lemme \ref{lemma-action-by-derivations}.
L'action d'une section locale $\theta$ sur $a'$ est parfois notée
$\theta \cdot a'$. Cela signifie simplement
que $(a')^\sharp$ et $(\theta \cdot a')^\sharp$ diffèrent d'une dérivation
$D$ liée à $\theta$ par l'équation (\ref{equation-D}).
\end{remark}

\begin{lemma}
\label{lemma-sheaf}
Soit $S$ un schéma.
Soient $X \subset X'$ et $Y \subset Y'$ des épaississements du premier ordre
sur $S$. Supposons donnés un morphisme $a : X \to Y$ et un homomorphisme
$A : a^*\mathcal{C}_{Y/Y'} \to \mathcal{C}_{X/X'}$ de
$\mathcal{O}_X$-modules. Pour un sous-schéma ouvert $U' \subset X'$,
considérons les morphismes $a' : U' \to Y'$ tels que
\begin{enumerate}
\item $a'$ soit un morphisme sur $S$,
\item $a'|_U = a|_U$, et
\item l'homomorphisme induit
$a^*\mathcal{C}_{Y/Y'}|_U \to \mathcal{C}_{X/X'}|_U$
soit la restriction de $A$ à $U$.
\end{enumerate}
Ici, $U = X \cap U'$. Alors la règle
\begin{equation}
\label{equation-sheaf}
U' \mapsto
\{a' : U' \to Y'\text{ vérifiant (1), (2), (3).}\}
\end{equation}
définit un faisceau d'ensembles sur $X'$.
\end{lemma}

\begin{proof}
Notons $\mathcal{F}$ la règle du lemme.
L'application de restriction $\mathcal{F}(U') \to \mathcal{F}(V')$, pour
$V' \subset U' \subset X'$,
de $\mathcal{F}$ est simplement l'application de restriction $a' \mapsto a'|_{V'}$.
Avec cette définition, il est clair que $\mathcal{F}$ est un
faisceau, puisque les morphismes se définissent localement.
\end{proof}

\noindent
Dans le lemme suivant, on identifie les faisceaux sur $X$ à ceux sur un épaississement
quelconque de $X$.

\begin{lemma}
\label{lemma-action-sheaf}
Les notations et hypothèses sont celles du lemme \ref{lemma-sheaf}.
Il existe une action du faisceau
$$
\SheafHom_{\mathcal{O}_X}(a^*\Omega_{Y/S}, \mathcal{C}_{X/X'})
$$
sur le faisceau (\ref{equation-sheaf}). De plus, cette action
est simplement transitive pour tout ouvert $U' \subset X'$ sur lequel le faisceau
(\ref{equation-sheaf}) possède une section.
\end{lemma}

\begin{proof}
Cela résulte de la combinaison des
lemmes \ref{lemma-difference-derivation},
\ref{lemma-action-by-derivations}
et \ref{lemma-sheaf}.
\end{proof}

\begin{remark}
\label{remark-special-case}
Un cas particulier des
lemmes \ref{lemma-difference-derivation},
\ref{lemma-action-by-derivations},
\ref{lemma-sheaf} et
\ref{lemma-action-sheaf}
est celui où $Y = Y'$. Dans ce cas, l'homomorphisme $A$ est toujours nul.
Le faisceau du
lemme \ref{lemma-sheaf}
est simplement donné par la règle
$$
U' \mapsto
\{a' : U' \to Y\text{ sur }S\text{ avec } a'|_U = a|_U\}
$$
et l'on fait agir sur lui le faisceau
$\SheafHom_{\mathcal{O}_X}(a^*\Omega_{Y/S}, \mathcal{C}_{X/X'})$.
\end{remark}

\begin{remark}
\label{remark-another-special-case}
Un autre cas particulier des
lemmes \ref{lemma-difference-derivation},
\ref{lemma-action-by-derivations},
\ref{lemma-sheaf} et
\ref{lemma-action-sheaf}
est celui où $S$ est lui-même un épaississement $Z \subset Z' = S$
et où $Y = Z \times_{Z'} Y'$. Voici le diagramme :
$$
\xymatrix{
(X \subset X') \ar@{..>}[rr]_{(a, ?)} \ar[rd]_{(g, g')} & &
(Y \subset Y') \ar[ld]^{(h, h')} \\
& (Z \subset Z')
}
$$
Dans ce cas, l'homomorphisme $A : a^*\mathcal{C}_{Y/Y'} \to \mathcal{C}_{X/X'}$
est déterminé par $a$ : l'homomorphisme
$h^*\mathcal{C}_{Z/Z'} \to \mathcal{C}_{Y/Y'}$ est surjectif (car
on a supposé $Y = Z \times_{Z'} Y'$),
donc son image inverse $g^*\mathcal{C}_{Z/Z'} = a^*h^*\mathcal{C}_{Z/Z'} \to
a^*\mathcal{C}_{Y/Y'}$ est surjective, et le composé
$g^*\mathcal{C}_{Z/Z'} \to a^*\mathcal{C}_{Y/Y'} \to \mathcal{C}_{X/X'}$
doit être l'homomorphisme canonique induit par $g'$. Ainsi, le faisceau du
lemme \ref{lemma-sheaf}
est simplement donné par la règle
$$
U' \mapsto
\{a' : U' \to Y'\text{ sur }Z'\text{ avec } a'|_U = a|_U\}
$$
et l'on fait agir sur lui le faisceau
$\SheafHom_{\mathcal{O}_X}(a^*\Omega_{Y/Z}, \mathcal{C}_{X/X'})$.
\end{remark}

\begin{lemma}
\label{lemma-omega-deformation}
Soit $S$ un schéma. Soit $X \subset X'$ un épaississement du premier ordre sur
$S$. Soit $Y$ un schéma sur $S$. Soient
$a', b' : X' \to Y$ deux morphismes sur $S$ tels que
$a = a'|_X = b'|_X$. On obtient ainsi un diagramme commutatif
$$
\xymatrix{
X \ar[r] \ar[d]_a & X' \ar[d]^{(b', a')} \\
Y \ar[r]^-{\Delta_{Y/S}} & Y \times_S Y
}
$$
Puisque les flèches horizontales sont des immersions de faisceaux conormaux
$\mathcal{C}_{X/X'}$ et $\Omega_{Y/S}$, d'après
Morphismes, lemme \ref{morphisms-lemma-conormal-functorial},
on obtient un homomorphisme $\theta : a^*\Omega_{Y/S} \to \mathcal{C}_{X/X'}$.
Alors cet homomorphisme $\theta$ et la dérivation $D$ du
lemme \ref{lemma-difference-derivation}
sont liés par l'équation (\ref{equation-D}).
\end{lemma}

\begin{proof}
Omise. Indication : l'égalité se vérifie sur les ouverts affines, où elle
résulte du calcul suivant. Si $f$ est une section locale de
$\mathcal{O}_Y$, alors $1 \otimes f - f \otimes 1$ est une section locale
de $\mathcal{C}_{Y/(Y \times_S Y)}$ correspondant à $\text{d}_{Y/S}(f)$.
Elle s'envoie sur la section locale $(a')^\sharp(f) - (b')^\sharp(f) = D(f)$
de $\mathcal{C}_{X/X'}$. Autrement dit, $\theta(\text{d}_{Y/S}(f)) = D(f)$.
\end{proof}

\noindent
Pour la suite, il nous faut un résultat affirmant, en substance, que la
construction du
lemme \ref{lemma-action-by-derivations}
est compatible avec la localisation étale.

\begin{lemma}
\label{lemma-sheaf-differentials-etale-localization}
Soit
$$
\xymatrix{
X_1 \ar[d] & X_2 \ar[l]^f \ar[d] \\
S_1 & S_2 \ar[l]
}
$$
un diagramme commutatif de schémas dans lequel $X_2 \to X_1$ et $S_2 \to S_1$
sont étales. Alors l'homomorphisme $c_f : f^*\Omega_{X_1/S_1} \to \Omega_{X_2/S_2}$ de
Morphismes, lemme \ref{morphisms-lemma-functoriality-differentials},
est un isomorphisme.
\end{lemma}

\begin{proof}
Rappelons qu'un morphisme étale $U \to V$ est un morphisme lisse
tel que $\Omega_{U/V} = 0$. En utilisant cela, on voit que
Morphismes, lemme \ref{morphisms-lemma-triangle-differentials},
entraîne $\Omega_{X_2/S_2} = \Omega_{X_2/S_1}$, et que
Morphismes, lemme \ref{morphisms-lemma-triangle-differentials-smooth},
entraîne que l'homomorphisme $f^*\Omega_{X_1/S_1} \to \Omega_{X_2/S_1}$
(pour le morphisme $f$ considéré comme morphisme sur $S_1$)
est un isomorphisme. Le lemme s'en déduit.
\end{proof}

\begin{lemma}
\label{lemma-action-by-derivations-etale-localization}
Considérons un diagramme commutatif d'épaississements du premier ordre
$$
\vcenter{
\xymatrix@C=1.8pc{
(T_2 \subset T_2') \ar[d]_{(h, h')} \ar[rr]_{(a_2, a_2')} & &
(X_2 \subset X_2') \ar[d]^{(f, f')} \\
(T_1 \subset T_1') \ar[rr]^{(a_1, a_1')} & &
(X_1 \subset X_1')
}
}
\quad
\begin{matrix}
\text{et un diagramme commutatif} \\
\text{de schémas}
\end{matrix}
\quad
\vcenter{
\xymatrix@C=1.8pc{
X_2' \ar[r] \ar[d] & S_2 \ar[d] \\
X_1' \ar[r] & S_1
}
}
$$
où $X_2 \to X_1$ et $S_2 \to S_1$ sont étales.
Pour tout homomorphisme $\mathcal{O}_{T_1}$-linéaire
$\theta_1 : a_1^*\Omega_{X_1/S_1} \to \mathcal{C}_{T_1/T'_1}$, notons
$\theta_2$ le composé
$$
\xymatrix{
a_2^*\Omega_{X_2/S_2} \ar@{=}[r] &
h^*a_1^*\Omega_{X_1/S_1} \ar[r]^-{h^*\theta_1} &
h^*\mathcal{C}_{T_1/T'_1} \ar[r] &
\mathcal{C}_{T_2/T'_2}
}
$$
(le signe d'égalité est expliqué dans la démonstration). Alors le diagramme
$$
\xymatrix{
T_2' \ar[rr]_{\theta_2 \cdot a_2'} \ar[d] & & X'_2 \ar[d] \\
T_1' \ar[rr]^{\theta_1 \cdot a_1'} & & X'_1
}
$$
est commutatif, où les actions $\theta_2 \cdot a_2'$ et $\theta_1 \cdot a_1'$
sont celles de la remarque \ref{remark-action-by-derivations}.
\end{lemma}

\begin{proof}
Le signe d'égalité provient de l'identification
$f^*\Omega_{X_1/S_1} = \Omega_{X_2/S_2}$ du
lemme \ref{lemma-sheaf-differentials-etale-localization}.
En effet, grâce à celle-ci, on a
$a_2^*\Omega_{X_2/S_2} = a_2^*f^*\Omega_{X_1/S_1} =
h^*a_1^*\Omega_{X_1/S_1}$, car $f \circ a_2 = a_1 \circ h$.
Cela étant, la commutativité du diagramme peut se vérifier
sur les ouverts affines. On peut donc supposer que les schémas du grand
diagramme initial sont affines. On obtient ainsi des diagrammes commutatifs
$$
\vcenter{
\xymatrix{
(B'_2, I_2) & & (A'_2, J_2) \ar[ll]^{a_2'} \\
(B'_1, I_1) \ar[u]^{h'} & & (A'_1, J_1) \ar[ll]_{a_1'} \ar[u]_{f'}
}
}
\quad\text{et}\quad
\vcenter{
\xymatrix{
A'_2 & & R_2 \ar[ll] \\
A'_1 \ar[u] & & R_1 \ar[ll] \ar[u]
}
}
$$
Cette notation signifie que $I_1, I_2, J_1, J_2$ sont des idéaux de carré
nul et que les morphismes de couples sont des homomorphismes d'anneaux envoyant les idéaux dans les idéaux.
Posons $A_1 = A'_1/J_1$, $A_2 = A'_2/J_2$, $B_1 = B'_1/I_1$ et
$B_2 = B'_2/I_2$. On sait que
$$
A_2 \otimes_{A_1} \Omega_{A_1/R_1} \longrightarrow \Omega_{A_2/R_2}
$$
est un isomorphisme. Alors
$\theta_1 : B_1 \otimes_{A_1} \Omega_{A_1/R_1} \to I_1$
est $B_1$-linéaire. Cela donne une $R_1$-dérivation
$D_1 = \theta_1 \circ \text{d}_{A_1/R_1} : A_1 \to I_1$.
De même, on voit que
$\theta_2 : B_2 \otimes_{A_2} \Omega_{A_2/R_2} \to I_2$
donne une $R_2$-dérivation
$D_2 = \theta_2 \circ \text{d}_{A_2/R_2} : A_2 \to I_2$.
La construction de $\theta_2$ entraîne la compatibilité suivante entre
$\theta_1$ et $\theta_2$ : pour tout $x \in A_1$, on a
$$
h'(D_1(x)) = D_2(f'(x))
$$
dans $I_2$. On peut considérer $D_1$ comme une application $A'_1 \to B'_1$
à l'aide de $A'_1 \to A_1 \xrightarrow{D_1} I_1 \to B_1$ ; de même,
on peut considérer $D_2$ comme une application $A'_2 \to B'_2$. L'égalité
affichée est alors valable pour $x \in A'_1$.
D'après la construction de l'action dans
le lemme \ref{lemma-action-by-derivations} et
la remarque \ref{remark-action-by-derivations},
on sait que $\theta_1 \cdot a_1'$ correspond à l'homomorphisme d'anneaux
$a_1' + D_1 : A'_1 \to B'_1$ et que $\theta_2 \cdot a_2'$ correspond
à l'homomorphisme d'anneaux $a_2' + D_2 : A'_2 \to B'_2$. L'égalité affichée
donne donc
$h' \circ (a_1' + D_1) = (a_2' + D_2) \circ f'$,
comme voulu.
\end{proof}

\begin{remark}
\label{remark-tiny-improvement}
Le lemme \ref{lemma-action-by-derivations-etale-localization}
peut être amélioré de la manière suivante.
Supposons donnés des diagrammes commutatifs comme dans
le lemme \ref{lemma-action-by-derivations-etale-localization},
sans supposer que $X_2 \to X_1$
et $S_2 \to S_1$ soient étales. Supposons ensuite que l'on dispose de
$\theta_1 : a_1^*\Omega_{X_1/S_1} \to \mathcal{C}_{T_1/T'_1}$
et de
$\theta_2 : a_2^*\Omega_{X_2/S_2} \to \mathcal{C}_{T_2/T'_2}$
tels que
$$
\xymatrix{
f_*\mathcal{O}_{X_2} \ar[rr]_{f_*D_2} & &
f_*a_{2, *}\mathcal{C}_{T_2/T_2'} \\
\mathcal{O}_{X_1} \ar[rr]^{D_1} \ar[u]^{f^\sharp} & &
a_{1, *}\mathcal{C}_{T_1/T_1'} \ar[u]_{\text{induit par }(h')^\sharp}
}
$$
soit commutatif, où $D_i$ correspond à $\theta_i$ comme dans
l'équation (\ref{equation-D}). On obtient alors la conclusion du
lemme \ref{lemma-action-by-derivations-etale-localization}.
L'intérêt de l'hypothèse que $X_2 \to X_1$ et
$S_2 \to S_1$ soient tous deux étales est qu'elle permet de construire un $\theta_2$
à partir de $\theta_1$.
\end{remark}










\section{Déformations infinitésimales des schémas}
\label{section-deform}

\noindent
Le lemme élémentaire suivant est souvent commode pour vérifier
qu'une déformation infinitésimale d'un morphisme est plate.

\begin{lemma}
\label{lemma-deform}
Soit $(f, f') : (X \subset X') \to (S \subset S')$ un morphisme
d'épaississements du premier ordre. Supposons $f$ plat.
Alors les conditions suivantes sont équivalentes :
\begin{enumerate}
\item $f'$ est plat et $X = S \times_{S'} X'$, et
\item l'homomorphisme canonique $f^*\mathcal{C}_{S/S'} \to \mathcal{C}_{X/X'}$
est un isomorphisme.
\end{enumerate}
\end{lemma}

\begin{proof}
Le problème étant local sur $X'$, on peut supposer que $X, X', S, S'$
sont des schémas affines. Écrivons $S' = \Spec(A')$, $X' = \Spec(B')$,
$S = \Spec(A)$, $X = \Spec(B)$, avec $A = A'/I$ et $B = B'/J$
pour certains idéaux de carré nul. On obtient alors le diagramme commutatif
suivant :
$$
\xymatrix{
0 \ar[r] &
J \ar[r] &
B' \ar[r] &
B \ar[r] & 0 \\
0 \ar[r] &
I \ar[r] \ar[u] &
A' \ar[r] \ar[u] &
A \ar[r] \ar[u] & 0
}
$$
dont les lignes sont exactes. L'homomorphisme canonique du lemme est l'homomorphisme
$$
I \otimes_A B = I \otimes_{A'} B' \longrightarrow J.
$$
L'hypothèse que $f$ est plat signifie que $A \to B$ est plat.

\medskip\noindent
Supposons (1). Alors $A' \to B'$ est plat et $J = IB'$. La platitude entraîne
$\text{Tor}_1^{A'}(B', A) = 0$ (voir
Algèbre, lemme \ref{algebra-lemma-characterize-flat}).
Cela signifie que $I \otimes_{A'} B' \to B'$ est injectif (voir
Algèbre, remarque \ref{algebra-remark-Tor-ring-mod-ideal}).
On voit donc que $I \otimes_A B \to J$ est un isomorphisme.

\medskip\noindent
Supposons (2). Il en résulte que $J = IB'$, donc que $X = S \times_{S'} X'$.
De plus, on obtient $\text{Tor}_1^{A'}(B', A'/I) = 0$ en renversant
les implications du paragraphe précédent. Ainsi, $B'$ est plat sur $A'$ d'après
Algèbre, lemme \ref{algebra-lemma-what-does-it-mean}.
\end{proof}

\noindent
Le lemme suivant est la version « nilpotente » du
« critère de platitude par fibres », voir
la section \ref{section-criterion-flat-fibres}.

\begin{lemma}
\label{lemma-flatness-morphism-thickenings}
Considérons un diagramme commutatif
$$
\xymatrix{
(X \subset X') \ar[rr]_{(f, f')} \ar[rd] & & (Y \subset Y') \ar[ld] \\
& (S \subset S')
}
$$
d'épaississements. Supposons que
\begin{enumerate}
\item $X'$ soit plat sur $S'$,
\item $f$ soit plat,
\item $S \subset S'$ soit un épaississement d'ordre fini, et
\item $X = S \times_{S'} X'$ et $Y = S \times_{S'} Y'$.
\end{enumerate}
Alors $f'$ est plat et $Y'$ est plat sur $S'$ en tout point de
l'image de $f'$.
\end{lemma}

\begin{proof}
C'est une conséquence immédiate du résultat d'Algèbre,
lemme \ref{algebra-lemma-criterion-flatness-fibre-nilpotent}.
\end{proof}

\noindent
De nombreuses propriétés des morphismes de schémas sont préservées par déformation
plate.

\begin{lemma}
\label{lemma-deform-property}
Considérons un diagramme commutatif
$$
\xymatrix{
(X \subset X') \ar[rr]_{(f, f')} \ar[rd] & & (Y \subset Y') \ar[ld] \\
& (S \subset S')
}
$$
d'épaississements. Supposons que $S \subset S'$ soit un épaississement d'ordre fini,
que $X'$ soit plat sur $S'$, que $X = S \times_{S'} X'$ et que
$Y = S \times_{S'} Y'$. Alors
\begin{enumerate}
\item $f$ est plat si et seulement si $f'$ est plat,
\label{item-flat}
\item $f$ est un isomorphisme si et seulement si $f'$ est un isomorphisme,
\label{item-isomorphism}
\item $f$ est une immersion ouverte si et seulement si $f'$ est une immersion ouverte,
\label{item-open-immersion}
\item $f$ est quasi-compact si et seulement si $f'$ est quasi-compact,
\label{item-quasi-compact}
\item $f$ est universellement fermé si et seulement si $f'$ est universellement fermé,
\label{item-universally-closed}
\item $f$ est (quasi-)séparé si et seulement si $f'$ est (quasi-)séparé,
\label{item-separated}
\item $f$ est un monomorphisme si et seulement si $f'$ est un monomorphisme,
\label{item-monomorphism}
\item $f$ est surjectif si et seulement si $f'$ est surjectif,
\label{item-surjective}
\item $f$ est universellement injectif si et seulement si $f'$ est universellement injectif,
\label{item-universally-injective}
\item $f$ est affine si et seulement si $f'$ est affine,
\label{item-affine}
\item
\label{item-finite-type}
$f$ est localement de type fini si et seulement si $f'$ est localement de type fini,
\item $f$ est localement quasi-fini si et seulement si $f'$ est localement quasi-fini,
\label{item-quasi-finite}
\item
\label{item-finite-presentation}
$f$ est localement de présentation finie si et seulement si $f'$ est localement de
présentation finie,
\item
\label{item-relative-dimension-d}
$f$ est localement de type fini et de dimension relative $d$ si et seulement si
$f'$ est localement de type fini et de dimension relative $d$,
\item $f$ est universellement ouvert si et seulement si $f'$ est universellement ouvert,
\label{item-universally-open}
\item $f$ est syntomique si et seulement si $f'$ est syntomique,
\label{item-syntomic}
\item $f$ est lisse si et seulement si $f'$ est lisse,
\label{item-smooth}
\item $f$ est non ramifié si et seulement si $f'$ est non ramifié,
\label{item-unramified}
\item $f$ est étale si et seulement si $f'$ est étale,
\label{item-etale}
\item $f$ est propre si et seulement si $f'$ est propre,
\label{item-proper}
\item $f$ est entier si et seulement si $f'$ est entier,
\label{item-integral}
\item $f$ est fini si et seulement si $f'$ est fini,
\label{item-finite}
\item
\label{item-finite-locally-free}
$f$ est fini localement libre (de rang $d$) si et seulement si $f'$
est fini localement libre (de rang $d$), et
\item ajouter d'autres propriétés ici.
\end{enumerate}
\end{lemma}

\begin{proof}
Les hypothèses sur $X$ et $Y$ signifient que $f$ est le changement de base de
$f'$ par $X \to X'$.
Les propriétés $\mathcal{P}$ énumérées en (1) -- (23) ci-dessus sont toutes stables
par changement de base ; si $f'$ possède la propriété $\mathcal{P}$,
il en va donc de même de $f$. Voir
Schémas, lemmes \ref{schemes-lemma-base-change-immersion},
\ref{schemes-lemma-quasi-compact-preserved-base-change},
\ref{schemes-lemma-separated-permanence} et
\ref{schemes-lemma-base-change-monomorphism},
et
Morphismes, lemmes
\ref{morphisms-lemma-base-change-surjective},
\ref{morphisms-lemma-base-change-universally-injective},
\ref{morphisms-lemma-base-change-affine},
\ref{morphisms-lemma-base-change-finite-type},
\ref{morphisms-lemma-base-change-quasi-finite},
\ref{morphisms-lemma-base-change-finite-presentation},
\ref{morphisms-lemma-base-change-relative-dimension-d},
\ref{morphisms-lemma-base-change-syntomic},
\ref{morphisms-lemma-base-change-smooth},
\ref{morphisms-lemma-base-change-unramified},
\ref{morphisms-lemma-base-change-etale},
\ref{morphisms-lemma-base-change-proper},
\ref{morphisms-lemma-base-change-finite} et
\ref{morphisms-lemma-base-change-finite-locally-free}.

\medskip\noindent
Dans chaque cas, le sens intéressant consiste donc à supposer
que $f$ possède la propriété et à en déduire que $f'$ la possède aussi.
Par récurrence sur l'ordre de l'épaississement, on peut
supposer que $S \subset S'$ est un épaississement du premier ordre, voir
la discussion qui suit immédiatement
la définition \ref{definition-thickening}.
Faisons quelques remarques générales que l'on utilisera sans autre
mention dans les arguments ci-dessous.
(I) Soient $W' \subset S'$ un ouvert affine et $U' \subset X'$
et $V' \subset Y'$ des ouverts affines au-dessus de $W'$ tels que $f'(U') \subset V'$.
Écrivons $W' = \Spec(R')$ et notons $I \subset R'$ l'idéal
définissant le sous-schéma fermé $W' \cap S$. Écrivons $U' = \Spec(B')$
et $V' = \Spec(A')$. On obtient alors un diagramme commutatif
$$
\xymatrix{
0 \ar[r] &
IB' \ar[r] &
B' \ar[r] &
B \ar[r] & 0 \\
0 \ar[r] &
IA' \ar[r] \ar[u] &
A' \ar[r] \ar[u] &
A \ar[r] \ar[u] & 0
}
$$
à lignes exactes. De plus, $IB' \cong I \otimes_R B$, voir la démonstration du
lemme \ref{lemma-deform}.
(II) Les morphismes $X \to X'$ et $Y \to Y'$ sont des homéomorphismes universels.
Les applications $f$ et $f'$ ont donc la même topologie (après tout changement de base).
(III) Si $f$ est plat, alors $f'$ est plat et
$Y' \to S'$ est plat en tout point de l'image de $f'$, voir
le lemme \ref{lemma-flatness-morphism-thickenings}.

\medskip\noindent
Pour (\ref{item-flat}). C'est la remarque générale (III).

\medskip\noindent
Pour (\ref{item-isomorphism}). Supposons que $f$ soit un isomorphisme.
D'après (III), on voit que $Y' \to S'$ est plat. Choisissons un
ouvert affine $V' \subset Y'$ et posons $U' = (f')^{-1}(V')$. Alors
$V = Y \cap V'$ est affine, ce qui entraîne que
$V \cong f^{-1}(V) = U = Y \times_{Y'} U'$ est affine. D'après
le lemme \ref{lemma-thickening-affine-scheme},
on voit que $U'$ est affine. On dispose donc d'un diagramme comme dans
la remarque générale (I), et de plus $IA \cong I \otimes_R A$, car
$R' \to A'$ est plat. Alors
$IB' \cong I \otimes_R B \cong I \otimes_R A \cong IA'$
et $A \cong B$. L'exactitude des lignes
du diagramme ci-dessus donne $A' \cong B'$,
c'est-à-dire $U' \cong V'$. Ainsi, $f'$ est un isomorphisme.

\medskip\noindent
Pour (\ref{item-open-immersion}). Supposons que $f$ soit une immersion ouverte.
Alors $f$ est un isomorphisme de $X$ sur un sous-schéma ouvert $V \subset Y$.
Soit $V' \subset Y'$ le sous-schéma ouvert dont l'espace topologique
sous-jacent est $V$. Alors $f'$ est un morphisme de $X'$ dans $V'$, qui est un isomorphisme d'après
(\ref{item-isomorphism}). Ainsi, $f'$ est une immersion ouverte.

\medskip\noindent
Pour (\ref{item-quasi-compact}). C'est immédiat d'après la remarque (II). Voir aussi
le lemme \ref{lemma-thicken-property-morphisms} pour un énoncé plus général.

\medskip\noindent
Pour (\ref{item-universally-closed}). C'est immédiat d'après la remarque (II). Voir aussi
le lemme \ref{lemma-thicken-property-morphisms} pour un énoncé plus général.

\medskip\noindent
Pour (\ref{item-separated}). Remarquons que
$X \times_Y X = Y \times_{Y'} (X' \times_{Y'} X')$, de sorte que
$X' \times_{Y'} X'$ est un épaississement de $X \times_Y X$.
Les applications $\Delta_{X/Y}$ et $\Delta_{X'/Y'}$ ont donc la même
topologie, ce qui permet de conclure. Voir aussi
le lemme \ref{lemma-thicken-property-morphisms} pour un énoncé plus général.

\medskip\noindent
Pour (\ref{item-monomorphism}). Supposons que $f$ soit un monomorphisme.
Considérons le morphisme diagonal $\Delta_{X'/Y'} : X' \to X' \times_{Y'} X'$.
Le changement de base de $\Delta_{X'/Y'}$ par $S \to S'$ est $\Delta_{X/Y}$,
qui est un isomorphisme par hypothèse. D'après (\ref{item-isomorphism}),
on en déduit que $\Delta_{X'/Y'}$ est un isomorphisme.

\medskip\noindent
Pour (\ref{item-surjective}). C'est clair. Voir aussi
le lemme \ref{lemma-thicken-property-morphisms} pour un énoncé plus général.

\medskip\noindent
Pour (\ref{item-universally-injective}). C'est immédiat d'après la remarque (II). Voir aussi
le lemme \ref{lemma-thicken-property-morphisms} pour un énoncé plus général.

\medskip\noindent
Pour (\ref{item-affine}). Supposons que $f$ soit affine. Choisissons un
ouvert affine $V' \subset Y'$ et posons $U' = (f')^{-1}(V')$.
Alors $V = Y \cap V'$ est affine, ce qui entraîne que
$U = Y \times_{Y'} U'$ est affine. D'après
le lemme \ref{lemma-thickening-affine-scheme},
on voit que $U'$ est affine. Ainsi, $f'$ est affine. Voir aussi
le lemme \ref{lemma-thicken-property-morphisms} pour un énoncé plus général.

\medskip\noindent
Pour (\ref{item-finite-type}). D'après la remarque (I), il s'agit de prouver que $A' \to B'$
est de type fini si $A \to B$ est de type fini. Supposons que
$x_1, \ldots, x_n \in B'$ soient des éléments dont les images dans $B$ engendrent $B$
comme $A$-algèbre. Alors $A'[x_1, \ldots, x_n] \to B$ est surjectif,
car $A'[x_1, \ldots, x_n] \to B$ est surjectif et
$I \otimes_R A[x_1, \ldots, x_n] \to I \otimes_R B$ est surjectif. Voir aussi
le lemme \ref{lemma-thicken-property-morphisms-cartesian}
pour un énoncé plus général.

\medskip\noindent
Pour (\ref{item-quasi-finite}). Cela résulte de (\ref{item-finite-type}) et du fait que
la quasi-finitude d'un morphisme de type fini se vérifie sur les fibres, voir
Morphismes, lemme \ref{morphisms-lemma-quasi-finite-at-point-characterize}.
Voir aussi le lemme \ref{lemma-thicken-property-morphisms-cartesian}
pour un énoncé plus général.

\medskip\noindent
Pour (\ref{item-finite-presentation}). D'après la remarque (I), il s'agit de prouver que
$A' \to B'$ est de présentation finie si $A \to B$ est de présentation finie.
On peut supposer que $B' = A'[x_1, \ldots, x_n]/K'$ pour un certain idéal $K'$ d'après
(\ref{item-finite-type}). On obtient une suite exacte courte
$$
0 \to K' \to A'[x_1, \ldots, x_n] \to B' \to 0
$$
Comme $B'$ est plat sur $R'$, on voit que $K' \otimes_{R'} R$ est le noyau de
la surjection $A[x_1, \ldots, x_n] \to B$. Par hypothèse sur $A \to B$,
il existe un nombre fini d'éléments $f'_1, \ldots, f'_m \in K'$ dont les images dans
$A[x_1, \ldots, x_n]$ engendrent ce noyau. Puisque $I$ est nilpotent, on voit
que $f'_1, \ldots, f'_m$ engendrent $K'$ par le lemme de Nakayama, voir
Algèbre, lemme \ref{algebra-lemma-NAK}.

\medskip\noindent
Pour (\ref{item-relative-dimension-d}). Cela résulte de (\ref{item-finite-type})
et de la remarque générale (II). Voir aussi
le lemme \ref{lemma-thicken-property-morphisms-cartesian}
pour un énoncé plus général.

\medskip\noindent
Pour (\ref{item-universally-open}). C'est immédiat d'après la remarque générale (II). Voir aussi
le lemme \ref{lemma-thicken-property-morphisms} pour un énoncé plus général.

\medskip\noindent
Pour (\ref{item-syntomic}). Supposons $f$ syntomique. D'après
(\ref{item-finite-presentation}), $f'$ est localement de présentation finie ;
d'après la remarque générale (III), $f'$ est plat, et les fibres de $f'$ sont celles
de $f$. Ainsi, $f'$ est syntomique d'après
Morphismes, lemme \ref{morphisms-lemma-syntomic-flat-fibres}.

\medskip\noindent
Pour (\ref{item-smooth}). Supposons $f$ lisse. D'après
(\ref{item-finite-presentation}), $f'$ est localement de présentation finie ;
d'après la remarque générale (III), $f'$ est plat, et les fibres de $f'$ sont les
fibres de $f$. Ainsi, $f'$ est lisse d'après
Morphismes, lemme \ref{morphisms-lemma-smooth-flat-smooth-fibres}.

\medskip\noindent
Pour (\ref{item-unramified}). Supposons $f$ non ramifié. D'après
(\ref{item-finite-type}), $f'$ est localement de type fini,
et les fibres de $f'$ sont celles de $f$.
Ainsi, $f'$ est non ramifié d'après
Morphismes, lemme \ref{morphisms-lemma-unramified-etale-fibres}. Voir aussi
le lemme \ref{lemma-thicken-property-morphisms-cartesian}
pour un énoncé plus général.

\medskip\noindent
Pour (\ref{item-etale}). Supposons $f$ étale. D'après
(\ref{item-finite-presentation}), $f'$ est localement de présentation finie ;
d'après la remarque générale (III), $f'$ est plat, et les fibres de $f'$ sont celles
de $f$. Ainsi, $f'$ est étale d'après
Morphismes, lemme \ref{morphisms-lemma-etale-flat-etale-fibres}.

\medskip\noindent
Pour (\ref{item-proper}). Cela résulte de la combinaison de
(\ref{item-separated}), (\ref{item-finite-type}), (\ref{item-quasi-compact})
et (\ref{item-universally-closed}). Voir aussi
le lemme \ref{lemma-thicken-property-morphisms-cartesian}
pour un énoncé plus général.

\medskip\noindent
Pour (\ref{item-integral}). Combiner (\ref{item-universally-closed}) et
(\ref{item-affine}) avec
Morphismes, lemme \ref{morphisms-lemma-integral-universally-closed}. Voir aussi
le lemme \ref{lemma-thicken-property-morphisms} pour un énoncé plus général.

\medskip\noindent
Pour (\ref{item-finite}). Combiner (\ref{item-integral})
et (\ref{item-finite-type}) avec
Morphismes, lemme \ref{morphisms-lemma-finite-integral}. Voir aussi
le lemme \ref{lemma-thicken-property-morphisms-cartesian}
pour un énoncé plus général.

\medskip\noindent
Pour (\ref{item-finite-locally-free}). Supposons $f$ fini localement libre. D'après
(\ref{item-finite}), on voit que $f'$ est fini ; d'après la remarque générale (III),
$f'$ est plat, et d'après (\ref{item-finite-presentation}), $f'$ est localement de
présentation finie. Ainsi, $f'$ est fini localement libre d'après
Morphismes, lemme \ref{morphisms-lemma-finite-flat}.
\end{proof}


\noindent
Le lemme suivant est la version « localement nilpotente » du
« critère de platitude par fibres », voir
la section \ref{section-criterion-flat-fibres}.

\begin{lemma}
\label{lemma-flatness-morphism-thickenings-fp-over-ft}
Considérons un diagramme commutatif
$$
\xymatrix{
(X \subset X') \ar[rr]_{(f, f')} \ar[rd] & & (Y \subset Y') \ar[ld] \\
& (S \subset S')
}
$$
d'épaississements. Supposons que
\begin{enumerate}
\item $Y' \to S'$ soit localement de type fini,
\item $X' \to S'$ soit plat et localement de présentation finie,
\item $f$ soit plat, et
\item $X = S \times_{S'} X'$ et $Y = S \times_{S'} Y'$.
\end{enumerate}
Alors $f'$ est plat et, pour tout $y' \in Y'$ dans l'image de $f'$,
l'anneau local $\mathcal{O}_{Y', y'}$ est
plat et essentiellement de présentation finie sur $\mathcal{O}_{S', s'}$.
\end{lemma}

\begin{proof}
C'est une conséquence immédiate du résultat d'Algèbre,
lemme \ref{algebra-lemma-criterion-flatness-fibre-locally-nilpotent}.
\end{proof}

\noindent
De nombreuses propriétés des morphismes de schémas sont préservées par les déformations
plates considérées dans le lemme précédent.

\begin{lemma}
\label{lemma-deform-property-fp-over-ft}
Considérons un diagramme commutatif
$$
\xymatrix{
(X \subset X') \ar[rr]_{(f, f')} \ar[rd] & & (Y \subset Y') \ar[ld] \\
& (S \subset S')
}
$$
d'épaississements. Supposons $Y' \to S'$ localement de type fini,
$X' \to S'$ plat et localement de présentation finie,
$X = S \times_{S'} X'$ et $Y = S \times_{S'} Y'$. Alors
\begin{enumerate}
\item $f$ est plat si et seulement si $f'$ est plat,
\label{item-flat-fp-over-ft}
\item $f$ est un isomorphisme si et seulement si $f'$ est un isomorphisme,
\label{item-isomorphism-fp-over-ft}
\item $f$ est une immersion ouverte si et seulement si $f'$ est une immersion ouverte,
\label{item-open-immersion-fp-over-ft}
\item $f$ est quasi-compact si et seulement si $f'$ est quasi-compact,
\label{item-quasi-compact-fp-over-ft}
\item $f$ est universellement fermé si et seulement si $f'$ est universellement fermé,
\label{item-universally-closed-fp-over-ft}
\item $f$ est (quasi-)séparé si et seulement si $f'$ est (quasi-)séparé,
\label{item-separated-fp-over-ft}
\item $f$ est un monomorphisme si et seulement si $f'$ est un monomorphisme,
\label{item-monomorphism-fp-over-ft}
\item $f$ est surjectif si et seulement si $f'$ est surjectif,
\label{item-surjective-fp-over-ft}
\item $f$ est universellement injectif si et seulement si $f'$ est universellement injectif,
\label{item-universally-injective-fp-over-ft}
\item $f$ est affine si et seulement si $f'$ est affine,
\label{item-affine-fp-over-ft}
\item $f$ est localement quasi-fini si et seulement si $f'$ est localement quasi-fini,
\label{item-quasi-finite-fp-over-ft}
\item
\label{item-relative-dimension-d-fp-over-ft}
$f$ est localement de type fini et de dimension relative $d$ si et seulement si
$f'$ est localement de type fini et de dimension relative $d$,
\item $f$ est universellement ouvert si et seulement si $f'$ est universellement ouvert,
\label{item-universally-open-fp-over-ft}
\item $f$ est syntomique si et seulement si $f'$ est syntomique,
\label{item-syntomic-fp-over-ft}
\item $f$ est lisse si et seulement si $f'$ est lisse,
\label{item-smooth-fp-over-ft}
\item $f$ est non ramifié si et seulement si $f'$ est non ramifié,
\label{item-unramified-fp-over-ft}
\item $f$ est étale si et seulement si $f'$ est étale,
\label{item-etale-fp-over-ft}
\item $f$ est propre si et seulement si $f'$ est propre,
\label{item-proper-fp-over-ft}
\item $f$ est fini si et seulement si $f'$ est fini,
\label{item-finite-fp-over-ft}
\item
\label{item-finite-locally-free-fp-over-ft}
$f$ est fini localement libre (de rang $d$) si et seulement si $f'$
est fini localement libre (de rang $d$), et
\item ajouter d'autres propriétés ici.
\end{enumerate}
\end{lemma}

\begin{proof}
Les hypothèses sur $X$ et $Y$ signifient que $f$ est le changement de base de
$f'$ par $X \to X'$.
Les propriétés $\mathcal{P}$ énumérées en (1) -- (20) ci-dessus sont toutes stables
par changement de base ; si $f'$ possède la propriété $\mathcal{P}$,
il en va donc de même de $f$. Voir
Schémas, lemmes \ref{schemes-lemma-base-change-immersion},
\ref{schemes-lemma-quasi-compact-preserved-base-change},
\ref{schemes-lemma-separated-permanence} et
\ref{schemes-lemma-base-change-monomorphism},
et
Morphismes, lemmes
\ref{morphisms-lemma-base-change-surjective},
\ref{morphisms-lemma-base-change-universally-injective},
\ref{morphisms-lemma-base-change-affine},
\ref{morphisms-lemma-base-change-quasi-finite},
\ref{morphisms-lemma-base-change-relative-dimension-d},
\ref{morphisms-lemma-base-change-syntomic},
\ref{morphisms-lemma-base-change-smooth},
\ref{morphisms-lemma-base-change-unramified},
\ref{morphisms-lemma-base-change-etale},
\ref{morphisms-lemma-base-change-proper},
\ref{morphisms-lemma-base-change-finite} et
\ref{morphisms-lemma-base-change-finite-locally-free}.

\medskip\noindent
Dans chaque cas, le sens intéressant consiste donc à supposer
que $f$ possède la propriété et à en déduire que $f'$ la possède aussi.
Faisons quelques remarques générales que l'on utilisera sans autre
mention dans les arguments ci-dessous.
(I) Soient $W' \subset S'$ un ouvert affine et $U' \subset X'$
et $V' \subset Y'$ des ouverts affines au-dessus de $W'$ tels que $f'(U') \subset V'$.
Écrivons $W' = \Spec(R')$ et notons $I \subset R'$ l'idéal
définissant le sous-schéma fermé $W' \cap S$. Écrivons $U' = \Spec(B')$
et $V' = \Spec(A')$. On obtient alors un diagramme commutatif
$$
\xymatrix{
0 \ar[r] &
IB' \ar[r] &
B' \ar[r] &
B \ar[r] & 0 \\
0 \ar[r] &
IA' \ar[r] \ar[u] &
A' \ar[r] \ar[u] &
A \ar[r] \ar[u] & 0
}
$$
à lignes exactes.
(II) Les morphismes $X \to X'$ et $Y \to Y'$ sont des homéomorphismes universels.
Les applications $f$ et $f'$ ont donc la même topologie (après tout changement de base).
(III) Si $f$ est plat, alors $f'$ est plat et $Y' \to S'$ est plat en tout
point de l'image de $f'$, voir
le lemme \ref{lemma-flatness-morphism-thickenings}.

\medskip\noindent
Pour (\ref{item-flat-fp-over-ft}). C'est la remarque générale (III).

\medskip\noindent
Pour (\ref{item-isomorphism-fp-over-ft}). Supposons que $f$ soit un isomorphisme.
Choisissons un ouvert affine $V' \subset Y'$ et posons $U' = (f')^{-1}(V')$.
Alors $V = Y \cap V'$ est affine, ce qui entraîne que
$V \cong f^{-1}(V) = U = Y \times_{Y'} U'$ est affine. D'après
le lemme \ref{lemma-thickening-affine-scheme},
on voit que $U'$ est affine. On dispose donc d'un diagramme comme dans
la remarque générale (I). D'après Algèbre, lemme
\ref{algebra-lemma-isomorphism-modulo-locally-nilpotent},
on voit que $A' \to B'$ est un isomorphisme, c'est-à-dire $U' \cong V'$.
Ainsi, $f'$ est un isomorphisme.

\medskip\noindent
Pour (\ref{item-open-immersion-fp-over-ft}). Supposons que $f$ soit une immersion ouverte.
Alors $f$ est un isomorphisme de $X$ sur un sous-schéma ouvert $V \subset Y$.
Soit $V' \subset Y'$ le sous-schéma ouvert dont l'espace topologique
sous-jacent est $V$. Alors $f'$ est un morphisme de $X'$ dans $V'$, qui est un isomorphisme d'après
(\ref{item-isomorphism-fp-over-ft}). Ainsi, $f'$ est une immersion ouverte.

\medskip\noindent
Pour (\ref{item-quasi-compact-fp-over-ft}). C'est immédiat d'après la remarque (II). Voir aussi
le lemme \ref{lemma-thicken-property-morphisms} pour un énoncé plus général.

\medskip\noindent
Pour (\ref{item-universally-closed-fp-over-ft}). C'est immédiat d'après la remarque (II). Voir
aussi le lemme \ref{lemma-thicken-property-morphisms} pour un énoncé plus général.

\medskip\noindent
Pour (\ref{item-separated-fp-over-ft}). Remarquons que
$X \times_Y X = Y \times_{Y'} (X' \times_{Y'} X')$, de sorte que
$X' \times_{Y'} X'$ est un épaississement de $X \times_Y X$.
Les applications $\Delta_{X/Y}$ et $\Delta_{X'/Y'}$ ont donc la même
topologie, ce qui permet de conclure. Voir aussi
le lemme \ref{lemma-thicken-property-morphisms} pour un énoncé plus général.

\medskip\noindent
Pour (\ref{item-monomorphism-fp-over-ft}). Supposons que $f$ soit un monomorphisme.
Considérons le morphisme diagonal $\Delta_{X'/Y'} : X' \to X' \times_{Y'} X'$.
Remarquons que $X' \times_{Y'} X' \to S'$ est localement de type fini.
Le changement de base de $\Delta_{X'/Y'}$ par $S \to S'$ est $\Delta_{X/Y}$,
qui est un isomorphisme par hypothèse. D'après (\ref{item-isomorphism-fp-over-ft}),
on en déduit que $\Delta_{X'/Y'}$ est un isomorphisme.

\medskip\noindent
Pour (\ref{item-surjective-fp-over-ft}). C'est clair. Voir aussi
le lemme \ref{lemma-thicken-property-morphisms} pour un énoncé plus général.

\medskip\noindent
Pour (\ref{item-universally-injective-fp-over-ft}).
C'est immédiat d'après la remarque (II). Voir aussi
le lemme \ref{lemma-thicken-property-morphisms} pour un énoncé plus général.

\medskip\noindent
Pour (\ref{item-affine-fp-over-ft}). Supposons que $f$ soit affine. Choisissons un
ouvert affine $V' \subset Y'$ et posons $U' = (f')^{-1}(V')$.
Alors $V = Y \cap V'$ est affine, ce qui entraîne que
$U = Y \times_{Y'} U'$ est affine. D'après
le lemme \ref{lemma-thickening-affine-scheme},
on voit que $U'$ est affine. Ainsi, $f'$ est affine. Voir aussi
le lemme \ref{lemma-thicken-property-morphisms} pour un énoncé plus général.

\medskip\noindent
Pour (\ref{item-quasi-finite-fp-over-ft}). Cela résulte du fait que $f'$
est localement de type fini
(d'après Morphismes, lemme \ref{morphisms-lemma-permanence-finite-type}) et du fait que
la quasi-finitude d'un morphisme de type fini se vérifie sur les fibres, voir
Morphismes, lemme \ref{morphisms-lemma-quasi-finite-at-point-characterize}.

\medskip\noindent
Pour (\ref{item-relative-dimension-d-fp-over-ft}).
Cela résulte de la remarque générale (II) et du fait que $f'$
est localement de type fini
(Morphismes, lemme \ref{morphisms-lemma-permanence-finite-type}).

\medskip\noindent
Pour (\ref{item-universally-open-fp-over-ft}).
C'est immédiat d'après la remarque générale (II). Voir aussi
le lemme \ref{lemma-thicken-property-morphisms} pour un énoncé plus général.

\medskip\noindent
Pour (\ref{item-syntomic-fp-over-ft}). Supposons $f$ syntomique. D'après
Morphismes, lemme \ref{morphisms-lemma-finite-presentation-permanence},
$f'$ est localement de présentation finie.
D'après la remarque générale (III), $f'$ est plat. Les fibres de $f'$ sont celles
de $f$. Ainsi, $f'$ est syntomique d'après
Morphismes, lemme \ref{morphisms-lemma-syntomic-flat-fibres}.

\medskip\noindent
Pour (\ref{item-smooth-fp-over-ft}). Supposons $f$ lisse. D'après
Morphismes, lemme \ref{morphisms-lemma-finite-presentation-permanence},
$f'$ est localement de présentation finie.
D'après la remarque générale (III), $f'$ est plat. Les fibres de $f'$ sont les
fibres de $f$. Ainsi, $f'$ est lisse d'après
Morphismes, lemme \ref{morphisms-lemma-smooth-flat-smooth-fibres}.

\medskip\noindent
Pour (\ref{item-unramified-fp-over-ft}). Supposons $f$ non ramifié. D'après
Morphismes, lemme \ref{morphisms-lemma-permanence-finite-type},
$f'$ est localement de type fini. Les fibres de $f'$ sont celles de $f$.
Ainsi, $f'$ est non ramifié d'après
Morphismes, lemme \ref{morphisms-lemma-unramified-etale-fibres}.

\medskip\noindent
Pour (\ref{item-etale-fp-over-ft}). Supposons $f$ étale. D'après
Morphismes, lemme \ref{morphisms-lemma-finite-presentation-permanence},
$f'$ est localement de présentation finie.
D'après la remarque générale (III), $f'$ est plat.
Les fibres de $f'$ sont celles de $f$. Ainsi, $f'$ est étale d'après
Morphismes, lemme \ref{morphisms-lemma-etale-flat-etale-fibres}.

\medskip\noindent
Pour (\ref{item-proper-fp-over-ft}). Cela résulte de la combinaison de
(\ref{item-separated-fp-over-ft}), du fait que $f$ est localement
de type fini (Morphismes, lemme \ref{morphisms-lemma-permanence-finite-type}),
de (\ref{item-quasi-compact-fp-over-ft})
et de (\ref{item-universally-closed-fp-over-ft}).

\medskip\noindent
Pour (\ref{item-finite-fp-over-ft}).
Combiner (\ref{item-universally-closed-fp-over-ft}),
(\ref{item-affine-fp-over-ft}),
Morphismes, lemme \ref{morphisms-lemma-integral-universally-closed},
le fait que $f$ est localement de type fini
(Morphismes, lemme \ref{morphisms-lemma-permanence-finite-type}) et
Morphismes, lemme \ref{morphisms-lemma-finite-integral}.

\medskip\noindent
Pour (\ref{item-finite-locally-free-fp-over-ft}).
Supposons $f$ fini localement libre. D'après
(\ref{item-finite-fp-over-ft}), on voit que $f'$ est fini.
D'après la remarque générale (III), $f'$ est plat.
D'après Morphismes, lemme \ref{morphisms-lemma-finite-presentation-permanence},
$f'$ est localement de présentation finie. Ainsi, $f'$ est fini localement libre d'après
Morphismes, lemme \ref{morphisms-lemma-finite-flat}.
\end{proof}


\begin{lemma}[Déformations des schémas projectifs]
\label{lemma-deform-projective}
Soit $f : X \to S$ un morphisme de schémas propre, plat et
de présentation finie. Soit $\mathcal{L}$ un faisceau $f$-ample. Supposons
$S$ quasi-compact. Il existe un entier $d_0 \geq 0$ tel que,
pour tout diagramme cartésien
$$
\vcenter{
\xymatrix{
X \ar[r]_{i'} \ar[d]_f & X' \ar[d]^{f'} \\
S \ar[r]^i & S'
}
}
\quad\text{et}\quad
\begin{matrix}
\text{un }\mathcal{O}_{X'}\text{-module inversible}\\
\mathcal{L}'\text{ avec }\mathcal{L} \cong (i')^*\mathcal{L}'
\end{matrix}
$$
où $S \subset S'$ est un épaississement et où $f'$ est
propre, plat et de présentation finie, on ait
\begin{enumerate}
\item $R^p(f')_*(\mathcal{L}')^{\otimes d} = 0$
pour tous $p > 0$ et $d \geq d_0$,
\item $\mathcal{A}'_d = (f')_*(\mathcal{L}')^{\otimes d}$
est localement libre de rang fini pour $d \geq d_0$,
\item $\mathcal{A}' =
\mathcal{O}_{S'} \oplus \bigoplus_{d \geq d_0} \mathcal{A}'_d$
est une $\mathcal{O}_{S'}$-algèbre quasi-cohérente de présentation finie,
\item il existe un isomorphisme canonique
$r' : X' \to \underline{\text{Proj}}_{S'}(\mathcal{A}')$, et
\item il existe un isomorphisme canonique
$\theta' : (r')^*\mathcal{O}_{\underline{\text{Proj}}_{S'}(\mathcal{A}')}(1)
\to \mathcal{L}'$.
\end{enumerate}
La construction de $\mathcal{A}'$, $r'$, $\theta'$
est fonctorielle en les données $(X', S', i, i', f', \mathcal{L}')$.
\end{lemma}

\begin{proof}
Décrivons d'abord les morphismes $r'$ et $\theta'$.
Remarquons que $\mathcal{L}'$ est $f'$-ample, voir
le lemme \ref{lemma-thicken-property-relatively-ample}.
Il existe un homomorphisme canonique de
$\mathcal{O}_{S'}$-algèbres graduées quasi-cohérentes
$\mathcal{A}' \to \bigoplus_{d \geq 0} (f')_*(\mathcal{L}')^{\otimes d}$
qui est un isomorphisme en degrés $\geq d_0$.
Il induit donc un isomorphisme des Proj relatifs,
compatible avec les faisceaux structuraux tordus de Serre, voir
Constructions, lemme
\ref{constructions-lemma-eventual-iso-graded-rings-map-relative-proj}.
On obtient ainsi le morphisme $r'$ par
Morphismes, lemme \ref{morphisms-lemma-characterize-relatively-ample}
(qui utilise à son tour la construction donnée dans
Constructions, lemme
\ref{constructions-lemma-invertible-map-into-relative-proj}),
et c'est un isomorphisme d'après
Morphismes, lemme \ref{morphisms-lemma-proper-ample-is-proj}.
L'homomorphisme $\theta'$ est fourni par Constructions, lemme
\ref{constructions-lemma-invertible-map-into-relative-proj}.
D'après Propriétés, lemme \ref{properties-lemma-ample-gcd-is-one},
on voit que $\theta'$ est un isomorphisme
(on utilise aussi le fait que le morphisme $r'$, au-dessus des ouverts
affines de $S'$, est le même que celui de
Propriétés, lemme \ref{properties-lemma-map-into-proj},
comme l'explique la démonstration
de Morphismes, lemme \ref{morphisms-lemma-proper-ample-is-proj}).

\medskip\noindent
En admettant l'annulation et la liberté locale énoncées en
(1) et (2), on peut voir la fonctorialité de la construction comme suit.
Soit $h : T \to S'$ un morphisme de schémas ; notons
$f_T : X'_T \to T$ le changement de base de $f'$ et
$\mathcal{L}_T$ l'image inverse de $\mathcal{L}$ sur $X'_T$.
Par les théorèmes de cohomologie et changement de base
(tels qu'ils sont formulés, par exemple, dans Catégories dérivées des schémas,
lemme \ref{perfect-lemma-compare-base-change}),
on a l'annulation correspondante sur $T$ et, de plus,
$h^*\mathcal{A}'_d = f_{T, *}\mathcal{L}_T^{\otimes d}$
(donc aussi la liberté locale des images directes,
ainsi que le caractère de type fini de la
$\mathcal{O}_T$-algèbre graduée correspondante $\mathcal{A}_T$).
Ainsi, le morphisme
$r_T : X_T \to
\underline{\text{Proj}}_T(\bigoplus f_{T, *}\mathcal{L}_T^{\otimes d})$
est simplement le changement de base de $r'$ à $T$, et l'image inverse de
$\theta'$ est l'homomorphisme $\theta_T$.

\medskip\noindent
Après ces observations, on voit qu'il suffit de prouver
(1), (2) et (3). Choisissons $d_0$ tel que
$R^pf_*\mathcal{L}^{\otimes d} = 0$ pour tous $d \geq d_0$ et $p > 0$, voir
Cohomologie des schémas, lemme \ref{coherent-lemma-coherent-proper-ample}.
On affirme que $d_0$ convient.

\medskip\noindent
Par les théorèmes de cohomologie et changement de base
(Catégories dérivées des schémas,
lemme \ref{perfect-lemma-flat-proper-perfect-direct-image-general}),
on voit que $E'_d = Rf'_*(\mathcal{L}')^{\otimes d}$
est un objet parfait de $D(\mathcal{O}_{S'})$
et que sa formation commute à tout changement de base.
En particulier, $E_d = Li^*E'_d = Rf_*\mathcal{L}^{\otimes d}$.
D'après Catégories dérivées des schémas, lemme
\ref{perfect-lemma-vanishing-implies-locally-free},
on voit que, pour $d \geq d_0$, le complexe $E_d$ est isomorphe au
$\mathcal{O}_S$-module localement libre de rang fini
$f_*\mathcal{L}^{\otimes d}$ placé
en degré cohomologique $0$. D'après
Catégories dérivées des schémas, lemme
\ref{perfect-lemma-open-where-cohomology-in-degree-i-rank-r},
on en déduit que $E'_d$ est isomorphe à un module localement libre de rang fini
placé en degré cohomologique $0$.
Bien entendu, cela signifie que $E'_d = \mathcal{A}'_d[0]$,
que $R^pf'_*(\mathcal{L}')^{\otimes d} = 0$ pour $p > 0$
et que $\mathcal{A}'_d$ est localement libre de rang fini.
Cela prouve (1) et (2).

\medskip\noindent
Il reste à montrer que
$\mathcal{A}'$ est de présentation finie comme faisceau de $\mathcal{O}_{S'}$-algèbres
(cette notion a été introduite dans Propriétés, section
\ref{properties-section-extending-quasi-coherent-sheaves}).
Soit $U' = \Spec(R') \subset S'$ un ouvert affine.
Alors $A' = \mathcal{A}'(U')$ est une $R'$-algèbre graduée
dont les composantes homogènes sont des $R'$-modules projectifs de type fini.
Il faut montrer que $A'$ est une $R'$-algèbre de présentation finie.
On le prouvera en se ramenant au cas noethérien.
En effet, on peut trouver une sous-$\mathbf{Z}$-algèbre de type fini
$R'_0 \subset R'$ et un couple\footnote{Possédant les mêmes propriétés
que $X' \to S'$ et $\mathcal{L}'$, c'est-à-dire
que $X'_0 \to \Spec(R'_0)$ est plat et propre et que $\mathcal{L}'_0$
est ample.} $(X'_0, \mathcal{L}'_0)$ sur $R'_0$
dont le changement de base est $(X'_{U'}, \mathcal{L}'|_{X'_{U'}})$, voir
Limites, lemmes
\ref{limits-lemma-descend-modules-finite-presentation},
\ref{limits-lemma-descend-invertible-modules},
\ref{limits-lemma-eventually-proper},
\ref{limits-lemma-descend-flat-finite-presentation} et
\ref{limits-lemma-limit-ample}.
Cohomologie des schémas, lemme \ref{coherent-lemma-coherent-proper-ample},
entraîne que
$A'_0 = \bigoplus_{d \geq 0} H^0(X'_0, (\mathcal{L}'_0)^{\otimes d})$
est une $R'_0$-algèbre graduée de type fini et
qu'il existe un entier $d'_0$ tel que
$H^p(X'_0, (\mathcal{L}'_0)^{\otimes d}) = 0$, $p > 0$, pour $d \geq d'_0$.
En appliquant les arguments ci-dessus à $X'_0 \to \Spec(R'_0)$ et
$\mathcal{L}'_0$, on voit que $(A'_0)_d$ est un $R'_0$-module projectif de type fini
et que
$$
A'_d = \mathcal{A}'_d(U') =
H^0(X'_{U'}, (\mathcal{L}')^{\otimes d}|_{X'_{U'}}) =
H^0(X'_0, (\mathcal{L}'_0)^{\otimes d}) \otimes_{R'_0} R' =
(A'_0)_d \otimes_{R'_0} R'
$$
pour $d \geq d'_0$. Une petite difficulté est que l'on ne sait pas
si l'on peut choisir $d'_0$ égal à $d_0$\footnote{En fait,
on peut se ramener à ce cas en utilisant davantage d'arguments de passage à la limite.}. Pour
la contourner, on achève la démonstration par la succession d'arguments
suivante :
\begin{enumerate}
\item[(a)] L'algèbre
$B = R'_0 \oplus \bigoplus_{d \geq \max(d_0, d'_0)} (A'_0)_d$
est une $R'_0$-algèbre de type fini : appliquer
le lemme d'Artin--Tate à $B \subset A'_0$, voir
Algèbre, lemme \ref{algebra-lemma-Artin-Tate}.
\item[(b)] Puisque $R'_0$ est noethérien, on voit que
$B$ est une $R'_0$-algèbre de présentation finie.
\item[(c)] Par exactitude à droite du produit tensoriel, on voit que
$B \otimes_{R'_0} R'$ est une $R'$-algèbre de présentation finie.
\item[(d)] D'après les égalités affichées, cela signifie exactement que
$C = R' \oplus \bigoplus_{d \geq \max(d_0, d'_0)} A'_d$
est une $R'$-algèbre de présentation finie.
\item[(e)] Le quotient $A'/C$ est la somme directe des
$R'$-modules projectifs de type fini $A'_d$, $d_0 \leq d \leq \max(d_0, d'_0)$ ;
il est donc de présentation finie comme $R'$-module.
\item[(f)] Le quotient $A'/C$ est de présentation finie
comme $C$-module d'après Algèbre, lemme
\ref{algebra-lemma-finitely-presented-over-subring}.
\item[(g)] Ainsi, $A'$ est de présentation finie comme $C$-module d'après
Algèbre, lemme \ref{algebra-lemma-extension}.
\item[(h)] D'après Algèbre, lemme \ref{algebra-lemma-finite-finite-type},
cela entraîne que $A'$ est de présentation finie comme $C$-algèbre.
\item[(i)] Enfin, d'après
Algèbre, lemme \ref{algebra-lemma-compose-finite-type},
appliqué à $R' \to C \to A'$,
cela entraîne que $A'$ est de présentation finie comme $R'$-algèbre.
\end{enumerate}
Cela achève la démonstration.
\end{proof}












\section{Morphismes formellement lisses}
\label{section-formally-smooth}

\noindent
Selon Michael Artin, les critères différentiels de lissité (par exemple,
Morphismes, lemme \ref{morphisms-lemma-smooth-at-point}) sont
essentiellement inutiles (en pratique). Dans cette section, on introduit la
notion de morphisme formellement lisse $X \to S$. Un tel morphisme est
caractérisé par la propriété que les points à valeurs dans $T$ de $X$ se relèvent
aux épaississements infinitésimaux de $T$, pourvu que $T$ soit affine.
Le résultat principal est qu'un morphisme formellement lisse et
localement de présentation finie est lisse, voir
le lemme \ref{lemma-smooth-formally-smooth}.
Il se trouve que ce critère est souvent plus facile à utiliser que les
critères différentiels mentionnés ci-dessus.

\medskip\noindent
Rappelons qu'un homomorphisme d'anneaux $R \to A$ est dit {\it formellement lisse}
(voir Algèbre, définition \ref{algebra-definition-formally-smooth})
si, pour tout diagramme commutatif de flèches pleines
$$
\xymatrix{
A \ar[r] \ar@{-->}[rd] & B/I \\
R \ar[r] \ar[u] & B \ar[u]
}
$$
où $I \subset B$ est un idéal de carré nul, il existe une flèche
pointillée rendant le diagramme commutatif. Cela conduit
à l'analogue suivant pour les morphismes de schémas.

\begin{definition}
\label{definition-formally-smooth}
Soit $f : X \to S$ un morphisme de schémas.
On dit que $f$ est {\it formellement lisse} si, pour tout diagramme commutatif de flèches pleines
$$
\xymatrix{
X \ar[d]_f & T \ar[d]^i \ar[l] \\
S & T' \ar[l] \ar@{-->}[lu]
}
$$
où $T \subset T'$ est un épaississement du premier ordre de schémas affines sur $S$,
il existe une flèche pointillée rendant le diagramme commutatif.
\end{definition}

\noindent
Dans le cas des morphismes formellement non ramifiés et formellement étales,
on pouvait omettre la condition que $T'$ soit affine, voir
les lemmes \ref{lemma-formally-unramified-not-affine} et
\ref{lemma-formally-etale-not-affine}.
Ce n'est plus vrai dans le cas des morphismes formellement lisses.
Une condition légèrement plus naturelle serait, en fait, que l'on puisse
compléter le diagramme par la flèche pointillée localement pour la topologie de Zariski sur $T'$.
L'analyse de la démonstration du
lemme \ref{lemma-formally-smooth}
montre que cette condition serait équivalente à la définition
actuelle. En particulier, la propriété d'être formellement lisse est
locale sur la source pour la topologie de Zariski (et même pour la topologie lisse ;
insérer ici une référence ultérieure).

\begin{lemma}
\label{lemma-composition-formally-smooth}
Un composé de morphismes formellement lisses est formellement lisse.
\end{lemma}

\begin{proof}
Omise.
\end{proof}

\begin{lemma}
\label{lemma-base-change-formally-smooth}
Tout changement de base d'un morphisme formellement lisse est formellement lisse.
\end{lemma}

\begin{proof}
Omise ; voir cependant Algèbre, lemme \ref{algebra-lemma-base-change-fs},
pour la version algébrique.
\end{proof}

\begin{lemma}
\label{lemma-formally-etale-unramified-smooth}
Soit $f : X \to S$ un morphisme de schémas.
Alors $f$ est formellement étale si et seulement si
$f$ est formellement lisse et formellement non ramifié.
\end{lemma}

\begin{proof}
Omise.
\end{proof}

\begin{lemma}
\label{lemma-formally-smooth-on-opens}
Soit $f : X \to S$ un morphisme de schémas.
Soient $U \subset X$ et $V \subset S$ des sous-schémas ouverts tels que
$f(U) \subset V$. Si $f$ est formellement lisse, il en va de même de $f|_U : U \to V$.
\end{lemma}

\begin{proof}
Considérons un diagramme de flèches pleines
$$
\xymatrix{
U \ar[d]_{f|_U} & T \ar[d]^i \ar[l]^a \\
V & T' \ar[l] \ar@{-->}[lu]
}
$$
comme dans la définition \ref{definition-formally-smooth}. Si $f$ est formellement
lisse, il existe un $S$-morphisme $a' : T' \to X$ tel que
$a'|_T = a$. Puisque $T$ et $T'$ ont le même ensemble sous-jacent,
on voit que $a'$ est un morphisme à valeurs dans $U$ (voir Schémas, section
\ref{schemes-section-open-immersion}). C'est manifestement aussi un $V$-morphisme.
On obtient donc la flèche pointillée voulue.
\end{proof}

\begin{lemma}
\label{lemma-affine-formally-smooth}
Soit $f : X \to S$ un morphisme de schémas.
Supposons que $X$ et $S$ soient affines.
Alors $f$ est formellement lisse si et seulement si
$\mathcal{O}_S(S) \to \mathcal{O}_X(X)$ est un homomorphisme d'anneaux
formellement lisse.
\end{lemma}

\begin{proof}
Cela résulte immédiatement des définitions
(définition \ref{definition-formally-smooth} et
Algèbre, définition \ref{algebra-definition-formally-smooth}),
par l'équivalence entre les catégories des anneaux et des schémas affines,
voir
Schémas, lemme \ref{schemes-lemma-category-affine-schemes}.
\end{proof}

\noindent
Le lemme suivant est le résultat principal de cette section. C'est une victoire du
point de vue fonctoriel : combiné avec
Limites,
proposition \ref{limits-proposition-characterize-locally-finite-presentation},
il permet de reconnaître si un morphisme $f : X \to S$ est lisse à l'aide de
propriétés « simples » du foncteur $h_X : \Sch/S \to \textit{Sets}$.

\begin{lemma}[Critère de relèvement infinitésimal]
\label{lemma-smooth-formally-smooth}
Soit $f : X \to S$ un morphisme de schémas.
Les conditions suivantes sont équivalentes :
\begin{enumerate}
\item Le morphisme $f$ est lisse, et
\item le morphisme $f$ est localement de présentation finie et
formellement lisse.
\end{enumerate}
\end{lemma}

\begin{proof}
Supposons $f : X \to S$ localement de présentation finie et formellement lisse.
Considérons deux ouverts affines $\Spec(A) = U \subset X$ et
$\Spec(R) = V \subset S$
tels que $f(U) \subset V$. D'après le lemme \ref{lemma-formally-smooth-on-opens},
on voit que $U \to V$ est formellement lisse. D'après le lemme
\ref{lemma-affine-formally-smooth}, on voit que $R \to A$ est formellement
lisse. D'après
Morphismes, lemme \ref{morphisms-lemma-locally-finite-presentation-characterize},
on voit que $R \to A$ est de présentation finie.
D'après Algèbre, proposition \ref{algebra-proposition-smooth-formally-smooth},
on voit que $R \to A$ est lisse.
La définition d'un morphisme lisse montre donc que $X \to S$ est lisse.

\medskip\noindent
Réciproquement, supposons que $f : X \to S$ soit lisse. Considérons un diagramme commutatif
de flèches pleines
$$
\xymatrix{
X \ar[d]_f & T \ar[d]^i \ar[l]^a \\
S & T' \ar[l] \ar@{-->}[lu]
}
$$
comme dans la définition \ref{definition-formally-smooth}.
On montrera que la flèche pointillée existe, ce qui
prouvera que $f$ est formellement lisse.

\medskip\noindent
Soit $\mathcal{F}$ le faisceau d'ensembles sur $T'$ du lemme \ref{lemma-sheaf},
dans le cas particulier étudié à la remarque \ref{remark-special-case}.
Soit
$$
\mathcal{H} =
\SheafHom_{\mathcal{O}_T}(a^*\Omega_{X/S}, \mathcal{C}_{T/T'})
$$
le faisceau de $\mathcal{O}_T$-modules muni de l'action
$\mathcal{H} \times \mathcal{F} \to \mathcal{F}$ donnée par
le lemme \ref{lemma-action-sheaf}. Il s'agit simplement
de montrer que $\mathcal{F}(T) \not = \emptyset$. Autrement dit, on
cherche à montrer que $\mathcal{F}$ est un $\mathcal{H}$-torseur trivial
sur $T$ (voir Cohomologie, section \ref{cohomology-section-h1-torsors}).
Il y a deux étapes : (I) Pour montrer que $\mathcal{F}$ est un torseur,
il faut montrer que $\mathcal{F}_t \not = \emptyset$ pour tout $t \in T$ (voir
Cohomologie, définition \ref{cohomology-definition-torsor}).
(II) Pour montrer que $\mathcal{F}$ est le torseur trivial, il suffit
de montrer que $H^1(T, \mathcal{H}) = 0$ (voir
Cohomologie, lemme \ref{cohomology-lemma-torsors-h1} --
on peut utiliser la cohomologie
de $\mathcal{H}$ comme faisceau abélien ou comme $\mathcal{O}_T$-module,
voir Cohomologie, lemme \ref{cohomology-lemma-modules-abelian}).

\medskip\noindent
Prouvons d'abord (I). Pour tout $t \in T$, on peut
choisir un voisinage ouvert affine $U \subset T$ de $t$
tel que $a(U)$ soit contenu
dans un ouvert affine $\Spec(A) = W \subset X$
qui s'envoie dans un ouvert affine $\Spec(R) = V \subset S$.
D'après Morphismes, lemme \ref{morphisms-lemma-smooth-characterize},
l'homomorphisme d'anneaux $R \to A$ est lisse.
D'après Algèbre, proposition \ref{algebra-proposition-smooth-formally-smooth},
l'homomorphisme d'anneaux $R \to A$ est donc formellement lisse.
Le lemme \ref{lemma-affine-formally-smooth}
entraîne à son tour que $W \to V$ est formellement lisse.
On peut donc relever $a|_U : U \to W$ en un $V$-morphisme
$a' : U' \to W \subset X$, ce qui montre que $\mathcal{F}(U) \not = \emptyset$.

\medskip\noindent
Prouvons enfin (II).
D'après Morphismes, lemme \ref{morphisms-lemma-finite-presentation-differentials},
on voit que $\Omega_{X/S}$ est de présentation finie
(il est même localement libre de rang fini d'après
Morphismes, lemme \ref{morphisms-lemma-smooth-omega-finite-locally-free}).
Ainsi, $a^*\Omega_{X/S}$ est de présentation finie (voir
Modules, lemme \ref{modules-lemma-pullback-finite-presentation}).
Le faisceau
$\mathcal{H} =
\SheafHom_{\mathcal{O}_T}(a^*\Omega_{X/S}, \mathcal{C}_{T/T'})$
est donc quasi-cohérent d'après la discussion de
Schémas, section \ref{schemes-section-quasi-coherent}.
D'après Cohomologie des schémas, lemme
\ref{coherent-lemma-quasi-coherent-affine-cohomology-zero},
on a donc $H^1(T, \mathcal{H}) = 0$, comme voulu.
\end{proof}

\noindent
Les modules quasi-cohérents localement projectifs sont définis dans
Propriétés, section \ref{properties-section-locally-projective}.

\begin{lemma}
\label{lemma-formally-smooth-sheaf-differentials}
Soit $f : X \to Y$ un morphisme de schémas formellement lisse.
Alors $\Omega_{X/Y}$ est localement projectif sur $X$.
\end{lemma}

\begin{proof}
Choisissons des ouverts affines $U \subset X$ et $V \subset Y$ tels que
$f(U) \subset V$. D'après
le lemme \ref{lemma-formally-smooth-on-opens},
$f|_U : U \to V$ est formellement lisse. Ainsi,
$\Gamma(V, \mathcal{O}_V) \to \Gamma(U, \mathcal{O}_U)$ est
un homomorphisme d'anneaux formellement lisse, voir
le lemme \ref{lemma-affine-formally-smooth}.
D'après
Algèbre, lemme \ref{algebra-lemma-characterize-formally-smooth-again},
le $\Gamma(U, \mathcal{O}_U)$-module
$\Omega_{\Gamma(U, \mathcal{O}_U)/\Gamma(V, \mathcal{O}_V)}$
est donc projectif. Ainsi, $\Omega_{U/V}$ est localement projectif, voir
Propriétés, section \ref{properties-section-locally-projective}.
\end{proof}

\begin{lemma}
\label{lemma-h1-is-zero}
Soit $T$ un schéma affine. Soient $\mathcal{F}$, $\mathcal{G}$ des
$\mathcal{O}_T$-modules quasi-cohérents. Considérons
$\mathcal{H} = \SheafHom_{\mathcal{O}_T}(\mathcal{F}, \mathcal{G})$.
Si $\mathcal{F}$ est localement projectif, alors $H^1(T, \mathcal{H}) = 0$.
\end{lemma}

\begin{proof}
D'après la définition d'un faisceau localement projectif sur un schéma (voir
Propriétés, définition \ref{properties-definition-locally-projective}),
on voit que $\mathcal{F}$ est un facteur direct d'un
$\mathcal{O}_T$-module libre. On peut donc supposer que
$\mathcal{F} = \bigoplus_{i \in I} \mathcal{O}_T$ est un module libre.
Dans ce cas, $\mathcal{H} = \prod_{i \in I} \mathcal{G}$ est
un produit de modules quasi-cohérents. D'après
Cohomologie, lemme \ref{cohomology-lemma-cohomology-products},
on en déduit que $H^1 = 0$, car la cohomologie d'un faisceau quasi-cohérent
sur un schéma affine est nulle, voir Cohomologie des schémas, lemme
\ref{coherent-lemma-quasi-coherent-affine-cohomology-zero}.
\end{proof}

\begin{lemma}
\label{lemma-formally-smooth}
Soit $f : X \to Y$ un morphisme de schémas. Les conditions suivantes sont équivalentes :
\begin{enumerate}
\item $f$ est formellement lisse,
\item pour tout $x \in X$, il existe des ouverts $x \in U \subset X$ et
$f(x) \in V \subset Y$ avec $f(U) \subset V$ tels que
$f|_U : U \to V$ soit formellement lisse,
\item pour tout couple d'ouverts affines $U \subset X$ et $V \subset Y$
avec $f(U) \subset V$, l'homomorphisme d'anneaux $\mathcal{O}_Y(V) \to \mathcal{O}_X(U)$
est formellement lisse, et
\item il existe un recouvrement ouvert affine $Y = \bigcup V_j$ et,
pour chaque $j$, un recouvrement ouvert affine $f^{-1}(V_j) = \bigcup U_{ji}$
tels que $\mathcal{O}_Y(V) \to \mathcal{O}_X(U)$ soit un homomorphisme d'anneaux
formellement lisse pour tous $j$ et $i$.
\end{enumerate}
\end{lemma}

\begin{proof}
Les implications (1) $\Rightarrow$ (2),
(1) $\Rightarrow$ (3) et (2) $\Rightarrow$ (4) résultent du
lemme \ref{lemma-formally-smooth-on-opens}.
L'implication (3) $\Rightarrow$ (4) est immédiate.

\medskip\noindent
Supposons (4). La démonstration du fait que $f$ est formellement lisse est la même
que la seconde partie de la démonstration du lemme \ref{lemma-smooth-formally-smooth}.
Considérons un diagramme commutatif de flèches pleines
$$
\xymatrix{
X \ar[d]_f & T \ar[d]^i \ar[l]^a \\
Y & T' \ar[l] \ar@{-->}[lu]
}
$$
comme dans la définition \ref{definition-formally-smooth}.
On montrera que la flèche pointillée existe, ce qui
prouvera que $f$ est formellement lisse.
Soit $\mathcal{F}$ le faisceau d'ensembles sur $T'$ du
lemme \ref{lemma-sheaf}, dans le cas particulier étudié à
la remarque \ref{remark-special-case}.
Soit
$$
\mathcal{H} =
\SheafHom_{\mathcal{O}_T}(a^*\Omega_{X/Y}, \mathcal{C}_{T/T'})
$$
le faisceau de $\mathcal{O}_T$-modules sur $T$
muni de l'action $\mathcal{H} \times \mathcal{F} \to \mathcal{F}$ du
lemme \ref{lemma-action-sheaf}.
L'action $\mathcal{H} \times \mathcal{F} \to \mathcal{F}$
fait de $\mathcal{F}$ un pseudo-$\mathcal{H}$-torseur, voir
Cohomologie, définition \ref{cohomology-definition-torsor}.
Il s'agit de montrer que $\mathcal{F}$ est un $\mathcal{H}$-torseur trivial.
Il y a deux étapes : (I) Pour montrer que $\mathcal{F}$ est un torseur,
il faut montrer que $\mathcal{F}$ possède localement une
section. (II) Pour montrer que $\mathcal{F}$ est le torseur trivial,
il suffit de montrer que $H^1(T, \mathcal{H}) = 0$, voir
Cohomologie, lemme \ref{cohomology-lemma-torsors-h1}.

\medskip\noindent
Prouvons d'abord (I). Pour tout $t \in T$, on peut
choisir un voisinage ouvert affine $W \subset T$ de $t$
tel que $a(W)$ soit contenu dans $U_{ji}$ pour certains $i, j$.
Soit $W' \subset T'$ le sous-schéma ouvert correspondant.
Par l'hypothèse (4), on peut relever $a|_W : W \to U_{ji}$
en un $V_j$-morphisme $a' : W' \to U_{ji}$, ce qui montre que
$\mathcal{F}(W')$ est non vide.

\medskip\noindent
Prouvons enfin (II). D'après
le lemme \ref{lemma-formally-smooth-sheaf-differentials},
on voit que $\Omega_{U_{ji}/V_j}$ est localement projectif.
Ainsi, $\Omega_{X/Y}$ est localement projectif, voir
Propriétés, lemme \ref{properties-lemma-locally-projective}.
Ainsi, $a^*\Omega_{X/Y}$ est localement projectif, voir
Propriétés, lemme \ref{properties-lemma-locally-projective-pullback}.
On a donc
$$
H^1(T, \mathcal{H}) =
H^1(T, \SheafHom_{\mathcal{O}_T}(a^*\Omega_{X/Y}, \mathcal{C}_{T/T'})) = 0
$$
d'après
le lemme \ref{lemma-h1-is-zero},
comme voulu.
\end{proof}

\begin{lemma}
\label{lemma-triangle-differentials-formally-smooth}
Soient $f : X \to Y$, $g : Y \to S$ des morphismes de schémas.
Supposons $f$ formellement lisse. Alors
$$
0 \to f^*\Omega_{Y/S} \to \Omega_{X/S} \to \Omega_{X/Y} \to 0
$$
(voir
Morphismes, lemme \ref{morphisms-lemma-triangle-differentials})
est une suite exacte courte.
\end{lemma}

\begin{proof}
La version algébrique de ce lemme est la suivante :
pour des homomorphismes d'anneaux $A \to B \to C$ avec $B \to C$ formellement lisse,
la suite
$$
0 \to C \otimes_B \Omega_{B/A} \to \Omega_{C/A} \to \Omega_{C/B} \to 0
$$
donnée dans
Algèbre, lemme \ref{algebra-lemma-exact-sequence-differentials},
est exacte. C'est
Algèbre, lemme \ref{algebra-lemma-ses-formally-smooth}.
\end{proof}

\begin{lemma}
\label{lemma-differentials-formally-unramified-formally-smooth}
Soit $h : Z \to X$ un morphisme formellement non ramifié de schémas sur $S$.
Supposons que $Z$ soit formellement lisse sur $S$. Alors la
suite exacte canonique
$$
0 \to \mathcal{C}_{Z/X} \to h^*\Omega_{X/S} \to \Omega_{Z/S} \to 0
$$
du
lemme \ref{lemma-universally-unramified-differentials-sequence}
est une suite exacte courte.
\end{lemma}

\begin{proof}
Soit $Z \to Z'$ l'épaississement universel du premier ordre de $Z$ sur $X$.
La démonstration du
lemme \ref{lemma-universally-unramified-differentials-sequence}
montre que notre suite s'identifie à la suite
$$
\mathcal{C}_{Z/Z'} \to \Omega_{Z'/S} \otimes \mathcal{O}_Z \to
\Omega_{Z/S} \to 0.
$$
Puisque $Z \to S$ est formellement lisse, on peut trouver, localement sur $Z'$,
un inverse à gauche $Z' \to Z$ sur $S$ de l'inclusion $Z \to Z'$.
La suite est donc localement scindée, voir
Morphismes, lemme \ref{morphisms-lemma-differentials-relative-immersion-section}.
\end{proof}

\begin{lemma}
\label{lemma-two-unramified-morphisms-formally-smooth}
Soit
$$
\xymatrix{
Z \ar[r]_i \ar[rd]_j & X \ar[d]^f \\
& Y
}
$$
un diagramme commutatif de schémas où $i$ et $j$ sont formellement
non ramifiés et où $f$ est formellement lisse. Alors la suite exacte canonique
$$
0 \to
\mathcal{C}_{Z/Y} \to
\mathcal{C}_{Z/X} \to
i^*\Omega_{X/Y} \to 0
$$
du
lemme \ref{lemma-two-unramified-morphisms}
est exacte et localement scindée.
\end{lemma}

\begin{proof}
Notons $Z \to Z'$ l'épaississement universel du premier ordre de $Z$ sur $X$.
Notons $Z \to Z''$ l'épaississement universel du premier ordre de $Z$ sur $Y$.
D'après
le lemme \ref{lemma-universally-unramified-differentials-sequence},
il existe un morphisme canonique $Z' \to Z''$ donnant un diagramme
commutatif
$$
\xymatrix{
Z \ar[r]_{i'} \ar[rd]_{j'} & Z' \ar[r]_a \ar[d]^k & X \ar[d]^f \\
& Z'' \ar[r]^b & Y
}
$$
Dans la démonstration du
lemme \ref{lemma-two-unramified-morphisms},
on a identifié la suite ci-dessus à la suite
$$
\mathcal{C}_{Z/Z''} \to
\mathcal{C}_{Z/Z'} \to
(i')^*\Omega_{Z'/Z''} \to 0
$$
Soit $U'' \subset Z''$ un ouvert affine. Notons $U \subset Z$ et
$U' \subset Z'$ les sous-schémas ouverts affines correspondants.
Comme $f$ est formellement lisse, il existe un morphisme $h : U'' \to X$
qui coïncide avec $i$ sur $U$ et tel que $f \circ h$ soit égal à $b|_{U''}$.
Puisque $Z'$ est l'épaississement universel du premier ordre, on obtient un unique
morphisme $g : U'' \to Z'$ tel que $g = a \circ h$. La propriété
universelle de $Z''$ entraîne que $k \circ g$ est l'inclusion
$U'' \to Z''$. Ainsi, $g$ est un inverse à droite de $k$. Voici le diagramme :
$$
\xymatrix{
U \ar[d] \ar[r] & Z' \ar[d]^k \\
U'' \ar[r] \ar[ru]^g & Z''
}
$$
Ainsi, $g$ induit un homomorphisme $\mathcal{C}_{Z/Z'}|_U \to \mathcal{C}_{Z/Z''}|_U$
qui est un inverse à gauche de l'homomorphisme
$\mathcal{C}_{Z/Z''} \to \mathcal{C}_{Z/Z'}$ sur $U$.
\end{proof}




































\section{Lissité sur une base noethérienne}
\label{section-smooth-Noetherian}

\noindent
Lorsque la base est noethérienne, on peut alléger
la formulation de la lissité formelle. En un certain sens, les lemmes suivants
constituent les premiers éléments de la théorie des déformations.

\begin{lemma}
\label{lemma-lifting-along-artinian-at-point}
Soit $f : X \to S$ un morphisme de schémas.
Soit $x \in X$.
Supposons $S$ localement noethérien et $f$ localement de type fini.
Les conditions suivantes sont équivalentes :
\begin{enumerate}
\item $f$ est lisse en $x$,
\item pour tout diagramme commutatif de flèches pleines
$$
\xymatrix{
X \ar[d]_f & \Spec(B) \ar[d]^i \ar[l]^-\alpha \\
S & \Spec(B') \ar[l]_-{\beta} \ar@{-->}[lu]
}
$$
où $B' \to B$ est une surjection d'anneaux locaux telle que
$\Ker(B' \to B)$ soit de carré nul et où $\alpha$ envoie le
point fermé de $\Spec(B)$ sur $x$, il existe
une flèche pointillée rendant le diagramme commutatif,
\item même condition qu'en (2), mais en faisant parcourir à $B' \to B$ les petites
extensions (voir Algèbre, définition \ref{algebra-definition-small-extension}),
et
\item même condition qu'en (2), mais en faisant parcourir à $B' \to B$ les petites
extensions telles que $\alpha$ induise un isomorphisme
$\kappa(x) \to \kappa(\mathfrak m)$, où $\mathfrak m \subset B$
est l'idéal maximal.
\end{enumerate}
\end{lemma}

\begin{proof}
Choisissons un voisinage affine $V \subset S$ de $f(x)$ et un
voisinage affine $U \subset X$ de $x$ tel que $f(U) \subset V$.
Pour tout diagramme « test » comme en (2), le morphisme $\alpha$ enverra
$\Spec(B)$ dans $U$ et le morphisme $\beta$ enverra $\Spec(B')$
dans $V$ (voir Schémas, section \ref{schemes-section-points}).
Le lemme se ramène donc au morphisme $f|_U : U \to V$ de schémas affines.
(En effet, $V$ est noethérien et $f|_U$ est de type fini, voir
Propriétés, lemme \ref{properties-lemma-locally-Noetherian}, et
Morphismes, lemme \ref{morphisms-lemma-locally-finite-type-characterize}.)
Dans ce cas affine, le lemme est identique à
Algèbre, lemme \ref{algebra-lemma-smooth-test-artinian}.
\end{proof}

\noindent
Il est parfois utile de savoir qu'il suffit de vérifier le
critère de relèvement pour les petites extensions « centrées » en des points
de type fini (voir
Morphismes, section \ref{morphisms-section-points-finite-type}).

\begin{lemma}
\label{lemma-lifting-along-artinian}
Soit $f : X \to S$ un morphisme de schémas.
Supposons $S$ localement noethérien et $f$ localement de type fini.
Les conditions suivantes sont équivalentes :
\begin{enumerate}
\item $f$ est lisse,
\item pour tout diagramme commutatif de flèches pleines
$$
\xymatrix{
X \ar[d]_f & \Spec(B) \ar[d]^i \ar[l]^-\alpha \\
S & \Spec(B') \ar[l]_-{\beta} \ar@{-->}[lu]
}
$$
où $B' \to B$ est une petite extension d'anneaux locaux artiniens
et où $\beta$ est de type fini (!), il existe une flèche pointillée rendant
le diagramme commutatif.
\end{enumerate}
\end{lemma}

\begin{proof}
Si $f$ est lisse, le critère de relèvement infinitésimal
(lemme \ref{lemma-smooth-formally-smooth}) affirme que
$f$ est formellement lisse, et (2) est vérifiée.

\medskip\noindent
Supposons (2). L'ensemble des points $x \in X$ où $f$ n'est pas lisse
forme un fermé $T$ de $X$. D'après la discussion de Morphismes,
section \ref{morphisms-section-points-finite-type}, si $T \not = \emptyset$,
il existe un point $x \in T \subset X$ tel que le morphisme
$$
\Spec(\kappa(x)) \to X \to S
$$
soit de type fini (il suffit de choisir un point $x$ de $T$ qui soit fermé
dans un ouvert affine de $X$). D'après
Morphismes, lemme \ref{morphisms-lemma-artinian-finite-type}, pour tout
anneau local artinien $B'$ de corps résiduel $\kappa(x)$, tout
morphisme $\beta : \Spec(B') \to S$ est de type fini. Ainsi,
on voit que tous les diagrammes utilisés dans
le lemme \ref{lemma-lifting-along-artinian-at-point} (4) correspondent
à des diagrammes du présent lemme (2). Donc $X \to S$ est lisse
en $x$, contradiction.
\end{proof}

\noindent
Voici une application utile.

\begin{lemma}
\label{lemma-check-smoothness-on-infinitesimal-nbhds}
Soit $f : X \to S$ un morphisme de type fini entre schémas localement noethériens.
Soit $Z \subset S$ un sous-schéma fermé de voisinage infinitésimal d'ordre $n$
$Z_n \subset S$. Posons $X_n = Z_n \times_S X$.
\begin{enumerate}
\item Si $X_n \to Z_n$ est lisse pour tout $n$, alors $f$
est lisse en tout point de $f^{-1}(Z)$.
\item Si $X_n \to Z_n$ est étale pour tout $n$, alors $f$
est étale en tout point de $f^{-1}(Z)$.
\end{enumerate}
\end{lemma}

\begin{proof}
Supposons $X_n \to Z_n$ lisse pour tout $n$.
Soit $x \in X$ un point au-dessus d'un point de $Z$.
Étant donnés une petite extension $B' \to B$ et des morphismes $\alpha$, $\beta$ comme dans
le lemme \ref{lemma-lifting-along-artinian-at-point} (3),
l'idéal maximal de $B'$ est nilpotent (puisque $B'$ est artinien) ;
le morphisme $\beta$ se factorise donc par $Z_n$ et $\alpha$
par $X_n$ pour un entier $n$ convenable. La propriété de relèvement de
$X_n \to Z_n$ fournit alors la flèche pointillée voulue.
Cela prouve (1). Le point (2) résulte de (1) et du fait qu'un morphisme
est étale si et seulement s'il est lisse de dimension relative $0$.
\end{proof}

\begin{lemma}
\label{lemma-check-flatness-on-infinitesimal-nbhds}
Soit $f : X \to S$ un morphisme de schémas localement noethériens.
Soit $Z \subset S$ un sous-schéma fermé de voisinage infinitésimal d'ordre $n$
$Z_n \subset S$. Posons $X_n = Z_n \times_S X$.
Si $X_n \to Z_n$ est plat pour tout $n$, alors $f$
est plat en tout point de $f^{-1}(Z)$.
\end{lemma}

\begin{proof}
C'est la traduction d'Algèbre, lemme \ref{algebra-lemma-flat-module-powers},
dans le langage des schémas.
\end{proof}










\section{Le complexe cotangent naïf}
\label{section-netherlander}

\noindent
Cette section poursuit
Modules, section \ref{modules-section-netherlander},
qui prolonge elle-même la discussion d'Algèbre,
section \ref{algebra-section-netherlander}.

\begin{definition}
\label{definition-netherlander}
Soit $f : X \to Y$ un morphisme de schémas.
Le {\it complexe cotangent naïf de $f$}
est le complexe défini dans Modules, définition
\ref{modules-definition-cotangent-complex-morphism-ringed-topoi}.
Notation : $\NL_f$ ou $\NL_{X/Y}$.
\end{definition}

\begin{lemma}
\label{lemma-NL-affine}
Soit $f : X \to Y$ un morphisme de schémas. Soient
$\Spec(A) = U \subset X$ et $\Spec(R) = V \subset S$
des ouverts affines tels que $f(U) \subset V$.
Il existe un morphisme canonique
$$
\widetilde{\NL_{A/R}} \longrightarrow \NL_{X/Y}|_U
$$
de complexes, qui est un isomorphisme dans $D(\mathcal{O}_U)$.
\end{lemma}

\begin{proof}
La construction de $\NL_{X/Y}$ dans
Modules, section \ref{modules-section-netherlander},
fournit un morphisme canonique de complexes
$\NL_{\mathcal{O}_X(U)/f^{-1}\mathcal{O}_Y(U)} \to \NL_{X/Y}(U)$
de $A = \mathcal{O}_X(U)$-modules, compatible
aux restrictions ultérieures. À l'aide de l'homomorphisme canonique
$R \to f^{-1}\mathcal{O}_Y(U)$, on obtient un morphisme canonique
$\NL_{A/R} \to \NL_{\mathcal{O}_X(U)/f^{-1}\mathcal{O}_Y(U)}$
de complexes de $A$-modules.
Par la propriété universelle du foncteur $\widetilde{\ }$
(voir Schémas, lemme \ref{schemes-lemma-compare-constructions}),
on obtient un morphisme comme dans l'énoncé du lemme.
Pour vérifier qu'il induit des isomorphismes sur les faisceaux de cohomologie,
on peut vérifier qu'il induit des isomorphismes sur leurs fibres.
Cela résulte des résultats d'Algèbre,
lemmes \ref{algebra-lemma-NL-localize-bottom} et
\ref{algebra-lemma-localize-NL},
et de
Modules, lemme \ref{modules-lemma-stalk-NL}
(ainsi que de la description des fibres de
$\mathcal{O}_X$ et $f^{-1}\mathcal{O}_Y$
en un point $\mathfrak p \in \Spec(A)$ comme $A_\mathfrak p$ et
$R_\mathfrak q$, où $\mathfrak q = R \cap \mathfrak p$ ; les références
utilisées sont Schémas, lemme \ref{schemes-lemma-spec-sheaves},
et
Faisceaux, lemme \ref{sheaves-lemma-stalk-pullback}).
\end{proof}

\begin{lemma}
\label{lemma-netherlander-quasi-coherent}
Soit $f : X \to Y$ un morphisme de schémas. Les faisceaux de cohomologie
du complexe $\NL_{X/Y}$ sont quasi-cohérents, nuls en dehors des
degrés $-1$, $0$, et égaux à $\Omega_{X/Y}$ en degré $0$.
\end{lemma}

\begin{proof}
Par construction du complexe cotangent naïf dans
Modules, section \ref{modules-section-netherlander},
$\NL_{X/Y}$ est un complexe placé en degrés $-1$, $0$,
et sa cohomologie en degré $0$ est $\Omega_{X/Y}$.
Le faisceau des différentielles est quasi-cohérent (d'après
Morphismes, lemme \ref{morphisms-lemma-differentials-diagonal}).
Pour achever la démonstration, il suffit de montrer que $H^{-1}(\NL_{X/Y})$
est quasi-cohérent. Cela se vérifie sur les ouverts affines
à l'aide du lemme \ref{lemma-NL-affine}.
\end{proof}

\begin{lemma}
\label{lemma-netherlander-fp}
Soit $f : X \to Y$ un morphisme de schémas. Si $f$ est localement de
présentation finie, alors $\NL_{X/Y}$ est, localement sur $X$, quasi-isomorphe à
un complexe
$$
\ldots \to 0 \to \mathcal{F}^{-1} \to \mathcal{F}^0 \to 0 \to \ldots
$$
de $\mathcal{O}_X$-modules quasi-cohérents,
où $\mathcal{F}^0$ est de présentation finie
et $\mathcal{F}^{-1}$ de type fini.
\end{lemma}

\begin{proof}
D'après le lemme \ref{lemma-NL-affine}, il suffit de montrer que $\NL_{A/R}$
a cette forme si $R \to A$ est un homomorphisme d'anneaux de présentation finie.
Écrivons $A = R[x_1, \ldots, x_n]/I$ avec $I$ de type fini.
Alors $I/I^2$ est un
$A$-module de type fini et $\NL_{A/R}$ est quasi-isomorphe à
$$
\ldots \to 0 \to I/I^2 \to
\bigoplus\nolimits_{i = 1, \ldots, n} A\text{d}x_i \to 0 \to \ldots
$$
d'après Algèbre, section \ref{algebra-section-netherlander},
et en particulier
Algèbre, lemme \ref{algebra-lemma-NL-homotopy}.
\end{proof}

\begin{lemma}
\label{lemma-NL-formally-smooth}
Soit $f : X \to Y$ un morphisme de schémas. Les conditions suivantes sont équivalentes :
\begin{enumerate}
\item $f$ est formellement lisse,
\item $H^{-1}(\NL_{X/Y}) = 0$ et $H^0(\NL_{X/Y}) = \Omega_{X/Y}$
est localement projectif.
\end{enumerate}
\end{lemma}

\begin{proof}
Cela résulte d'Algèbre, proposition
\ref{algebra-proposition-characterize-formally-smooth},
et du lemme \ref{lemma-formally-smooth}.
\end{proof}

\begin{lemma}
\label{lemma-NL-formally-etale}
Soit $f : X \to Y$ un morphisme de schémas. Les conditions suivantes sont équivalentes :
\begin{enumerate}
\item $f$ est formellement étale,
\item $H^{-1}(\NL_{X/Y}) = H^0(\NL_{X/Y}) = 0$.
\end{enumerate}
\end{lemma}

\begin{proof}
Un morphisme formellement étale est formellement lisse ; on a donc
$H^{-1}(\NL_{X/Y}) = 0$ d'après le lemme \ref{lemma-NL-formally-smooth}.
D'autre part, on a $\Omega_{X/Y} = 0$ d'après
le lemme \ref{lemma-characterize-formally-etale}.
Réciproquement, si (2) est vérifiée, alors $f$ est formellement lisse d'après
le lemme \ref{lemma-NL-formally-smooth}
et formellement non ramifié d'après
le lemme \ref{lemma-formally-unramified-differentials},
donc formellement étale d'après
le lemme \ref{lemma-formally-etale-unramified-smooth}.
\end{proof}

\begin{lemma}
\label{lemma-NL-smooth}
Soit $f : X \to Y$ un morphisme de schémas. Les conditions suivantes sont équivalentes :
\begin{enumerate}
\item $f$ est lisse, et
\item $f$ est localement de présentation finie,
$H^{-1}(\NL_{X/Y}) = 0$, et $H^0(\NL_{X/Y}) = \Omega_{X/Y}$
est localement libre de rang fini.
\end{enumerate}
\end{lemma}

\begin{proof}
Cela résulte de la définition d'un homomorphisme d'anneaux lisse
(Algèbre, définition \ref{algebra-definition-smooth}),
du lemme \ref{lemma-NL-affine} et
de la définition d'un morphisme de schémas lisse
(Morphismes, définition \ref{morphisms-definition-smooth}).
On utilise aussi le fait que, pour les modules sur un anneau,
être localement libre de rang fini équivaut à être projectif de type fini
(Algèbre, lemme \ref{algebra-lemma-finite-projective}).
\end{proof}

\begin{lemma}
\label{lemma-NL-etale}
Soit $f : X \to Y$ un morphisme de schémas. Les conditions suivantes sont équivalentes :
\begin{enumerate}
\item $f$ est étale, et
\item $f$ est localement de présentation finie et
$H^{-1}(\NL_{X/Y}) = H^0(\NL_{X/Y}) = 0$.
\end{enumerate}
\end{lemma}

\begin{proof}
Cela résulte de la définition d'un homomorphisme d'anneaux étale
(Algèbre, définition \ref{algebra-definition-etale}),
du lemme \ref{lemma-NL-affine} et
de la définition d'un morphisme de schémas étale
(Morphismes, définition \ref{morphisms-definition-etale}).
\end{proof}


\begin{lemma}
\label{lemma-NL-immersion}
Soit $i : Z \to X$ une immersion de schémas. Alors $\NL_{Z/X}$
est isomorphe à $\mathcal{C}_{Z/X}[1]$ dans $D(\mathcal{O}_Z)$,
où $\mathcal{C}_{Z/X}$ est le faisceau conormal de $Z$ dans $X$.
\end{lemma}

\begin{proof}
Cela résulte d'Algèbre, lemme \ref{algebra-lemma-NL-surjection},
de Morphismes, lemme \ref{morphisms-lemma-affine-conormal}, et
du lemme \ref{lemma-NL-affine}.
\end{proof}

\begin{lemma}
\label{lemma-exact-sequence-NL}
Soient $f : X \to Y$ et $g : Y \to Z$ des morphismes de schémas.
Il existe une suite exacte canonique à six termes
$$
\begin{aligned}
H^{-1}(f^*\NL_{Y/Z}) \to H^{-1}(\NL_{X/Z}) &\to H^{-1}(\NL_{X/Y}) \\
&\to f^*\Omega_{Y/Z} \to \Omega_{X/Z} \to \Omega_{X/Y} \to 0
\end{aligned}
$$
de faisceaux de cohomologie.
\end{lemma}

\begin{proof}
C'est un cas particulier de
Modules, lemme \ref{modules-lemma-exact-sequence-NL-ringed-topoi}.
\end{proof}

\begin{lemma}
\label{lemma-get-triangle-NL}
Soient $f : X \to Y$ et $Y \to Z$ des morphismes de schémas. Supposons que
$X \to Y$ soit un morphisme d'intersection complète. Il existe alors
un triangle distingué canonique
$$
f^*\NL_{Y/Z} \to \NL_{X/Z} \to \NL_{X/Y} \to f^*\NL_{Y/Z}[1]
$$
dans $D(\mathcal{O}_X)$, qui redonne la suite exacte à $6$ termes du
lemme \ref{lemma-exact-sequence-NL}.
\end{lemma}

\begin{proof}
Il suffit de montrer que le morphisme canonique
$$
f^*\NL_{Y/Z} \to \text{Cone}(\NL_{X/Z} \to \NL_{X/Y})[-1]
$$
de Modules, lemme \ref{modules-lemma-exact-sequence-NL-ringed-topoi},
est un isomorphisme dans $D(\mathcal{O}_X)$. Pour cela,
il suffit de montrer que la suite à $6$ termes commence par
un zéro à gauche, c'est-à-dire que $H^{-1}(f^*\NL_{Y/Z}) \to H^{-1}(\NL_{X/Z})$
est injectif. Localement sur les ouverts affines, cela résulte du
résultat algébrique correspondant dans Compléments d'algèbre, lemme
\ref{more-algebra-lemma-transitive-lci-at-end}.
Pour effectuer cette traduction algébrique, utiliser le lemme \ref{lemma-NL-affine}.
\end{proof}

\begin{lemma}
\label{lemma-smooth-etale-permanence}
Soient $X \to Y \to Z$ des morphismes de schémas. Supposons $X \to Z$ lisse
et $Y \to Z$ étale. Alors $X \to Y$ est lisse.
\end{lemma}

\begin{proof}
Le morphisme $X \to Y$ est localement de présentation finie d'après
Morphismes, lemme \ref{morphisms-lemma-finite-presentation-permanence}.
D'après le lemme \ref{lemma-NL-smooth}, on a $H^{-1}(\NL_{X/Z}) = 0$
et le module $\Omega_{X/Z}$ est localement libre de rang fini.
D'après le lemme \ref{lemma-NL-etale}, on a
$H^{-1}(\NL_{Y/Z}) = H^0(\NL_{Y/Z}) = 0$.
D'après le lemme \ref{lemma-exact-sequence-NL}, on obtient
$H^{-1}(\NL_{X/Y}) = 0$ et $\Omega_{X/Y} \cong \Omega_{X/Z}$
est localement libre de rang fini.
D'après le lemme \ref{lemma-NL-smooth}, le morphisme $X \to Y$ est lisse.
\end{proof}

\begin{lemma}
\label{lemma-get-NL}
Soit $f : X \to Y$ un morphisme de schémas qui se factorise
en $f = g \circ i$, où $i$ est une immersion et $g : P \to Y$
est formellement lisse (par exemple lisse). Il existe alors un isomorphisme canonique
$$
\NL_{X/Y} \cong \left(\mathcal{C}_{X/P} \to i^*\Omega_{P/Y}\right)
$$
dans $D(\mathcal{O}_X)$, où le faisceau conormal $\mathcal{C}_{X/P}$
est placé en degré $-1$.
\end{lemma}

\begin{proof}
(Pour l'assertion entre parenthèses, voir le lemme \ref{lemma-smooth-formally-smooth}.)
D'après les lemmes \ref{lemma-NL-immersion} et \ref{lemma-NL-formally-smooth}, on a
$\NL_{X/P} = \mathcal{C}_{X/P}[1]$ et $\NL_{P/Y} = \Omega_{P/Y}$, avec
$\Omega_{P/Y}$ localement projectif. Cela entraîne que
$i^*\NL_{P/Y} \to i^*\Omega_{P/Y}$ est aussi un quasi-isomorphisme
(un petit détail est omis ; la raison est que $i^*\NL_{P/Y}$ est
la même chose que $\tau_{\geq -1}Li^*\NL_{P/Y}$, voir Compléments d'algèbre, lemme
\ref{more-algebra-lemma-tensor-NL}).
Ainsi, le morphisme canonique
$$
i^*\NL_{P/Y} \to \text{Cone}(\NL_{X/Y} \to \NL_{X/P})[-1]
$$
de Modules, lemme \ref{modules-lemma-exact-sequence-NL-ringed-topoi},
est un isomorphisme dans $D(\mathcal{O}_X)$, car le groupe de cohomologie
$H^{-1}(i^*\NL_{P/Y})$ est nul d'après ce qui précède.
Autrement dit, on dispose d'un triangle distingué
$$
i^*\NL_{P/Y} \to \NL_{X/Y} \to \NL_{X/P} \to i^*\NL_{P/Y}[1]
$$
Cela signifie manifestement que $\NL_{X/Y}$ est le cône du morphisme
$\NL_{X/P}[-1] \to i^*\NL_{P/Y}$, ce qui équivaut
à l'énoncé du lemme, d'après le calcul effectué ci-dessus des faisceaux
de cohomologie de ces objets de la catégorie dérivée.
\end{proof}

\begin{lemma}
\label{lemma-base-change-NL}
Considérons un diagramme cartésien de schémas
$$
\xymatrix{
X' \ar[r]_{g'} \ar[d] & X \ar[d] \\
Y' \ar[r] & Y
}
$$
Le morphisme canonique $(g')^*\NL_{X/Y} \to \NL_{X'/Y'}$ induit
un isomorphisme sur $H^0$ et une surjection sur $H^{-1}$.
\end{lemma}

\begin{proof}
La traduction algébrique de cet énoncé est
Compléments d'algèbre, lemme \ref{more-algebra-lemma-base-change-NL}.
Pour effectuer la traduction, utiliser le lemme \ref{lemma-NL-affine}.
\end{proof}

\begin{lemma}
\label{lemma-flat-base-change-NL}
Considérons un diagramme cartésien de schémas
$$
\xymatrix{
X' \ar[d] \ar[r]_{g'} & X \ar[d] \\
Y' \ar[r] & Y
}
$$
Si $Y' \to Y$ est plat, alors le morphisme canonique
$(g')^*\NL_{X/Y} \to \NL_{X'/Y'}$ est un quasi-isomorphisme.
\end{lemma}

\begin{proof}
D'après le lemme \ref{lemma-NL-affine}, cela résulte d'Algèbre,
lemme \ref{algebra-lemma-change-base-NL}.
\end{proof}

\begin{lemma}
\label{lemma-base-change-NL-flat}
Considérons un diagramme cartésien de schémas
$$
\xymatrix{
X' \ar[r]_{g'} \ar[d] & X \ar[d] \\
Y' \ar[r] & Y
}
$$
Si $X \to Y$ est plat, alors le morphisme canonique
$(g')^*\NL_{X/Y} \to \NL_{X'/Y'}$ est un quasi-isomorphisme.
Si, de plus, $\NL_{X/Y}$ a une amplitude de Tor contenue dans $[-1, 0]$,
alors $L(g')^*\NL_{X/Y} \to \NL_{X'/Y'}$ est aussi un quasi-isomorphisme.
\end{lemma}

\begin{proof}
La traduction algébrique de cet énoncé est
Compléments d'algèbre, lemme \ref{more-algebra-lemma-base-change-NL-flat}.
Pour effectuer la traduction, utiliser le lemme \ref{lemma-NL-affine}
et Catégories dérivées des schémas, lemmes
\ref{perfect-lemma-affine-compare-bounded} et
\ref{perfect-lemma-tor-dimension-affine}.
\end{proof}














\section{Sommes amalgamées dans la catégorie des schémas, I}
\label{section-pushouts}

\noindent
Dans cette section, on construit les sommes amalgamées de $Y \leftarrow X \rightarrow X'$
lorsque $X \to Y$ est affine et que $X \to X'$ est un épaississement. Ce cas
jouera un rôle important par la suite, ce qui justifie une discussion détaillée.
Dans la section \ref{section-pushouts-II}, on traite un cas plus intéressant
et plus difficile. Voir Catégories, section
\ref{categories-section-pushouts}, pour une discussion générale
des sommes amalgamées dans une catégorie quelconque.

\begin{lemma}
\label{lemma-basic-example-pushout}
Soient $A' \to A$ une surjection d'anneaux et $B \to A$ un homomorphisme d'anneaux.
Soit $B' = B \times_A A'$ le produit fibré d'anneaux. Posons
$S = \Spec(A)$, $S' = \Spec(A')$, $T = \Spec(B)$ et $T' = \Spec(B')$.
Alors
$$
\vcenter{
\xymatrix{
S \ar[r]_i \ar[d]_f & S' \ar[d]^{f'} \\
T \ar[r]^{i'} & T'
}
}
\quad\text{correspondant à}\quad
\vcenter{
\xymatrix{
A & A' \ar[l] \\
B \ar[u] & B' \ar[l] \ar[u]
}
}
$$
est une somme amalgamée de schémas.
\end{lemma}

\begin{proof}
D'après Compléments d'algèbre, lemme \ref{more-algebra-lemma-points-of-fibre-product},
on a $T' = T \amalg_S S'$ comme espaces topologiques, c'est-à-dire que le diagramme
est une somme amalgamée dans la catégorie des espaces topologiques. Considérons ensuite
l'homomorphisme
$$
((i')^\sharp, (f')^\sharp) :
\mathcal{O}_{T'}
\longrightarrow
i'_*\mathcal{O}_T \times_{g_*\mathcal{O}_S} f'_*\mathcal{O}_{S'}
$$
où $g = i' \circ f = f' \circ i$. On affirme que cet homomorphisme est un isomorphisme de
faisceaux d'anneaux. En effet, on peut considérer les deux membres comme des
$\mathcal{O}_{T'}$-modules quasi-cohérents (utiliser
Schémas, lemme \ref{schemes-lemma-push-forward-quasi-coherent},
pour le membre de droite), et l'homomorphisme est $\mathcal{O}_{T'}$-linéaire.
Il suffit donc de montrer qu'il est un isomorphisme au niveau
des sections globales
(Schémas, lemme \ref{schemes-lemma-equivalence-quasi-coherent}).
Sur les sections globales, on retrouve l'identification
$B' \to B \times_A A'$ de l'énoncé du lemme (c'est ainsi
que l'on a choisi $B'$).

\medskip\noindent
Soit $X$ un schéma. Supposons donnés des morphismes de schémas
$m' : S' \to X$ et $n : T \to X$ tels que $m' \circ i = n \circ f$
(notons ce morphisme $m$). On obtient une unique application d'espaces topologiques
$n' : T' \to X$ compatible avec $m'$ et $n$, puisque
$T' = T \amalg_S S'$ (voir ci-dessus).
D'après la description de $\mathcal{O}_{T'}$ au paragraphe
précédent, on obtient un unique homomorphisme de faisceaux d'anneaux
$$
(n')^\sharp :
\mathcal{O}_X
\longrightarrow
(n')_*\mathcal{O}_{T'} =
m'_*\mathcal{O}_T \times_{m_*\mathcal{O}_T} n_*\mathcal{O}_S
$$
donné par $(m')^\sharp$ et $n^\sharp$.
Ainsi, $(n', (n')^\sharp)$ est l'unique morphisme
d'espaces annelés $T' \to X$ compatible avec $m'$ et $n$.
Pour achever la démonstration, il suffit de montrer que $n'$
est un morphisme de schémas, c'est-à-dire un morphisme d'espaces localement annelés.

\medskip\noindent
Soit $t' \in T'$ d'image $x \in X$.
Il faut montrer que $\mathcal{O}_{X, x} \to \mathcal{O}_{T', t'}$
est local. Si $t' \not \in T$, alors $t'$ est l'image d'un unique
point $s' \in S'$ et $\mathcal{O}_{T', t'} = \mathcal{O}_{S', s'}$.
En effet, $S' \setminus S \to T' \setminus T$ est un isomorphisme de
schémas, car $B' \to A'$ induit un isomorphisme
$\Ker(B' \to B) = \Ker(A' \to A)$.
Si $t'$ est l'image de $t \in T$,
on sait que le composé
$\mathcal{O}_{X, x} \to \mathcal{O}_{T', t'} \to \mathcal{O}_{T, t}$
est local, ce qui permet aussi de conclure.
\end{proof}

\begin{lemma}
\label{lemma-pushout-fpqc-local}
Soit $\mathcal{I} \to (\Sch/S)_{fppf}$, $i \mapsto X_i$, un diagramme de schémas.
Soit $(W, X_i \to W)$ un cocône sur ce diagramme dans la catégorie des schémas
(Catégories, remarque \ref{categories-remark-cones-and-cocones}).
S'il existe un recouvrement fpqc $\{W_a \to W\}_{a \in A}$ de schémas tel que
\begin{enumerate}
\item pour tout $a \in A$, on ait
$W_a = \colim X_i \times_W W_a$
dans la catégorie des schémas, et
\item pour tous $a, b \in A$, on ait
$W_a \times_W W_b = \colim X_i \times_W W_a \times_W W_b$
dans la catégorie des schémas,
\end{enumerate}
alors $W = \colim X_i$ dans la catégorie des schémas.
\end{lemma}

\begin{proof}
En effet, pour un schéma $T$, un morphisme $W \to T$ revient à se donner
une famille de morphismes $W_a \to T$, $a \in A$, qui coïncident sur les
intersections $W_a \times_W W_b$, voir
Descente, lemme \ref{descent-lemma-fpqc-universal-effective-epimorphisms}.
\end{proof}

\begin{lemma}
\label{lemma-pushout-along-thickening}
Soit $X \to X'$ un épaississement de schémas et soit $X \to Y$ un morphisme
affine de schémas. Il existe alors une somme amalgamée
$$
\xymatrix{
X \ar[r] \ar[d]_f
&
X' \ar[d]^{f'}
\\
Y \ar[r]
&
Y'
}
$$
dans la catégorie des schémas. De plus, $Y \subset Y'$ est un
épaississement, $X = Y \times_{Y'} X'$, et
$$
\mathcal{O}_{Y'} = \mathcal{O}_Y \times_{f_*\mathcal{O}_X} f'_*\mathcal{O}_{X'}
$$
comme faisceaux sur $|Y| = |Y'|$.
\end{lemma}

\begin{proof}
Construisons d'abord $Y'$ comme espace annelé. Comme espace
topologique, on prend $Y' = Y$. Notons $f' : X' \to Y'$ l'application d'espaces topologiques
égale à $f$. Comme faisceau structural $\mathcal{O}_{Y'}$, on prend
le membre de droite de l'égalité du lemme. Pour montrer que
$Y'$ est un schéma, il faut montrer que tout point possède un voisinage
affine. Puisque la formation du produit fibré de faisceaux
commute à la restriction aux ouverts, on peut supposer $Y$ affine.
Alors $X$ est affine (puisque $f$ est affine) et $X'$ l'est aussi
(voir le lemme \ref{lemma-thickening-affine-scheme}).
Supposons que $Y \leftarrow X \rightarrow X'$ corresponde
à $B \rightarrow A \leftarrow A'$. Posons $B' = B \times_A A'$ ;
c'est l'anneau des sections globales de $\mathcal{O}_{Y'}$. Comme $A' \to A$ est surjectif
à noyau localement nilpotent, on voit que $B' \to B$ est surjectif
à noyau localement nilpotent. Ainsi, $\Spec(B') = \Spec(B)$ (comme
espaces topologiques). On affirme que $Y' = \Spec(B')$. Pour cela,
on montrera que, pour $g' \in B'$ d'image $g \in B$,
on a $\mathcal{O}_{Y'}(D(g)) = B'_{g'}$. En effet, d'après
Compléments d'algèbre, lemme \ref{more-algebra-lemma-diagram-localize}, on a
$$
(B')_{g'} = B_g \times_{A_h} A'_{h'}
$$
où $h \in A$, $h' \in A'$ sont les images de $g'$. Puisque
$B_g$ (resp.\ $A_h$, resp.\ $A'_{h'}$) est égal à $\mathcal{O}_Y(D(g))$
(resp.\ $f_*\mathcal{O}_X(D(g))$, resp.\ $f'_*\mathcal{O}_{X'}(D(g))$),
l'assertion s'ensuit.

\medskip\noindent
Il reste à montrer que $Y'$ est la somme amalgamée.
La discussion ci-dessus montre que le schéma $Y'$
possède un recouvrement ouvert affine $Y' = \bigcup W'_i$
tel que les ouverts correspondants
$U'_i \subset X'$, $W_i \subset Y$ et
$U_i \subset X$ soient affines.
De plus, si $A'_i$, $B_i$, $A_i$ sont les anneaux correspondant à
$U'_i$, $W_i$, $U_i$, alors
$W'_i$ correspond à $B_i \times_{A_i} A'_i$.
On peut donc appliquer les lemmes \ref{lemma-basic-example-pushout} et
\ref{lemma-pushout-fpqc-local} pour conclure que notre construction est une somme amalgamée
dans la catégorie des schémas.
\end{proof}

\noindent
Dans le lemme suivant, on utilise le produit fibré de catégories
défini dans
Catégories, exemple \ref{categories-example-2-fibre-product-categories}.

\begin{lemma}
\label{lemma-equivalence-categories-schemes-over-pushout}
Soit $X \to X'$ un épaississement de schémas et soit $X \to Y$ un
morphisme affine de schémas. Soit $Y' = Y \amalg_X X'$ la somme amalgamée
(voir le lemme \ref{lemma-pushout-along-thickening}). Le changement de base donne
un foncteur
$$
F : (\Sch/Y') \longrightarrow (\Sch/Y) \times_{(\Sch/X)} (\Sch/X')
$$
défini par $V' \longmapsto (V' \times_{Y'} Y, V' \times_{Y'} X', 1)$,
qui admet un adjoint à gauche
$$
G : (\Sch/Y) \times_{(\Sch/X)} (\Sch/X') \longrightarrow (\Sch/Y')
$$
envoyant le triplet $(V, U', \varphi)$ sur la somme amalgamée
$V \amalg_{(V \times_Y X)} U'$. Enfin, $F \circ G$ est isomorphe au
foncteur identité.
\end{lemma}

\begin{proof}
Soit $(V, U', \varphi)$ un objet de la catégorie produit fibré.
Posons $U = U' \times_{X'} X$. Remarquons que $U \to U'$ est un épaississement.
Puisque $\varphi : V \times_Y X \to U' \times_{X'} X = U$ est un isomorphisme,
on dispose d'un morphisme $U \to V$ au-dessus de $X \to Y$ qui identifie $U$ au
produit fibré $X \times_Y V$. En particulier, $U \to V$ est affine, voir
Morphismes, lemme \ref{morphisms-lemma-base-change-affine}.
On peut donc appliquer le lemme \ref{lemma-pushout-along-thickening}
pour obtenir une somme amalgamée $V' = V \amalg_U U'$. Notons $V' \to Y'$ le morphisme
fourni par la propriété de somme amalgamée de $V'$ et par
les morphismes donnés $V \to Y$ et $U' \to X'$, qui coïncident sur $U$
comme morphismes à valeurs dans $Y'$. En posant $G(V, U', \varphi) = V'$,
on obtient le foncteur $G$.

\medskip\noindent
Prouvons que $G$ est adjoint à gauche de $F$. Soit $Z$ un schéma
sur $Y'$. Il faut montrer que
$$
\Mor(V', Z) = \Mor((V, U', \varphi), F(Z))
$$
où les ensembles de morphismes sont pris dans leurs catégories respectives.
Soit $g' : V' \to Z$ un morphisme. Notons $\tilde g$ (resp.\ $\tilde f'$)
le composé de $g'$ avec le morphisme $V \to V'$ (resp.\ $U' \to V'$).
Effectuons le changement de base de $\tilde g$ (resp.\ $\tilde f'$) par $Y \to Y'$ (resp.\ $X' \to Y'$)
pour obtenir un morphisme $g : V \to Z \times_{Y'} Y$
(resp.\ $f' : U' \to Z \times_{Y'} X'$). Alors $(g, f')$ est un élément
du membre de droite de l'égalité ci-dessus (détails omis).
Réciproquement, supposons que $(g, f') : (V, U', \varphi) \to F(Z)$ soit un
élément du membre de droite.
On peut considérer le composé $\tilde g : V \to Z$
(resp.\ $\tilde f' : U' \to Z$) de $g$ (resp.\ $f$) avec
$Z \times_{Y'} X' \to Z$ (resp.\ $Z \times_{Y'} Y \to Z$).
Alors $\tilde g$ et $\tilde f'$ coïncident comme morphismes de $U$ dans $Z$.
Par la propriété universelle de la somme amalgamée, on obtient un morphisme
$g' : V' \to Z$, c'est-à-dire un élément du membre de gauche.
On omet de vérifier que ces constructions sont inverses l'une de l'autre.

\medskip\noindent
Pour prouver que $F \circ G$ est isomorphe à l'identité, il faut
montrer que le morphisme d'adjonction
$(V, U', \varphi) \to F(G(V, U', \varphi))$ est un isomorphisme.
On peut pour cela travailler localement sur des ouverts affines. Écrivons $X = \Spec(A)$, $X' = \Spec(A')$
et $Y = \Spec(B)$. Alors $A' \to A$ et $B \to A$ sont des homomorphismes d'anneaux comme dans
Compléments d'algèbre, lemme \ref{more-algebra-lemma-module-over-fibre-product},
et $Y' = \Spec(B')$, avec $B' = B \times_A A'$. Supposons ensuite que
$V = \Spec(D)$, $U' = \Spec(C')$ et que $\varphi$ soit donné par un
isomorphisme de $A$-algèbres $D \otimes_B A \to C' \otimes_{A'} A = C'/IC'$.
Posons $D' = D \times_{C'/IC'} C'$. Dans ce cas, il faut
prouver que $D' \otimes_{B'} B \cong D$ et $D' \otimes_{B'} A' \cong C'$.
C'est un cas particulier de Compléments d'algèbre, lemme
\ref{more-algebra-lemma-module-over-fibre-product}.
\end{proof}

\begin{lemma}
\label{lemma-scheme-over-pushout-flat-modules}
Soit $X \to X'$ un épaississement de schémas et soit $X \to Y$ un
morphisme affine de schémas. Soit $Y' = Y \amalg_X X'$ la somme amalgamée
(voir le lemme \ref{lemma-pushout-along-thickening}). Soit $V' \to Y'$
un morphisme de schémas. Posons
$V = Y \times_{Y'} V'$, $U' = X' \times_{Y'} V'$ et $U = X \times_{Y'} V'$.
Il existe une équivalence de catégories entre
\begin{enumerate}
\item les $\mathcal{O}_{V'}$-modules quasi-cohérents plats sur $Y'$, et
\item la catégorie des triplets $(\mathcal{G}, \mathcal{F}', \varphi)$, où
\begin{enumerate}
\item $\mathcal{G}$ est un $\mathcal{O}_V$-module quasi-cohérent plat sur $Y$,
\item $\mathcal{F}'$ est un $\mathcal{O}_{U'}$-module quasi-cohérent plat
sur $X'$, et
\item $\varphi : (U \to V)^*\mathcal{G} \to (U \to U')^*\mathcal{F}'$
est un isomorphisme de $\mathcal{O}_U$-modules.
\end{enumerate}
\end{enumerate}
L'équivalence envoie $\mathcal{G}'$ sur
$((V \to V')^*\mathcal{G}', (U' \to V')^*\mathcal{G}', can)$.
Supposons que $\mathcal{G}'$ corresponde au triplet
$(\mathcal{G}, \mathcal{F}', \varphi)$. Alors
\begin{enumerate}
\item[(a)] $\mathcal{G}'$ est un $\mathcal{O}_{V'}$-module de type fini si et
seulement si $\mathcal{G}$ et $\mathcal{F}'$ sont des
$\mathcal{O}_Y$- et $\mathcal{O}_{U'}$-modules de type fini.
\item[(b)] si $V' \to Y'$ est localement de présentation finie, alors
$\mathcal{G}'$ est un $\mathcal{O}_{V'}$-module de
présentation finie si et seulement si $\mathcal{G}$ et $\mathcal{F}'$ sont des
$\mathcal{O}_Y$- et $\mathcal{O}_{U'}$-modules de présentation finie.
\end{enumerate}
\end{lemma}

\begin{proof}
Un foncteur quasi-inverse associe au triplet
$(\mathcal{G}, \mathcal{F}', \varphi)$ le produit fibré
$$
(V \to V')_*\mathcal{G}
\times_{(U \to V')_*\mathcal{F}}
(U' \to V')_*\mathcal{F}'
$$
où $\mathcal{F} = (U \to U')^*\mathcal{F}'$. Cette construction convient, car,
sur les ouverts affines, on retrouve l'équivalence de Compléments d'algèbre, lemme
\ref{more-algebra-lemma-relative-flat-module-over-fibre-product}.
Certains détails sont omis.

\medskip\noindent
Les points (a) et (b) résultent de
Compléments d'algèbre, lemmes
\ref{more-algebra-lemma-relative-finite-module-over-fibre-product} et
\ref{more-algebra-lemma-relative-finitely-presented-module-over-fibre-product}.
\end{proof}

\begin{lemma}
\label{lemma-equivalence-categories-schemes-over-pushout-flat}
Dans la situation du
lemme \ref{lemma-equivalence-categories-schemes-over-pushout},
si $V' = G(V, U', \varphi)$ pour un certain triplet $(V, U', \varphi)$, alors
\begin{enumerate}
\item $V' \to Y'$ est localement de type fini si et seulement si $V \to Y$ et
$U' \to X'$ sont localement de type fini,
\item $V' \to Y'$ est plat si et seulement si $V \to Y$ et $U' \to X'$ sont plats,
\item $V' \to Y'$ est plat et localement de présentation finie si et seulement si
$V \to Y$ et $U' \to X'$ sont plats et localement de présentation finie,
\item $V' \to Y'$ est lisse si et seulement si $V \to Y$ et $U' \to X'$ sont lisses,
\item $V' \to Y'$ est étale si et seulement si $V \to Y$ et $U' \to X'$
sont étales, et
\item ajouter ici d'autres propriétés selon les besoins.
\end{enumerate}
Si $W'$ est plat sur $Y'$, alors le morphisme d'adjonction
$G(F(W')) \to W'$ est un isomorphisme. Ainsi, $F$ et $G$ définissent des foncteurs
quasi-inverses l'un de l'autre entre la catégorie des schémas plats sur $Y'$
et la catégorie des triplets $(V, U', \varphi)$ tels que $V \to Y$
et $U' \to X'$ soient plats.
\end{lemma}

\begin{proof}
En se plaçant sur des ouverts affines, les assertions de ce lemme
sont équivalentes aux assertions correspondantes de
Compléments d'algèbre, lemme
\ref{more-algebra-lemma-properties-algebras-over-fibre-product}.
\end{proof}


\section{Ouverture du lieu de platitude}
\label{section-open-flat}

\noindent
La démonstration de ce résultat demande un certain travail et mérite
(peut-être) une section à elle seule. La voici.

\begin{theorem}
\label{theorem-openness-flatness}
\begin{reference}
\cite[IV, théorème 11.3.1]{EGA}
\end{reference}
Soit $S$ un schéma.
Soit $f : X \to S$ un morphisme localement de présentation finie.
Soit $\mathcal{F}$ un $\mathcal{O}_X$-module quasi-cohérent
localement de présentation finie. Alors
$$
U = \{x \in X \mid \mathcal{F}\text{ est plat sur }S\text{ en }x\}
$$
est ouvert dans $X$.
\end{theorem}

\begin{proof}
L'ouverture se vérifie localement sur $X$ ; on peut donc supposer
que $f$ est un morphisme de schémas affines. Dans ce cas, le
théorème est exactement le
théorème \ref{algebra-theorem-openness-flatness} d'Algèbre.
\end{proof}

\begin{lemma}
\label{lemma-flat-locus-base-change}
Soit $S$ un schéma.
Soit
$$
\xymatrix{
X' \ar[r]_{g'} \ar[d]_{f'} & X \ar[d]^f \\
S' \ar[r]^g & S
}
$$
un diagramme cartésien de schémas.
Soit $\mathcal{F}$ un $\mathcal{O}_X$-module quasi-cohérent.
Soit $x' \in X'$ d'images
$x = g'(x')$ et $s' = f'(x')$.
\begin{enumerate}
\item Si $\mathcal{F}$ est plat sur $S$ en $x$, alors
$(g')^*\mathcal{F}$ est plat sur $S'$ en $x'$.
\item Si $g$ est plat en $s'$ et si $(g')^*\mathcal{F}$ est plat sur $S'$ en
$x'$, alors $\mathcal{F}$ est plat sur $S$ en $x$.
\end{enumerate}
En particulier, si $g$ est plat, si $f$ est localement de présentation finie
et si $\mathcal{F}$ est localement de présentation finie,
alors la formation de l'ouvert du
théorème \ref{theorem-openness-flatness}
commute au changement de base.
\end{lemma}

\begin{proof}
Considérons le diagramme commutatif d'anneaux locaux
$$
\xymatrix{
\mathcal{O}_{X', x'} & \mathcal{O}_{X, x} \ar[l] \\
\mathcal{O}_{S', s'} \ar[u] & \mathcal{O}_{S, s} \ar[l] \ar[u]
}
$$
Notons que $\mathcal{O}_{X', x'}$
est un localisé de
$\mathcal{O}_{X, x} \otimes_{\mathcal{O}_{S, s}} \mathcal{O}_{S', s'}$
et que $((g')^*\mathcal{F})_{x'}$ est égal à
$\mathcal{F}_x \otimes_{\mathcal{O}_{X, x}} \mathcal{O}_{X', x'}$.
Le lemme résulte donc du
lemme \ref{algebra-lemma-base-change-flat-up-down} d'Algèbre.
\end{proof}






\section{Critère de platitude par fibres}
\label{section-criterion-flat-fibres}

\noindent
Considérons un diagramme commutatif de schémas (diagramme de gauche)
$$
\xymatrix{
X \ar[rr]_f \ar[dr] & & Y \ar[dl] \\
& S
}
\quad
\xymatrix{
X_s \ar[rr]_{f_s} \ar[rd] & & Y_s \ar[dl] \\
& \Spec(\kappa(s))
}
$$
et un $\mathcal{O}_X$-module quasi-cohérent $\mathcal{F}$.
Étant donné un point $x \in X$ au-dessus de $s \in S$, d'image $y = f(x)$,
on se pose la question : $\mathcal{F}$ est-il plat
sur $Y$ en $x$ ? Si $\mathcal{F}$ est plat sur $S$ en $x$, le
théorème affirme que cette question est étroitement liée à celle
de savoir si la restriction de $\mathcal{F}$ à la fibre
$$
\mathcal{F}_s = (X_s \to X)^*\mathcal{F}
$$
est plate sur $Y_s$ en $x$. On trouvera ci-dessous une version
« noethérienne » et une version « de présentation finie » ; une version
« nilpotente » a déjà été traitée, voir le
lemme \ref{lemma-flatness-morphism-thickenings}.

\begin{theorem}
\label{theorem-criterion-flatness-fibre-Noetherian}
Soit $S$ un schéma.
Soit $f : X \to Y$ un morphisme de schémas sur $S$.
Soit $\mathcal{F}$ un $\mathcal{O}_X$-module quasi-cohérent.
Soit $x \in X$. Posons $y = f(x)$ et notons $s \in S$ l'image de $x$ dans $S$.
Supposons $S$, $X$, $Y$ localement noethériens,
$\mathcal{F}$ cohérent et $\mathcal{F}_x \not = 0$.
Alors les conditions suivantes sont équivalentes :
\begin{enumerate}
\item $\mathcal{F}$ est plat sur $S$ en $x$ et
$\mathcal{F}_s$ est plat sur $Y_s$ en $x$, et
\item $Y$ est plat sur $S$ en $y$ et $\mathcal{F}$ est
plat sur $Y$ en $x$.
\end{enumerate}
\end{theorem}

\begin{proof}
Considérons les homomorphismes d'anneaux
$$
\mathcal{O}_{S, s} \longrightarrow
\mathcal{O}_{Y, y} \longrightarrow \mathcal{O}_{X, x}
$$
et le module $\mathcal{F}_x$. La fibre de $\mathcal{F}_s$ en $x$
est le module $\mathcal{F}_x/\mathfrak m_s \mathcal{F}_x$ et
l'anneau local de $Y_s$ en $y$ est
$\mathcal{O}_{Y, y}/\mathfrak m_s \mathcal{O}_{Y, y}$.
L'implication (1) $\Rightarrow$ (2) est donc le
lemme \ref{algebra-lemma-criterion-flatness-fibre-Noetherian} d'Algèbre.
Si (2) est vérifiée, le premier homomorphisme d'anneaux est fidèlement plat
et $\mathcal{F}_x$ est plat sur $\mathcal{O}_{Y, y}$ ; le
lemme \ref{algebra-lemma-composition-flat} d'Algèbre
montre donc que $\mathcal{F}_x$ est plat sur $\mathcal{O}_{S, s}$.
De plus, $\mathcal{F}_x/\mathfrak m_s \mathcal{F}_x$ s'obtient
à partir du module plat $\mathcal{F}_x$ par le changement de base
$\mathcal{O}_{Y, y} \to \mathcal{O}_{Y, y}/\mathfrak m_s \mathcal{O}_{Y, y}$ ;
il est donc plat d'après le
lemme \ref{algebra-lemma-flat-base-change} d'Algèbre.
\end{proof}

\noindent
Voici la version non noethérienne.

\begin{theorem}
\label{theorem-criterion-flatness-fibre}
Soit $S$ un schéma.
Soit $f : X \to Y$ un morphisme de schémas sur $S$.
Soit $\mathcal{F}$ un $\mathcal{O}_X$-module quasi-cohérent.
Supposons que
\begin{enumerate}
\item $X$ soit localement de présentation finie sur $S$,
\item $\mathcal{F}$ soit un $\mathcal{O}_X$-module de présentation finie, et
\item $Y$ soit localement de type fini sur $S$.
\end{enumerate}
Soit $x \in X$. Posons $y = f(x)$ et soit $s \in S$ l'image de $x$ dans $S$.
Si $\mathcal{F}_x \not = 0$, les conditions suivantes sont équivalentes :
\begin{enumerate}
\item $\mathcal{F}$ est plat sur $S$ en $x$ et
$\mathcal{F}_s$ est plat sur $Y_s$ en $x$, et
\item $Y$ est plat sur $S$ en $y$ et $\mathcal{F}$ est
plat sur $Y$ en $x$.
\end{enumerate}
De plus, l'ensemble des points $x$ où (1) et (2) sont vérifiées est ouvert dans
$\text{Supp}(\mathcal{F})$.
\end{theorem}

\begin{proof}
Considérons les homomorphismes d'anneaux
$$
\mathcal{O}_{S, s} \longrightarrow
\mathcal{O}_{Y, y} \longrightarrow \mathcal{O}_{X, x}
$$
et le module $\mathcal{F}_x$. La fibre de $\mathcal{F}_s$ en $x$
est le module $\mathcal{F}_x/\mathfrak m_s \mathcal{F}_x$ et
l'anneau local de $Y_s$ en $y$ est
$\mathcal{O}_{Y, y}/\mathfrak m_s \mathcal{O}_{Y, y}$.
L'implication (1) $\Rightarrow$ (2) est donc le
lemme \ref{algebra-lemma-criterion-flatness-fibre-fp-over-ft} d'Algèbre.
Si (2) est vérifiée, le premier homomorphisme d'anneaux est fidèlement plat
et $\mathcal{F}_x$ est plat sur $\mathcal{O}_{Y, y}$ ; le
lemme \ref{algebra-lemma-composition-flat} d'Algèbre
montre donc que $\mathcal{F}_x$ est plat sur $\mathcal{O}_{S, s}$.
De plus, $\mathcal{F}_x/\mathfrak m_s \mathcal{F}_x$ s'obtient
à partir du module plat $\mathcal{F}_x$ par le changement de base
$\mathcal{O}_{Y, y} \to \mathcal{O}_{Y, y}/\mathfrak m_s \mathcal{O}_{Y, y}$ ;
il est donc plat d'après le
lemme \ref{algebra-lemma-flat-base-change} d'Algèbre.

\medskip\noindent
D'après
Morphismes, lemme \ref{morphisms-lemma-finite-presentation-permanence},
le morphisme $f$ est localement de présentation finie.
Considérons l'ensemble
\begin{equation}
\label{equation-open}
U = \{x \in X \mid \mathcal{F} \text{ plat en }x
\text{ à la fois sur }Y\text{ et }S\}.
\end{equation}
Cet ensemble est ouvert dans $X$ d'après le
théorème \ref{theorem-openness-flatness}.
Notons que si $x \in U$, alors $\mathcal{F}_s$ est plat en
$x$ sur $Y_s$, comme changement de base d'un module plat par le
morphisme $Y_s \to Y$, voir
Morphismes, lemme \ref{morphisms-lemma-base-change-module-flat}.
Ainsi, en tout point de $U \cap \text{Supp}(\mathcal{F})$,
la condition (1) est vérifiée. D'autre part, il est
clair que si $x \in \text{Supp}(\mathcal{F})$ vérifie
(1) et (2), alors $x \in U$. L'ouvert recherché est donc
$U \cap \text{Supp}(\mathcal{F})$.
\end{proof}

\noindent
On utilise souvent ces théorèmes sous les formes simplifiées suivantes.
Nous ne donnons que les énoncés globaux ; il existe bien entendu aussi des
versions ponctuelles.

\begin{lemma}
\label{lemma-morphism-between-flat-Noetherian}
Soit $S$ un schéma.
Soit $f : X \to Y$ un morphisme de schémas sur $S$.
Supposons que
\begin{enumerate}
\item $S$, $X$, $Y$ soient localement noethériens,
\item $X$ soit plat sur $S$,
\item pour tout $s \in S$, le morphisme
$f_s : X_s \to Y_s$ soit plat.
\end{enumerate}
Alors $f$ est plat. Si $f$ est en outre surjectif, $Y$ est plat sur $S$.
\end{lemma}

\begin{proof}
C'est un cas particulier du
théorème \ref{theorem-criterion-flatness-fibre-Noetherian}.
\end{proof}

\begin{lemma}
\label{lemma-morphism-between-flat}
Soit $S$ un schéma.
Soit $f : X \to Y$ un morphisme de schémas sur $S$.
Supposons que
\begin{enumerate}
\item $X$ soit localement de présentation finie sur $S$,
\item $X$ soit plat sur $S$,
\item pour tout $s \in S$, le morphisme
$f_s : X_s \to Y_s$ soit plat, et
\item $Y$ soit localement de type fini sur $S$.
\end{enumerate}
Alors $f$ est plat. Si $f$ est en outre surjectif, $Y$ est plat sur $S$.
\end{lemma}

\begin{proof}
C'est un cas particulier du
théorème \ref{theorem-criterion-flatness-fibre}.
\end{proof}

\begin{lemma}
\label{lemma-base-change-criterion-flatness-fibre}
Soit $S$ un schéma. Soit $f : X \to Y$ un morphisme de schémas sur $S$.
Soit $\mathcal{F}$ un $\mathcal{O}_X$-module quasi-cohérent.
Supposons que
\begin{enumerate}
\item $X$ soit localement de présentation finie sur $S$,
\item $\mathcal{F}$ soit un $\mathcal{O}_X$-module de présentation finie,
\item $\mathcal{F}$ soit plat sur $S$, et
\item $Y$ soit localement de type fini sur $S$.
\end{enumerate}
Alors l'ensemble
$$
U = \{x \in X \mid \mathcal{F} \text{ plat en }x \text{ sur }Y\}.
$$
est ouvert dans $X$ et sa formation commute à tout changement de base :
si $S' \to S$ est un morphisme de schémas et si $U'$ est l'ensemble des points
de $X' = X \times_S S'$ où $\mathcal{F}' = \mathcal{F} \times_S S'$
est plat sur $Y' = Y \times_S S'$, alors $U' = U \times_S S'$.
\end{lemma}

\begin{proof}
D'après
Morphismes, lemme \ref{morphisms-lemma-finite-presentation-permanence},
le morphisme $f$ est localement de présentation finie.
Ainsi, $U$ est ouvert d'après le
théorème \ref{theorem-openness-flatness}.
Puisque l'on a supposé $\mathcal{F}$ plat sur $S$, le
théorème \ref{theorem-criterion-flatness-fibre}
donne
$$
U = \{x \in X \mid \mathcal{F}_s \text{ plat en }x \text{ sur }Y_s\}.
$$
où $s$ désigne toujours l'image de $x$ dans $S$. (Cette description
reste trivialement valable lorsque $\mathcal{F}_x = 0$.) De plus, les hypothèses
du lemme restent vérifiées pour le morphisme $f' : X' \to Y'$
et le faisceau $\mathcal{F}'$. Ainsi, $U'$ admet une description analogue.
Autrement dit, il suffit de prouver que, pour
$s' \in S'$ d'image $s \in S$, l'ensemble
$$
\{x' \in X'_{s'} \mid
\mathcal{F}'_{s'} \text{ plat en }x' \text{ sur }Y'_{s'}\}
$$
est l'image réciproque du lieu correspondant dans $X_s$.
C'est vrai d'après le
lemme \ref{lemma-flat-locus-base-change},
car, dans le diagramme cartésien
$$
\xymatrix{
X'_{s'} \ar[d] \ar[r] & X_s \ar[d] \\
Y'_{s'} \ar[r] & Y_s
}
$$
les morphismes horizontaux sont plats, puisqu'ils s'obtiennent par changement de
base par le morphisme plat $\Spec(\kappa(s')) \to \Spec(\kappa(s))$.
\end{proof}

\begin{lemma}
\label{lemma-base-change-flatness-fibres}
Soit $S$ un schéma. Soit $f : X \to Y$ un morphisme de schémas sur $S$.
Supposons que
\begin{enumerate}
\item $X$ soit localement de présentation finie sur $S$,
\item $X$ soit plat sur $S$, et
\item $Y$ soit localement de type fini sur $S$.
\end{enumerate}
Alors l'ensemble
$$
U = \{x \in X \mid X\text{ plat en }x \text{ sur }Y\}.
$$
est ouvert dans $X$ et sa formation commute à tout changement de base.
\end{lemma}

\begin{proof}
C'est un cas particulier du
lemme \ref{lemma-base-change-criterion-flatness-fibre}.
\end{proof}

\noindent
Le lemme suivant est une variante du
lemme \ref{algebra-lemma-free-fibre-flat-free} d'Algèbre.
Notons que l'hypothèse
que $(\mathcal{F}_s)_x$ est un $\mathcal{O}_{X_s, x}$-module plat
signifie que $(\mathcal{F}_s)_x$ est un $\mathcal{O}_{X_s, x}$-module libre,
ce qui est toujours le cas si $x \in X_s$ est un point générique d'une
composante irréductible de $X_s$ et si $X_s$ est réduit (en effet,
dans ce cas $\mathcal{O}_{X_s, x}$ est un corps, voir
Algèbre, lemme \ref{algebra-lemma-minimal-prime-reduced-ring}).

\begin{lemma}
\label{lemma-flat-and-free-at-point-fibre}
Soit $f : X \to S$ un morphisme de schémas de présentation finie.
Soit $\mathcal{F}$ un $\mathcal{O}_X$-module de présentation finie.
Soit $x \in X$ d'image $s \in S$.
Si $\mathcal{F}$ est plat en $x$ sur $S$ et si $(\mathcal{F}_s)_x$ est un
$\mathcal{O}_{X_s, x}$-module plat, alors $\mathcal{F}$
est libre de rang fini dans un voisinage de $x$.
\end{lemma}

\begin{proof}
Si $\mathcal{F}_x \otimes \kappa(x)$ est nul, alors $\mathcal{F}_x = 0$
d'après le lemme de Nakayama (Algèbre, lemme \ref{algebra-lemma-NAK}), donc
$\mathcal{F}$ est nul dans un voisinage de $x$
(Modules, lemme \ref{modules-lemma-finite-type-stalk-zero})
et le lemme est établi. On peut donc supposer $\mathcal{F}_x \otimes \kappa(x)$
non nul, et l'on voit que le
théorème \ref{theorem-criterion-flatness-fibre}
s'applique avec $f = \text{id} : X \to X$. On en conclut que $\mathcal{F}_x$
est plat sur $\mathcal{O}_{X, x}$. Ainsi, $\mathcal{F}_x$ est libre, voir
par exemple Algèbre, lemme \ref{algebra-lemma-finite-flat-local}.
Choisissons un voisinage ouvert $x \in U \subset X$ et des sections
$s_1, \ldots, s_r \in \mathcal{F}(U)$ dont les images forment une base de
$\mathcal{F}_x$. Le morphisme correspondant
$\psi : \mathcal{O}_U^{\oplus r} \to \mathcal{F}|_U$ est surjectif après
restriction de $U$ (Modules, lemme \ref{modules-lemma-finite-type-stalk-zero}).
Alors $\Ker(\psi)$ est de type fini (voir Modules, lemme
\ref{modules-lemma-kernel-surjection-finite-free-onto-finite-presentation})
et $\Ker(\psi)_x = 0$. En restreignant encore $U$,
$\psi$ devient donc un isomorphisme.
\end{proof}

\begin{lemma}
\label{lemma-finite-free-open}
Soit $f : X \to S$ un morphisme de schémas
localement de présentation finie.
Soit $\mathcal{F}$ un $\mathcal{O}_X$-module de présentation finie
plat sur $S$. Alors l'ensemble
$$
\{x \in X : \mathcal{F}\text{ libre dans un voisinage de }x\}
$$
est ouvert dans $X$ et sa formation commute à tout changement de base
$S' \to S$.
\end{lemma}

\begin{proof}
L'ouverture est immédiate. Soit $x \in X$ d'image $s \in S$.
D'après le lemme \ref{lemma-flat-and-free-at-point-fibre},
$x$ appartient à notre ensemble si et seulement si
$\mathcal{F}|_{X_s}$ est plat en $x$ sur $X_s$.
Ceci équivaut aussi à la platitude de $\mathcal{F}$
en $x$ sur $X$ (car cette propriété est
impliquée par la liberté de $\mathcal{F}_x$ et implique
la platitude de $\mathcal{F}|_{X_s}$ en $x$ sur $X_s$).
L'assertion sur le changement de base résulte donc du
lemme \ref{lemma-base-change-criterion-flatness-fibre}
appliqué à $\text{id} : X \to X$ sur $S$.
\end{proof}
\section{Immersions fermées entre schémas lisses}
\label{section-closed-immersions-smooth}

\noindent
Quelques résultats qui ne trouvent guère leur place ailleurs.

\begin{lemma}
\label{lemma-etale-local-structure}
Soit $S$ un schéma. Soit $Y \to X$ une immersion fermée de schémas
lisses sur $S$. Pour tout $y \in Y$, il existe des entiers
$0 \leq m, n$ et un diagramme commutatif
$$
\xymatrix{
Y \ar[d] &
V \ar[l] \ar[d] \ar[r] &
\mathbf{A}^m_S
\ar[d]^{(a_1, \ldots, a_m) \mapsto (a_1, \ldots, a_m, 0 \ldots, 0)} \\
X &
U \ar[l] \ar[r]^-\pi &
\mathbf{A}^{m + n}_S
}
$$
où $U \subset X$ est ouvert, $V = Y \cap U$,
$\pi$ est étale, $V = \pi^{-1}(\mathbf{A}^m_S)$ et $y \in V$.
\end{lemma}

\begin{proof}
La question est locale sur $X$ ; on peut donc remplacer $X$ par
un voisinage ouvert de $y$. Puisque $Y \to X$ est une immersion régulière
d'après Diviseurs, lemme
\ref{divisors-lemma-immersion-smooth-into-smooth-regular-immersion},
on peut supposer $X = \Spec(A)$ affine et qu'il existe une suite régulière
$f_1, \ldots, f_n \in A$ telle que $Y = V(f_1, \ldots, f_n)$.
En restreignant encore $X$ (et donc $Y$),
on peut supposer qu'il existe un morphisme étale $Y \to \mathbf{A}^m_S$, voir
Morphismes, lemme \ref{morphisms-lemma-smooth-etale-over-affine-space}.
Soient $\overline{g}_1, \ldots, \overline{g}_m$ dans $\mathcal{O}_Y(Y)$
les fonctions coordonnées de ce morphisme étale.
Choisissons des relèvements $g_1, \ldots, g_m \in A$ de ces fonctions
et considérons le morphisme
$$
(g_1, \ldots, g_m, f_1, \ldots, f_n) :
X
\longrightarrow
\mathbf{A}^{m + n}_S
$$
sur $S$. C'est un morphisme entre schémas localement de présentation finie
sur $S$ ; il est donc localement de présentation finie
(Morphismes, lemme \ref{morphisms-lemma-finite-presentation-permanence}).
La restriction de ce morphisme à
$\mathbf{A}^m_S \subset \mathbf{A}^{m + n}_S$
est étale par construction. Par conséquent, pour montrer que
$X \to \mathbf{A}^{m + n}_S$ est étale en $y$,
il suffit de montrer que $X \to \mathbf{A}^{m + n}_S$ est plat en $y$,
voir Morphismes, lemme \ref{morphisms-lemma-etale-at-point}.
Soit $s \in S$ l'image de $y$. Il suffit de vérifier que
$X_s \to \mathbf{A}^{m + n}_s$ est plat en $y$, voir le
théorème \ref{theorem-criterion-flatness-fibre}.
Soit $z \in \mathbf{A}^{m + n}_s$ l'image de $y$.
L'homomorphisme local d'anneaux
$$
\mathcal{O}_{\mathbf{A}^{m + n}_s, z}
\longrightarrow
\mathcal{O}_{X_s, y}
$$
est plat d'après Algèbre, lemme \ref{algebra-lemma-CM-over-regular-flat}.
En effet, les schémas lisses sur un corps sont réguliers et les anneaux réguliers
sont de Cohen-Macaulay, voir Variétés, lemme \ref{varieties-lemma-smooth-regular}
et Algèbre, lemme \ref{algebra-lemma-regular-ring-CM}.
La source et le but sont donc des anneaux locaux réguliers (et donc de Cohen-Macaulay).
Ils ont la même dimension : en effet, on a
$\dim(\mathcal{O}_{Y_s, y}) = \dim(\mathcal{O}_{\mathbf{A}^m_s, z})$
d'après Compléments d'algèbre, lemme \ref{more-algebra-lemma-dimension-etale-extension},
on a $\dim(\mathcal{O}_{\mathbf{A}^{m + n}_s, z}) = n +
\dim(\mathcal{O}_{\mathbf{A}^m_s, z})$, et l'on a
$\dim(\mathcal{O}_{X_s, y}) = n + \dim(\mathcal{O}_{Y_s, y})$,
car $\mathcal{O}_{Y_s, y}$ est le quotient de
$\mathcal{O}_{X_s, y}$ par la suite régulière $f_1, \ldots, f_n$
de longueur $n$ (voir
Diviseurs, remarque \ref{divisors-remark-relative-regular-immersion-elements}).
Enfin, l'anneau de la fibre de la flèche ci-dessus est fini sur $\kappa(z)$,
puisque $Y_s \to \mathbf{A}^m_s$ est étale en $y$.
Ceci achève la démonstration.
\end{proof}

\begin{remark}
\label{remark-relative-blowup}
Fixons un anneau $R$ et posons $S = \Spec(R)$. Fixons des entiers $0 \leq m$ et
$1 \leq n$. Considérons l'immersion fermée
$$
Z = \mathbf{A}^m_S \longrightarrow \mathbf{A}^{m + n}_S = X,\quad
(a_1, \ldots, a_m) \mapsto (a_1, \ldots, a_m, 0, \ldots 0).
$$
Nous allons considérer l'éclatement $X'$ de $X$
le long du sous-schéma fermé $Z$. Écrivons
$$
X = \Spec(A)
\quad\text{avec}\quad
A = R[x_1, \ldots, x_m, y_1, \ldots, y_n]
$$
Alors $X'$ est le Proj de l'algèbre de Rees de $A$ relativement à
l'idéal $(y_1, \ldots, y_n)$. Cette algèbre de Rees est égale à
$B = A[T_1, \ldots, T_n]/(y_iT_j - y_jT_i)$ ; les détails sont omis.
Ainsi, $X' = \text{Proj}(B)$ est lisse sur $S$, car il est
recouvert par les ouverts affines
\begin{align*}
D_+(T_i)
& =
\Spec(B_{(T_i)}) \\
& =
\Spec(A[t_1, \ldots, \hat t_i, \ldots t_n]/(y_j - y_i t_j)) \\
& =
\Spec(R[x_1, \ldots, x_m, y_i, t_1, \ldots, \hat t_i, \ldots, t_n])
\end{align*}
qui sont isomorphes à $\mathbf{A}^{n + m}_S$.
Dans cette carte, le diviseur exceptionnel est défini par
$y_i = 0$ ; il est donc lui aussi lisse
sur $S$.
\end{remark}

\begin{lemma}
\label{lemma-blowup}
Soit $S$ un schéma. Soit $Z \to X$ une immersion fermée de schémas
lisses sur $S$. Soit $b : X' \to X$ l'éclatement de $Z$,
de diviseur exceptionnel $E \subset X'$. Alors $X'$ et $E$ sont lisses
sur $S$. Le morphisme $p : E \to Z$ est canoniquement isomorphe
au fibré projectif
$$
\mathbf{P}(\mathcal{I}/\mathcal{I}^2) \longrightarrow Z
$$
où $\mathcal{I} \subset \mathcal{O}_X$ est le faisceau d'idéaux
de $Z$. Le faisceau relatif $\mathcal{O}_E(1)$ provenant de la structure
de fibré projectif est isomorphe à la restriction de
$\mathcal{O}_{X'}(-E)$ à $E$.
\end{lemma}

\begin{proof}
D'après Diviseurs, lemme
\ref{divisors-lemma-immersion-smooth-into-smooth-regular-immersion},
l'immersion $Z \to X$ est régulière ; le faisceau d'idéaux
$\mathcal{I}$ est donc de type fini, et $b$ est un morphisme projectif
admettant pour faisceau inversible relativement ample
$\mathcal{O}_{X'}(1) = \mathcal{O}_{X'}(-E)$, voir
Diviseurs, lemmes
\ref{divisors-lemma-blowing-up-gives-effective-Cartier-divisor} et
\ref{divisors-lemma-blowing-up-projective}.
Le morphisme canonique $\mathcal{I} \to b_*\mathcal{O}_{X'}(1)$
donne une immersion fermée
$$
X' \longrightarrow
\mathbf{P}\left(\bigoplus\nolimits_{n \geq 0}
\text{Sym}^n_{\mathcal{O}_X}(\mathcal{I})\right)
$$
par la construction même de l'éclatement. La restriction de ce morphisme
à $E$ donne un morphisme canonique
$$
E \longrightarrow
\mathbf{P}\left(\bigoplus\nolimits_{n \geq 0}
\text{Sym}^n_{\mathcal{O}_Z}(\mathcal{I}/\mathcal{I}^2)\right)
$$
sur $Z$. Puisque $\mathcal{I}/\mathcal{I}^2$ est localement libre de rang fini,
si ce morphisme canonique est un isomorphisme, la dernière assertion du
lemme est vérifiée. Après ces observations, la question est locale
sur $X$ pour la topologie étale. En effet, l'éclatement commute au changement
de base plat d'après Diviseurs, lemme \ref{divisors-lemma-flat-base-change-blowing-up},
et la lissité peut se vérifier après précomposition par un morphisme
étale surjectif. Ainsi, la structure locale étale d'une
immersion fermée de schémas sur $S$ donnée au lemme
\ref{lemma-etale-local-structure} nous ramène au
cas traité dans la remarque \ref{remark-relative-blowup}.
\end{proof}
\section{Modules plats et assassins relatifs}
\label{section-flat-relative-assassin}

\noindent
Dans cette section, nous montrerons que, dans les situations géométriques,
le support d'un module plat est, en un certain sens, équidimensionnel sur la base.
Pour le cas noethérien, nous renvoyons à
\cite[IV Proposition 12.1.1.5]{EGA}.
Commençons par deux lemmes auxiliaires.

\begin{lemma}
\label{lemma-mod-injective-valuation-ring}
Soit $A$ un anneau de valuation. Soit $A \to B$ un homomorphisme local
d'anneaux locaux essentiellement de type fini.
Soit $u : N \to M$ un homomorphisme de $B$-modules de type fini.
Supposons $M$ plat sur $A$ et
$\overline{u} : N/\mathfrak m_A N \to M/\mathfrak m_A M$ injectif.
Alors $u$ est injectif et $M/u(N)$ est plat sur $A$.
\end{lemma}

\begin{proof}
Nous déduirons ce lemme du
lemme \ref{algebra-lemma-mod-injective-general} d'Algèbre
(notons que les rôles de $M$ et $N$ ont été intervertis).
Pour cela, nous utiliserons Compléments d'algèbre, lemme
\ref{more-algebra-lemma-valuation-ring-flat-essentially-finite-type},
afin de nous ramener au cas de présentation finie.

\medskip\noindent
Par hypothèse, on peut écrire $B$ comme quotient du localisé d'une
algèbre de polynômes $P = A[x_1, \ldots, x_n]$ en un idéal premier
$\mathfrak q$. On peut alors considérer $u : N \to M$ comme un homomorphisme
de $P_\mathfrak q$-modules de type fini. On peut donc supposer, et l'on suppose,
que $B$ est essentiellement de présentation finie sur $A$.

\medskip\noindent
Le $B$-module $N$ est de type fini, mais pas nécessairement de présentation finie.
Écrivons $N = \colim N_\lambda$ comme limite inductive filtrante de
$B$-modules de présentation finie, avec des morphismes de transition surjectifs.
Par exemple, choisissons une présentation
$0 \to K \to B^{\oplus r} \to N \to 0$,
écrivons $K$ comme réunion de ses sous-modules de type fini $K_\lambda$ et
posons $N_\lambda = \Coker(K_\lambda \to B^{\oplus r})$.
Le module $N/\mathfrak m_A N$ est de présentation finie
sur l'anneau noethérien $B/\mathfrak m_A B$.
Par conséquent, pour $\lambda$ assez grand, on a
$N_\lambda/\mathfrak m_A N_\lambda = N/\mathfrak m_A N$.
Si l'on peut démontrer le lemme pour le composé
$u_\lambda : N_\lambda \to M$, on en déduit que
$N_\lambda = N$ et que le résultat vaut pour $u$.
On peut donc supposer, et l'on suppose, que $N$ est de présentation finie sur $B$.

\medskip\noindent
D'après Compléments d'algèbre, lemme
\ref{more-algebra-lemma-valuation-ring-flat-essentially-finite-type},
le module $M$ est de présentation finie sur $B$.
Toutes les hypothèses du
lemme \ref{algebra-lemma-mod-injective-general} d'Algèbre
sont donc vérifiées, ce qui permet de conclure.
\end{proof}

\begin{lemma}
\label{lemma-make-base-change}
\begin{reference}
Ce résultat figure dans la démonstration de
\cite[IV Proposition 12.1.1.5]{EGA}
\end{reference}
Soit $f : X \to S$ un morphisme de schémas. Soit $y \in X$ un point
d'image $t \in S$. Notons $Y \subset X$ l'adhérence de $\{y\}$,
munie de sa structure de sous-schéma fermé intègre de $X$. Soit $s \in S$ et soit
$x \in Y_s$ un point générique d'une composante irréductible de $Y_s$.
Il existe un diagramme cartésien
$$
\xymatrix{
X' \ar[r]_{g'} \ar[d]_{f'} & X \ar[d]^f \\
S' \ar[r]^g & S
}
$$
ayant les propriétés suivantes :
\begin{enumerate}
\item $S'$ est le spectre d'un anneau de valuation
de point générique $t'$ et de point fermé $s'$,
\item $g(t') = t$ et $g(s') = s$,
\item il existe un point $y' \in X'_{t'}$ qui est
un point générique d'une composante irréductible de
$(S' \times_S Y)_{t'} = Y_t \times_t t'$
et vérifie $g'(y') = y$,
\item si l'on note $Y' \subset X'$ l'adhérence de $\{y'\}$,
munie de sa structure de sous-schéma fermé intègre de $X'$,
il existe un point $x' \in Y'_{s'}$ qui est un point générique
d'une composante irréductible de $Y'_{s'}$
et vérifie $g'(x') = x$.
\end{enumerate}
\end{lemma}

\begin{proof}
Choisissons un anneau de valuation $R$, posons $S' = \Spec(R)$, de
point générique $t'$ et de point fermé $s'$, et choisissons un morphisme
$h : S' \to X$ tel que $h(t') = y$ et $h(s') = x$.
Voir Schémas, lemme \ref{schemes-lemma-points-specialize}.
Posons $g = f \circ h$, de sorte que $g(t') = t$ et $g(s') = s$.
Considérons le changement de base
$$
\xymatrix{
X' \ar[r]_{g'} \ar[d] & X \ar[d] \\
S' \ar@/^1em/[u]^\sigma \ar[r]^-g & S
}
$$
On obtient une section $\sigma$ du morphisme changé de base telle que
$h = g' \circ \sigma$.

\medskip\noindent
Bien entendu, $\sigma$ se factorise par le changement de base $S' \times_S Y$ de $Y$,
puisque $h$ se factorise par $Y$. Soit $y' \in X'_{t'} \subset X'$ le
point générique d'une composante irréductible de la fibre
$$
(S' \times_S Y)_{t'} = Y_t \times_t t'
$$
contenant le point $\sigma(t')$, c'est-à-dire telle que
$y' \leadsto \sigma(t')$.
Puisque $g'(y') \in Y_t$ et $g(y') \leadsto g(\sigma(t')) = y$,
on obtient $g'(y') = y$, car $y$ est le point générique de
la fibre $Y_t$.
Notons $Y' \subset X'$ l'adhérence de $\{y'\}$ dans $X'$,
munie de sa structure de sous-schéma fermé intègre.
Alors $\sigma$ se factorise par $Y'$, puisque $\sigma(t') \in Y'$.
Choisissons un point générique $x' \in Y'_{s'}$ d'une composante
irréductible de $Y'_{s'}$ contenant $\sigma(s')$ ; on
a donc $x' \leadsto \sigma(s')$, puis $g'(x') \leadsto g'(\sigma(s')) = x$.
À nouveau, puisque $x$ est par hypothèse un point générique d'une
composante irréductible de $Y_s$
et puisque $g'(Y') \subset Y$, on en conclut que $g'(x') = x$.
\end{proof}

\begin{lemma}
\label{lemma-associated-point-specializes}
Soit $f : X \to S$ un morphisme de schémas localement de type fini.
Soit $\mathcal{F}$ un $\mathcal{O}_X$-module quasi-cohérent de type fini.
Soit $y \in \text{Ass}_{X/S}(\mathcal{F})$ d'image $t \in S$.
Notons $Y \subset X$ l'adhérence de $\{y\}$ dans $X$,
munie de sa structure de sous-schéma fermé intègre. Soit $s \in S$
et soit $x \in Y_s$ un point générique d'une composante irréductible
de $Y_s$. Si $\mathcal{F}$ est plat sur $S$ en $x$, alors
$x \in \text{Ass}_{X/S}(\mathcal{F})$ et
$\dim_x(Y_s) = \dim(Y_t)$.
\end{lemma}

\begin{proof}
Choisissons un diagramme comme au lemme \ref{lemma-make-base-change}.
Posons $\mathcal{F}' = (g')^*\mathcal{F}$.
Diviseurs, lemme \ref{divisors-lemma-base-change-relative-assassin},
implique que $y' \in \text{Ass}_{X'/S'}(\mathcal{F}')$.
Le choix de $y'$ donne aussi $\dim(Y'_{t'}) = \dim(Y_t)$,
voir par exemple Algèbre, lemme
\ref{algebra-lemma-inequalities-under-field-extension}.
D'après Algèbre, lemme
\ref{algebra-lemma-finite-type-domain-over-valuation-ring-dim-fibres},
$Y'_{s'}$ est équidimensionnel de dimension égale à $\dim(Y_t)$.
Puisque $\mathcal{F}$ est plat en $x$ sur $S$,
$\mathcal{F}'$ est plat en $x'$ sur $S'$, voir
Morphismes, lemme \ref{morphisms-lemma-base-change-module-flat}.

\medskip\noindent
Supposons que l'on puisse montrer $x' \in \text{Ass}_{X'/S}(\mathcal{F}')$. Alors
Diviseurs, lemme \ref{divisors-lemma-base-change-relative-assassin},
implique que $x \in \text{Ass}_{X/S}(\mathcal{F})$ et que
la composante irréductible $C'$ de $Y'_{s'}$ contenant $x'$
est une composante irréductible de $C \times_s s'$, où $C \subset Y_s$
est la composante irréductible contenant $x$.
D'où $\dim(C) = \dim(C') = \dim(Y_t)$ (voir ci-dessus), ce qui achève
la démonstration.
On est ainsi ramené au cas traité dans le paragraphe suivant.

\medskip\noindent
Supposons $S = \Spec(A)$, où $A$ est un anneau de valuation,
et où $t$ et $s$ sont les points générique et fermé de $S$.
Nous allons supposer $x \not \in \text{Ass}_{X/S}(\mathcal{F})$
afin d'aboutir à une contradiction. Autrement dit, on suppose
$x \not \in \text{Ass}_{X_s}(\mathcal{F}_s)$, où
$\mathcal{F}_s$ est l'image réciproque de $\mathcal{F}$ sur $X_s$.
Considérons l'homomorphisme d'anneaux
$$
A \longrightarrow \mathcal{O}_{X, x} = B
$$
et le module $N = \mathcal{F}_x$ sur $B = \mathcal{O}_{X, x}$.
Alors $B/\mathfrak m_A B = \mathcal{O}_{X_s, x}$ et
$N/\mathfrak m_A N$ est la fibre de $\mathcal{F}_s$ au point $x$.
Notons $\mathfrak q \subset B$ l'idéal premier correspondant au
point $y$, voir Schémas, lemme \ref{schemes-lemma-specialize-points}.
Puisque $x$ est un point générique de $Y_s$, le radical de
$\mathfrak q + \mathfrak m_A B$ est $\mathfrak m_B$.
Alors $\text{Ass}_{B/\mathfrak m_A B}(N/\mathfrak m_A N)$
est un ensemble fini d'idéaux premiers (Algèbre, lemme \ref{algebra-lemma-finite-ass})
qui ne contient pas l'idéal maximal de $B/\mathfrak m_A B$,
puisque $x \not \in \text{Ass}_{X/S}(\mathcal{F})$.
L'image de $\mathfrak q$ dans $B/\mathfrak m_A B$ n'est donc contenue
dans aucun de ces idéaux premiers. Par conséquent,
le lemme d'évitement des idéaux premiers (Algèbre, lemme \ref{algebra-lemma-silly})
permet de trouver un élément $g \in \mathfrak q$
dont l'image dans $B/\mathfrak m_A B$ est non diviseur de zéro sur
$N/\mathfrak m_A N$ (on utilise ici la description des diviseurs de zéro
d'Algèbre, lemme \ref{algebra-lemma-ass-zero-divisors}). Puisque
$N = \mathcal{F}_x$ est $A$-plat, le
lemme \ref{lemma-mod-injective-valuation-ring} montre que
$$
g : N \longrightarrow N
$$
est injectif. En particulier, si $K = \text{Frac}(A)$ est le
corps des fractions de $A$, alors
$$
g : N \otimes_A K \longrightarrow N \otimes_A K
$$
est injectif. Observons que $\mathfrak q$ correspond à un
idéal premier de $B \otimes_A K$. Notons $\mathcal{F}_t$ la
restriction de $\mathcal{F}$ à la fibre générique $X_t$. On a
$(B \otimes_A K)_{\mathfrak q} = \mathcal{O}_{X_t, y}$
et $(N \otimes_A K)_\mathfrak q$
est la fibre en $y$ de $\mathcal{F}_t$. On obtient donc
que $g \in \mathfrak m_y \subset \mathcal{O}_{X_t, y}$
est non diviseur de zéro sur la fibre $(\mathcal{F}_t)_y$,
ce qui contredit l'hypothèse $y \in \text{Ass}_{X/S}(\mathcal{F})$.
\end{proof}

\begin{lemma}
\label{lemma-relative-dimension-support-flat}
Soit $f : X \to S$ un morphisme de schémas localement de type fini.
Soit $\mathcal{F}$ un $\mathcal{O}_X$-module quasi-cohérent de type fini,
plat sur $S$. Supposons $S$ irréductible, de point générique $\eta$.
Si $\dim(\text{Supp}(\mathcal{F}_\eta)) \leq r$, alors, pour tout $s \in S$,
on a $\dim(\text{Supp}(\mathcal{F}_s)) \leq r$.
\end{lemma}

\begin{proof}
Soit $x \in \text{Supp}(\mathcal{F}_s)$ un point générique d'une composante
irréductible de $\text{Supp}(\mathcal{F}_s)$. D'après
Algèbre, lemme \ref{algebra-lemma-going-down-flat-module},
on peut trouver une spécialisation $y \leadsto x$ dans $\text{Supp}(\mathcal{F})$
avec $f(y) = \eta$. On peut bien entendu supposer que $y$ est un point générique
d'une composante irréductible de $\text{Supp}(\mathcal{F}_\eta)$.
Le lemme \ref{lemma-associated-point-specializes}
permet de conclure que la dimension de $\overline{\{x\}}$ est au plus $r$.
\end{proof}

\begin{lemma}
\label{lemma-flat-associated-equidimensional}
Soit $f : X \to S$ un morphisme de schémas localement de type fini.
Soit $\mathcal{F}$ un $\mathcal{O}_X$-module quasi-cohérent de type fini.
Soit $y \in \text{Ass}_{X/S}(\mathcal{F})$.
Notons $Y \subset X$ l'adhérence de $\{y\}$ dans $X$,
munie de sa structure de sous-schéma fermé intègre. Notons $T \subset S$ l'adhérence
de $\{f(y)\}$, munie de sa structure de sous-schéma fermé intègre. On obtient
un diagramme commutatif
$$
\xymatrix{
Y \ar[r] \ar[d] & X \ar[d] \\
T \ar[r] & S
}
$$
où $Y \to T$ est dominant. Supposons $\mathcal{F}$ plat sur $S$
en tous les points génériques des composantes irréductibles des fibres de $Y \to T$
(par exemple si $\mathcal{F}$ est plat sur $S$). Alors
\begin{enumerate}
\item si $s \in S$ et si $x \in Y_s$ est le point générique d'une
composante irréductible de $Y_s$, alors $x \in \text{Ass}_{X/S}(\mathcal{F})$, et
\item il existe un entier $d \geq 0$ tel que
$Y \to T$ soit de dimension relative $d$, voir
Morphismes, définition \ref{morphisms-definition-relative-dimension-d}.
\end{enumerate}
\end{lemma}

\begin{proof}
Cela résulte immédiatement de la version ponctuelle,
le lemme \ref{lemma-associated-point-specializes}.
Notons que, pour calculer la dimension des schémas
localement algébriques $Y_s$, il suffit de se placer au voisinage
des points génériques, voir Variétés, section
\ref{varieties-section-algebraic-schemes}.
\end{proof}

\begin{remark}
\label{remark-flat-equidimensional}
Voici quelques cas où les résultats ci-dessus, en particulier le
lemme \ref{lemma-flat-associated-equidimensional},
permettent de conclure qu'un morphisme de schémas $f : X \to S$
est de dimension relative $d$ au sens de
Morphismes, définition \ref{morphisms-definition-relative-dimension-d}.
C'est par exemple le cas si
\begin{enumerate}
\item $X$ est intègre, de point générique $\xi$,
\item le degré de transcendance de $\kappa(\xi)$ sur $\kappa(f(\xi))$ est $d$,
\item $f$ est localement de type fini, et
\item il existe un $\mathcal{O}_X$-module quasi-cohérent $\mathcal{F}$
de type fini, plat sur $S$ et tel que $\text{Supp}(\mathcal{F}) = X$.
\end{enumerate}
Voici un autre ensemble d'hypothèses qui convient :
\begin{enumerate}
\item $S$ est irréductible, de point générique $\eta$,
\item $X_\eta$ est dense dans $X$,
\item toute composante irréductible de $X_\eta$ est de dimension $d$,
\item $f$ est localement de type fini, et
\item il existe un $\mathcal{O}_X$-module quasi-cohérent $\mathcal{F}$
de type fini, plat sur $S$ et tel que $\text{Supp}(\mathcal{F}) = X$.
\end{enumerate}
On peut bien entendu affaiblir la condition de platitude sur $\mathcal{F}$
et demander seulement que $\mathcal{F}$ soit plat sur $S$ en codimension $0$,
c'est-à-dire que $\mathcal{F}$ soit plat sur $S$ en tout point générique de
toute fibre. Si nous avons un jour besoin de ces résultats, nous les énoncerons
et les démontrerons ici avec soin.
\end{remark}
\section{Retour sur la normalisation}
\label{section-normalization}

\noindent
La normalisation commute au changement de base lisse.

\begin{lemma}
\label{lemma-integral-closure-smooth-pullback}
Soit $f : Y \to X$ un morphisme lisse de schémas. Soit $\mathcal{A}$ un
faisceau quasi-cohérent de $\mathcal{O}_X$-algèbres. La clôture intégrale
de $\mathcal{O}_Y$ dans $f^*\mathcal{A}$ est égale à $f^*\mathcal{A}'$,
où $\mathcal{A}' \subset \mathcal{A}$ est la clôture intégrale de
$\mathcal{O}_X$ dans $\mathcal{A}$.
\end{lemma}

\begin{proof}
C'est la traduction du
lemme \ref{algebra-lemma-integral-closure-commutes-smooth} d'Algèbre
dans le langage des schémas. Les détails sont omis.
\end{proof}

\begin{lemma}[La normalisation commute au changement de base lisse]
\label{lemma-normalization-smooth-localization}
Soit
$$
\xymatrix{
Y_2 \ar[r] \ar[d]_{f_2} & Y_1 \ar[d]^{f_1} \\
X_2 \ar[r]^\varphi & X_1
}
$$
un carré cartésien dans la catégorie des schémas. Supposons $f_1$ quasi-compact
et quasi-séparé, et $\varphi$ lisse.
Soit $Y_i \to X_i' \to X_i$ la normalisation de $X_i$ dans $Y_i$.
Alors $X_2' \cong X_2 \times_{X_1} X_1'$.
\end{lemma}

\begin{proof}
Le changement de base de la factorisation $Y_1 \to X_1' \to X_1$ à $X_2$
est une factorisation $Y_2 \to X_2 \times_{X_1} X_1' \to X_2$, et
$X_2 \times_{X_1} X_1' \to X_2$ est entier
(Morphismes, lemme \ref{morphisms-lemma-base-change-finite}).
On obtient donc un morphisme
$h : X_2' \to X_2 \times_{X_1} X_1'$ par la propriété universelle de
Morphismes, lemme \ref{morphisms-lemma-characterize-normalization}.
Observons que $X_2'$ est le spectre relatif de la clôture intégrale
de $\mathcal{O}_{X_2}$ dans $f_{2, *}\mathcal{O}_{Y_2}$.
Si $\mathcal{A}' \subset f_{1, *}\mathcal{O}_{Y_1}$ désigne la clôture
intégrale de $\mathcal{O}_{X_1}$, alors $X_2 \times_{X_1} X_1'$ est le
spectre relatif de $\varphi^*\mathcal{A}'$, voir
Constructions, lemme \ref{constructions-lemma-spec-properties}.
D'après
Cohomologie des schémas, lemme \ref{coherent-lemma-flat-base-change-cohomology},
on sait que $f_{2, *}\mathcal{O}_{Y_2} = \varphi^*f_{1, *}\mathcal{O}_{Y_1}$.
Le résultat découle donc du
lemme \ref{lemma-integral-closure-smooth-pullback}.
\end{proof}

\begin{lemma}[Normalisation et morphismes lisses]
\label{lemma-normalization-and-smooth}
Soit $X \to Y$ un morphisme lisse de schémas. Supposons que tout ouvert
quasi-compact de $Y$ ait un nombre fini de composantes irréductibles. Il en est alors
de même de $X$, et il existe un unique isomorphisme $X^\nu = X \times_Y Y^\nu$
sur $X$, où $X^\nu$, $Y^\nu$ sont les normalisations de $X$, $Y$.
\end{lemma}

\begin{proof}
D'après Descente, lemme \ref{descent-lemma-locally-finite-nr-irred-local-fppf},
tout ouvert quasi-compact de $X$ a un nombre fini de composantes irréductibles.
Notons que $X_{red} = X \times_Y Y_{red}$, car un schéma lisse sur un schéma
réduit est réduit, voir
Descente, lemme \ref{descent-lemma-reduced-local-smooth}.
On peut donc supposer que $X$ et $Y$ sont réduits (car la normalisation
d'un schéma est, par définition, égale à celle de son réduit).
Notons ensuite que $X' = X \times_Y Y^\nu$ est un schéma normal
d'après Descente, lemme \ref{descent-lemma-normal-local-smooth}.
Le morphisme $X' \to Y^\nu$ est lisse (donc plat) ; les points génériques
des composantes irréductibles de $X'$ sont donc au-dessus des points génériques
des composantes irréductibles de $Y^\nu$. Puisque $Y^\nu \to Y$ est birationnel,
on en conclut que $X' \to X$ est aussi birationnel (car $X' \to Y^\nu$
induit un isomorphisme sur les fibres au-dessus des points génériques de $Y$).
On en déduit qu'il existe une factorisation
$X^\nu \to X' \to X$, voir
Morphismes, lemme \ref{morphisms-lemma-normalization-normal},
dont le premier morphisme est un isomorphisme, puisque $X'$ est normal et entier sur $X$.
\end{proof}

\begin{lemma}[Normalisation et hensélisation]
\label{lemma-normalization-and-henselization}
Soit $X$ un schéma localement noethérien. Soit $\nu : X^\nu \to X$
le morphisme de normalisation. Alors, pour tout point $x \in X$,
le changement de base
$$
X^\nu \times_X \Spec(\mathcal{O}_{X, x}^h) \to \Spec(\mathcal{O}_{X, x}^h),
\quad\text{resp.}\quad
X^\nu \times_X \Spec(\mathcal{O}_{X, x}^{sh}) \to \Spec(\mathcal{O}_{X, x}^{sh})
$$
est la normalisation de $\Spec(\mathcal{O}_{X, x}^h)$,
resp.\ $\Spec(\mathcal{O}_{X, x}^{sh})$.
\end{lemma}

\begin{proof}
Soient $\eta_1, \ldots, \eta_r$ les points génériques des composantes
irréductibles de $X$ passant par $x$. Le changement de base de la normalisation à
$\Spec(\mathcal{O}_{X, x})$ est le spectre de la
clôture intégrale de $\mathcal{O}_{X, x}$ dans $\prod \kappa(\eta_i)$.
Cela résulte de notre construction de la normalisation de $X$ dans
Morphismes, définition \ref{morphisms-definition-normalization},
et de Morphismes, lemme \ref{morphisms-lemma-integral-closure} ;
on peut aussi utiliser la description de la normalisation
dans Morphismes, lemme \ref{morphisms-lemma-description-normalization}.
On est ainsi ramené au problème d'algèbre suivant.
Soit $A$ un anneau local noethérien ; rappelons que son
hensélisé $A^h$ et son hensélisé strict $A^{sh}$
sont alors eux aussi noethériens (Compléments d'algèbre, lemme
\ref{more-algebra-lemma-henselization-noetherian}).
Soient $\mathfrak p_1, \ldots, \mathfrak p_r$
ses idéaux premiers minimaux. Soit $A'$ la clôture intégrale
de $A$ dans $\prod \kappa(\mathfrak p_i)$.
Le problème est de montrer que
$A' \otimes_A A^h$, resp.\ $A' \otimes_A A^{sh}$, s'obtient
à partir de l'anneau local noethérien $A^h$, resp.\ $A^{sh}$,
de la même manière.

\medskip\noindent
Puisque $A^h$, resp.\ $A^{sh}$, sont des limites inductives de $A$-algèbres étales,
les idéaux premiers minimaux de $A$ et $A^{sh}$ sont exactement
les idéaux premiers de $A^h$, resp.\ $A^{sh}$, au-dessus des idéaux premiers minimaux
de $A$ (par descente des idéaux premiers, voir
Algèbre, lemmes \ref{algebra-lemma-flat-going-down} et
\ref{algebra-lemma-minimal-prime-image-minimal-prime}).
Ainsi, Compléments d'algèbre, lemme \ref{more-algebra-lemma-fibres-henselization},
montre que
$A^h \otimes_A \prod \kappa(\mathfrak p_i)$,
resp.
$A^{sh} \otimes_A \prod \kappa(\mathfrak p_i)$,
est le produit des corps résiduels aux idéaux premiers minimaux
de $A^h$, resp.\ $A^{sh}$. On sait que la formation de la
clôture intégrale dans un suranneau commute au changement de base
étale, voir
Algèbre, lemme \ref{algebra-lemma-integral-closure-commutes-etale}.
En écrivant $A^h$ et $A^{sh}$ comme limites inductives de $A$-algèbres étales,
on voit qu'il en est de même pour le changement de base à
$A^h$ et $A^{sh}$ (on peut aussi utiliser le résultat plus général
d'Algèbre, lemme \ref{algebra-lemma-integral-closure-commutes-colim-smooth}).
\end{proof}








\section{Morphismes normaux}
\label{section-normal}

\noindent
L'article \cite{DM} de Deligne et Mumford mentionne la notion de morphisme
normal. Il s'agit d'un type parmi toute une série\footnote{
Les autres types sont coprof $\leq k$, Cohen-Macaulay, $(S_k)$,
régulier, $(R_k)$ et réduit. Voir \cite[IV, définition 6.8.1.]{EGA}.
Les morphismes de Gorenstein seront définis dans
Dualité pour les schémas, section \ref{duality-section-gorenstein}.}
de morphismes qui peuvent tous être définis de manière analogue. Nous les
introduirons au besoin, chacun dans sa propre section.

\begin{definition}
\label{definition-normal}
Soit $f : X \to Y$ un morphisme de schémas.
Supposons que toutes les fibres $X_y$ soient des schémas localement noethériens.
\begin{enumerate}
\item Soient $x \in X$ et $y = f(x)$. On dit que $f$ est {\it normal en $x$}
si $f$ est plat en $x$ et si le schéma $X_y$ est géométriquement
normal en $x$ sur $\kappa(y)$ (voir
Variétés, définition \ref{varieties-definition-geometrically-normal}).
\item On dit que $f$ est un {\it morphisme normal} si $f$ est normal
en tout point de $X$.
\end{enumerate}
\end{definition}

\noindent
La normalité du morphisme $X \to Y$ est donc une condition
plus forte que celle demandant seulement que toutes les
fibres soient des schémas normaux localement noethériens.

\begin{lemma}
\label{lemma-normal}
Soit $f : X \to Y$ un morphisme de schémas.
Supposons toutes les fibres de $f$ localement noethériennes.
Les conditions suivantes sont équivalentes :
\begin{enumerate}
\item $f$ est normal, et
\item $f$ est plat et ses fibres sont des schémas géométriquement normaux.
\end{enumerate}
\end{lemma}

\begin{proof}
Cela résulte directement des définitions.
\end{proof}

\begin{lemma}
\label{lemma-smooth-normal}
Un morphisme lisse est normal.
\end{lemma}

\begin{proof}
Soit $f : X \to Y$ un morphisme lisse.
Puisque $f$ est localement de présentation finie, voir
Morphismes, lemme \ref{morphisms-lemma-smooth-locally-finite-presentation},
les fibres $X_y$ sont localement de type fini sur un corps, donc
localement noethériennes. De plus, $f$ est plat, voir
Morphismes, lemme \ref{morphisms-lemma-smooth-flat}.
Enfin, les fibres $X_y$ sont lisses sur un corps (d'après
Morphismes, lemme \ref{morphisms-lemma-base-change-smooth}),
donc géométriquement normales d'après
Variétés, lemme \ref{varieties-lemma-smooth-geometrically-normal}.
Ainsi, $f$ est normal d'après le
lemme \ref{lemma-normal}.
\end{proof}

\noindent
Nous voulons montrer que cette notion est locale à la source et au but
pour la topologie lisse. Traitons d'abord la propriété d'avoir des
fibres localement noethériennes.

\begin{lemma}
\label{lemma-locally-Noetherian-fibres-fppf-local-source-and-target}
La propriété $\mathcal{P}(f)=$« les fibres de $f$ sont localement noethériennes »
est locale pour la topologie fppf à la source et au but.
\end{lemma}

\begin{proof}
Soit $f : X \to Y$ un morphisme de schémas.
Soit $\{\varphi_i : Y_i \to Y\}_{i \in I}$ un recouvrement fppf de $Y$.
Notons $f_i : X_i \to Y_i$ le changement de base de $f$ par $\varphi_i$.
Soit $i \in I$ et soit $y_i \in Y_i$ un point.
Posons $y = \varphi_i(y_i)$. Notons que
$$
X_{i, y_i} = \Spec(\kappa(y_i)) \times_{\Spec(\kappa(y))} X_y.
$$
De plus, puisque $\varphi_i$ est de présentation finie, l'extension de corps
$\kappa(y_i)/\kappa(y)$ est de type fini.
Dans cette situation, $X_y$ est donc localement noethérien si et
seulement si $X_{i, y_i}$ est localement noethérien, voir
Variétés, lemme \ref{varieties-lemma-locally-Noetherian-base-change}.
Ce fait donne le caractère local au but.

\medskip\noindent
Soit $\{X_i \to X\}$ un recouvrement fppf de $X$.
Soit $y \in Y$. Dans ce cas, $\{X_{i, y} \to X_y\}$ est un
recouvrement fppf de la fibre. Le caractère local à la source
résulte donc de Descente, lemme \ref{descent-lemma-Noetherian-local-fppf}.
\end{proof}

\begin{lemma}
\label{lemma-normal-fppf-local-source-and-target}
La propriété
$\mathcal{P}(f)=$« les fibres de $f$ sont localement noethériennes et $f$ est normal »
est locale pour la topologie fppf au but et
locale pour la topologie lisse à la source.
\end{lemma}

\begin{proof}
On a
$\mathcal{P}(f) =
\mathcal{P}_1(f) \wedge \mathcal{P}_2(f) \wedge \mathcal{P}_3(f)$,
où
$\mathcal{P}_1(f)=$« les fibres de $f$ sont localement noethériennes »,
$\mathcal{P}_2(f)=$« $f$ est plat », et
$\mathcal{P}_3(f)=$« les fibres de $f$ sont géométriquement normales ».
On a déjà vu que $\mathcal{P}_1$ et $\mathcal{P}_2$
sont locales pour la topologie fppf à la source et au but, voir le
lemme \ref{lemma-locally-Noetherian-fibres-fppf-local-source-and-target},
et Descente, lemmes \ref{descent-lemma-descending-property-flat} et
\ref{descent-lemma-flat-fpqc-local-source}. Il reste donc à traiter
$\mathcal{P}_3$.

\medskip\noindent
Soit $f : X \to Y$ un morphisme de schémas.
Soit $\{\varphi_i : Y_i \to Y\}_{i \in I}$ un recouvrement fpqc de $Y$.
Notons $f_i : X_i \to Y_i$ le changement de base de $f$ par $\varphi_i$.
Soit $i \in I$ et soit $y_i \in Y_i$ un point.
Posons $y = \varphi_i(y_i)$. Notons que
$$
X_{i, y_i} = \Spec(\kappa(y_i)) \times_{\Spec(\kappa(y))} X_y.
$$
Dans cette situation, $X_y$ est donc géométriquement normal si et
seulement si $X_{i, y_i}$ est géométriquement normal, voir
Variétés, lemme \ref{varieties-lemma-geometrically-normal-upstairs}.
Ce fait montre que $\mathcal{P}_3$ est locale pour la topologie fpqc au but.

\medskip\noindent
Soit $\{X_i \to X\}$ un recouvrement lisse de $X$.
Soit $y \in Y$. Dans ce cas, $\{X_{i, y} \to X_y\}$ est un
recouvrement lisse de la fibre. Le caractère local de $\mathcal{P}_3$
pour la topologie lisse à la source résulte donc de
Descente, lemme \ref{descent-lemma-normal-local-smooth}.
En combinant ces observations, on obtient le lemme.
\end{proof}
\section{Morphismes réguliers}
\label{section-regular}

\noindent
Comparer avec la section \ref{section-normal}. La version algébrique de
cette notion est étudiée dans
Compléments d'algèbre, section \ref{more-algebra-section-regular}.

\begin{definition}
\label{definition-regular}
Soit $f : X \to Y$ un morphisme de schémas.
Supposons que toutes les fibres $X_y$ soient des schémas localement noethériens.
\begin{enumerate}
\item Soient $x \in X$ et $y = f(x)$. On dit que $f$ est {\it régulier en $x$}
si $f$ est plat en $x$ et si le schéma $X_y$ est géométriquement
régulier en $x$ sur $\kappa(y)$ (voir
Variétés, définition \ref{varieties-definition-geometrically-regular}).
\item On dit que $f$ est un {\it morphisme régulier} si $f$ est régulier
en tout point de $X$.
\end{enumerate}
\end{definition}

\noindent
La régularité du morphisme $X \to Y$ est une condition
plus forte que celle demandant seulement que toutes les
fibres soient des schémas réguliers localement noethériens.

\begin{lemma}
\label{lemma-regular}
Soit $f : X \to Y$ un morphisme de schémas.
Supposons toutes les fibres de $f$ localement noethériennes.
Les conditions suivantes sont équivalentes :
\begin{enumerate}
\item $f$ est régulier,
\item $f$ est plat et ses fibres sont des schémas géométriquement réguliers,
\item pour toute paire d'ouverts affines $U \subset X$, $V \subset Y$
telle que $f(U) \subset V$, l'homomorphisme d'anneaux $\mathcal{O}_Y(V) \to \mathcal{O}_X(U)$
est régulier,
\item il existe un recouvrement ouvert $Y = \bigcup_{j \in J} V_j$
et des recouvrements ouverts $f^{-1}(V_j) = \bigcup_{i \in I_j} U_i$ tels
que chacun des morphismes $U_i \to V_j$ soit régulier, et
\item il existe un recouvrement ouvert affine $Y = \bigcup_{j \in J} V_j$
et des recouvrements ouverts affines $f^{-1}(V_j) = \bigcup_{i \in I_j} U_i$ tels
que les homomorphismes d'anneaux $\mathcal{O}_Y(V_j) \to \mathcal{O}_X(U_i)$ soient réguliers.
\end{enumerate}
\end{lemma}

\begin{proof}
L'équivalence de (1) et (2) résulte immédiatement des définitions.
Soit $x \in X$ et posons $y = f(x)$. Par définition, $f$ est plat en $x$
si et seulement si $\mathcal{O}_{Y, y} \to \mathcal{O}_{X, x}$ est un
homomorphisme d'anneaux plat, et $X_y$ est géométriquement régulier en $x$ sur
$\kappa(y)$ si et seulement si
$\mathcal{O}_{X_y, x} = \mathcal{O}_{X, x}/\mathfrak m_y\mathcal{O}_{X, x}$
est une algèbre géométriquement régulière sur $\kappa(y)$.
La régularité de $f$ en $x$
ne dépend donc que de l'homomorphisme local d'anneaux locaux
$\mathcal{O}_{Y, y} \to \mathcal{O}_{X, x}$.
L'équivalence de (1) et (4) est donc claire.

\medskip\noindent
Rappelons (Compléments d'algèbre, définition \ref{more-algebra-definition-regular})
qu'un homomorphisme d'anneaux $A \to B$ est régulier si et seulement s'il est plat
et si les anneaux des fibres $B \otimes_A \kappa(\mathfrak p)$ sont noethériens
et géométriquement réguliers pour tout idéal premier $\mathfrak p \subset A$.
D'après Variétés, lemme \ref{varieties-lemma-geometrically-regular},
cela équivaut à ce que $\Spec(B \otimes_A \kappa(\mathfrak p))$
soit un schéma géométriquement régulier sur $\kappa(\mathfrak p)$.
Ainsi, (2) implique (3). Il est clair que (3) implique
(5). Enfin, supposons (5). Alors $f$ est plat
(voir Morphismes, lemme \ref{morphisms-lemma-flat-characterize}).
De plus, si $y \in Y$, alors $y \in V_j$ pour un certain $j$, et
$X_y = \bigcup_{i \in I_j} U_{i, y}$, chaque $U_{i, y}$
étant géométriquement régulier sur $\kappa(y)$ d'après
Variétés, lemme \ref{varieties-lemma-geometrically-regular}.
Une nouvelle application de
Variétés, lemme \ref{varieties-lemma-geometrically-regular},
montre que $X_y$ est géométriquement régulier. Ainsi, (2) est vérifiée,
ce qui achève la démonstration du lemme.
\end{proof}

\begin{lemma}
\label{lemma-smooth-regular}
Un morphisme lisse est régulier.
\end{lemma}

\begin{proof}
Soit $f : X \to Y$ un morphisme lisse.
Puisque $f$ est localement de présentation finie, voir
Morphismes, lemme \ref{morphisms-lemma-smooth-locally-finite-presentation},
les fibres $X_y$ sont localement de type fini sur un corps, donc
localement noethériennes. De plus, $f$ est plat, voir
Morphismes, lemme \ref{morphisms-lemma-smooth-flat}.
Enfin, les fibres $X_y$ sont lisses sur un corps (d'après
Morphismes, lemme \ref{morphisms-lemma-base-change-smooth}),
donc géométriquement régulières d'après
Variétés, lemme \ref{varieties-lemma-smooth-geometrically-normal}.
Ainsi, $f$ est régulier d'après le
lemme \ref{lemma-regular}.
\end{proof}

\begin{lemma}
\label{lemma-regular-fppf-local-source-and-target}
La propriété $\mathcal{P}(f)=$« les fibres de $f$ sont
localement noethériennes et $f$ est régulier »
est locale pour la topologie fppf au but et
locale pour la topologie lisse à la source.
\end{lemma}

\begin{proof}
On a
$\mathcal{P}(f) =
\mathcal{P}_1(f) \wedge \mathcal{P}_2(f) \wedge \mathcal{P}_3(f)$,
où
$\mathcal{P}_1(f)=$« les fibres de $f$ sont localement noethériennes »,
$\mathcal{P}_2(f)=$« $f$ est plat », et
$\mathcal{P}_3(f)=$« les fibres de $f$ sont géométriquement régulières ».
On a déjà vu que $\mathcal{P}_1$ et $\mathcal{P}_2$
sont locales pour la topologie fppf à la source et au but, voir le
lemme \ref{lemma-locally-Noetherian-fibres-fppf-local-source-and-target},
et Descente, lemmes \ref{descent-lemma-descending-property-flat} et
\ref{descent-lemma-flat-fpqc-local-source}. Il reste donc à traiter
$\mathcal{P}_3$.

\medskip\noindent
Soit $f : X \to Y$ un morphisme de schémas.
Soit $\{\varphi_i : Y_i \to Y\}_{i \in I}$ un recouvrement fpqc de $Y$.
Notons $f_i : X_i \to Y_i$ le changement de base de $f$ par $\varphi_i$.
Soit $i \in I$ et soit $y_i \in Y_i$ un point.
Posons $y = \varphi_i(y_i)$. Notons que
$$
X_{i, y_i} = \Spec(\kappa(y_i)) \times_{\Spec(\kappa(y))} X_y.
$$
Dans cette situation, $X_y$ est donc géométriquement régulier si et
seulement si $X_{i, y_i}$ est géométriquement régulier, voir
Variétés, lemme \ref{varieties-lemma-geometrically-regular-upstairs}.
Ce fait montre que $\mathcal{P}_3$ est locale pour la topologie fpqc au but.

\medskip\noindent
Soit $\{X_i \to X\}$ un recouvrement lisse de $X$.
Soit $y \in Y$. Dans ce cas, $\{X_{i, y} \to X_y\}$ est un
recouvrement lisse de la fibre. Le caractère local de $\mathcal{P}_3$
pour la topologie lisse à la source résulte donc de
Descente, lemme \ref{descent-lemma-regular-local-smooth}.
En combinant ces observations, on obtient le lemme.
\end{proof}
\section{Morphismes de Cohen-Macaulay}
\label{section-CM}

\noindent
Comparer avec la section \ref{section-normal}.
Notons que, comme il est signalé dans
Algèbre, section \ref{algebra-section-geometrically-CM},
et dans
Variétés, section \ref{varieties-section-CM},
« géométriquement de Cohen-Macaulay » équivaut simplement à « de Cohen-Macaulay ».

\begin{definition}
\label{definition-CM}
Soit $f : X \to Y$ un morphisme de schémas.
Supposons que toutes les fibres $X_y$ soient des schémas localement noethériens.
\begin{enumerate}
\item Soient $x \in X$ et $y = f(x)$. On dit que $f$ est
{\it de Cohen-Macaulay en $x$} si $f$ est plat en $x$ et si
l'anneau local du schéma $X_y$ en $x$ est de Cohen-Macaulay.
\item On dit que $f$ est un {\it morphisme de Cohen-Macaulay} si $f$ est
de Cohen-Macaulay en tout point de $X$.
\end{enumerate}
\end{definition}

\noindent
Voici une reformulation.

\begin{lemma}
\label{lemma-CM}
Soit $f : X \to Y$ un morphisme de schémas.
Supposons toutes les fibres de $f$ localement noethériennes.
Les conditions suivantes sont équivalentes :
\begin{enumerate}
\item $f$ est de Cohen-Macaulay, et
\item $f$ est plat et ses fibres sont des schémas de Cohen-Macaulay.
\end{enumerate}
\end{lemma}

\begin{proof}
Cela résulte directement des définitions.
\end{proof}

\begin{lemma}
\label{lemma-CM-dimension}
Soit $f : X \to Y$ un morphisme de schémas localement noethériens,
localement de type fini et de Cohen-Macaulay.
Pour tout point $x$ de $X$ d'image $y$ dans $Y$,
$$
\dim_x(X) = \dim_y(Y) + \dim_x(X_y),
$$
où $X_y$ désigne la fibre au-dessus de $y$.
\end{lemma}

\begin{proof}
Après remplacement de $X$ par un voisinage ouvert de $x$,
il existe un entier naturel $d$ tel que toutes les fibres
de $X \to Y$ soient de dimension $d$ en tout point, voir
Morphismes, lemme
\ref{morphisms-lemma-flat-finite-presentation-CM-fibres-relative-dimension}.
Alors $f$ est plat, localement de type fini
et de dimension relative $d$. Le résultat découle donc de
Morphismes, lemme \ref{morphisms-lemma-rel-dimension-dimension}.
\end{proof}

\begin{lemma}
\label{lemma-composition-CM}
Soient $f : X \to Y$ et $g : Y \to Z$ des morphismes de schémas. Supposons que les
fibres de $f$, $g$ et $g \circ f$ soient localement noethériennes.
Soit $x \in X$ d'images $y \in Y$ et $z \in Z$.
\begin{enumerate}
\item Si $f$ est de Cohen-Macaulay en $x$ et si $g$ est de Cohen-Macaulay
en $f(x)$, alors $g \circ f$ est de Cohen-Macaulay en $x$.
\item Si $f$ et $g$ sont de Cohen-Macaulay, alors $g \circ f$ est de Cohen-Macaulay.
\item Si $g \circ f$ est de Cohen-Macaulay en $x$ et si $f$ est plat en $x$,
alors $f$ est de Cohen-Macaulay en $x$ et $g$ est de Cohen-Macaulay en $f(x)$.
\item Si $g \circ f$ est de Cohen-Macaulay et si $f$ est plat, alors
$f$ est de Cohen-Macaulay et $g$ est de Cohen-Macaulay en tout point
de l'image de $f$.
\end{enumerate}
\end{lemma}

\begin{proof}
Considérons l'homomorphisme d'anneaux locaux noethériens
$$
\mathcal{O}_{Y_z, y} \to \mathcal{O}_{X_z, x}
$$
et observons que sa fibre est
$$
\mathcal{O}_{X_z, x}/\mathfrak m_{Y_z, y}\mathcal{O}_{X_z, x} =
\mathcal{O}_{X_y, x}
$$
Le lemme résulte donc du
lemme \ref{algebra-lemma-CM-goes-up} d'Algèbre.
\end{proof}

\begin{lemma}
\label{lemma-flat-morphism-from-CM-scheme}
Soit $f : X \to Y$ un morphisme plat de schémas localement noethériens.
Si $X$ est de Cohen-Macaulay, alors $f$ est de Cohen-Macaulay et
$\mathcal{O}_{Y, f(x)}$ est de Cohen-Macaulay pour tout $x \in X$.
\end{lemma}

\begin{proof}
Après traduction en algèbre, cela résulte du
lemme \ref{algebra-lemma-CM-goes-up} d'Algèbre.
\end{proof}

\begin{lemma}
\label{lemma-base-change-CM}
Soit $f : X \to Y$ un morphisme de schémas.
Supposons que toutes les fibres $X_y$ soient des schémas localement noethériens.
Soit $Y' \to Y$ un morphisme localement de type fini. Soit $f' : X' \to Y'$
le changement de base de $f$.
Soit $x' \in X'$ un point d'image $x \in X$.
\begin{enumerate}
\item Si $f$ est de Cohen-Macaulay en $x$, alors
$f' : X' \to Y'$ est de Cohen-Macaulay en $x'$.
\item Si $f$ est plat en $x$ et si $f'$ est de Cohen-Macaulay en $x'$, alors $f$
est de Cohen-Macaulay en $x$.
\item Si $Y' \to Y$ est plat en $f'(x')$ et si $f'$ est de Cohen-Macaulay en
$x'$, alors $f$ est de Cohen-Macaulay en $x$.
\end{enumerate}
\end{lemma}

\begin{proof}
Notons que l'hypothèse sur $Y' \to Y$ implique que, pour $y' \in Y'$
d'image $y \in Y$, l'extension de corps $\kappa(y')/\kappa(y)$
est de type fini. Toutes les fibres
$X'_{y'} = (X_y)_{\kappa(y')}$ sont donc elles aussi localement noethériennes, voir
Variétés, lemme \ref{varieties-lemma-locally-Noetherian-base-change}.
Le lemme a donc un sens. Posons $y' = f'(x')$ et $y = f(x)$.
On obtient le diagramme commutatif d'anneaux locaux suivant :
$$
\xymatrix{
\mathcal{O}_{X', x'} & \mathcal{O}_{X, x} \ar[l] \\
\mathcal{O}_{Y', y'} \ar[u] & \mathcal{O}_{Y, y} \ar[l] \ar[u]
}
$$
où l'anneau en haut à gauche est un localisé du produit tensoriel des anneaux
en haut à droite et en bas à gauche sur l'anneau en bas à droite.

\medskip\noindent
Supposons $f$ de Cohen-Macaulay en $x$.
La platitude de $\mathcal{O}_{Y, y} \to \mathcal{O}_{X, x}$
implique celle de $\mathcal{O}_{Y', y'} \to \mathcal{O}_{X', x'}$, voir
Algèbre, lemme \ref{algebra-lemma-base-change-flat-up-down}.
Le fait que $\mathcal{O}_{X, x}/\mathfrak m_y\mathcal{O}_{X, x}$
soit de Cohen-Macaulay implique que
$\mathcal{O}_{X', x'}/\mathfrak m_{y'}\mathcal{O}_{X', x'}$
est de Cohen-Macaulay, voir
Variétés, lemme \ref{varieties-lemma-CM-base-change}. Ainsi, $f'$
est de Cohen-Macaulay en $x'$.

\medskip\noindent
Supposons $f$ plat en $x$ et $f'$ de Cohen-Macaulay en $x'$.
Le fait que $\mathcal{O}_{X', x'}/\mathfrak m_{y'}\mathcal{O}_{X', x'}$
soit de Cohen-Macaulay implique que
$\mathcal{O}_{X, x}/\mathfrak m_y\mathcal{O}_{X, x}$
est de Cohen-Macaulay, voir
Variétés, lemme \ref{varieties-lemma-CM-base-change}.
Ainsi, $f$ est de Cohen-Macaulay en $x$.

\medskip\noindent
Supposons $Y' \to Y$ plat en $y'$ et $f'$ de Cohen-Macaulay en
$x'$. La platitude de $\mathcal{O}_{Y', y'} \to \mathcal{O}_{X', x'}$
et de $\mathcal{O}_{Y, y} \to \mathcal{O}_{Y', y'}$ implique celle
de $\mathcal{O}_{Y, y} \to \mathcal{O}_{X, x}$, voir
Algèbre, lemme \ref{algebra-lemma-base-change-flat-up-down}.
Le fait que $\mathcal{O}_{X', x'}/\mathfrak m_{y'}\mathcal{O}_{X', x'}$
soit de Cohen-Macaulay implique que
$\mathcal{O}_{X, x}/\mathfrak m_y\mathcal{O}_{X, x}$
est de Cohen-Macaulay, voir
Variétés, lemme \ref{varieties-lemma-CM-base-change}. Ainsi, $f$
est de Cohen-Macaulay en $x$.
\end{proof}

\begin{lemma}
\label{lemma-flat-finite-presentation-CM-open}
\begin{reference}
\cite[IV, corollaire 12.1.7(iii)]{EGA}
\end{reference}
Soit $f : X \to S$ un morphisme de schémas plat et localement
de présentation finie. Posons
$$
W = \{x \in X \mid f\text{ est de Cohen-Macaulay en }x\}
$$
Alors
\begin{enumerate}
\item $W = \{x \in X \mid \mathcal{O}_{X_{f(x)}, x}\text{ est de Cohen-Macaulay}\}$,
\item $W$ est ouvert dans $X$,
\item $W$ est dense dans toute fibre de $X \to S$,
\item la formation de $W$ commute à tout changement de base de $f$ :
pour tout morphisme $g : S' \to S$, considérons
le changement de base $f' : X' \to S'$ de $f$ et la
projection $g' : X' \to X$. Alors l'ensemble
$W'$ correspondant au morphisme $f'$ est égal à $W' = (g')^{-1}(W)$.
\end{enumerate}
\end{lemma}

\begin{proof}
Puisque $f$ est plat à fibres localement noethériennes, l'égalité de (1)
vaut par définition. Les assertions (2) et (3) résultent du
lemme \ref{algebra-lemma-generic-CM-flat-finite-presentation} d'Algèbre.
L'assertion (4) résulte soit du
lemme \ref{algebra-lemma-CM-locus-commutes-base-change} d'Algèbre,
soit de
Variétés, lemme \ref{varieties-lemma-CM-base-change}.
\end{proof}

\begin{lemma}
\label{lemma-flat-finite-presentation-characterize-CM}
Soit $f : X \to S$ un morphisme de schémas plat et localement
de présentation finie. Soit $x \in X$ d'image $s \in S$.
Posons $d = \dim_x(X_s)$.
Les conditions suivantes sont équivalentes :
\begin{enumerate}
\item $f$ est de Cohen-Macaulay en $x$,
\item il existe un voisinage ouvert $U \subset X$ de $x$
et un morphisme localement quasi-fini $U \to \mathbf{A}^d_S$ sur $S$
qui est plat en $x$,
\item il existe un voisinage ouvert $U \subset X$ de $x$
et un morphisme plat localement quasi-fini $U \to \mathbf{A}^d_S$ sur $S$,
\item pour tout $S$-morphisme $g : U \to \mathbf{A}^d_S$
d'un voisinage ouvert $U \subset X$ de $x$, on a :
$g$ est quasi-fini en $x$ $\Rightarrow$ $g$ est plat en $x$.
\end{enumerate}
\end{lemma}

\begin{proof}
L'ouverture du lieu de platitude montre que (2) et (3)
sont équivalentes, voir le théorème \ref{theorem-openness-flatness}.

\medskip\noindent
Choisissons des ouverts affines $U = \Spec(A) \subset X$ avec $x \in U$ et
$V = \Spec(R) \subset S$ avec $f(U) \subset V$. Alors $R \to A$
est un homomorphisme d'anneaux plat de présentation finie. Soit $\mathfrak p \subset A$
l'idéal premier correspondant à $x$. Après remplacement de $A$ par un
localisé principal, on peut supposer qu'il existe un homomorphisme quasi-fini
$R[x_1, \ldots, x_d]  \to A$, voir
Algèbre, lemme \ref{algebra-lemma-quasi-finite-over-polynomial-algebra}.
Il existe donc au moins un couple $(U, g)$ formé d'un
voisinage ouvert $U \subset X$ de $x$ et d'un morphisme
localement\footnote{Si $S$ est quasi-séparé, alors $g$ sera quasi-fini.} quasi-fini
$g : U \to \mathbf{A}^d_S$.

\medskip\noindent
Affirmation : étant donnés $R \to A$ plat et de présentation finie, un idéal premier
$\mathfrak p \subset A$ et $\varphi : R[x_1, \ldots, x_d] \to A$
quasi-fini en $\mathfrak p$, on a : $\Spec(\varphi)$
est plat en $\mathfrak p$ si et seulement si $\Spec(A) \to \Spec(R)$
est de Cohen-Macaulay en $\mathfrak p$. En effet, d'après le
théorème \ref{theorem-criterion-flatness-fibre},
la platitude peut se vérifier sur les fibres. Il en est de même
de la propriété d'être de Cohen-Macaulay (puisque $A$ est déjà supposé plat sur $R$).
L'affirmation résulte donc du
lemme \ref{algebra-lemma-where-CM} d'Algèbre.

\medskip\noindent
Cette affirmation montre que (1) équivaut à (4) ; compte tenu
de la construction d'un couple convenable $(U, g)$ au deuxième
paragraphe, elle montre aussi que (1) équivaut à (2).
\end{proof}

\begin{lemma}
\label{lemma-flat-finite-presentation-CM-pieces}
Soit $f : X \to S$ un morphisme de schémas plat et localement
de présentation finie. Pour $d \geq 0$, il existe des ouverts $U_d \subset X$
ayant les propriétés suivantes :
\begin{enumerate}
\item $W = \bigcup_{d \geq 0} U_d$ est dense dans toute fibre de $f$, et
\item $U_d \to S$ est de dimension relative $d$ (voir
Morphismes, définition \ref{morphisms-definition-relative-dimension-d}).
\end{enumerate}
\end{lemma}

\begin{proof}
Cela résulte de la combinaison du
lemme \ref{lemma-flat-finite-presentation-CM-open}
avec
Morphismes, lemme
\ref{morphisms-lemma-flat-finite-presentation-CM-fibres-relative-dimension}.
\end{proof}

\begin{lemma}
\label{lemma-flat-finite-presentation-specialization-dimension}
Soit $f : X \to S$ un morphisme de schémas plat et localement
de présentation finie.
Supposons que $x' \leadsto x$ soit une spécialisation de points de $X$,
d'image $s' \leadsto s$ dans $S$. Si $x$ est un point générique d'une
composante irréductible de $X_s$, alors $\dim_{x'}(X_{s'}) = \dim_x(X_s)$.
\end{lemma}

\begin{proof}
Le point $x$ appartient à $U_d$ pour un certain $d$, où $U_d$ est comme au
lemme \ref{lemma-flat-finite-presentation-CM-pieces}.
\end{proof}

\begin{lemma}
\label{lemma-CM-local-source-and-target}
La propriété
$\mathcal{P}(f)=$« les fibres de $f$ sont localement noethériennes et $f$ est
de Cohen-Macaulay » est locale pour la topologie fppf au but et
locale pour la topologie syntomique à la source.
\end{lemma}

\begin{proof}
On a
$\mathcal{P}(f) =
\mathcal{P}_1(f) \wedge \mathcal{P}_2(f)$,
où
$\mathcal{P}_1(f)=$« $f$ est plat », et
$\mathcal{P}_2(f)=$« les fibres de $f$ sont localement noethériennes
et de Cohen-Macaulay ».
On sait que $\mathcal{P}_1$ est
locale pour la topologie fppf à la source et au but, voir
Descente, lemmes \ref{descent-lemma-descending-property-flat} et
\ref{descent-lemma-flat-fpqc-local-source}. Il reste donc à traiter
$\mathcal{P}_2$.

\medskip\noindent
Soit $f : X \to Y$ un morphisme de schémas.
Soit $\{\varphi_i : Y_i \to Y\}_{i \in I}$ un recouvrement fppf de $Y$.
Notons $f_i : X_i \to Y_i$ le changement de base de $f$ par $\varphi_i$.
Soit $i \in I$ et soit $y_i \in Y_i$ un point.
Posons $y = \varphi_i(y_i)$. Notons que
$$
X_{i, y_i} = \Spec(\kappa(y_i)) \times_{\Spec(\kappa(y))} X_y.
$$
et que $\kappa(y_i)/\kappa(y)$ est une extension de corps
de type fini. Ainsi, si $X_y$ est localement noethérien,
$X_{i, y_i}$ est localement noethérien, voir
Variétés, lemme \ref{varieties-lemma-locally-Noetherian-base-change}.
Si, de plus, $X_y$ est de Cohen-Macaulay,
alors $X_{i, y_i}$ est de Cohen-Macaulay, voir
Variétés, lemme \ref{varieties-lemma-CM-base-change}.
Ainsi, $\mathcal{P}_2$ est locale pour la topologie fppf au but.

\medskip\noindent
Soit $\{X_i \to X\}$ un recouvrement syntomique de $X$.
Soit $y \in Y$. Dans ce cas, $\{X_{i, y} \to X_y\}$ est un
recouvrement syntomique de la fibre. Le caractère local de $\mathcal{P}_2$
pour la topologie syntomique à la source résulte donc de
Descente, lemme \ref{descent-lemma-CM-local-syntomic}.
En combinant ces observations, on obtient le lemme.
\end{proof}

\section{Sections des morphismes de Cohen-Macaulay}
\label{section-slicing-CM}

\noindent
Les résultats de cette section conduisent finalement à l'assertion selon laquelle
la topologie fppf est la même que la topologie
« de présentation finie, plate et quasi-finie ».
Le lemme suivant est très proche de
Diviseurs, lemme \ref{divisors-lemma-fibre-Cartier}.

\begin{lemma}
\label{lemma-slice-given-element}
Soit $f : X \to S$ un morphisme de schémas.
Soit $x \in X$ un point d'image $s \in S$.
Soit $h \in \mathfrak m_x \subset \mathcal{O}_{X, x}$.
Supposons que
\begin{enumerate}
\item $f$ soit localement de présentation finie,
\item $f$ soit plat en $x$, et
\item l'image $\overline{h}$ de $h$ dans
$\mathcal{O}_{X_s, x} = \mathcal{O}_{X, x}/\mathfrak m_s\mathcal{O}_{X, x}$
soit non diviseur de zéro.
\end{enumerate}
Alors il existe un voisinage ouvert affine $U \subset X$ de $x$
tel que $h$ provienne d'un élément $h \in \Gamma(U, \mathcal{O}_U)$ et
tel que $D = V(h)$ soit un diviseur de Cartier effectif de $U$, avec $x \in D$ et
$D \to S$ plat et localement de présentation finie.
\end{lemma}

\begin{proof}
Nous allons le démontrer en nous ramenant au cas noethérien.
Par ouverture du lieu de platitude (voir le
théorème \ref{theorem-openness-flatness}),
on peut supposer, après remplacement de $X$ par un
voisinage ouvert de $x$, que $X \to S$ est plat.
On peut aussi supposer $X$ et $S$ affines.
Après restriction éventuelle de $X$, on peut supposer qu'il existe
un élément $h \in \Gamma(X, \mathcal{O}_X)$ dont l'image est l'élément $h$ donné.

\medskip\noindent
On peut écrire $S = \Spec(A)$ et $A = \colim_i A_i$
comme limite inductive filtrante de $\mathbf{Z}$-algèbres de type fini.
Alors, d'après
Algèbre, lemme \ref{algebra-lemma-flat-finite-presentation-limit-flat},
ou
Limites, lemmes \ref{limits-lemma-descend-finite-presentation},
\ref{limits-lemma-descend-affine-finite-presentation} et
\ref{limits-lemma-descend-finite-presentation},
on peut trouver un diagramme cartésien
$$
\xymatrix{
X \ar[r] \ar[d]_f & X_0 \ar[d]^{f_0} \\
S \ar[r] & S_0
}
$$
avec $f_0$ plat et de présentation finie, $X_0$ affine et
$S_0$ affine et noethérien. Soit $x_0 \in X_0$, resp.\ $s_0 \in S_0$,
l'image de $x$, resp.\ $s$. On peut aussi supposer qu'il existe un élément
$h_0 \in \Gamma(X_0, \mathcal{O}_{X_0})$ qui se restreint à $h$ sur $X$.
(Si l'on utilise la référence algébrique ci-dessus, c'est clair ; si l'on utilise
les références au chapitre sur les limites, cela résulte de
Limites, lemme \ref{limits-lemma-descend-finite-presentation},
en considérant $h$ comme un morphisme $X \to \mathbf{A}^1_S$.)
Notons que $\mathcal{O}_{X_s, x}$ est un localisé de
$\mathcal{O}_{(X_0)_{s_0}, x_0} \otimes_{\kappa(s_0)} \kappa(s)$, de sorte que
$\mathcal{O}_{(X_0)_{s_0}, x_0} \to \mathcal{O}_{X_s, x}$ est un homomorphisme
local plat, donc fidèlement plat. Par conséquent, l'image
$\overline{h}_0 \in \mathcal{O}_{(X_0)_{s_0}, x_0}$
appartient à $\mathfrak m_{(X_0)_{s_0}, x_0}$ et est non diviseur de zéro.
Nous affirmons qu'après remplacement de $X_0$ par un voisinage ouvert principal de
$x_0$, l'élément $h_0$ est non diviseur de zéro dans
$B_0 = \Gamma(X_0, \mathcal{O}_{X_0})$ et $B_0/h_0B_0$ est plat
sur $A_0 = \Gamma(S_0, \mathcal{O}_{S_0})$.
Dans ce cas,
$$
0 \to B_0 \xrightarrow{h_0} B_0 \to B_0/h_0B_0 \to 0
$$
est une suite exacte courte de $A_0$-modules plats. Elle reste donc exacte
après tensorisation par $A$ (d'après
Algèbre, lemme \ref{algebra-lemma-flat-tor-zero}),
et le lemme s'ensuit.

\medskip\noindent
Il reste à démontrer l'affirmation ci-dessus. L'énoncé algébrique correspondant
est le suivant (on omet ici l'indice ${}_0$) :
soit $A \to B$ un homomorphisme d'anneaux noethériens, plat et de type fini.
Soit $\mathfrak q \subset B$ un idéal premier au-dessus de $\mathfrak p \subset A$.
Supposons que $h \in \mathfrak q$ ait pour image un non-diviseur de zéro dans
$B_{\mathfrak q}/\mathfrak p B_{\mathfrak q}$.
Il s'agit de montrer qu'après remplacement éventuel de $B$ par $B_g$ pour un
$g \in B$, $g \not \in \mathfrak q$, l'élément $h$ devient non diviseur de zéro
et $B/hB$ devient plat sur $A$. D'après
Algèbre, lemme \ref{algebra-lemma-grothendieck},
$h$ est non diviseur de zéro dans $B_{\mathfrak q}$ et
$B_{\mathfrak q}/hB_{\mathfrak q}$ est plat sur $A$.
Par ouverture du lieu de platitude, voir
Algèbre, théorème \ref{algebra-theorem-openness-flatness},
ou le
théorème \ref{theorem-openness-flatness},
$B/hB$ est plat sur $A$ après remplacement de $B$ par $B_g$ pour un
$g \in B$, $g \not \in \mathfrak q$. Enfin, soit
$I = \{b \in B \mid hb = 0\}$ l'annulateur de $h$. Alors
$IB_{\mathfrak q} = 0$, puisque $h$ est non diviseur de zéro dans $B_{\mathfrak q}$.
De plus, $I$ est de type fini, car $B$ est noethérien.
Il existe donc un $g \in B$, $g \not \in \mathfrak q$, tel que $IB_g = 0$.
Après remplacement de $B$ par $B_g$, l'élément $h$ est non diviseur de zéro.
\end{proof}

\begin{lemma}
\label{lemma-slice-given-elements}
Soit $f : X \to S$ un morphisme de schémas.
Soit $x \in X$ un point d'image $s \in S$.
Soient $h_1, \ldots, h_r \in \mathcal{O}_{X, x}$.
Supposons que
\begin{enumerate}
\item $f$ soit localement de présentation finie,
\item $f$ soit plat en $x$, et
\item les images de $h_1, \ldots, h_r$ dans
$\mathcal{O}_{X_s, x} = \mathcal{O}_{X, x}/\mathfrak m_s\mathcal{O}_{X, x}$
forment une suite régulière.
\end{enumerate}
Alors il existe un voisinage ouvert affine $U \subset X$ de $x$
tel que $h_1, \ldots, h_r$ proviennent d'éléments
$h_1, \ldots, h_r \in \Gamma(U, \mathcal{O}_U)$ et
tel que $Z = V(h_1, \ldots, h_r) \to U$ soit une immersion régulière, avec
$x \in Z$ et $Z \to S$ plat et localement de présentation finie.
De plus, le changement de base $Z_{S'} \to U_{S'}$ est une immersion régulière
pour tout schéma $S'$ sur $S$.
\end{lemma}

\begin{proof}
(Nos conventions sur les suites régulières impliquent que $h_i \in \mathfrak m_x$
pour tout $i$.) Le cas $r = 1$ résulte du
lemme \ref{lemma-slice-given-element}
combiné avec
Diviseurs, lemme \ref{divisors-lemma-relative-Cartier},
qui montre que $V(h_1)$ reste un diviseur de Cartier effectif après changement de base.
Le cas $r > 1$ résulte d'une récurrence immédiate sur $r$ (en appliquant
le résultat pour $r = 1$ exactement $r$ fois ; les détails sont omis).

\medskip\noindent
Une autre démonstration du lemme utilise les résultats de
Diviseurs, section \ref{divisors-section-relative-regular-immersion}.
En effet, d'abord, par ouverture du lieu de platitude (voir le
théorème \ref{theorem-openness-flatness}),
on peut supposer, après remplacement de $X$ par un
voisinage ouvert de $x$, que $X \to S$ est plat.
On peut aussi supposer $X$ et $S$ affines.
Après restriction éventuelle de $X$, on peut supposer que l'on a
$h_1, \ldots, h_r \in \Gamma(X, \mathcal{O}_X)$. Posons
$Z = V(h_1, \ldots, h_r)$. Notons que $X_s$ est un schéma noethérien
(car c'est un $\kappa(s)$-schéma algébrique, voir
Variétés, section \ref{varieties-section-algebraic-schemes})
et que la topologie de $X_s$ est induite par celle de $X$
(voir
Schémas, lemme \ref{schemes-lemma-fibre-topological}).
Ainsi, en restreignant encore $X$,
on peut supposer que $Z_s \subset X_s$ est une immersion régulière
définie par les $r$ éléments $h_i|_{X_s}$, voir
Diviseurs, lemme \ref{divisors-lemma-Noetherian-scheme-regular-ideal}
et sa démonstration. Il est aussi clair que $r = \dim_x(X_s) - \dim_x(Z_s)$, car
\begin{align*}
\dim_x(X_s) & = \dim(\mathcal{O}_{X_s, x}) +
\text{trdeg}_{\kappa(s)}(\kappa(x)), \\
\dim_x(Z_s) & = \dim(\mathcal{O}_{Z_s, x}) +
\text{trdeg}_{\kappa(s)}(\kappa(x)), \\
\dim(\mathcal{O}_{X_s, x}) & = \dim(\mathcal{O}_{Z_s, x}) + r
\end{align*}
les deux premières égalités résultant du
lemme \ref{algebra-lemma-dimension-at-a-point-finite-type-field} d'Algèbre,
et la troisième de $r$ applications du
lemme \ref{algebra-lemma-one-equation} d'Algèbre.
Ainsi,
Diviseurs, lemme \ref{divisors-lemma-fibre-quasi-regular}, assertion (3),
s'applique et montre qu'après restriction de $X$ pour la topologie de Zariski, le morphisme
$Z \to X$ est une immersion régulière à laquelle s'applique
Diviseurs, lemme
\ref{divisors-lemma-relative-regular-immersion-flat-in-neighbourhood}
(ce qui donne la platitude et l'assertion sur le changement de base).
\end{proof}

\begin{lemma}
\label{lemma-slice-once}
Soit $f : X \to S$ un morphisme de schémas.
Soit $x \in X$ un point d'image $s \in S$.
Supposons que
\begin{enumerate}
\item $f$ soit localement de présentation finie,
\item $f$ soit plat en $x$, et
\item $\mathcal{O}_{X_s, x}$ vérifie $\text{depth} \geq 1$.
\end{enumerate}
Alors il existe un voisinage ouvert affine $U \subset X$ de $x$
et un diviseur de Cartier effectif $D \subset U$ contenant $x$ tels que
$D \to S$ soit plat et de présentation finie.
\end{lemma}

\begin{proof}
Choisir un élément $h \in \mathfrak m_x \subset \mathcal{O}_{X, x}$ dont
l'image soit non diviseur de zéro dans $\mathcal{O}_{X_s, x}$ et appliquer le
lemme \ref{lemma-slice-given-element}.
\end{proof}

\begin{lemma}
\label{lemma-slice-CM}
\begin{reference}
\cite[IV Proposition 17.16.1]{EGA}
\end{reference}
Soit $f : X \to S$ un morphisme de schémas.
Soit $x \in X$ un point d'image $s \in S$.
Supposons que
\begin{enumerate}
\item $f$ soit localement de présentation finie,
\item $f$ soit de Cohen-Macaulay en $x$, et
\item $x$ soit un point fermé de $X_s$.
\end{enumerate}
Alors il existe une immersion régulière $Z \to X$ contenant $x$ telle que
\begin{enumerate}
\item[(a)] $Z \to S$ soit plat et localement de présentation finie,
\item[(b)] $Z \to S$ soit localement quasi-fini, et
\item[(c)] $Z_s = \{x\}$ ensemblistement.
\end{enumerate}
\end{lemma}

\begin{proof}
On peut remplacer, et l'on remplace, $S$ par un voisinage ouvert affine de $s$.
Nous démontrerons le lemme pour $S$ affine par récurrence sur $d = \dim_x(X_s)$.

\medskip\noindent
Cas $d = 0$. Dans ce cas, nous montrons que l'on peut prendre pour $Z$
un voisinage ouvert de $x$. (Notons qu'une immersion ouverte est
une immersion régulière.) En effet, si $d = 0$, alors $X \to S$
est quasi-fini en $x$, voir
Morphismes, lemme \ref{morphisms-lemma-locally-quasi-finite-rel-dimension-0}.
Il existe donc un voisinage ouvert affine $U \subset X$ tel
que $U \to S$ soit quasi-fini, voir
Morphismes, lemme \ref{morphisms-lemma-quasi-finite-points-open}.
Ainsi, après remplacement de $X$ par $U$, la fibre $X_s$ est un ensemble
fini discret. En remplaçant $X$ par un nouveau voisinage ouvert affine
de $X$, on obtient donc $f^{-1}(\{s\}) = \{x\}$ (car la topologie
de $X_s$ est induite par celle de $X$, voir
Schémas, lemme \ref{schemes-lemma-fibre-topological}).
Cela démontre le lemme dans ce cas.

\medskip\noindent
Supposons maintenant $d > 0$. Puisque $x$ est un point fermé de sa
fibre, l'extension $\kappa(x)/\kappa(s)$ est finie (d'après le
théorème des zéros de Hilbert, voir
Morphismes, lemme \ref{morphisms-lemma-closed-point-fibre-locally-finite-type}).
On a donc
$$
\text{depth}(\mathcal{O}_{X_s, x}) = \dim(\mathcal{O}_{X_s, x}) = d > 0
$$
la première égalité provenant de ce que $\mathcal{O}_{X_s, x}$ est de Cohen-Macaulay,
et la seconde de
Morphismes, lemme \ref{morphisms-lemma-dimension-fibre-at-a-point}.
On peut donc appliquer le
lemme \ref{lemma-slice-once}
pour trouver un diagramme
$$
\xymatrix{
D \ar[r] \ar[rrd] & U \ar[r] \ar[rd] & X \ar[d] \\
& & S
}
$$
avec $x \in D$. Notons que
$\mathcal{O}_{D_s, x} = \mathcal{O}_{X_s, x}/(\overline{h})$ pour un certain
non-diviseur de zéro $\overline{h}$, voir
Diviseurs, lemme \ref{divisors-lemma-relative-Cartier}.
Ainsi, $\mathcal{O}_{D_s, x}$ est de Cohen-Macaulay, de dimension
inférieure d'une unité à celle de $\mathcal{O}_{X_s, x}$, voir
Algèbre, lemme \ref{algebra-lemma-reformulate-CM},
par exemple. Le morphisme $D \to S$ est donc plat,
localement de présentation finie et de Cohen-Macaulay en $x$, avec
$\dim_x(D_s) = \dim_x(X_s) - 1 = d - 1$. Par l'hypothèse de récurrence,
on peut trouver une immersion régulière $Z \to D$ ayant les propriétés (a), (b), (c).
Puisque les deux morphismes $Z \to D \to U$ sont des immersions régulières,
$Z \to U$ est aussi une immersion régulière d'après
Diviseurs, lemme \ref{divisors-lemma-composition-regular-immersion}.
Ceci achève la démonstration.
\end{proof}

\begin{lemma}
\label{lemma-qf-fp-flat-neighbourhood-dominates-fppf}
Soit $f : X \to S$ un morphisme plat de schémas
localement de présentation finie. Soit $s \in S$ un point de l'image de $f$.
Alors il existe un diagramme commutatif
$$
\xymatrix{
S' \ar[rr] \ar[rd]_g & & X \ar[ld]^f \\
& S
}
$$
où $g : S' \to S$ est plat, localement de présentation finie,
localement quasi-fini, et où $s \in g(S')$.
\end{lemma}

\begin{proof}
La fibre $X_s$ est non vide par hypothèse. Il existe donc un point
fermé $x \in X_s$ où $f$ est de Cohen-Macaulay, voir le
lemme \ref{lemma-flat-finite-presentation-CM-open}.
Appliquer le
lemme \ref{lemma-slice-CM}
et poser $S' = Z$.
\end{proof}

\noindent
Le lemme suivant montre que les faisceaux pour la topologie fppf
sont exactement les faisceaux pour la topologie
« quasi-finie, plate et de présentation finie ».

\begin{lemma}
\label{lemma-qf-fp-flat-dominates-fppf}
Soit $S$ un schéma. Soit $\mathcal{U} = \{S_i \to S\}_{i \in I}$ un recouvrement
fppf de $S$, voir
Topologies, définition \ref{topologies-definition-fppf-covering}.
Alors il existe un recouvrement fppf $\mathcal{V} = \{T_j \to S\}_{j \in J}$
qui raffine (voir
Sites, définition \ref{sites-definition-morphism-coverings})
$\mathcal{U}$ et tel que chaque morphisme $T_j \to S$ soit localement quasi-fini.
\end{lemma}

\begin{proof}
Pour tout $s \in S$, il existe un $i \in I$ tel que $s$ appartienne à
l'image de $S_i \to S$. D'après le
lemme \ref{lemma-qf-fp-flat-neighbourhood-dominates-fppf},
on peut trouver un morphisme $g_s : T_s \to S$ tel que $s \in g_s(T_s)$,
plat, localement de présentation finie et localement quasi-fini,
et tel que $g_s$ se factorise par $S_i \to S$. Ainsi,
$\{T_s \to S\}$ est le recouvrement recherché de $S$ raffinant $\mathcal{U}$.
\end{proof}

\section{Fibres génériques}
\label{section-generic}

\noindent
Quelques résultats sur les rapports entre les fibres génériques et les
fibres voisines.

\begin{lemma}
\label{lemma-empty-generic-fibre}
Soit $f : X \to Y$ un morphisme de type fini de schémas. Supposons que
$Y$ soit irréductible, de point générique $\eta$. Si $X_\eta = \emptyset$,
alors il existe un ouvert non vide $V \subset Y$ tel que
$X_V = V \times_Y X = \emptyset$.
\end{lemma}

\begin{proof}
Cela résulte immédiatement du résultat plus général
Morphismes,
lemme \ref{morphisms-lemma-quasi-compact-generic-point-not-in-image}.
\end{proof}

\begin{lemma}
\label{lemma-nonempty-generic-fibre}
Soit $f : X \to Y$ un morphisme de type fini de schémas. Supposons que
$Y$ soit irréductible, de point générique $\eta$. Si $X_\eta \not = \emptyset$,
alors il existe un ouvert non vide $V \subset Y$ tel que
$X_V = V \times_Y X \to V$ soit surjectif.
\end{lemma}

\begin{proof}
Cela résulte, en prenant des ouverts affines, de
Algèbre, lemme \ref{algebra-lemma-characterize-image-finite-type}.
(Bien entendu, cela résulte aussi de la platitude générique.)
\end{proof}

\begin{lemma}
\label{lemma-nowhere-dense-generic-fibre}
Soit $f : X \to Y$ un morphisme de type fini de schémas. Supposons que
$Y$ soit irréductible, de point générique $\eta$.
Si $Z \subset X$ est un sous-ensemble fermé tel que $Z_\eta$ soit nulle part dense
dans $X_\eta$, alors il existe un ouvert non vide $V \subset Y$ tel
que $Z_y$ soit nulle part dense dans $X_y$ pour tout $y \in V$.
\end{lemma}

\begin{proof}
Soit $Y' \subset Y$ le réduit de $Y$.
Posons $X' = Y' \times_Y X$ et $Z' = Y' \times_Y Z$.
Comme $Y' \to Y$ est un homéomorphisme universel d'après
Morphismes, lemme \ref{morphisms-lemma-reduction-universal-homeomorphism},
il suffit de démontrer le lemme pour $Z' \subset X' \to Y'$.
On peut donc supposer que $Y$ est intègre, voir
Propriétés, lemme \ref{properties-lemma-characterize-integral}.
D'après
Morphismes, proposition \ref{morphisms-proposition-generic-flatness},
il existe un ouvert affine non vide $V \subset Y$ tel que
$X_V \to V$ et $Z_V \to V$ soient plats et de présentation finie.
Nous affirmons que $V$ convient.
Choisissons $y \in V$. Si $Z_y$ est d'intérieur non vide, alors $Z_y$ contient
un point générique $\xi$ d'une composante irréductible de $X_y$.
Notons que $\eta \leadsto f(\xi)$. Comme $Z_V \to V$ est plat, on peut
choisir une spécialisation $\xi' \leadsto \xi$, $\xi' \in Z$,
avec $f(\xi') = \eta$, voir
Morphismes, lemme \ref{morphisms-lemma-generalizations-lift-flat}.
D'après le
lemme \ref{lemma-flat-finite-presentation-specialization-dimension},
on a
$$
\dim_{\xi'}(Z_\eta) = \dim_{\xi}(Z_y) = \dim_{\xi}(X_y) = \dim_{\xi'}(X_\eta).
$$
Ainsi, une composante irréductible de $Z_\eta$ passant par $\xi'$ est
de dimension $\dim_{\xi'}(X_\eta)$, ce qui contredit l'hypothèse selon laquelle
$Z_\eta$ est nulle part dense dans $X_\eta$, et le résultat est démontré.
\end{proof}

\begin{lemma}
\label{lemma-scheme-theoretically-dense-generic-fibre}
Soit $f : X \to Y$ un morphisme de type fini de schémas.
Supposons que $Y$ soit irréductible, de point générique $\eta$.
Soit $U \subset X$ un sous-schéma ouvert tel que $U_\eta$ soit
schématiquement dense dans $X_\eta$.
Alors il existe un ouvert non vide $V \subset Y$ tel
que $U_y$ soit schématiquement dense dans $X_y$ pour tout $y \in V$.
\end{lemma}

\begin{proof}
Soit $Y' \subset Y$ le réduit de $Y$.
Posons $X' = Y' \times_Y X$ et $U' = Y' \times_Y U$.
Comme $Y' \to Y$ induit une bijection sur les points, et comme
$U' \to U$ et $X' \to X$ induisent des isomorphismes sur les fibres schématiques,
on peut remplacer $Y$ par $Y'$ et $X$ par $X'$.
On peut donc supposer que $Y$ est intègre, voir
Propriétés, lemme \ref{properties-lemma-characterize-integral}.
On peut aussi remplacer $Y$ par un ouvert affine non vide. Autrement dit, on
peut supposer que $Y = \Spec(A)$, où $A$ est un anneau intègre de corps des
fractions $K$.

\medskip\noindent
Puisque $f$ est de type fini, $X$ est quasi-compact.
Écrivons $X = X_1 \cup \ldots \cup X_n$, où les $X_i$ sont des ouverts affines. D'après
Morphismes, définition \ref{morphisms-definition-scheme-theoretically-dense},
$U_i = X_i \cap U$ est un sous-schéma ouvert de $X_i$ tel que
$U_{i, \eta}$ soit schématiquement dense dans $X_{i, \eta}$.
Il suffit donc de démontrer le résultat pour les couples $(X_i, U_i)$ ;
autrement dit, on peut supposer que $X$ est affine.

\medskip\noindent
Écrivons $X = \Spec(B)$. Notons que $B_K$ est noethérien, puisque c'est une
$K$-algèbre de type fini. Ainsi, $U_\eta$ est quasi-compact. On peut donc
trouver un nombre fini d'éléments $g_1, \ldots, g_m \in B$ tels que $D(g_j) \subset U$
et tels que $U_\eta = D(g_1)_\eta \cup \ldots \cup D(g_m)_\eta$.
Le fait que $U_\eta$ soit schématiquement dense dans
$X_\eta$ signifie que $B_K \to \bigoplus_j (B_K)_{g_j}$
est injectif, voir
Morphismes, exemple \ref{morphisms-example-scheme-theoretic-closure}.
D'après
Algèbre, lemme \ref{algebra-lemma-when-injective-covering},
cela équivaut à l'injectivité de
$B_K \to \bigoplus\nolimits_{j = 1, \ldots, m} B_K$,
$b \mapsto (g_1b, \ldots, g_mb)$. Soit $M$ le conoyau de cette
application sur $A$, c'est-à-dire que l'on a une suite exacte
$$
0 \to I \to B \xrightarrow{(g_1, \ldots, g_m)}
\bigoplus\nolimits_{j = 1, \ldots, m} B \to M \to 0
$$
Quitte à remplacer $A$ par $A_h$ pour un certain $h$ non nul, on peut supposer que $B$
est une $A$-algèbre plate et de présentation finie, et que $M$
est plat sur $A$, voir
Algèbre, lemme \ref{algebra-lemma-generic-flatness}.
La platitude de $B$ sur $A$ implique que $B$ est sans torsion comme
$A$-module, voir
Compléments sur l'algèbre, lemme \ref{more-algebra-lemma-flat-torsion-free}.
Ainsi $B \subset B_K$. Par hypothèse, $I_K = 0$, ce qui implique que $I = 0$
(car $I \subset B \subset B_K$ est inclus dans $I_K$). On a donc maintenant
une suite exacte courte
$$
0 \to B \xrightarrow{(g_1, \ldots, g_m)}
\bigoplus\nolimits_{j = 1, \ldots, m} B \to M \to 0
$$
où $M$ est plat sur $A$. Pour tout homomorphisme $A \to \kappa$, où
$\kappa$ est un corps, on obtient donc une suite exacte courte
$$
0 \to B \otimes_A \kappa \xrightarrow{(g_1 \otimes 1, \ldots, g_m \otimes 1)}
\bigoplus\nolimits_{j = 1, \ldots, m} B \otimes_A \kappa \to
M \otimes_A \kappa \to 0
$$
voir
Algèbre, lemme \ref{algebra-lemma-flat-tor-zero}.
En reprenant en sens inverse les arguments ci-dessus,
cela signifie que $\bigcup D(g_j \otimes 1)$ est schématiquement
dense dans $\Spec(B \otimes_A \kappa)$.
Comme $\bigcup D(g_j \otimes 1) = \bigcup D(g_j)_\kappa \subset U_\kappa$,
on en déduit que $U_\kappa$ est schématiquement dense dans $X_\kappa$,
ce qu'il fallait démontrer.
\end{proof}

\noindent
Soient un morphisme de schémas $f : X \to Y$ et
un point $y \in Y$. Rappelons que la fibre $X_y$ est homéomorphe
au sous-ensemble $f^{-1}(\{y\})$ de $X$ muni de la topologie induite, voir
Schémas, lemme \ref{schemes-lemma-fibre-topological}.
Soit un sous-ensemble fermé $T(y) \subset X_y$.
Soit $T$ l'adhérence de $T(y)$ dans $X$.
Munissons $T$ de la structure de schéma réduit induite.
Alors $T$ est un sous-schéma fermé de $X$ tel
que $T_y = T(y)$ ensemblistement. En fait, $T$ est le plus petit
sous-schéma fermé de $X$ ayant cette propriété. Il est donc « sans inconvénient »
de désigner un sous-ensemble fermé de $X_y$ par $T_y$ si on le souhaite.
Dans le lemme suivant, on applique ceci à la fibre générique de $f$.

\begin{lemma}
\label{lemma-cover-generic-fibre-neighbourhood}
Soit $f : X \to Y$ un morphisme de type fini de schémas. Supposons que
$Y$ soit irréductible, de point générique $\eta$. Soit
$X_\eta = Z_{1, \eta} \cup \ldots \cup Z_{n, \eta}$ un recouvrement de
la fibre générique par des sous-ensembles fermés de $X_\eta$.
Soit $Z_i$ l'adhérence de $Z_{i, \eta}$ dans $X$ (voir la discussion ci-dessus).
Alors il existe un ouvert non vide $V \subset Y$ tel
que $X_y = Z_{1, y} \cup \ldots \cup Z_{n, y}$ pour tout $y \in V$.
\end{lemma}

\begin{proof}
Si $Y$ est noethérien, alors $U = X \setminus (Z_1 \cup \ldots \cup Z_n)$
est de type fini sur $Y$, et on peut appliquer directement le
lemme \ref{lemma-empty-generic-fibre}
pour obtenir $U_V = \emptyset$ pour un ouvert non vide $V \subset Y$.
Dans le cas général, raisonnons comme suit. Puisque la question est topologique,
on peut remplacer $Y$ par son réduit. Ainsi $Y$ est intègre, voir
Propriétés, lemme \ref{properties-lemma-characterize-integral}.
Quitte à rétrécir $Y$, on peut supposer que $X \to Y$ est plat, voir
Morphismes, proposition \ref{morphisms-proposition-generic-flatness}.
Dans ce cas, tout point $x$ de $X_y$ est une spécialisation d'un point
$x' \in X_\eta$, d'après
Morphismes, lemme \ref{morphisms-lemma-generalizations-lift-flat}.
Comme les $Z_i$ sont fermés dans $X$ et recouvrent la fibre générique, ceci
implique que $X_y = \bigcup Z_{i, y}$ pour $y \in Y$, comme désiré.
\end{proof}

\noindent
Le lemme suivant affirme que les fibres génériques de morphismes dont la source est
réduite sont réduites.

\begin{lemma}
\label{lemma-reduction-generic-fibre}
Soit $f : X \to Y$ un morphisme de schémas. Soit $\eta \in Y$ un point générique
d'une composante irréductible de $Y$. Alors
$(X_\eta)_{red} = (X_{red})_\eta$.
\end{lemma}

\begin{proof}
Choisissons un voisinage affine $\Spec(A) \subset Y$ de $\eta$.
Choisissons un ouvert affine $\Spec(B) \subset X$ envoyé dans
$\Spec(A)$ par le morphisme $f$. Soit $\mathfrak p \subset A$ l'idéal premier minimal
correspondant à $\eta$. Soit $B_{red}$ le quotient de $B$ par
le nilradical $\sqrt{(0)}$. Le contenu algébrique du lemme est que
$C = B_{red} \otimes_A \kappa(\mathfrak p)$ est réduit.
Notons $I \subset A$ le nilradical, de sorte que $A_{red} = A/I$.
Notons $\mathfrak p_{red} = \mathfrak p/I$ ;
c'est un idéal premier minimal de $A_{red}$ et
$\kappa(\mathfrak p) = \kappa(\mathfrak p_{red})$.
Comme $A \to B_{red}$ et $A \to \kappa(\mathfrak p)$
se factorisent tous deux par $A \to A_{red}$, on a
$C = B_{red} \otimes_{A_{red}} \kappa(\mathfrak p_{red})$.
Or $\kappa(\mathfrak p_{red}) = (A_{red})_{\mathfrak p_{red}}$
est une localisation d'après
Algèbre, lemme \ref{algebra-lemma-minimal-prime-reduced-ring}.
Ainsi $C$ est une localisation de $B_{red}$
(Algèbre, lemme \ref{algebra-lemma-tensor-localization})
et est donc réduit.
\end{proof}

\begin{lemma}
\label{lemma-make-generic-fibre-geometrically-reduced}
Soit $f : X \to Y$ un morphisme de schémas.
Supposons que $Y$ soit irréductible et que $f$ soit de type fini.
Il existe un diagramme
$$
\xymatrix{
X' \ar[d]_{f'} \ar[r]_{g'} & X_V \ar[r] \ar[d] & X \ar[d]^f \\
Y' \ar[r]^g & V \ar[r] & Y
}
$$
où
\begin{enumerate}
\item $V$ est un ouvert non vide de $Y$,
\item $X_V = V \times_Y X$,
\item $g : Y' \to V$ est un homéomorphisme universel fini,
\item $X' = (Y' \times_Y X)_{red} = (Y' \times_V X_V)_{red}$,
\item $g'$ est un homéomorphisme universel fini,
\item $Y'$ est un schéma affine intègre,
\item $f'$ est plat et de présentation finie, et
\item la fibre générique de $f'$ est géométriquement réduite.
\end{enumerate}
\end{lemma}

\begin{proof}
Soit $V = \Spec(A)$ un ouvert affine non vide de $Y$.
Par hypothèse, le nilradical $\mathfrak p = \sqrt{(0)}$ de $A$
est un idéal premier.
Soit $K = \kappa(\mathfrak p)$.
Soit $p$ la caractéristique de $K$ si elle est positive, et $1$
si elle est nulle. D'après
Variétés, lemme \ref{varieties-lemma-finite-extension-geometrically-reduced},
il existe une extension de corps finie purement inséparable
$K'/K$ telle que $(X_{K'})_{red}$ soit géométriquement réduit sur $K'$.
Choisissons des éléments $x_1, \ldots, x_n \in K'$ qui engendrent $K'$ sur
$K$ et tels qu'une puissance d'exposant une puissance de $p$ de $x_i$ appartienne à $A/\mathfrak p$.
Soit $A' \subset K'$ la $A$-sous-algèbre finie de $K'$ engendrée par
$x_1, \ldots, x_n$. Notons que $A'$ est un anneau intègre de corps des fractions $K'$. D'après
Algèbre, lemme \ref{algebra-lemma-p-ring-map},
l'application $A \to A'$ induit un homéomorphisme universel sur les spectres.
Posons $Y' = \Spec(A')$. Posons $X' = (Y' \times_Y X)_{red}$.
La fibre générique de $X' \to Y'$ est $(X_{K'})_{red}$ d'après le
lemme \ref{lemma-reduction-generic-fibre},
et elle est géométriquement réduite par construction.
Notons que $X' \to X_V$ est un homéomorphisme universel fini, car c'est la
composée du morphisme de réduction $X' \to Y' \times_Y X$ (voir
Morphismes, lemme \ref{morphisms-lemma-reduction-universal-homeomorphism})
et du changement de base de $g$.
À ce stade, toutes les propriétés du lemme sont satisfaites, à l'exception
possible de (7). On peut obtenir cette dernière en rétrécissant $Y'$, et donc $V$, voir
Morphismes, proposition \ref{morphisms-proposition-generic-flatness}.
\end{proof}

\begin{lemma}
\label{lemma-make-components-generic-fibre-geometrically-irreducible}
Soit $f : X \to Y$ un morphisme de schémas.
Supposons que $Y$ soit irréductible et que $f$ soit de type fini.
Il existe un diagramme
$$
\xymatrix{
X' \ar[d]_{f'} \ar[r]_{g'} & X_V \ar[r] \ar[d] & X \ar[d]^f \\
Y' \ar[r]^g & V \ar[r] & Y
}
$$
où
\begin{enumerate}
\item $V$ est un ouvert non vide de $Y$,
\item $X_V = V \times_Y X$,
\item $g : Y' \to V$ est fini étale surjectif,
\item $X' = Y' \times_Y X = Y' \times_V X_V$,
\item $g'$ est fini étale surjectif,
\item $Y'$ est un schéma affine irréductible, et
\item toutes les composantes irréductibles de la fibre générique de $f'$
sont géométriquement irréductibles.
\end{enumerate}
\end{lemma}

\begin{proof}
Soit $V = \Spec(A)$ un ouvert affine non vide de $Y$.
Par hypothèse, le nilradical $\mathfrak p = \sqrt{(0)}$ de $A$ est un idéal premier.
Soit $K = \kappa(\mathfrak p)$. D'après
Variétés, lemme
\ref{varieties-lemma-finite-extension-geometrically-irreducible-components},
il existe une extension de corps finie séparable
$K'/K$ telle que toutes les composantes irréductibles de $X_{K'}$ soient
géométriquement irréductibles sur $K'$.
Choisissons un élément $\alpha \in K'$ qui engendre $K'$ sur
$K$, voir
Corps, lemme \ref{fields-lemma-primitive-element}.
Soit $P(T) \in K[T]$ le polynôme minimal de $\alpha$ sur $K$.
Quitte à remplacer $\alpha$ par $f \alpha$ pour un certain
$f \in A$, $f \not \in \mathfrak p$,
on peut supposer qu'il existe un polynôme unitaire
$T^d + a_1T^{d - 1} + \ldots + a_d \in A[T]$ dont l'image est
$P(T) \in K[T]$ par l'application $A[T] \to K[T]$.
Posons $A' = A[T]/(P)$. Alors $A \to A'$ est un homomorphisme d'anneaux fini libre
tel qu'il existe un unique idéal premier $\mathfrak q$ au-dessus de
$\mathfrak p$, que
$K = \kappa(\mathfrak p) \subset \kappa(\mathfrak q) = K'$
est une extension finie séparable, et que $\mathfrak pA'_{\mathfrak q}$
est l'idéal maximal de $A'_{\mathfrak q}$.
Ainsi $g : Y' = \Spec(A') \to V = \Spec(A)$
est étale en $\mathfrak q$, voir
Algèbre, lemme \ref{algebra-lemma-characterize-etale}.
Cela signifie qu'il existe un ouvert $W \subset \Spec(A')$ tel
que $g|_W : W \to \Spec(A)$ soit étale.
Comme $g$ est fini et comme $\mathfrak q$ est le seul point au-dessus de
$\mathfrak p$, on voit que $Z = g(Y' \setminus W)$ est un sous-ensemble fermé de $V$
qui ne contient pas $\mathfrak p$. Ainsi, quitte à remplacer $V$ par un ouvert affine
principal de $V$ qui ne rencontre pas $Z$, on obtient que $g$ est fini
étale.
\end{proof}

\section{Assassins relatifs}
\label{section-assassin}

\begin{lemma}
\label{lemma-relative-assassin-in-neighbourhood}
Soit $f : X \to S$ un morphisme de schémas.
Soit $\mathcal{F}$ un $\mathcal{O}_X$-module quasi-cohérent.
Soit $\xi \in \text{Ass}_{X/S}(\mathcal{F})$ et posons
$Z = \overline{\{\xi\}} \subset X$.
Si $f$ est localement de type fini et si $\mathcal{F}$ est un
$\mathcal{O}_X$-module de type fini, alors il existe un ouvert non vide
$V \subset Z$ tel que, pour tout $s \in f(V)$, les points génériques
de $V_s$ appartiennent à $\text{Ass}_{X/S}(\mathcal{F})$.
\end{lemma}

\begin{proof}
On peut remplacer $S$ par un voisinage ouvert affine de $f(\xi)$
et $X$ par un voisinage ouvert affine de $\xi$. On peut donc supposer
$S = \Spec(A)$, $X = \Spec(B)$ et que $f$ est donné par
l'homomorphisme d'anneaux de type fini $A \to B$, voir
Morphismes, lemme \ref{morphisms-lemma-locally-finite-type-characterize}.
De plus, on peut écrire $\mathcal{F} = \widetilde{M}$ pour un certain
$B$-module de type fini $M$, voir
Propriétés, lemme \ref{properties-lemma-finite-type-module}.
Soit $\mathfrak q \subset B$ l'idéal premier correspondant à $\xi$, et
soit $\mathfrak p \subset A$ l'idéal premier correspondant de $A$.
Par hypothèse, $\mathfrak q \in \text{Ass}_B(M \otimes_A \kappa(\mathfrak p))$,
voir
Algèbre, remarque \ref{algebra-remark-bourbaki}
et
Diviseurs, lemme \ref{divisors-lemma-associated-affine-open}.
Avec ces notations, $Z = V(\mathfrak q) \subset \Spec(B)$.
En particulier, $f(Z) \subset V(\mathfrak p)$. Il suffit donc clairement de
démontrer le lemme après avoir remplacé $A$, $B$ et $M$ par
$A/\mathfrak pA$, $B/\mathfrak pB$ et $M/\mathfrak pM$.
Autrement dit, on peut supposer que $A$ est un anneau intègre de corps des fractions
$K$ et que $\mathfrak q \subset B$ est un idéal premier associé de $M \otimes_A K$.

\medskip\noindent
On peut maintenant utiliser la platitude générique. En effet, d'après
Algèbre, lemme \ref{algebra-lemma-generic-flatness},
il existe un élément non nul $g \in A$ tel que $M_g$ soit plat
comme $A_g$-module. Après avoir remplacé $A$ par $A_g$, on peut supposer que
$M$ est plat comme $A$-module.

\medskip\noindent
Dans ce cas, d'après
Algèbre, lemme \ref{algebra-lemma-post-bourbaki},
$\mathfrak q$ est aussi un idéal premier associé de $M$.
On obtient donc une application injective de $B$-modules $B/\mathfrak q \to M$.
Soit $Q$ son conoyau ; on obtient ainsi une suite exacte courte
$$
0 \to B/\mathfrak q \to M \to Q \to 0
$$
de $B$-modules de type fini. En appliquant une nouvelle fois la platitude générique,
Algèbre, lemme \ref{algebra-lemma-generic-flatness},
cette fois au $B$-module $Q$, on peut supposer que $Q$
est un $A$-module plat. En particulier, on peut supposer que la suite exacte
courte ci-dessus est universellement injective, voir
Algèbre, lemme \ref{algebra-lemma-flat-tor-zero}.
Dans cette situation,
$(B/\mathfrak q) \otimes_A \kappa(\mathfrak p')
\subset M \otimes_A \kappa(\mathfrak p')$
pour tout idéal premier $\mathfrak p'$ de $A$. Le lemme en résulte, car tout idéal
premier minimal $\mathfrak q'$ du support de
$(B/\mathfrak q) \otimes_A \kappa(\mathfrak p')$
est un idéal premier associé de $(B/\mathfrak q) \otimes_A \kappa(\mathfrak p')$ d'après
Diviseurs, lemme \ref{divisors-lemma-minimal-support-in-ass}.
\end{proof}

\begin{lemma}
\label{lemma-bad-case}
Soit $f : X \to Y$ un morphisme de schémas. Soit $\mathcal{F}$ un
$\mathcal{O}_X$-module quasi-cohérent. Soit $U \subset X$ un sous-schéma
ouvert. Supposons que
\begin{enumerate}
\item $f$ soit de type fini,
\item $\mathcal{F}$ soit de type fini,
\item $Y$ soit irréductible, de point générique $\eta$, et
\item $\text{Ass}_{X_\eta}(\mathcal{F}_\eta)$ ne soit pas contenu dans $U_\eta$.
\end{enumerate}
Alors il existe un sous-schéma ouvert non vide $V \subset Y$ tel que,
pour tout $y \in V$, l'ensemble $\text{Ass}_{X_y}(\mathcal{F}_y)$ ne soit pas
contenu dans $U_y$.
\end{lemma}

\begin{proof}
Soit $Z \subset X$ le support schématique de $\mathcal{F}$, voir
Morphismes, définition \ref{morphisms-definition-scheme-theoretic-support}.
Alors $Z_\eta$ est le support schématique de $\mathcal{F}_\eta$
(Morphismes, lemme \ref{morphisms-lemma-flat-pullback-support}).
Ainsi, les points génériques des composantes irréductibles de $Z_\eta$
appartiennent à $\text{Ass}_{X_\eta}(\mathcal{F}_\eta)$ d'après
Diviseurs, lemme \ref{divisors-lemma-minimal-support-in-ass}.
On voit donc que $Z_\eta \cap U_\eta = \emptyset$.
Ainsi $T = Z \setminus U$ est un sous-ensemble fermé de $Z$ tel que $T_\eta = \emptyset$.
Si l'on munit $T$ de la structure de schéma réduit induite, alors
$T \to Y$ est un morphisme de type fini. D'après le
lemme \ref{lemma-empty-generic-fibre},
il existe un ouvert non vide $V \subset Y$ tel que $T_V = \emptyset$.
Alors $V$ convient.
\end{proof}

\begin{lemma}
\label{lemma-good-case}
Soit $f : X \to Y$ un morphisme de schémas. Soit $\mathcal{F}$ un
$\mathcal{O}_X$-module quasi-cohérent. Soit $U \subset X$ un sous-schéma
ouvert. Supposons que
\begin{enumerate}
\item $f$ soit de type fini,
\item $\mathcal{F}$ soit de type fini,
\item $Y$ soit irréductible, de point générique $\eta$, et
\item $\text{Ass}_{X_\eta}(\mathcal{F}_\eta) \subset U_\eta$.
\end{enumerate}
Alors il existe un sous-schéma ouvert non vide $V \subset Y$ tel que,
pour tout $y \in V$, on ait $\text{Ass}_{X_y}(\mathcal{F}_y) \subset U_y$.
\end{lemma}

\begin{proof}
(Cette démonstration est identique à celle du
lemme \ref{lemma-scheme-theoretically-dense-generic-fibre}.
Nous invitons le lecteur à lire d'abord cette dernière.)
Puisque l'énoncé porte sur les fibres, on peut clairement remplacer
$Y$ par son réduit. On peut donc supposer que $Y$ est intègre, voir
Propriétés, lemme \ref{properties-lemma-characterize-integral}.
On peut aussi supposer que $Y = \Spec(A)$ est affine. Alors $A$
est un anneau intègre de corps des fractions $K$.

\medskip\noindent
Puisque $f$ est de type fini, $X$ est quasi-compact.
Écrivons $X = X_1 \cup \ldots \cup X_n$, où les $X_i$ sont des ouverts affines,
et posons $\mathcal{F}_i = \mathcal{F}|_{X_i}$. Par
hypothèse, la fibre générique de $U_i = X_i \cap U$ contient
$\text{Ass}_{X_{i, \eta}}(\mathcal{F}_{i, \eta})$.
Il suffit donc de démontrer le résultat pour les triplets
$(X_i, \mathcal{F}_i, U_i)$ ;
autrement dit, on peut supposer que $X$ est affine.

\medskip\noindent
Écrivons $X = \Spec(B)$. Soit $N$ un $B$-module de type fini tel que
$\mathcal{F} = \widetilde{N}$.
Notons que $B_K$ est noethérien, puisque c'est une
$K$-algèbre de type fini. Ainsi $U_\eta$ est quasi-compact. On peut donc
trouver un nombre fini d'éléments $g_1, \ldots, g_m \in B$ tels que $D(g_j) \subset U$
et tels que $U_\eta = D(g_1)_\eta \cup \ldots \cup D(g_m)_\eta$.
Puisque $\text{Ass}_{X_\eta}(\mathcal{F}_\eta) \subset U_\eta$,
l'application $N_K \to \bigoplus_j (N_K)_{g_j}$ est injective. D'après
Algèbre, lemme \ref{algebra-lemma-when-injective-covering},
cela équivaut à l'injectivité de
$N_K \to \bigoplus\nolimits_{j = 1, \ldots, m} N_K$,
$n \mapsto (g_1n, \ldots, g_mn)$. Soient $I$ et $M$ le noyau et le
conoyau de cette application sur $A$, c'est-à-dire que l'on a une suite exacte
$$
0 \to I \to N \xrightarrow{(g_1, \ldots, g_m)}
\bigoplus\nolimits_{j = 1, \ldots, m} N \to M \to 0
$$
Quitte à remplacer $A$ par $A_h$ pour un certain $h$ non nul, on peut supposer que
$B$ est une $A$-algèbre plate et de présentation finie et que
$M$ et $N$ sont tous deux plats sur $A$, voir
Algèbre, lemme \ref{algebra-lemma-generic-flatness}.
La platitude de $N$ sur $A$ implique que $N$ est sans torsion comme
$A$-module, voir
Compléments sur l'algèbre, lemme \ref{more-algebra-lemma-flat-torsion-free}.
Ainsi $N \subset N_K$. Par construction, $I_K = 0$, ce qui implique que $I = 0$
(car $I \subset N \subset N_K$ est inclus dans $I_K$). On a donc maintenant
une suite exacte courte
$$
0 \to N \xrightarrow{(g_1, \ldots, g_m)}
\bigoplus\nolimits_{j = 1, \ldots, m} N \to M \to 0
$$
où $M$ est plat sur $A$. Pour tout homomorphisme $A \to \kappa$, où
$\kappa$ est un corps, on obtient donc une suite exacte courte
$$
0 \to N \otimes_A \kappa \xrightarrow{(g_1 \otimes 1, \ldots, g_m \otimes 1)}
\bigoplus\nolimits_{j = 1, \ldots, m} N \otimes_A \kappa \to
M \otimes_A \kappa \to 0
$$
voir
Algèbre, lemme \ref{algebra-lemma-flat-tor-zero}.
En reprenant en sens inverse les arguments ci-dessus,
cela signifie que $\bigcup D(g_j \otimes 1)$ contient
$\text{Ass}_{B \otimes_A \kappa}(N \otimes_A \kappa)$.
Comme $\bigcup D(g_j \otimes 1) = \bigcup D(g_j)_\kappa \subset U_\kappa$,
on en déduit que $U_\kappa$ contient
$\text{Ass}_{X \otimes \kappa}(\mathcal{F} \otimes \kappa)$,
ce qu'il fallait démontrer.
\end{proof}

\begin{lemma}
\label{lemma-base-change-assassin-in-U}
Soit $f : X \to S$ un morphisme localement de type fini.
Soit $\mathcal{F}$ un $\mathcal{O}_X$-module quasi-cohérent
de type fini. Soit $U \subset X$ un sous-schéma ouvert.
Soit $g : S' \to S$ un morphisme de schémas, soit
$f' : X' = X_{S'} \to S'$ le changement de base de $f$,
soit $g' : X' \to X$ la projection, posons
$\mathcal{F}' = (g')^*\mathcal{F}$, et posons
$U' = (g')^{-1}(U)$. Enfin, soit $s' \in S'$ d'image $s = g(s')$.
Dans ce cas,
$$
\text{Ass}_{X_s}(\mathcal{F}_s) \subset U_s
\Leftrightarrow
\text{Ass}_{X'_{s'}}(\mathcal{F}'_{s'}) \subset U'_{s'}.
$$
\end{lemma}

\begin{proof}
Cela résulte immédiatement de
Diviseurs, lemme \ref{divisors-lemma-base-change-relative-assassin}.
Voir aussi
Diviseurs, remarque \ref{divisors-remark-base-change-relative-assassin}.
\end{proof}

\begin{lemma}
\label{lemma-relative-assassin-constructible}
Soit $f : X \to Y$ un morphisme de présentation finie.
Soit $\mathcal{F}$ un $\mathcal{O}_X$-module quasi-cohérent
de présentation finie. Soit $U \subset X$ un sous-schéma ouvert
tel que $U \to Y$ soit quasi-compact. Alors l'ensemble
$$
E = \{y \in Y \mid \text{Ass}_{X_y}(\mathcal{F}_y) \subset U_y\}
$$
est localement constructible dans $Y$.
\end{lemma}

\begin{proof}
Soit $y \in Y$. Il faut montrer qu'il existe un voisinage ouvert
$V$ de $y$ dans $Y$ tel que $E \cap V$ soit constructible dans $V$. On peut donc
supposer que $Y$ est affine. Écrivons $Y = \Spec(A)$ et
$A = \colim A_i$ comme limite inductive filtrante de
$\mathbf{Z}$-algèbres de type fini. D'après
Limites, lemme \ref{limits-lemma-descend-finite-presentation},
on peut trouver un $i$ et un morphisme $f_i : X_i \to \Spec(A_i)$ de
présentation finie dont le changement de base à $Y$ redonne $f$.
Quitte à augmenter $i$, on peut supposer qu'il existe un
$\mathcal{O}_{X_i}$-module quasi-cohérent $\mathcal{F}_i$ de présentation
finie dont l'image inverse sur $X$ est isomorphe à $\mathcal{F}$, voir
Limites, lemme \ref{limits-lemma-descend-modules-finite-presentation}.
Quitte à augmenter encore une fois $i$, on peut supposer qu'il existe
un sous-schéma ouvert $U_i \subset X_i$ dont l'image inverse dans $X$
est $U$, voir
Limites, lemme \ref{limits-lemma-descend-opens}.
D'après le
lemme \ref{lemma-base-change-assassin-in-U},
il suffit de démontrer le lemme pour $f_i$. On se ramène donc au
cas où $Y$ est le spectre d'un anneau noethérien.

\medskip\noindent
Nous allons utiliser le critère de
Topologie, lemme \ref{topology-lemma-characterize-constructible-Noetherian},
pour démontrer que $E$ est constructible lorsque $Y$ est un schéma noethérien.
Pour le voir, soit $Z \subset Y$ un sous-schéma fermé irréductible.
Il faut montrer que $E \cap Z$ contient un ouvert non vide
ou n'est pas dense dans $Z$. Cela résulte des
lemmes \ref{lemma-bad-case} et
\ref{lemma-good-case},
appliqués au changement de base $(X, \mathcal{F}, U) \times_Y Z$ sur $Z$.
\end{proof}

\section{Fibres réduites}
\label{section-reduced}

\begin{lemma}
\label{lemma-nonreduced-in-neighbourhood}
Soit $f : X \to Y$ un morphisme de schémas. Supposons que $Y$ soit irréductible, de
point générique $\eta$, et que $f$ soit de type fini. Si $X_\eta$ n'est pas réduit,
alors il existe un ouvert non vide $V \subset Y$
tel que, pour tout $y \in V$, la fibre $X_y$ ne soit pas réduite.
\end{lemma}

\begin{proof}
Soit $Y' \subset Y$ le réduit de $Y$. Soit $X' \to Y'$
le changement de base de $f$. Notons que $Y' \to Y$
induit une bijection sur les points et que $X' \to X$ identifie les fibres.
On peut donc supposer que $Y'$ est réduit, c'est-à-dire intègre, voir
Propriétés, lemme \ref{properties-lemma-characterize-integral}.
On peut aussi remplacer $Y$ par un ouvert affine. On peut donc supposer que
$Y = \Spec(A)$, où $A$ est un anneau intègre. Notons $K$ le
corps des fractions de $A$.
Choisissons un ouvert affine $\Spec(B) = U \subset X$ et une section
$h_\eta \in \Gamma(U_\eta, \mathcal{O}_{U_\eta}) = B_K$
qui soit non nulle et nilpotente.
Quitte à rétrécir $Y$, on peut supposer que $h$ provient de
$h \in \Gamma(U, \mathcal{O}_U) = B$. Quitte à rétrécir encore un peu $Y$,
on peut supposer que $h$ est nilpotent. Soit
$I = \{b \in B \mid hb = 0\}$ l'annulateur de $h$.
Alors $C = B/I$ est une $A$-algèbre de type fini dont la fibre générique
$(B/I)_K$ est non nulle (car $h_\eta \not = 0$). Appliquons
la platitude générique à $A \to C$ et $A \to B/hB$, voir
Algèbre, lemme \ref{algebra-lemma-generic-flatness} ;
on obtient un $g \in A$, $g \not = 0$ tel que $C_g$ soit libre comme
$A_g$-module et que $(B/hB)_g$ soit plat comme $A_g$-module.
Remplaçons $Y$ par $D(g) \subset Y$. On a maintenant la suite exacte courte
$$
0 \to C \to B \to B/hB \to 0.
$$
où $B/hB$ est plat sur $A$ et où $C$ est un $A$-module libre non nul.
Il en résulte que, pour tout homomorphisme $A \to \kappa$ vers un corps,
l'anneau $C \otimes_A \kappa$ est non nul et la suite
$$
0 \to C \otimes_A \kappa \to B \otimes_A \kappa \to
B/hB \otimes_A \kappa \to 0
$$
est exacte, voir
Algèbre, lemme \ref{algebra-lemma-flat-tor-zero}.
Notons que
$B/hB \otimes_A \kappa = (B \otimes_A \kappa) / h(B \otimes_A \kappa)$
par exacte à droite du produit tensoriel. On en conclut que
la multiplication par $h$ n'est pas nulle sur $B \otimes_A \kappa$.
Cela signifie clairement que, pour tout point $y \in Y$, l'élément $h$
se restreint en un élément non nul de $U_y$, d'où $X_y$ n'est pas réduit.
\end{proof}

\begin{lemma}
\label{lemma-base-change-fibres-geometrically-reduced}
Soit $f : X \to Y$ un morphisme de schémas.
Soit $g : Y' \to Y$ un morphisme quelconque, et notons
$f' : X' \to Y'$ le changement de base de $f$.
Alors
\begin{align*}
\{y' \in Y' \mid X'_{y'}\text{ est géométriquement réduit}\} \\
= g^{-1}(\{y \in Y \mid X_y\text{ est géométriquement réduit}\}).
\end{align*}
\end{lemma}

\begin{proof}
Cela revient à dire que, pour $y' \in Y'$ d'image
$y \in Y$, la fibre $X'_{y'} = X_y \times_y y'$ est géométriquement
réduite sur $\kappa(y')$ si et seulement si $X_y$ est géométriquement
réduit sur $\kappa(y)$. Cela résulte de
Variétés, lemme \ref{varieties-lemma-geometrically-reduced-upstairs}.
\end{proof}

\begin{lemma}
\label{lemma-not-geometrically-reduced-in-neighbourhood}
Soit $f : X \to Y$ un morphisme de schémas. Supposons que $Y$ soit irréductible, de
point générique $\eta$, et que $f$ soit de type fini. Si $X_\eta$ n'est pas
géométriquement réduit, alors il existe un ouvert non vide $V \subset Y$
tel que, pour tout $y \in V$, la fibre $X_y$ ne soit pas géométriquement réduite.
\end{lemma}

\begin{proof}
Appliquons le
lemme \ref{lemma-make-generic-fibre-geometrically-reduced}
pour obtenir
$$
\xymatrix{
X' \ar[d]_{f'} \ar[r]_{g'} & X_V \ar[r] \ar[d] & X \ar[d]^f \\
Y' \ar[r]^g & V \ar[r] & Y
}
$$
avec toutes les propriétés mentionnées dans ce lemme.
Soit $\eta'$ le point générique de $Y'$.
Considérons le morphisme $X' \to X_{Y'}$ (qui est le morphisme de
réduction) et le morphisme de fibres génériques qui en résulte,
$X'_{\eta'} \to X_{\eta'}$.
Puisque $X'_{\eta'}$ est géométriquement réduit, tandis que $X_\eta$
ne l'est pas, ce morphisme ne peut être un isomorphisme, voir
Variétés, lemme \ref{varieties-lemma-geometrically-reduced-upstairs}.
Ainsi $X_{\eta'}$ n'est pas réduit. D'après le
lemme \ref{lemma-nonreduced-in-neighbourhood},
les fibres de $X_{Y'} \to Y'$ ne sont donc pas réduites en tout point $y' \in V'$
d'un ouvert non vide $V' \subset Y'$. Puisque $g : Y' \to V$ est un homéomorphisme, le
lemme \ref{lemma-base-change-fibres-geometrically-reduced}
montre que $g(V')$ est l'ouvert recherché.
\end{proof}

\begin{lemma}
\label{lemma-geometrically-reduced-generic-fibre}
Soit $f : X \to Y$ un morphisme de schémas.
Supposons que
\begin{enumerate}
\item $Y$ soit irréductible, de point générique $\eta$,
\item $X_\eta$ soit géométriquement réduit, et
\item $f$ soit de type fini.
\end{enumerate}
Alors il existe un sous-schéma ouvert non vide $V \subset Y$
tel que $X_V \to V$ ait des fibres géométriquement réduites.
\end{lemma}

\begin{proof}
Soit $Y' \subset Y$ le réduit de $Y$. Soit $X' \to Y'$
le changement de base de $f$. Notons que $Y' \to Y$
induit une bijection sur les points et que $X' \to X$ identifie les fibres.
On peut donc supposer que $Y'$ est réduit, c'est-à-dire intègre, voir
Propriétés, lemme \ref{properties-lemma-characterize-integral}.
On peut aussi remplacer $Y$ par un ouvert affine. On peut donc supposer que
$Y = \Spec(A)$, où $A$ est un anneau intègre. Notons $K$ le
corps des fractions de $A$. Quitte à rétrécir un peu $Y$, on peut également supposer que
$X \to Y$ est plat et de présentation finie, voir
Morphismes, proposition \ref{morphisms-proposition-generic-flatness}.

\medskip\noindent
Puisque $X_\eta$ est géométriquement réduit, il existe un sous-ensemble ouvert dense
$V \subset X_\eta$ tel que $V \to \Spec(K)$ soit lisse, voir
Variétés, lemme \ref{varieties-lemma-geometrically-reduced-dense-smooth-open}.
Soit $U \subset X$ l'ensemble des points où $f$ est lisse. D'après
Morphismes, lemme \ref{morphisms-lemma-set-points-where-fibres-smooth},
on a $V \subset U_\eta$. Ainsi, la fibre générique de $U$ est dense
dans la fibre générique de $X$. Puisque $X_\eta$ est réduit, il en résulte
que $U_\eta$ est schématiquement dense dans $X_\eta$, voir
Morphismes, lemme \ref{morphisms-lemma-reduced-scheme-theoretically-dense}.
Notons que, puisque $U \to Y$ est lisse, toutes les fibres de $U \to Y$
sont géométriquement réduites. Il suffit donc de montrer que, après avoir
rétréci $Y$, pour tout $y \in Y$, le schéma $U_y$ est schématiquement
dense dans $X_y$, voir
Morphismes, lemme \ref{morphisms-lemma-reduced-subscheme-closure}.
Cela résulte du
lemme \ref{lemma-scheme-theoretically-dense-generic-fibre}.
\end{proof}

\begin{lemma}
\label{lemma-geometrically-reduced-constructible}
Soit $f : X \to Y$ un morphisme quasi-compact et
localement de présentation finie. Alors l'ensemble
$$
E = \{y \in Y \mid X_y\text{ est géométriquement réduit}\}
$$
est localement constructible dans $Y$.
\end{lemma}

\begin{proof}
Soit $y \in Y$. Nous devons montrer qu'il existe un voisinage ouvert
$V$ de $y$ dans $Y$ tel que $E \cap V$ soit constructible dans $V$. On peut donc
supposer que $Y$ est affine. Alors $X$ est quasi-compact.
Choisissons un recouvrement ouvert affine fini $X = U_1 \cup \ldots \cup U_n$.
Les fibres des $U_i \to Y$ en $y$ forment alors un recouvrement ouvert affine
de la fibre de $X \to Y$ en $y$. On peut donc aussi supposer que $X$ est affine.
Écrivons $Y = \Spec(A)$.
Écrivons $A = \colim A_i$ comme limite inductive d'algèbres
de type fini sur $\mathbf{Z}$. D'après
Limites, lemme \ref{limits-lemma-descend-finite-presentation},
on peut trouver un $i$ et un morphisme $f_i : X_i \to \Spec(A_i)$ de
présentation finie dont le changement de base à $Y$ redonne $f$. D'après le
lemme \ref{lemma-base-change-fibres-geometrically-reduced},
il suffit de démontrer le lemme pour $f_i$. On se ramène ainsi au
cas où $Y$ est le spectre d'un anneau noethérien.

\medskip\noindent
Nous utiliserons le critère du
Topologie, lemme \ref{topology-lemma-characterize-constructible-Noetherian},
pour montrer que $E$ est constructible lorsque $Y$ est un schéma noethérien.
Pour cela, soit $Z \subset Y$ un sous-schéma fermé irréductible
de point générique $\xi$.
Nous devons montrer que $E \cap Z$ contient un ouvert non vide
ou n'est pas dense dans $Z$. Si $X_\xi$ est géométriquement réduit, alors le
lemme \ref{lemma-geometrically-reduced-generic-fibre}
(appliqué au morphisme $X_Z \to Z$)
implique que toutes les fibres $X_y$ sont géométriquement réduites
pour un ouvert non vide $V \subset Z$.
Si $X_\xi$ n'est pas géométriquement réduit, alors le
lemme \ref{lemma-not-geometrically-reduced-in-neighbourhood}
(appliqué au morphisme $X_Z \to Z$)
implique que toutes les fibres $X_y$ ne sont pas géométriquement réduites
pour un ouvert non vide $V \subset Z$. Cela conclut.
\end{proof}

\begin{lemma}
\label{lemma-proper-flat-over-dvr-reduced-fibre}
Soit $f : X \to \Spec(R)$ un morphisme propre, où $R$ est un
anneau de valuation discrète. Supposons que toute composante irréductible
de $X$ domine $\Spec(R)$ (par exemple si $f$ est plat) et
que la fibre spéciale soit réduite. Alors
$X$ et la fibre générique $X_\eta$ sont réduits.
\end{lemma}

\begin{proof}
(Le fait qu'un morphisme plat satisfasse l'hypothèse sur les composantes
irréductibles résulte par exemple du
Diviseurs, lemme \ref{divisors-lemma-bourbaki},
ou de la descente des idéaux premiers pour les homomorphismes d'anneaux plats.)
Supposons que la fibre spéciale $X_s$ soit réduite.
Soit $x \in X$ un point quelconque, et montrons que $\mathcal{O}_{X, x}$
est réduit ; cela montrera que $X$ et $X_\eta$ sont réduits.
Soit $x \leadsto x'$ une spécialisation telle que $x'$
appartienne à la fibre spéciale ; une telle spécialisation existe
puisqu'un morphisme propre est fermé. Considérons l'anneau local
$A = \mathcal{O}_{X, x'}$. Alors $\mathcal{O}_{X, x}$
est un localisé de $A$, de sorte qu'il suffit de montrer que
$A$ est réduit. Soit $\pi \in R$ une uniformisante.
Si $a \in A$ est non nul, il existe un $n \geq 0$ et un élément
$a' \in A$ tels que $a = \pi^n a'$ et $a' \not \in \pi A$.
Cela résulte du théorème d'intersection de Krull
(Algèbre, lemme \ref{algebra-lemma-intersect-powers-ideal-module-zero}).
Si $a$ est nilpotent, alors $a$ appartient à tous les idéaux premiers minimaux de $A$,
et il en va donc de même de $a'$, car $\pi$ n'appartient à aucun
des idéaux premiers minimaux de $A$ en vertu de notre hypothèse sur les composantes
irréductibles de $X$. Ainsi $a'$ est nilpotent.
Mais $a'$ s'envoie sur un élément non nul de l'anneau réduit
$A/\pi A = \mathcal{O}_{X_s, x'}$.
C'est une contradiction à moins que $A$ ne soit réduit, ce qui
est ce que nous voulions montrer.
\end{proof}

\begin{lemma}
\label{lemma-geometrically-reduced-open}
Soit $f : X \to Y$ un morphisme plat et propre de présentation finie.
Alors l'ensemble $\{y \in Y \mid X_y\text{ est géométriquement réduit}\}$
est ouvert dans $Y$.
\end{lemma}

\begin{proof}
On peut supposer que $Y$ est affine. Alors $Y$ est une limite projective filtrante
de schémas affines de type fini sur $\mathbf{Z}$.
On peut donc supposer que $X \to Y$ est le
changement de base de $X_0 \to Y_0$, où $Y_0$ est le spectre d'une algèbre de type
fini sur $\mathbf{Z}$ et $X_0 \to Y_0$ est plat et propre.
Voir Limites, lemme \ref{limits-lemma-descend-finite-presentation},
\ref{limits-lemma-descend-flat-finite-presentation}, et
\ref{limits-lemma-eventually-proper}. Puisque la formation de
l'ensemble des points où les fibres sont géométriquement
réduites commute au changement de base
(lemme \ref{lemma-base-change-fibres-geometrically-reduced}),
on peut supposer que la base est noethérienne.

\medskip\noindent
Supposons que $Y$ soit noethérien. L'ensemble est constructible d'après le
lemme \ref{lemma-geometrically-reduced-constructible}.
Il suffit donc de montrer que l'ensemble est stable par générisation
(Topologie, lemme \ref{topology-lemma-characterize-closed-Noetherian}). D'après
Propriétés, lemme \ref{properties-lemma-locally-Noetherian-specialization-dvr},
on se ramène au cas où $Y = \Spec(R)$, où $R$ est un anneau de
valuation discrète, et où la fibre fermée $X_y$ est géométriquement
réduite. Il faut montrer que la fibre générique $X_\eta$ est géométriquement réduite.

\medskip\noindent
Sinon, il existe une extension finie $L$ du corps des fractions
de $R$ telle que $X_L$ ne soit pas réduit, voir
Variétés, lemme \ref{varieties-lemma-geometrically-reduced}.
Il existe un anneau de valuation discrète
$R' \subset L$ de corps des fractions $L$ et dominant $R$, voir
Algèbre, lemme \ref{algebra-lemma-integral-closure-Dedekind}.
Après avoir remplacé $R$ par $R'$, on se ramène au
lemme \ref{lemma-proper-flat-over-dvr-reduced-fibre}.
\end{proof}

\section{Composantes irréductibles des fibres}
\label{section-irreducible}

\begin{lemma}
\label{lemma-irreducible-components-in-neighbourhood}
Soit $f : X \to Y$ un morphisme de schémas. Supposons $Y$ irréductible de
point générique $\eta$ et $f$ de type fini. Si $X_\eta$ a $n$
composantes irréductibles, alors il existe un ouvert non vide $V \subset Y$
tel que, pour tout $y \in V$, la fibre $X_y$ ait au moins $n$
composantes irréductibles.
\end{lemma}

\begin{proof}
La question étant purement topologique, on peut remplacer $X$ et $Y$ par
leurs réduits. En particulier, cela implique que $Y$ est intègre, voir
Propriétés, lemme \ref{properties-lemma-characterize-integral}.
Soit $X_\eta = X_{1, \eta} \cup \ldots \cup X_{n, \eta}$
la décomposition de $X_\eta$ en composantes irréductibles.
Soit $X_i \subset X$ le sous-schéma fermé réduit dont la fibre
générique est $X_{i, \eta}$. Notons que $Z_{i, j} = X_i \cap X_j$
est un sous-ensemble fermé de $X_i$ dont la fibre générique $Z_{i, j, \eta}$
est nulle part dense dans $X_{i, \eta}$. Après avoir rétréci $Y$, on peut donc
supposer que $Z_{i, j, y}$
est nulle part dense dans $X_{i, y}$ pour tout $y \in Y$, voir
lemme \ref{lemma-nowhere-dense-generic-fibre}.
Après avoir encore rétréci $Y$, on peut supposer que
$X_y = \bigcup X_{i, y}$ pour $y \in Y$, voir
lemme \ref{lemma-cover-generic-fibre-neighbourhood}.
De plus, après avoir rétréci $Y$, on peut supposer que chaque $X_i \to Y$
est plat et de présentation finie, voir
Morphismes, proposition \ref{morphisms-proposition-generic-flatness}.
Les morphismes $X_i \to Y$ sont ouverts, voir
Morphismes, lemme \ref{morphisms-lemma-fppf-open}.
Il existe donc un voisinage ouvert $V$ de $\eta$ qui est contenu
dans $f(X_i)$ pour chaque $i$.
Pour tout $y \in V$, les schémas $X_{i, y}$ sont des
sous-ensembles fermés non vides de $X_y$, on a $X_y = \bigcup X_{i, y}$,
et les intersections $Z_{i, j, y} = X_{i, y} \cap X_{j, y}$
ne sont pas denses dans $X_{i, y}$. Cela implique clairement que
$X_y$ a au moins $n$ composantes irréductibles.
\end{proof}

\begin{lemma}
\label{lemma-base-change-fibres-geometrically-irreducible}
Soit $f : X \to Y$ un morphisme de schémas.
Soit $g : Y' \to Y$ un morphisme quelconque, et notons
$f' : X' \to Y'$ le changement de base de $f$.
Alors
\begin{align*}
\{y' \in Y' \mid X'_{y'}\text{ est géométriquement irréductible}\} \\
= g^{-1}(\{y \in Y \mid X_y\text{ est géométriquement irréductible}\}).
\end{align*}
\end{lemma}

\begin{proof}
Cela se ramène à l'énoncé suivant : pour $y' \in Y'$ d'image
$y \in Y$, la fibre $X'_{y'} = X_y \times_y y'$ est géométriquement
irréductible sur $\kappa(y')$ si et seulement si $X_y$ est géométriquement
irréductible sur $\kappa(y)$. Cela résulte de
Variétés,
lemme \ref{varieties-lemma-geometrically-irreducible-check-after-extension}.
\end{proof}

\begin{lemma}
\label{lemma-base-change-fibres-nr-geometrically-irreducible-components}
Soit $f : X \to Y$ un morphisme de schémas. Soit
$$
n_{X/Y} : Y \to \{0, 1, 2, 3, \ldots, \infty\}
$$
la fonction qui associe à $y \in Y$ le nombre de composantes irréductibles
de $(X_y)_K$, où $K$ est une extension séparablement close
de $\kappa(y)$. Cette fonction est bien définie et, si $g : Y' \to Y$ est un morphisme,
alors
$$
n_{X'/Y'} = n_{X/Y} \circ g
$$
où $X' \to Y'$ est le changement de base de $f$.
\end{lemma}

\begin{proof}
Supposons que $y' \in Y'$ ait pour image $y \in Y$.
Supposons que $K \supset \kappa(y)$ et $K' \supset \kappa(y')$ soient des extensions
séparablement closes. On peut alors choisir un diagramme commutatif
$$
\xymatrix{
K \ar[r] & K'' & K' \ar[l] \\
\kappa(y) \ar[u] \ar[rr] & & \kappa(y') \ar[u]
}
$$
de corps. Le résultat s'ensuit, car les morphismes de schémas
$$
\xymatrix{
(X'_{y'})_{K'} &
(X'_{y'})_{K''} = (X_y)_{K''} \ar[l] \ar[r] &
(X_y)_K
}
$$
induisent des bijections entre les composantes irréductibles, voir
Variétés,
lemme \ref{varieties-lemma-separably-closed-field-irreducible-components}.
\end{proof}

\begin{lemma}
\label{lemma-irreducible-polynomial-over-domain}
Soit $A$ un anneau intègre de corps des fractions $K$.
Soit $P \in A[x_1, \ldots, x_n]$.
Notons $\overline{K}$ la clôture algébrique de $K$.
Supposons que $P$ soit irréductible dans $\overline{K}[x_1, \ldots, x_n]$.
Alors il existe un $f \in A$ tel que
$P^\varphi \in \kappa[x_1, \ldots, x_n]$ soit irréductible pour tout
homomorphisme $\varphi : A_f \to \kappa$ vers un corps.
\end{lemma}

\begin{proof}
Il existe un automorphisme $\Psi$ de $A[x_1, \ldots, x_n]$ sur $A$
tel que $\Psi(P) = ax_n^d +$ des termes de degré inférieur en $x_n$, avec
$a \not = 0$, voir
Algèbre, lemme \ref{algebra-lemma-helper-polynomial}.
On peut remplacer $P$ par $\Psi(P)$ et $A$ par $A_a$.
On peut ainsi supposer que $P$ est unitaire en $x_n$ de degré $d > 0$.
Pour $i = 1, \ldots, n - 1$, soit $d_i$ le degré de $P$ en $x_i$.
Notons que cela implique que $P^\varphi$ est unitaire de degré $d$ en $x_n$
et de degré $\leq d_i$ en $x_i$ pour tout homomorphisme
$\varphi : A \to \kappa$, où $\kappa$ est un corps.
Ainsi, si $P^\varphi$ est réductible, on peut écrire
$$
P^\varphi = Q_1 Q_2
$$
avec $Q_1, Q_2$ unitaires de degrés $e_1, e_2 \geq 0$ en $x_n$, où
$e_1 + e_2 = d$, et de degré $\leq d_i$ en $x_i$ pour
$i = 1, \ldots, n - 1$. Autrement dit, on peut écrire
\begin{equation}
\label{equation-factors}
Q_j = x_n^{e_j} + \sum\nolimits_{0 \leq l < e_j}
\left( \sum\nolimits_{L \in \mathcal{L}} a_{j, l, L} x^L \right) x_n^l
\end{equation}
où la somme porte sur l'ensemble $\mathcal{L}$ des multi-indices $L$
de la forme $L = (l_1, \ldots, l_{n - 1})$ avec $0 \leq l_i \leq d_i$.
Pour tous $e_1, e_2 \geq 0$ tels que $e_1 + e_2 = d$, considérons la $A$-algèbre
$$
B_{e_1, e_2} =
A[\{a_{1, l, L}\}_{0 \leq l < e_1, L \in \mathcal{L}},
\{a_{2, l, L}\}_{0 \leq l < e_2, L \in \mathcal{L}}]/(\text{relations})
$$
où $(\text{relations})$ est l'idéal engendré par les coefficients
du polynôme
$$
P - Q_1Q_2 \in
A[\{a_{1, l, L}\}_{0 \leq l < e_1, L \in \mathcal{L}},
\{a_{2, l, L}\}_{0 \leq l < e_2, L \in \mathcal{L}}][x_1, \ldots, x_n]
$$
où $Q_1$ et $Q_2$ sont définis comme dans (\ref{equation-factors}).
L'hypothèse que $P$ est irréductible sur $\overline{K}$ implique qu'il
n'existe aucun homomorphisme de $A$-algèbres
$B_{e_1, e_2} \to \overline{K}$. D'après le Nullstellensatz de Hilbert, voir
Algèbre, théorème \ref{algebra-theorem-nullstellensatz},
cela signifie que $B_{e_1, e_2} \otimes_A K = 0$.
Comme $B_{e_1, e_2}$ est une $A$-algèbre de type fini, cela signifie que
l'on peut trouver un $f_{e_1, e_2} \in A$ tel que
$(B_{e_1, e_2})_{f_{e_1, e_2}} = 0$. Par construction, cela signifie que,
si $\varphi : A_{f_{e_1, e_2}} \to \kappa$ est un homomorphisme vers un corps,
alors $P^\varphi$ n'a pas de factorisation $P^\varphi = Q_1 Q_2$
avec $Q_1$ de degré $e_1$ en $x_n$ et $Q_2$ de degré $e_2$ en $x_n$.
Ainsi, en prenant
$f = \prod_{e_1, e_2 \geq 0, e_1 + e_2 = d} f_{e_1, e_2}$, on conclut.
\end{proof}

\begin{lemma}
\label{lemma-geom-irreducible-generic-fibre}
Soit $f : X \to Y$ un morphisme de schémas.
Supposons que
\begin{enumerate}
\item $Y$ soit irréductible, de point générique $\eta$,
\item $X_\eta$ soit géométriquement irréductible, et
\item $f$ soit de type fini.
\end{enumerate}
Alors il existe un sous-schéma ouvert non vide $V \subset Y$
tel que $X_V \to V$ ait des fibres géométriquement irréductibles.
\end{lemma}

\begin{proof}[Première démonstration du lemme \ref{lemma-geom-irreducible-generic-fibre}]
Nous donnons deux démonstrations du lemme. Elles sont essentiellement équivalentes ;
la seconde est plus autonome, mais un peu plus longue.
Choisissons un diagramme
$$
\xymatrix{
X' \ar[d]_{f'} \ar[r]_{g'} & X_V \ar[r] \ar[d] & X \ar[d]^f \\
Y' \ar[r]^g & V \ar[r] & Y
}
$$
comme dans le
lemme \ref{lemma-make-generic-fibre-geometrically-reduced}.
Notons que la fibre générique de $f'$ est le réduit de la
fibre générique de $f$ (voir le
lemme \ref{lemma-reduction-generic-fibre})
et est donc géométriquement irréductible.
Supposons que le lemme soit vrai pour le morphisme $f'$. Après avoir rétréci
$V$, toutes les fibres de $f'$ sont alors géométriquement irréductibles.
Comme $X' = (Y' \times_V X_V)_{red}$, cela implique que toutes les fibres
de $Y' \times_V X_V$ sont géométriquement irréductibles. D'après le
lemme \ref{lemma-base-change-fibres-geometrically-irreducible},
toutes les fibres de $X_V \to V$ sont donc géométriquement irréductibles et
on conclut. De cette manière, on voit que l'on peut supposer que la fibre générique
est à la fois géométriquement réduite et géométriquement irréductible,
et l'on peut supposer $Y = \Spec(A)$ avec $A$ un anneau intègre.

\medskip\noindent
Soit $x \in X_\eta$ le point générique. Puisque $X_\eta$ est géométriquement
irréductible et réduit, on voit que $L = \kappa(x)$ est une extension de type fini
de $K = \kappa(\eta)$ qui est géométriquement réduite et
géométriquement irréductible, voir
Variétés, lemmes \ref{varieties-lemma-geometrically-reduced-at-point} et
\ref{varieties-lemma-geometrically-irreducible-function-field}.
En particulier, l'extension de corps $L/K$ est séparable, voir
Algèbre, lemme \ref{algebra-lemma-characterize-separable-field-extensions}.
On peut donc trouver $x_1, \ldots, x_{r + 1} \in L$ qui engendrent $L$
sur $K$ et tels que $x_1, \ldots, x_r$ soient une base de transcendance de
$L$ sur $K$, voir
Algèbre, lemme
\ref{algebra-lemma-generating-finitely-generated-separable-field-extensions}.
Soit $P \in K(x_1, \ldots, x_r)[T]$ le polynôme minimal de
$x_{r + 1}$. En chassant les dénominateurs, on peut supposer que
$P$ a ses coefficients dans $A[x_1, \ldots, x_r]$.
Notons que, puisque $L$ est géométriquement réduit et géométriquement irréductible
sur $K$, le polynôme $P$ est irréductible dans
$\overline{K}[x_1, \ldots, x_r, T]$, où $\overline{K}$ est la
clôture algébrique de $K$. Notons
$$
B' = A[x_1, \ldots, x_{r + 1}]/(P(x_{r + 1}))
$$
et posons $X' = \Spec(B')$. Par construction, le corps des fractions de $B'$
est isomorphe à $L = \kappa(x)$ comme extension de $K$. Il existe donc un
ouvert $U \subset X$, un ouvert $U' \subset X'$ et un $Y$-isomorphisme
$U \to U'$, voir
Morphismes, lemme \ref{morphisms-lemma-common-open}.
Voici un diagramme :
$$
\xymatrix{
X \ar[rd] &
U \ar[l] \ar@{=}[r] \ar[d] &
U' \ar[r] \ar[d] &
X' \ar[ld] \ar@{=}[r] & \Spec(B') \\
& Y \ar@{=}[r] & Y &
}
$$
Notons que $U_\eta \subset X_\eta$ et $U'_\eta \subset X'_\eta$ sont des
ouverts denses. Ainsi, après avoir rétréci $Y$ en appliquant le
lemme \ref{lemma-nowhere-dense-generic-fibre},
on obtient que $U_y$ est dense dans $X_y$ et que $U'_y$ est dense dans $X'_y$
pour tout $y \in Y$. Il suffit donc de démontrer le lemme pour
$X' \to Y$, ce qui est le contenu du
lemme \ref{lemma-irreducible-polynomial-over-domain}.
\end{proof}

\begin{proof}[Seconde démonstration du lemme \ref{lemma-geom-irreducible-generic-fibre}]
Soit $Y' \subset Y$ le réduit de $Y$. Soit $X' \to X$ le morphisme de réduction
de $X$. Notons que $X' \to X  \to Y$ se factorise par $Y'$, voir
Schémas, lemme \ref{schemes-lemma-map-into-reduction}.
Comme $Y' \to Y$ et $X' \to X$ sont des homéomorphismes
universels d'après
Morphismes, lemme \ref{morphisms-lemma-reduction-universal-homeomorphism},
on voit qu'il suffit de démontrer le lemme pour $X' \to Y'$. On peut donc
supposer que $X$ et $Y$ sont réduits. En particulier, $Y$ est intègre, voir
Propriétés, lemme \ref{properties-lemma-characterize-integral}.
Ainsi, d'après
Morphismes, proposition \ref{morphisms-proposition-generic-flatness},
il existe un ouvert affine non vide $V \subset Y$ tel que $X_V \to V$ soit
plat et de présentation finie. Après avoir remplacé $Y$ par $V$, on peut
supposer, outre (1), (2), (3), que $Y$ est affine intègre, que $X$
est réduit et que $f$ est plat et de présentation finie. En particulier,
$f$ est universellement ouvert, voir
Morphismes, lemme \ref{morphisms-lemma-fppf-open}.

\medskip\noindent
Choisissons un ouvert affine non vide $U \subset X$. Alors $U \to Y$ est plat et de
présentation finie, à fibre générique géométriquement irréductible.
Le complémentaire $X_\eta \setminus U_\eta$ est nulle part dense. Ainsi, après
avoir rétréci $Y$, on peut supposer $U_y \subset X_y$ ouvert dense pour tout
$y \in Y$, voir le
lemme \ref{lemma-nowhere-dense-generic-fibre}.
On peut donc remplacer $X$ par $U$ et on se ramène au
cas où $Y$ est affine intègre et où $X$ est affine réduit, plat et de présentation
finie sur $Y$, à fibre générique $X_\eta$ géométriquement irréductible.

\medskip\noindent
Écrivons $X = \Spec(B)$ et $Y = \Spec(A)$. Alors $A$ est un anneau intègre,
$B$ est réduit, $A \to B$ est plat et de présentation finie, et $B_K$ est
géométriquement irréductible sur le corps des fractions $K$ de $A$.
En particulier, on voit que $B_K$ est un anneau intègre. Soit $L$ le corps des fractions
de $B_K$. Notons que
$L$ est une extension de corps de type fini de $K$, puisque $B$ est une $A$-algèbre
de présentation finie. Soit $K'/K$ une extension finie purement inséparable
telle que $(L \otimes_K K')_{red}$ soit une extension de corps séparablement engendrée,
voir
Algèbre, lemme \ref{algebra-lemma-make-separable}.
Choisissons $x_1, \ldots, x_n \in K'$ qui engendrent l'extension de corps
$K'$ sur $K$, et tels que $x_i^{q_i} \in A$ pour une certaine puissance d'un nombre premier
$q_i$ (démonstration de l'existence des $x_i$ omise). Soit $A'$ la sous-$A$-algèbre
de $K'$ engendrée par $x_1, \ldots, x_n$. Alors $A'$ est une
sous-$A$-algèbre finie $A' \subset K'$ dont le corps des fractions est $K'$. Notons que
$\Spec(A') \to \Spec(A)$ est un homéomorphisme universel, voir
Algèbre, lemme \ref{algebra-lemma-p-ring-map}.
Il suffit donc de démontrer le résultat après changement de base à $\Spec(A')$.
Nous allons remplacer $A$ par $A'$ et $B$ par $(B \otimes_A A')_{red}$
pour arriver à la situation où $L$ est une extension de corps séparablement engendrée
de $K$. Bien sûr, il peut arriver que $(B \otimes_A A')_{red}$ ne soit plus
plat ou de présentation finie sur $A'$, mais on peut y remédier en
remplaçant $A'$ par $A'_f$ pour un $f \in A'$ convenable, voir
Algèbre, lemme \ref{algebra-lemma-generic-flatness}.

\medskip\noindent
À ce stade, on sait que $A$ est un anneau intègre, que $B$ est réduit, que $A \to B$
est plat et de présentation finie, et que $B_K$ est un anneau intègre dont le corps des fractions
$L$ est une extension de corps séparablement engendrée du corps des fractions $K$
de $A$. D'après Algèbre, lemme
\ref{algebra-lemma-generating-finitely-generated-separable-field-extensions},
on peut écrire $L = K(x_1, \ldots, x_{r + 1})$
où $x_1, \ldots, x_r$ sont algébriquement indépendants sur $K$, et où
$x_{r + 1}$ est séparable sur $K(x_1, \ldots, x_r)$. Après avoir chassé
les dénominateurs, on peut supposer que le polynôme minimal
$P \in K(x_1, \ldots, x_r)[T]$ de $x_{r + 1}$ sur $K(x_1, \ldots, x_r)$
a ses coefficients dans $A[x_1, \ldots, x_r]$. Notons que, puisque
$L/K$ est séparable et que $L$ est géométriquement irréductible sur
$K$, le polynôme $P$ est irréductible sur la clôture algébrique
$\overline{K}$ de $K$. Notons
$$
B' = A[x_1, \ldots, x_{r + 1}]/(P(x_{r + 1})).
$$
Par construction, les corps des fractions de $B$ et $B'$ sont isomorphes comme
extensions de $K$. Il existe donc un isomorphisme de $A$-algèbres
$B_h \cong B'_{h'}$ pour des $h \in B$ et $h' \in B'$ convenables, voir
Morphismes, lemme \ref{morphisms-lemma-common-open}.
Autrement dit, $X$ et $X' = \Spec(B')$ ont un ouvert affine commun $U$.
Voici un diagramme :
$$
\xymatrix{
X = \Spec(B) \ar[rd] &
U \ar[l] \ar[r] \ar[d] &
\Spec(B') = X' \ar[ld] \\
& Y = \Spec(A) &
}
$$
Après avoir encore rétréci $Y$ (en appliquant le
lemme \ref{lemma-nowhere-dense-generic-fibre}
à $Z = X \setminus U$ dans $X$ et à $Z' = X' \setminus U$ dans $X'$),
on voit que $U_y$ est dense dans $X_y$ et que $U_y$ est dense dans $X'_y$
pour tout $y \in Y$. Il suffit donc de démontrer le lemme pour
$X' \to Y$, ce qui est le contenu du
lemme \ref{lemma-irreducible-polynomial-over-domain}.
\end{proof}

\begin{lemma}
\label{lemma-nr-geom-irreducible-components-good}
Soit $f : X \to Y$ un morphisme de schémas. Soit
$n_{X/Y}$ la fonction sur $Y$ qui compte le nombre de composantes
géométriquement irréductibles des fibres de $f$, introduite dans le
lemme \ref{lemma-base-change-fibres-nr-geometrically-irreducible-components}.
Supposons $f$ de type fini.
Soit $y \in Y$ un point. Alors il existe un ouvert non vide
$V \subset \overline{\{y\}}$ tel que $n_{X/Y}|_V$ soit constant.
\end{lemma}

\begin{proof}
Soit $Z$ le schéma muni de la structure induite réduite sur $\overline{\{y\}}$.
Soit $f_Z : X_Z \to Z$ le changement de base de $f$. Il suffit clairement de démontrer
le lemme pour $f_Z$ et le point générique de $Z$. On peut donc supposer que
$Y$ est un schéma intègre, voir
Propriétés, lemme \ref{properties-lemma-characterize-integral}.
Notre but, dans ce cas, est de produire un ouvert non vide $V \subset Y$ tel que
$n_{X/Y}|_V$ soit constant.

\medskip\noindent
Appliquons le
lemme \ref{lemma-make-components-generic-fibre-geometrically-irreducible}
à $f : X \to Y$ ; on obtient $g : Y' \to V \subset Y$. Comme $g : Y' \to V$ est
fini étale surjectif, donc en particulier ouvert (voir
Morphismes, lemme \ref{morphisms-lemma-etale-open}),
il suffit de démontrer qu'il existe un ouvert $V' \subset Y'$
tel que $n_{X'/Y'}|_{V'}$ soit constant, voir
lemme \ref{lemma-base-change-fibres-nr-geometrically-irreducible-components}.
On voit ainsi que l'on peut supposer que toutes les composantes irréductibles de
la fibre générique $X_\eta$ sont géométriquement irréductibles sur $\kappa(\eta)$.

\medskip\noindent
À ce stade, supposons que
$X_\eta = X_{1, \eta} \bigcup \ldots \bigcup X_{n, \eta}$
soit la décomposition de la fibre générique en composantes
(géométriquement) irréductibles.
En particulier, $n_{X/Y}(\eta) = n$.
Soit $X_i$ l'adhérence de
$X_{i, \eta}$ dans $X$. Après avoir rétréci $Y$, on peut supposer que
$X = \bigcup X_i$, voir
lemme \ref{lemma-cover-generic-fibre-neighbourhood}.
Après avoir encore rétréci $Y$, on voit que chaque fibre de
$f$ a au moins $n$ composantes irréductibles, voir
lemme \ref{lemma-irreducible-components-in-neighbourhood}.
Ainsi $n_{X/Y}(y) \geq n$ pour tout $y \in Y$.
Après avoir encore rétréci $Y$, on obtient que $X_{i, y}$
est géométriquement irréductible pour chaque $i$ et tout $y \in Y$, voir
lemme \ref{lemma-geom-irreducible-generic-fibre}.
Comme $X_y = \bigcup X_{i, y}$,
cela montre que $n_{X/Y}(y) \leq n$ et achève la démonstration.
\end{proof}

\begin{lemma}
\label{lemma-nr-geom-irreducible-components-constructible}
Soit $f : X \to Y$ un morphisme de schémas. Soit
$n_{X/Y}$ la fonction sur $Y$ qui compte le nombre de composantes
géométriquement irréductibles des fibres de $f$, introduite dans le
lemme \ref{lemma-base-change-fibres-nr-geometrically-irreducible-components}.
Supposons $f$ de présentation finie. Alors les ensembles de niveau
$$
E_n = \{y \in Y \mid n_{X/Y}(y) = n\}
$$
de $n_{X/Y}$ sont localement constructibles dans $Y$.
\end{lemma}

\begin{proof}
Fixons $n$. Soit $y \in Y$. Nous devons montrer qu'il existe un voisinage ouvert
$V$ de $y$ dans $Y$ tel que $E_n \cap V$ soit constructible dans $V$. On peut donc
supposer que $Y$ est affine. Écrivons $Y = \Spec(A)$ et
$A = \colim A_i$ comme limite inductive d'algèbres
de type fini sur $\mathbf{Z}$. D'après
Limites, lemme \ref{limits-lemma-descend-finite-presentation},
on peut trouver un $i$ et un morphisme $f_i : X_i \to \Spec(A_i)$ de
présentation finie dont le changement de base à $Y$ redonne $f$. D'après le
lemme \ref{lemma-base-change-fibres-nr-geometrically-irreducible-components},
il suffit de démontrer le lemme pour $f_i$. On se ramène ainsi au
cas où $Y$ est le spectre d'un anneau noethérien.

\medskip\noindent
Nous utiliserons le critère du
Topologie, lemme \ref{topology-lemma-characterize-constructible-Noetherian},
pour montrer que $E_n$ est constructible lorsque $Y$ est un schéma noethérien.
Pour cela, soit $Z \subset Y$ un sous-schéma fermé irréductible.
Nous devons montrer que $E_n \cap Z$ contient un ouvert non vide
ou n'est pas dense dans $Z$. Soit $\xi \in Z$ le point générique. Alors le
lemme \ref{lemma-nr-geom-irreducible-components-good}
montre que $n_{X/Y}$ est constant dans un voisinage de $\xi$ dans $Z$.
Cela implique clairement le résultat voulu.
\end{proof}

\section{Composantes connexes des fibres}
\label{section-connected}

\begin{lemma}
\label{lemma-connected-components-in-neighbourhood}
Soit $f : X \to Y$ un morphisme de schémas. Supposons $Y$ irréductible de
point générique $\eta$ et $f$ de type fini. Si $X_\eta$ a $n$
composantes connexes, alors il existe un ouvert non vide $V \subset Y$
tel que, pour tout $y \in V$, la fibre $X_y$ ait au moins $n$
composantes connexes.
\end{lemma}

\begin{proof}
La question étant purement topologique, on peut remplacer $X$ et $Y$ par
leurs réduits. En particulier, cela implique que $Y$ est intègre, voir
Propriétés, lemme \ref{properties-lemma-characterize-integral}.
Soit $X_\eta = X_{1, \eta} \cup \ldots \cup X_{n, \eta}$
la décomposition de $X_\eta$ en composantes connexes.
Soit $X_i \subset X$ le sous-schéma fermé réduit dont la fibre
générique est $X_{i, \eta}$. Notons que $Z_{i, j} = X_i \cap X_j$
est un sous-ensemble fermé de $X$ dont la fibre générique $Z_{i, j, \eta}$ est vide.
Après avoir rétréci $Y$, on peut donc supposer que $Z_{i, j} = \emptyset$, voir
lemme \ref{lemma-empty-generic-fibre}.
Après avoir encore rétréci $Y$, on peut supposer que
$X_y = \bigcup X_{i, y}$ pour $y \in Y$, voir
lemme \ref{lemma-cover-generic-fibre-neighbourhood}.
De plus, après avoir rétréci $Y$, on peut supposer que chaque $X_i \to Y$
est plat et de présentation finie, voir
Morphismes, proposition \ref{morphisms-proposition-generic-flatness}.
Les morphismes $X_i \to Y$ sont ouverts, voir
Morphismes, lemme \ref{morphisms-lemma-fppf-open}.
Il existe donc un voisinage ouvert $V$ de $\eta$ qui est contenu
dans $f(X_i)$ pour chaque $i$.
Pour tout $y \in V$, les schémas $X_{i, y}$ sont des
sous-ensembles fermés non vides de $X_y$, on a $X_y = \bigcup X_{i, y}$,
et les intersections $Z_{i, j, y} = X_{i, y} \cap X_{j, y}$
sont vides ! Cela implique clairement que
$X_y$ a au moins $n$ composantes connexes.
\end{proof}

\begin{lemma}
\label{lemma-base-change-fibres-geometrically-connected}
Soit $f : X \to Y$ un morphisme de schémas.
Soit $g : Y' \to Y$ un morphisme quelconque, et notons
$f' : X' \to Y'$ le changement de base de $f$.
Alors
\begin{align*}
\{y' \in Y' \mid X'_{y'}\text{ est géométriquement connexe}\} \\
= g^{-1}(\{y \in Y \mid X_y\text{ est géométriquement connexe}\}).
\end{align*}
\end{lemma}

\begin{proof}
Cela se ramène à l'énoncé suivant : pour $y' \in Y'$ d'image
$y \in Y$, la fibre $X'_{y'} = X_y \times_y y'$ est géométriquement
connexe sur $\kappa(y')$ si et seulement si $X_y$ est géométriquement connexe
sur $\kappa(y)$. Cela résulte de
Variétés,
lemme \ref{varieties-lemma-geometrically-connected-check-after-extension}.
\end{proof}

\begin{lemma}
\label{lemma-base-change-fibres-nr-geometrically-connected-components}
Soit $f : X \to Y$ un morphisme de schémas. Soit
$$
n_{X/Y} : Y \to \{0, 1, 2, 3, \ldots, \infty\}
$$
la fonction qui associe à $y \in Y$ le nombre de composantes connexes
de $(X_y)_K$, où $K$ est une extension séparablement close
de $\kappa(y)$. Cette fonction est bien définie et, si $g : Y' \to Y$ est un morphisme,
alors
$$
n_{X'/Y'} = n_{X/Y} \circ g
$$
où $X' \to Y'$ est le changement de base de $f$.
\end{lemma}

\begin{proof}
Supposons que $y' \in Y'$ ait pour image $y \in Y$.
Supposons que $K \supset \kappa(y)$ et $K' \supset \kappa(y')$ soient des extensions
séparablement closes. On peut alors choisir un diagramme commutatif
$$
\xymatrix{
K \ar[r] & K'' & K' \ar[l] \\
\kappa(y) \ar[u] \ar[rr] & & \kappa(y') \ar[u]
}
$$
de corps. Le résultat s'ensuit, car les morphismes de schémas
$$
\xymatrix{
(X'_{y'})_{K'} &
(X'_{y'})_{K''} = (X_y)_{K''} \ar[l] \ar[r] &
(X_y)_K
}
$$
induisent des bijections entre les composantes connexes, voir
Variétés,
lemme \ref{varieties-lemma-separably-closed-field-connected-components}.
\end{proof}

\begin{lemma}
\label{lemma-geometrically-connected-generic-fibre}
Soit $f : X \to Y$ un morphisme de schémas.
Supposons que
\begin{enumerate}
\item $Y$ soit irréductible, de point générique $\eta$,
\item $X_\eta$ soit géométriquement connexe, et
\item $f$ soit de type fini.
\end{enumerate}
Alors il existe un sous-schéma ouvert non vide $V \subset Y$
tel que $X_V \to V$ ait des fibres géométriquement connexes.
\end{lemma}

\begin{proof}
Choisissons un diagramme
$$
\xymatrix{
X' \ar[d]_{f'} \ar[r]_{g'} & X_V \ar[r] \ar[d] & X \ar[d]^f \\
Y' \ar[r]^g & V \ar[r] & Y
}
$$
comme dans le
lemme \ref{lemma-make-components-generic-fibre-geometrically-irreducible}.
Notons que la fibre générique de $f'$ est géométriquement connexe
(par exemple d'après le
lemme \ref{lemma-base-change-fibres-nr-geometrically-connected-components}).
Supposons que le lemme soit vrai pour le morphisme $f'$. Cela signifie qu'il
existe un ouvert non vide $W \subset Y'$ tel que toute fibre de
$X' \to Y'$ au-dessus de $W$ soit géométriquement connexe.
Alors, comme $g$ est un morphisme ouvert d'après
Morphismes, lemme \ref{morphisms-lemma-etale-open},
toutes les fibres de $f$ aux points de l'ouvert non vide $V = g(W)$ sont
géométriquement connexes, voir
lemme \ref{lemma-base-change-fibres-nr-geometrically-connected-components}.
On voit ainsi que l'on peut supposer que les composantes irréductibles
de la fibre générique $X_\eta$ sont géométriquement irréductibles.

\medskip\noindent
Soit $Y'$ le réduit de $Y$, et posons $X' = Y' \times_Y X$.
Il suffit alors de démontrer le lemme pour le morphisme $X' \to Y'$
(par exemple d'après le
lemme \ref{lemma-base-change-fibres-nr-geometrically-connected-components}
une fois encore). Comme la fibre générique de $X' \to Y'$ est la même que la
fibre générique de $X \to Y$, on voit que l'on peut supposer que $Y$ est
irréductible et réduit (c'est-à-dire intègre, voir
Propriétés, lemme \ref{properties-lemma-characterize-integral})
et que les composantes irréductibles de la fibre générique $X_\eta$
sont géométriquement irréductibles.

\medskip\noindent
À ce stade, supposons que
$X_\eta = X_{1, \eta} \bigcup \ldots \bigcup X_{n, \eta}$
soit la décomposition de la fibre générique en composantes
(géométriquement) irréductibles.
Soit $X_i$ l'adhérence de $X_{i, \eta}$ dans $X$.
Après avoir rétréci $Y$, on peut supposer que
$X = \bigcup X_i$, voir
lemme \ref{lemma-cover-generic-fibre-neighbourhood}.
Soit $Z_{i, j} = X_i \cap X_j$.
Soit
$$
\{1, \ldots, n\} \times \{1, \ldots, n\} = I \amalg J
$$
où $(i, j) \in I$ si $Z_{i, j, \eta} = \emptyset$ et
$(i, j) \in J$ si $Z_{i, j, \eta} \not = \emptyset$.
Après avoir rétréci $Y$, on peut supposer que $Z_{i, j} = \emptyset$
pour tout $(i, j) \in I$, voir
lemme \ref{lemma-empty-generic-fibre}.
Après avoir rétréci $Y$, on obtient que $X_{i, y}$
est géométriquement irréductible pour chaque $i$ et tout $y \in Y$, voir
lemme \ref{lemma-geom-irreducible-generic-fibre}.
Après avoir encore rétréci $Y$, on arrive à la situation où
chaque $Z_{i, j} \to Y$ est plat et de présentation finie pour
tout $(i, j) \in J$, voir
Morphismes, proposition \ref{morphisms-proposition-generic-flatness}.
Cela signifie que $f(Z_{i, j}) \subset Y$ est ouvert, voir
Morphismes, lemme \ref{morphisms-lemma-fppf-open}.
Nous affirmons que
$$
V  = \bigcap\nolimits_{(i, j) \in J} f(Z_{i, j})
$$
convient, c'est-à-dire que $X_y$ est géométriquement connexe pour tout
$y \in V$. En effet, le fait que $X_\eta$ soit connexe implique que
la relation d'équivalence engendrée par les paires de $J$ n'a qu'une
seule classe d'équivalence. Or, si $y \in V$ et si $K \supset \kappa(y)$
est une extension séparablement close, alors les composantes irréductibles
de $(X_y)_K$ sont les fibres $(X_{i, y})_K$. De plus, la construction et
$y \in V$ montrent que $(X_{i, y})_K$ rencontre $(X_{j, y})_K$
si et seulement si $(i, j) \in J$. La remarque sur les classes d'équivalence
montre donc que $(X_y)_K$ est connexe, et on conclut.
\end{proof}

\begin{lemma}
\label{lemma-nr-geom-connected-components-good}
Soit $f : X \to Y$ un morphisme de schémas. Soit
$n_{X/Y}$ la fonction sur $Y$ qui compte le nombre de composantes
géométriquement connexes des fibres de $f$, introduite dans le
lemme \ref{lemma-base-change-fibres-nr-geometrically-connected-components}.
Supposons $f$ de type fini.
Soit $y \in Y$ un point. Alors il existe un ouvert non vide
$V \subset \overline{\{y\}}$ tel que $n_{X/Y}|_V$ soit constant.
\end{lemma}

\begin{proof}
Soit $Z$ le schéma muni de la structure induite réduite sur $\overline{\{y\}}$.
Soit $f_Z : X_Z \to Z$ le changement de base de $f$. Il suffit clairement de démontrer
le lemme pour $f_Z$ et le point générique de $Z$. On peut donc supposer que
$Y$ est un schéma intègre, voir
Propriétés, lemme \ref{properties-lemma-characterize-integral}.
Notre but, dans ce cas, est de produire un ouvert non vide $V \subset Y$ tel que
$n_{X/Y}|_V$ soit constant.

\medskip\noindent
Appliquons le
lemme \ref{lemma-make-components-generic-fibre-geometrically-irreducible}
à $f : X \to Y$ ; on obtient $g : Y' \to V \subset Y$. Comme $g : Y' \to V$ est
fini étale surjectif, donc en particulier ouvert (voir
Morphismes, lemme \ref{morphisms-lemma-etale-open}),
il suffit de démontrer qu'il existe un ouvert $V' \subset Y'$
tel que $n_{X'/Y'}|_{V'}$ soit constant, voir
lemme \ref{lemma-base-change-fibres-nr-geometrically-irreducible-components}.
On voit ainsi que l'on peut supposer que toutes les composantes irréductibles de
la fibre générique $X_\eta$ sont géométriquement irréductibles sur $\kappa(\eta)$.
D'après
Variétés, lemme
\ref{varieties-lemma-irreducible-components-geometrically-irreducible},
cela implique que les composantes connexes de $X_\eta$ sont elles aussi
géométriquement connexes.

\medskip\noindent
À ce stade, supposons que
$X_\eta = X_{1, \eta} \bigcup \ldots \bigcup X_{n, \eta}$
soit la décomposition de la fibre générique en composantes
(géométriquement) connexes.
En particulier, $n_{X/Y}(\eta) = n$.
Soit $X_i$ l'adhérence de
$X_{i, \eta}$ dans $X$. Après avoir rétréci $Y$, on peut supposer que
$X = \bigcup X_i$, voir
lemme \ref{lemma-cover-generic-fibre-neighbourhood}.
Après avoir encore rétréci $Y$, on voit que chaque fibre de
$f$ a au moins $n$ composantes connexes, voir
lemme \ref{lemma-connected-components-in-neighbourhood}.
Ainsi $n_{X/Y}(y) \geq n$ pour tout $y \in Y$.
Après avoir encore rétréci $Y$, on obtient que $X_{i, y}$
est géométriquement connexe pour chaque $i$ et tout $y \in Y$, voir
lemme \ref{lemma-geometrically-connected-generic-fibre}.
Comme $X_y = \bigcup X_{i, y}$,
cela montre que $n_{X/Y}(y) \leq n$ et achève la démonstration.
\end{proof}

\begin{lemma}
\label{lemma-nr-geom-connected-components-constructible}
Soit $f : X \to Y$ un morphisme de schémas. Soit
$n_{X/Y}$ la fonction sur $Y$ qui compte le nombre de composantes
géométriquement connexes des fibres de $f$, introduite dans le
lemme \ref{lemma-base-change-fibres-nr-geometrically-connected-components}.
Supposons $f$ de présentation finie. Alors les ensembles de niveau
$$
E_n = \{y \in Y \mid n_{X/Y}(y) = n\}
$$
de $n_{X/Y}$ sont localement constructibles dans $Y$.
\end{lemma}

\begin{proof}
Fixons $n$. Soit $y \in Y$. Nous devons montrer qu'il existe un voisinage ouvert
$V$ de $y$ dans $Y$ tel que $E_n \cap V$ soit constructible dans $V$. On peut donc
supposer que $Y$ est affine. Écrivons $Y = \Spec(A)$ et
$A = \colim A_i$ comme limite inductive d'algèbres
de type fini sur $\mathbf{Z}$. D'après
Limites, lemme \ref{limits-lemma-descend-finite-presentation},
on peut trouver un $i$ et un morphisme $f_i : X_i \to \Spec(A_i)$ de
présentation finie dont le changement de base à $Y$ redonne $f$. D'après le
lemme \ref{lemma-base-change-fibres-nr-geometrically-connected-components},
il suffit de démontrer le lemme pour $f_i$. On se ramène ainsi au
cas où $Y$ est le spectre d'un anneau noethérien.

\medskip\noindent
Nous utiliserons le critère du
Topologie, lemme \ref{topology-lemma-characterize-constructible-Noetherian},
pour montrer que $E_n$ est constructible lorsque $Y$ est un schéma noethérien.
Pour cela, soit $Z \subset Y$ un sous-schéma fermé irréductible.
Nous devons montrer que $E_n \cap Z$ contient un ouvert non vide
ou n'est pas dense dans $Z$. Soit $\xi \in Z$ le point générique. Alors le
lemme \ref{lemma-nr-geom-connected-components-good}
montre que $n_{X/Y}$ est constant dans un voisinage de $\xi$ dans $Z$.
Cela implique clairement le résultat voulu.
\end{proof}

\begin{lemma}
\label{lemma-connected-flat-over-dvr}
\begin{slogan}
Une dégénérescence plate d'un schéma non connexe est soit non connexe,
soit non réduite.
\end{slogan}
Soit $f : X \to S$ un morphisme de schémas.
Supposons que
\begin{enumerate}
\item $S$ soit le spectre d'un anneau de valuation discrète,
\item $f$ soit plat,
\item $X$ soit connexe,
\item la fibre fermée $X_s$ soit réduite.
\end{enumerate}
Alors la fibre générique $X_\eta$ est connexe.
\end{lemma}

\begin{proof}
Écrivons $S = \Spec(R)$ et soit $\pi \in R$ une uniformisante.
Pour obtenir une contradiction, supposons que $X_\eta$ ne soit pas connexe.
Cela signifie qu'il existe un idempotent non trivial
$e \in \Gamma(X_\eta, \mathcal{O}_{X_\eta})$.
Soit $U = \Spec(A)$ un ouvert affine quelconque de $X$.
Notons que $\pi$ est un non-diviseur de zéro dans $A$, puisque $A$ est plat sur $R$, voir
Plus sur l'algèbre, lemme \ref{more-algebra-lemma-flat-torsion-free},
par exemple. Alors $e|_{U_\eta}$ correspond à un élément $e \in A[1/\pi]$.
Soit $z \in A$ un élément tel que $e = z/\pi^n$, avec $n \geq 0$ minimal.
Notons que $z^2 = \pi^nz$. Cela signifie que $z \bmod \pi A$ est nilpotent
si $n > 0$. Par hypothèse, $A/\pi A$ est réduit, et la minimalité de
$n$ implique donc $n = 0$. On conclut ainsi que $e \in A$ ! Autrement dit,
$e \in \Gamma(X, \mathcal{O}_X)$. Comme $X$ est connexe, il s'ensuit
que $e$ est un idempotent trivial, ce qui est une contradiction.
\end{proof}

\section{Composantes connexes rencontrant une section}
\label{section-connected-components}

\noindent
Les résultats de cette section s'appliquent en particulier à un schéma en groupes
$G \to S$ et à sa section neutre $e : S \to G$.

\begin{situation}
\label{situation-connected-along-section}
Ici, $f : X \to Y$ est un morphisme de schémas, et
$s : Y \to X$ est une section de $f$.
Pour tout $y \in Y$, on note $X^0_y$ la composante connexe de $X_y$
contenant $s(y)$. Enfin, on pose $X^0 = \bigcup_{y \in Y} X^0_y$.
\end{situation}

\begin{lemma}
\label{lemma-base-change-connected-along-section}
Soient $f : X \to Y$, $s : Y \to X$ comme dans la
situation \ref{situation-connected-along-section}.
Si $g : Y' \to Y$ est un morphisme quelconque, considérons le diagramme de changement de base
$$
\xymatrix{
X' \ar[r]_{g'} \ar[d]^{f'} & X \ar[d]_f \\
Y' \ar@/^1pc/[u]^{s'} \ar[r]^g & Y \ar@/_1pc/[u]_s
}
$$
de sorte que l'on obtient $(X')^0 \subset X'$.
Alors $(X')^0 = (g')^{-1}(X^0)$.
\end{lemma}

\begin{proof}
Soit $y' \in Y'$ d'image $y \in Y$. On peut considérer
$X^0_y$ comme un sous-schéma fermé de $X_y$, voir par exemple
Morphismes,
définition \ref{morphisms-definition-scheme-structure-connected-component}.
Comme $s(y) \in X^0_y$, il résulte de
Variétés, lemme
\ref{varieties-lemma-geometrically-connected-if-connected-and-point}
que $X_y^0$ est un schéma géométriquement connexe sur $\kappa(y)$.
Ainsi $X_y^0 \times_y y' \to X'_{y'}$ est un sous-schéma fermé connexe
qui contient $s'(y')$. Donc $X_y^0 \times_y y' \subset (X'_{y'})^0$.
L'autre inclusion $X_y^0 \times_y y' \supset (X'_{y'})^0$ est claire,
car l'image de $(X'_{y'})^0$ dans $X_y$ est un sous-ensemble connexe de $X_y$ qui
contient $s(y)$.
\end{proof}

\begin{lemma}
\label{lemma-connected-along-section-good}
Soient $f : X \to Y$, $s : Y \to X$ comme dans la
situation \ref{situation-connected-along-section}.
Supposons $f$ de type fini. Soit $y \in Y$ un point.
Alors il existe un ouvert non vide $V \subset \overline{\{y\}}$ tel que
l'image réciproque de $X^0$ dans le changement de base $X_V$ soit ouverte et fermée dans
$X_V$.
\end{lemma}

\begin{proof}
Soit $Z \subset Y$ le sous-schéma fermé muni de la structure induite réduite
sur $\overline{\{y\}}$. Soient $f_Z : X_Z \to Z$ et $s_Z : Z \to X_Z$
les changements de base de $f$ et $s$. D'après le
lemme \ref{lemma-base-change-connected-along-section},
on a $(X_Z)^0 = (X^0)_Z$. Il suffit donc de démontrer le lemme pour
le morphisme $X_Z \to Z$ et le point $x \in X_Z$ qui s'envoie sur le point générique
de $Z$. Autrement dit, on s'est ramené au cas
où $Y$ est un schéma intègre (voir
Propriétés, lemme \ref{properties-lemma-characterize-integral})
de point générique $\eta$. Notre but est de montrer qu'après avoir rétréci
$Y$, le sous-ensemble $X^0$ devient un sous-ensemble ouvert et fermé de $X$.

\medskip\noindent
Notons que le schéma $X_\eta$ est de type fini sur un corps, donc noethérien.
Ses composantes connexes sont ainsi à la fois ouvertes et fermées. On peut donc écrire
$X_\eta = X_\eta^0 \amalg T_\eta$ pour un certain sous-ensemble ouvert et fermé
$T_\eta$ de $X_\eta$. Soit ensuite $T \subset X$ l'adhérence de $T_\eta$,
et soit $X^{00} \subset X$ l'adhérence de $X_\eta^0$. Notons que
$T_\eta$, respectivement $X^0_\eta$, est la fibre générique de $T$, respectivement de $X^{00}$,
voir la discussion précédant le
lemme \ref{lemma-cover-generic-fibre-neighbourhood}.
De plus, ce lemme implique qu'après avoir rétréci $Y$, on peut supposer que
$X = X^{00} \cup T$ (au sens des ensembles).
Notons que $(T \cap X^{00})_\eta = T_\eta \cap X^0_\eta = \emptyset$.
Après avoir rétréci $Y$, on peut donc supposer que $T \cap X^{00} = \emptyset$, voir
lemme \ref{lemma-empty-generic-fibre}.
En particulier, $X^{00}$ est ouvert dans $X$. Notons que $X^0_\eta$ est connexe
et possède un point rationnel, à savoir $s(\eta)$ ; il est donc géométriquement
connexe, voir
Variétés,
lemme \ref{varieties-lemma-geometrically-connected-if-connected-and-point}.
Ainsi, après avoir rétréci $Y$, on peut supposer que toutes les fibres de $X^{00} \to Y$
sont géométriquement connexes, voir
lemme \ref{lemma-geometrically-connected-generic-fibre}.
Il s'ensuit alors que les fibres $X^{00}_y$ sont des
sous-ensembles ouverts, fermés et connexes de $X_y$ contenant $s(y)$.
On en déduit que $X^0 = X^{00}$, et on conclut.
\end{proof}

\begin{lemma}
\label{lemma-connected-along-section-locally-constructible}
Soient $f : X \to Y$, $s : Y \to X$ comme dans la
situation \ref{situation-connected-along-section}.
Si $f$ est de présentation finie, alors $X^0$ est localement constructible
dans $X$.
\end{lemma}

\begin{proof}
Soit $x \in X$. Nous devons montrer qu'il existe un voisinage ouvert
$U$ de $x$ tel que $X^0 \cap U$ soit constructible dans $U$.
Cela nous ramène au cas où $Y$ est affine.
Écrivons $Y = \Spec(A)$ et $A = \colim A_i$ comme limite
inductive d'algèbres de type fini sur $\mathbf{Z}$. D'après
Limites, lemme \ref{limits-lemma-descend-finite-presentation},
on peut trouver un $i$ et un morphisme $f_i : X_i \to \Spec(A_i)$ de
présentation finie, muni d'une section $s_i : \Spec(A_i) \to X_i$
dont le changement de base à $Y$ redonne $f$ et la section $s$. D'après le
lemme \ref{lemma-base-change-connected-along-section},
il suffit de démontrer le lemme pour $f_i, s_i$. On se ramène ainsi au
cas où $Y$ est le spectre d'un anneau noethérien.

\medskip\noindent
Supposons que $Y$ soit un schéma affine noethérien. Puisque $f$ est de présentation finie,
c'est-à-dire de type fini, on voit que $X$ est lui aussi un schéma noethérien, voir
Morphismes, lemme \ref{morphisms-lemma-finite-type-noetherian}.
Pour démontrer le lemme dans
ce cas, il suffit de montrer que, pour tout sous-ensemble fermé irréductible
$Z \subset X$, l'intersection $Z \cap X^0$ contient un ouvert non vide
de $Z$ ou n'est pas dense dans $Z$, voir
Topologie, lemme \ref{topology-lemma-characterize-constructible-Noetherian}.
Soit $x \in Z$ le point générique, et soit $y = f(x)$. D'après le
lemme \ref{lemma-connected-along-section-good},
il existe un sous-ensemble ouvert non vide $V \subset \overline{\{y\}}$ tel
que $X^0 \cap X_V$ soit ouvert et fermé dans $X_V$. Comme
$f(Z) \subset \overline{\{y\}}$ et $f(x) = y \in V$, on voit que
$W = f^{-1}(V) \cap Z$ est un sous-ensemble ouvert non vide de $Z$. Il s'ensuit que
$X^0 \cap W$ est ouvert et fermé dans $W$. Puisque $W$ est irréductible,
on voit que $X^0 \cap W$ est soit vide, soit égal à $W$.
Cela démontre le lemme.
\end{proof}

\begin{lemma}
\label{lemma-connected-along-section-open-neighbourhood}
Soient $f : X \to Y$, $s : Y \to X$ comme dans la
situation \ref{situation-connected-along-section}.
Soit $y \in Y$ un point.
Supposons que
\begin{enumerate}
\item $f$ soit de présentation finie et plat, et
\item la fibre $X_y$ soit géométriquement réduite.
\end{enumerate}
Alors $X^0$ est un voisinage de $X^0_y$ dans $X$.
\end{lemma}

\begin{proof}
On peut remplacer $Y$ par un voisinage ouvert affine de $y$.
Écrivons $Y = \Spec(A)$ et $A = \colim A_i$ comme limite
inductive d'algèbres de type fini sur $\mathbf{Z}$. D'après
Limites, lemme \ref{limits-lemma-descend-finite-presentation},
on peut trouver un $i$ et un morphisme $f_i : X_i \to \Spec(A_i)$ de
présentation finie, muni d'une section $s_i : \Spec(A_i) \to X_i$
dont le changement de base à $Y$ redonne $f$ et la section $s$.
Après avoir éventuellement augmenté $i$, on peut aussi supposer que $f_i$ est plat, voir
Limites, lemme \ref{limits-lemma-descend-flat-finite-presentation}.
Soit $y_i$ l'image de $y$ dans $Y_i$. Notons que
$X_y = (X_{i, y_i}) \times_{y_i} y$. Ainsi $X_{i, y_i}$ est géométriquement
réduit, voir
Variétés, lemme \ref{varieties-lemma-geometrically-reduced-upstairs}.
D'après le
lemme \ref{lemma-base-change-connected-along-section},
il suffit de démontrer le lemme pour le système $f_i, s_i, y_i \in Y_i$.
On se ramène ainsi au cas où $Y$ est le spectre d'un anneau noethérien.

\medskip\noindent
Supposons que $Y$ soit le spectre d'un anneau noethérien.
Puisque $f$ est de présentation finie,
c'est-à-dire de type fini, on voit que $X$ est lui aussi un schéma noethérien, voir
Morphismes, lemme \ref{morphisms-lemma-finite-type-noetherian}.
Soit $x \in X^0$ un point au-dessus de $y$. D'après
Topologie, lemme \ref{topology-lemma-constructible-neighbourhood-Noetherian},
il suffit de démontrer que, pour tout fermé irréductible $Z \subset X$
passant par $x$, l'intersection $X^0 \cap Z$ est dense dans $Z$.
En particulier, il suffit de démontrer que le point générique $x' \in Z$
appartient à $X^0$. D'après
Propriétés, lemme \ref{properties-lemma-locally-Noetherian-specialization-dvr},
on peut trouver un anneau de valuation discrète $R$ et un morphisme
$\Spec(R) \to X$ qui envoie le point spécial sur $x$ et
le point générique sur $x'$. Nous considérerons $\Spec(R)$
comme un schéma sur $Y$ via la composée $\Spec(R) \to X \to Y$. D'après le
lemme \ref{lemma-base-change-connected-along-section},
on a que $(X_R)^0$ est l'image réciproque de $X^0$.
Par construction, on dispose d'une seconde section $t : \Spec(R) \to X_R$
(outre le changement de base $s_R$ de $s$)
du morphisme structural $X_R \to \Spec(R)$ telle que
$t(\eta_R)$ soit un point de $X_R$ qui s'envoie sur $x'$ et
$t(0_R)$ un point de $X_R$ qui s'envoie sur $x$. Notons que
$t(0_R)$ appartient à $(X_R)^0$ et que $t(\eta_R) \leadsto t(0_R)$.
Il suffit donc de démontrer que cela implique $t(\eta_R) \in (X_R)^0$.
Il suffit par conséquent de démontrer le lemme lorsque $Y$
est le spectre d'un anneau de valuation discrète et $y$ son point fermé.

\medskip\noindent
Supposons que $Y$ soit le spectre d'un anneau de valuation discrète et que $y$ soit son point
fermé. Notre but est de démontrer que $X^0$ est un voisinage de $X^0_y$.
Notons que $X^0_y$ est ouvert et fermé dans $X_y$, puisque $X_y$ possède un nombre fini
de composantes irréductibles. Le complémentaire $C = X_y \setminus X_y^0$
est donc fermé dans $X$. Ainsi $U = X \setminus C$ est un voisinage ouvert de
$X^0_y$, et $U^0 = X^0$. Il suffit donc de démontrer le résultat pour le
morphisme $U \to Y$. Autrement dit, on peut supposer que
$X_y$ est connexe. Supposons que $X$ ne soit pas connexe, et écrivons
$X = X_1 \amalg \ldots \amalg X_n$ comme une décomposition en composantes
connexes. Alors $s(Y)$ est entièrement contenu dans l'un des $X_i$.
Disons $s(Y) \subset X_1$. Alors $X^0 \subset X_1$. On peut donc remplacer
$X$ par $X_1$ et supposer que $X$ est connexe. À ce stade, le
lemme \ref{lemma-connected-flat-over-dvr}
implique que $X_\eta$ est connexe, c'est-à-dire $X^0 = X$, et on conclut.
\end{proof}

\begin{lemma}
\label{lemma-connected-along-section-open}
Soient $f : X \to Y$, $s : Y \to X$ comme dans la
situation \ref{situation-connected-along-section}.
Supposons que
\begin{enumerate}
\item $f$ soit de présentation finie et plat, et
\item toutes les fibres de $f$ soient géométriquement réduites.
\end{enumerate}
Alors $X^0$ est ouvert dans $X$.
\end{lemma}

\begin{proof}
C'est une conséquence immédiate du
lemme \ref{lemma-connected-along-section-open-neighbourhood}.
\end{proof}

\section{Dimension des fibres}
\label{section-dimension}

\begin{lemma}
\label{lemma-dimension-in-neighbourhood}
Soit $f : X \to Y$ un morphisme de schémas. Supposons $Y$ irréductible de
point générique $\eta$ et $f$ de type fini. Si $X_\eta$ est de dimension $n$,
alors il existe un ouvert non vide $V \subset Y$
tel que, pour tout $y \in V$, la fibre $X_y$ soit de dimension $n$.
\end{lemma}

\begin{proof}
Soit $Z = \{x \in X \mid \dim_x(X_{f(x)}) > n \}$. D'après
Morphismes, lemme \ref{morphisms-lemma-openness-bounded-dimension-fibres},
c'est un sous-ensemble fermé de $X$. Par hypothèse, $Z_\eta = \emptyset$.
Ainsi, d'après le
lemme \ref{lemma-empty-generic-fibre},
on peut rétrécir $Y$ et supposer que $Z = \emptyset$. Soit
$Z' = \{x \in X \mid \dim_x(X_{f(x)}) > n - 1 \} =
\{x \in X \mid \dim_x(X_{f(x)}) = n\}$. Comme précédemment, c'est un sous-ensemble fermé
de $X$. Par hypothèse, on a $Z'_\eta \not = \emptyset$. Ainsi, après
avoir rétréci $Y$, on peut supposer que $Z' \to Y$ est surjectif, voir
lemme \ref{lemma-nonempty-generic-fibre}.
On conclut.
\end{proof}

\begin{lemma}
\label{lemma-base-change-dimension-fibres}
Soit $f : X \to Y$ un morphisme de type fini. Soit
$$
n_{X/Y} : Y \to \{0, 1, 2, 3, \ldots, \infty\}
$$
la fonction qui associe à $y \in Y$ la dimension de $X_y$.
Si $g : Y' \to Y$ est un morphisme, alors
$$
n_{X'/Y'} = n_{X/Y} \circ g
$$
où $X' \to Y'$ est le changement de base de $f$.
\end{lemma}

\begin{proof}
Cela résulte de
Morphismes, lemme \ref{morphisms-lemma-dimension-fibre-after-base-change}.
\end{proof}

\begin{lemma}
\label{lemma-dimension-fibres-constructible}
Soit $f : X \to Y$ un morphisme de schémas. Soit
$n_{X/Y}$ la fonction sur $Y$ donnant la dimension des fibres de $f$,
introduite dans le
lemme \ref{lemma-base-change-dimension-fibres}.
Supposons $f$ de présentation finie. Alors les ensembles de niveau
$$
E_n = \{y \in Y \mid n_{X/Y}(y) = n\}
$$
de $n_{X/Y}$ sont localement constructibles dans $Y$.
\end{lemma}

\begin{proof}
Fixons $n$. Soit $y \in Y$. Nous devons montrer qu'il existe un voisinage ouvert
$V$ de $y$ dans $Y$ tel que $E_n \cap V$ soit constructible dans $V$. On peut donc
supposer que $Y$ est affine. Écrivons $Y = \Spec(A)$ et
$A = \colim A_i$ comme limite inductive d'algèbres
de type fini sur $\mathbf{Z}$. D'après
Limites, lemme \ref{limits-lemma-descend-finite-presentation},
on peut trouver un $i$ et un morphisme $f_i : X_i \to \Spec(A_i)$ de
présentation finie dont le changement de base à $Y$ redonne $f$. D'après le
lemme \ref{lemma-base-change-dimension-fibres},
il suffit de démontrer le lemme pour $f_i$. On se ramène ainsi au
cas où $Y$ est le spectre d'un anneau noethérien.

\medskip\noindent
Nous utiliserons le critère de
Topologie, lemme \ref{topology-lemma-characterize-constructible-Noetherian},
pour démontrer que $E_n$ est constructible lorsque $Y$ est un schéma noethérien.
Pour cela, soit $Z \subset Y$ un sous-schéma fermé irréductible.
Nous devons montrer que $E_n \cap Z$ contient un sous-ensemble ouvert non vide
ou n'est pas dense dans $Z$. Soit $\xi \in Z$ le point générique. Alors le
lemme \ref{lemma-dimension-in-neighbourhood}
montre que $n_{X/Y}$ est constante dans un voisinage de $\xi$ dans $Z$.
Cela donne le résultat voulu.
\end{proof}

\begin{lemma}
\label{lemma-dimension-fibres-flat}
Soit $f : X \to Y$ un morphisme plat de schémas, de présentation finie. Soit
$n_{X/Y}$ la fonction sur $Y$ donnant la dimension des fibres de $f$,
introduite dans le lemme \ref{lemma-base-change-dimension-fibres}.
Alors $n_{X/Y}$ est semi-continue inférieurement.
\end{lemma}

\begin{proof}
Soit $W \subset X$, $W = \coprod_{d \geq 0} U_d$ l'ouvert construit dans les
lemmes \ref{lemma-flat-finite-presentation-CM-open} et
\ref{lemma-flat-finite-presentation-CM-pieces}.
Soit $y \in Y$ un point. Si $n_{X/Y}(y) = \dim(X_y) = n$, alors
$y$ appartient à l'image de $U_n \to Y$.
D'après Morphismes, lemme \ref{morphisms-lemma-fppf-open},
on voit que $f(U_n)$ est ouvert dans $Y$.
Il existe donc un voisinage ouvert de $y$ où
$n_{X/Y}$ est $\geq n$.
\end{proof}

\begin{lemma}
\label{lemma-dimension-fibres-proper}
Soit $f : X \to Y$ un morphisme propre de schémas. Soit
$n_{X/Y}$ la fonction sur $Y$ donnant la dimension des fibres de $f$,
introduite dans le lemme \ref{lemma-base-change-dimension-fibres}.
Alors $n_{X/Y}$ est semi-continue supérieurement.
\end{lemma}

\begin{proof}
Soit $Z_d = \{x \in X \mid \dim_x(X_{f(x)}) > d\}$.
Alors $Z_d$ est un sous-ensemble fermé de $X$ d'après
Morphismes, lemme \ref{morphisms-lemma-openness-bounded-dimension-fibres}.
Puisque $f$ est propre, $f(Z_d)$ est fermé.
Comme $y \in f(Z_d) \Leftrightarrow n_{X/Y}(y) > d$,
on voit que le lemme est vrai.
\end{proof}

\begin{lemma}
\label{lemma-dimension-fibres-proper-flat}
Soit $f : X \to Y$ un morphisme propre et plat de schémas, de présentation finie.
Soit $n_{X/Y}$ la fonction sur $Y$ donnant la dimension des fibres de $f$,
introduite dans le lemme \ref{lemma-base-change-dimension-fibres}.
Alors $n_{X/Y}$ est localement constante.
\end{lemma}

\begin{proof}
C'est une conséquence immédiate des
lemmes \ref{lemma-dimension-fibres-flat} et
\ref{lemma-dimension-fibres-proper}.
\end{proof}

\section{Normalisation de Noether relative faible}
\label{section-relative-noether-normalization}

\noindent
Le but de cette section est de démontrer le
lemme \ref{lemma-weak-relative-noether-normalization}.

\begin{lemma}
\label{lemma-silly}
Soit $R$ un anneau. Soient $\mathfrak p_1, \ldots, \mathfrak p_r$
des idéaux premiers de $R$ tels que $\mathfrak p_i \not \subset \mathfrak p_j$
si $i \not = j$. Soient $k_i \subset \kappa(\mathfrak p_i)$ des
sous-corps tels que les extensions $\kappa(\mathfrak p_i)/k_i$
ne soient pas algébriques. Soit $J \subset R$ un idéal qui n'est contenu
dans aucun des $\mathfrak p_i$. Alors il existe un élément $x \in J$
tel que l'image de $x$ dans $\kappa(\mathfrak p_i)$
est transcendante sur $k_i$ pour $i = 1, \ldots, r$.
\end{lemma}

\begin{proof}
L'idéal
$J_i = J \mathfrak p_1 \ldots \hat{\mathfrak p}_i \ldots \mathfrak p_r$ n'est pas
contenu dans $\mathfrak p_i$, voir
Algèbre, lemme \ref{algebra-lemma-product-ideals-in-prime}.
Il s'ensuit que tout élément $\xi$ de
$\kappa(\mathfrak p_i) = \text{Frac}(B/\mathfrak p_i)$
est de la forme $\xi = a/b$, avec $a, b \in J_i$
et $b \not \in \mathfrak p_i$. En choisissant $\xi$ transcendant
sur $k_i$, on voit que soit $a$, soit $b$ s'envoie sur un élément
de $\kappa(\mathfrak p_i)$ transcendant sur $k_i$.
On conclut que, pour tout $i = 1, \ldots, r$,
on peut trouver un élément
$x_i \in J_i = J \mathfrak p_1 \ldots \hat{\mathfrak p}_i \ldots \mathfrak p_r$
qui s'envoie sur un élément de $\kappa(\mathfrak p_i)$
transcendant sur $k_i$. Alors $x = x_1 + \ldots + x_r$ convient.
\end{proof}

\begin{lemma}
\label{lemma-sum-is-ok}
Soit $R \to S$ un homomorphisme d'anneaux de type fini. Soient $d \geq 0$ et $a, b \in S$.
Supposons que les fibres de
$$
f_a : \Spec(S) \longrightarrow \mathbf{A}^1_R
$$
donné par l'homomorphisme de $R$-algèbres $R[x] \to S$ qui envoie $x$ sur $a$
soient de dimension $\leq d$. Alors il existe un $n_0$ tel que, pour
$n \geq n_0$, les fibres de
$$
f_{a^n + b} : \Spec(S) \longrightarrow \mathbf{A}^1_R
$$
donné par l'homomorphisme de $R$-algèbres $R[x] \to S$ qui envoie $x$ sur $a^n + b$
soient de dimension $\leq d$.
\end{lemma}

\begin{proof}
Dans ce paragraphe, on se ramène au cas où $R \to S$ est de présentation
finie. En effet, écrivons $S = R[A, B, x_1, \ldots, x_n]/J$ pour un idéal
$J \subset R[x_1, \ldots, x_n]$, où $A$ et $B$ s'envoient sur $a$ et $b$ dans $S$.
Alors $J$ est la réunion de ses idéaux de type fini $J_\lambda \subset J$.
Posons $S_\lambda = R[A, B, x_1, \ldots, x_n]/J_\lambda$ et
notons $a_\lambda, b_\lambda \in S_\lambda$ les images de $A$ et $B$.
Alors, pour un certain $\lambda$, les fibres de
$$
f_{a_\lambda} : \Spec(S_\lambda) \longrightarrow \mathbf{A}^1_R
$$
sont de dimension $\leq d$, voir Limites, lemme \ref{limits-lemma-limit-dimension}.
Fixons un tel $\lambda$. Si l'on peut trouver un $n_0$ qui convient pour $R \to S_\lambda$,
$a_\lambda$, $b_\lambda$, alors $n_0$ convient pour $R \to S$. En effet, les fibres
de $f_{a_\lambda^n + b_\lambda} : \Spec(S_\lambda) \to \mathbf{A}^1_R$
contiennent les fibres de $f_{a^n + b} : \Spec(S) \to \mathbf{A}^1_R$.
On est ainsi ramené au cas traité dans le paragraphe suivant.

\medskip\noindent
Supposons que $R \to S$ soit de présentation finie. Dans ce paragraphe, on se ramène
au cas où $R$ est de type fini sur $\mathbf{Z}$. D'après
Algèbre, lemme \ref{algebra-lemma-limit-module-finite-presentation},
on peut trouver un ensemble ordonné filtrant $\Lambda$ et un système d'homomorphismes d'anneaux
$R_\lambda \to S_\lambda$ indexé par $\Lambda$, dont la colimite est
$R \to S$, tels que $S_\mu = S_\lambda \otimes_{R_\lambda} R_\mu$
pour $\mu \geq \lambda$ et que chacun des $R_\lambda$ et $S_\lambda$
soit de type fini sur $\mathbf{Z}$. Choisissons $\lambda_0 \in \Lambda$ et des
éléments $a_{\lambda_0}, b_{\lambda_0} \in S_{\lambda_0}$ qui s'envoient sur
$a, b \in S$. Pour $\lambda \geq \lambda_0$, notons
$a_\lambda, b_\lambda \in S_\lambda$
l'image de $a_{\lambda_0}, b_{\lambda_0}$.
Alors, pour $\lambda \geq \lambda_0$ assez grand, les fibres de
$$
f_{a_\lambda} : 
\Spec(S_\lambda) \longrightarrow \mathbf{A}^1_{R_\lambda}
$$
sont de dimension $\leq d$, voir
Limites, lemme \ref{limits-lemma-descend-dimension-d}.
Fixons un tel $\lambda$. Si l'on peut trouver un $n_0$ qui convient pour
$R_\lambda \to S_\lambda$, $a_\lambda$, $b_\lambda$, alors $n_0$
convient pour $R \to S$. En effet, toute fibre de
$f_{a^n + b} : \Spec(S) \to \mathbf{A}^1_R$
a la même dimension qu'une fibre de
$f_{a_\lambda^n + b_\lambda} : \Spec(S_\lambda) \to \mathbf{A}^1_{R_\lambda}$
d'après Morphismes, lemme \ref{morphisms-lemma-dimension-fibre-after-base-change}.
On est ainsi ramené au cas traité dans le paragraphe suivant.

\medskip\noindent
Supposons que $R$ et $S$ soient de type fini sur $\mathbf{Z}$. En particulier,
la dimension de $R$ est finie, et on peut raisonner par récurrence sur $\dim(R)$.
On peut donc supposer le résultat vrai pour toutes les situations où
$R' \to S'$, $a$, $b$ sont comme dans le lemme, avec $R'$ et $S'$ de type fini
sur $\mathbf{Z}$, mais où $\dim(R') < \dim(R)$.

\medskip\noindent
Puisque l'énoncé ne concerne que la topologie du spectre de $S$,
on peut supposer $S$ réduit. Soit $S^\nu$ la normalisation de $S$.
Alors $S \subset S^\nu$ est une extension finie, puisque $S$ est excellent, voir
Algèbre, proposition \ref{algebra-proposition-ubiquity-nagata},
et Morphismes, lemme \ref{morphisms-lemma-nagata-normalization}.
Ainsi $\Spec(S^\nu) \to \Spec(S)$ est surjectif et fini
(Algèbre, lemme \ref{algebra-lemma-integral-overring-surjective}).
Il s'ensuit que, si le résultat est vrai pour $R \to S^\nu$ et les
images de $a$, $b$ dans $S^\nu$, alors il est vrai pour $R \to S$, $a$, $b$.
(Un petit détail est omis.) On est ainsi ramené au cas traité dans le
paragraphe suivant.

\medskip\noindent
Supposons que $R$ et $S$ soient de type fini sur $\mathbf{Z}$ et que
$S$ soit normal. Alors $S = S_1 \times \ldots \times S_r$ pour certains
anneaux intègres normaux $S_i$. Si le résultat est vrai pour chaque $R \to S_i$
et pour les images de $a$, $b$ dans $S_i$, alors il est vrai pour
$R \to S$, $a$, $b$. (Un petit détail est omis.)
On est ainsi ramené au cas traité dans le paragraphe suivant.

\medskip\noindent
Supposons que $R$ et $S$ soient de type fini sur $\mathbf{Z}$ et que
$S$ soit un anneau intègre normal. On peut remplacer $R$ par l'image de $R$ dans $S$
(cela n'augmente pas la dimension de $R$).
On est ainsi ramené au cas traité dans le paragraphe suivant.

\medskip\noindent
Supposons que $R \subset S$ soient de type fini sur $\mathbf{Z}$ et que
$S$ soit un anneau intègre normal. Considérons le morphisme
$$
f_a : \Spec(S) \to \mathbf{A}^1_R
$$
L'hypothèse nous dit que $f_a$ a des fibres de dimension $\leq d$.
Les fibres de $f : \Spec(S) \to \Spec(R)$ sont donc de dimension $\leq d + 1$
(Morphismes, lemme \ref{morphisms-lemma-dimension-fibre-at-a-point-additive}).
Considérons le morphisme de schémas intègres
$$
\phi : \Spec(S) \to \mathbf{A}^2_R = \Spec(R[x, y])
$$
correspondant à l'homomorphisme de $R$-algèbres $R[x, y] \to S$ qui envoie $x$ sur $a$
et $y$ sur $b$. Il y a deux cas à considérer :
\begin{enumerate}
\item $\phi$ est dominant, et
\item $\phi$ n'est pas dominant.
\end{enumerate}
Nous affirmons que, dans les deux cas, il existe un entier $n_0$
et un ouvert non vide $V \subset \Spec(R)$ tels que, pour $n \geq n_0$,
les fibres de $f_{a^n + b}$ aux points $q \in \mathbf{A}^1_V$
soient de dimension $\leq d$.

\medskip\noindent
Démontrons l'affirmation dans le cas (1). On a $f_{a^n + b} = \pi_n \circ \phi$,
où
$$
\pi_n :  \mathbf{A}^2_R \to \mathbf{A}^1_R
$$
est le morphisme plat correspondant à l'homomorphisme de $R$-algèbres $R[x] \to R[x, y]$
qui envoie $x$ sur $x^n + y$. Puisque $\phi$ est dominant, il existe un ouvert dense
$U \subset \Spec(S)$ tel que $\phi|_U : U \to \mathbf{A}^2_R$ soit plat
(cela résulte par exemple de la platitude générique, voir
Morphismes, proposition \ref{morphisms-proposition-generic-flatness}).
Alors la composée
$$
f_{a^n + b}|_U :
U \xrightarrow{\phi|_U}
\mathbf{A}^2_R
\xrightarrow{\pi_n}
\mathbf{A}^1_R
$$
est elle aussi plate. Les fibres de ce morphisme sont donc de codimension au moins
$1$ dans les fibres de $f|_U : U \to \Spec(R)$ d'après
Morphismes, lemme \ref{morphisms-lemma-dimension-fibre-at-a-point-additive}.
Autrement dit, les fibres de $f_{a^n + b}|_U$ sont de dimension $\leq d$.
D'autre part, puisque $U$ est dense dans $\Spec(S)$, on peut trouver un ouvert non vide
$V \subset \Spec(R)$ tel que $U \cap f^{-1}(p) \subset f^{-1}(p)$
soit dense pour tout $p \in V$ (voir par exemple le lemme
\ref{lemma-nowhere-dense-generic-fibre}).
Ainsi $\dim(f^{-1}(p) \setminus U \cap f^{-1}(p)) \leq d$, et
on conclut que notre affirmation est vraie (puisque toute fibre de
$f_{a^n + b} : \Spec(S) \to \mathbf{A}^1_R$
est contenue dans une fibre de $f$).

\medskip\noindent
Cas (2). Dans ce cas, on peut trouver un élément non nul $g = \sum c_{ij} x^i y^j$
de $R[x, y]$ tel que $\Im(\phi) \subset V(g)$. On peut en fait supposer
$g$ irréductible sur $\text{Frac}(R)$. Si $g \in R[x]$, disons de coefficient
dominant $c$, alors, sur $V = D(c) \subset \Spec(R)$, les fibres de
$f$ sont déjà de dimension $\leq d$ (car l'image de $f_a$ est contenue
dans $V(g) \subset \mathbf{A}^1_R$, qui a des fibres finies sur $V$). On peut donc
supposer que $g$ n'appartient pas à $R[x]$.
Soit $s \geq 1$ le degré de $g$ comme polynôme en $y$, et
soit $t$ le degré de $\sum c_{is} x^i$ comme polynôme en $x$.
Alors $c_{ts}$ est non nul et
$$
g(x, -x^n) = (-1)^s c_{ts} x^{t + sn} + l.o.t.
$$
pourvu que $n$ soit supérieur au degré de $g$ comme polynôme en $x$
(un petit détail est omis). Pour un tel $n$, le polynôme $g(x, -x^n)$ 
est un polynôme non nul en $x$ et s'envoie sur un polynôme non nul de
$\kappa(\mathfrak p)[x]$ pour tout $\mathfrak p \subset R$,
$c_{st} \not \in \mathfrak p$. On conclut que notre affirmation est
vraie pour $V$ égal à l'ouvert principal $D(c_{ts})$ de $\Spec(R)$.

\medskip\noindent
Nous pouvons maintenant raisonner par récurrence sur $\dim(R)$. En effet, soit
$I \subset R$ un idéal tel que $V(I) = \Spec(R) \setminus V$.
Observons que $\dim(R/I) < \dim(R)$, puisque $R$ est un anneau intègre.
Soit $n'_0$ l'entier fourni par l'hypothèse de récurrence sur $\dim(R)$
pour $R/I \to S/IS$ et les images de $a$ et $b$ dans $S/IS$.
Alors $\max(n_0, n'_0)$ convient.
\end{proof}

\begin{lemma}
\label{lemma-weak-relative-noether-normalization}
Soit $R \to S$ un homomorphisme d'anneaux de type fini. Soit $d$ le maximum
des dimensions des fibres de $\Spec(S) \to \Spec(R)$.
Alors il existe un homomorphisme d'anneaux quasi-fini
$R[t_1, \ldots, t_d] \to S$.
\end{lemma}

\begin{proof}
Dans ce paragraphe, on se ramène au cas où $R \to S$ est de présentation
finie. En effet, écrivons $S = R[x_1, \ldots, x_n]/J$ pour un idéal
$J \subset R[x_1, \ldots, x_n]$.
Alors $J$ est la réunion de ses idéaux de type fini $J_\lambda \subset J$.
Posons $S_\lambda = R[x_1, \ldots, x_n]/J_\lambda$.
Alors, pour un certain $\lambda$, les fibres de
$\Spec(S_\lambda) \to \Spec(R)$
sont de dimension $\leq d$, voir Limites, lemme \ref{limits-lemma-limit-dimension}.
Fixons un tel $\lambda$. Si l'on peut trouver un homomorphisme quasi-fini
$R[t_1, \ldots, t_d] \to S_\lambda$, alors la composée
$R[t_1, \ldots, t_d] \to S$ est bien sûr quasi-finie.
On est ainsi ramené au cas traité dans le paragraphe suivant.

\medskip\noindent
Supposons que $R \to S$ soit de présentation finie. Dans ce paragraphe, on se ramène
au cas où $R$ est de type fini sur $\mathbf{Z}$. D'après
Algèbre, lemme \ref{algebra-lemma-limit-module-finite-presentation},
on peut trouver un ensemble ordonné filtrant $\Lambda$ et un système d'homomorphismes d'anneaux
$R_\lambda \to S_\lambda$ indexé par $\Lambda$, dont la colimite est
$R \to S$, tels que $S_\mu = S_\lambda \otimes_{R_\lambda} R_\mu$
pour $\mu \geq \lambda$ et
que chacun des $R_\lambda$ et $S_\lambda$ soit de type
fini sur $\mathbf{Z}$.
Alors, pour $\lambda$ assez grand, les fibres de
$\Spec(S_\lambda) \to \Spec(R_\lambda)$
sont de dimension $\leq d$, voir
Limites, lemme \ref{limits-lemma-descend-dimension-d}.
Fixons un tel $\lambda$. Si l'on peut trouver un homomorphisme d'anneaux quasi-fini
$R_\lambda[t_1, \ldots, t_d] \to S_\lambda$, alors le changement de base
$R[t_1, \ldots, t_d] \to S$ est lui aussi quasi-fini
(Algèbre, lemme \ref{algebra-lemma-quasi-finite-base-change}).
On est ainsi ramené au cas traité dans le paragraphe suivant.

\medskip\noindent
Supposons que $R$ et $S$ soient de type fini sur $\mathbf{Z}$. Si $d = 0$, alors
l'homomorphisme d'anneaux est quasi-fini et la démonstration est achevée. Supposons $d > 0$.
Nous allons trouver un élément $a \in S$ tel que
les fibres de l'homomorphisme de $R$-algèbres $R[x] \to S$, $x \mapsto a$,
soient de dimension $< d$. Cela achèvera la démonstration par récurrence.

\medskip\noindent
Nous allons démontrer l'existence de $a$ par récurrence sur $\dim(R)$.

\medskip\noindent
Soient $\mathfrak q_1, \ldots, \mathfrak q_r \subset S$ ceux
des idéaux premiers minimaux de $S$ tels que $\dim_{\mathfrak q_i}(S/R) = d$.
Pour la notation, voir
Algèbre, définition \ref{algebra-definition-relative-dimension}.
Disons que $\mathfrak q_i$ est au-dessus de l'idéal premier $\mathfrak p_i \subset R$.
On a $\text{trdeg}_{\kappa(\mathfrak p_i)}(\kappa(\mathfrak q_i)) = d$,
car $\mathfrak q_i$ est un point générique de sa fibre ; on peut par exemple
appliquer Algèbre, lemme \ref{algebra-lemma-dimension-at-a-point-finite-type-field},
à $S \otimes_R \kappa(\mathfrak p_i)$.
Ainsi, d'après le lemme \ref{lemma-silly}, on peut trouver un élément $a \in S$ tel que
l'image de $a$ dans $\kappa(\mathfrak q_i)$ soit transcendante sur
$\kappa(\mathfrak p_i)$ pour $i = 1, \ldots, r$.
Considérons le morphisme
$$
f_a : \Spec(S) \longrightarrow \mathbf{A}^1_R
$$
correspondant à l'homomorphisme de $R$-algèbres $R[x] \to S$
qui envoie $x$ sur $a$. Soit $U \subset \Spec(S)$ l'ouvert
où les fibres sont de dimension $\leq d - 1$, voir
Morphismes, lemme \ref{morphisms-lemma-openness-bounded-dimension-fibres}.
Par construction, $U$ contient tous les points génériques de $\Spec(S)$.
En particulier, on voit que $U$ contient tous les points génériques de
toutes les fibres génériques de $\Spec(S) \to \Spec(R)$, puisque de tels points sont
nécessairement des points génériques de $\Spec(S)$. Posons $T = \Spec(S) \setminus U$,
considéré comme un sous-schéma fermé réduit de $\Spec(S)$.
Il résulte de ce que l'on vient de dire et de l'hypothèse
$\dim(S/R) \leq d$ que
les fibres génériques de $T \to \Spec(R)$ sont de dimension $\leq d - 1$.
Ainsi, d'après le lemme \ref{lemma-dimension-in-neighbourhood},
appliqué plusieurs fois pour produire des voisinages ouverts des points
génériques de $\Spec(R)$, on peut trouver un ouvert dense $V \subset \Spec(R)$ tel que
$T_V \to V$ ait des fibres de dimension $\leq d - 1$.
On conclut que, pour $q \in \mathbf{A}^1_V$, la fibre de
$f_a$ au-dessus de $q$ est de dimension $\leq d - 1$ (car on a borné la dimension
de la fibre de $f_a|_U$ et de celle de $f_a|_T$).

\medskip\noindent
Par évitement des idéaux premiers, on peut supposer que $V = D(f)$ pour un certain $f \in R$.
On voit alors que l'homomorphisme d'anneaux $R_f[x] \to S_f$, $x \mapsto a$,
a des fibres de dimension $\leq d - 1$.
On peut remplacer $a$ par $fa$ et supposer $a \in (f)$.
Par récurrence sur $\dim(R)$, on peut trouver un élément
$\overline{b} \in S/fS$
tel que les fibres de $\Spec(S/fS) \to \Spec(R/fR[x])$,
$x \mapsto \overline{b}$, soient de dimension $\leq d - 1$.
Soit $b \in S$ un relèvement de $\overline{b}$.
D'après le lemme \ref{lemma-sum-is-ok}, il existe un $n > 0$ tel
que $a^n + b$ convienne encore pour $R_f \to S_f$.
D'autre part, l'image de $a^n + b$ dans $S/fS$ est $\overline{b}$,
et la démonstration est achevée.
\end{proof}

\section{Théorèmes de Bertini}
\label{section-bertini}

\noindent
Nous poursuivons la discussion commencée dans
Variétés, section \ref{varieties-section-bertini}.
Dans cette section, nous démontrons que les sections hyperplanes générales
des variétés géométriquement irréductibles sont géométriquement irréductibles,
en suivant l'argument remarquable donné dans \cite{Jou}.

\begin{lemma}
\label{lemma-amazing-bertini-lemma}
\begin{reference}
Voir les pages 71 et 72 de \cite{Jou}
\end{reference}
Soit $K/k$ une extension de corps géométriquement irréductible et de type fini.
Soit $n \geq 1$.
Soient $g_1, \ldots, g_n \in K$ des éléments tels qu'il
existe $c_1, \ldots, c_n \in k$ pour lesquels les éléments
$$
x_1, \ldots, x_n, \sum g_ix_i, \sum c_ig_i \in K(x_1, \ldots, x_n)
$$
soient algébriquement indépendants sur $k$. Alors
$K(x_1, \ldots, x_n)$ est géométriquement irréductible sur
$k(x_1, \ldots, x_n, \sum g_ix_i)$.
\end{lemma}

\begin{proof}
Soient $c_1, \ldots, c_n \in k$ comme dans l'énoncé du lemme.
Écrivons $\xi = \sum g_ix_i$ et $\delta = \sum c_ig_i$.
Pour $a \in k$, considérons l'automorphisme $\sigma_a$
de $K(x_1, \ldots, x_n)$ qui est l'identité sur $K$ et qui est donné par les règles
$$
\sigma_a(x_i) = x_i + a c_i
$$
Observons que $\sigma_a(\xi) = \xi + a \delta$ et $\sigma_a(\delta) = \delta$.
Considérons la tour d'extensions de corps
$$
K_0 = k(x_1, \ldots, x_n) \subset
K_1 = K_0(\xi) \subset
K_2 = K_0(\xi, \delta) \subset K(x_1, \ldots, x_n) = \Omega
$$
Observons que $\sigma_a(K_0) = K_0$ et $\sigma_a(K_2) = K_2$.
Soit $\theta \in \Omega$ un élément algébrique séparable sur $K_1$.
Nous devons montrer que $\theta \in K_1$, voir Algèbre, lemme
\ref{algebra-lemma-geometrically-irreducible-separable-elements}.

\medskip\noindent
Notons $K'_2$ la fermeture algébrique séparable de $K_2$ dans $\Omega$.
Comme $K'_2/K_2$ est finie (Algèbre, lemme
\ref{algebra-lemma-make-geometrically-irreducible}) et
séparable, il n'y a qu'un nombre fini de corps intermédiaires entre
$K'_2$ et $K_2$ (Corps, lemme \ref{fields-lemma-primitive-element}).
Si $k$ est infini\footnote{Nous traiterons le cas d'un corps fini dans le
dernier paragraphe de la démonstration.}, on peut alors trouver des éléments distincts
$a_1, a_2$ de $k$ tels que
$$
K_2(\sigma_{a_1}(\theta)) = K_2(\sigma_{a_2}(\theta))
$$
comme sous-corps de $\Omega$. Écrivons $\theta_i = \sigma_{a_i}(\theta)$
et $\xi_i = \sigma_{a_i}(\xi) = \xi + a_i \delta$. Observons que
$$
K_2 = K_0(\xi_1, \xi_2)
$$
car on a $\xi_i = \xi + a_i \delta$,
$\xi = (a_2 \xi_1 - a_1 \xi_2)/(a_2 - a_1)$, et
$\delta = (\xi_1 - \xi_2)/(a_1 - a_2)$.
Puisque $K_2/K_0$ est purement transcendante de degré $2$, on conclut
que $\xi_1$ et $\xi_2$ sont algébriquement indépendants sur $K_0$.
Comme $\theta_1$ est algébrique sur $K_0(\xi_1)$, on conclut que
$\xi_2$ est transcendant sur $K_0(\xi_1, \theta_1)$.

\medskip\noindent
Par hypothèse, $K/k$ est géométriquement irréductible. Cela implique
que $K(x_1, \ldots, x_n)/K_0$ est géométriquement irréductible
(Algèbre, lemme
\ref{algebra-lemma-geometrically-irreducible-base-change-transcendental}).
Cela implique à son tour que $K_0(\xi_1, \theta_1)/K_0$
est géométriquement irréductible en tant que sous-extension
(Algèbre, lemme \ref{algebra-lemma-subalgebra-geometrically-irreducible}).
Puisque $\xi_2$ est transcendant sur $K_0(\xi_1, \theta_1)$,
on conclut que $K_0(\xi_1, \xi_2, \theta_1)/K_0(\xi_2)$
est géométriquement irréductible (Algèbre, lemme
\ref{algebra-lemma-geometrically-irreducible-add-transcendental}).
Par notre choix de $a_1, a_2$ ci-dessus, on a
$$
K_0(\xi_1, \xi_2, \theta_1) =
K_2(\sigma_{a_1}(\theta)) =
K_2(\sigma_{a_2}(\theta)) =
K_0(\xi_1, \xi_2, \theta_2)
$$
Comme $\theta_2$ est algébrique séparable sur $K_0(\xi_2)$,
on conclut à nouveau, d'après Algèbre, lemme
\ref{algebra-lemma-geometrically-irreducible-separable-elements}, que
$\theta_2 \in K_0(\xi_2)$. En appliquant $\sigma_{a_2}^{-1}$
à cette relation, on obtient $\theta \in K_0(\xi) = K_1$, comme souhaité.

\medskip\noindent
Cela achève la démonstration lorsque $k$ est infini. Si $k$ est fini,
on peut choisir une variable $t$ et considérer l'extension
$K(t)/k(t)$, qui est géométriquement irréductible d'après
Algèbre, lemme
\ref{algebra-lemma-geometrically-irreducible-base-change-transcendental}.
Comme il reste vrai que
$x_1, \ldots, x_n, \sum g_ix_i, \sum c_ig_i$
dans $K(t, x_1, \ldots, x_n)$ sont algébriquement indépendants sur $k(t)$,
on conclut que $K(t, x_1, \ldots, x_n)$
est géométriquement irréductible sur
$k(t, x_1, \ldots, x_n, \sum g_ix_i)$
par l'argument déjà donné.
Il suffit alors d'utiliser encore une fois Algèbre, lemme
\ref{algebra-lemma-geometrically-irreducible-base-change-transcendental},
pour achever la démonstration.
\end{proof}

\begin{lemma}
\label{lemma-algebraically-independent}
Soit $A$ un anneau intègre de type fini sur un corps $k$. Soit $n \geq 2$.
Soient $g_1, \ldots, g_n \in A$ des éléments tels que $V(g_1, g_2)$
ait une composante irréductible de dimension $\dim(A) - 2$.
Alors il existe $c_1, \ldots, c_n \in k$ tels que les éléments
$$
x_1, \ldots, x_n, \sum g_ix_i, \sum c_ig_i \in
\text{Frac}(A)(x_1, \ldots, x_n)
$$
soient algébriquement indépendants sur $k$.
\end{lemma}

\begin{proof}
L'indépendance algébrique sur $k$ signifie que le morphisme
$$
T = \Spec(A[x_1, \ldots, x_n])
\longrightarrow
\Spec(k[x_1, \ldots, x_n, y, z]) = S
$$
donné par $y = \sum g_ix_i$ et $z = \sum c_ig_i$ est dominant.
Posons $d = \dim(A)$. Si $T \to S$ n'est pas dominant, alors son image
est de dimension $< n + 2$, et toute composante irréductible
de toute fibre est donc de dimension $> d + n - (n + 2) = d - 2$, voir
Variétés, lemme
\ref{varieties-lemma-dimension-fibres-locally-algebraic}.
Choisissons un point fermé $u \in V(g_1, g_2)$ contenu dans une composante irréductible
de dimension $d - 2$ et dans aucune autre composante de $V(g_1, g_2)$.
Considérons le point fermé $t = (u, 1, 0, \ldots 0)$
de $T$ situé au-dessus de $u$. Posons $(c_1, \ldots, c_n) = (0, 1, 0, \ldots, 0)$.
Alors $t$ s'envoie sur le point $s = (1, 0, \ldots, 0)$ de $S$.
La fibre de $T \to S$ au-dessus de $s$ est définie par
$$
x_1 - 1, x_2, \ldots, x_n, \sum x_ig_i, g_2
$$
et est donc, de manière équivalente, définie par
$$
x_1 - 1, x_2, \ldots, x_n, g_1, g_2
$$
D'après notre condition sur $g_1, g_2$, ce sous-schéma possède une composante irréductible
de dimension $d - 2$.
\end{proof}

\begin{lemma}
\label{lemma-bertini-irreducible}
\begin{reference}
\cite[théorème 6.3, partie 4)]{Jou}
\end{reference}
Dans Variétés, situation \ref{varieties-situation-family-divisors}, supposons que
\begin{enumerate}
\item $X$ soit de type fini sur $k$,
\item $X$ soit géométriquement irréductible sur $k$,
\item il existe $v_1, v_2, v_3 \in V$ et une composante irréductible
$Z$ de $H_{v_2} \cap H_{v_3}$ telles que $Z \not \subset H_{v_1}$ et
$\text{codim}(Z, X) = 2$, et
\item toute composante irréductible $Y$ de $\bigcap_{v \in V} H_v$
ait $\text{codim}(Y, X) \geq 2$.
\end{enumerate}
Alors, pour $v \in V \otimes_k k'$ général,
le schéma $H_v$ est géométriquement irréductible sur $k'$.
\end{lemma}

\begin{proof}
Pour que l'hypothèse (3) soit satisfaite, les éléments $v_1, v_2, v_3$
doivent être linéairement indépendants sur $k$ dans $V$ (un petit détail est omis).
On peut donc choisir une base $v_1, \ldots, v_r$ de $V$ contenant
ces éléments comme ses $3$ premiers éléments. Rappelons que
$H_{univ} \subset \mathbf{A}^r_k \times_k X$ est le « diviseur universel ».
Considérons la projection $q : H_{univ} \to \mathbf{A}^r_k$
dont les fibres schématiques sont les diviseurs $H_v$. D'après le
lemme \ref{lemma-geom-irreducible-generic-fibre},
il suffit de montrer que la fibre générique de $q$ est géométriquement
irréductible.
Pour cela, on peut remplacer $X$ par son réduit ; on peut donc
supposer que $X$ est un schéma intègre de type fini sur $k$.

\medskip\noindent
Soit $U \subset X$ un ouvert affine non vide tel que
$\mathcal{L}|_U \cong \mathcal{O}_U$. Écrivons $U = \Spec(A)$.
Notons $f_i \in A$ l'élément correspondant à la section $\psi(v_i)|_U$
via l'isomorphisme $\mathcal{L}|_U \cong \mathcal{O}_U$.
Alors $H_{univ} \cap (\mathbf{A}^r_k \times_k U)$ est donné par
$$
H_U = \Spec(A[x_1, \ldots, x_r]/(x_1f_1 + \ldots + x_rf_r))
$$
Par notre choix de la base, on voit que $f_1$ ne peut pas être nul, car cela
signifierait que $v_1 = 0$ et donc que $H_{v_1} = X$, ce qui contredit
l'hypothèse (3). Ainsi $\sum x_if_i$ est un non-diviseur de zéro dans
$A[x_1, \ldots, x_r]$.
Il s'ensuit que toute composante irréductible de $H_U$ est de dimension
$d + r - 1$, où $d = \dim(X) = \dim(A)$. Si $U' = U \cap D(f_1)$,
alors on voit que
$$
\begin{aligned}
H_{U'} &= \Spec(A_{f_1}[x_1, \ldots, x_r]/(x_1f_1 + \ldots + x_rf_r)) \\
&\cong \Spec(A_{f_1}[x_2, \ldots x_r]) = \mathbf{A}^{r - 1}_k \times_k U'
\end{aligned}
$$
est irréductible. D'autre part, on a
$$
H_U \setminus H_{U'} =
\Spec(A/(f_1)[x_1, \ldots, x_r]/(x_2f_2 + \ldots + x_rf_r))
$$
qui est de dimension au plus $d + r - 2$. En effet, pour $i \not = 1$,
le schéma $(H_U \setminus H_{U'}) \times_U D(f_i)$ est soit vide
(si $f_i = 0$), soit, par le même argument que ci-dessus, isomorphe à un espace affine de dimension $r - 1$
sur un ouvert de $\Spec(A/(f_1))$, et est donc de
dimension au plus $d + r - 2$.
D'autre part, $(H_U \setminus H_{U'}) \times_U V(f_2, \ldots, f_r)$
est un espace affine de dimension $r$ sur $\Spec(A/(f_1, \ldots, f_r))$ ;
l'hypothèse (4) nous dit donc qu'il est de dimension au plus
$d + r - 2$. On conclut que $H_U$ est irréductible pour tout $U$ comme ci-dessus.
Il s'ensuit que $H_{univ}$ est irréductible.

\medskip\noindent
Il suffit donc de montrer que le point générique de $H_{univ}$
est géométriquement irréductible sur le point générique de $\mathbf{A}^r_k$, voir
Variétés, lemme
\ref{varieties-lemma-geometrically-irreducible-function-field}.
Choisissons un ouvert affine non vide $U = \Spec(A)$ de $X$,
contenu dans $X \setminus H_{v_1}$ et qui rencontre la composante
irréductible $Z$ de $H_{v_2} \cap H_{v_3}$ dont l'existence est affirmée dans
l'hypothèse (3). Avec les notations ci-dessus, il faut démontrer que
l'extension de corps
$$
\text{Frac}(A[x_1, \ldots, x_r]/(x_1f_1 + \ldots + x_rf_r))/
k(x_1, \ldots , x_r)
$$
est géométriquement irréductible. Observons que $f_1$ est inversible dans $A$
par notre choix de $U$.
Posons $K = \text{Frac}(A)$, le corps des fractions de $A$.
En éliminant la variable $x_1$ comme ci-dessus,
on trouve qu'il faut montrer que l'extension de corps
$$
K(x_2, \ldots, x_r)/
k(x_2, \ldots, x_r, -\sum\nolimits_{i = 2, \ldots, r} f_1^{-1}f_i x_i)
$$
est géométriquement irréductible. D'après le lemme \ref{lemma-amazing-bertini-lemma},
il suffit de montrer que, pour certains $c_2, \ldots, c_r \in k$, les éléments
$$
x_2, \ldots, x_r, \sum\nolimits_{i = 2, \ldots, r} f_1^{-1}f_i x_i,
\sum\nolimits_{i = 2, \ldots, r} c_if_1^{-1}f_i
$$
sont algébriquement indépendants sur $k$ dans le corps des fractions de
$A[x_2, \ldots, x_r]$. Cela résulte du
lemme \ref{lemma-algebraically-independent}
et du fait que $Z \cap U$ est une composante irréductible
de $V(f_1^{-1}f_2, f_1^{-1}f_3) \subset U$.
\end{proof}

\begin{remark}
\label{remark-geometric-proof-bertini-irreducible}
Esquissons une démonstration « géométrique » d'un cas particulier du
lemme \ref{lemma-bertini-irreducible}. Supposons que $k$
soit un corps algébriquement clos et que $X \subset \mathbf{P}^n_k$ soit lisse et
irréductible de dimension $\geq 2$. Nous affirmons qu'il existe un
hyperplan $H \subset \mathbf{P}^n_k$ tel que $X \cap H$ soit
lisse et irréductible. En effet, d'après
Variétés, lemme \ref{varieties-lemma-bertini},
pour $v \in V = kT_0 \oplus \ldots \oplus kT_n$ général,
la section hyperplane correspondante $X \cap H_v$ est lisse.
D'autre part, d'après Enriques--Severi--Zariski, le schéma
$X \cap H_v$ est connexe, voir
Variétés, lemme \ref{varieties-lemma-connectedness-ample-divisor}.
Ainsi $X \cap H_v$ est lisse et irréductible.
\end{remark}

\section{Théorème du cube}
\label{section-theorem-cube}

\noindent
Le lemme suivant nous dit que la diagonale du foncteur de Picard
est représentable par des immersions localement fermées sous
les hypothèses faites dans le lemme.

\begin{lemma}
\label{lemma-diagonal-picard-flat-proper}
Soit $f : X \to S$ un morphisme plat et propre de présentation finie.
Soit $\mathcal{E}$ un $\mathcal{O}_X$-module localement libre de rang fini.
Pour un morphisme $g : T \to S$, considérons le diagramme de changement de base
$$
\xymatrix{
X_T \ar[d]_p \ar[r]_q & X \ar[d]^f \\
T \ar[r]^g & S
}
$$
Supposons que $\mathcal{O}_T \to p_*\mathcal{O}_{X_T}$ soit un
isomorphisme pour tout $g : T \to S$. Alors il existe une
immersion $j : Z \to S$ de présentation finie telle qu'un
morphisme $g : T \to S$ se factorise par $Z$ si et seulement s'il
existe un $\mathcal{O}_T$-module localement libre de rang fini $\mathcal{N}$
tel que $p^*\mathcal{N} \cong q^*\mathcal{E}$.
\end{lemma}

\begin{proof}
Observons que les fibres $X_s$ de $f$ sont connexes par notre hypothèse
$H^0(X_s, \mathcal{O}_{X_s}) = \kappa(s)$. Le rang de
$\mathcal{E}$ est donc constant sur les fibres. Puisque $f$ est ouvert
(Morphismes, lemme \ref{morphisms-lemma-fppf-open}) et fermé,
on conclut qu'il existe une décomposition $S = \coprod S_r$
de $S$ en sous-schémas ouverts et fermés telle que $\mathcal{E}$
soit de rang constant $r$ sur l'image réciproque de $S_r$.
On peut donc supposer que $\mathcal{E}$ est de rang constant $r$.
On notera $\mathcal{E}^\vee = \SheafHom(\mathcal{E}, \mathcal{O}_X)$
le module dual, de rang $r$.

\medskip\noindent
D'après la cohomologie et le changement de base (plus précisément, d'après
Catégories dérivées de schémas, lemme
\ref{perfect-lemma-flat-proper-perfect-direct-image-general}),
on voit que $E = Rf_*\mathcal{E}$ est un objet parfait de la
catégorie dérivée de $S$ et que sa formation commute au
changement de base arbitraire. Il en va de même pour $E' = Rf_*\mathcal{E}^\vee$.
Comme il n'y a jamais de cohomologie en degrés $< 0$, on voit que
$E$ et $E'$ ont (localement) une tor-amplitude dans $[0, b]$ pour un certain $b$.
Observons que, pour tout $g : T \to S$, on a
$p_*(q^*\mathcal{E}) = H^0(Lg^*E)$ et
$p_*(q^*\mathcal{E}^\vee) = H^0(Lg^*E')$.
Soient $j : Z \to S$ et $j' : Z' \to S$ les immersions
de présentation finie construites dans Catégories dérivées de schémas, lemme
\ref{perfect-lemma-locally-closed-where-H0-locally-free},
pour $E$ et $E'$, avec $a = 0$ et $r = r$ ; elles sont, en gros,
caractérisées par la propriété que $H^0(Lj^*E)$ et $H^0((j')^*E')$
sont des modules localement libres de rang fini compatibles à l'image réciproque.

\medskip\noindent
Soit $g : T \to S$ un morphisme. S'il existe un $\mathcal{N}$
comme dans le lemme, alors, en utilisant la formule de projection
Cohomologie, lemme \ref{cohomology-lemma-projection-formula},
on voit que les modules
$$
p_*(q^*\mathcal{E}) \cong
p_*(p^*\mathcal{N}) \cong
\mathcal{N} \otimes_{\mathcal{O}_T} p_*\mathcal{O}_{X_T} \cong
\mathcal{N}\quad\text{et de même }\quad
p_*(q^*\mathcal{E}^\vee) \cong \mathcal{N}^\vee
$$
sont des modules localement libres de rang $r$
et le restent après tout changement de base ultérieur $T' \to T$.
Dans ce cas, $T \to S$ se factorise donc par $j$ et par $j'$.
On peut ainsi remplacer $S$ par $Z \times_S Z'$ et supposer que
$f_*\mathcal{E}$ et $f_*\mathcal{E}^\vee$ sont des
$\mathcal{O}_S$-modules localement libres de rang $r$
dont la formation commute au changement de base arbitraire
(un petit détail est omis).

\medskip\noindent
Dans cette situation, si $g : T \to S$ est un morphisme et s'il existe un
$\mathcal{N}$ comme dans le lemme, alors l'application (le cup-produit en degré $0$)
$$
p_*(q^*\mathcal{E})
\otimes_{\mathcal{O}_T}
p_*(q^*\mathcal{E}^\vee)
\longrightarrow \mathcal{O}_T
$$
est un accouplement parfait. Réciproquement, si cette application de cup-produit est un
accouplement parfait, on voit que, localement sur $T$, on peut choisir une
base de sections
$\sigma_1, \ldots, \sigma_r$ dans $p_*(q^*\mathcal{E})$
et $\tau_1, \ldots, \tau_r$ dans $p_*(q^*\mathcal{E}^\vee)$
dont les produits vérifient $\sigma_i \tau_j = \delta_{ij}$.
En considérant $\sigma_i$ comme une section de $q^*\mathcal{E}$ sur $X_T$
et $\tau_j$ comme une section de $q^*\mathcal{E}^\vee$ sur $X_T$,
on conclut que
$$
\sigma_1, \ldots, \sigma_r :
\mathcal{O}_{X_T}^{\oplus r}
\longrightarrow
q^*\mathcal{E}
$$
est un isomorphisme dont l'inverse est donné par
$$
\tau_1, \ldots, \tau_r :
q^*\mathcal{E}
\longrightarrow
\mathcal{O}_{X_T}^{\oplus r}
$$
Autrement dit, on voit que $p^*p_*q^*\mathcal{E} \cong q^*\mathcal{E}$.
Or, la condition que le cup-produit soit non dégénéré définit
un sous-schéma ouvert rétrocompact (à savoir, le lieu où un déterminant
convenable est non nul), et la démonstration est achevée.
\end{proof}

\noindent
Le lemme précédent nous dit en particulier que, si un fibré vectoriel est
trivial sur les fibres d'une famille propre et plate d'espaces propres, alors
il est l'image réciproque d'un fibré vectoriel. Détaillons un peu cela.

\begin{lemma}
\label{lemma-trivial-on-fibres}
Soit $f : X \to S$ un morphisme plat et propre de présentation finie
tel que $f_*\mathcal{O}_X = \mathcal{O}_S$, et que cela reste
vrai après tout changement de base. Soit $\mathcal{E}$ un
$\mathcal{O}_X$-module localement libre de rang fini. Supposons que
\begin{enumerate}
\item $\mathcal{E}|_{X_s}$ soit isomorphe à
$\mathcal{O}_{X_s}^{\oplus r_s}$ pour tout $s \in S$, et
\item $S$ soit réduit.
\end{enumerate}
Alors $\mathcal{E} = f^*\mathcal{N}$ pour un certain
$\mathcal{O}_S$-module localement libre de rang fini $\mathcal{N}$.
\end{lemma}

\begin{proof}
En effet, dans ce cas, l'immersion localement fermée $j : Z \to S$ du
lemme \ref{lemma-diagonal-picard-flat-proper}
est bijective, et est donc une
immersion fermée. Mais puisque $S$ est réduit, $j$ est un isomorphisme.
\end{proof}

\begin{lemma}
\label{lemma-triviality-generic-fibre-valuation-ring}
Soit $f : X \to S$ un morphisme propre et plat de présentation finie.
Soit $\mathcal{L}$ un $\mathcal{O}_X$-module inversible.
Supposons que
\begin{enumerate}
\item $S$ soit le spectre d'un anneau de valuation,
\item $\mathcal{L}$ soit trivial sur la fibre générique $X_\eta$ de $f$,
\item la fibre fermée $X_0$ de $f$ soit intègre,
\item $H^0(X_\eta, \mathcal{O}_{X_\eta})$ soit égal au corps des fonctions de $S$.
\end{enumerate}
Alors $\mathcal{L}$ est trivial.
\end{lemma}

\begin{proof}
Écrivons $S = \Spec(A)$. Nous allons d'abord démontrer le lemme lorsque $A$ est un
anneau de valuation discrète (car c'est le cas le plus souvent utilisé en pratique).
Soit $\pi \in A$ une uniformisante.
Prenons une section trivialisante $s \in \Gamma(X_\eta, \mathcal{L}_\eta)$.
Après avoir remplacé $s$ par $\pi^n s$ si nécessaire,
on peut supposer que $s \in \Gamma(X, \mathcal{L})$.
Si $s|_{X_0} = 0$, on voit alors
que $s$ est divisible par $\pi$ (parce que $X_0$ est la fibre
schématique et que $X$ est plat sur $A$). On peut donc supposer que $s|_{X_0}$
est non nul. Le lieu des zéros $Z(s)$ de $s$ est alors contenu dans $X_0$,
mais ne contient pas le point générique de $X_0$ (parce que $X_0$ est intègre).
Cela signifie que $Z(s)$ est de codimension $\geq 2$ dans $X$, ce qui contredit
Diviseurs, lemme \ref{divisors-lemma-effective-Cartier-codimension-1},
à moins que $Z(s) = \emptyset$, comme souhaité.

\medskip\noindent
Démonstration dans le cas général. Puisque l'anneau de valuation $A$ est
cohérent (Algèbre, exemple \ref{algebra-example-valuation-ring-coherent}),
on voit que $H^0(X, \mathcal{L})$ est un $A$-module cohérent, voir
Catégories dérivées de schémas, lemme
\ref{perfect-lemma-cohomology-over-coherent-ring}.
De manière équivalente, $H^0(X, \mathcal{L})$ est un
$A$-module de présentation finie (Algèbre, lemme \ref{algebra-lemma-coherent-ring}).
Puisque $H^0(X, \mathcal{L})$ est sans torsion (par platitude de $X$ sur $A$),
on déduit de Plus sur l'algèbre, lemme
\ref{more-algebra-lemma-generalized-valuation-ring-modules},
que $H^0(X, \mathcal{L}) = A^{\oplus n}$ pour un certain $n$.
Par changement de base plat (Cohomologie des schémas, lemme
\ref{coherent-lemma-flat-base-change-cohomology}),
on a
$$
K = H^0(X_\eta, \mathcal{O}_{X_\eta}) \cong
H^0(X_\eta, \mathcal{L}_\eta) =
H^0(X, \mathcal{L}) \otimes_A K
$$
où $K$ est le corps des fractions de $A$. Ainsi $n = 1$.
Choisissons un générateur $s \in H^0(X, \mathcal{L})$.
Soit $\mathfrak m \subset A$ l'idéal maximal.
Alors $\kappa = A/\mathfrak m = \colim A/\pi$, où
il s'agit de la colimite filtrante sur les $\pi \in \mathfrak m$ non nuls
(on utilise ici le fait que $A$ est un anneau de valuation).
Ainsi $X_0 = \lim X \times_S \Spec(A/\pi)$.
Si $s|_{X_0}$ est nul, alors, pour un certain $\pi$,
on voit que la restriction de $s$ à $X \times_S \Spec(A/\pi)$ est nulle, voir
Limites, lemme \ref{limits-lemma-descend-section}.
Mais si cela se produit, alors $\pi^{-1} s$ est
une section globale de $\mathcal{L}$, ce qui contredit
le fait que $s$ soit un générateur de $H^0(X, \mathcal{L})$.
Ainsi $s|_{X_0}$ n'est pas nul. Soit $Z(s) \subset X$ le schéma des zéros de $s$.
Puisque $s|_{X_0}$ n'est pas nul et que $X_0$ est intègre,
on voit que $Z(s)_0 \subset X_0$ est un diviseur de Cartier effectif.
Puisque $f$ est propre et que $S$ est local, tout point de $Z(s)$
se spécialise en un point de $Z(s)_0$. Ainsi, d'après
Diviseurs, lemme \ref{divisors-lemma-fibre-Cartier}, partie (3),
on voit que $Z(s)$ est un diviseur de Cartier effectif relatif ;
en particulier, $Z(s) \to S$ est plat.
Par conséquent, si $Z(s)$ était non vide, alors $Z(s)_\eta$ serait non vide,
ce qui contredit le fait que $s|_{X_\eta}$ soit une trivialisation
de $\mathcal{L}_\eta$. Ainsi $Z(s) = \emptyset$, comme souhaité.
\end{proof}

\begin{lemma}
\label{lemma-get-a-closed}
Soient $f : X \to S$ et $\mathcal{E}$ comme dans le
lemme \ref{lemma-diagonal-picard-flat-proper},
et supposons en outre que $\mathcal{E}$ soit un $\mathcal{O}_X$-module inversible.
Si, de plus, les fibres géométriques de $f$ sont
intègres, alors $Z$ est fermé dans $S$.
\end{lemma}

\begin{proof}
Puisque $j : Z \to S$ est de présentation finie, il suffit
de montrer ceci : pour tout morphisme $g : \Spec(A) \to S$, où $A$ est un
anneau de valuation de corps des fractions $K$, tel que $g(\Spec(K)) \in j(Z)$,
on a $g(\Spec(A)) \subset j(Z)$. Voir
Morphismes, lemme \ref{morphisms-lemma-reach-points-scheme-theoretic-image}.
Cela résulte du lemme \ref{lemma-triviality-generic-fibre-valuation-ring}
et de la caractérisation de $j : Z \to S$ dans le
lemme \ref{lemma-diagonal-picard-flat-proper}.
\end{proof}

\begin{lemma}
\label{lemma-H1-O-picard-flat-proper}
Considérons un diagramme commutatif de schémas
$$
\xymatrix{
X' \ar[rr] \ar[dr]_{f'} & & X \ar[dl]^f \\
& S
}
$$
où $f' : X' \to S$ et $f : X \to S$ satisfont les hypothèses du
lemme \ref{lemma-diagonal-picard-flat-proper}.
Soit $\mathcal{L}$ un $\mathcal{O}_X$-module inversible,
et soit $\mathcal{L}'$ son image réciproque sur $X'$. Soit $Z \subset S$,
resp.\ $Z' \subset S$, le sous-schéma localement fermé construit
dans le lemme \ref{lemma-diagonal-picard-flat-proper}
pour $(f, \mathcal{L})$, resp.\ $(f', \mathcal{L}')$,
de sorte que $Z \subset Z'$. Si $s \in Z$ et si
$$
H^1(X_s, \mathcal{O}) \longrightarrow H^1(X'_s, \mathcal{O})
$$
est injectif, alors $Z \cap U = Z' \cap U$ pour un certain voisinage ouvert
$U$ de $s$.
\end{lemma}

\begin{proof}
On peut remplacer $S$ par $Z'$. Après avoir remplacé $S$ par un voisinage ouvert affine
de $s$, on peut supposer que $\mathcal{L}' = \mathcal{O}_{X'}$.
Soient $E = Rf_*\mathcal{L}$ et $E' = Rf'_*\mathcal{L}' = Rf'_*\mathcal{O}_{X'}$.
Ce sont des complexes parfaits dont la formation commute à tout
changement de base (Catégories dérivées de schémas, lemme
\ref{perfect-lemma-flat-proper-perfect-direct-image-general}).
En particulier, on voit que
$$
E \otimes_{\mathcal{O}_S}^\mathbf{L} \kappa(s) =
R\Gamma(X_s, \mathcal{L}_s) = R\Gamma(X_s, \mathcal{O}_{X_s})
$$
La seconde égalité vient de ce que $s \in Z$. Posons
$h_i = \dim_{\kappa(s)} H^i(X_s, \mathcal{O}_{X_s})$.
Après avoir rétréci $S$, on peut représenter $E$ par un complexe
$$
\mathcal{O}_S \to \mathcal{O}_S^{\oplus h_1} \to
\mathcal{O}_S^{\oplus h_2} \to \ldots
$$
voir Plus sur l'algèbre, lemme
\ref{more-algebra-lemma-lift-perfect-from-residue-field}
(à strictement parler, cela utilise aussi
Catégories dérivées de schémas, lemmes
\ref{perfect-lemma-affine-compare-bounded} et
\ref{perfect-lemma-perfect-affine}). De même, on peut supposer que $E'$
est représenté par un complexe
$$
\mathcal{O}_S \to \mathcal{O}_S^{\oplus h'_1} \to
\mathcal{O}_S^{\oplus h'_2} \to \ldots
$$
où $h'_i = \dim_{\kappa(s)} H^i(X'_s, \mathcal{O}_{X'_s})$.
Par fonctorialité de la cohomologie, on dispose d'un morphisme
$$
E \longrightarrow E'
$$
dans $D(\mathcal{O}_S)$ dont la formation commute à tout changement de base.
Puisque le complexe représentant $E$ est un complexe fini de modules libres
de type fini et que $S$ est affine, on peut choisir un morphisme de complexes
$$
\xymatrix{
\mathcal{O}_S \ar[r]_d \ar[d]_a &
\mathcal{O}_S^{\oplus h_1} \ar[r] \ar[d]_b &
\mathcal{O}_S^{\oplus h_2} \ar[r] \ar[d]_c & \ldots \\
\mathcal{O}_S \ar[r]^{d'} &
\mathcal{O}_S^{\oplus h'_1} \ar[r] &
\mathcal{O}_S^{\oplus h'_2} \ar[r] & \ldots
}
$$
représentant le morphisme donné $E \to E'$. Puisque $s \in Z$, on voit que
la section trivialisante de $\mathcal{L}_s$ donne par image réciproque une section trivialisante
de $\mathcal{L}'_s = \mathcal{O}_{X'_s}$. Ainsi
$a \otimes \kappa(s)$ est un isomorphisme ; après avoir rétréci $S$,
on voit donc que $a$ est un isomorphisme. Enfin, on utilise l'hypothèse
que $H^1(X_s, \mathcal{O}) \to H^1(X'_s, \mathcal{O})$
est injectif pour voir qu'il existe un mineur $h_1 \times h_1$ de la
matrice définissant $b$ dont l'image est un élément non nul
de $\kappa(s)$. Après avoir rétréci $S$, on peut donc supposer
que $b$ est injectif. Or, puisque $\mathcal{L}' = \mathcal{O}_{X'}$,
on voit que $d' = 0$. Il s'ensuit que $d = 0$. On voit ainsi
que la section trivialisante de $\mathcal{L}_s$ se relève en une section
de $\mathcal{L}$ sur $X$. Un argument topologique immédiat (omis)
montre que cela signifie que $\mathcal{L}$ est trivial après avoir éventuellement
encore un peu rétréci $S$.
\end{proof}

\begin{lemma}
\label{lemma-H1-O-multiple-picard-flat-proper}
Considérons $n$ diagrammes commutatifs de schémas
$$
\xymatrix{
X_i \ar[rr] \ar[dr]_{f_i} & & X \ar[dl]^f \\
& S
}
$$
où $f_i : X_i \to S$ et $f : X \to S$ satisfont les hypothèses du
lemme \ref{lemma-diagonal-picard-flat-proper}.
Soit $\mathcal{L}$ un $\mathcal{O}_X$-module inversible,
et soit $\mathcal{L}_i$ son image réciproque sur $X_i$. Soit $Z \subset S$,
resp.\ $Z_i \subset S$, le sous-schéma localement fermé construit
dans le lemme \ref{lemma-diagonal-picard-flat-proper}
pour $(f, \mathcal{L})$, resp.\ $(f_i, \mathcal{L}_i)$,
de sorte que $Z \subset \bigcap_{i = 1, \ldots, n} Z_i$. Si $s \in Z$ et si
$$
H^1(X_s, \mathcal{O}) \longrightarrow
\bigoplus\nolimits_{i = 1, \ldots, n} H^1(X_{i, s}, \mathcal{O})
$$
est injectif, alors $Z \cap U = (\bigcap_{i = 1, \ldots, n} Z_i) \cap U$
(intersection schématique) pour un certain voisinage ouvert $U$ de $s$.
\end{lemma}

\begin{proof}
Ce lemme est une variante du lemme \ref{lemma-H1-O-picard-flat-proper},
et nous recommandons vivement au lecteur de lire d'abord cette démonstration ;
la présente démonstration en est essentiellement une copie, avec de légères modifications.
Il résulte de la description des points (à valeurs dans des schémas) de $Z$ et des $Z_i$
que $Z \subset \bigcap_{i = 1, \ldots, n} Z_i$, où l'intersection
est schématique. On peut donc remplacer $S$ par l'intersection
schématique $\bigcap_{i = 1, \ldots, n} Z_i$. Après avoir remplacé
$S$ par un voisinage ouvert affine de $s$, on peut supposer que
$\mathcal{L}_i = \mathcal{O}_{X_i}$ pour $i = 1, \ldots, n$.
Soient $E = Rf_*\mathcal{L}$ et
$E_i = Rf_{i, *}\mathcal{L}_i = Rf_{i, *}\mathcal{O}_{X_i}$.
Ce sont des complexes parfaits dont la formation commute à tout
changement de base (Catégories dérivées de schémas, lemme
\ref{perfect-lemma-flat-proper-perfect-direct-image-general}).
En particulier, on voit que
$$
E \otimes_{\mathcal{O}_S}^\mathbf{L} \kappa(s) =
R\Gamma(X_s, \mathcal{L}_s) = R\Gamma(X_s, \mathcal{O}_{X_s})
$$
La seconde égalité vient de ce que $s \in Z$. Posons
$h_j = \dim_{\kappa(s)} H^j(X_s, \mathcal{O}_{X_s})$.
Après avoir rétréci $S$, on peut représenter $E$ par un complexe
$$
\mathcal{O}_S \to \mathcal{O}_S^{\oplus h_1} \to
\mathcal{O}_S^{\oplus h_2} \to \ldots
$$
voir Plus sur l'algèbre, lemme
\ref{more-algebra-lemma-lift-perfect-from-residue-field}
(à strictement parler, cela utilise aussi
Catégories dérivées de schémas, lemmes
\ref{perfect-lemma-affine-compare-bounded} et
\ref{perfect-lemma-perfect-affine}). De même, on peut supposer que $E_i$
est représenté par un complexe
$$
\mathcal{O}_S \to \mathcal{O}_S^{\oplus h_{i, 1}} \to
\mathcal{O}_S^{\oplus h_{i, 2}} \to \ldots
$$
où $h_{i, j} = \dim_{\kappa(s)} H^j(X_{i, s}, \mathcal{O}_{X_{i, s}})$.
Par fonctorialité de la cohomologie, on dispose d'un morphisme
$$
E \longrightarrow E_i
$$
dans $D(\mathcal{O}_S)$ dont la formation commute à tout changement de base.
Puisque le complexe représentant $E$ est un complexe fini de modules libres
de type fini et que $S$ est affine, on peut choisir un morphisme de complexes
$$
\xymatrix{
\mathcal{O}_S \ar[r]_d \ar[d]_{a_i} &
\mathcal{O}_S^{\oplus h_1} \ar[r] \ar[d]_{b_i} &
\mathcal{O}_S^{\oplus h_2} \ar[r] \ar[d]_{c_i} & \ldots \\
\mathcal{O}_S \ar[r]^{d_i} &
\mathcal{O}_S^{\oplus h_{i, 1}} \ar[r] &
\mathcal{O}_S^{\oplus h_{i, 2}} \ar[r] & \ldots
}
$$
représentant le morphisme donné $E \to E_i$. Puisque $s \in Z$, on voit que
la section trivialisante de $\mathcal{L}_s$ donne par image réciproque une section trivialisante
de $\mathcal{L}_{i, s} = \mathcal{O}_{X_{i, s}}$. Ainsi
$a_i \otimes \kappa(s)$ est un isomorphisme ; après avoir rétréci $S$,
on voit donc que $a_i$ est un isomorphisme. Enfin, on utilise l'hypothèse
que $H^1(X_s, \mathcal{O}) \to
\bigoplus_{i = 1, \ldots, n} H^1(X_{i, s}, \mathcal{O})$
est injectif pour voir qu'il existe un mineur $h_1 \times h_1$ de la
matrice définissant $\oplus b_i$ dont l'image est un élément non nul
de $\kappa(s)$. Après avoir rétréci $S$, on peut donc supposer
que $(b_1, \ldots, b_n) : \mathcal{O}_S^{h_1}
\to \bigoplus_{i = 1, \ldots, n} \mathcal{O}_S^{h_{i, 1}}$
est injectif. Or, puisque $\mathcal{L}_i = \mathcal{O}_{X_i}$,
on voit que $d_i = 0$ pour $i = 1, \ldots n$. Il s'ensuit que $d = 0$,
car $(b_1, \ldots, b_n) \circ d = (\oplus d_i) \circ (a_1, \ldots, a_n)$.
On voit ainsi
que la section trivialisante de $\mathcal{L}_s$ se relève en une section
de $\mathcal{L}$ sur $X$. Un argument topologique immédiat (omis)
montre que cela signifie que $\mathcal{L}$ est trivial après avoir éventuellement
encore un peu rétréci $S$.
\end{proof}

\begin{lemma}
\label{lemma-pic-of-product}
Soient $f : X \to S$ et $g : Y \to S$ des morphismes de schémas
satisfaisant les hypothèses du lemme \ref{lemma-diagonal-picard-flat-proper}.
Soient $\sigma : S \to X$ et $\tau : S \to Y$ des sections de
$f$ et de $g$. Soit $s \in S$.
Soit $\mathcal{L}$ un faisceau inversible sur $X \times_S Y$.
Si $(1 \times \tau)^*\mathcal{L}$ sur $X$, $(\sigma \times 1)^*\mathcal{L}$
sur $Y$ et $\mathcal{L}|_{(X \times_S Y)_s}$ sont triviales, alors
il existe un voisinage ouvert $U$ de $s$ tel que
$\mathcal{L}$ soit trivial sur $(X \times_S Y)_U$.
\end{lemma}

\begin{proof}
Par la formule de K\"unneth (Variétés, lemme \ref{varieties-lemma-kunneth}),
le morphisme
$$
H^1(X_s \times_{\Spec(\kappa(s))} Y_s, \mathcal{O}) \to
H^1(X_s, \mathcal{O}) \oplus H^1(Y_s, \mathcal{O})
$$
est injectif. On peut donc
appliquer le lemme \ref{lemma-H1-O-multiple-picard-flat-proper}
aux deux morphismes
$$
1 \times \tau : X \to X \times_S Y
\quad\text{et}\quad
\sigma \times 1 : Y \to X \times_S Y
$$
pour conclure.
\end{proof}

\begin{theorem}[Théorème du cube]
\label{theorem-of-the-cube}
Soit $S$ un schéma. Soient $X$, $Y$ et $Z$ des schémas sur $S$.
Soient $x : S \to X$ et $y : S \to Y$ des sections des morphismes structuraux.
Soit $\mathcal{L}$ un module inversible sur $X \times_S Y \times_S Z$. Si
\begin{enumerate}
\item $X \to S$ et $Y \to S$ sont des morphismes plats, propres
et de présentation finie, à fibres géométriquement intègres,
\item les images réciproques de $\mathcal{L}$ par
$x \times \text{id}_Y \times \text{id}_Z$ et
$\text{id}_X \times y \times \text{id}_Z$
sont triviales sur $Y \times_S Z$ et $X \times_S Z$,
\item il existe un point $z \in Z$ tel que la restriction de $\mathcal{L}$
à $X \times_S Y \times_S z$ soit triviale, et
\item $Z$ est connexe,
\end{enumerate}
alors $\mathcal{L}$ est trivial.
\end{theorem}

\noindent
Un cas particulier fréquemment utilisé est le suivant.
Soit $k$ un corps. Soient $X, Y, Z$ des variétés munies respectivement de
points $k$-rationnels $x, y, z$. Soit $\mathcal{L}$ un module inversible
sur $X \times Y \times Z$. Si
\begin{enumerate}
\item $\mathcal{L}$ est trivial sur
$x \times Y \times Z$, $X \times y \times Z$ et $X \times Y \times z$, et
\item $X$ et $Y$ sont géométriquement intègres et propres sur $k$,
\end{enumerate}
alors $\mathcal{L}$ est trivial.

\begin{proof}
Remarquons que le morphisme $X \times_S Y \to S$ est plat, propre
et de présentation finie, à fibres géométriquement intègres
(voir Variétés, lemmes \ref{varieties-lemma-geometrically-integral},
\ref{varieties-lemma-bijection-irreducible-components}
et \ref{varieties-lemma-geometrically-reduced-any-base-change} pour
l'assertion concernant les fibres). Par Catégories dérivées de schémas, lemme
\ref{perfect-lemma-proper-flat-geom-red-connected},
on voit que l'image directe du faisceau structural par $X \to S$, $Y \to S$ ou
$X \times_S Y \to S$ est le faisceau structural de $S$, et que cela reste vrai
après tout changement de base. On peut donc appliquer
le lemme \ref{lemma-diagonal-picard-flat-proper} au morphisme
$$
p : X \times_S Y \times_S Z \longrightarrow Z
$$
et au module inversible $\mathcal{L}$ pour obtenir un sous-schéma localement fermé
« universel » $Z' \subset Z$ tel que $\mathcal{L}|_{X \times_S Y \times_S Z'}$
soit l'image réciproque d'un module inversible $\mathcal{N}$ sur $Z'$.
L'existence de $z$ montre que $Z'$ est non vide. D'après le
lemme \ref{lemma-get-a-closed},
on voit que $Z' \subset Z$ est un sous-schéma fermé.
Soit $z' \in Z'$ un point.
Remarquons que l'on peut écrire $p$ comme le morphisme produit
$$
(X \times_S Z) \times_Z (Y \times_S Z) \longrightarrow Z
$$
On peut donc appliquer le lemme \ref{lemma-pic-of-product}
au morphisme $p$, au point $z'$ et aux sections
$\sigma : Z \to X \times_S Z$ et $\tau : Z \to Y \times_S Z$
données par $x$ et $y$. On en déduit que $Z'$ est ouvert.
Ainsi $Z' = Z$ et $\mathcal{L} = p^*\mathcal{N}$
pour un certain module inversible $\mathcal{N}$ sur $Z$.
En prenant l'image réciproque par
$x \times y \times \text{id}_Z : Z  \to X \times_S Y \times_S Z$,
on obtient d'une part $\mathcal{N}$ et, d'autre part,
le module inversible trivial d'après l'hypothèse (2).
Ainsi $\mathcal{N} = \mathcal{O}_Z$, ce qui achève la démonstration.
\end{proof}









\section{Arguments de passage à la limite}
\label{section-limits}

\noindent
Quelques lemmes faisant intervenir des limites de schémas, ainsi que l'approximation
noethérienne. Nous nous en tenons principalement au cas affine. Certains de ces lemmes sont des
cas particuliers de lemmes du chapitre consacré aux limites.

\begin{lemma}
\label{lemma-Noetherian-approximation}
Soit $f : X \to S$ un morphisme de schémas affines, de
présentation finie. Alors il existe un diagramme cartésien
$$
\xymatrix{
X_0 \ar[d]_{f_0} & X \ar[l]^g \ar[d]^f \\
S_0 & S \ar[l]
}
$$
tel que
\begin{enumerate}
\item $X_0$, $S_0$ soient des schémas affines,
\item $S_0$ soit de type fini sur $\mathbf{Z}$,
\item $f_0$ soit de type fini.
\end{enumerate}
\end{lemma}

\begin{proof}
Écrivons $S = \Spec(A)$ et $X = \Spec(B)$.
Comme $f$ est de présentation finie, on voit que
$B$ est une $A$-algèbre de présentation finie, voir
Morphismes,
Lemme \ref{morphisms-lemma-locally-finite-presentation-characterize}.
Le lemme résulte donc de
Algèbre, Lemme \ref{algebra-lemma-limit-module-finite-presentation}.
\end{proof}

\begin{lemma}
\label{lemma-Noetherian-approximation-module}
Soit $f : X \to S$ un morphisme de schémas affines, de
présentation finie. Soit $\mathcal{F}$ un $\mathcal{O}_X$-module quasi-cohérent
de présentation finie. Alors il existe un diagramme comme dans le
Lemme \ref{lemma-Noetherian-approximation}
tel qu'il existe un $\mathcal{O}_{X_0}$-module cohérent $\mathcal{F}_0$
vérifiant $g^*\mathcal{F}_0 = \mathcal{F}$.
\end{lemma}

\begin{proof}
Écrivons $S = \Spec(A)$, $X = \Spec(B)$ et
$\mathcal{F} = \widetilde{M}$. Comme $f$ est de présentation finie, on voit que
$B$ est une $A$-algèbre de présentation finie, voir
Morphismes,
Lemme \ref{morphisms-lemma-locally-finite-presentation-characterize}.
Comme $\mathcal{F}$ est de présentation finie sur $\mathcal{O}_X$, on voit que
$M$ est un $B$-module de présentation finie, voir
Propriétés, Lemme \ref{properties-lemma-finite-presentation-module}.
Le lemme résulte donc de
Algèbre, Lemme \ref{algebra-lemma-limit-module-finite-presentation}.
\end{proof}

\begin{lemma}
\label{lemma-Noetherian-approximation-flat-module}
Soit $f : X \to S$ un morphisme de schémas affines, de
présentation finie. Soit $\mathcal{F}$ un $\mathcal{O}_X$-module quasi-cohérent
de présentation finie et plat sur $S$. On peut alors choisir un diagramme comme dans le
Lemme \ref{lemma-Noetherian-approximation-module}
et un faisceau $\mathcal{F}_0$ tels qu'en outre $\mathcal{F}_0$
soit plat sur $S_0$.
\end{lemma}

\begin{proof}
Écrivons $S = \Spec(A)$, $X = \Spec(B)$ et
$\mathcal{F} = \widetilde{M}$. Comme $f$ est de présentation finie, on voit que
$B$ est une $A$-algèbre de présentation finie, voir
Morphismes,
Lemme \ref{morphisms-lemma-locally-finite-presentation-characterize}.
Comme $\mathcal{F}$ est de présentation finie sur $\mathcal{O}_X$, on voit que
$M$ est un $B$-module de présentation finie, voir
Propriétés, Lemme \ref{properties-lemma-finite-presentation-module}.
Comme $\mathcal{F}$ est plat sur $S$, on voit que $M$ est plat sur $A$, voir
Morphismes, Lemme \ref{morphisms-lemma-flat-module-characterize}.
Le lemme résulte donc de
Algèbre, Lemme \ref{algebra-lemma-flat-finite-presentation-limit-flat}.
\end{proof}

\begin{lemma}
\label{lemma-Noetherian-approximation-flat}
Soit $f : X \to S$ un morphisme de schémas affines, de
présentation finie et plat. Alors il existe un diagramme comme dans le
Lemme \ref{lemma-Noetherian-approximation}
tel qu'en outre $f_0$ soit plat.
\end{lemma}

\begin{proof}
C'est un cas particulier du
Lemme \ref{lemma-Noetherian-approximation-flat-module}.
\end{proof}

\begin{lemma}
\label{lemma-Noetherian-approximation-smooth}
Soit $f : X \to S$ un morphisme lisse de schémas affines.
Alors il existe un diagramme comme dans le
Lemme \ref{lemma-Noetherian-approximation}
tel qu'en outre $f_0$ soit lisse.
\end{lemma}

\begin{proof}
Écrivons $S = \Spec(A)$, $X = \Spec(B)$ ;
puisque $f$ est lisse, on voit que $B$ est une $A$-algèbre lisse, voir
Morphismes,
Lemme \ref{morphisms-lemma-smooth-characterize}.
Le lemme résulte donc de
Algèbre, Lemme \ref{algebra-lemma-finite-presentation-fs-Noetherian}.
\end{proof}

\begin{lemma}
\label{lemma-Noetherian-approximation-geometrically-reduced}
Soit $f : X \to S$ un morphisme de schémas affines, de
présentation finie et dont les fibres sont géométriquement réduites.
Alors il existe un diagramme comme dans le
Lemme \ref{lemma-Noetherian-approximation}
tel qu'en outre les fibres de $f_0$ soient géométriquement réduites.
\end{lemma}

\begin{proof}
Appliquons le
Lemme \ref{lemma-Noetherian-approximation}
pour obtenir un diagramme cartésien
$$
\xymatrix{
X_0 \ar[d]_{f_0} & X \ar[l]^g \ar[d]^f \\
S_0 & S \ar[l]_h
}
$$
de schémas affines, où $X_0 \to S_0$ est un morphisme de type fini entre
schémas de type fini sur $\mathbf{Z}$. D'après le
Lemme \ref{lemma-geometrically-reduced-constructible},
l'ensemble $E \subset S_0$ des points en lesquels la fibre de
$f_0$ est géométriquement réduite est une partie constructible. D'après le
Lemme \ref{lemma-base-change-fibres-geometrically-reduced},
on a $h(S) \subset E$. Écrivons $S_0 = \Spec(A_0)$ et
$S = \Spec(A)$. Écrivons $A = \colim_i A_i$ comme une
colimite filtrante de $A_0$-algèbres de type fini. D'après
Limites, Lemme \ref{limits-lemma-limit-contained-in-constructible},
on voit que l'image de $\Spec(A_i) \to S_0$ est contenue dans $E$
pour un certain $i$. En remplaçant $S_0$ par $\Spec(A_i)$ et
$X_0$ par $X_0 \times_{S_0} \Spec(A_i)$, on voit que
toutes les fibres de $f_0$ sont géométriquement réduites.
\end{proof}

\begin{lemma}
\label{lemma-Noetherian-approximation-geometrically-irreducible}
Soit $f : X \to S$ un morphisme de schémas affines, de
présentation finie et dont les fibres sont géométriquement irréductibles.
Alors il existe un diagramme comme dans le
Lemme \ref{lemma-Noetherian-approximation}
tel qu'en outre les fibres de $f_0$ soient géométriquement irréductibles.
\end{lemma}

\begin{proof}
Appliquons le
Lemme \ref{lemma-Noetherian-approximation}
pour obtenir un diagramme cartésien
$$
\xymatrix{
X_0 \ar[d]_{f_0} & X \ar[l]^g \ar[d]^f \\
S_0 & S \ar[l]_h
}
$$
de schémas affines, où $X_0 \to S_0$ est un morphisme de type fini entre
schémas de type fini sur $\mathbf{Z}$. D'après le
Lemme \ref{lemma-nr-geom-irreducible-components-constructible},
l'ensemble $E \subset S_0$ des points en lesquels la fibre de
$f_0$ est géométriquement irréductible est une partie constructible. D'après le
Lemme \ref{lemma-base-change-fibres-geometrically-irreducible},
on a $h(S) \subset E$. Écrivons $S_0 = \Spec(A_0)$ et
$S = \Spec(A)$. Écrivons $A = \colim_i A_i$ comme une
colimite filtrante de $A_0$-algèbres de type fini. D'après
Limites, Lemme \ref{limits-lemma-limit-contained-in-constructible},
on voit que l'image de $\Spec(A_i) \to S_0$ est contenue dans $E$
pour un certain $i$. En remplaçant $S_0$ par $\Spec(A_i)$ et
$X_0$ par $X_0 \times_{S_0} \Spec(A_i)$, on voit que
toutes les fibres de $f_0$ sont géométriquement irréductibles.
\end{proof}

\begin{lemma}
\label{lemma-Noetherian-approximation-geometrically-connected}
Soit $f : X \to S$ un morphisme de schémas affines, de
présentation finie et dont les fibres sont géométriquement connexes.
Alors il existe un diagramme comme dans le
Lemme \ref{lemma-Noetherian-approximation}
tel qu'en outre les fibres de $f_0$ soient géométriquement connexes.
\end{lemma}

\begin{proof}
Appliquons le
Lemme \ref{lemma-Noetherian-approximation}
pour obtenir un diagramme cartésien
$$
\xymatrix{
X_0 \ar[d]_{f_0} & X \ar[l]^g \ar[d]^f \\
S_0 & S \ar[l]_h
}
$$
de schémas affines, où $X_0 \to S_0$ est un morphisme de type fini entre
schémas de type fini sur $\mathbf{Z}$. D'après le
Lemme \ref{lemma-nr-geom-connected-components-constructible},
l'ensemble $E \subset S_0$ des points en lesquels la fibre de
$f_0$ est géométriquement connexe est une partie constructible. D'après le
Lemme \ref{lemma-base-change-fibres-geometrically-connected},
on a $h(S) \subset E$. Écrivons $S_0 = \Spec(A_0)$ et
$S = \Spec(A)$. Écrivons $A = \colim_i A_i$ comme une
colimite filtrante de $A_0$-algèbres de type fini. D'après
Limites, Lemme \ref{limits-lemma-limit-contained-in-constructible},
on voit que l'image de $\Spec(A_i) \to S_0$ est contenue dans $E$
pour un certain $i$. En remplaçant $S_0$ par $\Spec(A_i)$ et
$X_0$ par $X_0 \times_{S_0} \Spec(A_i)$, on voit que
toutes les fibres de $f_0$ sont géométriquement connexes.
\end{proof}

\begin{lemma}
\label{lemma-Noetherian-approximation-dimension-d}
Soit $d \geq 0$ un entier.
Soit $f : X \to S$ un morphisme de schémas affines, de
présentation finie, dont toutes les fibres sont de dimension $d$.
Alors il existe un diagramme comme dans le
Lemme \ref{lemma-Noetherian-approximation}
tel qu'en outre toutes les fibres de $f_0$ soient de dimension $d$.
\end{lemma}

\begin{proof}
Appliquons le
Lemme \ref{lemma-Noetherian-approximation}
pour obtenir un diagramme cartésien
$$
\xymatrix{
X_0 \ar[d]_{f_0} & X \ar[l]^g \ar[d]^f \\
S_0 & S \ar[l]_h
}
$$
de schémas affines, où $X_0 \to S_0$ est un morphisme de type fini entre
schémas de type fini sur $\mathbf{Z}$. D'après le
Lemme \ref{lemma-dimension-fibres-constructible},
l'ensemble $E \subset S_0$ des points en lesquels la fibre de
$f_0$ est de dimension $d$ est une partie constructible. D'après le
Lemme \ref{lemma-base-change-dimension-fibres},
on a $h(S) \subset E$. Écrivons $S_0 = \Spec(A_0)$ et
$S = \Spec(A)$. Écrivons $A = \colim_i A_i$ comme une
colimite filtrante de $A_0$-algèbres de type fini. D'après
Limites, Lemme \ref{limits-lemma-limit-contained-in-constructible},
on voit que l'image de $\Spec(A_i) \to S_0$ est contenue dans $E$
pour un certain $i$. En remplaçant $S_0$ par $\Spec(A_i)$ et
$X_0$ par $X_0 \times_{S_0} \Spec(A_i)$, on voit que
toutes les fibres de $f_0$ sont de dimension $d$.
\end{proof}

\begin{lemma}
\label{lemma-Noetherian-approximation-standard-syntomic}
Soit $f : X \to S$ un morphisme de schémas affines qui est
syntomique standard (voir
Morphismes, Définition \ref{morphisms-definition-syntomic}).
Alors il existe un diagramme comme dans le
Lemme \ref{lemma-Noetherian-approximation}
tel qu'en outre $f_0$ soit syntomique standard.
\end{lemma}

\begin{proof}
Ce lemme est une copie de
Algèbre,
Lemme \ref{algebra-lemma-relative-global-complete-intersection-Noetherian}.
\end{proof}

\begin{lemma}
\label{lemma-Noetherian-approximation-combine}
(Approximation noethérienne et combinaison de propriétés.)
Soient $P$, $Q$ des propriétés de morphismes de schémas qui sont stables
par changement de base. Soit $f : X \to S$ un morphisme de présentation finie
entre schémas affines. Supposons que l'on puisse trouver des diagrammes cartésiens
$$
\vcenter{
\xymatrix{
X_1 \ar[d]_{f_1} & X \ar[l] \ar[d]^f \\
S_1 & S \ar[l]
}
}
\quad\text{et}\quad
\vcenter{
\xymatrix{
X_2 \ar[d]_{f_2} & X \ar[l] \ar[d]^f \\
S_2 & S \ar[l]
}
}
$$
de schémas affines, où $S_1$, $S_2$ sont de type fini sur $\mathbf{Z}$
et $f_1$, $f_2$ sont de type fini, tels que $f_1$ ait la propriété $P$
et $f_2$ ait la propriété $Q$. Alors on peut trouver un diagramme cartésien
$$
\xymatrix{
X_0 \ar[d]_{f_0} & X \ar[l] \ar[d]^f \\
S_0 & S \ar[l]
}
$$
de schémas affines, où $S_0$ est de type fini sur $\mathbf{Z}$
et $f_0$ est de type fini, tel que $f_0$ ait à la fois la propriété $P$ et
la propriété $Q$.
\end{lemma}

\begin{proof}
Les deux diagrammes donnés correspondent à des diagrammes cocartésiens d'anneaux
$$
\vcenter{
\xymatrix{
B_1 \ar[r] & B \\
A_1 \ar[u] \ar[r] & A \ar[u]
}
}
\quad\text{et}\quad
\vcenter{
\xymatrix{
B_2 \ar[r] & B \\
A_2 \ar[u] \ar[r] & A \ar[u]
}
}
$$
Soit $A_0 \subset A$ une sous-$\mathbf{Z}$-algèbre de type fini de $A$
contenant les images de $A_1 \to A$ et de $A_2 \to A$. Une telle sous-algèbre
existe car, par hypothèse, $A_1$ et $A_2$ sont toutes deux de type fini sur
$\mathbf{Z}$. Remarquons que les anneaux $B_{0, 1} = B_1 \otimes_{A_1} A_0$
et $B_{0, 2} = B_2 \otimes_{A_2} A_0$ sont des $A_0$-algèbres de type fini
qui vérifient
$B_{0, 1} \otimes_{A_0} A \cong B \cong B_{0, 2} \otimes_{A_0} A$
en tant que $A$-algèbres. Comme $A$ est la colimite filtrante de ses
sous-$A_0$-algèbres de type fini, d'après
Limites, Lemme \ref{limits-lemma-descend-finite-presentation},
on peut supposer, quitte à agrandir $A_0$, qu'il existe un isomorphisme
$B_{0, 1} \cong B_{0, 2}$ de $A_0$-algèbres. Puisque les propriétés $P$ et $Q$
sont supposées stables par changement de base, on conclut qu'en posant
$S_0 = \Spec(A_0)$ et
$$
X_0 = X_1 \times_{S_1} S_0 =
\Spec(B_{0, 1}) \cong \Spec(B_{0, 2}) = X_2 \times_{S_2} S_0
$$
on obtient ce que l'on veut.
\end{proof}












\section{Voisinages \'etales}
\label{section-etale-neighbourhoods}

\noindent
Il s'avère que certaines propriétés des morphismes sont plus faciles à étudier
après un changement de base \'etale. Il est commode d'introduire la
terminologie suivante.

\begin{definition}
\label{definition-etale-neighbourhood}
Soit $S$ un schéma. Soit $s \in S$ un point.
\begin{enumerate}
\item Un {\it voisinage \'etale de $(S, s)$} est un
couple $(U, u)$ muni d'un morphisme \'etale
de schémas $\varphi : U \to S$ tel que $\varphi(u) = s$.
\item Un {\it morphisme de voisinages \'etales} $f : (V, v) \to (U, u)$
de $(S, s)$ est simplement un morphisme de $S$-schémas $f : V \to U$ tel
que $f(v) = u$.
\item Un {\it voisinage \'etale élémentaire} est un voisinage \'etale
$\varphi : (U, u) \to (S, s)$ tel que $\kappa(s) = \kappa(u)$.
\end{enumerate}
\end{definition}

\noindent
La notion de voisinage \'etale élémentaire porte de nombreux noms
dans la littérature ; ces objets sont par exemple parfois appelés
``voisinages \'etales'' (\cite[Page 36]{Milne} ou
``fortement \'etales'' (\cite[Page 108]{KPR}).
Nous suivons ici la convention de l'article
\cite{GruRay} en les appelant voisinages \'etales élémentaires.

\medskip\noindent
Si $f : (V, v) \to (U, u)$ est un morphisme de voisinages \'etales,
alors $f$ est automatiquement \'etale, voir
Morphismes, Lemme \ref{morphisms-lemma-etale-permanence}.
Ainsi, $(V, v)$ devient un voisinage \'etale de
$(U, u)$. Bien entendu, puisque la composée de morphismes \'etales
est \'etale (Morphismes, Lemme \ref{morphisms-lemma-composition-etale}),
on voit réciproquement que tout voisinage \'etale $(V, v)$ de
$(U, u)$ est aussi un voisinage \'etale de $(S, s)$.
Remarquons également que, si $U \subset S$ est un voisinage ouvert
de $s$, alors $(U, s) \to (S, s)$ est un voisinage \'etale.
Cela résulte du fait qu'une immersion ouverte est
\'etale (Morphismes, Lemme \ref{morphisms-lemma-open-immersion-etale}).
Nous utiliserons ces remarques sans autre mention dans toute cette
section.

\medskip\noindent
Remarquons que $\kappa(u)/\kappa(s)$ est une extension finie séparable
si $(U, u) \to (S, s)$ est un voisinage \'etale,
voir Morphismes, Lemme \ref{morphisms-lemma-etale-at-point}.

\begin{lemma}
\label{lemma-realize-prescribed-residue-field-extension-etale}
Soit $S$ un schéma.
Soit $s \in S$.
Soit $k/\kappa(s)$ une extension finie séparable.
Alors il existe un voisinage \'etale $(U, u) \to (S, s)$
tel que l'extension de corps $\kappa(u)/\kappa(s)$ soit
isomorphe à $k/\kappa(s)$.
\end{lemma}

\begin{proof}
On peut supposer $S$ affine.
Dans ce cas, le lemme résulte de
Algèbre, Lemme \ref{algebra-lemma-make-etale-map-prescribed-residue-field}.
\end{proof}

\begin{lemma}
\label{lemma-etale-neighbourhoods-not-quite-filtered}
Soient $S$ un schéma et $s$ un point de $S$.
La catégorie des voisinages \'etales possède les propriétés suivantes :
\begin{enumerate}
\item Soient $(U_i, u_i)_{i=1, 2}$ deux voisinages \'etales de
$s$ dans $S$. Alors il existe un troisième voisinage \'etale
$(U, u)$ et des morphismes
$(U, u) \to (U_i, u_i)$, $i = 1, 2$.
\item Soient $h_1, h_2: (U, u) \to (U', u')$ deux
morphismes entre voisinages \'etales de $s$.
Supposons que $h_1$, $h_2$ induisent la même application $\kappa(u') \to \kappa(u)$ entre corps
résiduels. Alors il existe un voisinage \'etale $(U'', u'')$ et un morphisme
$h : (U'', u'') \to (U, u)$
qui égalise $h_1$ et $h_2$, c'est-à-dire tel que
$h_1 \circ h = h_2 \circ h$.
\end{enumerate}
\end{lemma}

\begin{proof}
Pour (1), considérons le produit fibré $U = U_1 \times_S U_2$.
Il est \'etale sur $U_1$ et sur $U_2$, car les morphismes \'etales sont
préservés par changement de base, voir
Morphismes, Lemme \ref{morphisms-lemma-base-change-etale}.
Il existe un point de $U$ qui s'envoie sur $u_1$ et sur $u_2$, par exemple
d'après la description des points d'un produit fibré dans
Schémas, Lemme \ref{schemes-lemma-points-fibre-product}.
Pour (2), définissons $U''$ comme le produit fibré
$$
\xymatrix{
U'' \ar[r] \ar[d] & U \ar[d]^{(h_1, h_2)} \\
U' \ar[r]^-\Delta & U' \times_S U'.
}
$$
Puisque $h_1$ et $h_2$ induisent la même application de corps résiduels
$\kappa(u') \to \kappa(u)$, il existe un point $u'' \in U''$
au-dessus de $u'$ tel que $\kappa(u'') = \kappa(u')$.
En particulier, $U'' \not = \emptyset$.
De plus, puisque $U'$ est \'etale sur $S$, le produit fibré
$U'\times_S U'$ l'est aussi (voir
Morphismes, Lemmes \ref{morphisms-lemma-base-change-etale} et
\ref{morphisms-lemma-composition-etale}).
La flèche verticale $(h_1, h_2)$ est donc \'etale d'après
Morphismes, Lemme \ref{morphisms-lemma-etale-permanence}.
Par conséquent, $U''$ est \'etale sur $U'$ par changement de base, et donc aussi
\'etale sur $S$ (car les composées de morphismes \'etales sont \'etales).
Ainsi, $(U'', u'')$ est une solution au problème.
\end{proof}

\begin{lemma}
\label{lemma-elementary-etale-neighbourhoods}
Soient $S$ un schéma et $s$ un point de $S$.
La catégorie des voisinages \'etales élémentaires de $(S, s)$
est cofiltrante (voir
Catégories, Définition \ref{categories-definition-codirected}).
\end{lemma}

\begin{proof}
Cela résulte immédiatement des définitions et du
Lemme \ref{lemma-etale-neighbourhoods-not-quite-filtered}.
\end{proof}

\begin{lemma}
\label{lemma-describe-henselization}
Soit $S$ un schéma. Soit $s \in S$. Alors on a
$$
\mathcal{O}_{S, s}^h =
\colim_{(U, u)} \mathcal{O}(U)
$$
où la colimite est prise sur la catégorie filtrante qui est l'opposée de la
catégorie des voisinages \'etales élémentaires $(U, u)$ de $(S, s)$.
\end{lemma}

\begin{proof}
Soit $\Spec(A) \subset S$ un voisinage affine de $s$.
Soit $\mathfrak p \subset A$ l'idéal premier correspondant à $s$.
Avec ces choix, on dispose d'isomorphismes canoniques
$\mathcal{O}_{S, s} = A_{\mathfrak p}$ et $\kappa(s) = \kappa(\mathfrak p)$.
Un système cofinal de voisinages \'etales élémentaires est donné par les
voisinages \'etales élémentaires $(U, u)$ tels que $U$ soit affine et que
$U \to S$ se factorise par $\Spec(A)$. Autrement dit, on voit que
le membre de droite est égal à $\colim_{(B, \mathfrak q)} B$,
où la colimite porte sur les $A$-algèbres \'etales $B$ munies d'un idéal premier
$\mathfrak q$ au-dessus de $\mathfrak p$ et telles que
$\kappa(\mathfrak p) = \kappa(\mathfrak q)$.
Le lemme résulte donc de
Algèbre, Lemme \ref{algebra-lemma-henselization-different}.
\end{proof}

\noindent
On peut relever à l'espace total les voisinages \'etales de points
des fibres.

\begin{lemma}
\label{lemma-lift-etale-neighbourhood-fibre}
\begin{slogan}
Relever à l'espace total un voisinage \'etale d'un point de la fibre.
\end{slogan}
Soit $X \to S$ un morphisme de schémas. Soit $x \in X$ d'image $s \in S$.
Soit $(V, v) \to (X_s, x)$ un voisinage \'etale.
Alors il existe un voisinage \'etale $(U, u) \to (X, x)$
tel qu'il existe un morphisme $(U_s, u) \to (V, v)$
de voisinages \'etales de $(X_s, x)$ qui soit une immersion ouverte.
\end{lemma}

\begin{proof}
On peut supposer $X$, $V$ et $S$ affines. Supposons que le morphisme
$X \to S$ soit donné par $A \to B$, le point $x$ par un idéal premier
$\mathfrak q \subset B$, le point $s$ par $\mathfrak p = A \cap \mathfrak q$,
et le morphisme $V \to X_s$ par $B \otimes_A \kappa(\mathfrak p) \to C$.
Puisque $\kappa(\mathfrak p)$ est un localisé de
$A/\mathfrak p$, il existe $f \in A$, $f \not \in \mathfrak p$
et un morphisme d'anneaux \'etale $B \otimes_A (A/\mathfrak p)_f \to D$
tel que
$$
C = (B \otimes_A \kappa(\mathfrak p))
\otimes_{B \otimes_A (A/\mathfrak p)_f} D
$$
Voir Algèbre, Lemme \ref{algebra-lemma-etale}, partie (9).
Après avoir remplacé $A$ par $A_f$ et $B$ par $B_f$, on peut supposer que
$D$ est \'etale sur $B \otimes_A A/\mathfrak p = B/\mathfrak p B$.
On peut alors appliquer Algèbre, Lemme \ref{algebra-lemma-lift-etale}.
Cela démontre le lemme.
\end{proof}














\section{Voisinages \'etales et branches}
\label{section-etale-nbhds-branches}

\noindent
Le nombre de branches (géométriques) d'un schéma en un point a été
défini dans Propriétés, Section \ref{properties-section-unibranch}.
Dans Variétés, Section \ref{varieties-section-number-of-branches},
nous l'avons relié aux fibres du morphisme de normalisation.
Dans cette section, nous donnons une caractérisation de ce nombre au moyen de
voisinages \'etales.

\begin{lemma}
\label{lemma-nr-minimal-primes}
Soit $R = \colim R_i$ la colimite d'un système inductif d'anneaux
dont les morphismes de transition sont fidèlement plats.
Alors le nombre d'idéaux premiers minimaux de $R$,
considéré comme un élément de $\{0, 1, 2, \ldots, \infty\}$,
est la borne supérieure des nombres d'idéaux premiers minimaux des $R_i$.
\end{lemma}

\begin{proof}
Si $A \to B$ est un morphisme plat d'anneaux, alors $\Spec(B) \to \Spec(A)$
envoie les idéaux premiers minimaux sur des idéaux premiers minimaux en vertu du théorème de descente
(Algèbre, lemme \ref{algebra-lemma-flat-going-down}).
Si $A \to B$ est fidèlement plat, alors tout idéal premier
minimal est l'image d'un idéal premier minimal (d'après
Algèbre, lemmes \ref{algebra-lemma-ff-rings} et
\ref{algebra-lemma-minimal-prime-image-minimal-prime}).
Par conséquent, le nombre d'idéaux premiers minimaux de $R_i$ est
$\geq$ au nombre d'idéaux premiers minimaux de $R_{i'}$ si $i \leq i'$.
D'après Algèbre, lemme \ref{algebra-lemma-colimit-faithfully-flat},
chacun des morphismes $R_i \to R$ est
fidèlement plat, et l'on voit aussi que
le nombre d'idéaux premiers minimaux de $R$ est
$\geq$ au nombre d'idéaux premiers minimaux de $R_i$.
Enfin, supposons que $\mathfrak q_1, \ldots, \mathfrak q_n$
soient des idéaux premiers minimaux de $R$ deux à deux distincts. On peut alors
trouver un $i$ tel que $R_i \cap \mathfrak q_1, \ldots, R_i \cap \mathfrak q_n$
soient deux à deux distincts (comme ensembles, et donc comme idéaux premiers).
Ceci démontre le lemme.
\end{proof}

\begin{lemma}
\label{lemma-nr-branches}
Soient $X$ un schéma et $x \in X$ un point. Alors
\begin{enumerate}
\item le nombre de branches de $X$ en $x$ est égal à
la borne supérieure du nombre de composantes irréductibles de $U$
passant par $u$, pour les voisinages \'etales élémentaires
$(U, u) \to (X, x)$,
\item le nombre de branches géométriques de $X$ en $x$ est égal à
la borne supérieure du nombre de composantes irréductibles de $U$
passant par $u$, pour les voisinages \'etales
$(U, u) \to (X, x)$,
\item $X$ est unibranche en $x$ si et seulement si, pour tout
voisinage \'etale élémentaire $(U, u) \to (X, x)$, il existe
exactement une composante irréductible de $U$ passant par $u$, et
\item $X$ est géométriquement unibranche en $x$ si et seulement si, pour tout
voisinage \'etale $(U, u) \to (X, x)$, il existe
exactement une composante irréductible de $U$ passant par $u$.
\end{enumerate}
\end{lemma}

\begin{proof}
Les assertions (3) et (4) se déduisent des assertions (1) et (2) au moyen de
Propriétés, lemme \ref{properties-lemma-number-of-branches-1}.

\medskip\noindent
Démonstration de (1). Soit $\Spec(A)$ un voisinage ouvert affine
de $x$, et soit $\mathfrak p \subset A$ l'idéal premier
correspondant à $x$. On peut remplacer $X$ par $\Spec(A)$, et
il suffit, dans la borne supérieure, de considérer les voisinages \'etales élémentaires affines
$(U, u)$, puisqu'ils forment un sous-système cofinal.
Rappelons que l'hensélisé $A_\mathfrak p^h$
est la colimite des anneaux $B_\mathfrak q$ sur la catégorie
des couples $(B, \mathfrak q)$, où $B$ est une $A$-algèbre \'etale
et $\mathfrak q$ un idéal premier au-dessus de $\mathfrak p$ tel que
$\kappa(\mathfrak q) = \kappa(\mathfrak p)$ ; voir
Algèbre, lemme \ref{algebra-lemma-henselization-different}.
Ces couples $(B, \mathfrak q)$ correspondent exactement aux
voisinages \'etales élémentaires affines $(U, u)$
par la correspondance entre anneaux et schémas affines.
Observons que les composantes irréductibles de $\Spec(B)$
passant par $\mathfrak q$ correspondent exactement aux idéaux
premiers minimaux de $B_\mathfrak q$. Le nombre d'idéaux premiers
minimaux de $A_\mathfrak p^h$ est le nombre de branches
de $X$ en $x$, d'après Propriétés, définition
\ref{properties-definition-number-of-branches}.
Observons que les morphismes de transition
$B_\mathfrak q \to B'_{\mathfrak q'}$ du système
sont tous plats. Puisqu'un morphisme local plat d'anneaux est fidèlement plat
(Algèbre, lemme \ref{algebra-lemma-local-flat-ff}),
on voit que l'assertion découle du
lemme \ref{lemma-nr-minimal-primes}.

\medskip\noindent
Démonstration de (2). La démonstration est la même que celle de (1), si ce n'est que l'on utilise
Algèbre, lemme \ref{algebra-lemma-strict-henselization-different}.
Il existe une petite différence : étant donnée une clôture algébrique séparable
$\kappa^{sep}$ de $\kappa(x)$, pour tout voisinage \'etale
$(U, u)$ on peut choisir un plongement de $\kappa(x)$-algèbres
$\phi : \kappa(u) \to \kappa^{sep}$,
car l'extension $\kappa(u)/\kappa(x)$ est finie et séparable
(Morphismes, lemme \ref{morphisms-lemma-etale-at-point}).
On peut donc prendre la borne supérieure sur tous les triplets
$(U, u, \phi)$, où $(U, u) \to (X, x)$ est un voisinage
\'etale affine et $\phi : \kappa(u) \to \kappa^{sep}$
est un plongement de $\kappa(x)$-algèbres. Ces triplets correspondent
exactement à ceux de
Algèbre, lemme \ref{algebra-lemma-strict-henselization-different},
et le reste de la démonstration est exactement le même.
\end{proof}

\noindent
Nous aurons besoin d'une variante relative du lemme précédent.

\begin{lemma}
\label{lemma-nr-branches-fibre}
Soit $X \to S$ un morphisme de schémas et soit $x \in X$ un point d'image $s$.
Alors
\begin{enumerate}
\item le nombre de branches de la fibre $X_s$ en $x$ est égal à
la borne supérieure du nombre de composantes irréductibles de la fibre $U_s$
passant par $u$, pour les voisinages \'etales élémentaires
$(U, u) \to (X, x)$,
\item le nombre de branches géométriques de la fibre $X_s$ en $x$ est égal à
la borne supérieure du nombre de composantes irréductibles de la fibre $U_s$
passant par $u$, pour les voisinages \'etales
$(U, u) \to (X, x)$,
\item la fibre $X_s$ est unibranche en $x$ si et seulement si, pour tout
voisinage \'etale élémentaire $(U, u) \to (X, x)$, il existe
exactement une composante irréductible de la fibre $U_s$ passant par $u$, et
\item $X$ est géométriquement unibranche en $x$ si et seulement si, pour tout
voisinage \'etale $(U, u) \to (X, x)$, il existe
exactement une composante irréductible de $U_s$ passant par $u$.
\end{enumerate}
\end{lemma}

\begin{proof}
Combiner les lemmes \ref{lemma-nr-branches} et
\ref{lemma-lift-etale-neighbourhood-fibre}.
\end{proof}

\begin{lemma}
\label{lemma-number-of-branches-and-smooth}
Soit $X \to S$ un morphisme lisse de schémas.
Soit $x \in X$ d'image $s \in S$.
Alors
\begin{enumerate}
\item Le nombre de branches géométriques de $X$ en $x$
est égal au nombre de branches géométriques de $S$ en $s$.
\item Si $\kappa(x)/\kappa(s)$ est une extension de corps purement inséparable\footnote{
En fait, il suffirait que $\kappa(x)$ soit géométriquement
irréductible sur $\kappa(s)$. Si nous en avons un jour besoin, nous ajouterons une démonstration détaillée.},
alors le nombre de branches de $X$ en $x$
est égal au nombre de branches de $S$ en $s$.
\end{enumerate}
\end{lemma}

\begin{proof}
Ceci résulte immédiatement de Plus sur l'algèbre, lemme
\ref{more-algebra-lemma-invariance-number-branches-smooth},
et des définitions.
\end{proof}




\section{Morphismes non ramifiés et \'etales}
\label{section-unramified-is-etale}

\noindent
Les morphismes non ramifiés sont parfois automatiquement \'etales.

\begin{lemma}
\label{lemma-unramified-dominant-unibranch-is-etale}
Soit $f : X \to Y$ un morphisme de schémas. Soit $x \in X$
d'image $y \in Y$. Supposons que
\begin{enumerate}
\item $Y$ soit intègre et géométriquement unibranche en $y$,
\item $f$ soit localement de type fini,
\item il existe une spécialisation $x' \leadsto x$ telle que $f(x')$
soit le point générique de $Y$,
\item $f$ soit non ramifié en $x$.
\end{enumerate}
Alors $f$ est \'etale en $x$.
\end{lemma}

\begin{proof}
On peut remplacer $X$ et $Y$ par des voisinages ouverts affines convenables
de $x$ et $y$. Alors $Y$ est le spectre d'un anneau intègre $A$ et
$X$ est le spectre d'une $A$-algèbre de type fini $B$.
Soient $\mathfrak q \subset B$ l'idéal premier correspondant
à $x$ et $\mathfrak p \subset A$ l'idéal premier correspondant à $y$.
L'anneau local $A_\mathfrak p = \mathcal{O}_{Y, y}$ est géométriquement unibranche.
Le morphisme d'anneaux $A \to B$ est non ramifié en $\mathfrak q$.
En outre, le point $x'$ de (3) correspond à un idéal premier
$\mathfrak q' \subset \mathfrak q$ tel que $A \cap \mathfrak q' = (0)$.
Il s'ensuit que $A_\mathfrak p \to B_\mathfrak q$ est injectif.
On conclut par
Plus sur l'algèbre, lemme
\ref{more-algebra-lemma-local-unramified-extension-unibranch-domain-is-etale}.
\end{proof}

\begin{lemma}
\label{lemma-global-unramified-dominant-unibranch-is-etale}
\begin{reference}
\cite[Exposé I, Corollaire 9.11]{SGA1}
\end{reference}
Soit $f : X \to Y$ un morphisme de schémas. Supposons que
\begin{enumerate}
\item $Y$ soit intègre et géométriquement unibranche,
\item au moins une composante irréductible de $X$ domine $Y$,
\item $f$ soit non ramifié, et
\item $X$ soit connexe.
\end{enumerate}
Alors $f$ est \'etale et $X$ est irréductible.
\end{lemma}

\begin{proof}
Soit $X' \subset X$ la composante irréductible qui domine $Y$.
Cela signifie que le point générique de $X'$ s'envoie sur le point générique de $Y$
(voir par exemple Morphismes, lemme
\ref{morphisms-lemma-dominant-finite-number-irreducible-components}).
D'après le lemme \ref{lemma-unramified-dominant-unibranch-is-etale},
$f$ est \'etale en tout point de $X'$. En particulier,
le sous-schéma ouvert $U \subset X$ où $f$ est \'etale contient $X'$.
Notons que tout ouvert quasi-compact de $U$ possède un nombre fini de composantes
irréductibles ; voir
Descente, lemme \ref{descent-lemma-locally-finite-nr-irred-local-fppf}.
D'autre part, puisque $Y$ est géométriquement unibranche et que $U$ est \'etale
sur $Y$, le schéma $U$ est géométriquement unibranche. En particulier, par
chaque point de $U$ passe au plus une composante irréductible.
Un argument topologique simple montre alors que $X' \subset U$
est à la fois ouvert et fermé. Ainsi, $X'$ est ouvert et fermé dans $X$ et, par
connexité, on obtient $X = U = X'$, comme souhaité.
\end{proof}

\begin{lemma}
\label{lemma-weird-permanence-etale}
Soient $f : X \to Y$ et $g : Y \to Z$ des morphismes de schémas.
Soit $x \in X$ d'image $y \in Y$. Supposons que
\begin{enumerate}
\item $Y$ soit intègre et géométriquement unibranche en $y$,
\item $g \circ f$ soit \'etale en $x$,
\item il existe une spécialisation $x' \leadsto x$ telle que $f(x')$
soit le point générique de $Y$.
\end{enumerate}
Alors $f$ est \'etale en $x$ et $g$ est \'etale en $y$.
\end{lemma}

\begin{proof}
Après avoir remplacé $X$ par un voisinage ouvert de $x$,
on peut supposer que $g \circ f$ est \'etale.
On trouve alors que $f$ est non ramifié par Morphismes, lemmes
\ref{morphisms-lemma-unramified-permanence}
et \ref{morphisms-lemma-etale-smooth-unramified}.
Ainsi, $f$ est \'etale en $x$ d'après le
lemme \ref{lemma-unramified-dominant-unibranch-is-etale}.
Par descente \'etale, on voit alors que $g$ est \'etale en $y$ ; voir
par exemple Descente, lemme
\ref{descent-lemma-syntomic-smooth-etale-permanence}.
\end{proof}

\begin{lemma}
\label{lemma-global-weird-permanence-etale}
Soient $f : X \to Y$ et $g : Y \to Z$ des morphismes de schémas.
Supposons que
\begin{enumerate}
\item $Y$ soit intègre et géométriquement unibranche,
\item $g \circ f$ soit \'etale,
\item toute composante irréductible de $X$ domine $Y$.
\end{enumerate}
Alors $f$ est \'etale et $g$ est \'etale en tout point de
l'image de $f$.
\end{lemma}

\begin{proof}
Ceci résulte immédiatement de la version ponctuelle,
lemme \ref{lemma-weird-permanence-etale}.
\end{proof}










\section{Découpage des morphismes lisses}
\label{section-etale-over-smooth}

\noindent
Dans cette section, nous expliquons un résultat qui affirme, en gros, que
les recouvrements lisses d'un schéma $S$ peuvent être raffinés par des recouvrements \'etales.
La technique de démonstration repose sur un argument de découpage.

\begin{lemma}
\label{lemma-slice-smooth-given-element}
Soit $f : X \to S$ un morphisme de schémas.
Soit $x \in X$ un point d'image $s \in S$.
Soit $h \in \mathfrak m_x \subset \mathcal{O}_{X, x}$.
Supposons que
\begin{enumerate}
\item $f$ soit lisse en $x$, et
\item l'image $\text{d}\overline{h}$ de $\text{d}h$ dans
$$
\Omega_{X_s/s, x} \otimes_{\mathcal{O}_{X_s, x}} \kappa(x) =
\Omega_{X/S, x} \otimes_{\mathcal{O}_{X, x}} \kappa(x)
$$
soit non nulle.
\end{enumerate}
Alors il existe un voisinage ouvert affine $U \subset X$ de $x$
tel que $h$ provienne d'un $h \in \Gamma(U, \mathcal{O}_U)$ et
que $D = V(h)$ soit un diviseur de Cartier effectif dans $U$, avec $x \in D$, et que
$D \to S$ soit lisse.
\end{lemma}

\begin{proof}
Puisque $f$ est lisse en $x$, on peut supposer, après avoir remplacé $X$ par un voisinage
ouvert de $x$, que $f$ est lisse. En particulier, on voit que
$f$ est plat et localement de présentation finie. D'après le
lemme \ref{lemma-slice-given-element},
on sait déjà qu'il existe un voisinage ouvert $U \subset X$ de $x$
tel que $h$ provienne d'un $h \in \Gamma(U, \mathcal{O}_U)$ et
que $D = V(h)$ soit un diviseur de Cartier effectif dans $U$, avec $x \in D$, et que
$D \to S$ soit plat et de présentation finie. D'après
Morphismes, lemme \ref{morphisms-lemma-differentials-relative-immersion},
on dispose d'une suite exacte à droite
$$
\mathcal{C}_{D/U} \to i^*\Omega_{U/S} \to \Omega_{D/S} \to 0
$$
où $i : D \to U$ est l'immersion fermée et $\mathcal{C}_{D/U}$
est le faisceau conormal de $D$ dans $U$. Puisque $D$ est un diviseur de Cartier
effectif défini par $h \in \Gamma(U, \mathcal{O}_U)$, on voit que
$\mathcal{C}_{D/U} = h \cdot \mathcal{O}_S$. Comme $U \to S$ est lisse,
le faisceau $\Omega_{U/S}$ est localement libre de type fini ; son image réciproque
$i^*\Omega_{U/S}$ est donc elle aussi localement libre de type fini. La première flèche de
la suite envoie le générateur libre $h$ sur la section $\text{d}h|_D$
de $i^*\Omega_{U/S}$, dont la valeur dans la fibre
$\Omega_{U/S, x} \otimes \kappa(x)$ est non nulle par hypothèse. Par exactitude à droite
de $\otimes \kappa(x)$, on conclut que
$$
\dim_{\kappa(x)} \left( \Omega_{D/S, x} \otimes \kappa(x) \right)
=
\dim_{\kappa(x)} \left( \Omega_{U/S, x} \otimes \kappa(x) \right) - 1.
$$
D'après
Morphismes, lemme \ref{morphisms-lemma-smooth-at-point},
on voit que $\Omega_{U/S, x} \otimes \kappa(x)$ peut être engendré par
au plus $\dim_x(U_s)$ éléments. D'après la formule affichée, on voit que
$\Omega_{D/S, x} \otimes \kappa(x)$ peut être engendré par au plus
$\dim_x(U_s) - 1$ éléments. Notons que
$\dim_x(D_s) = \dim_x(U_s) - 1$, par exemple parce que
$\dim(\mathcal{O}_{D_s, x}) = \dim(\mathcal{O}_{U_s, x}) - 1$ d'après
Algèbre, lemme \ref{algebra-lemma-one-equation}
(de plus, $D_s \subset U_s$ est un diviseur de Cartier effectif ; voir
Diviseurs, lemme \ref{divisors-lemma-relative-Cartier}),
puis en utilisant
Morphismes, lemme \ref{morphisms-lemma-dimension-fibre-at-a-point}.
On conclut ainsi que $\Omega_{D/S, x} \otimes \kappa(x)$ peut être engendré
par au plus $\dim_x(D_s)$ éléments, et l'on conclut que $D \to S$
est lisse en $x$ en appliquant de nouveau
Morphismes, lemme \ref{morphisms-lemma-smooth-at-point}.
Après avoir rétréci $U$, on obtient que $D \to S$ est lisse, ce qui conclut.
\end{proof}

\begin{lemma}
\label{lemma-slice-smooth-once}
Soit $f : X \to S$ un morphisme de schémas.
Soit $x \in X$ un point d'image $s \in S$.
Supposons que
\begin{enumerate}
\item $f$ soit lisse en $x$, et
\item le morphisme
$$
\Omega_{X_s/s, x} \otimes_{\mathcal{O}_{X_s, x}} \kappa(x)
\longrightarrow
\Omega_{\kappa(x)/\kappa(s)}
$$
ait un noyau non nul.
\end{enumerate}
Alors il existe un voisinage ouvert affine $U \subset X$ de $x$
et un diviseur de Cartier effectif $D \subset U$ contenant $x$, tels que
$D \to S$ soit lisse.
\end{lemma}

\begin{proof}
Posons $k = \kappa(s)$ et $R = \mathcal{O}_{X_s, x}$.
Notons $\mathfrak m$ l'idéal maximal de $R$ et posons
$\kappa = R/\mathfrak m$, de sorte que $\kappa = \kappa(x)$.
Comme la formation des modules de différentielles commute à la localisation (voir
Algèbre, lemme \ref{algebra-lemma-differentials-localize}),
on a $\Omega_{X_s/s, x} = \Omega_{R/k}$. D'après
Algèbre, lemme \ref{algebra-lemma-differential-seq},
il existe une suite exacte
$$
\mathfrak m/\mathfrak m^2 \xrightarrow{\text{d}}
\Omega_{R/k} \otimes_R \kappa \to
\Omega_{\kappa/k} \to 0.
$$
Par conséquent, si (2) est satisfaite, il existe un élément $\overline{h} \in \mathfrak m$
tel que $\text{d}\overline{h}$ soit non nul. Choisissons un relèvement
$h \in \mathcal{O}_{X, x}$ de $\overline{h}$ et appliquons le
lemme \ref{lemma-slice-smooth-given-element}.
\end{proof}

\begin{remark}
\label{remark-necessary-condition-slice-smooth}
La seconde condition du
lemme \ref{lemma-slice-smooth-once}
est nécessaire, même lorsque $x$ est un point fermé d'une fibre de dimension
strictement positive. En voici un exemple. Soit $k$ un corps
imparfait de caractéristique $p > 0$. Soit $a \in k$ un
élément qui n'est pas une puissance $p$-ième. Posons
$\mathfrak m = (x, y^p - a) \subset k[x, y]$. Ceci correspond à un point fermé
$w$ de $X = \mathbf{A}^2_k$. Posons $S = \mathbf{A}^1_k$ et
soit $f : X \to S$ le morphisme correspondant à $k[x] \to k[x, y]$.
Il n'existe alors aucun diagramme commutatif
$$
\xymatrix{
S' \ar[rr]_h \ar[rd]_g & & X \ar[ld]^f \\
& S
}
$$
avec $g$ \'etale et $w$ dans l'image de $h$. En effet, l'extension de corps
résiduels $\kappa(w)/\kappa(f(w))$ est purement inséparable,
tandis que, pour tout $s' \in S'$ tel que $g(s') = f(w)$, l'extension
$\kappa(s')/\kappa(f(w))$ serait séparable.
\end{remark}

\noindent
Si l'on suppose que l'extension de corps résiduels est séparable, le
phénomène de la
remarque \ref{remark-necessary-condition-slice-smooth}
ne se produit pas. Voici l'énoncé précis.

\begin{lemma}
\label{lemma-slice-smooth-once-separable-residue-field-extension}
Soit $f : X \to S$ un morphisme de schémas.
Soit $x \in X$ un point d'image $s \in S$.
Supposons que
\begin{enumerate}
\item $f$ soit lisse en $x$,
\item l'extension de corps résiduels $\kappa(x)/\kappa(s)$
soit séparable, et
\item $x$ ne soit pas un point générique de $X_s$.
\end{enumerate}
Alors il existe un voisinage ouvert affine $U \subset X$ de $x$
et un diviseur de Cartier effectif $D \subset U$ contenant $x$, tels que
$D \to S$ soit lisse.
\end{lemma}

\begin{proof}
Posons $k = \kappa(s)$ et $R = \mathcal{O}_{X_s, x}$.
Notons $\mathfrak m$ l'idéal maximal de $R$ et posons
$\kappa = R/\mathfrak m$, de sorte que $\kappa = \kappa(x)$.
Comme la formation des modules de différentielles commute à la localisation (voir
Algèbre, lemme \ref{algebra-lemma-differentials-localize}),
on a $\Omega_{X_s/s, x} = \Omega_{R/k}$. D'après l'hypothèse (2) et
Algèbre, lemme \ref{algebra-lemma-computation-differential},
le morphisme
$$
\text{d} :
\mathfrak m/\mathfrak m^2
\longrightarrow
\Omega_{R/k} \otimes_R \kappa(\mathfrak m)
$$
est injectif. L'hypothèse (3) implique que
$\mathfrak m/\mathfrak m^2 \not = 0$.
Il existe donc un élément $\overline{h} \in \mathfrak m$
tel que $\text{d}\overline{h}$ soit non nul. Choisissons un relèvement
$h \in \mathcal{O}_{X, x}$ de $\overline{h}$ et appliquons le
lemme \ref{lemma-slice-smooth-given-element}.
\end{proof}

\noindent
Le sous-schéma $Z$ construit dans le lemme suivant est en fait une intersection
complète dans un voisinage ouvert affine de $x$. Si nous en avons un jour besoin,
nous formulerons explicitement un lemme séparé énonçant ce fait.

\begin{lemma}
\label{lemma-slice-smooth}
Soit $f : X \to S$ un morphisme de schémas.
Soit $x \in X$ un point d'image $s \in S$.
Supposons que
\begin{enumerate}
\item $f$ soit lisse en $x$, et
\item $x$ soit un point fermé de $X_s$ et que $\kappa(s) \subset \kappa(x)$
soit séparable.
\end{enumerate}
Alors il existe une immersion $Z \to X$ dont l'image contient $x$ telle que
\begin{enumerate}
\item $Z \to S$ soit \'etale, et
\item $Z_s = \{x\}$ ensemblistement.
\end{enumerate}
\end{lemma}

\begin{proof}
On peut, et l'on va, remplacer $S$ par un voisinage ouvert affine de $s$.
On peut, et l'on va, remplacer $X$ par un voisinage ouvert affine de $x$
tel que $X \to S$ soit lisse.
Nous allons démontrer le lemme pour les morphismes lisses entre schémas affines
par récurrence sur $d = \dim_x(X_s)$.

\medskip\noindent
Le cas $d = 0$. Dans ce cas, nous montrons que l'on peut prendre pour $Z$
un voisinage ouvert de $x$. En effet, si $d = 0$, alors $X \to S$
est quasi-fini en $x$ ; voir
Morphismes, lemme \ref{morphisms-lemma-locally-quasi-finite-rel-dimension-0}.
Il existe donc un voisinage ouvert affine $U \subset X$
tel que $U \to S$ soit quasi-fini ; voir
Morphismes, lemme \ref{morphisms-lemma-quasi-finite-points-open}.
Ainsi, après avoir remplacé $X$ par $U$, on voit que
$X$ est quasi-fini et lisse sur $S$, donc
lisse de dimension relative $0$ sur $S$, et donc
\'etale sur $S$. De plus, la fibre $X_s$ est un ensemble fini
discret. Après avoir remplacé $X$ par un voisinage ouvert affine supplémentaire
de $X$, on voit donc que $f^{-1}(\{s\}) = \{x\}$ (car la topologie
de $X_s$ est induite par celle de $X$ ; voir
Schémas, lemme \ref{schemes-lemma-fibre-topological}).
Ceci démontre le lemme dans ce cas.

\medskip\noindent
Supposons maintenant que $d > 0$. Notons que, puisque $x$ est un point fermé de sa
fibre, l'extension $\kappa(x)/\kappa(s)$ est finie (d'après le
théorème des zéros de Hilbert ; voir
Morphismes, lemme \ref{morphisms-lemma-closed-point-fibre-locally-finite-type}).
On a donc $\Omega_{\kappa(x)/\kappa(s)} = 0$, puisque ceci vaut pour les
extensions de corps algébriques séparables.
On peut donc appliquer le
lemme \ref{lemma-slice-smooth-once}
pour obtenir un diagramme
$$
\xymatrix{
D \ar[r] \ar[rrd] & U \ar[r] \ar[rd] & X \ar[d] \\
& & S
}
$$
avec $x \in D$. Notons que
$\dim_x(D_s) = \dim_x(X_s) - 1$, par exemple parce que
$\dim(\mathcal{O}_{D_s, x}) = \dim(\mathcal{O}_{X_s, x}) - 1$ d'après
Algèbre, lemme \ref{algebra-lemma-one-equation}
(de plus, $D_s \subset X_s$ est un diviseur de Cartier effectif ; voir
Diviseurs, lemme \ref{divisors-lemma-relative-Cartier}),
puis en utilisant
Morphismes, lemme \ref{morphisms-lemma-dimension-fibre-at-a-point}.
Ainsi, le morphisme $D \to S$ est lisse et
$\dim_x(D_s) = \dim_x(X_s) - 1 = d - 1$. Par l'hypothèse de récurrence,
on peut trouver une immersion $Z \to D$ comme souhaité, ce qui achève la démonstration.
\end{proof}

\begin{lemma}
\label{lemma-etale-nbhd-dominates-smooth}
\begin{slogan}
Les morphismes lisses admettent des sections locales pour la topologie \'etale.
\end{slogan}
\begin{reference}
Voir \cite[Corollaire 17.16.3 (ii)]{EGA4}.
\end{reference}
Soit $f : X \to S$ un morphisme lisse de schémas.
Soit $s \in S$ un point de l'image de $f$.
Alors il existe un voisinage \'etale $(S', s') \to (S, s)$
et un $S$-morphisme $S' \to X$.
\end{lemma}

\begin{proof}[Première démonstration du lemme \ref{lemma-etale-nbhd-dominates-smooth}]
Par hypothèse, $X_s \not = \emptyset$. D'après
Variétés, lemme \ref{varieties-lemma-smooth-separable-closed-points-dense},
il existe un point fermé $x \in X_s$ tel que $\kappa(x)$
soit une extension finie séparable de $\kappa(s)$.
Par conséquent, d'après le
lemme \ref{lemma-slice-smooth},
il existe une immersion $Z \to X$ telle que $Z \to S$ soit \'etale et
que $x \in Z$. Prenons $(S' , s') = (Z, x)$.
\end{proof}

\begin{proof}[Seconde démonstration du lemme \ref{lemma-etale-nbhd-dominates-smooth}]
Choisissons un point $x \in X$ tel que $f(x) = s$.
Choisissons un diagramme
$$
\xymatrix{
X \ar[d] & U \ar[l] \ar[d] \ar[r]_-\pi & \mathbf{A}^d_V \ar[ld] \\
S & V \ar[l]
}
$$
avec $\pi$ \'etale, $x \in U$ et $V = \Spec(R)$ affine ; voir
Morphismes, lemme \ref{morphisms-lemma-smooth-etale-over-affine-space}.
En particulier, $s \in V$. Le morphisme
$\pi : U \to \mathbf{A}^d_V$ est ouvert ; voir
Morphismes, lemme \ref{morphisms-lemma-etale-open}.
Ainsi, $W = \pi(U) \cap \mathbf{A}^d_s$ est un ouvert non vide de
$\mathbf{A}^d_s$. Soit $w \in W$ un point tel que l'extension $\kappa(s) \subset \kappa(w)$
soit finie et séparable ; voir
Variétés, lemme \ref{varieties-lemma-affine-space-over-field}.
D'après
Algèbre, lemme \ref{algebra-lemma-dim-affine-space},
il existe $d$ éléments
$\overline{f}_1, \ldots, \overline{f}_d \in \kappa(s)[x_1, \ldots, x_d]$
qui engendrent l'idéal maximal correspondant à $w$ dans
$\kappa(s)[x_1, \ldots, x_d]$.
Après avoir remplacé $R$ par une localisation principale,
on peut supposer qu'il existe des $f_1, \ldots, f_d \in R[x_1, \ldots, x_d]$
dont les images sont
$\overline{f}_1, \ldots, \overline{f}_d \in \kappa(s)[x_1, \ldots, x_d]$.
Considérons la $R$-algèbre
$$
R' = R[x_1, \ldots, x_d]/(f_1, \ldots, f_d)
$$
et posons $S' = \Spec(R')$. Par construction, nous avons une
immersion fermée $j : S' \to \mathbf{A}^d_V$ au-dessus de $V$.
Par construction, la fibre de $S' \to V$ au-dessus de $s$ est constituée d'un unique
point $s'$ dont le corps résiduel est une extension finie séparable de $\kappa(s)$.
Soit $\mathfrak q' \subset R'$ l'idéal premier correspondant. D'après
Algèbre, lemme \ref{algebra-lemma-localize-relative-complete-intersection},
on voit que $(R')_g$ est une intersection complète globale relative sur $R$
pour un certain $g \in R'$, $g \not \in \mathfrak q$.
Ainsi, $S' \to V$ est plat et de présentation finie dans un
voisinage de $s'$ ; voir
Algèbre, lemme \ref{algebra-lemma-relative-global-complete-intersection}.
Par construction, la fibre schématique
de $S' \to V$ au-dessus de $s$ est $\Spec(\kappa(s'))$. Il résulte donc de
Morphismes, lemme \ref{morphisms-lemma-etale-at-point},
que $S' \to S$ est \'etale en $s'$. Posons
$$
S'' = U \times_{\pi, \mathbf{A}^d_V, j} S'.
$$
Par construction, il existe un point $s'' \in S''$ qui s'envoie sur
$s'$ par la projection $p : S'' \to S'$. Notons que $p$ est \'etale
comme changement de base du morphisme \'etale $\pi$ ; voir
Morphismes, lemme \ref{morphisms-lemma-base-change-etale}.
Choisissons un petit voisinage ouvert affine $S''' \subset S''$ de $s''$
qui s'envoie dans le voisinage ouvert non vide de $s' \in S'$
où le morphisme $S' \to S$ est \'etale. Alors le voisinage \'etale
$(S''', s'') \to (S, s)$ est une solution au problème posé par le lemme.
\end{proof}

\noindent
Le lemme suivant montre que les faisceaux pour la topologie lisse sont
la même chose que les faisceaux pour la topologie \'etale.

\begin{lemma}
\label{lemma-etale-dominates-smooth}
Soit $S$ un schéma. Soit $\mathcal{U} = \{S_i \to S\}_{i \in I}$ un recouvrement
lisse de $S$ ; voir
Topologies, définition \ref{topologies-definition-smooth-covering}.
Alors il existe un recouvrement \'etale $\mathcal{V} = \{T_j \to S\}_{j \in J}$
(voir
Topologies, définition \ref{topologies-definition-etale-covering})
qui raffine (voir
Sites, définition \ref{sites-definition-morphism-coverings})
$\mathcal{U}$.
\end{lemma}

\begin{proof}
Pour tout $s \in S$, il existe un $i \in I$ tel que $s$ appartienne à
l'image de $S_i \to S$. D'après le
lemme \ref{lemma-etale-nbhd-dominates-smooth},
on peut trouver un morphisme \'etale $g_s : T_s \to S$ tel que $s \in g_s(T_s)$
et tel que $g_s$ se factorise par $S_i \to S$. Ainsi,
$\{T_s \to S\}$ est un recouvrement \'etale de $S$ qui raffine $\mathcal{U}$.
\end{proof}

\begin{lemma}
\label{lemma-cover-smooth-by-special}
Soit $f : X \to S$ un morphisme lisse de schémas. Alors il existe un
recouvrement \'etale $\{U_i \to X\}_{i \in I}$ tel que $U_i \to S$
se factorise sous la forme $U_i \to V_i \to S$, où $V_i \to S$ est \'etale et
$U_i \to V_i$ est un morphisme lisse de schémas affines, qui
admet une section et dont les fibres sont géométriquement connexes.
\end{lemma}

\begin{proof}
Soit $s \in S$. D'après
Variétés, lemme \ref{varieties-lemma-smooth-separable-closed-points-dense},
l'ensemble des points fermés $x \in X_s$ tels que $\kappa(x)/\kappa(s)$
soit séparable est dense dans $X_s$. Il suffit donc de construire un
morphisme \'etale $U \to X$ dont l'image contient $x$
et tel que $U \to S$ se factorise de la manière décrite dans le lemme.
Pour ce faire, choisissons une immersion $Z \to X$ passant par $x$
telle que $Z \to S$ soit \'etale (lemme \ref{lemma-slice-smooth}).
Après avoir remplacé $S$ par $Z$ et $X$ par $Z \times_S X$,
on voit que l'on peut supposer que $X \to S$ admet une section $\sigma : S \to X$
avec $\sigma(s) = x$. On peut alors d'abord remplacer $S$ par un voisinage
ouvert affine de $s$, puis remplacer $X$ par un voisinage ouvert affine
de $x$. Enfin, considérons le sous-ensemble
$X^0 \subset X$ de la section \ref{section-connected-components}.
D'après les lemmes \ref{lemma-connected-along-section-open} et
\ref{lemma-connected-along-section-locally-constructible},
c'est un sous-schéma ouvert rétrocompact contenant $\sigma$
tel que les fibres de $X^0 \to S$ soient géométriquement connexes.
Si $X^0$ n'est pas affine, choisissons un ouvert affine $U \subset X^0$
contenant $x$. Comme $X^0 \to S$ est lisse, l'image de $U$
est ouverte. Choisissons un voisinage ouvert affine $V \subset S$ de $s$
contenu dans $\sigma^{-1}(U)$ et dans l'image de $U \to S$.
Enfin, le lecteur vérifiera que $U \cap f^{-1}(V) \to V$
possède toutes les propriétés souhaitées. Par exemple, $U \cap f^{-1}(V)$,
qui est égal à $U \times_S V$, est affine comme produit fibré de schémas
affines. En outre, les fibres géométriques de $U \cap f^{-1}(V) \to V$ sont des
sous-schémas ouverts non vides des fibres irréductibles de $X^0 \to S$
et sont donc connexes. Certains détails sont omis.
\end{proof}











\section{Voisinages \'etales et approximation d'Artin}
\label{section-etale-nbhds-artin}

\noindent
Dans cette section, nous démontrons des résultats du type suivant : si deux
schémas pointés ont des anneaux locaux complets isomorphes, alors
ils ont des voisinages \'etales isomorphes. Nous nous appuierons
sur le théorème de Popescu ; voir
Lissage des morphismes d'anneaux, théorème \ref{smoothing-theorem-popescu}.

\begin{lemma}
\label{lemma-map-approximation}
Soit $S$ un schéma localement noethérien. Soient $X$, $Y$ des
schémas localement de type
fini sur $S$. Soient $x \in X$ et $y \in Y$ des points situés au-dessus du
même point $s \in S$. Supposons que $\mathcal{O}_{S, s}$ soit un G-anneau.
Supposons en outre donné un homomorphisme local de $\mathcal{O}_{S, s}$-algèbres
$$
\varphi : \mathcal{O}_{Y, y} \longrightarrow \mathcal{O}_{X, x}^\wedge
$$
Pour tout $N \geq 1$,
il existe un voisinage \'etale élémentaire
$(U, u) \to (X, x)$ et un $S$-morphisme
$f : U \to Y$ envoyant $u$ sur $y$ tels que le diagramme
$$
\xymatrix{
\mathcal{O}_{X, x}^\wedge \ar[r] &
\mathcal{O}_{U, u}^\wedge \\
\mathcal{O}_{Y, y} \ar[r]^{f^\sharp_u} \ar[u]^\varphi &
\mathcal{O}_{U, u} \ar[u]
}
$$
soit commutatif modulo $\mathfrak m_u^N$.
\end{lemma}

\begin{proof}
La question est locale sur $X$ ; on peut donc supposer que $X$, $Y$, $S$ sont affines.
Écrivons $S = \Spec(R)$, $X = \Spec(A)$, $Y = \Spec(B)$.
Écrivons $B = R[x_1, \ldots, x_n]/(f_1, \ldots, f_m)$.
Soit $\mathfrak p \subset A$ l'idéal premier correspondant à $x$.
L'anneau local $\mathcal{O}_{X, x} = A_\mathfrak p$ est un G-anneau d'après
Compléments d'algèbre, proposition
\ref{more-algebra-proposition-finite-type-over-G-ring}.
L'homomorphisme $\varphi$ est un homomorphisme
$$
B_\mathfrak q^\wedge \longrightarrow A_\mathfrak p^\wedge
$$
où $\mathfrak q \subset B$ est l'idéal premier correspondant à $y$.
Soient $a_1, \ldots, a_n \in A_\mathfrak p^\wedge$ les images
de $x_1, \ldots, x_n$ par
$R[x_1, \ldots, x_n] \to B \to B_\mathfrak q^\wedge \to A_\mathfrak p^\wedge$.
On peut alors appliquer Lissage des morphismes d'anneaux, lemme
\ref{smoothing-lemma-approximation-property-variant},
pour obtenir un homomorphisme \'etale d'anneaux $A \to A'$, un idéal premier
$\mathfrak p' \subset A'$ et des éléments $b_1, \ldots, b_n \in A'$ tels que
$\kappa(\mathfrak p) = \kappa(\mathfrak p')$,
$a_i - b_i \in (\mathfrak p')^N(A'_{\mathfrak p'})^\wedge$, et
$f_j(b_1, \ldots, b_n) = 0$ pour $j = 1, \ldots, n$.
Ceci détermine un homomorphisme de $R$-algèbres $B \to A'$ envoyant la
classe de $x_i$ sur $b_i \in A'$. Ceci achève la démonstration
en prenant $U = \Spec(A') \to \Spec(B)$ pour le morphisme $f$
et en prenant $u = \mathfrak p'$.
\end{proof}

\begin{lemma}
\label{lemma-isomorphism-approximation}
Soit $S$ un schéma localement noethérien. Soient $X$, $Y$ des
schémas localement de type
fini sur $S$. Soient $x \in X$ et $y \in Y$ des points situés au-dessus du
même point $s \in S$. Supposons que $\mathcal{O}_{S, s}$ soit un G-anneau.
Supposons donné un isomorphisme de $\mathcal{O}_{S, s}$-algèbres
$$
\varphi : \mathcal{O}_{Y, y}^\wedge \longrightarrow \mathcal{O}_{X, x}^\wedge
$$
entre les anneaux locaux complets. Alors, pour tout $N \geq 1$,
il existe des morphismes
$$
(X, x) \leftarrow (U, u) \rightarrow (Y, y)
$$
de schémas pointés au-dessus de $S$ tels que les deux flèches définissent des
voisinages \'etales élémentaires et tels que le diagramme
$$
\xymatrix{
& \mathcal{O}_{U, u}^\wedge \\
\mathcal{O}_{Y, y}^\wedge \ar[rr]^\varphi \ar[ru] & &
\mathcal{O}_{X, x}^\wedge \ar[lu]
}
$$
soit commutatif modulo $\mathfrak m_u^N$.
\end{lemma}

\begin{proof}
On peut supposer $N \geq 2$. Appliquons le lemme \ref{lemma-map-approximation} pour obtenir
$(U, u) \to (X, x)$ et $f : (U, u) \to (Y, y)$.
Nous affirmons que $f$ est \'etale en $u$, ce qui achèvera la démonstration.
En fait, nous allons montrer que l'homomorphisme induit
$\mathcal{O}_{Y, y}^\wedge \to \mathcal{O}_{U, u}^\wedge$
est un isomorphisme. Une fois ceci démontré, le
lemme \ref{lemma-lifting-along-artinian-at-point}
montrera que $f$ est lisse en $u$ et, bien entendu,
$f$ est également non ramifié en $u$ ; ainsi,
Morphismes, lemme \ref{morphisms-lemma-etale-smooth-unramified},
nous dit que $f$ est \'etale en $u$.
Pour un anneau local $(R, \mathfrak m)$, posons
$\text{Gr}_\mathfrak m(R) =
\bigoplus_{n \geq 0} \mathfrak m^n/\mathfrak m^{n + 1}$.
Pour démontrer l'affirmation, considérons le diagramme induit
d'anneaux gradués
$$
\xymatrix{
& \text{Gr}_{\mathfrak m_u}(\mathcal{O}_{U, u}) \\
\text{Gr}_{\mathfrak m_y}(\mathcal{O}_{Y, y}) \ar[rr]^\varphi \ar[ru] & &
\text{Gr}_{\mathfrak m_x}(\mathcal{O}_{X, x}) \ar[lu]
}
$$
Puisque $N \geq 2$, ce diagramme est effectivement commutatif, car les
algèbres graduées affichées sont engendrées en degré $1$ !
Par hypothèse, la flèche inférieure est un isomorphisme.
D'après Compléments d'algèbre, lemme \ref{more-algebra-lemma-flat-unramified}
(par exemple), l'homomorphisme
$\mathcal{O}_{X, x}^\wedge \to \mathcal{O}_{U, u}^\wedge$
est un isomorphisme ; la flèche nord-ouest
du diagramme est donc un isomorphisme. On en conclut que
$f$ induit un isomorphisme
$\text{Gr}_{\mathfrak m_x}(\mathcal{O}_{X, x}) \to
\text{Gr}_{\mathfrak m_y}(\mathcal{O}_{U, u})$.
Par récurrence, en utilisant les suites exactes courtes
$$
0 \to \text{Gr}^n_{\mathfrak m}(R) \to R/\mathfrak m^{n + 1} \to
R/\mathfrak m^n \to 0
$$
pour les deux anneaux locaux, on conclut (par le lemme du serpent)
que $f$ induit des isomorphismes
$\mathcal{O}_{Y, y}/\mathfrak m_y^n \to \mathcal{O}_{U, u}/\mathfrak m_u^n$
pour tout $n$, ce que nous voulions montrer.
\end{proof}

\begin{lemma}
\label{lemma-relative-map-approximation-pre}
Donnons-nous $X \to S$, $Y \to T$, $x$, $s$, $y$, $t$, $\sigma$, $y_\sigma$ et
$\varphi$ comme suit : nous avons des morphismes de schémas
$$
\vcenter{
\xymatrix{
X \ar[d] & Y \ar[d] \\
S & T
}
}
\quad\text{avec les points}\quad
\vcenter{
\xymatrix{
x \ar[d] & y \ar[d] \\
s & t
}
}
$$
Ici, $S$ est localement noethérien et $T$ est de type fini sur $\mathbf{Z}$.
Les morphismes $X \to S$ et $Y \to T$ sont localement de type fini.
L'anneau local $\mathcal{O}_{S, s}$ est un G-anneau. L'homomorphisme
$$
\sigma : \mathcal{O}_{T, t} \longrightarrow \mathcal{O}_{S, s}^\wedge
$$
est un homomorphisme local. Posons
$Y_\sigma = Y \times_{T, \sigma} \Spec(\mathcal{O}_{S, s}^\wedge)$.
Ensuite, $y_\sigma$ est un point de $Y_\sigma$ qui s'envoie sur $y$ et sur
le point fermé de $\Spec(\mathcal{O}_{S, s}^\wedge)$. Enfin,
$$
\varphi :
\mathcal{O}_{X, x}^\wedge
\longrightarrow
\mathcal{O}_{Y_\sigma, y_\sigma}^\wedge
$$
est un isomorphisme de $\mathcal{O}_{S, s}^\wedge$-algèbres.
Dans cette situation, il existe un diagramme commutatif
$$
\xymatrix{
X \ar[d] &
W \ar[l] \ar[rd] \ar[rr] & &
Y \times_{T, \tau} V \ar[r] \ar[ld] & Y \ar[d] \\
S & &
V \ar[ll] \ar[rr]^\tau & &
T
}
$$
de schémas, ainsi que des points $w \in W$, $v \in V$, tels que
\begin{enumerate}
\item $(V, v) \to (S, s)$ est un voisinage \'etale élémentaire,
\item $(W, w) \to (X, x)$ est un voisinage \'etale élémentaire, et
\item $\tau(v) = t$.
\end{enumerate}
Soit $y_\tau \in Y \times_T V$ le point qui correspond à $y_\sigma$
par l'identification $(Y_\sigma)_s = (Y \times_T V)_v$.
Alors
\begin{enumerate}
\item[(4)] $(W, w) \to (Y \times_{T, \tau} V, y_\tau)$ est un voisinage
\'etale élémentaire.
\end{enumerate}
\end{lemma}

\begin{proof}
Notons $X_\sigma = X \times_S \Spec(\mathcal{O}_{S, s}^\wedge)$
et $x_\sigma \in X_\sigma$ l'unique point situé au-dessus de $x$.
Observons que $\mathcal{O}_{S, s}^\wedge$ est un G-anneau d'après
Compléments d'algèbre, proposition
\ref{more-algebra-proposition-Noetherian-complete-G-ring}.
D'après le lemme \ref{lemma-isomorphism-approximation},
on peut choisir
$$
(X_\sigma, x_\sigma) \leftarrow (U, u) \rightarrow (Y_\sigma, y_\sigma)
$$
où les deux flèches sont des voisinages \'etales élémentaires.

\medskip\noindent
Après avoir remplacé $S$ par un voisinage ouvert de $s$, on peut
supposer que $S = \Spec(R)$ est affine. Puisque $\mathcal{O}_{S, s}$
est un G-anneau, le théorème \ref{smoothing-theorem-popescu} de Lissage des morphismes d'anneaux montre que
l'anneau $\mathcal{O}_{S, s}^\wedge$ est une colimite filtrante de
$R$-algèbres lisses. On peut donc écrire
$$
\Spec(\mathcal{O}_{S, s}^\wedge) = \lim S_i
$$
comme limite filtrante de schémas affines $S_i$ lisses sur $S$.
Notons $s_i \in S_i$ l'image du point fermé de
$\Spec(\mathcal{O}_{S, s}^\wedge)$. Observons que $\kappa(s) = \kappa(s_i)$.
Posons $X_i = X \times_S S_i$ et notons $x_i \in X_i$ l'unique
point s'envoyant sur $x$. Notons que $\kappa(x) = \kappa(x_i)$.
Comme $T$ est de type fini sur $\mathbf{Z}$, d'après Limites, proposition
\ref{limits-proposition-characterize-locally-finite-presentation},
on peut choisir un $i$ et un morphisme $\sigma_i : (S_i, s_i) \to (T, t)$
de schémas pointés dont la composée avec
$\Spec(\mathcal{O}_{S, s}^\wedge) \to S_i$ est égale à $\sigma$.
Posons $Y_i = Y \times_T S_i$ et notons $y_i$ l'image de
$y_\sigma$. Notons que $\kappa(y_i) = \kappa(y_\sigma)$.
D'après Limites, lemme \ref{limits-lemma-descend-finite-presentation},
on peut choisir un $i$ et un diagramme
$$
\xymatrix{
X_i \ar[rd] &
U_i \ar[l] \ar[d] \ar[r] &
Y_i \ar[ld] \\
& S_i
}
$$
dont le changement de base à $\Spec(\mathcal{O}_{S, s}^\wedge)$
redonne $X_\sigma \leftarrow U \rightarrow Y_\sigma$.
D'après Limites, lemme \ref{limits-lemma-descend-etale},
quitte à augmenter $i$, on peut supposer que les morphismes
$X_i \leftarrow U_i \rightarrow Y_i$ sont \'etales.
Soit $u_i \in U_i$ l'image de $u$. Alors $u_i \mapsto x_i$,
d'où
$\kappa(x) = \kappa(x_\sigma) = \kappa(u) \supset \kappa(u_i) \supset
\kappa(x_i) = \kappa(x)$, et l'on voit que $\kappa(u_i) = \kappa(x_i)$.
Ainsi, $(X_i, x_i) \leftarrow (U_i, u_i)$ est un voisinage
\'etale élémentaire. Comme, de plus, $\kappa(y_i) = \kappa(y_\sigma) = \kappa(u)$,
on voit que $(U_i, u_i) \to (Y_i, y_i)$ est également un voisinage
\'etale élémentaire.

\medskip\noindent
À ce stade, nous avons construit un diagramme
$$
\xymatrix{
X \ar[d] &
X \times_S S_i \ar[l] \ar[rd] &
U_i \ar[l] \ar[r] \ar[d] &
Y \times_T S_i \ar[r] \ar[ld] &
Y \ar[d] \\
S & &
S_i \ar[ll] \ar[rr] & &
T
}
$$
comme dans l'énoncé du lemme, à ceci près que $S_i \to S$ est lisse.
D'après le lemme \ref{lemma-slice-smooth}, et après avoir rétréci $S_i$,
on peut supposer qu'il existe un sous-schéma fermé $V \subset S_i$ passant
par $s_i$ tel que $V \to S$ soit \'etale.
En prenant pour $W$ l'image réciproque schématique de $V$
dans $U_i$, on conclut.
\end{proof}

\noindent
Nous recommandons vivement au lecteur de ne pas lire le reste de cette section.

\begin{lemma}
\label{lemma-relative-map-approximation}
Considérons un diagramme
$$
\vcenter{
\xymatrix{
X \ar[d] & Y \ar[d] \\
S & T \ar[l]
}
}
\quad\text{avec les points}\quad
\vcenter{
\xymatrix{
x \ar[d] & y \ar[d] \\
s & t \ar[l]
}
}
$$
où $S$ est un schéma localement noethérien et les morphismes sont
localement de type fini. Supposons que $\mathcal{O}_{S, s}$ soit un G-anneau.
Supposons en outre donné un homomorphisme local de $\mathcal{O}_{S, s}$-algèbres
$$
\sigma : \mathcal{O}_{T, t} \longrightarrow \mathcal{O}_{S, s}^\wedge
$$
et un homomorphisme local de $\mathcal{O}_{S, s}$-algèbres
$$
\varphi :
\mathcal{O}_{X, x}
\longrightarrow
\mathcal{O}_{Y_\sigma, y_\sigma}^\wedge
$$
où $Y_\sigma = Y \times_{T, \sigma} \Spec(\mathcal{O}_{S, s}^\wedge)$
et $y_\sigma$ est l'unique point de $Y_\sigma$ situé au-dessus de $y$.
Pour tout $N \geq 1$, il existe un diagramme commutatif
$$
\xymatrix{
X \ar[d] & X \times_S V \ar[l] \ar[rd] &
W \ar[l]^-f \ar[r] \ar[d] &
Y \times_{T, \tau} V \ar[r] \ar[ld] & Y \ar[d] \\
S & & V \ar[ll] \ar[rr]^\tau & & T
}
$$
de schémas au-dessus de $S$, ainsi que des points $w \in W$, $v \in V$, tels que
\begin{enumerate}
\item $v \mapsto s$, $\tau(v) = t$, $f(w) = (x, v)$, et $w \mapsto (y, v)$,
\item $(V, v) \to (S, s)$ est un voisinage \'etale élémentaire,
\item le diagramme
$$
\xymatrix{
\mathcal{O}_{S, s}^\wedge \ar[r] & \mathcal{O}_{V, v}^\wedge \\
\mathcal{O}_{T, t} \ar[r]^{\tau^\sharp_v} \ar[u]_\sigma &
\mathcal{O}_{V, v} \ar[u]
}
$$
est commutatif modulo $\mathfrak m_v^N$,
\item $(W, w) \to (Y \times_{T, \tau} V, (y, v))$ est un
voisinage \'etale élémentaire,
\item le diagramme
$$
\xymatrix{
\mathcal{O}_{X, x} \ar[r]_\varphi &
\mathcal{O}_{Y_\sigma, y_\sigma}^\wedge \ar[r] &
\mathcal{O}_{Y_\sigma, y_\sigma}/\mathfrak m_{y_\sigma}^N \ar@{=}[r] &
\mathcal{O}_{Y \times_{T, \tau} V, (y, v)}/\mathfrak m_{(y, v)}^N
\ar[d]_{\cong} \\
\mathcal{O}_{X, x} \ar[r] \ar@{=}[u] &
\mathcal{O}_{X \times_S V, (x, v)} \ar[r]^{f^\sharp_w} &
\mathcal{O}_{W, w} \ar[r] &
\mathcal{O}_{W, w}/\mathfrak m_w^N
}
$$
est commutatif. L'égalité vient du fait que
$Y_\sigma$ et $Y \times_{T, \tau} V$ sont canoniquement isomorphes au-dessus de
$\mathcal{O}_{V, v}/\mathfrak m_v^N = \mathcal{O}_{S, s}/\mathfrak m_s^N$
d'après les parties (2) et (3).
\end{enumerate}
\end{lemma}

\begin{proof}
Après avoir remplacé $X$, $S$, $T$, $Y$ par des sous-schémas ouverts affines, on
peut supposer que le diagramme de l'énoncé du lemme provient de
l'application de $\Spec$ à un diagramme
$$
\vcenter{
\xymatrix{
A & B \\
R \ar[u] \ar[r] & C \ar[u]
}
}
\quad\text{avec les idéaux premiers}\quad
\vcenter{
\xymatrix{
\mathfrak p_A & \mathfrak p_B \\
\mathfrak p_R \ar@{-}[u] \ar@{-}[r] & \mathfrak p_C \ar@{-}[u]
}
}
$$
d'anneaux noethériens et d'homomorphismes d'anneaux de type fini.
Dans cette démonstration, tout anneau $E$ sera une $R$-algèbre noethérienne munie
d'un idéal premier $\mathfrak p_E$ au-dessus de $\mathfrak p_R$, et tous les homomorphismes d'anneaux
seront des homomorphismes de $R$-algèbres compatibles avec les idéaux premiers donnés. De plus,
si nous écrivons $E^\wedge$, nous entendons le complété du localisé
de $E$ en $\mathfrak p_E$. Nous utiliserons aussi sans autre mention
qu'un homomorphisme \'etale d'anneaux $E_1 \to E_2$ tel que
$\kappa(\mathfrak p_{E_1}) = \kappa(\mathfrak p_{E_2})$ induit
un isomorphisme $E_1^\wedge = E_2^\wedge$ d'après
Compléments d'algèbre, lemme \ref{more-algebra-lemma-flat-unramified}.

\medskip\noindent
Avec ces notations, $\sigma$ et $\varphi$ correspondent aux homomorphismes d'anneaux
$$
\sigma : C \to R^\wedge
\quad\text{et}\quad
\varphi : A \longrightarrow (B \otimes_{C, \sigma} R^\wedge)^\wedge
$$
Voici un diagramme :
$$
\xymatrix{
A \ar@/^1em/[rrr]^\varphi &
B \ar[r] &
B \otimes_{C, \sigma} R^\wedge \ar[r] &
(B \otimes_{C, \sigma} R^\wedge)^\wedge \\
R \ar[r] \ar[u] &
C \ar[r]^\sigma \ar[u] & R^\wedge \ar[u] \ar[ru]
}
$$
Observons que $R^\wedge$ est un G-anneau d'après
Compléments d'algèbre, proposition
\ref{more-algebra-proposition-Noetherian-complete-G-ring}.
Ainsi, $B \otimes_{C, \sigma} R^\wedge$ est un G-anneau d'après
Compléments d'algèbre, proposition
\ref{more-algebra-proposition-finite-type-over-G-ring}.
D'après le lemme \ref{lemma-map-approximation} (sous sa forme algébrique),
il existe un homomorphisme \'etale d'anneaux $B \otimes_{C, \sigma} R^\wedge \to B'$,
qui induit un isomorphisme
$\kappa(\mathfrak p_{B \otimes_{C, \sigma} R^\wedge})
\to \kappa(\mathfrak p_{B'})$,
et un homomorphisme de $R$-algèbres $A \to B'$ tel que la composée
$$
A \to B' \to (B')^\wedge = (B \otimes_{C, \sigma} R^\wedge)^\wedge
$$
soit égale à $\varphi$ modulo
$(\mathfrak p_{(B \otimes_{C, \sigma} R^\wedge)^\wedge})^N$.
On peut donc remplacer $\varphi$ par cette composée, car la seule
façon dont $\varphi$ intervient dans la conclusion est la condition de
commutativité de la partie (5) de l'énoncé du lemme.
Diagramme :
$$
\xymatrix{
& & B' \ar[r] & (B')^\wedge \ar@{=}[d] \\
A \ar[rru] &
B \ar[r] &
B \otimes_{C, \sigma} R^\wedge \ar[r] \ar[u] &
(B \otimes_{C, \sigma} R^\wedge)^\wedge \\
R \ar[r] \ar[u] &
C \ar[r]^\sigma \ar[u] & R^\wedge \ar[u] \ar[ru]
}
$$
Ensuite, on utilise le fait que $R^\wedge$ est une colimite filtrante de
$R$-algèbres lisses (Lissage des morphismes d'anneaux, théorème \ref{smoothing-theorem-popescu}),
car $R_{\mathfrak p_R}$ est un G-anneau par hypothèse. Comme $C$ est de
présentation finie sur $R$, on obtient une factorisation
$$
C \to R' \to R^\wedge
$$
pour un certain homomorphisme lisse $R \to R'$ ; voir
Algèbre, lemme \ref{algebra-lemma-characterize-finite-presentation}.
Quitte à remplacer $R'$ par un terme ultérieur, on peut supposer qu'il existe
une algèbre \'etale sur $B \otimes_C R'$, notée $B''$, dont le changement de base
à $B \otimes_{C, \sigma} R^\wedge$ est $B'$ ; voir
Algèbre, lemme \ref{algebra-lemma-etale}.
Alors $B'$ est la colimite filtrante de ces $B''$, et l'on
conclut que, quitte à remplacer $R'$ par un terme ultérieur, on peut supposer qu'il existe
un homomorphisme de $R$-algèbres $A \to B''$ tel que $A \to B'' \to B'$ soit
l'homomorphisme construit précédemment (même référence que ci-dessus). Diagramme :
$$
\xymatrix{
& & B'' \ar[r] & B' \ar[r] & (B')^\wedge \ar@{=}[d] \\
A \ar[rru] &
B \ar[r] &
B \otimes_C R' \ar[r] \ar[u] &
B \otimes_{C, \sigma} R^\wedge \ar[r] \ar[u] &
(B \otimes_{C, \sigma} R^\wedge)^\wedge \\
R \ar[r] \ar[u] &
C \ar[r] \ar[u] &
R' \ar[r] \ar[u] &
R^\wedge \ar[u] \ar[ru]
}
$$
et
$$
B' = B'' \otimes_{(B \otimes_C R')} (B \otimes_{C, \sigma} R^\wedge)
$$
Ceci signifie que l'on peut remplacer $C$ par $R'$, $\sigma : C \to R^\wedge$ par
$R' \to R^\wedge$ et $B$ par $B''$, de sorte que l'on se ramène au diagramme
$$
\xymatrix{
A \ar[r] &
B \ar[r] &
B \otimes_{C, \sigma} R^\wedge \\
R \ar[r] \ar[u] &
C \ar[r]^\sigma \ar[u] & R^\wedge \ar[u]
}
$$
où $\varphi$ est égal à la composée des flèches horizontales
suivie de l'homomorphisme canonique de $B \otimes_{C, \sigma} R^\wedge$
vers son complété.
La dernière étape de la démonstration consiste à appliquer le lemme \ref{lemma-map-approximation}
(ou sa démonstration)
une fois de plus à $\Spec(C)$ et $\Spec(R)$ au-dessus de $\Spec(R)$ et à l'homomorphisme
$C \to R^\wedge$. Le lemme fournit un homomorphisme d'anneaux $C \to D$
tel que $R \to D$ soit \'etale, que
$\kappa(\mathfrak p_R) = \kappa(\mathfrak p_D)$, et que
$$
C \to D \to D^\wedge = R^\wedge
$$
coïncide avec $\sigma : C \to R^\wedge$ modulo $(\mathfrak p_{R^\wedge})^N$.
Nous pouvons alors prendre
$$
V = \Spec(D)
\quad\text{et}\quad
W = \Spec(B \otimes_C D)
$$
comme solution au problème posé par le lemme. En effet, le diagramme
$$
\xymatrix{
A \ar[r] &
B \otimes_{C, \sigma} R^\wedge \ar[r] &
B \otimes_{C, \sigma} R^\wedge/(\mathfrak p_{R^\wedge})^N \ar@{=}[r] &
B \otimes_C D/(\mathfrak p_D)^N \\
A \ar@{=}[u] \ar[r] &
A \otimes_R D \ar[r] &
B \otimes_R D \ar[r] &
B \otimes_C D/(\mathfrak p_D)^N \ar@{=}[u]
}
$$
est commutatif parce que la composée $C \to D \to D^\wedge = R^\wedge$ coïncide avec
$\sigma$ modulo $(\mathfrak p_{R^\wedge})^N$. Ceci démontre la partie (5),
et les autres propriétés résultent immédiatement de la construction.
\end{proof}

\begin{lemma}
\label{lemma-control-agreement}
Soit $T \to S$ un morphisme de type fini entre schémas noethériens.
Soit $t \in T$ d'image $s \in S$, et soit
$\sigma : \mathcal{O}_{T, t} \to \mathcal{O}_{S, s}^\wedge$
un morphisme local de $\mathcal{O}_{S, s}$-algèbres. Pour tout $N \geq 1$,
il existe un morphisme de type fini $(T', t') \to (T, t)$
tel que $\sigma$ se factorise par
$\mathcal{O}_{T, t} \to \mathcal{O}_{T', t'}$
et tel que, pour tout morphisme local de $\mathcal{O}_{S, s}$-algèbres
$\sigma' : \mathcal{O}_{T, t} \to \mathcal{O}_{S, s}^\wedge$
qui se factorise par $\mathcal{O}_{T, t} \to \mathcal{O}_{T', t'}$,
les morphismes $\sigma$ et $\sigma'$ coïncident modulo $\mathfrak m_s^N$.
\end{lemma}

\begin{proof}
Nous pouvons supposer $S$ et $T$ affines. Écrivons $S = \Spec(R)$ et
$T = \Spec(C)$. Soient $c_1, \ldots, c_n \in C$ des générateurs
de $C$ comme $R$-algèbre. Soit $\mathfrak p \subset R$ l'idéal
premier correspondant à $s$. Écrivons $\mathfrak p = (f_1, \ldots, f_m)$.
Après avoir remplacé $R$ par un localisé principal
(afin de chasser les dénominateurs dans $R_\mathfrak p$),
nous pouvons supposer qu'il existe
$r_1, \ldots, r_n \in R$ et $a_{i, I} \in \mathcal{O}_{S, s}^\wedge$,
où $I = (i_1, \ldots, i_m)$ avec $\sum i_j = N$, tels que
$$
\sigma(c_i) = r_i + \sum\nolimits_I a_{i, I} f_1^{i_1} \ldots f_m^{i_m}
$$
dans $\mathcal{O}_{S, s}^\wedge$. Considérons alors
$$
C' = C[t_{i, I}]/
\left(c_i - r_i - \sum\nolimits_I t_{i, I} f_1^{i_1} \ldots f_m^{i_m}\right)
$$
avec $\mathfrak p' = \mathfrak pC' + (t_{i, I})$, et la factorisation
de $\sigma : C \to \mathcal{O}_{S, s}^\wedge$ par $C'$
obtenue en envoyant $t_{i, I}$ sur $a_{i, I}$.
Le choix $T' = \Spec(C')$ convient, car tout $\sigma'$ comme dans l'énoncé
du lemme envoie $c_i$ sur $r_i$ modulo la
puissance $N$-ième de l'idéal maximal.
\end{proof}

\begin{lemma}
\label{lemma-control-graded}
Soient $Y \to T \to S$ des morphismes de type fini entre schémas noethériens.
Soit $t \in T$ d'image $s \in S$, et soit
$\sigma : \mathcal{O}_{T, t} \to \mathcal{O}_{S, s}^\wedge$
un morphisme local de $\mathcal{O}_{S, s}$-algèbres.
Il existe un morphisme de type fini $(T', t') \to (T, t)$
tel que $\sigma$ se factorise par
$\mathcal{O}_{T, t} \to \mathcal{O}_{T', t'}$
et tel que, pour tout morphisme local de $\mathcal{O}_{S, s}$-algèbres
$\sigma' : \mathcal{O}_{T, t} \to \mathcal{O}_{S, s}^\wedge$
qui se factorise par $\mathcal{O}_{T, t} \to \mathcal{O}_{T', t'}$,
les immersions fermées
$$
Y \times_{T, \sigma} \Spec(\mathcal{O}_{S, s}^\wedge) = Y_\sigma
\longleftarrow Y_t \longrightarrow
Y_{\sigma'} =
Y \times_{T, \sigma'} \Spec(\mathcal{O}_{S, s}^\wedge)
$$
ont des algèbres conormales isomorphes.
\end{lemma}

\begin{proof}
Une observation utile est que $\kappa(s) = \kappa(t)$ par l'existence de
$\sigma$. Remarquons que l'énoncé a un sens, puisque les fibres de $Y_\sigma$
et de $Y_{\sigma'}$ au-dessus de $s \in \Spec(\mathcal{O}_{S, s}^\wedge)$
sont toutes deux canoniquement isomorphes à $Y_t$. Nous considérerons la propriété
«~$\sigma'$ se factorise par $\mathcal{O}_{T, t} \to \mathcal{O}_{T', t'}$~»
comme une contrainte sur $\sigma'$. Si nous avons plusieurs telles contraintes,
disons données par $(T'_i, t'_i) \to (T, t)$, $i = 1, \ldots, n$,
nous pouvons les combiner en considérant
$(T'_1 \times_T \ldots \times_T T'_n, (t'_1, \ldots, t'_n)) \to
(T, t)$. Nous utiliserons ceci sans autre mention dans la suite.

\medskip\noindent
D'après le lemme \ref{lemma-control-agreement}, nous pouvons supposer que tout $\sigma'$
comme dans l'énoncé du lemme est égal à $\sigma$ modulo
$\mathfrak m_s^2$. Remarquons que l'algèbre conormale de $Y_t$ dans
$Y_\sigma$ n'est autre que la $\mathcal{O}_{Y_t}$-algèbre graduée quasi-cohérente
$$
\bigoplus\nolimits_{n \geq 0}
\mathfrak m_s^n\mathcal{O}_{Y_\sigma}/
\mathfrak m_s^{n + 1}\mathcal{O}_{Y_\sigma}
$$
et de même pour $Y_{\sigma'}$. Puisque $\sigma$ et $\sigma'$
coïncident modulo $\mathfrak m_s^2$, nous voyons que ces deux algèbres
sont les mêmes en degrés $0$ et $1$. D'autre part, ces
algèbres conormales sont engendrées en degré $1$ sur leur partie de degré $0$.
Par conséquent, s'il existe un isomorphisme prolongeant l'isomorphisme
qui vient d'être construit en degrés $0$ et $1$, alors il est unique.

\medskip\noindent
Nous pouvons supposer $S$ et $T$ affines. Soit $Y = Y_1 \cup \ldots \cup Y_n$
un recouvrement ouvert affine. Si nous pouvons construire $(T_i', t'_i) \to (T, t)$
comme dans le lemme de telle sorte que l'isomorphisme souhaité (voir le paragraphe précédent)
existe pour $Y_i \to T \to S$ et $\sigma$, alors ces isomorphismes se recollent par
unicité et donnent le résultat pour $Y \to T$. Nous pouvons donc supposer $Y$ affine.

\medskip\noindent
Écrivons $S = \Spec(R)$, $T = \Spec(C)$ et $Y = \Spec(B)$.
Choisissons une présentation $B = C[x_1, \ldots, x_n]/(f_1, \ldots, f_m)$.
Notons $R^\wedge = \mathcal{O}_{S, s}^\wedge$.
Soient $a_{kj} \in R^\wedge[x_1, \ldots, x_n]$ des polynômes
tels que
$$
\sum\nolimits_{j = 1, \ldots, m} a_{kj}\sigma(f_j) = 0,\quad
\text{pour }k = 1, \ldots, K
$$
forment un système de générateurs du module des relations entre
les $\sigma(f_j) \in R^\wedge[x_1, \ldots, x_n]$.
Nous avons donc une suite exacte
\begin{equation}
\label{equation-resolution}
R^\wedge[x_1, \ldots, x_n]^{\oplus K} \to
R^\wedge[x_1, \ldots, x_n]^{\oplus m} \to
R^\wedge[x_1, \ldots, x_n] \to B \otimes_{C, \sigma} R^\wedge \to 0
\end{equation}
Soit $c$ un entier pour lequel le lemme d'Artin--Rees s'applique
à la fois au premier et au second morphisme de cette suite, ainsi qu'à l'idéal
$\mathfrak m_{R^\wedge}R^\wedge[x_1, \ldots, x_n]$, comme dans
Compléments d'algèbre, section \ref{more-algebra-section-artin-rees}.
Écrivons
$$
a_{kj} = \sum\nolimits_{I \in \Omega} a_{kj, I} x^I
\quad\text{et}\quad
f_j = \sum\nolimits_{I \in \Omega} f_{j, I} x^I
$$
en notation multi-indicielle, où $a_{kj, I} \in R^\wedge$, $f_{j, I} \in C$,
et $\Omega$ est un ensemble fini de multi-indices. Nous voyons alors que
$$
\sum\nolimits_{j = 1, \ldots, m,\ I, I' \in \Omega,\ I + I' = I''}
a_{kj, I} \sigma(f_{j, I'}) = 0,\quad
I''\text{ un multi-indice}
$$
dans $R^\wedge$. Prenons donc
$$
C' = C[t_{jk, I}]/
\left(
\sum\nolimits_{j = 1, \ldots, m,\ I, I' \in \Omega,\ I + I' = I''}
t_{kj, I} f_{j, I'},\ I''\text{ un multi-indice}\right)
$$
Alors $\sigma$ se factorise par un morphisme $\tilde\sigma : C' \to R^\wedge$
envoyant $t_{kj, I}$ sur $a_{jk, I}$.
Ainsi, $T' = \Spec(C')$ possède un point $t' \in T'$ tel que
$\sigma$ se factorise par $\mathcal{O}_{T, t} \to \mathcal{O}_{T', t'}$.
Soit $t_{kj} = \sum t_{kj, I} x^I$ dans $C'[x_1, \ldots, x_n]$.
Nous voyons alors que nous avons un complexe
\begin{equation}
\label{equation-resolution-new}
C'[x_1, \ldots, x_n]^{\oplus K} \to
C'[x_1, \ldots, x_n]^{\oplus m} \to
C'[x_1, \ldots, x_n] \to B \otimes_C C' \to 0
\end{equation}
qui est exact en $C'[x_1, \ldots, x_n]$ et dont le changement de base
par $\tilde\sigma$ donne (\ref{equation-resolution}).

\medskip\noindent
D'après le lemme \ref{lemma-control-agreement},
nous pouvons trouver un morphisme supplémentaire $(T'', t'') \to (T', t')$
tel que $\tilde\sigma$
se factorise par $\mathcal{O}_{T', t'} \to \mathcal{O}_{T'', t''}$
et tel que, si $\sigma' : C \to R^\wedge$
se factorise par $\mathcal{O}_{T'', t''}$, alors le morphisme induit
$\tilde \sigma' : C' \to R^\wedge$
coïncide modulo $\mathfrak m_s^{c + 1}$ avec $\tilde \sigma$.
Ainsi, si $\sigma'$ est un tel morphisme, nous obtenons un complexe
$$
R^\wedge[x_1, \ldots, x_n]^{\oplus K} \to
R^\wedge[x_1, \ldots, x_n]^{\oplus m} \to
R^\wedge[x_1, \ldots, x_n] \to B \otimes_{C, \sigma'} R^\wedge \to 0
$$
sur $R^\wedge[x_1, \ldots, x_n]$ en appliquant $\tilde\sigma'$
aux polynômes $t_{kj}$ et $f_j$. Autrement dit, il
s'agit du changement de base du complexe (\ref{equation-resolution-new})
par $\tilde\sigma'$. Les matrices définissant ce complexe
sont congruentes modulo $\mathfrak m_s^{c + 1}$ aux matrices
définissant le complexe (\ref{equation-resolution}), car
$\tilde \sigma$ et $\tilde \sigma'$ sont congrus modulo
$\mathfrak m_s^{c + 1}$. Puisque (\ref{equation-resolution}) est exacte,
nous pouvons appliquer
Compléments d'algèbre, lemme \ref{more-algebra-lemma-approximate-complex-graded}
pour conclure que
$$
\text{Gr}_{\mathfrak m_s}(B \otimes_{C, \sigma'} R^\wedge)
\cong
\text{Gr}_{\mathfrak m_s}(B \otimes_{C, \sigma} R^\wedge)
$$
comme souhaité.
\end{proof}

\begin{lemma}
\label{lemma-relative-isomorphism-approximation}
Avec les notations et les hypothèses du lemme \ref{lemma-relative-map-approximation},
supposons que $\varphi$ induise un isomorphisme sur les complétés.
Nous pouvons alors choisir notre diagramme de telle sorte que $f$ soit \'etale.
\end{lemma}

\begin{proof}
Nous pouvons supposer $N \geq 2$ et remplacer $(T, t)$ par $(T', t')$ comme dans le
lemme \ref{lemma-control-graded}. Puisque $(V, v) \to (S, s)$ est un voisinage
\'etale élémentaire, il en va de même de $(X \times_S V, (x, v)) \to (X, x)$.
Ainsi, $\mathcal{O}_{X, x} \to \mathcal{O}_{X \times_S V, (x, v)}$
induit un isomorphisme sur les complétés d'après
Compléments d'algèbre, lemme \ref{more-algebra-lemma-flat-unramified}.
Nous affirmons que $\mathcal{O}_{X, x} \to \mathcal{O}_{W, w}$ induit
un isomorphisme sur les complétés. Une fois ceci démontré, le
lemme \ref{lemma-lifting-along-artinian-at-point}
montrera que $f$ est lisse en $w$ et, bien entendu,
$f$ est aussi non ramifié en $u$; par conséquent,
Morphismes, lemme \ref{morphisms-lemma-etale-smooth-unramified}
nous dit que $f$ est \'etale en $w$.

\medskip\noindent
Nous utilisons d'abord la commutativité de la partie (5) du
lemme \ref{lemma-relative-map-approximation}
pour voir que, pour $i \leq N$, il existe un diagramme commutatif
$$
\xymatrix{
\text{Gr}^i_{\mathfrak m_x}(\mathcal{O}_{X, x})
\ar[r]_-\varphi &
\text{Gr}^i_{\mathfrak m_{y_\sigma}}(\mathcal{O}_{Y_\sigma, y_\sigma}^\wedge)
\ar@{=}[r] &
\text{Gr}^i_{\mathfrak m_{(y, v)}}(\mathcal{O}_{Y \times_{T, \tau} V, (y, v)})
\ar[d]_{\cong} \\
\text{Gr}^i_{\mathfrak m_x}(\mathcal{O}_{X, x})
\ar[r]^-{\cong} \ar@{=}[u] &
\text{Gr}^i_{\mathfrak m_{(x, v)}}(\mathcal{O}_{X \times_S V, (x, v)})
\ar[r]^{f^\sharp_w} &
\text{Gr}^i_{\mathfrak m_w}(\mathcal{O}_{W, w})
}
$$
Ceci implique que $f^\sharp_w$ définit un isomorphisme
$\kappa(x) \to \kappa(w)$ sur les corps résiduels et un isomorphisme
$\mathfrak m_x/\mathfrak m_x^2 \to \mathfrak m_w/\mathfrak m_w^2$
sur les espaces cotangents. Par conséquent, $f^\sharp_w$ définit une surjection
$\mathcal{O}_{X, x}^\wedge \to \mathcal{O}_{W, w}^\wedge$
entre les anneaux locaux complets.

\medskip\noindent
D'après le lemme \ref{lemma-control-graded}, il existe un isomorphisme de
$\text{Gr}_{\mathfrak m_s}(\mathcal{O}_{(Y \times_{T, \tau} V, (y, v)})$
sur
$\text{Gr}_{\mathfrak m_s}(\mathcal{O}_{Y_\sigma, y_\sigma})$.
Ceci s'obtient en prenant les germes de l'isomorphisme des faisceaux conormaux
au point $y$. Puisque nos anneaux locaux sont noethériens,
la formation du gradué associé relativement à $\mathfrak m_s$
commute à la complétion, car la complétion relativement à un idéal
est un foncteur exact sur les modules finis sur les anneaux noethériens.
Ceci fournit l'isomorphisme vertical de droite dans le diagramme d'anneaux gradués
$$
\xymatrix{
\text{Gr}_{\mathfrak m_s}(\mathcal{O}_{W, w}^\wedge) &
\text{Gr}_{\mathfrak m_s}
(\mathcal{O}_{(Y \times_{T, \tau} V, (y, v)}^\wedge) \ar[l] \\
\text{Gr}_{\mathfrak m_s}(\mathcal{O}_{X, x}^\wedge)
\ar[r]^\varphi \ar[u] &
\text{Gr}_{\mathfrak m_s}(\mathcal{O}_{Y_\sigma, y_\sigma}^\wedge)
\ar[u]_{\cong}
}
$$
Nous n'affirmons pas que le diagramme soit commutatif. D'après le résultat du paragraphe
précédent, la flèche de gauche est surjective. Les trois autres flèches
sont des isomorphismes. Il s'ensuit que la flèche de gauche est un morphisme surjectif
entre des anneaux noethériens isomorphes. C'est donc un isomorphisme
d'après Algèbre, lemme \ref{algebra-lemma-surjective-endo-noetherian-ring-is-iso}
(on peut également le démontrer directement au moyen des fonctions de Hilbert).
En particulier, $\mathcal{O}_{X, x}^\wedge \to \mathcal{O}_{W, w}^\wedge$
doit être injectif aussi bien que surjectif, ce qui achève la démonstration.
\end{proof}















\section{Les morphismes finis localement libres dominent localement les morphismes \'etales}
\label{section-finite-free-over-etale}

\noindent
Dans cette section, nous expliquons un résultat qui énonce, en substance, que
les recouvrements \'etales d'un schéma $S$ peuvent être raffinés par des recouvrements
de Zariski de schémas finis localement libres sur $S$.

\begin{lemma}
\label{lemma-dominate-etale-neighbourhood-finite-flat}
Soit $S$ un schéma. Soit $s \in S$.
Soit $f : (U, u) \to (S, s)$ un voisinage \'etale.
Il existe un voisinage ouvert affine $s \in V \subset S$
et un morphisme surjectif fini localement libre $\pi : T \to V$
tels que, pour tout $t \in \pi^{-1}(s)$, il existe un
voisinage ouvert $t \in W_t \subset T$ et un
diagramme commutatif
$$
\xymatrix{
T \ar[d]^\pi & W_t \ar[l] \ar[rr]_{h_t} \ar[rd] & & U \ar[dl] \\
V \ar[rr] & & S
}
$$
avec $h_t(t) = u$.
\end{lemma}

\begin{proof}
Le problème est local sur $S$; nous pouvons donc remplacer $S$ par un
voisinage ouvert quelconque de $s$.
Nous pouvons aussi remplacer $U$ par un voisinage ouvert de $u$.
Ainsi, d'après Morphismes, lemme \ref{morphisms-lemma-etale-locally-standard-etale},
nous pouvons supposer que
$U \to S$ est un morphisme \'etale standard de schémas affines.
Dans ce cas, le lemme (avec $V = S$) résulte de
Algèbre, lemme \ref{algebra-lemma-standard-etale-finite-flat-Zariski}.
\end{proof}

\begin{lemma}
\label{lemma-dominate-etale-affine-finite-flat}
Soit $f : U \to S$ un morphisme \'etale surjectif de schémas affines.
Il existe un morphisme surjectif fini localement libre
$\pi : T \to S$ et un recouvrement ouvert fini
$T = T_1 \cup \ldots \cup T_n$ tels que chacun des morphismes
$T_i \to S$ se factorise par $U \to S$. Diagramme :
$$
\xymatrix{
& \coprod T_i  \ar[rd] \ar[ld] & \\
T \ar[rd]^\pi & & U \ar[ld]_f \\
& S &
}
$$
où la flèche sud-ouest est un recouvrement de Zariski.
\end{lemma}

\begin{proof}
C'est une reformulation de
Algèbre, lemme \ref{algebra-lemma-etale-finite-flat-zariski}.
\end{proof}

\begin{remark}
\label{remark-topologies}
En termes de topologies, les
lemmes \ref{lemma-dominate-etale-neighbourhood-finite-flat} et
\ref{lemma-dominate-etale-affine-finite-flat} signifient ce qui suit.
Soit $S$ un schéma quelconque. Soit $\{f_i : U_i \to S\}$ un recouvrement \'etale
de $S$. Il existe un recouvrement ouvert de Zariski $S = \bigcup V_j$,
pour tout $j$ un morphisme fini localement libre et surjectif
$W_j \to V_j$, et pour tout $j$ un recouvrement ouvert de Zariski
$\{W_{j, k} \to W_j\}$ tel que la famille
$\{W_{j, k} \to S\}$ raffine le recouvrement \'etale donné
$\{f_i : U_i \to S\}$. Que signifie ceci en pratique ?
Par exemple, supposons que nous ayons un problème de descente que nous
sachions résoudre pour les recouvrements de Zariski et pour les recouvrements fppf
de la forme $\{\pi : T \to S\}$, où $\pi$ est fini localement libre
et surjectif. Alors ce problème de descente admet également une réponse
affirmative pour les recouvrements \'etales. Cette astuce a été utilisée par
Gabber dans sa démonstration de $\text{Br}(X) = \text{Br}'(X)$
pour un schéma affine $X$; voir \cite{Hoobler}.
\end{remark}




\section{Localisation \'etale des morphismes quasi-finis}
\label{section-etale-localization}

\noindent
Nous abordons maintenant une série de lemmes autour du thème
«~les morphismes quasi-finis deviennent finis après localisation \'etale~».
L'idée générale est la suivante. Soient un morphisme
de schémas $f : X \to S$ et un point $s \in S$. Soit
$\varphi : (U, u) \to (S, s)$ un voisinage \'etale de $s$ dans $S$.
Considérons le produit fibré $X_U = U \times_S X$ et le
diagramme fondamental
\begin{equation}
\label{equation-basic-diagram}
\vcenter{
\xymatrix{
V \ar[r] \ar[dr] & X_U \ar[d] \ar[r] & X \ar[d]^f \\
& U \ar[r]^\varphi & S
}
}
\end{equation}
où $V \subset X_U$ est ouvert.
Existe-t-il un modèle standard pour le morphisme $f_U : X_U \to U$, ou pour
le morphisme $V \to U$ lorsque $V$ est un ouvert convenable ?
Bien entendu, la réponse est non en général. Mais, pour les morphismes quasi-finis,
nous pouvons dire quelque chose.

\begin{lemma}
\label{lemma-etale-makes-quasi-finite-finite-at-point}
Soit $f : X \to S$ un morphisme de schémas.
Soit $x \in X$. Posons $s = f(x)$.
Supposons que
\begin{enumerate}
\item $f$ soit localement de type fini, et
\item $x \in X_s$ soit isolé\footnote{En présence de (1),
cela signifie que $f$ est
quasi-fini en $x$; voir
Morphismes, lemme \ref{morphisms-lemma-quasi-finite-at-point-characterize}.}.
\end{enumerate}
Alors il existe
\begin{enumerate}
\item[(a)] un voisinage \'etale élémentaire $(U, u) \to (S, s)$,
\item[(b)] un sous-schéma ouvert $V \subset X_U$
(voir \ref{equation-basic-diagram})
\end{enumerate}
tels que
\begin{enumerate}
\item[(\romannumeral1)] $V \to U$ soit un morphisme fini,
\item[(\romannumeral2)] il existe un unique point $v$ de $V$
d'image $u$ dans $U$, et
\item[(\romannumeral3)] le point $v$ s'envoie sur $x$
par le morphisme $X_U \to X$, qui induit $\kappa(x) = \kappa(v)$.
\end{enumerate}
De plus, pour tout voisinage \'etale élémentaire $(U', u') \to (U, u)$,
si l'on pose $V' = U' \times_U V \subset X_{U'}$, le triplet $(U', u', V')$
satisfait encore les propriétés
(\romannumeral1), (\romannumeral2) et (\romannumeral3).
\end{lemma}

\begin{proof}
Soient $Y \subset X$, $W \subset S$ des ouverts affines tels que
$f(Y) \subset W$ et $x \in Y$. Remarquons que $x$ est
aussi un point isolé de la fibre du morphisme $f|_Y : Y \to W$.
Si nous pouvons démontrer le résultat pour $f|_Y : Y \to W$, alors le
résultat en découle pour $f$. Nous nous ramenons donc au cas où
$f$ est un morphisme de schémas affines. Ce cas est traité dans
Algèbre, lemme \ref{algebra-lemma-etale-makes-quasi-finite-finite-one-prime}.
\end{proof}

\noindent
Dans le lemme précédent et dans le suivant, nous ne supposons pas que le morphisme
$f$ soit séparé. Cela signifie que les ouverts $V$, $V_i$ qui y sont construits
ne sont pas nécessairement fermés dans $X_U$. De plus, si nous choisissons
le voisinage $U$ affine, alors chaque $V_i$ est affine, mais
les intersections $V_i \cap V_j$ ne sont pas nécessairement affines (dans le cas non
séparé).

\begin{lemma}
\label{lemma-etale-makes-quasi-finite-finite-multiple-points}
Soit $f : X \to S$ un morphisme de schémas.
Soient $x_1, \ldots, x_n \in X$ des points de même image $s$ dans $S$.
Supposons que
\begin{enumerate}
\item $f$ soit localement de type fini, et
\item $x_i \in X_s$ soit isolé pour $i = 1, \ldots, n$.
\end{enumerate}
Alors il existe
\begin{enumerate}
\item[(a)] un voisinage \'etale élémentaire $(U, u) \to (S, s)$,
\item[(b)] pour chaque $i$, un sous-schéma ouvert $V_i \subset X_U$,
\end{enumerate}
tels que, pour chaque $i$, les propriétés suivantes soient vérifiées :
\begin{enumerate}
\item[(\romannumeral1)] $V_i \to U$ est un morphisme fini,
\item[(\romannumeral2)] il existe un unique point $v_i$ de $V_i$
d'image $u$ dans $U$, et
\item[(\romannumeral3)] le point $v_i$ s'envoie sur $x_i$ dans $X$ et
$\kappa(x_i) = \kappa(v_i)$.
\end{enumerate}
\end{lemma}

\begin{proof}
Nous raisonnons par récurrence sur $n$.
Supposons en effet que $(U, u) \to (S, s)$ et $V_i \subset X_U$,
$i = 1, \ldots, n - 1$, conviennent pour $x_1, \ldots, x_{n - 1}$. Puisque
$\kappa(s) = \kappa(u)$, la fibre $(X_U)_u = X_s$. Il existe donc
un unique point $x'_n \in X_u \subset X_U$ correspondant à
$x_n \in X_s$. De plus, $x'_n$ est isolé dans $X_u$. Ainsi, d'après le
lemme \ref{lemma-etale-makes-quasi-finite-finite-at-point}, il
existe un voisinage \'etale élémentaire $(U', u') \to (U, u)$
et un ouvert $V_n \subset X_{U'}$ qui convient pour $x'_n$, et donc
pour $x_n$.
D'après la dernière assertion du
lemme \ref{lemma-etale-makes-quasi-finite-finite-at-point},
les sous-schémas ouverts $V'_i = U'\times_U V_i$, pour $i = 1, \ldots, n - 1$, conviennent encore
relativement à $x_1, \ldots, x_{n - 1}$. D'où le résultat.
\end{proof}

\noindent
Si nous autorisons une extension de corps non triviale $\kappa(u)/\kappa(s)$, c'est-à-dire
des voisinages \'etales généraux, nous pouvons scinder les points comme suit.

\begin{lemma}
\label{lemma-etale-makes-quasi-finite-finite-multiple-points-var}
Soit $f : X \to S$ un morphisme de schémas.
Soient $x_1, \ldots, x_n \in X$ des points de même image $s$ dans $S$.
Supposons que
\begin{enumerate}
\item $f$ soit localement de type fini, et
\item $x_i \in X_s$ soit isolé pour $i = 1, \ldots, n$.
\end{enumerate}
Alors il existe
\begin{enumerate}
\item[(a)] un voisinage \'etale $(U, u) \to (S, s)$,
\item[(b)] pour chaque $i$, un entier $m_i$ et des
sous-schémas ouverts $V_{i, j} \subset X_U$, $j = 1, \ldots, m_i$,
\end{enumerate}
tels que l'on ait
\begin{enumerate}
\item[(\romannumeral1)] chaque $V_{i, j} \to U$ soit un morphisme fini,
\item[(\romannumeral2)] il existe un unique point $v_{i, j}$ de $V_{i, j}$
d'image $u$ dans $U$, tel que $\kappa(u) \subset \kappa(v_{i, j})$ soit
une extension finie purement inséparable,
\item[(\romannumeral4)] si $v_{i, j} = v_{i', j'}$, alors $i = i'$ et
$j = j'$, et
\item[(\romannumeral3)] les points $v_{i, j}$ s'envoient sur $x_i$ dans $X$ et
aucun autre point de $(X_U)_u$ ne s'envoie sur $x_i$.
\end{enumerate}
\end{lemma}

\begin{proof}
Cette démonstration est une variante de la démonstration de
Algèbre, lemme \ref{algebra-lemma-etale-makes-quasi-finite-finite-variant},
formulée dans le langage des schémas.
D'après Morphismes, lemme \ref{morphisms-lemma-quasi-finite-at-point-characterize},
le morphisme $f$ est quasi-fini en chacun des points $x_i$.
Ainsi, l'extension $\kappa(s) \subset \kappa(x_i)$ est finie pour tout $i$
(Morphismes, lemme \ref{morphisms-lemma-residue-field-quasi-finite}).
Pour chaque $i$, soit $\kappa(s) \subset L_i \subset \kappa(x_i)$
le sous-corps tel que $L_i/\kappa(s)$ soit séparable et que
$\kappa(x_i)/L_i$ soit purement inséparable. Choisissons une extension finie
galoisienne $L/\kappa(s)$ telle qu'il existe des
$\kappa(s)$-plongements $L_i \to L$ pour $i = 1, \ldots, n$.
Choisissons un voisinage \'etale $(U, u) \to (S, s)$ tel que
$L \cong \kappa(u)$ comme extensions de $\kappa(s)$
(lemme \ref{lemma-realize-prescribed-residue-field-extension-etale}).

\medskip\noindent
Soient $y_{i, j}$, $j = 1, \ldots, m_i$, les points de $X_U$
au-dessus de $x_i \in X$ et de $u \in U$. D'après
Schémas, lemme \ref{schemes-lemma-points-fibre-product},
ces points $y_{i, j}$ correspondent exactement aux idéaux premiers des anneaux
$\kappa(u) \otimes_{\kappa(s)} \kappa(x_i)$. Ceci explique également
pourquoi ils sont en nombre fini; en fait,
$m_i = [L_i : \kappa(s)]$, mais nous n'en avons pas besoin.
Par notre choix de
$L$ (et par la théorie élémentaire des corps),
nous voyons que l'extension $\kappa(u) \subset \kappa(y_{i, j})$ est
finie purement inséparable pour chaque paire $i, j$.
De plus, d'après Morphismes, lemme \ref{morphisms-lemma-base-change-quasi-finite},
par exemple, le morphisme
$X_U \to U$ est quasi-fini aux points $y_{i, j}$ pour
tous $i, j$.

\medskip\noindent
Appliquons le lemme \ref{lemma-etale-makes-quasi-finite-finite-multiple-points}
au morphisme $X_U \to U$, au point $u \in U$
et aux points $y_{i, j} \in (X_U)_u$. Nous obtenons un voisinage \'etale
$(U', u') \to (U, u)$ avec $\kappa(u) = \kappa(u')$ et des
ouverts $V_{i, j} \subset X_{U'}$ possédant les propriétés
(\romannumeral1), (\romannumeral2) et (\romannumeral3)
de ce lemme. Nous affirmons que le voisinage \'etale
$(U', u') \to (S, s)$ et les ouverts $V_{i, j} \subset X_{U'}$
constituent une solution au problème posé par le lemme.
Nous omettons les vérifications.
\end{proof}

\begin{lemma}
\label{lemma-etale-splits-off-quasi-finite-part-technical}
Soit $f : X \to S$ un morphisme de schémas.
Soit $s \in S$. Soient $x_1, \ldots, x_n \in X_s$. Supposons que
\begin{enumerate}
\item $f$ soit localement de type fini,
\item $f$ soit séparé, et
\item $x_1, \ldots, x_n$ soient des points de $X_s$, isolés et deux à deux distincts.
\end{enumerate}
Il existe alors un voisinage \'etale élémentaire $(U, u) \to (S, s)$
et une décomposition
$$
U \times_S X = W \amalg V_1 \amalg \ldots \amalg V_n
$$
en sous-schémas ouverts et fermés telle que les morphismes
$V_i \to U$ soient finis, que les fibres de $V_i \to U$ au-dessus de $u$ soient
les singletons $\{v_i\}$, que chaque $v_i$ s'envoie sur $x_i$ avec
$\kappa(x_i) = \kappa(v_i)$, et que la fibre de $W \to U$
au-dessus de $u$ ne contienne aucun point s'envoyant sur l'un des $x_i$.
\end{lemma}

\begin{proof}
Choisissons $(U, u) \to (S, s)$ et $V_i \subset X_U$ comme dans le
lemme \ref{lemma-etale-makes-quasi-finite-finite-multiple-points}.
Puisque $X_U \to U$ est séparé
(Schémas, lemme \ref{schemes-lemma-separated-permanence})
et que $V_i \to U$ est fini, donc propre
(Morphismes, lemme \ref{morphisms-lemma-finite-proper}),
on voit que $V_i \subset X_U$ est fermé d'après
Morphismes, lemme \ref{morphisms-lemma-image-proper-scheme-closed}.
Ainsi, $V_i \cap V_j$ est un sous-ensemble fermé de $V_i$ qui
ne contient pas $v_i$. Par conséquent, l'image de $V_i \cap V_j$
dans $U$ est un fermé (car $V_i \to U$ est propre) qui ne
contient pas $u$. Après avoir rétréci $U$, on peut donc supposer
que $V_i \cap V_j = \emptyset$ pour tous $i, j$. Cela donne la
décomposition de l'énoncé.
\end{proof}

\noindent
Voici la variante où l'on se ramène à des extensions de corps
purement inséparables.

\begin{lemma}
\label{lemma-etale-splits-off-quasi-finite-part-technical-variant}
Soit $f : X \to S$ un morphisme de schémas.
Soient $s \in S$ et $x_1, \ldots, x_n \in X_s$. Supposons que
\begin{enumerate}
\item $f$ soit localement de type fini,
\item $f$ soit séparé, et
\item $x_1, \ldots, x_n$ soient des points isolés deux à deux distincts de $X_s$.
\end{enumerate}
Il existe alors un voisinage \'etale $(U, u) \to (S, s)$
et une décomposition
$$
U \times_S X =
W \amalg
\ \coprod\nolimits_{i = 1, \ldots, n}
\ \coprod\nolimits_{j = 1, \ldots, m_i}
V_{i, j}
$$
en sous-schémas ouverts et fermés telle que les morphismes
$V_{i, j} \to U$ soient finis, que les fibres de $V_{i, j} \to U$ au-dessus de $u$ soient
les singletons $\{v_{i, j}\}$, que chaque $v_{i, j}$ s'envoie sur $x_i$,
que $\kappa(u) \subset \kappa(v_{i, j})$ soit purement inséparable,
et que la fibre de $W \to U$ au-dessus de $u$ ne contienne aucun point s'envoyant
sur l'un des $x_i$.
\end{lemma}

\begin{proof}
On procède exactement comme dans la démonstration du
lemme \ref{lemma-etale-splits-off-quasi-finite-part-technical}, si ce n'est qu'on
utilise le lemme \ref{lemma-etale-makes-quasi-finite-finite-multiple-points-var}
au lieu du lemme \ref{lemma-etale-makes-quasi-finite-finite-multiple-points}.
\end{proof}

\noindent
La version suivante est peut-être un peu plus facile à lire.

\begin{lemma}
\label{lemma-etale-splits-off-quasi-finite-part}
Soit $f : X \to S$ un morphisme de schémas.
Soit $s \in S$. Supposons que
\begin{enumerate}
\item $f$ soit localement de type fini,
\item $f$ soit séparé, et
\item $X_s$ ait au plus un nombre fini de points isolés.
\end{enumerate}
Il existe alors un voisinage \'etale élémentaire $(U, u) \to (S, s)$
et une décomposition
$$
U \times_S X = W \amalg V
$$
en sous-schémas ouverts et fermés telle que le morphisme
$V \to U$ soit fini et que la fibre $W_u$ du
morphisme $W \to U$ ne contienne aucun point isolé.
En particulier, si $f^{-1}(s)$ est un ensemble fini, alors $W_u = \emptyset$.
\end{lemma}

\begin{proof}
Cela résulte immédiatement du
lemme \ref{lemma-etale-splits-off-quasi-finite-part-technical},
en choisissant pour $x_1, \ldots, x_n$ l'ensemble de tous les
points isolés de $X_s$ et en posant $V = \bigcup V_i$.
\end{proof}






\section{Localisation \'etale des morphismes entiers}
\label{section-etale-localization-integral}

\noindent
Voici quelques variantes des résultats de la section \ref{section-etale-localization}
dans le cas des morphismes entiers.

\begin{lemma}
\label{lemma-etale-makes-integral-split}
Soit $R \to S$ un homomorphisme d'anneaux entier. Soit $\mathfrak p \subset R$ un idéal
premier. Supposons
\begin{enumerate}
\item qu'il n'y ait qu'un nombre fini d'idéaux premiers $\mathfrak q_1, \ldots, \mathfrak q_n$
au-dessus de $\mathfrak p$, et
\item que, pour tout $i$, la sous-extension séparable maximale
$\kappa(\mathfrak q)/\kappa(\mathfrak q_i)_{sep}/\kappa(\mathfrak p)$
(Corps, lemme \ref{fields-lemma-separable-first})
soit finie sur $\kappa(\mathfrak p)$.
\end{enumerate}
Il existe alors un homomorphisme d'anneaux \'etale $R \to R'$ et un idéal premier
$\mathfrak p'$ au-dessus de $\mathfrak p$ tels que
$$
S \otimes_R R' = A_1 \times \ldots \times A_m
$$
où $R' \to A_j$ est entier et possède un unique idéal premier $\mathfrak r_j$
au-dessus de $\mathfrak p'$ tel que $\kappa(\mathfrak r_j)/\kappa(\mathfrak p')$
soit purement inséparable.
\end{lemma}

\begin{proof}[Première démonstration]
Cette démonstration utilise
Algèbre, lemme \ref{algebra-lemma-etale-makes-quasi-finite-finite-variant}.
Choisissons en effet un générateur $\theta_i \in \kappa(\mathfrak q_i)_{sep}$
de ce corps sur $\kappa(\mathfrak p)$
(Corps, lemme \ref{fields-lemma-primitive-element}).
Le spectre de l'anneau fibre $S \otimes_R \kappa(\mathfrak p)$
est fini et discret, ses points correspondant à
$\mathfrak q_1, \ldots, \mathfrak q_n$.
D'après le théorème chinois des restes
(Algèbre, lemme \ref{algebra-lemma-chinese-remainder}),
on voit que $S \otimes_R \kappa(\mathfrak p) \to \prod \kappa(\mathfrak q_i)$
est surjectif. Par conséquent, après avoir remplacé $R$ par $R_g$ pour un certain
$g \in R$, $g \not \in \mathfrak p$, on peut supposer que
$(0, \ldots, 0, \theta_i, 0, \ldots, 0) \in \prod \kappa(\mathfrak q_i)$
est l'image d'un élément $x_i \in S$.
Soit $S' \subset S$ la $R$-sous-algèbre engendrée par les éléments $x_i$.
Puisque $\Spec(S) \to \Spec(S')$ est surjectif
(Algèbre, lemme \ref{algebra-lemma-integral-overring-surjective}),
on conclut que
$\mathfrak q_i' = S' \cap \mathfrak q_i$ sont les idéaux premiers de
$S'$ au-dessus de $\mathfrak p$. Le choix des $x_i$ montre que
ces idéaux premiers sont distincts et que
$\kappa(\mathfrak q'_i)_{sep} = \kappa(\mathfrak q_i)_{sep}$.
En particulier, les extensions de corps
$\kappa(\mathfrak q_i)/\kappa(\mathfrak q'_i)$ sont purement
inséparables.
Puisque $R \to S'$ est fini, on peut appliquer
Algèbre, lemme \ref{algebra-lemma-etale-makes-quasi-finite-finite-variant},
et l'on obtient $R \to R'$ et $\mathfrak p'$, ainsi qu'une décomposition
$$
S' \otimes_R R' = A'_1 \times \ldots \times A'_m \times B'
$$
où $R' \to A'_j$ est entier et possède un unique idéal premier
$\mathfrak r'_j$ au-dessus de $\mathfrak p'$ tel que
$\kappa(\mathfrak r'_j)/\kappa(\mathfrak p')$
soit purement inséparable, et où $B'$ ne possède aucun idéal premier
au-dessus de $\mathfrak p'$. Comme $R' \to B'$ est fini
(car $R \to S'$ est fini), on peut,
après avoir localisé $R'$ en un certain $g' \in R'$, $g' \not \in \mathfrak p'$,
supposer que $B' = 0$. Le morphisme $S' \otimes_R R' \to S \otimes_R R'$
nous fournit la décomposition correspondante pour $S$.
\end{proof}

\begin{proof}[Seconde démonstration]
Cette démonstration utilise l'hensélisation stricte. Supposons d'abord que $R$ soit strictement
hensélien, d'idéal maximal $\mathfrak p$. Alors $S/\mathfrak p S$
possède un nombre fini d'idéaux premiers correspondant à
$\mathfrak q_1, \ldots, \mathfrak q_n$, tous maximaux,
et dont les corps résiduels sont purement inséparables sur $\kappa(\mathfrak p)$.
Ainsi, $S/\mathfrak p S$ est égal à $\prod (S/\mathfrak p S)_{\mathfrak p_i}$.
D'après Plus sur l'algèbre, lemme \ref{more-algebra-lemma-characterize-henselian-pair},
on peut relever cette décomposition en produit en une décomposition en produit de $S$
comme dans l'énoncé.

\medskip\noindent
Dans le cas général, soit $R^{sh}$ l'hensélisé strict de
$R_\mathfrak p$. On peut alors appliquer le résultat du premier paragraphe
à $R^{sh} \to S \otimes_R R^{sh}$. Considérons les $m$ idempotents deux à deux orthogonaux
de $S \otimes_R R^{sh}$ correspondant à la décomposition en produit.
Puisque $R^{sh}$ est une colimite filtrante de morphismes d'anneaux \'etales
$(R, \mathfrak p) \to (R', \mathfrak p')$ d'après
Algèbre, lemme \ref{algebra-lemma-strict-henselization-different},
on voit que ces idempotents sont définis sur un certain $R'$, comme voulu.
\end{proof}







\section{Théorème principal de Zariski}
\label{section-application-etale-neighbourhoods}

\noindent
Dans cette section, on démontre le théorème principal de Zariski sous la forme donnée par Grothendieck.
Souvent, lorsque l'on dit « théorème principal de Zariski » dans ce contexte, on entend l'un des
lemmes \ref{lemma-finite-type-separated},
\ref{lemma-quasi-finite-separated-quasi-affine} ou
\ref{lemma-quasi-finite-separated-pass-through-finite}.
Dans la plupart des textes, c'est au dernier de ces énoncés que l'on donne le nom de
théorème principal de Zariski.

\medskip\noindent
On a déjà démontré la version algébrique dans
Algèbre, théorème \ref{algebra-theorem-main-theorem},
et on a déjà reformulé cette version algébrique
dans le langage des schémas ; voir
Morphismes, théorème \ref{morphisms-theorem-main-theorem}.
La version de cette section est plus subtile ; pour obtenir le résultat
complet, on utilise les techniques de localisation \'etale
de la section \ref{section-etale-localization} afin de se ramener au
cas algébrique.

\begin{lemma}
\label{lemma-finite-type-separated}
Soit $f : X \to S$ un morphisme de schémas.
Supposons que $f$ soit de type fini et séparé.
Soit $S'$ la normalisation de $S$ dans $X$ ; voir
Morphismes, définition \ref{morphisms-definition-normalization-X-in-Y}.
Diagramme :
$$
\xymatrix{
X \ar[rd]_f \ar[rr]_{f'} & & S' \ar[ld]^\nu \\
& S &
}
$$
Il existe alors un sous-schéma ouvert $U' \subset S'$ tel que
\begin{enumerate}
\item $(f')^{-1}(U') \to U'$ soit un isomorphisme, et
\item $(f')^{-1}(U') \subset X$ soit l'ensemble des points où
$f$ est quasi-fini.
\end{enumerate}
\end{lemma}

\begin{proof}
D'après Morphismes, lemme \ref{morphisms-lemma-quasi-finite-points-open},
le sous-ensemble $U \subset X$ des points où $f$ est quasi-fini est ouvert.
Le lemme équivaut aux trois assertions suivantes :
\begin{enumerate}
\item[(a)] $U' = f'(U) \subset S'$ est ouvert,
\item[(b)] $U = (f')^{-1}(U')$, et
\item[(c)] $U \to U'$ est un isomorphisme.
\end{enumerate}
Soit $x \in U$ arbitraire. On affirme qu'il existe un voisinage ouvert
$f'(x) \in V \subset S'$ tel que $(f')^{-1}V \to V$ soit un
isomorphisme. Montrons d'abord que cette assertion implique le lemme.
En effet, $(f')^{-1}V \cong V$ est alors à la fois localement de type
fini sur $S$ (comme sous-schéma ouvert de $X$), et, pour $v \in V$, l'extension de corps
résiduels $\kappa(v)/\kappa(\nu(v))$ est algébrique (car
$V \subset S'$ et $S'$ est entier sur $S$). Les fibres
de $V \to S$ sont donc discrètes (Morphismes, lemme
\ref{morphisms-lemma-algebraic-residue-field-extension-closed-point-fibre})
et $(f')^{-1}V \to S$ est localement quasi-fini
(Morphismes, lemme \ref{morphisms-lemma-locally-quasi-finite-fibres}).
Il s'ensuit que $(f')^{-1}V \subset U$ et $V \subset U'$. Puisque $x$ était
arbitraire, les assertions (a), (b) et (c) sont vraies.

\medskip\noindent
Posons $s = f(x)$. Soit $(T, t) \to (S, s)$ un voisinage \'etale
élémentaire. On note par un indice ${}_T$ le changement de base à $T$.
Soit $y = (x, t) \in X_T$ l'unique point de
la fibre $X_t$ situé au-dessus de $x$. Remarquons que $U_T \subset X_T$
est l'ensemble des points où $f_T$ est quasi-fini ; voir
Morphismes, lemme \ref{morphisms-lemma-base-change-quasi-finite}.
Remarquons que
$$
X_T \xrightarrow{f'_T} S'_T \xrightarrow{\nu_T} T
$$
est la normalisation de $T$ dans $X_T$ ; voir
lemme \ref{lemma-normalization-smooth-localization}.
Supposons que l'assertion soit vraie pour $y \in U_T \subset X_T \to S'_T \to T$, c'est-à-dire
que l'on puisse trouver un voisinage ouvert
$f'_T(y) \in V' \subset S'_T$ tel que $(f'_T)^{-1}V' \to V'$ soit un
isomorphisme. Le morphisme $S'_T \to S'$ est \'etale ; l'image
$V \subset S'$ de $V'$ est donc ouverte. On a $f'(x) \in V$ puisque $f'_T(y) \in V'$.
Observons que
$$
\xymatrix{
(f'_T)^{-1}V' \ar[r] \ar[d] & (f')^{-1}(V) \ar[d] \\
V' \ar[r] & V
}
$$
est un carré cartésien (car $S'_T \times_{S'} X = X_T$).
Puisque la flèche verticale de gauche est un isomorphisme
et que $\{V' \to V\}$ est un recouvrement \'etale, on conclut que la flèche verticale de droite
est un isomorphisme d'après
Descente, lemme \ref{descent-lemma-descending-property-isomorphism}.
Autrement dit, l'assertion est vraie pour $x \in U \subset X \to S' \to S$.

\medskip\noindent
D'après le paragraphe précédent, on peut remplacer $S$ par un
voisinage \'etale élémentaire de $s = f(x)$ pour démontrer l'assertion.
On peut donc supposer qu'il existe une décomposition
$$
X = V \amalg W
$$
en sous-schémas ouverts et fermés, où $V \to S$ est fini et $x \in V$ ;
voir le lemme \ref{lemma-etale-splits-off-quasi-finite-part-technical}.
Puisque $X$ est la réunion disjointe de $V$ et $W$ au-dessus de $S$ et que
$V \to S$ est fini, on voit que la
normalisation de $S$ dans $X$ est le morphisme
$$
X = V \amalg W \longrightarrow V \amalg W' \longrightarrow S
$$
où $W'$ est la normalisation de $S$ dans $W$ ; voir
Morphismes, lemmes \ref{morphisms-lemma-normalization-in-disjoint-union},
\ref{morphisms-lemma-finite-integral} et
\ref{morphisms-lemma-normalization-in-integral}.
L'assertion s'ensuit, ce qui achève la démonstration.
\end{proof}

\begin{lemma}
\label{lemma-quasi-finite-separated-quasi-affine}
\begin{slogan}
Les morphismes quasi-finis et séparés sont quasi-affines
\end{slogan}
Soit $f : X \to S$ un morphisme de schémas.
Supposons que $f$ soit quasi-fini et séparé.
Soit $S'$ la normalisation de $S$ dans $X$ ; voir
Morphismes, définition \ref{morphisms-definition-normalization-X-in-Y}.
Diagramme :
$$
\xymatrix{
X \ar[rd]_f \ar[rr]_{f'} & & S' \ar[ld]^\nu \\
& S &
}
$$
Alors $f'$ est une immersion ouverte quasi-compacte et $\nu$ est entier.
En particulier, $f$ est quasi-affine.
\end{lemma}

\begin{proof}
Cela résulte du lemme \ref{lemma-finite-type-separated}. En effet, d'après
ce lemme, il existe un sous-schéma ouvert $U' \subset S'$ tel que
$(f')^{-1}(U') = X$ et que $X \to U'$ soit un isomorphisme. Autrement
dit, $f'$ est une immersion ouverte. Remarquons que $f'$ est quasi-compact puisque
$f$ est quasi-compact et que $\nu : S' \to S$ est séparé
(Schémas, lemme \ref{schemes-lemma-quasi-compact-permanence}).
Il s'ensuit que $f$ est quasi-affine d'après
Morphismes, lemme \ref{morphisms-lemma-characterize-quasi-affine}.
\end{proof}

\begin{lemma}[Théorème principal de Zariski]
\label{lemma-quasi-finite-separated-pass-through-finite}
\begin{reference}
\cite[IV Corollaire 18.12.13]{EGA}
\end{reference}
Soit $f : X \to S$ un morphisme de schémas.
Supposons que $f$ soit quasi-fini et séparé, et que
$S$ soit quasi-compact et quasi-séparé. Il existe alors
une factorisation
$$
\xymatrix{
X \ar[rd]_f \ar[rr]_j & & T \ar[ld]^\pi \\
& S &
}
$$
où $j$ est une immersion ouverte quasi-compacte et $\pi$ est fini.
\end{lemma}

\begin{proof}
Soit $X \to S' \to S$ la factorisation fournie par
le lemme \ref{lemma-quasi-finite-separated-quasi-affine}.
D'après
Propriétés, lemme
\ref{properties-lemma-integral-algebra-directed-colimit-finite},
on peut écrire
$\nu_*\mathcal{O}_{S'} = \colim_{i \in I} \mathcal{A}_i$ comme une
colimite filtrante de $\mathcal{O}_S$-algèbres quasi-cohérentes finies
$\mathcal{A}_i \subset \nu_*\mathcal{O}_{S'}$. Alors
$\pi_i : T_i = \underline{\Spec}_S(\mathcal{A}_i) \to S$
est un morphisme fini pour tout $i$.
Remarquons que les morphismes de transition $T_{i'} \to T_i$ sont affines
et que $S' = \lim T_i$.

\medskip\noindent
D'après Limites, lemme \ref{limits-lemma-descend-opens},
il existe un $i$ et un ouvert quasi-compact
$U_i \subset T_i$ dont l'image réciproque dans $S'$ est égale à
$f'(X)$. Pour $i' \geq i$, soit $U_{i'}$ l'image réciproque
de $U_i$ dans $T_{i'}$. Alors $X \cong f'(X) = \lim_{i' \geq i} U_{i'}$ ; voir
Limites, lemme \ref{limits-lemma-directed-inverse-system-has-limit}.
D'après Limites, lemme \ref{limits-lemma-finite-type-eventually-closed},
on voit que $X \to U_{i'}$ est une immersion fermée pour un certain $i' \geq i$.
(En fait, $X \cong U_{i'}$ pour tout indice suffisamment
grand $i'$, mais on n'en a pas besoin.) Ainsi, $X \to T_{i'}$ est une immersion. D'après
Morphismes, lemme \ref{morphisms-lemma-factor-quasi-compact-immersion},
on peut la factoriser en $X \to T \to T_{i'}$, où la première flèche
est une immersion ouverte et la seconde une immersion fermée. Cela achève la démonstration.
\end{proof}

\begin{lemma}
\label{lemma-quasi-finite-separated-pass-through-finite-addendum}
Avec les notations et sous les hypothèses du
lemme \ref{lemma-quasi-finite-separated-pass-through-finite},
supposons de plus que $f$ soit localement de présentation finie. On peut alors
choisir la factorisation de sorte que $T$ soit fini et de
présentation finie sur $S$.
\end{lemma}

\begin{proof}
D'après Limites, lemme
\ref{limits-lemma-finite-in-finite-and-finite-presentation}, on peut écrire
$T = \lim T_i$, où tous les $T_i$ sont finis et de présentation finie
sur $S$ et où les morphismes de transition $T_{i'} \to T_i$ sont des immersions
fermées. D'après
Limites, lemme \ref{limits-lemma-descend-opens},
il existe un $i$ et un sous-schéma ouvert $U_i \subset T_i$ dont l'image
réciproque dans $T$ est $X$. D'après
Limites, lemme
\ref{limits-lemma-finite-type-eventually-closed},
on voit que $X \cong U_i$ pour $i$ suffisamment grand.
Remplacer $T$ par $T_i$ achève la démonstration.
\end{proof}








\section{Applications du théorème principal de Zariski, I}
\label{section-applications-zmt}

\noindent
Une première application est la caractérisation des morphismes finis
comme les morphismes propres à fibres finies.

\begin{lemma}
\label{lemma-characterize-finite}
Soit $f : X \to S$ un morphisme de schémas.
Les assertions suivantes sont équivalentes :
\begin{enumerate}
\item $f$ est fini,
\item $f$ est propre à fibres finies,
\item $f$ est propre et localement quasi-fini,
\item $f$ est universellement fermé, séparé, localement de type fini
et à fibres finies.
\end{enumerate}
\end{lemma}

\begin{proof}
(1) implique (2) d'après
Morphismes, lemmes \ref{morphisms-lemma-finite-proper},
\ref{morphisms-lemma-quasi-finite}
et \ref{morphisms-lemma-finite-quasi-finite}.
(2) implique (3) d'après Morphismes, lemme \ref{morphisms-lemma-finite-fibre}.
(3) implique (4) d'après la définition des morphismes propres et
Morphismes, lemmes \ref{morphisms-lemma-quasi-finite-locally-quasi-compact} et
\ref{morphisms-lemma-quasi-finite}.

\medskip\noindent
Supposons (4). Choisissons $s \in S$. D'après
Morphismes, lemme \ref{morphisms-lemma-finite-fibre}, on
voit que les points, en nombre fini, de $X_s$ sont tous isolés dans $X_s$.
Choisissons un voisinage \'etale élémentaire $(U, u) \to (S, s)$
et une décomposition $X_U = V \amalg W$ comme dans le
lemme \ref{lemma-etale-splits-off-quasi-finite-part}.
Remarquons que $W_u = \emptyset$ car tous les points de $X_s$ sont isolés.
Puisque $f$ est universellement fermé, on voit que
l'image de $W$ dans $U$ est un fermé qui ne contient pas $u$.
Après avoir rétréci $U$, on peut supposer que $W = \emptyset$.
Autrement dit, on voit que $X_U = V$ est fini sur $U$.
Comme $s \in S$ était arbitraire,
il existe une famille $\{U_i \to S\}$
de morphismes \'etales dont les images recouvrent $S$ et telle que
les changements de base $X_{U_i} \to U_i$ soient finis.
Remarquons que $\{U_i \to S\}$ est un recouvrement \'etale ;
voir Topologies, définition \ref{topologies-definition-etale-covering}.
C'est donc un recouvrement fpqc ; voir
Topologies,
lemme \ref{topologies-lemma-zariski-etale-smooth-syntomic-fppf-fpqc}.
On conclut alors que $f$ est fini d'après
Descente, lemme \ref{descent-lemma-descending-property-finite}.
\end{proof}

\noindent
On en déduit les résultats utiles suivants.

\begin{lemma}
\label{lemma-proper-finite-fibre-finite-in-neighbourhood}
Soit $f : X \to S$ un morphisme de schémas.
Soit $s \in S$.
Supposons que $f$ soit propre et que $f^{-1}(\{s\})$ soit un ensemble fini.
Il existe alors un voisinage ouvert $V \subset S$ de $s$
tel que $f|_{f^{-1}(V)} : f^{-1}(V) \to V$ soit fini.
\end{lemma}

\begin{proof}
Le morphisme $f$ est quasi-fini en tous les points de $f^{-1}(\{s\})$
d'après Morphismes, lemme \ref{morphisms-lemma-finite-fibre}.
D'après Morphismes, lemme \ref{morphisms-lemma-quasi-finite-points-open}, l'ensemble
des points où $f$ est quasi-fini est un ouvert $U \subset X$.
Posons $Z = X \setminus U$. Alors $s \not \in f(Z)$. Puisque $f$ est propre,
l'ensemble $f(Z) \subset S$ est fermé. Choisissons un voisinage ouvert quelconque
$V \subset S$ de $s$ tel que $f(Z) \cap V = \emptyset$. Alors
$f^{-1}(V) \to V$ est localement quasi-fini et propre.
Il est donc quasi-fini
(Morphismes, lemme \ref{morphisms-lemma-quasi-finite-locally-quasi-compact}),
donc à fibres finies
(Morphismes, lemme \ref{morphisms-lemma-quasi-finite}), et il
est ainsi fini d'après le lemme \ref{lemma-characterize-finite}.
\end{proof}

\begin{lemma}
\label{lemma-flat-proper-family-cannot-collapse-fibre}
Considérons un diagramme commutatif de schémas
$$
\xymatrix{
X \ar[rr]_h \ar[rd]_f & & Y \ar[ld]^g \\
& S
}
$$
Soit $s \in S$. Supposons
\begin{enumerate}
\item que $X \to S$ soit un morphisme propre,
\item que $Y \to S$ soit séparé et localement de type fini, et
\item que l'image de $X_s \to Y_s$ soit finie.
\end{enumerate}
Il existe alors un ouvert
$U \subset S$ contenant $s$ tel que $X_U \to Y_U$
se factorise par un sous-schéma fermé $Z \subset Y_U$ fini sur $U$.
\end{lemma}

\begin{proof}
Soit $Z \subset Y$ l'image schématique de $h$ ; voir
Morphismes, section \ref{morphisms-section-scheme-theoretic-image}.
D'après Morphismes, lemme \ref{morphisms-lemma-scheme-theoretic-image-is-proper},
le morphisme $X \to Z$ est surjectif et $Z \to S$ est propre.
Ainsi, $X_s \to Z_s$ est surjectif. On voit alors que
(3) implique que $Z_s$ est fini.
Par conséquent, $Z \to S$ est fini dans un voisinage ouvert de $s$ d'après le
lemme \ref{lemma-proper-finite-fibre-finite-in-neighbourhood}.
\end{proof}





\section{Applications du théorème principal de Zariski, II}
\label{section-applications-zmt-II}

\noindent
Dans cette section, on donne quelques conséquences supplémentaires du théorème principal de Zariski
concernant la structure des morphismes quasi-finis.

\begin{lemma}
\label{lemma-separated-locally-quasi-finite-over-affine}
Soit $f : X \to Y$ un morphisme séparé et localement quasi-fini,
où $Y$ est affine. Alors tout ensemble fini de points de $X$ est contenu
dans un ouvert affine de $X$.
\end{lemma}

\begin{proof}
Soient $x_1, \ldots, x_n \in X$. Choisissons un ouvert quasi-compact
$U \subset X$ tel que $x_i \in U$. Alors $U \to Y$ est quasi-affine d'après le
lemme \ref{lemma-quasi-finite-separated-quasi-affine}.
Il existe donc un ouvert affine $V \subset U$ contenant
$x_1, \ldots, x_n$ d'après
Propriétés, lemme \ref{properties-lemma-ample-finite-set-in-affine}.
\end{proof}

\begin{lemma}
\label{lemma-quasi-finite-finite-over-dense-open}
Soit $f : Y \to X$ un morphisme quasi-fini.
Il existe un ouvert dense $U \subset X$ tel que
$f|_{f^{-1}(U)} : f^{-1}(U) \to U$ soit fini.
\end{lemma}

\begin{proof}
Si $U_i \subset X$, $i \in I$, est une famille d'ouverts telle que les
restrictions $f|_{f^{-1}(U_i)} : f^{-1}(U_i) \to U_i$ soient finies,
alors, pour $U = \bigcup U_i$, la restriction $f|_{f^{-1}(U)} : f^{-1}(U) \to U$
est finie ; voir
Morphismes, lemme \ref{morphisms-lemma-finite-local}.
Le problème est donc local sur $X$, et l'on peut supposer que $X$ est affine.

\medskip\noindent
Supposons $X$ affine.
Écrivons $Y = \bigcup_{j = 1, \ldots, m} V_j$, avec $V_j$ affine.
C'est possible puisque $f$ est quasi-fini et donc, en
particulier, quasi-compact. Chaque $V_j \to X$ est quasi-fini
et séparé. Soit $\eta \in X$ un point générique d'une composante irréductible
de $X$. On déduit de
Morphismes, lemmes
\ref{morphisms-lemma-quasi-finite} et \ref{morphisms-lemma-generically-finite}
qu'il existe un voisinage ouvert $\eta \in U_\eta$ tel que
$f^{-1}(U_\eta) \cap V_j \to U_\eta$ soit fini. On peut choisir $U_\eta$ de sorte
qu'il convienne pour chaque $j = 1, \ldots, m$.
Remarquons que l'ensemble des points génériques de $X$ est dense dans $X$.
On voit donc qu'il existe un ouvert dense $W = \bigcup_\eta U_\eta$
tel que chaque morphisme $f^{-1}(W) \cap V_j \to W$ soit fini.
Il suffit de montrer qu'il existe un ouvert dense $U \subset W$
tel que $f|_{f^{-1}(U)} : f^{-1}(U) \to U$ soit fini.
Ainsi, on peut remplacer $X$ par un sous-schéma ouvert affine de $W$ et
supposer que chaque $V_j \to X$ est fini.

\medskip\noindent
Supposons $X$ affine, $Y = \bigcup_{j = 1, \ldots, m} V_j$ avec les $V_j$ affines,
et les restrictions $f|_{V_j} : V_j \to X$ finies.
Posons
$$
\Delta_{ij} =
\Big(\overline{V_i \cap V_j} \setminus V_i \cap V_j\Big) \cap V_j.
$$
C'est un fermé nulle part dense de $V_j$, car c'est la frontière
du sous-ensemble ouvert $V_i \cap V_j$ dans $V_j$. D'après
Morphismes, lemme \ref{morphisms-lemma-image-nowhere-dense-finite},
l'image $f(\Delta_{ij})$ est un fermé nulle part dense de $X$. D'après
Topologie, lemme \ref{topology-lemma-nowhere-dense},
la réunion $T = \bigcup f(\Delta_{ij})$ est un fermé nulle part dense
de $X$. Ainsi, $U = X \setminus T$ est un ouvert dense de $X$.
Nous affirmons que $f|_{f^{-1}(U)} : f^{-1}(U) \to U$ est fini.
Pour le voir, soit $U' \subset U$ un ouvert affine.
Posons $Y' = f^{-1}(U') = U' \times_X Y$,
$V_j' = Y' \cap V_j = U' \times_X V_j$. Considérons la restriction
$$
f' = f|_{Y'} : Y' \longrightarrow U'
$$
de $f$. Ce morphisme possède alors la propriété suivante :
$Y' = \bigcup_{j = 1, \ldots, m} V'_j$ est un recouvrement ouvert affine,
chaque $V'_j \to U'$ est fini, et $V_i' \cap V_j'$ est (ouvert et) fermé
à la fois dans $V'_i$ et dans $V'_j$. Ainsi, $V_i' \cap V_j'$ est affine, et l'application
$$
\mathcal{O}(V'_i) \otimes_{\mathbf{Z}} \mathcal{O}(V'_j)
\longrightarrow
\mathcal{O}(V'_i \cap V'_j)
$$
est surjective. Cela entraîne que $Y'$ est séparé, voir
Schémas, lemme \ref{schemes-lemma-characterize-separated}.
Enfin, considérons le diagramme commutatif
$$
\xymatrix{
\coprod_{j = 1, \ldots, m} V'_j \ar[rd] \ar[rr] & & Y' \ar[ld] \\
& U' &
}
$$
La flèche sud-est est finie, donc propre, la flèche horizontale est
surjective, et la flèche sud-ouest est séparée. Par conséquent, d'après
Morphismes, lemme \ref{morphisms-lemma-image-proper-is-proper},
on conclut que $Y' \to U'$ est propre. Comme il est aussi quasi-fini,
on voit qu'il est fini d'après le lemme \ref{lemma-characterize-finite},
ce qui conclut.
\end{proof}

\begin{lemma}
\label{lemma-stratify-flat-fp-lqf-universally-bounded}
Soit $f : X \to S$ un morphisme plat, localement de présentation finie, séparé,
localement quasi-fini, à fibres universellement bornées. Il existe alors
des sous-ensembles fermés
$$
\emptyset = Z_{-1} \subset Z_0 \subset Z_1 \subset Z_2 \subset
\ldots \subset Z_n = S
$$
tels que, si $S_r = Z_r \setminus Z_{r - 1}$, la stratification
$S = \coprod_{r = 0, \ldots, n} S_r$ soit caractérisée par la propriété
universelle suivante : étant donné $g : T \to S$, la projection
$X \times_S T \to T$ est finie localement libre
de degré $r$ si et seulement si $g(T) \subset S_r$ (ensemblistement).
\end{lemma}

\begin{proof}
Soit $n$ un entier qui majore le degré des fibres de $X \to S$.
D'après Morphismes, lemme \ref{morphisms-lemma-base-change-universally-bounded},
on voit que les degrés des fibres de tout changement de base sont encore majorés par $n$.
En particulier, tous les entiers $r$ qui interviennent dans l'énoncé du lemme
seront $\leq n$. Nous allons démontrer le lemme par récurrence sur $n$. Le cas
initial $n = 0$ est évident.

\medskip\noindent
Nous affirmons que l'ensemble des points $s \in S$
tels que $\deg_{\kappa(s)}(X_s) = n$ est un sous-ensemble ouvert $S_n \subset S$
et que $X \times_S S_n \to S_n$ est fini localement libre de degré $n$.
En effet, soit $s \in S$ un tel point. Choisissons un morphisme étale
\'elementaire $(U, u) \to (S, s)$ et une décomposition
$U \times_S X = W \amalg V$ comme dans le
lemme \ref{lemma-etale-splits-off-quasi-finite-part}.
Puisque $V \to U$ est fini, plat et localement de présentation finie,
on voit que $V \to U$ est fini localement libre, voir
Morphismes, lemme \ref{morphisms-lemma-finite-flat}.
Après avoir rétréci $U$ à un voisinage plus petit de $u$,
on peut supposer que $V \to U$ est fini localement libre d'un certain degré $d$, voir
Morphismes, lemme \ref{morphisms-lemma-finite-locally-free}.
Comme $u \mapsto s$ et $W_u = \emptyset$, on voit que $d = n$. Puisque $n$
est le degré maximal d'une fibre, on voit que $W = \emptyset$ !
Ainsi, $U \times_S X \to U$ est fini localement libre de degré $n$.
D'après Descente, lemme \ref{descent-lemma-descending-property-finite-locally-free},
on conclut que $X \to S$ est fini localement libre de degré $n$
au-dessus de $\Im(U \to S)$, qui est un voisinage ouvert de $s$
(Morphismes, lemme \ref{morphisms-lemma-etale-open}).
Ceci démontre l'affirmation.

\medskip\noindent
Soit $S' = S \setminus S_n$, muni de la structure de schéma réduit induite,
et posons $X' = X \times_S S'$. Remarquons que les degrés des fibres
de $X' \to S'$ sont universellement majorés par $n - 1$. Par récurrence, on trouve une
stratification $S' = S_0 \amalg \ldots \amalg S_{n - 1}$ adaptée
au morphisme $X' \to S'$. Nous affirmons que $S = \coprod_{r = 0, \ldots, n} S_r$
convient au morphisme $X \to S$. Soit $g : T \to S$ un morphisme de schémas
et supposons que $X \times_S T \to T$ soit fini localement libre de degré $r$.
Comme on l'a remarqué plus haut, cela implique que $r \leq n$. Si $r = n$, il est
clair que $T \to S$ se factorise par $S_n$. Si $r < n$, alors
$g(T) \subset S' = S \setminus S_d$ (ensemblistement), donc
$T_{red} \to S$ se factorise par $S'$, voir
Schémas, lemme \ref{schemes-lemma-map-into-reduction}.
Remarquons que $X \times_S T_{red} \to T_{red}$ est lui aussi
fini localement libre de degré $r$, comme changement de base.
Par la propriété universelle de la stratification
$S' = \coprod_{r = 0, \ldots, n - 1} S_r$, on voit que $g(T) = g(T_{red})$
est contenu dans $S_r$.
Réciproquement, supposons donné $g : T \to S$ tel que
$g(T) \subset S_r$ (ensemblistement).
Si $r = n$, alors $g$ se factorise par $S_n$ et
il est clair que $X \times_S T \to T$
est fini localement libre de degré $n$, comme changement de base.
Si $r < n$, alors $X \times_S T \to T$ est un morphisme
séparé, plat et localement de présentation finie, tel que
sa restriction à $T_{red}$ soit finie localement libre de degré $r$.
Puisque $T_{red} \to T$ est un homéomorphisme universel, on conclut
que $X \times_S T_{red} \to X \times_S T$ est aussi un homéomorphisme universel
et donc que $X \times_S T \to T$ est universellement fermé (puisque ceci
est vrai pour le morphisme fini $X \times_S T_{red} \to T_{red}$).
Il s'ensuit que $X \times_S T \to T$ est fini, par exemple d'après le
lemme \ref{lemma-characterize-finite}. On peut alors utiliser
Morphismes, lemme \ref{morphisms-lemma-finite-flat},
pour voir que $X \times_S T \to T$ est fini localement libre.
Enfin, son degré est $r$, puisque toutes ses fibres sont de degré $r$.
\end{proof}

\begin{lemma}
\label{lemma-stratify-flat-fp-qf}
Soit $f : X \to S$ un morphisme de schémas plat, localement de
présentation finie, séparé et quasi-fini. Il existe alors
des sous-ensembles fermés
$$
\emptyset = Z_{-1} \subset Z_0 \subset Z_1 \subset Z_2 \subset
\ldots \subset S
$$
tels que, si $S_r = Z_r \setminus Z_{r - 1}$, la stratification
$S = \coprod S_r$ soit caractérisée par la propriété universelle suivante :
pour tout morphisme $g : T \to S$, la projection $X \times_S T \to T$ est
finie localement libre de degré $r$ si et seulement si $g(T) \subset S_r$
(ensemblistement). De plus, les morphismes d'inclusion $S_r \to S$ sont
quasi-compacts.
\end{lemma}

\begin{proof}
La question est locale sur $S$ ; on peut donc supposer $S$ affine.
D'après Morphismes, lemme
\ref{morphisms-lemma-locally-quasi-finite-qc-source-universally-bounded},
les fibres de $f$ sont alors universellement bornées.
L'existence de la stratification résulte donc du
lemme \ref{lemma-stratify-flat-fp-lqf-universally-bounded}.

\medskip\noindent
Nous allons montrer que $U_r = S \setminus Z_r \to S$ est quasi-compact pour
chaque $r \geq 0$. Cela démontrera la dernière assertion par topologie élémentaire.
Comme une composée d'applications quasi-compactes est quasi-compacte,
il suffit de démontrer que $U_r \to U_{r - 1}$ est quasi-compact.
Choisissons un ouvert affine $W \subset U_{r - 1}$. Écrivons $W = \Spec(A)$.
Alors $Z_r \cap W = V(I)$ pour un certain idéal $I \subset A$,
et $X \times_S \Spec(A/I) \to \Spec(A/I)$ est fini localement
libre de degré $r$. Remarquons que $A/I = \colim A/I_i$, où $I_i \subset I$
parcourt les idéaux de type fini. D'après
Limites, lemme \ref{limits-lemma-descend-finite-locally-free},
on voit que $X \times_S \Spec(A/I_i) \to \Spec(A/I_i)$
est fini localement libre de degré $r$ pour un certain $i$. (On utilise ici
que $X \to S$ est de présentation finie, puisqu'il est localement de présentation
finie, séparé et quasi-compact.)
Ainsi, $\Spec(A/I_i) \to \Spec(A) = W$ se factorise (ensemblistement)
par $Z_r \cap W$. Il s'ensuit que $Z_r \cap W = V(I_i)$ est l'ensemble des zéros
d'une famille finie d'éléments de $A$. Cela signifie que
$W \setminus Z_r$ est une réunion finie d'ouverts principaux, donc est quasi-compact,
comme souhaité.
\end{proof}

\begin{lemma}
\label{lemma-stratify-flat-fp-lqf}
Soit $f : X \to S$ un morphisme de schémas plat, localement de présentation finie,
séparé et localement quasi-fini. Il existe alors
des sous-schémas ouverts
$$
S = U_0 \supset U_1 \supset U_2 \supset \ldots
$$
tels qu'un morphisme $\Spec(k) \to S$, où $k$ est un corps,
se factorise par $U_d$ si et
seulement si $X \times_S \Spec(k)$ est de degré $\geq d$ sur $k$.
\end{lemma}

\begin{proof}
L'énoncé signifie simplement que l'ensemble des points où le degré
de la fibre est $\geq d$ est ouvert. On peut donc travailler localement sur $S$ et
supposer $S$ affine. Dans ce cas, pour tout ouvert quasi-compact $W \subset X$,
l'ensemble $U_d(W)$ des points où les fibres de $W \to S$ sont de
degré $\geq d$ est ouvert d'après le lemme \ref{lemma-stratify-flat-fp-qf}.
Puisque $U_d = \bigcup_W U_d(W)$, le résultat s'ensuit.
\end{proof}

\begin{lemma}
\label{lemma-go-down-with-annihilators}
Soit $f : X \to S$ un morphisme de schémas plat, localement de
présentation finie et localement quasi-fini. Soit
$g \in \Gamma(X, \mathcal{O}_X)$ non nul. Il existe alors
un ouvert $V \subset X$ tel que $g|_V \not = 0$, un ouvert
$U \subset S$ s'insérant dans un diagramme commutatif
$$
\xymatrix{
V \ar[r] \ar[d]_\pi & X \ar[d]^f \\
U \ar[r] & S,
}
$$
un sous-faisceau quasi-cohérent $\mathcal{F} \subset \mathcal{O}_U$, un entier
$r > 0$, et un homomorphisme injectif de $\mathcal{O}_U$-modules
$\mathcal{F}^{\oplus r} \to \pi_*\mathcal{O}_V$
dont l'image contient $g|_V$.
\end{lemma}

\begin{proof}
On peut supposer $X$ et $S$ affines. On obtient une filtration
$\emptyset = Z_{-1} \subset Z_0 \subset Z_1 \subset Z_2 \subset \ldots
\subset Z_n = S$ comme dans les
lemmes \ref{lemma-stratify-flat-fp-lqf-universally-bounded} et
\ref{lemma-stratify-flat-fp-qf}.
Soit $T \subset X$ le support schématique du module fini
sur $\mathcal{O}_X$ donné par $\Im(g : \mathcal{O}_X \to \mathcal{O}_X)$.
Remarquons que $T$ est le support de $g$ comme section de $\mathcal{O}_X$
(Modules, définition \ref{modules-definition-support}) et que,
pour tout ouvert $V \subset X$, on a $g|_V \not = 0$ si et seulement si
$V \cap T \not = \emptyset$.
Soit $r$ le plus petit entier tel que $f(T) \subset Z_r$
ensemblistement. Soit $\xi \in T$ un point générique d'une composante irréductible
de $T$ tel que $f(\xi) \not \in Z_{r - 1}$ (et donc
$f(\xi) \in Z_r$). On peut remplacer $S$ par un voisinage affine de
$f(\xi)$ contenu dans $S \setminus Z_{r - 1}$. Écrivons $S = \Spec(A)$
et soit $I = (a_1, \ldots, a_m) \subset A$ un idéal de type fini
tel que $V(I) = Z_r$ (ensemblistement, voir
Algèbre, lemme \ref{algebra-lemma-qc-open}).
Puisque le support de $g$ est contenu dans $f^{-1}V(I)$ par notre choix de $r$,
on voit qu'il existe un entier $N$ tel que
$a_j^N g = 0$ pour $j = 1, \ldots, m$. En remplaçant $a_j$ par $a_j^r$,
on peut supposer que $Ig = 0$. Pour tout $A$-module $M$, notons
$M[I]$ le sous-module des éléments annulés par $I$ dans $M$, c'est-à-dire $M[I] = \{m \in M \mid Im = 0\}$.
Écrivons $X = \Spec(B)$, de sorte que $g \in B[I]$. Comme $A \to B$ est plat, on
voit que
$$
B[I] = A[I] \otimes_A B \cong A[I] \otimes_{A/I} B/IB
$$
Par notre choix de $Z_r$, le $A/I$-module $B/IB$ est
fini localement libre de rang $r$. Ainsi, après avoir remplacé $S$ par
un voisinage ouvert affine plus petit de $f(\xi)$, on peut supposer
que $B/IB \cong (A/IA)^{\oplus r}$ comme $A/I$-modules.
Choisissons une application $\psi : A^{\oplus r} \to B$ dont la réduction modulo $I$ est
l'isomorphisme de la phrase précédente. On voit alors que
l'application induite
$$
A[I]^{\oplus r} \longrightarrow B[I]
$$
est un isomorphisme. Le lemme s'ensuit en prenant pour $\mathcal{F}$ le
faisceau quasi-cohérent associé au $A$-module $A[I]$, et pour
l'application $\mathcal{F}^{\oplus r} \to \pi_*\mathcal{O}_V$ celle
qui correspond à $A[I]^{\oplus r} \subset A^{\oplus r} \to B$.
\end{proof}

\begin{lemma}
\label{lemma-there-is-a-scheme-integral-over}
Soit $U \to X$ un morphisme \'etale surjectif de schémas. Supposons $X$
quasi-compact et quasi-séparé. Il existe alors un morphisme entier surjectif
$Y \to X$ tel que, pour
tout $y \in Y$, il existe un voisinage ouvert $V \subset Y$
tel que $V \to X$ se factorise par $U$. En fait, on peut supposer
$Y \to X$ fini et de présentation finie.
\end{lemma}

\begin{proof}
Puisque $X$ est quasi-compact, il existe un nombre fini d'ouverts affines
$U_i \subset U$ tels que $U' = \coprod U_i \to X$ soit surjectif.
Après avoir remplacé $U$ par $U'$, on voit que l'on peut supposer $U$ affine.
En particulier, $U \to X$ est séparé
(Schémas, lemme \ref{schemes-lemma-affine-separated}).
Il existe alors un entier $d$ majorant le degré des fibres géométriques
de $U \to X$ (voir Morphismes, lemme
\ref{morphisms-lemma-locally-quasi-finite-qc-source-universally-bounded}).
Nous allons démontrer le lemme par récurrence sur $d$ pour tous les schémas
quasi-compacts et séparés $U$ se projetant de manière \'etale et surjective sur $X$.
Si $d = 1$, alors $U = X$ et le résultat vaut avec $Y = U$.
Supposons $d > 1$.

\medskip\noindent
On applique le lemme \ref{lemma-quasi-finite-separated-quasi-affine}
et on obtient une factorisation
$$
\xymatrix{
U \ar[rr]_j \ar[rd] & & Y \ar[ld]^\pi \\
& X
}
$$
où $\pi$ est entier et $j$ une immersion ouverte quasi-compacte. On peut, et l'on va,
supposer que $j(U)$ est schématiquement dense dans $Y$. Remarquons que
$$
U \times_X Y = U \amalg W
$$
où le premier terme est l'image de $U \to U \times_X Y$
(qui est fermée d'après
Schémas, lemme \ref{schemes-lemma-semi-diagonal},
et ouverte parce qu'elle est \'etale comme morphisme entre schémas \'etales sur $Y$),
et le second facteur est le complémentaire (ouvert et fermé).
L'image $V \subset Y$ de $W$ est un sous-schéma ouvert contenant
$Y \setminus U$.

\medskip\noindent
Le morphisme \'etale $W \to Y$ a des fibres géométriques de cardinal $< d$.
En effet, cela est clair par inspection pour les points géométriques de $U \subset Y$.
Puisque $U \subset Y$ est dense, cela vaut pour tous les points géométriques de $Y$,
par exemple d'après le lemme
\ref{lemma-stratify-flat-fp-lqf-universally-bounded}
(le degré des fibres d'un morphisme \'etale séparé quasi-compact
ne croît pas par spécialisation). On peut donc appliquer l'hypothèse de récurrence
à $W \to V$ et trouver un morphisme entier surjectif
$Z \to V$, avec $Z$ un schéma, qui se factorise localement pour la topologie de Zariski par $W$.
Choisissons une factorisation $Z \to Z' \to Y$ où $Z' \to Y$ est entier et
$Z \to Z'$ une immersion ouverte
(lemme \ref{lemma-quasi-finite-separated-quasi-affine}).
Après avoir remplacé $Z'$ par l'adhérence schématique de $Z$ dans $Z'$,
on peut supposer $Z$ schématiquement dense dans $Z'$.
Après cette opération, on a $Z' \times_Y V = Z$. Enfin,
soit $T \subset Y$ le sous-schéma fermé réduit induit
sur $Y \setminus V$. Considérons le morphisme
$$
Z' \amalg T \longrightarrow X
$$
C'est, par construction, un morphisme entier surjectif.
Puisque $T \subset U$, il est clair que le morphisme $T \to X$
se factorise par $U$. D'autre part, soit $z \in Z'$
un point. Si $z \not \in Z$, alors $z$ s'envoie sur un point de
$Y \setminus V \subset U$, et l'on trouve un voisinage de $z$
sur lequel le morphisme se factorise par $U$.
Si $z \in Z$, alors il existe un voisinage $\Omega \subset Z$
qui se factorise par $W \subset U \times_X Y$, et donc par $U$.
Ceci démontre l'existence.

\medskip\noindent
Supposons avoir trouvé un morphisme entier surjectif $Y \to X$ qui se factorise
localement pour la topologie de Zariski par $U$. Choisissons un recouvrement ouvert affine fini
$Y = \bigcup V_j$ tel que $V_j \to X$ se factorise par $U$. On peut
écrire $Y = \lim Y_i$, où $Y_i \to X$ est fini et de présentation
finie, voir Limites, lemme
\ref{limits-lemma-integral-limit-finite-and-finite-presentation}.
Pour $i$ assez grand, on peut trouver des ouverts affines $V_{i, j} \subset Y_i$
dont les images réciproques dans $Y$ redonnent les $V_j$, voir
Limites, lemme \ref{limits-lemma-descend-opens}.
Pour $i$ encore plus grand, les morphismes $V_j \to U$ au-dessus de $X$
proviennent de morphismes $V_{i, j} \to U$ au-dessus de $X$, voir
Limites, proposition
\ref{limits-proposition-characterize-locally-finite-presentation}.
Ceci achève la démonstration.
\end{proof}







\section{Application aux morphismes à fibres connexes}
\label{section-application-geometrically-connected}

\noindent
Dans cette section, nous démontrons quelques lemmes qui produisent des morphismes
dont toutes les fibres sont géométriquement connexes ou géométriquement intègres.
Ils nous seront utiles plus loin dans notre étude de la structure locale des morphismes
de type fini.

\begin{lemma}
\label{lemma-descent-connected-fibres}
Considérons un diagramme de morphismes de schémas
$$
\xymatrix{
Z \ar[r]_{\sigma} \ar[rd] & X \ar[d] \\
& Y
}
$$
et un point $y \in Y$. Supposons que
\begin{enumerate}
\item $X \to Y$ soit de présentation finie et plat,
\item $Z \to Y$ soit fini localement libre,
\item $Z_y \not = \emptyset$,
\item toutes les fibres de $X \to Y$ soient géométriquement réduites, et
\item $X_y$ soit géométriquement connexe sur $\kappa(y)$.
\end{enumerate}
Il existe alors un ouvert quasi-compact $X^0 \subset X$ tel que $X^0_y = X_y$
et tel que toutes les fibres non vides de $X^0 \to Y$ soient géométriquement connexes.
\end{lemma}

\begin{proof}
Dans cette démonstration, nous utiliserons que les propriétés d'être plat, d'être de présentation
finie et d'être fini localement libre sont préservées par changement de base et par composition.
Nous utiliserons aussi qu'un morphisme fini localement libre est à la fois ouvert et
fermé. On trouvera ces faits dans
Morphismes, lemmes
\ref{morphisms-lemma-base-change-flat},
\ref{morphisms-lemma-base-change-finite-presentation},
\ref{morphisms-lemma-base-change-finite-locally-free},
\ref{morphisms-lemma-composition-flat},
\ref{morphisms-lemma-composition-finite-presentation},
\ref{morphisms-lemma-composition-finite-locally-free},
\ref{morphisms-lemma-fppf-open}, et
\ref{morphisms-lemma-finite-proper}.

\medskip\noindent
Remarquons que $X_Z \to Z$ est un morphisme plat de présentation finie
qui possède une section $s$ provenant de $\sigma$. Notons $X_Z^0$
le sous-ensemble de $X_Z$ défini dans la
situation \ref{situation-connected-along-section}.
D'après le
lemme \ref{lemma-connected-along-section-open},
c'est un sous-ensemble ouvert de $X_Z$.

\medskip\noindent
L'image réciproque $X_{Z \times_Y Z}$ de $X$ sur $Z \times_Y Z$ est munie
de deux sections $s_0, s_1$, à savoir les changements de base de $s$ par
$\text{pr}_0, \text{pr}_1 : Z \times_Y Z \to Z$. La construction de la
situation \ref{situation-connected-along-section}
donne deux sous-ensembles $(X_{Z \times_Y Z})_{s_0}^0$ et
$(X_{Z \times_Y Z})_{s_1}^0$. D'après le
lemme \ref{lemma-base-change-connected-along-section},
ce sont les images réciproques de $X_Z^0$ par les morphismes
$1_X \times \text{pr}_0, 1_X \times \text{pr}_1 : X_{Z \times_Y Z} \to X_Z$.
En particulier, ces sous-ensembles sont ouverts.

\medskip\noindent
Soit $(Z \times_Y Z)_y = \{z_1, \ldots, z_n\}$.
Comme $X_y$ est géométriquement connexe, on voit que les fibres de
$(X_{Z \times_Y Z})_{s_0}^0$ et de $(X_{Z \times_Y Z})_{s_1}^0$
au-dessus de chaque $z_i$ coïncident (elles sont égales à la fibre entière). Une autre
manière d'exprimer ceci est
$$
s_0(z_i) \in (X_{Z \times_Y Z})_{s_1}^0
\quad\text{et}\quad
s_1(z_i) \in (X_{Z \times_Y Z})_{s_0}^0.
$$
Puisque les ensembles $(X_{Z \times_Y Z})_{s_0}^0$ et $(X_{Z \times_Y Z})_{s_1}^0$
sont ouverts dans $X_{Z \times_Y Z}$, il existe un voisinage ouvert
$W \subset Z \times_Y Z$ de $(Z \times_Y Z)_y$ tel que
$$
s_0(W) \subset (X_{Z \times_Y Z})_{s_1}^0
\quad\text{et}\quad
s_1(W) \subset (X_{Z \times_Y Z})_{s_0}^0.
$$
Il résulte alors directement de la construction de la
situation \ref{situation-connected-along-section}
que
$$
p^{-1}(W) \cap (X_{Z \times_Y Z})_{s_0}^0
=
p^{-1}(W) \cap (X_{Z \times_Y Z})_{s_1}^0
$$
où $p : X_{Z \times_Y Z} \to Z \times_W Z$ est la projection.
Puisque $Z \times_Y Z \to Y$ est fini localement libre, donc ouvert et fermé,
il existe un voisinage ouvert affine $V \subset Y$ de $y$ tel que
$q^{-1}(V) \subset W$, où $q : Z \times_Y Z \to Y$ est le
morphisme structural. Pour démontrer le lemme, on peut remplacer $Y$ par $V$.
Après ce remplacement, on voit que $X_Z^0 \subset Y_Z$ est un ouvert tel que
$$
(1_X \times \text{pr}_0)^{-1}(X_Z^0) =
(1_X \times \text{pr}_1)^{-1}(X_Z^0).
$$
Cela signifie que l'image $X^0 \subset X$ de $X_Z^0$ est un ouvert tel
que $(X_Z \to X)^{-1}(X^0) = X_Z^0$, voir
Descente, lemme \ref{descent-lemma-open-fpqc-covering}.
Enfin, $X^0$ est quasi-compact parce que $X_Z^0$ est quasi-compact
d'après le lemme \ref{lemma-connected-along-section-locally-constructible}
(on utilise qu'à ce stade $Y$ est affine, donc $X$ est quasi-compact et
quasi-séparé, de sorte que localement constructible équivaut à constructible
et entraîne en particulier quasi-compact ; détails omis).
On voit ainsi que $X^0$ possède toutes les propriétés voulues.
\end{proof}

\begin{lemma}
\label{lemma-fibre-geometrically-connected-reduced}
Soit $h : Y \to S$ un morphisme de schémas.
Soit $s \in S$ un point.
Soit $T \subset Y_s$ un sous-schéma ouvert.
Supposons que
\begin{enumerate}
\item $h$ soit plat et de présentation finie,
\item toutes les fibres de $h$ soient géométriquement réduites, et
\item $T$ soit géométriquement connexe sur $\kappa(s)$.
\end{enumerate}
On peut alors trouver un voisinage \'etale élémentaire affine
$(S', s') \to (S, s)$
et un ouvert quasi-compact $V \subset Y_{S'}$ tels que
\begin{enumerate}
\item[(a)] toutes les fibres de $V \to S'$ soient géométriquement connexes,
\item[(b)] $V_{s'} = T \times_s s'$.
\end{enumerate}
\end{lemma}

\begin{proof}
Le problème est manifestement local sur $S$ ; on peut donc remplacer $S$ par un
voisinage ouvert affine de $s$.
La topologie de $Y_s$ est induite par celle de $Y$, voir
Schémas, lemme \ref{schemes-lemma-fibre-topological}.
On peut donc trouver un ouvert quasi-compact $V \subset Y$ tel que $V_s = T$.
La restriction de $h$ à $V$ est quasi-compacte (car $S$ est affine et $V$
quasi-compact), quasi-séparée, localement de présentation finie et
plate, donc plate de présentation finie.
Ainsi, après avoir remplacé $Y$ par $V$, on peut supposer, en plus
de (1) et (2), que $Y_s = T$ et que $S$ est affine.

\medskip\noindent
Choisissons un point fermé $y \in Y_s$ tel que $h$ soit de Cohen--Macaulay en $y$, voir
le lemme \ref{lemma-flat-finite-presentation-CM-open}.
D'après le
lemme \ref{lemma-slice-CM},
il existe un diagramme
$$
\xymatrix{
Z \ar[r] \ar[rd] & Y \ar[d] \\
& S
}
$$
tel que $Z \to S$ soit plat, localement de présentation finie, localement
quasi-fini, avec $Z_s = \{y\}$. Appliquons le
lemme \ref{lemma-etale-makes-quasi-finite-finite-at-point}
pour trouver un voisinage élémentaire $(S', s') \to (S, s)$ et un ouvert
$Z' \subset Z_{S'} = S' \times_S Z$ tel que $Z' \to S'$ soit fini et possède un unique
point $z' \in Z'$ au-dessus de $s$. Remarquons que $Z' \to S'$ est aussi
localement de présentation finie et plat (comme ouvert du changement de base
de $Z \to S$), donc $Z' \to S'$ est fini localement libre, voir
Morphismes, lemme \ref{morphisms-lemma-finite-flat}.
Remarquons que $Y_{S'} \to S'$ est plat et de présentation finie,
à fibres géométriquement réduites, comme changement de base de $h$.
De plus, $Y_{s'} = Y_s$ est géométriquement connexe.
Appliquons le lemme \ref{lemma-descent-connected-fibres}
à $Z' \to Y_{S'}$ au-dessus de $S'$ pour obtenir un ouvert quasi-compact $V \subset Y_{S'}$
satisfaisant (2), dont les fibres sur $S'$ sont soit vides, soit
géométriquement connexes. Comme $V \to S'$ est ouvert
(Morphismes, lemme \ref{morphisms-lemma-fppf-open}), après avoir remplacé
$S'$ par un voisinage ouvert affine de $s'$,
on peut supposer $V \to S'$ surjectif, d'où (1).
\end{proof}

\begin{lemma}
\label{lemma-cover-by-geometrically-connected}
Soit $f : X \to S$ un morphisme de schémas
localement de présentation finie et plat, à fibres géométriquement
réduites. Il existe alors
un recouvrement \'etale $\{X_i \to X\}_{i \in I}$
tel que $X_i \to S$ se factorise en $X_i \to S_i \to S$,
où $S_i \to S$ est \'etale et $X_i \to S_i$ est
plat de présentation finie, à fibres géométriquement connexes
et géométriquement réduites.
\end{lemma}

\begin{proof}
Choisissons un point $x \in X$ d'image $s \in S$. Nous allons construire
un diagramme
$$
\xymatrix{
X' \ar[r] \ar[rd] & S' \times_S X \ar[r] \ar[d] & X \ar[d] \\
& S' \ar[r] & S
}
$$
et des points $s' \in S'$, $x' \in X'$, $y \in S' \times_S X$
tels que $x'$ s'envoie sur $x$, que $(S', s') \to (S, s)$
soit un voisinage \'etale, que $(X', x') \to (S' \times_S X, y)$
soit un voisinage \'etale\footnote{La démonstration donne en fait
un ouvert $X' \subset S' \times_S X$.}, et que
$X' \to S'$ ait des fibres géométriquement
connexes. Si l'on peut faire cela pour tout $x \in X$, alors
le lemme s'ensuit (les membres du recouvrement étant donnés par la
famille des morphismes \'etales $X' \to X$ ainsi obtenus).
La première étape consiste à remplacer $X$ et $S$ par des voisinages ouverts affines
de $x$ et de $s$, ce qui nous ramène au cas où $X$ et $S$ sont affines
(et donc où $f$ est de présentation finie).

\medskip\noindent
Choisissons une extension algébrique séparable $\overline{k}$ de $\kappa(s)$.
Notons $X_{\overline{k}}$ le changement de base de $X_s$.
Choisissons un point $\overline{x}$ de $X_{\overline{k}}$ s'envoyant sur $x \in X_s$.
Choisissons un voisinage ouvert connexe quasi-compact
$\overline{V} \subset X_{\overline{k}}$
de $\overline{x}$. (Cela est possible parce que tout schéma
localement de type fini sur un corps est localement connexe
comme espace topologique localement noethérien.)
D'après Variétés, lemme \ref{varieties-lemma-Galois-action-quasi-compact-open},
on peut trouver une extension finie séparable $k'/\kappa(s)$
et un ouvert quasi-compact $V' \subset X_{k'}$ dont le
changement de base est $\overline{V}$. En particulier, $V'$ est
géométriquement connexe sur $k'$, voir
Variétés, lemme \ref{varieties-lemma-characterize-geometrically-connected}. D'après
le lemme \ref{lemma-realize-prescribed-residue-field-extension-etale},
on peut trouver un voisinage \'etale $(S', s') \to (S, s)$
tel que $\kappa(s')$ soit isomorphe à $k'$ comme extension
de $\kappa(s)$.
Notons $x' \in (S' \times_S X)_{s'} = X_{k'}$ l'image de $\overline{x}$.
Ainsi, après avoir remplacé $(S, s)$ par $(S', s')$ et $(X, x)$ par
$(S' \times_S X, x')$, on se ramène au cas traité au
paragraphe suivant.

\medskip\noindent
Supposons qu'il existe un ouvert quasi-compact $V \subset X_s$
qui contienne $x$ et soit géométriquement connexe.
On peut alors appliquer le lemme \ref{lemma-fibre-geometrically-connected-reduced}
pour trouver un voisinage \'etale affine $(S', s') \to (S, s)$
et un ouvert quasi-compact $X' \subset S' \times_S X$ tels que
$X' \to S'$ ait des fibres géométriquement connexes
et que $X'$ contienne un point au-dessus de $x$.
Ceci achève la démonstration.
\end{proof}

\begin{lemma}
\label{lemma-normal-morphism-irreducible}
Soit $h : Y \to S$ un morphisme de schémas.
Soit $s \in S$ un point.
Soit $T \subset Y_s$ un sous-schéma ouvert.
Supposons que
\begin{enumerate}
\item $h$ soit de présentation finie,
\item $h$ soit normal, et
\item $T$ soit géométriquement irréductible sur $\kappa(s)$.
\end{enumerate}
Alors on peut trouver un voisinage \'etale élémentaire affine
$(S', s') \to (S, s)$ et un ouvert quasi-compact $V \subset Y_{S'}$ tels que
\begin{enumerate}
\item[(a)] toutes les fibres de $V \to S'$ soient géométriquement intègres,
\item[(b)] $V_{s'} = T \times_s s'$.
\end{enumerate}
\end{lemma}

\begin{proof}
Appliquons le
lemme \ref{lemma-fibre-geometrically-connected-reduced}
pour trouver un voisinage \'etale élémentaire affine $(S', s') \to (S, s)$ et
un ouvert quasi-compact $V \subset Y_{S'}$ tels que toutes les fibres de
$V \to S'$ soient géométriquement connexes et $V_{s'} = T \times_s s'$.
Comme $V$ est un ouvert du changement de base de $h$, toutes les fibres de $V \to S'$
sont géométriquement normales, voir le lemme \ref{lemma-normal}.
En particulier, elles sont géométriquement réduites. Pour achever la démonstration,
il reste à montrer qu'elles sont géométriquement irréductibles. Or, si $t \in S'$,
alors $V_t$ est de type fini sur $\kappa(t)$ et donc
$V_t \times_{\kappa(t)} \overline{\kappa(t)}$ est de type fini
sur $\overline{\kappa(t)}$, donc noethérien. Par le choix de $S' \to S$,
le schéma $V_t \times_{\kappa(t)} \overline{\kappa(t)}$ est connexe.
Ainsi $V_t \times_{\kappa(t)} \overline{\kappa(t)}$ est irréductible d'après
Propriétés, lemme \ref{properties-lemma-normal-Noetherian},
ce qui achève la démonstration.
\end{proof}








\section{Application à la structure des morphismes de type fini}
\label{section-application-finite-type}

\noindent
Le résultat de cette section se trouve dans \cite{GruRay}.
En termes informels, il affirme qu'un morphisme de type fini est, localement pour la topologie \'etale
sur la source et le but, la composée d'un
morphisme fini suivi d'un morphisme lisse à fibres géométriquement connexes
de dimension relative égale à la dimension de la fibre du morphisme
initial.

\begin{lemma}
\label{lemma-local-structure-finite-type}
Soit $f : X \to S$ un morphisme. Soit $x \in X$ et posons $s = f(x)$.
Supposons que $f$ soit localement de type fini et que $n = \dim_x(X_s)$.
Alors il existe un diagramme commutatif
$$
\xymatrix{
X \ar[dd] & X' \ar[l]^g \ar[d]^\pi & x \ar@{|->}[dd] &
x' \ar@{|->}[l]  \ar@{|->}[d] \\
& Y \ar[d]^h & & y \ar@{|->}[d] \\
S \ar@{=}[r] & S & s & s \ar@{=}[l]
}
$$
et un point $x' \in X'$ avec $g(x') = x$ tels que, si $y = \pi(x')$,
on ait
\begin{enumerate}
\item $h : Y \to S$ soit lisse de dimension relative $n$,
\item $g : (X', x') \to (X, x)$ soit un voisinage \'etale élémentaire,
\item $\pi$ soit fini, et $\pi^{-1}(\{y\}) = \{x'\}$, et
\item $\kappa(y)$ soit une extension purement transcendante de $\kappa(s)$.
\end{enumerate}
De plus, si $f$ est localement de présentation finie, alors $\pi$ est
de présentation finie.
\end{lemma}

\begin{proof}
Le problème est local sur $X$ et $S$ ; on peut donc supposer $X$ et
$S$ affines. D'après
Algèbre, lemme \ref{algebra-lemma-refined-quasi-finite-over-polynomial-algebra},
après avoir remplacé $X$ par un voisinage ouvert standard de $x$ dans $X$,
on peut supposer qu'il existe une factorisation
$$
\xymatrix{
X \ar[r]^\pi & \mathbf{A}^n_S \ar[r] & S
}
$$
telle que $\pi$ soit quasi-fini et que $\kappa(\pi(x))$
soit une extension purement transcendante de $\kappa(s)$. D'après le
lemme \ref{lemma-etale-makes-quasi-finite-finite-at-point},
il existe un voisinage \'etale élémentaire
$$
(Y, y) \to (\mathbf{A}^n_S, \pi(x))
$$
et un ouvert $X' \subset X \times_{\mathbf{A}^n_S} Y$ qui contient un
unique point $x'$ au-dessus de $y$, tel que $X' \to Y$ soit fini.
Ceci démontre les assertions (1) -- (4). Pour la dernière assertion, utiliser
Morphismes, lemme \ref{morphisms-lemma-finite-presentation-permanence}.
\end{proof}

\begin{lemma}
\label{lemma-local-local-structure-finite-type}
\begin{slogan}
Un morphisme de type fini est, dans des voisinages \'etales, fini au-dessus d'un
morphisme lisse.
\end{slogan}
Soit $f : X \to S$ un morphisme. Soit $x \in X$ et posons $s = f(x)$.
Supposons que $f$ soit localement de type fini et que $n = \dim_x(X_s)$.
Alors il existe un diagramme commutatif
$$
\xymatrix{
X \ar[dd] & X' \ar[l]^g \ar[d]^\pi & x \ar@{|->}[dd] &
x' \ar@{|->}[l] \ar@{|->}[d] \\
& Y' \ar[d]^h & & y' \ar@{|->}[d] \\
S & S' \ar[l]_e & s & s' \ar@{|->}[l]
}
$$
et un point $x' \in X'$ avec $g(x') = x$ tels que, si $y' = \pi(x')$,
$s' = h(y')$, on ait
\begin{enumerate}
\item $h : Y' \to S'$ soit lisse de dimension relative $n$,
\item toutes les fibres de $Y' \to S'$ soient géométriquement intègres,
\item $g : (X', x') \to (X, x)$ soit un voisinage \'etale élémentaire,
\item $\pi$ soit fini, et $\pi^{-1}(\{y'\}) = \{x'\}$,
\item $\kappa(y')$ soit une extension purement transcendante de $\kappa(s')$, et
\item $e : (S', s') \to (S, s)$ soit un voisinage \'etale élémentaire.
\end{enumerate}
De plus, si $f$ est localement de présentation finie, alors $\pi$ est
de présentation finie.
\end{lemma}

\begin{proof}
La question est locale sur $S$ ; on peut donc remplacer $S$ par un voisinage
ouvert affine de $s$. Appliquons ensuite le
lemme \ref{lemma-local-structure-finite-type}
pour obtenir un diagramme commutatif
$$
\xymatrix{
X \ar[dd] & X' \ar[l]^g \ar[d]^\pi & x \ar@{|->}[dd] &
x' \ar@{|->}[l] \ar@{|->}[d] \\
& Y \ar[d]^h & & y \ar@{|->}[d] \\
S \ar@{=}[r] & S & s & s \ar@{=}[l]
}
$$
où $h$ est lisse de dimension relative $n$ et $\kappa(y)$ est
une extension purement transcendante de $\kappa(s)$. Puisque la question est
aussi locale sur $X$, on peut remplacer $Y$ par un voisinage affine de $y$
(et $X'$ par son image réciproque par $\pi$). Comme $S$ est affine,
ceci assure que $Y \to S$ est quasi-compact, séparé et lisse,
en particulier de présentation finie.
Soit $T$ la composante connexe de $Y_s$ contenant $y$.
Comme $Y_s$ est noethérien, on voit que $T$ est ouvert.
On voit également que $T$ est géométriquement connexe sur $\kappa(s)$ d'après
Variétés,
lemme \ref{varieties-lemma-geometrically-connected-if-connected-and-point}.
Puisque $T$ est aussi lisse sur $\kappa(s)$, il est géométriquement normal, voir
Variétés, lemme \ref{varieties-lemma-smooth-geometrically-normal}.
On en conclut que $T$ est géométriquement irréductible sur $\kappa(s)$ (car un
schéma noethérien normal connexe est irréductible, voir
Propriétés, lemme \ref{properties-lemma-normal-Noetherian}).
Enfin, observons que le morphisme lisse $h$ est normal d'après le
lemme \ref{lemma-smooth-normal}.
Nous avons alors vérifié que toutes les hypothèses du
lemme \ref{lemma-normal-morphism-irreducible}
sont satisfaites pour le morphisme $h : Y \to S$ et l'ouvert $T \subset Y_s$.
En appliquant le
lemme \ref{lemma-normal-morphism-irreducible},
on obtient $e : S' \to S$, $s' \in S'$ et $Y'$ comme
dans le diagramme commutatif
$$
\xymatrix{
X \ar[dd] & X' \ar[l]^g \ar[d]^\pi & X' \times_Y Y' \ar[l] \ar[d] &
x \ar@{|->}[dd] & x' \ar@{|->}[l] \ar@{|->}[d] &
(x', s') \ar@{|->}[l] \ar@{|->}[d] \\
& Y \ar[d]^h & Y' \ar[d] \ar[l] & & y \ar@{|->}[d] &
(y, s') \ar@{|->}[l] \ar@{|->}[d] \\
S \ar@{=}[r] & S & S' \ar[l]_e & s & s \ar@{=}[l] & s' \ar@{|->}[l]
}
$$
où $e : (S', s') \to (S, s)$ est un voisinage \'etale élémentaire,
et où $Y' \subset Y_{S'}$ est un voisinage ouvert dont toutes les fibres
sur $S'$ sont géométriquement irréductibles, tel que $Y'_{s'} = T$ via
l'identification $Y_s = Y_{S', s'}$. Soit $(y, s') \in Y'$ le point
correspondant à $y \in T$ ; c'est aussi l'unique point de $Y \times_S S'$
au-dessus de $y$, de corps résiduel égal à $\kappa(y)$ et dont l'image est $s'$
dans $S'$. De même, soit $(x', s') \in X' \times_Y Y' \subset X' \times_S S'$
l'unique point au-dessus de $x'$, de corps résiduel égal à $\kappa(x')$
et situé au-dessus de $s'$. La partie extérieure de ce diagramme fournit alors une solution au
problème posé dans le lemme. Quelques détails mineurs sont omis.
\end{proof}

\begin{lemma}
\label{lemma-local-local-structure-finite-type-affine}
Hypothèses et notations comme dans le
lemme \ref{lemma-local-local-structure-finite-type}.
Outre les propriétés (1) -- (6), on peut imposer en plus que
\begin{enumerate}
\item[(7)] $S'$, $Y'$, $X'$ soient affines.
\end{enumerate}
\end{lemma}

\begin{proof}
Observons que si $Y'$ est affine, alors $X'$ est affine puisque $\pi$ est fini.
Choisissons un voisinage ouvert affine $U' \subset S'$ de $s'$.
Choisissons un voisinage ouvert affine $V' \subset h^{-1}(U')$ de $y'$.
Posons $W' = h(V')$. C'est un voisinage ouvert de $s'$ dans $S'$, voir
Morphismes, lemme \ref{morphisms-lemma-smooth-open},
contenu dans $U'$. Choisissons un voisinage ouvert affine $U'' \subset W'$
de $s'$. Alors $h^{-1}(U'') \cap V'$ est affine, car il est égal à
$U'' \times_{U'} V'$. Par construction,
$h^{-1}(U'') \cap V' \to U''$ est un morphisme lisse et surjectif dont les
fibres sont des sous-schémas ouverts non vides de fibres géométriquement intègres
de $Y' \to S'$, et sont donc géométriquement intègres. On peut ainsi remplacer
$S'$ par $U''$ et $Y'$ par $h^{-1}(U'') \cap V'$.
\end{proof}

\noindent
La portée de la propriété $\pi^{-1}(\{y'\}) = \{x'\}$ est en partie
expliquée par le lemme suivant.

\begin{lemma}
\label{lemma-finite-morphism-single-point-in-fibre}
Soit $\pi : X \to Y$ un morphisme fini.
Soit $x \in X$, avec $y = \pi(x)$, tel que $\pi^{-1}(\{y\}) = \{x\}$.
Alors
\begin{enumerate}
\item Pour tout voisinage $U \subset X$ de $x$ dans $X$, il
existe un voisinage $V \subset Y$ de $y$ tel que
$\pi^{-1}(V) \subset U$.
\item Le morphisme d'anneaux $\mathcal{O}_{Y, y} \to \mathcal{O}_{X, x}$
est fini.
\item Si $\pi$ est de présentation finie, alors
$\mathcal{O}_{Y, y} \to \mathcal{O}_{X, x}$ est de présentation finie.
\item Pour tout $\mathcal{O}_X$-module quasi-cohérent $\mathcal{F}$,
on a $\mathcal{F}_x = \pi_*\mathcal{F}_y$ comme
$\mathcal{O}_{Y, y}$-modules.
\end{enumerate}
\end{lemma}

\begin{proof}
La première assertion est purement topologique ; utiliser le fait que
$\pi$ est une application continue et fermée telle que $\pi^{-1}(\{y\}) = \{x\}$.
Pour démontrer les deuxième et troisième assertions, on peut supposer que
$X = \Spec(B)$ et $Y = \Spec(A)$. Alors
$A \to B$ est un morphisme fini d'anneaux et $y$ correspond à un idéal premier
$\mathfrak p$ de $A$ tel qu'il existe un unique idéal premier $\mathfrak q$ de
$B$ au-dessus de $\mathfrak p$. Alors
$B_{\mathfrak q} = B_{\mathfrak p}$, voir
Algèbre, lemme \ref{algebra-lemma-unique-prime-over-localize-below}.
Autrement dit, le morphisme $A_{\mathfrak p} \to B_{\mathfrak q}$
est égal au morphisme $A_{\mathfrak p} \to B_{\mathfrak p}$ que l'on obtient
en localisant $A \to B$ en $\mathfrak p$.
Ainsi, (2) et (3) résultent de propriétés élémentaires de la localisation
(certains détails sont omis). Pour la dernière assertion, supposons que
$\mathcal{F} = \widetilde M$ pour un $B$-module $M$.
Alors $\mathcal{F} = M_{\mathfrak q}$ et
$\pi_*\mathcal{F}_y = M_{\mathfrak p}$. D'après ce qui précède, ces
localisations coïncident. On peut aussi utiliser la partie (1) et
la définition des germes pour voir que $\mathcal{F}_x = \pi_*\mathcal{F}_y$
directement.
\end{proof}






\section{Application à la topologie fppf}
\label{section-application-fppf}

\noindent
On peut utiliser les techniques de localisation \'etale ci-dessus pour démontrer le
résultat suivant, qui décrit la topologie fppf comme la topologie
« engendrée par » les recouvrements de Zariski et les recouvrements de la
forme $\{f : T \to S\}$, où $f$ est fini localement libre et surjectif.

\begin{lemma}
\label{lemma-dominate-fppf-etale-locally}
Soit $S$ un schéma. Soit $\{S_i \to S\}_{i \in I}$ un recouvrement fppf.
Alors il existe
\begin{enumerate}
\item un recouvrement \'etale $\{S'_a \to S\}$,
\item des morphismes finis localement libres et surjectifs $V_a \to S'_a$,
\end{enumerate}
tels que le recouvrement fppf $\{V_a \to S\}$ raffine le
recouvrement donné $\{S_i \to S\}$.
\end{lemma}

\begin{proof}
On peut supposer que chaque $S_i \to S$ est localement quasi-fini, voir le
lemme \ref{lemma-qf-fp-flat-dominates-fppf}.

\medskip\noindent
Fixons un point $s \in S$. Choisissons $i \in I$ et un point
$s_i \in S_i$ au-dessus de $s$. Choisissons un voisinage \'etale élémentaire
$(S', s) \to (S, s)$ tel qu'il existe un ouvert
$$
S_i \times_S S'  \supset V
$$
qui contienne un unique point $v \in V$ au-dessus de $s \in S'$
et tel que $V \to S'$ soit fini, voir le
lemme \ref{lemma-etale-makes-quasi-finite-finite-at-point}.
Alors $V \to S'$ est fini localement libre, puisqu'il est fini
et que $S_i \times_S S' \to S'$ est plat et localement de présentation finie
comme changement de base du morphisme $S_i \to S$, voir
Morphismes, lemmes \ref{morphisms-lemma-base-change-finite-presentation},
\ref{morphisms-lemma-base-change-flat} et
\ref{morphisms-lemma-finite-flat}.
Ainsi $V \to S'$ est ouvert et, après avoir rétréci $S'$,
on peut supposer que $V \to S'$ est fini localement libre et surjectif.
Comme on peut procéder ainsi pour tout point de $S$, on conclut que
$\{S_i \to S\}$ peut être raffiné par un recouvrement de la forme
$\{V_a \to S\}_{a \in A}$ où chaque $V_a \to S$ se factorise en
$V_a \to S'_a \to S$, avec $S'_a \to S$ \'etale et $V_a \to S'_a$ fini
localement libre et surjectif.
\end{proof}

\begin{lemma}
\label{lemma-dominate-fppf}
Soit $S$ un schéma. Soit $\{S_i \to S\}_{i \in I}$ un recouvrement fppf.
Alors il existe
\begin{enumerate}
\item un recouvrement ouvert de Zariski $S = \bigcup U_j$,
\item des morphismes finis localement libres et surjectifs $W_j \to U_j$,
\item des recouvrements ouverts de Zariski $W_j = \bigcup_k W_{j, k}$,
\item des morphismes finis localement libres et surjectifs $T_{j, k} \to W_{j, k}$
\end{enumerate}
tels que le recouvrement fppf $\{T_{j, k} \to S\}$ raffine le
recouvrement donné $\{S_i \to S\}$.
\end{lemma}

\begin{proof}
Soit $\{V_a \to S\}_{a \in A}$ le recouvrement fppf obtenu dans le
lemme \ref{lemma-dominate-fppf-etale-locally}.
Autrement dit, ce recouvrement raffine
$\{S_i \to S\}$
et chaque $V_a \to S$ se factorise en
$V_a \to S'_a \to S$, avec $S'_a \to S$ \'etale et $V_a \to S'_a$
fini localement libre et surjectif.

\medskip\noindent
D'après la
remarque \ref{remark-topologies},
il existe un recouvrement ouvert de Zariski $S = \bigcup U_j$,
pour chaque $j$ un morphisme fini localement libre et surjectif
$W_j \to U_j$, et pour chaque $j$ un recouvrement ouvert de Zariski
$\{W_{j, k} \to W_j\}$ tel que la famille
$\{W_{j, k} \to S\}$ raffine le recouvrement \'etale
$\{S'_a \to S\}$, c'est-à-dire que, pour chaque paire $j, k$, il existe
un $a(j, k)$ et une factorisation $W_{j, k} \to S'_a \to S$
du morphisme $W_{j, k} \to S$. Posons
$T_{j, k} = W_{j, k} \times_{S'_a} V_a$ ; tout est alors clair.
\end{proof}

\begin{lemma}
\label{lemma-extend-integral-surjective-morphisms}
Soit $S$ un schéma. Si $U \subset S$ est ouvert et $V \to U$ est un morphisme
intégral surjectif, alors il existe un morphisme intégral surjectif
$\overline{V} \to S$ tel que $\overline{V} \times_S U$ soit
isomorphe à $V$ comme schéma sur $U$.
\end{lemma}

\begin{proof}
Soit $V' \to S$ la normalisation de $S$ dans $U$, voir
Morphismes, section \ref{morphisms-section-normalization-X-in-Y}.
Par construction, $V' \to S$ est intégral.
D'après Morphismes, lemmes \ref{morphisms-lemma-normalization-localization} et
\ref{morphisms-lemma-normalization-in-integral}, on voit que
l'image réciproque de $U$ dans $V'$ est $V$. Soit $Z$ la structure de sous-schéma réduit
induite sur $S \setminus U$. Alors
$\overline{V} = V' \amalg Z$ convient.
\end{proof}

\begin{lemma}
\label{lemma-extend-finite-surjective-morphisms}
Soit $S$ un schéma quasi-compact et quasi-séparé.
Si $U \subset S$ est un ouvert quasi-compact
et $V \to U$ un morphisme fini surjectif, alors il existe un
morphisme fini surjectif $\overline{V} \to S$ tel que $\overline{V} \times_S U$
soit isomorphe à $V$ comme schéma sur $U$.
\end{lemma}

\begin{proof}
D'après le théorème principal de Zariski
(lemme \ref{lemma-quasi-finite-separated-pass-through-finite}),
on peut supposer que $V$ est un ouvert quasi-compact d'un schéma $V'$
fini sur $S$. Après avoir remplacé $V'$ par l'image schématique
de $V$, on peut supposer que $V$ est dense dans $V'$.
Il s'ensuit que $V' \times_S U = V$, car $V \to V' \times_S U$
est fermé puisque $V$ est fini sur $U$. Soit $Z$ la structure de sous-schéma réduit
induite sur $S \setminus U$. Alors
$\overline{V} = V' \amalg Z$ convient.
\end{proof}

\begin{lemma}
\label{lemma-dominate-fppf-integral}
Soit $S$ un schéma. Soit $\{S_i \to S\}_{i \in I}$ un recouvrement fppf.
Alors il existe un morphisme intégral surjectif $S' \to S$ et un
recouvrement ouvert $S' = \bigcup U'_\alpha$ tels que, pour tout $\alpha$, le
morphisme $U'_\alpha \to S$ se factorise par $S_i \to S$ pour un certain $i$.
\end{lemma}

\begin{proof}
Choisissons $S = \bigcup U_j$, $W_j \to U_j$, $W_j = \bigcup W_{j, k}$ et
$T_{j, k} \to W_{j, k}$ comme dans le lemme \ref{lemma-dominate-fppf}.
D'après le lemme \ref{lemma-extend-integral-surjective-morphisms},
on peut prolonger $W_j \to U_j$ en un morphisme intégral surjectif
$\overline{W}_j \to S$.
On peut ensuite prolonger $T_{j, k} \to W_{j, k}$ en un morphisme
intégral surjectif $\overline{T}_{j, k} \to \overline{W}_j$.
On définit $\overline{T}_j$ comme le produit de tous les schémas
$\overline{T}_{j, k}$ sur $\overline{W}_j$
(Limites, lemme \ref{limits-lemma-infinite-product}).
Puis on définit $S'$ comme le produit de tous les schémas
$\overline{T}_j$ sur $S$.
Si $x \in S'$, il existe un $j$ tel que l'image de $x$ dans
$S$ appartienne à $U_j$. Il existe donc un $k$ tel que l'image de
$x$ par la projection $S' \to \overline{W}_j$ appartienne à $W_{j, k}$.
Ainsi, par la projection $S' \to \overline{T}_j \to \overline{T}_{j, k}$,
l'image du point $x$ appartient à $T_{j, k}$. Or $T_{j, k} \to S$
se factorise par $S_i$ pour un certain $i$.
Enfin, le morphisme $S' \to S$ est intégral et surjectif
d'après Limites, lemmes \ref{limits-lemma-infinite-product-integral} et
\ref{limits-lemma-infinite-product-surjective}.
\end{proof}

\begin{lemma}
\label{lemma-dominate-fppf-finite}
Soit $S$ un schéma quasi-compact et quasi-séparé.
Soit $\{S_i \to S\}_{i \in I}$ un recouvrement fppf.
Alors il existe un morphisme fini surjectif $S' \to S$
de présentation finie et un
recouvrement ouvert $S' = \bigcup U'_\alpha$ tel que, pour tout $\alpha$, le
morphisme $U'_\alpha \to S$ se factorise par $S_i \to S$ pour un certain $i$.
\end{lemma}

\begin{proof}
Soit $Y \to X$ le morphisme intégral surjectif obtenu dans le
lemme \ref{lemma-dominate-fppf-integral}.
Choisissons un recouvrement ouvert affine fini $Y = \bigcup V_j$
tel que $V_j \to X$ se factorise par $S_{i(j)}$.
On peut écrire $Y = \lim Y_\lambda$, avec
$Y_\lambda \to X$ fini et de présentation finie, voir
Limites, lemme \ref{limits-lemma-integral-limit-finite-and-finite-presentation}.
Pour $\lambda$ assez grand, on peut trouver des ouverts affines
$V_{\lambda, j} \subset Y_\lambda$
dont les images réciproques dans $Y$ redonnent les $V_j$, voir
Limites, lemme \ref{limits-lemma-descend-opens}.
Pour $\lambda$ encore plus grand, les morphismes $V_j \to S_{i(j)}$
sur $X$ proviennent de morphismes $V_{\lambda, j} \to S_{i(j)}$ sur
$X$, voir
Limites, proposition
\ref{limits-proposition-characterize-locally-finite-presentation}.
Poser $S' = Y_\lambda$ pour ce $\lambda$ achève la démonstration.
\end{proof}

\begin{lemma}
\label{lemma-fppf-ph}
Un recouvrement fppf de schémas est un recouvrement ph.
\end{lemma}

\begin{proof}
Soit $\{T_i \to T\}$ un recouvrement fppf de schémas, voir
Topologies, définition \ref{topologies-definition-fppf-covering}.
Observons que $T_i \to T$ est localement de type fini.
Soit $U \subset T$ un ouvert affine.
Il suffit de montrer que $\{T_i \times_T U \to U\}$
peut être raffiné par un recouvrement ph standard, voir
Topologies, définition \ref{topologies-definition-ph-covering}.
Ceci résulte immédiatement du lemme \ref{lemma-dominate-fppf-finite}
et du fait qu'un morphisme fini est propre
(Morphismes, lemme \ref{morphisms-lemma-finite-proper}).
\end{proof}

\begin{remark}
\label{remark-change-topologies}
Comme conséquence du lemme \ref{lemma-fppf-ph}, on obtient un morphisme de comparaison
$$
\epsilon : (\Sch/S)_{ph} \longrightarrow (\Sch/S)_{fppf}
$$
C'est le morphisme de sites donné par le foncteur identité
sur les catégories sous-jacentes (avec des choix convenables de sites
comme dans Topologies, remarque \ref{topologies-remark-choice-sites}).
Le foncteur $\epsilon_*$ est l'identité sur les préfaisceaux sous-jacents
et le foncteur $\epsilon^{-1}$ associe à un faisceau fppf
son faisceau ph associé.
Par composition, on peut en outre comparer la topologie ph
aux topologies syntomique, lisse, \'etale et de Zariski.
\end{remark}




\section{Schémas quasi-projectifs}
\label{section-quasi-projective}

\noindent
Le terme « schéma quasi-projectif » n'a pas encore été défini.
Une définition possible serait : un schéma qui possède un faisceau inversible
ample. Toutefois, si $X$ est un schéma sur un schéma de base $S$, alors
on dit que {\it $X$ est quasi-projectif sur $S$} si le morphisme
$X \to S$ est quasi-projectif
(Morphismes, définition \ref{morphisms-definition-quasi-projective}).
Comme le morphisme identité de tout schéma est quasi-projectif, on
voit qu'un schéma quasi-projectif sur $S$ ne possède pas nécessairement
un faisceau inversible ample. Pour cette raison, il semble préférable de
ne pas définir le terme « schéma quasi-projectif ».

\begin{lemma}
\label{lemma-quasi-projective}
Soit $S$ un schéma qui possède un faisceau inversible ample.
Soit $f : X \to S$ un morphisme de schémas. Les conditions suivantes sont
équivalentes :
\begin{enumerate}
\item $X \to S$ est quasi-projectif,
\item $X \to S$ est H-quasi-projectif,
\item il existe une immersion ouverte quasi-compacte $X \to X'$ de schémas
sur $S$ telle que $X' \to S$ soit projectif,
\item $X \to S$ est de type fini et $X$ possède un faisceau inversible
ample, et
\item $X \to S$ est de type fini et il existe un
faisceau inversible $f$-très ample.
\end{enumerate}
\end{lemma}

\begin{proof}
L'implication (2) $\Rightarrow$ (1) est donnée par
Morphismes, lemme \ref{morphisms-lemma-H-quasi-projective-quasi-projective}.
L'implication (1) $\Rightarrow$ (2) est donnée par
Morphismes, lemme
\ref{morphisms-lemma-projective-over-quasi-projective-is-H-projective}.
L'implication (2) $\Rightarrow$ (3) est donnée par
Morphismes, lemme \ref{morphisms-lemma-H-quasi-projective-open-H-projective}.

\medskip\noindent
Supposons que $X \subset X'$ soit comme en (3). En particulier, $X \to S$ est
de type fini. D'après
Morphismes, lemme \ref{morphisms-lemma-H-quasi-projective-open-H-projective},
le morphisme $X \to S$ est H-projectif.
Il existe donc une immersion quasi-compacte $i : X \to \mathbf{P}^n_S$.
Ainsi $\mathcal{L} = i^*\mathcal{O}_{\mathbf{P}^n_S}(1)$
est $f$-très ample. Comme $X \to S$ est quasi-compact, on déduit de
Morphismes, lemme \ref{morphisms-lemma-ample-very-ample},
que $\mathcal{L}$ est $f$-ample. Ainsi $X \to S$ est quasi-projectif
par définition.

\medskip\noindent
L'implication (4) $\Rightarrow$ (2) est donnée par
Morphismes, lemme \ref{morphisms-lemma-quasi-projective-finite-type-over-S}.

\medskip\noindent
L'implication (5) $\Rightarrow$ (4) résulte de
Morphismes, lemme \ref{morphisms-lemma-pullback-ample-tensor-relatively-ample}.

\medskip\noindent
L'équivalence de (1) et (5) résulte de
Morphismes, lemmes \ref{morphisms-lemma-ample-very-ample} et
\ref{morphisms-lemma-finite-type-ample-very-ample}.
\end{proof}

\begin{lemma}
\label{lemma-category-quasi-projective}
Soit $S$ un schéma qui possède un faisceau inversible ample.
Soit $\text{QP}_S$ la sous-catégorie pleine de la
catégorie des schémas sur $S$ qui satisfont aux conditions
équivalentes du lemme \ref{lemma-quasi-projective}.
\begin{enumerate}
\item si $S' \to S$ est un morphisme de schémas et que $S'$ possède
un faisceau inversible ample, alors le changement de base définit
un foncteur $\text{QP}_S \to \text{QP}_{S'}$,
\item si $X \in \text{QP}_S$ et $Y \in \text{QP}_X$, alors $Y \in \text{QP}_S$,
\item la catégorie $\text{QP}_S$ est stable par produits fibrés,
\item la catégorie $\text{QP}_S$ est stable par
réunions disjointes finies,
\item si $X \to S$ est projectif, alors $X \in \text{QP}_S$,
\item si $X \to S$ est quasi-affine de type fini, alors
$X$ appartient à $\text{QP}_S$,
\item si $X \to S$ est quasi-fini et séparé, alors
$X \in \text{QP}_S$,
\item si $X \to S$ est une immersion quasi-compacte, alors
$X \in \text{QP}_S$,
\item ajouter d'autres assertions ici.
\end{enumerate}
\end{lemma}

\begin{proof}
L'assertion (1) résulte de Morphismes, lemme
\ref{morphisms-lemma-base-change-quasi-projective}.

\medskip\noindent
L'assertion (2) résulte de la quatrième caractérisation du
lemme \ref{lemma-quasi-projective}.

\medskip\noindent
Si $X \to S$ et $Y \to S$ sont quasi-projectifs, alors
$X \times_S Y \to Y$ est quasi-projectif d'après
Morphismes, lemme \ref{morphisms-lemma-base-change-quasi-projective}.
Ainsi, (3) résulte de (2).

\medskip\noindent
Si $X = Y \amalg Z$ est une réunion disjointe de schémas
et si $\mathcal{L}$ est un $\mathcal{O}_X$-module inversible
tel que $\mathcal{L}|_Y$ et $\mathcal{L}|_Z$ soient amples, alors
$\mathcal{L}$ est ample (détails omis). Ainsi,
l'assertion (4) résulte de la quatrième caractérisation du
lemme \ref{lemma-quasi-projective}.

\medskip\noindent
L'assertion (5) résulte de
Morphismes, lemme \ref{morphisms-lemma-projective-quasi-projective}.

\medskip\noindent
L'assertion (6) résulte de
Morphismes, lemme
\ref{morphisms-lemma-quasi-affine-finite-type-quasi-projective}.

\medskip\noindent
L'assertion (7) résulte de l'assertion (6) et du
lemme \ref{lemma-quasi-finite-separated-quasi-affine}.

\medskip\noindent
L'assertion (8) résulte de l'assertion (7) et de
Morphismes, lemme \ref{morphisms-lemma-immersion-locally-quasi-finite}.
\end{proof}

\noindent
Le lemme suivant n'a pas vraiment sa place dans cette section, mais
il ne semble exister aucun meilleur endroit où le placer.

\begin{lemma}
\label{lemma-integral-over-quasi-affine}
Soit $X$ un schéma quasi-affine. Soit $f : U \to X$ un morphisme
entier. Alors $U$ est quasi-affine et le diagramme
$$
\xymatrix{
U \ar[r] \ar[d] & \Spec(\Gamma(U, \mathcal{O}_U)) \ar[d] \\
X \ar[r] & \Spec(\Gamma(X, \mathcal{O}_X))
}
$$
est cartésien.
\end{lemma}

\begin{proof}
Le schéma $U$ est quasi-affine parce que les morphismes entiers sont affines,
les morphismes affines sont quasi-affines, un schéma est quasi-affine si et seulement si
le morphisme structural vers $\Spec(\mathbf{Z})$ est quasi-affine, et que
les composés de morphismes quasi-affines sont quasi-affines.
Les deux premières assertions résultent immédiatement de la définition,
et la troisième résulte de
Morphismes, lemme \ref{morphisms-lemma-composition-quasi-affine}.
Posons $U' =
X \times_{\Spec(\Gamma(X, \mathcal{O}_X))} \Spec(\Gamma(U, \mathcal{O}_U))$
et considérons le diagramme prolongé
$$
\xymatrix{
U \ar[r]_j \ar[rd] & U' \ar[d] \ar[r] &
\Spec(\Gamma(U, \mathcal{O}_U)) \ar[d] \\
& X \ar[r] & \Spec(\Gamma(X, \mathcal{O}_X))
}
$$
Le morphisme $j$ est fermé en vertu de
Morphismes, lemme \ref{morphisms-lemma-image-proper-scheme-closed},
ainsi que du fait qu'un morphisme entier est universellement fermé
(Morphismes, lemme \ref{morphisms-lemma-integral-universally-closed}),
et du fait que les flèches verticales du diagramme sont séparées.
D'autre part, $j$ est ouvert parce que les flèches
horizontales du diagramme du lemme sont ouvertes d'après
Propriétés, lemme \ref{properties-lemma-quasi-affine}.
Ainsi, $j$ identifie $U$ à un sous-schéma ouvert et fermé de $U'$.
Si $U \not = U'$, alors $U$ n'est pas dense dans $U'$ et, a fortiori,
n'est pas dense dans le spectre de $\Gamma(U, \mathcal{O}_U)$.
Cependant, l'image schématique de
$U$ dans $\Spec(\Gamma(U, \mathcal{O}_U))$ est $\Spec(\Gamma(U, \mathcal{O}_U))$,
car tout idéal de $\Gamma(U, \mathcal{O}_U)$
définissant un sous-schéma fermé par lequel $U$
se factoriserait devrait être nul.
Par conséquent, $U$ est dense dans $\Spec(\Gamma(U, \mathcal{O}_U))$, par exemple d'après
Morphismes, lemme \ref{morphisms-lemma-quasi-compact-scheme-theoretic-image}.
Ainsi, $U = U'$, ce qui conclut.
\end{proof}









\section{Schémas projectifs}
\label{section-projective}

\noindent
Cette section est l'analogue de la section \ref{section-quasi-projective}
pour les morphismes projectifs.

\begin{lemma}
\label{lemma-projective}
Soit $S$ un schéma qui possède un faisceau inversible ample.
Soit $f : X \to S$ un morphisme de schémas. Les conditions suivantes sont
équivalentes :
\begin{enumerate}
\item $X \to S$ est projectif,
\item $X \to S$ est H-projectif,
\item $X \to S$ est quasi-projectif et propre,
\item $X \to S$ est H-quasi-projectif et propre,
\item $X \to S$ est propre et $X$ possède un faisceau inversible ample,
\item $X \to S$ est propre et il existe un faisceau inversible $f$-ample,
\item $X \to S$ est propre et il existe un faisceau inversible $f$-très ample,
\item il existe une $\mathcal{O}_S$-algèbre graduée quasi-cohérente $\mathcal{A}$
engendrée par $\mathcal{A}_1$ sur $\mathcal{A}_0$, où $\mathcal{A}_1$ est un
$\mathcal{O}_S$-module de type fini, telle que
$X = \underline{\text{Proj}}_S(\mathcal{A})$.
\end{enumerate}
\end{lemma}

\begin{proof}
Observons d'abord que, dans chacun des cas, le morphisme $f$ est propre, voir
Morphismes, lemmes \ref{morphisms-lemma-H-projective} et
\ref{morphisms-lemma-locally-projective-proper}.
Il suffit donc de démontrer l'équivalence des notions dans le
cas où $f$ est un morphisme propre. Nous l'utiliserons dans la suite sans
autre mention.

\medskip\noindent
Les équivalences (1) $\Leftrightarrow$ (3) et
(2) $\Leftrightarrow$ (4) sont données par
Morphismes, lemme \ref{morphisms-lemma-projective-is-quasi-projective-proper}.

\medskip\noindent
L'implication (2) $\Rightarrow$ (1) est donnée par
Morphismes, lemme \ref{morphisms-lemma-H-projective}.

\medskip\noindent
Les implications (1) $\Rightarrow$ (2) et (3) $\Rightarrow$ (4) sont données par
Morphismes, lemme
\ref{morphisms-lemma-projective-over-quasi-projective-is-H-projective}.

\medskip\noindent
L'implication (1) $\Rightarrow$ (7) résulte immédiatement de
Morphismes, définitions \ref{morphisms-definition-projective} et
\ref{morphisms-definition-very-ample}.

\medskip\noindent
Les conditions (3) et (6) sont équivalentes d'après
Morphismes, définition \ref{morphisms-definition-quasi-projective}.

\medskip\noindent
Ainsi, (1) -- (4) et (6) sont équivalentes et entraînent (7). D'après le
lemme \ref{lemma-quasi-projective},
les conditions (3), (5) et (7) sont équivalentes.
On voit donc que (1) -- (7) sont équivalentes.

\medskip\noindent
D'après Diviseurs, lemme \ref{divisors-lemma-relative-proj-projective},
on voit que (8) entraîne (1). Réciproquement, si (2) est satisfaite, alors
on peut choisir une immersion fermée
$$
i :
X
\longrightarrow
\mathbf{P}^n_S = \underline{\text{Proj}}_S(\mathcal{O}_S[T_0, \ldots, T_n]).
$$
Voir Constructions, lemme \ref{constructions-lemma-projective-space-bundle}
pour cette égalité. D'après
Diviseurs, lemme \ref{divisors-lemma-closed-subscheme-proj},
on voit que $X$ est le Proj relatif d'une algèbre quotient graduée
quasi-cohérente $\mathcal{A}$ de $\mathcal{O}_S[T_0, \ldots, T_n]$.
Alors $\mathcal{A}$ satisfait aux conditions de (8).
\end{proof}

\begin{lemma}
\label{lemma-category-projective}
Soit $S$ un schéma qui possède un faisceau inversible ample.
Soit $\text{P}_S$ la sous-catégorie pleine de la
catégorie des schémas sur $S$ qui satisfont aux conditions
équivalentes du lemme \ref{lemma-projective}.
\begin{enumerate}
\item si $S' \to S$ est un morphisme de schémas et que $S'$ possède
un faisceau inversible ample, alors le changement de base définit
un foncteur $\text{P}_S \to \text{P}_{S'}$,
\item si $X \in \text{P}_S$ et $Y \in \text{P}_X$, alors $Y \in \text{P}_S$,
\item la catégorie $\text{P}_S$ est stable par produits fibrés,
\item la catégorie $\text{P}_S$ est stable par
réunions disjointes finies,
\item si $X \to S$ est fini, alors $X$ appartient à $\text{P}_S$,
\item ajouter d'autres assertions ici.
\end{enumerate}
\end{lemma}

\begin{proof}
L'assertion (1) résulte de Morphismes, lemme
\ref{morphisms-lemma-base-change-projective}.

\medskip\noindent
L'assertion (2) résulte de la cinquième caractérisation du
lemme \ref{lemma-projective} et du fait que les composés
de morphismes propres sont propres
(Morphismes, lemme \ref{morphisms-lemma-composition-proper}).

\medskip\noindent
Si $X \to S$ et $Y \to S$ sont projectifs, alors
$X \times_S Y \to Y$ est projectif d'après
Morphismes, lemme \ref{morphisms-lemma-base-change-projective}.
Ainsi, (3) résulte de (2).

\medskip\noindent
Si $X = Y \amalg Z$ est une réunion disjointe de schémas
et si $\mathcal{L}$ est un $\mathcal{O}_X$-module inversible
tel que $\mathcal{L}|_Y$ et $\mathcal{L}|_Z$ soient amples, alors
$\mathcal{L}$ est ample (détails omis). Ainsi,
l'assertion (4) résulte de la cinquième caractérisation du
lemme \ref{lemma-projective}.

\medskip\noindent
L'assertion (5) résulte de
Morphismes, lemme \ref{morphisms-lemma-finite-projective}.
\end{proof}

\noindent
Voici un résultat d'un type légèrement différent.

\begin{lemma}
\label{lemma-ample-in-neighbourhood}
\begin{reference}
\cite[IV Corollaire 9.6.4]{EGA}
\end{reference}
Soit $f : X \to Y$ un morphisme propre de schémas.
Soit $\mathcal{L}$ un $\mathcal{O}_X$-module inversible.
Soit $y \in Y$ un point tel que $\mathcal{L}_y$ soit ample
sur $X_y$. Alors il existe un voisinage ouvert $V \subset Y$
de $y$ tel que $\mathcal{L}|_{f^{-1}(V)}$ soit ample sur $f^{-1}(V)/V$.
\end{lemma}

\begin{proof}
On peut supposer $Y$ affine. On trouve alors un ensemble ordonné filtrant $I$
et un système projectif de morphismes $X_i \to Y_i$ de schémas,
où $Y_i$ est de type fini sur $\mathbf{Z}$, dont les
morphismes de transition $X_i \to X_{i'}$ et $Y_i \to Y_{i'}$ sont affines,
où $X_i \to Y_i$ est propre, tels que $X \to Y = \lim (X_i \to Y_i)$.
Voir Limites, lemme
\ref{limits-lemma-proper-limit-of-proper-finite-presentation-noetherian}.
Après avoir restreint $I$, on peut supposer disposer d'un système compatible de
$\mathcal{O}_{X_i}$-modules inversibles $\mathcal{L}_i$
dont les images réciproques donnent $\mathcal{L}$, voir
Limites, lemme \ref{limits-lemma-descend-invertible-modules}.
Soit $y_i \in Y_i$ l'image de $y$.
Alors $\kappa(y) = \colim \kappa(y_i)$.
Par conséquent, pour un certain $i$, $\mathcal{L}_{i, y_i}$
est ample sur $X_{i, y_i}$ d'après
Limites, lemme \ref{limits-lemma-limit-ample}.
D'après Cohomologie des schémas, lemme \ref{coherent-lemma-ample-in-neighbourhood},
on trouve un voisinage ouvert
$V_i \subset Y_i$ de $y_i$ tel que
la restriction de $\mathcal{L}_i$ à $f_i^{-1}(V_i)$
soit ample relativement à $V_i$.
En prenant pour $V \subset Y$ l'image réciproque de
$V_i$, on achève la démonstration (indications : utiliser
Morphismes, lemme \ref{morphisms-lemma-ample-base-change},
le fait que $X \to Y \times_{Y_i} X_i$ est affine,
et le fait que l'image réciproque d'un
faisceau inversible ample par un morphisme affine est ample d'après
Morphismes, lemme \ref{morphisms-lemma-pullback-ample-tensor-relatively-ample}).
\end{proof}







\section{Proj et Spec}
\label{section-proj-spec}

\noindent
Dans cette section, nous précisons la relation entre le Proj et le spectre
d'un anneau gradué.

\medskip\noindent
Soit $R$ un anneau. Soit $A$ une $R$-algèbre graduée, voir
Algèbre, section \ref{algebra-section-graded}.
Pour $m \geq 0$, nous notons $A_{\geq m} = \bigoplus_{d \geq m} A_d$.
Considérons l'anneau gradué
$$
B = \bigoplus\nolimits_{d \geq 0} A_{\geq d}
$$
Pour $d' \geq d$ et $a \in A_{d'}$, notons $a^{(d)} \in B$
l'élément de $B_d$ correspondant à $a$.
Notons $\sigma : A \to B$ et $\psi : A \to B$
les deux homomorphismes évidents d'anneaux : si $a \in A_d$, alors
$\sigma(a) = a^{(0)}$ et $\psi(a) = a^{(d)}$.
Alors $\psi$ est un homomorphisme d'anneaux gradués et $\sigma$ munit $B$
d'une structure d'algèbre graduée sur $A$. Il existe aussi un
homomorphisme surjectif d'anneaux gradués $\tau : B \to A$
qui, pour $d' \geq d$ et $a \in A_{d'}$, envoie
$a^{(d)}$ sur $0$ si $d' > d$, et sur $a$ si $d' = d$.

\medskip\noindent
Schémas affines et spectres.
Posons $X = \Spec(A)$. L'idéal non pertinent $A_+$ définit un sous-schéma fermé
$Z = V(A_+) = \Spec(A/A_+) = \Spec(A_0)$. Posons $U = X \setminus Z$.
$$
U \longrightarrow X \longrightarrow Z
$$
Schémas projectifs et Proj. Posons $P = \text{Proj}(A)$. Nous pouvons
et allons considérer $P$ comme un schéma sur $\Spec(A_0) = Z$.
Posons $L = \text{Proj}(B)$. Nous pouvons et allons considérer $L$ comme un schéma
sur $\Spec(B_0) = \Spec(A) = X$ ; observons que l'identification
de $B_0$ avec $A$ est donnée par $\sigma$.
La surjection $\tau$ définit une immersion fermée $0 : P \to L$.
Comme $A \xrightarrow{\sigma} B \to A$ est égal à l'homomorphisme $A \to A_0 \to A$,
on conclut que
$$
\xymatrix{
P \ar[d] \ar[r]_0 & L \ar[d] \\
Z \ar[r] & X
}
$$
est commutatif.

\medskip\noindent
Nous affirmons que $\psi$ définit un morphisme $L \to P$.
Pour le voir, d'après Constructions, lemme \ref{constructions-lemma-morphism-proj},
il suffit de vérifier que $\psi(A_+) \not \subset \mathfrak p$ pour
tout idéal premier homogène
$\mathfrak p \subset B$ tel que $B_+ \not \subset \mathfrak p$.
Choisissons d'abord $g \in B_+$ homogène tel que $g \not \in \mathfrak p$.
On peut alors écrire $g$ comme une somme finie $g = \sum a_i^{(d)}$
avec $a_i \in A_{d_i}$ pour certains $d_i \geq d$.
On conclut qu'il existe $d' \geq d$ et $a \in A_{d'}$
tels que $a^{(d)} \not \in \mathfrak p$.
Alors
$$
(a^{(d)})^{d'} =
(a^{d'})^{(d'd)} =
a^{(d)} (a^{d' - 1})^{(d(d' - 1))} =
\psi(a) (a^{d' - 1})^{(d(d' - 1))}
$$
(la notation laisse quelque peu à désirer) n'appartient pas à $\mathfrak p$.
Ainsi $\psi(a) \not \in \mathfrak p$, ce qui démontre l'assertion.
On peut donc prolonger notre diagramme ci-dessus en un diagramme commutatif
$$
\xymatrix{
P \ar[d] \ar[r]_0 & L \ar[d] \ar[r]_\pi & P \ar[d] \\
Z \ar[r] & X \ar[r] & Z
}
$$
où $X \to Z$ est donné par $A_0 \to A$.
Comme $\tau \circ \psi = \text{id}_A$, on voit que $\pi \circ 0 = \text{id}_P$.

\medskip\noindent
Observons que $\pi$ est un morphisme affine. Cela ressort clairement de la construction
donnée dans Constructions, lemme \ref{constructions-lemma-morphism-proj}. En effet, si
$f \in A_d$ pour un certain $d > 0$, alors, en
posant $g = \psi(f)$, on a $\pi^{-1}(D_+(f)) = D_+(g)$.
Dans ce cas, on a l'égalité suivante entre composantes homogènes :
$$
(B[1/g])_{m'} = \bigoplus\nolimits_{m \geq m'} (A[1/f])_m
$$
Cet isomorphisme est compatible avec les localisations ultérieures.
En prenant $m' = 0$, on voit que $\pi_*\mathcal{O}_L$ est la
somme directe des $\mathcal{O}_P(m)$ pour $m \geq 0$\footnote{On obtient de même
que $\pi_*\mathcal{O}_L(i) = \bigoplus_{m \geq -i} \mathcal{O}_P(m)$.}.
On en conclut que $L$ s'identifie au spectre relatif :
$$
L = \underline{\Spec}_P
\left(
\bigoplus\nolimits_{m \geq 0} \mathcal{O}_P(m)
\right)
$$
En particulier, $L \to P$ est un cône\footnote{Souvent, $L$ est un fibré en droites
sur $P$, voir ci-dessous.}, voir
Constructions, section \ref{constructions-section-cone}.
De plus, il est clair que $0 : P \to L$ est le sommet du cône.

\medskip\noindent
Soient $f \in A_d$ pour un certain $d > 0$ et $g = \psi(f) \in B_d$,
comme au paragraphe précédent. En examinant la structure des homomorphismes d'anneaux
$$
\xymatrix{
A_0 \ar[r] \ar[d] & A \ar[d]^\sigma \ar[r] & A_0 \ar[d] \\
(A[1/f])_0 \ar[r]^-\psi &
(B[1/g])_0 = \bigoplus\nolimits_{m \geq 0} (A[1/f])_m
\ar[r]^-\tau &
(A[1/f])_0
}
$$
quelques calculs\footnote{Les assertions (1) et (2) sont claires.
Pour voir (3), notons que si $a \in A_d$, alors
$\sigma(a) = \sigma(f) \psi(a/f)$. Pour (4), notons que
$b/g^m$ appartient au noyau de $\tau$ si et seulement si
$b \in A_{\geq md}$ s'envoie sur zéro dans $A_{md}$.
Il suffit donc de montrer que, si $m' > md$ et $a \in A_{m'}$,
alors une certaine puissance de $a^{(md)}/g^m$ appartient à l'idéal engendré par
$\sigma(f)$. Prenons $e$ tel que $em' - emd \geq d$. Alors
$$
(a^{(md)}/g^m)^e = (a^e)^{(emd)}/g^{em} =
(fa^e)^{(emd + d)}/g^{em + 1} = \sigma(f) \cdot (a^e)^{(emd + d)}/g^{em + 1}
$$
comme voulu (toutes nos excuses pour cette notation épouvantable). Pour voir (5), raisonnons comme
précédemment et notons que $a^{(md)}/g^m = \sigma(f) \cdot a^{(md + 1)}/g^{m + 1}$
si $d = 1$.}
dans les anneaux gradués montrent que
\begin{enumerate}
\item $\sigma(A_+)(B[1/g])_0 \subset \Ker(\tau : (B[1/g])_0 \to (A[1/f])_0)$,
\item $\sigma(f) \in (B[1/g])_0$ est un élément non diviseur de zéro,
\item $\sigma(f) (B[1/g])_0 = \sigma(A_d) (B[1/g])_0$ en tant qu'idéaux,
\item $\sigma(f) (B[1/g])_0$ et
$\Ker(\tau : (B[1/g])_0 \to (A[1/f])_0)$ ont le même radical,
\item si $d = 1$, alors
$\sigma(f) (B[1/g])_0 = \Ker(\tau : (B[1/g])_0 \to (A[1/f])_0)$.
\end{enumerate}
On voit en particulier que
$$
0(D_+(f)) = V(\sigma(f)) \subset D_+(g) = \Spec((B[1/g])_0)
$$
ensemblistement. En d'autres termes, l'idéal engendré par
$\sigma(A_d)$ définit un diviseur de Cartier effectif sur
$D_+(g)$, dont l'ensemble sous-jacent est
égal à l'image de l'immersion fermée $0 : P \to L$.

\medskip\noindent
Nous affirmons que $L \to X$ est un isomorphisme au-dessus de $U$.
En effet, si $f \in A_d$ pour un certain $d > 0$, alors
$$
\Spec(A_f) \times_X L = \text{Proj}(A_f \otimes_A B) =
\text{Proj}(B_{\sigma(f)})
$$
Pour tout $e$, on a
$(B_{\sigma(f)})_e = A_f \otimes_B B_e = A_f \otimes_A A_{\geq e} = A_f$,
la dernière égalité étant induite par l'injection $A_{\geq e} \subset A$.
Par conséquent, $B_{\sigma(f)} \cong A_f[T]$, où $T$ est de degré $1$.
Cela démontre l'assertion, puisque $\text{Proj}(A_f[T]) \to \Spec(A_f)$
est un isomorphisme. Dorénavant, nous identifions $U$ à l'ouvert correspondant
de $L$.

\medskip\noindent
L'identification effectuée au paragraphe précédent nous permet de
considérer la restriction $\pi|_U : U \to P$. Choisissons
$f \in A_d$ pour un certain $d > 0$ et $g = \psi(f) \in B_d$,
comme nous l'avons fait plusieurs fois ci-dessus. Alors
$$
U \cap \pi^{-1}(D_+(f)) = U \cap D_+(g)
$$
est le complémentaire du lieu des zéros de $\sigma(f) \in (B[1/g])_0$
via l'identification de $D_+(g)$ avec le spectre de $(B[1/g])_0$.
C'est l'assertion (4) ci-dessus. Par conséquent,
$U \cap D_+(g)$ est affine et
$$
\mathcal{O}_L(U \cap D_+(g)) = (B[1/g])_0[1/\sigma(f)] =
\bigoplus\nolimits_{m \in \mathbf{Z}} (A[1/f])_m
$$
où le dernier signe d'égalité est le prolongement naturel de
l'identification $(B[1/g])_0 = \bigoplus_{m \geq 0} (A[1/f])_m$
effectuée ci-dessus. Exactement comme précédemment pour $\pi : L \to P$,
on conclut que $\pi|_U : U \to P$ est affine et que
$$
U = \underline{\Spec}_P
\left(
\bigoplus\nolimits_{m \in \mathbf{Z}} \mathcal{O}_P(m)
\right)
$$
comme schémas sur $P$.

\medskip\noindent
En résumé, nos constructions donnent un diagramme commutatif
\begin{equation}
\label{equation-proj-and-spec}
\vcenter{
\xymatrix{
\underline{\Spec}_P \left(
\bigoplus\nolimits_{m \in \mathbf{Z}} \mathcal{O}_P(m)
\right)
\ar[r] \ar@{=}[d] &
L =
\underline{\Spec}_P \left(
\bigoplus\nolimits_{m \geq 0} \mathcal{O}_P(m)
\right) \ar[d]^\sigma \ar[r]_-\pi & P \ar[d] \\
U \ar[r] & X \ar[r] & Z
}
}
\end{equation}
de schémas, où $\pi$ est un cône dont la section nulle $0 : P \to L$
a pour image ensembliste l'image réciproque de $Z$ dans $L$.

\medskip\noindent
Soit $W \subset P$ le plus grand ouvert tel que $\mathcal{O}_P(1)|_W$
soit inversible et que les morphismes naturels induisent des isomorphismes
$\mathcal{O}_P(m)|_W \cong \mathcal{O}_P(1)^{\otimes m}|_W$
pour tout $m \in \mathbf{Z}$, c'est-à-dire l'ouvert de
Constructions, lemme \ref{constructions-lemma-where-invertible} pour $d = 1$.
On voit alors que $L|_W = \pi^{-1}(W) \to W$ est un fibré vectoriel
(Constructions, section \ref{constructions-section-vector-bundle})
de rang $1$, à savoir
$$
L|_W = \mathbf{V}(\mathcal{O}_P(1)|_W)
$$
selon la notation grothendieckienne. Cela résulte immédiatement de ce qui précède, qui montre
que $L|_W$ est égal au spectre relatif de l'algèbre symétrique du
$\mathcal{O}_W$-module $\mathcal{O}_P(1)|_W$. Il est alors clair que le
morphisme $0|_W : W \to L|_W$ est la section nulle de ce
fibré vectoriel. En particulier, $0(W)$ est un diviseur de Cartier effectif
sur $L|_W$. De plus, l'ouvert $U|_W = (\pi|_U)^{-1}(W)$
est le complémentaire de la section nulle.

\medskip\noindent
Si $A$ est engendrée par $f_1, \ldots, f_r \in A_1$ sur $A_0$, alors
$(f_1, \ldots, f_r)^m = A_{\geq m}$ pour tout $m \geq 0$ et, par conséquent, notre
$B$ ci-dessus est l'algèbre de Rees de $A_+ = (f_1, \ldots, f_r)$. Ainsi, dans ce
cas, $L \to X$ est le morphisme d'éclatement de centre $Z$ et $W = P$, où $W$ désigne l'ouvert défini au
paragraphe précédent.

\medskip\noindent
Si $P$ est quasi-compact, alors, pour $d$ suffisamment divisible, le
sous-schéma fermé $D \subset L$ défini par $\sigma(A_d)\mathcal{O}_L$
est un diviseur de Cartier effectif, $0 : P \to L$ se factorise par $D$,
et $0(P) = D$ ensemblistement. Cela résulte de
Constructions, lemme \ref{constructions-lemma-proj-quasi-compact},
et des assertions (1), (2), (3) et (4) démontrées ci-dessus. (On peut prendre pour $d$ tout multiple du
ppcm des degrés des éléments trouvés dans le lemme.)

\medskip\noindent
Nous continuons de supposer $P$ quasi-compact.
Soit $\mathcal{F}$ un $\mathcal{O}_P$-module quasi-cohérent.
Posons $\mathcal{F}_U = \pi^*\mathcal{F}|_U$. On a alors
\begin{equation}
\label{equation-cohomology-torsor}
R\Gamma(U, \mathcal{F}_U) =
\bigoplus\nolimits_{m \in \mathbf{Z}}
R\Gamma(P, \mathcal{F} \otimes_{\mathcal{O}_P} \mathcal{O}_P(m))
\end{equation}
De plus, cette décomposition en somme directe est fonctorielle en $\mathcal{F}$
et la structure de $A$-module induite sur le membre de droite est la même que la
structure de $A$-module du membre de gauche provenant de $U \subset X$.
Pour démontrer la formule, puisque $\pi|_U$ est affine et que
$(\pi|_U)_*\mathcal{O}_U = \bigoplus_{m \in \mathbf{Z}} \mathcal{O}_P(m)$,
on obtient
\begin{align*}
R(\pi|_U)_*\mathcal{F}_U
& =
(\pi|_U)_*\mathcal{F}_U \\
& =
(\pi|_U)_*(\pi|_U)^*\mathcal{F} \\
& =
\mathcal{F} \otimes_{\mathcal{O}_P} 
\bigoplus\nolimits_{m \in \mathbf{Z}} \mathcal{O}_P(m) \\
& =
\bigoplus\nolimits_{m \in \mathbf{Z}}
\mathcal{F} \otimes_{\mathcal{O}_P} \mathcal{O}_P(m)
\end{align*}
Par Leray, on trouve que
$R\Gamma(U, \mathcal{F}_U) =
R\Gamma(P, R(\pi|_U)_*\mathcal{F}_U)$, voir
Cohomologie, lemme \ref{cohomology-lemma-apply-Leray}.
La démonstration est achevée, car la prise de cohomologie commute aux
sommes directes dans ce cas, voir Catégories dérivées des schémas, lemme
\ref{perfect-lemma-quasi-coherence-pushforward-direct-sums}.
C'est ici que l'on utilise que $P$ est quasi-compact ; $P$ est séparé
d'après Constructions, lemme \ref{constructions-lemma-proj-separated}.

\begin{lemma}
\label{lemma-apply-proj-spec}
Soit $R$ un anneau. Soit $P$ un schéma propre sur $R$ et soit
$\mathcal{L}$ un $\mathcal{O}_P$-module inversible ample.
Posons $A = \bigoplus_{m \geq 0} \Gamma(P, \mathcal{L}^{\otimes m})$.
Alors $P = \text{Proj}(A)$ et le diagramme (\ref{equation-proj-and-spec})
devient le diagramme
$$
\xymatrix{
\underline{\Spec}_P \left(
\bigoplus\nolimits_{m \in \mathbf{Z}} \mathcal{L}^{\otimes m}
\right)
\ar[r] \ar@{=}[d] &
L =
\underline{\Spec}_P \left(
\bigoplus\nolimits_{m \geq 0} \mathcal{L}^{\otimes m}
\right) \ar[d]^\sigma \ar[r]_-\pi & P \ar[d] \\
U \ar[r] & X \ar[r] & Z
}
$$
ayant les propriétés expliquées ci-dessus.
\end{lemma}

\begin{proof}
On a $P = \text{Proj}(A)$ d'après
Morphismes, lemme \ref{morphisms-lemma-proper-ample-is-proj}.
De plus, d'après Propriétés, lemme \ref{properties-lemma-ample-gcd-is-one},
via cette identification, on a $\mathcal{O}_P(m) = \mathcal{L}^{\otimes m}$
pour tout $m \in \mathbf{Z}$.
\end{proof}









\section{Points fermés dans les fibres}
\label{section-closed-points-fibres}

\noindent
Une partie du contenu de cette section est tirée de la prépublication
\cite{Osserman-Payne}.

\begin{lemma}
\label{lemma-locally-principal-vertical}
Soit $f : X \to S$ un morphisme de schémas.
Soit $Z \subset X$ un sous-schéma fermé.
Soit $s \in S$.
Supposons que
\begin{enumerate}
\item $S$ est irréductible, de point générique $\eta$,
\item $X$ est irréductible,
\item $f$ est dominant,
\item $f$ est localement de type fini,
\item $\dim(X_s) \leq \dim(X_\eta)$,
\item $Z$ est localement principal dans $X$, et
\item $Z_\eta = \emptyset$.
\end{enumerate}
Alors la fibre $Z_s$ est (ensemblistement) une réunion de
composantes irréductibles de $X_s$.
\end{lemma}

\begin{proof}
Notons $X_{red}$ la réduction de $X$. Alors $Z \cap X_{red}$ est
un sous-schéma fermé localement principal de $X_{red}$, voir
Diviseurs, lemme \ref{divisors-lemma-pullback-locally-principal}.
On peut donc supposer que $X$ est réduit. Autrement dit, $X$ est intègre, voir
Propriétés, lemme \ref{properties-lemma-characterize-integral}.
Dans ce cas, le morphisme $X \to S$ se factorise par $S_{red}$, voir
Schémas, lemme \ref{schemes-lemma-map-into-reduction}.
Ainsi, on peut remplacer $S$ par $S_{red}$ et supposer que $S$ est également intègre.

\medskip\noindent
Dire que $f$ est dominant signifie que le point générique de $X$
s'envoie sur $\eta$, voir
Morphismes,
lemme \ref{morphisms-lemma-dominant-finite-number-irreducible-components}.
De plus, le schéma $X_\eta$ est un schéma intègre qui est localement de
type fini sur le corps $\kappa(\eta)$. Ainsi,
$d = \dim(X_\eta) \geq 0$ est égal à $\dim_\xi(X_\eta)$ pour
tout point $\xi$ de $X_\eta$, voir
Algèbre, lemmes \ref{algebra-lemma-dimension-spell-it-out} et
\ref{algebra-lemma-dimension-at-a-point-finite-type-over-field}.
Au vu de
Morphismes, lemme \ref{morphisms-lemma-openness-bounded-dimension-fibres}
et de la condition (5), on conclut que $\dim_x(X_s) = d$
pour tout $x \in X_s$.

\medskip\noindent
Dans le cas noethérien, l'assertion se démontre comme suit.
Si le lemme est faux, il existe $x \in Z_s$ qui est un point générique
d'une composante irréductible de $Z_s$, mais pas un point générique
d'une composante irréductible de $X_s$. On voit alors que
$\dim_x(Z_s) \leq d - 1$, car $\dim_x(X_s) = d$ et, dans un voisinage
de $x$ dans $X_s$, le sous-schéma fermé $Z_s$ ne contient aucune des
composantes irréductibles de $X_s$. Ainsi, après avoir remplacé $X$ par un
voisinage ouvert de $x$, on peut supposer que
$\dim_z(Z_{f(z)}) \leq d - 1$ pour tout $z \in Z$, voir
Morphismes, lemme \ref{morphisms-lemma-openness-bounded-dimension-fibres}.
Soit $\xi' \in Z$ un point générique d'une composante irréductible de $Z$
et posons $s' = f(\xi)$. Comme $Z \not = X$ et que ce sous-schéma est localement principal, on voit que
$\dim(\mathcal{O}_{X, \xi}) = 1$, voir
Algèbre, lemme \ref{algebra-lemma-minimal-over-1}
(c'est ici que l'on utilise que $X$ est noethérien).
Soit $\xi \in X$ le point générique de $X$ et
soit $\xi_1$ un point générique d'une composante irréductible quelconque
de $X_{s'}$ qui contient $\xi'$. On voit alors qu'on a
les spécialisations
$$
\xi \leadsto \xi_1 \leadsto \xi'.
$$
Comme $\dim(\mathcal{O}_{X, \xi}) = 1$, l'une des deux spécialisations
doit être une égalité.
Par hypothèse, $s' \not = \eta$ ; la première spécialisation
n'est donc pas une égalité.
Ainsi, $\xi' = \xi_1$ est un point générique d'une composante irréductible de
$X_{s'}$. En appliquant
Morphismes, lemme \ref{morphisms-lemma-openness-bounded-dimension-fibres}
une fois encore, ceci entraîne
$\dim_{\xi'}(Z_{s'}) = \dim_{\xi'}(X_{s'}) \geq \dim(X_\eta) = d$
ce qui donne la contradiction recherchée.

\medskip\noindent
Dans le cas général, on se ramène au cas noethérien comme suit.
Si le lemme est faux, alors il existe un point
$x \in X$ au-dessus de $s$ tel que $x$ soit un point générique d'une
composante irréductible de $Z_s$, mais
ne soit le point générique d'aucune des composantes irréductibles de $X_s$.
Soit $U \subset S$ un voisinage affine de $s$ et soit
$V \subset X$ un voisinage affine de $x$ tel que $f(V) \subset U$.
Écrivons $U = \Spec(A)$ et $V = \Spec(B)$, de sorte que $f|_V$
est donné par un homomorphisme d'anneaux $A \to B$. Soient $\mathfrak q \subset B$,
resp.\ $\mathfrak p \subset A$, les idéaux premiers correspondant à $x$, resp.\ $s$.
Quitte à rétrécir $V$, on peut supposer que $Z \cap V$ est défini par
un élément $g \in B$. Notons $K$ le corps des fractions de $A$.
Ce que nous savons à ce stade est ce qui suit :
\begin{enumerate}
\item $A \subset B$ est une extension de type fini d'anneaux intègres,
\item l'élément $g \otimes 1$ est inversible dans $B \otimes_A K$,
\item $d = \dim(B \otimes_A K) = \dim(B \otimes_A \kappa(\mathfrak p))$,
\item $g \otimes 1$ n'est pas une unité de $B \otimes_A \kappa(\mathfrak p)$, et
\item $g \otimes 1$ n'appartient à aucun des idéaux premiers minimaux de
$B \otimes_A \kappa(\mathfrak p)$.
\end{enumerate}
Nous cherchons une contradiction.

\medskip\noindent
Choisissons des éléments $x_1, \ldots, x_n \in B$ qui engendrent $B$ sur $A$.
Pour une $\mathbf{Z}$-algèbre de type fini $A_0 \subset A$,
soit $B_0 \subset B$ la sous-$A_0$-algèbre engendrée par
$x_1, \ldots, x_n$, notons $K_0$ le corps des fractions de $A_0$, et posons
$\mathfrak p_0 = A_0 \cap \mathfrak p$.
Nous affirmons que, lorsque $A_0$ est assez grande, (1) -- (5) sont encore satisfaites pour
le système $(A_0 \subset B_0, g, \mathfrak p_0)$.

\medskip\noindent
Démontrons successivement chacune des conditions. La condition (1) est satisfaite par construction.
Pour la partie (2), écrivons $(g \otimes 1) h = 1$ pour un certain
$h \otimes 1/a \in B \otimes_A K$. Écrivons
$g = \sum a_I x^I$, $h = \sum a'_I x^I$ (notation multi-indicielle)
pour certains coefficients $a_I, a'_I \in A$. Dès que $A_0$ contient
$a$ et les $a_I, a'_I$, alors (2) est vraie parce que
$B_0 \otimes_{A_0} K_0 \subset B \otimes_A K$ (comme localisations de
l'homomorphisme injectif $B_0 \to B$).
Pour obtenir (3), considérons la suite exacte
$$
0 \to I \to A[X_1, \ldots, X_n] \to B \to 0
$$
qui définit $I$, où le second homomorphisme envoie $X_i$ sur $x_i$. Comme $\otimes$
est exact à droite, on voit que $I \otimes_A K$ et, respectivement,
$I \otimes_A \kappa(\mathfrak p)$ sont les noyaux des surjections
$K[X_1, \ldots, X_n] \to B \otimes_A K$ et, respectivement,
$\kappa(\mathfrak p)[X_1, \ldots, X_n] \to B \otimes_A \kappa(\mathfrak p)$.
Comme un anneau de polynômes sur un corps est noethérien,
il existe un nombre fini d'éléments $h_j \in I$, $j = 1, \ldots, m$,
qui engendrent $I \otimes_A K$ et $I \otimes_A \kappa(\mathfrak p)$.
Écrivons $h_j = \sum a_{j, I}X^I$. Dès que
$A_0$ contient tous les $a_{j, I}$, on se trouve dans la situation où
$$
B_0 \otimes_{A_0} K_0 \otimes_{K_0} K = B \otimes_A K
\quad\text{et}\quad
B_0 \otimes_{A_0} \kappa(\mathfrak p_0)
\otimes_{\kappa(\mathfrak p_0)} \kappa(\mathfrak p)
=
B \otimes_A \kappa(\mathfrak p).
$$
En appliquant au choix
Morphismes, lemme \ref{morphisms-lemma-dimension-fibre-after-base-change}
ou
Algèbre, lemme \ref{algebra-lemma-dimension-preserved-field-extension},
on voit que les égalités de dimension de (3) sont satisfaites.
La partie (4) est immédiate. Comme
$B_0 \otimes_{A_0} \kappa(\mathfrak p_0) \subset
B \otimes_A \kappa(\mathfrak p)$, tout idéal premier minimal de
$B_0 \otimes_{A_0} \kappa(\mathfrak p_0)$ est la contraction d'un idéal
premier minimal de $B \otimes_A \kappa(\mathfrak p)$ d'après
Algèbre, lemme \ref{algebra-lemma-image-dense-generic-points}.
Cela montre que (5) est satisfaite.
On réduit ainsi le problème au cas noethérien, qui a été
traité ci-dessus.
\end{proof}

\noindent
Voici une application algébrique du lemme précédent.
La quatrième hypothèse du lemme est satisfaite si $A \to B$ est plat, voir
le lemme \ref{lemma-equality-dimensions}.

\begin{lemma}
\label{lemma-horizontal}
Soit $A \to B$ un homomorphisme local d'anneaux locaux, et soit
$g \in \mathfrak m_B$. Supposons que
\begin{enumerate}
\item $A$ et $B$ sont des anneaux intègres et $A \subset B$,
\item $B$ est essentiellement de type fini sur $A$,
\item $g$ n'appartient à aucun idéal premier minimal contenant $\mathfrak m_AB$, et
\item $\dim(B/\mathfrak m_AB) +
\text{trdeg}_{\kappa(\mathfrak m_A)}(\kappa(\mathfrak m_B)) =
\text{trdeg}_A(B)$.
\end{enumerate}
Alors $A \subset B/gB$, c'est-à-dire que le point générique de $\Spec(A)$
appartient à l'image du morphisme $\Spec(B/gB) \to \Spec(A)$.
\end{lemma}

\begin{proof}
Notons que les deux assertions sont équivalentes d'après
Algèbre, lemme \ref{algebra-lemma-image-dense-generic-points}.
Pour commencer la démonstration, soit $C$ une $A$-algèbre de type fini
et soit $\mathfrak q$ un idéal premier de $C$ tel que $B = C_{\mathfrak q}$.
On peut bien sûr supposer que $C$ est un anneau intègre et que $g \in C$.
Après avoir remplacé $C$ par un localisé, on voit que
$\dim(C/\mathfrak m_AC) = \dim(B/\mathfrak m_AB) +
\text{trdeg}_{\kappa(\mathfrak m_A)}(\kappa(\mathfrak m_B))$, voir
Morphismes, lemme \ref{morphisms-lemma-dimension-fibre-at-a-point}.
En posant $K$ égal au corps des fractions de $A$,
on voit par la même référence que
$\dim(C \otimes_A K) = \text{trdeg}_A(B)$. Ainsi, l'hypothèse
(4) signifie que les fibres générique et fermée du morphisme
$\Spec(C) \to \Spec(A)$ ont même dimension.

\medskip\noindent
Supposons le lemme faux. Alors $(B/gB) \otimes_A K = 0$.
Cela signifie que $g \otimes 1$ est inversible dans $B \otimes_A K
= C_{\mathfrak q} \otimes_A K$. Comme $C_{\mathfrak q}$ est une limite
de localisés principaux, on conclut que $g \otimes 1$
est inversible dans $C_h \otimes_A K$ pour un certain
$h \in C$, $h \not \in \mathfrak q$. Ainsi, après avoir remplacé $C$
par $C_h$, on peut supposer que $(C/gC) \otimes_A K = 0$.
On remplace encore une fois $C$ afin que les idéaux premiers minimaux
de $C/\mathfrak m_AC$ correspondent bijectivement aux idéaux premiers minimaux
de $B/\mathfrak m_AB$. À ce stade, on applique
le lemme \ref{lemma-locally-principal-vertical}
à $X = \Spec(C) \to \Spec(A) = S$ et au sous-schéma localement fermé
$Z = \Spec(C/gC)$. Comme $Z_K = \emptyset$, on voit que
$Z \otimes \kappa(\mathfrak m_A)$ doit contenir une composante irréductible
de
$X \otimes \kappa(\mathfrak m_A) = \Spec(C/\mathfrak m_AC)$.
Mais ceci contredit l'hypothèse que $g$ n'appartient
à aucun idéal premier minimal contenant $\mathfrak m_AB$. Le lemme en résulte.
\end{proof}

\begin{lemma}
\label{lemma-equality-dimensions}
Soit $A \to B$ un homomorphisme local d'anneaux locaux. Supposons que
\begin{enumerate}
\item $A$ et $B$ sont des anneaux intègres et $A \subset B$,
\item $B$ est essentiellement de type fini sur $A$, et
\item $B$ est plat sur $A$.
\end{enumerate}
Alors on a
$$
\dim(B/\mathfrak m_AB) +
\text{trdeg}_{\kappa(\mathfrak m_A)}(\kappa(\mathfrak m_B)) =
\text{trdeg}_A(B).
$$
\end{lemma}

\begin{proof}
Soit $C$ une $A$-algèbre de type fini et soit $\mathfrak q$ un idéal premier de $C$
tel que $B = C_{\mathfrak q}$. On peut supposer que $C$ est un anneau intègre.
On a
$\dim_{\mathfrak q}(C/\mathfrak m_AC) = \dim(B/\mathfrak m_AB) +
\text{trdeg}_{\kappa(\mathfrak m_A)}(\kappa(\mathfrak m_B))$, voir
Morphismes, lemme \ref{morphisms-lemma-dimension-fibre-at-a-point}.
En posant $K$ égal au corps des fractions de $A$,
on voit par la même référence que
$\dim(C \otimes_A K) = \text{trdeg}_A(B)$.
Ainsi, il s'agit en fait de démontrer que
$\dim_{\mathfrak q}(C/\mathfrak m_AC) = \dim(C \otimes_A K)$.
Choisissons un anneau de valuation $A'$ de $K$ dominant $A$, voir
Algèbre, lemme \ref{algebra-lemma-dominate}.
Posons $C' = C \otimes_A A'$.
Choisissons un idéal premier $\mathfrak q'$ de $C'$ au-dessus de $\mathfrak q$ ; un tel
idéal premier existe car
$$
C'/\mathfrak m_{A'}C' =
C/\mathfrak m_AC \otimes_{\kappa(\mathfrak m_A)} \kappa(\mathfrak m_{A'})
$$
ce qui démontre que $C/\mathfrak m_AC \to C'/\mathfrak m_{A'}C'$ est fidèlement
plat. Ceci démontre également que
$\dim_{\mathfrak q}(C/\mathfrak m_AC) =
\dim_{\mathfrak q'}(C'/\mathfrak m_{A'}C')$, voir
Algèbre,
lemme \ref{algebra-lemma-dimension-at-a-point-preserved-field-extension}.
Notons que $B' = C'_{\mathfrak q'}$ est un localisé de $B \otimes_A A'$.
Ainsi, $B'$ est plat sur $A'$. La fibre générique $B' \otimes_{A'} K$
est un localisé de $B \otimes_A K$. Ainsi, $B'$ est un anneau intègre.
Si l'on démontre le lemme pour $A' \subset B'$, on obtient alors l'égalité
$\dim_{\mathfrak q'}(C'/\mathfrak m_{A'}C') = \dim(C' \otimes_{A'} K)$
ce qui implique l'égalité recherchée
$\dim_{\mathfrak q}(C/\mathfrak m_AC) = \dim(C \otimes_A K)$
d'après ce qui précède. Cela ramène le
lemme au cas où $A$ est un anneau de valuation.

\medskip\noindent
Soit $A \subset B$ comme dans le lemme, avec $A$ un anneau de valuation.
Comme précédemment, écrivons $B = C_{\mathfrak q}$ pour un anneau intègre $C$ de type
fini sur $A$. D'après
Algèbre,
lemme \ref{algebra-lemma-finite-type-domain-over-valuation-ring-dim-fibres},
on obtient $\dim(C/\mathfrak m_AC) = \dim(C \otimes_A K)$, ce qui achève la démonstration.
\end{proof}

\begin{lemma}
\label{lemma-closed-point-nearby-fibre}
Soit $f : X \to S$ un morphisme de schémas.
Soit $x \leadsto x'$ une spécialisation de points de $X$.
Posons $s = f(x)$ et $s' = f(x')$.
Supposons que
\begin{enumerate}
\item $x'$ est un point fermé de $X_{s'}$, et
\item $f$ est localement de type fini.
\end{enumerate}
Alors l'ensemble
$$
\{x_1 \in X
\text{ tel que }
f(x_1) = s
\text{ et }
x_1\text{ est fermé dans }X_s
\text{ et }
x \leadsto x_1 \leadsto x'
\}
$$
est dense dans l'adhérence de $x$ dans $X_s$.
\end{lemma}

\begin{proof}
Appliquons
Schémas, lemme \ref{schemes-lemma-points-specialize},
à la spécialisation $x \leadsto x'$.
Cela fournit un morphisme $\varphi : \Spec(B) \to X$,
où $B$ est un anneau de valuation tel que $\varphi$ envoie le
point générique sur $x$ et le point fermé sur $x'$. On peut aussi
supposer que $\kappa(x)$ est le corps des fractions de $B$.
Soit $A = B \cap \kappa(s)$. Notons que c'est un anneau de valuation (voir
Algèbre, lemme \ref{algebra-lemma-valuation-ring-cap-field})
qui domine l'image de $\mathcal{O}_{S, s'} \to \kappa(s)$.
Considérons le diagramme commutatif
$$
\xymatrix{
\Spec(B) \ar[rd] \ar[r] &
X_A \ar[d] \ar[r] & X \ar[d] \\
& \Spec(A) \ar[r] & S
}
$$
Le point générique (resp.\ fermé) de $B$ s'envoie sur un point $x_A$
(resp.\ $x'_A$) de $X_A$ situé au-dessus du point générique (resp.\ fermé)
de $\Spec(A)$. Notons que $x'_A$ est un point fermé
de la fibre spéciale de $X_A$ d'après
Morphismes,
lemme \ref{morphisms-lemma-base-change-closed-point-fibre-locally-finite-type}.
Notons que la fibre générique de $X_A \to \Spec(A)$ est isomorphe
à $X_s$. Ainsi, nous sommes ramenés au cas où $S$ est
le spectre d'un anneau de valuation, $s = \eta \in S$ est le point générique, et
$s' \in S$ est le point fermé.

\medskip\noindent
Nous allons démontrer le lemme par récurrence sur $\dim_x(X_\eta)$.
Si $\dim_x(X_\eta) = 0$, alors il n'existe aucun autre point de $X_\eta$
se spécialisant en $x$, et $x$ est fermé dans sa fibre, voir
Morphismes, lemme \ref{morphisms-lemma-quasi-finite-at-point-characterize},
et le résultat en découle. Supposons que $\dim_x(X_\eta) > 0$.

\medskip\noindent
Soit $X' \subset X$ le sous-schéma fermé irréductible
$\overline{\{x\}}$ de $X$, muni de sa structure de schéma induite réduite, voir
Schémas, définition \ref{schemes-definition-reduced-induced-scheme}.
Pour démontrer le lemme, on peut remplacer $X$ par $X'$, puisque ceci ne fait que diminuer
$\dim_x(X_\eta)$. On peut donc aussi supposer que $X$ est un schéma intègre
et que $x$ est son point générique. En outre, on peut remplacer $X$ par un
voisinage affine de $x'$. Ainsi, on a $X = \Spec(B)$, où
$A \subset B$ est une extension de type fini d'anneaux intègres. Notons que, dans
ce cas, $\dim_x(X_\eta) = \dim(X_\eta) = \dim(X_{s'})$, et qu'en fait
$X_{s'}$ est équidimensionnel, voir
Algèbre,
lemme \ref{algebra-lemma-finite-type-domain-over-valuation-ring-dim-fibres}.

\medskip\noindent
Soit $W \subset X_\eta$ un fermé propre (c'est le
sous-ensemble que l'on veut « éviter »). Comme $X_s$ est de type fini sur un corps,
on voit que $W$ possède un nombre fini de composantes irréductibles
$W = W_1 \cup \ldots \cup W_n$. Soient
$\mathfrak q_j \subset B$, $j = 1, \ldots, r$
les idéaux premiers correspondants. Soit $\mathfrak q \subset B$
l'idéal maximal correspondant au point $x'$.
Soient $\mathfrak p_1, \ldots, \mathfrak p_s \subset B$ les
idéaux premiers minimaux contenant $\mathfrak m_AB$. Il y en a un nombre fini,
puisqu'ils correspondent aux composantes irréductibles du
schéma noethérien $X_{s'}$. De plus, chacune de ces composantes irréductibles
est de dimension $> 0$ (voir ci-dessus) ; on voit donc que
$\mathfrak p_i \not = \mathfrak q$ pour tout $i$.
Choisissons maintenant un élément $g \in \mathfrak q$ tel que
$g \not \in \mathfrak q_j$ pour tout $j$ et $g \not \in \mathfrak p_i$
pour tout $i$, voir
Algèbre, lemme \ref{algebra-lemma-silly}.
Notons $Z \subset X$ le sous-schéma fermé localement principal défini par $g$.
Écrivons $Z_\eta = Z_{1, \eta} \cup \ldots \cup Z_{n, \eta}$, $n \geq 0$,
pour la décomposition de la fibre générique de $Z$ en composantes irréductibles
(en nombre fini puisque la fibre générique est noethérienne).
Notons $Z_i \subset X$ l'adhérence de $Z_{i, \eta}$.
Après avoir remplacé $X$ par un voisinage affine plus petit,
on peut supposer que $x' \in Z_i$ pour tout $i = 1, \ldots, n$.
Par construction, $Z \cap X_{s'}$ ne contient aucune composante irréductible
de $X_{s'}$. Ainsi, d'après
le lemme \ref{lemma-locally-principal-vertical},
on conclut que $Z_\eta \not = \emptyset$ ! Autrement dit,
$n \geq 1$. Soit $x_1 \in Z_1$ le point générique ; on voit
que $x_1 \leadsto x'$ et $f(x_1) = \eta$.
De plus, par construction, $Z_{1, \eta} \cap W_j \subset W_j$
est un fermé propre. Ainsi, toute composante irréductible de
$Z_{1, \eta} \cap W_j$ est de codimension $\geq 2$ dans $X_\eta$,
tandis que $\text{codim}(Z_{1, \eta}, X_\eta) = 1$ d'après
Algèbre, lemme \ref{algebra-lemma-minimal-over-1}.
Ainsi, $W \cap Z_{1, \eta}$ est un fermé propre.
À ce stade, on voit que l'hypothèse de récurrence s'applique à
$Z_1 \to S$ et à la spécialisation $x_1 \leadsto x'$.
Ceci fournit un point fermé $x_2$ de $Z_{1, \eta}$ qui n'appartient pas
à $W$ et qui se spécialise en $x'$. Ainsi, on obtient
$x \leadsto x_2 \leadsto x'$ ; le point $x_2$ est fermé dans $X_\eta$,
et $x_2 \not \in W$, comme souhaité.
\end{proof}

\begin{remark}
\label{remark-full-specialization-sequence}
La démonstration du
lemme \ref{lemma-closed-point-nearby-fibre}
montre en fait qu'il existe une suite de spécialisations
$$
x \leadsto x_1 \leadsto x_2 \leadsto \ldots \leadsto x_d \leadsto x'
$$
où tous les $x_i$ appartiennent à la fibre $X_s$, chaque spécialisation est
immédiate, et $x_d$ est un point fermé de $X_s$. L'entier
$d = \text{trdeg}_{\kappa(s)}(\kappa(x)) = \dim(\overline{\{x\}})$
où l'adhérence est prise dans $X_s$. De plus, les points
$x_i$ peuvent être choisis de manière à éviter tout fermé de $X_s$ qui
ne contient pas le point $x$.
\end{remark}

\noindent
Exemples, section \ref{examples-section-topology-finite-type},
montre que le lemme suivant est faux si $A$ n'est pas supposé
noethérien.

\begin{lemma}
\label{lemma-quasi-finite-quasi-section-meeting-nearby-open}
Soit $\varphi : A \to B$ un homomorphisme local d'anneaux locaux.
Soit $V \subset \Spec(B)$ un sous-schéma ouvert
qui contient au moins un idéal premier ne se trouvant pas au-dessus de $\mathfrak m_A$.
Supposons que $A$ soit noethérien, que $\varphi$ soit essentiellement de type fini, et que
l'extension $A/\mathfrak m_A \subset B/\mathfrak m_B$ soit finie.
Alors il existe un $\mathfrak q \in V$,
$\mathfrak m_A \not = \mathfrak q \cap A$, tel que
$A \to B/\mathfrak q$ soit le localisé d'un homomorphisme d'anneaux quasi-fini.
\end{lemma}

\begin{proof}
Comme $A$ est noethérien et que $A \to B$ est essentiellement de type fini,
on sait que $B$ est également noethérien. D'après
Propriétés, lemme \ref{properties-lemma-complement-closed-point-Jacobson},
l'espace topologique $\Spec(B) \setminus \{\mathfrak m_B\}$
est jacobsonien. On peut donc choisir un point fermé $\mathfrak q$
qui appartient à l'ouvert non vide
$$
V \setminus \{\mathfrak q \subset B \mid \mathfrak m_A = \mathfrak q \cap A\}.
$$
(Non vide par hypothèse, ouvert parce que $\{\mathfrak m_A\}$ est un fermé
de $\Spec(A)$.)
Alors $\Spec(B/\mathfrak q)$ possède deux points, à savoir $\mathfrak m_B$
et $\mathfrak q$, et $\mathfrak q$ ne se trouve pas au-dessus de $\mathfrak m_A$.
Écrivons $B/\mathfrak q = C_{\mathfrak m}$ pour une $A$-algèbre de type fini
$C$ et un idéal premier $\mathfrak m$. Alors $A \to C$ est quasi-fini en
$\mathfrak m$ d'après
Algèbre, lemme \ref{algebra-lemma-isolated-point-fibre} (2).
Par conséquent, d'après
Algèbre, lemme \ref{algebra-lemma-quasi-finite-open},
on voit qu'après avoir remplacé $C$ par un localisé principal, l'homomorphisme
$A \to C$ est quasi-fini.
\end{proof}

\begin{lemma}
\label{lemma-quasi-finite-quasi-section-meeting-nearby-open-X}
Soit $f : X \to S$ un morphisme de schémas.
Soit $x \in X$ d'image $s \in S$.
Soit $U \subset X$ un sous-schéma ouvert.
Supposons que $f$ soit localement de type fini, que $S$ soit localement noethérien, et que $x$ soit un point fermé
de $X_s$, et qu'il existe un point $x' \in U$ tel que
$x' \leadsto x$ et $f(x') \not = s$. Alors il existe un sous-schéma fermé
$Z \subset X$ tel que (a) $x \in Z$, (b) $f|_Z : Z \to S$ soit
quasi-fini en $x$, et (c) il existe un $z \in Z$, $z \in U$,
tel que $z \leadsto x$ et $f(z) \not = s$.
\end{lemma}

\begin{proof}
C'est une reformulation du
lemme \ref{lemma-quasi-finite-quasi-section-meeting-nearby-open}.
En effet, posons $A = \mathcal{O}_{S, s}$ et $B = \mathcal{O}_{X, x}$.
Notons $V \subset \Spec(B)$ l'image réciproque de $U$.
L'homomorphisme d'anneaux $f^\sharp : A \to B$ est essentiellement de type fini.
Par hypothèse, il existe au moins un point de $V$ qui ne
s'envoie pas sur le point fermé de $\Spec(A)$. Toutes les hypothèses du
lemme \ref{lemma-quasi-finite-quasi-section-meeting-nearby-open}
sont donc satisfaites, et l'on obtient un idéal premier $\mathfrak q \subset B$ qui ne
se trouve pas au-dessus de $\mathfrak m_A$ et tel que $A \to B/\mathfrak q$ soit
le localisé d'un homomorphisme d'anneaux quasi-fini. Soit $z \in X$
l'image du point $\mathfrak q$ par le morphisme canonique
$\Spec(B) \to X$. Posons $Z = \overline{\{z\}}$
muni de la structure de schéma induite réduite. Comme $z \leadsto x$,
on voit que $x \in Z$ et $\mathcal{O}_{Z, x} = B/\mathfrak q$.
Par construction, $Z \to S$ est quasi-fini en $x$.
\end{proof}

\begin{remark}
\label{remark-alternative-closed-point-nearby-fibre}
On peut utiliser le
lemme \ref{lemma-quasi-finite-quasi-section-meeting-nearby-open}
ou sa variante, le
lemme \ref{lemma-quasi-finite-quasi-section-meeting-nearby-open-X},
pour donner une autre démonstration du
lemme \ref{lemma-closed-point-nearby-fibre}
lorsque $S$ est localement noethérien.
En voici une esquisse.
On remplace d'abord $S$ par
le spectre de l'anneau local en $s'$. Puis on peut raisonner par récurrence
sur $\dim(S)$. Le cas $\dim(S) = 0$ est trivial, car alors $s' = s$.
Remplaçons $X$ par $\overline{\{x\}}$, muni de sa structure de schéma induite réduite.
Appliquons le
lemme \ref{lemma-quasi-finite-quasi-section-meeting-nearby-open-X}
à $X \to S$ et à $x' \mapsto s'$, ainsi qu'à tout
ouvert non vide $U \subset X$ contenant $x$. On obtient un sous-schéma fermé
$x' \in Z \subset X$ et un point $z \in Z$
tels que $Z \to S$ soit quasi-fini en $x'$ et que $f(z) \not = s'$.
Alors $z$ est un point fermé de $X_{f(z)}$, et $z \leadsto x'$.
Comme $f(z) \not = s'$, on a $\dim(\mathcal{O}_{S, f(z)}) < \dim(S)$.
Puisque $x$ est le point générique de $X$, on a $x \leadsto z$, donc
$s = f(x) \leadsto f(z)$.
Appliquons l'hypothèse de récurrence à $s \leadsto f(z)$ et à $z \mapsto f(z)$
pour conclure.
\end{remark}

\begin{lemma}
\label{lemma-change-hypotheses}
Supposons que $f : X \to S$ soit localement de type fini et que $S$ soit localement noethérien,
que $x \in X$ soit un point fermé de sa fibre $X_s$ et que $U \subset X$ soit un sous-schéma
ouvert tel que $U \cap X_s = \emptyset$ et $x \in \overline{U}$ ; alors
les conclusions du
lemme \ref{lemma-quasi-finite-quasi-section-meeting-nearby-open-X}
sont valables.
\end{lemma}

\begin{proof}
En effet, on peut se ramener au lemme cité comme suit : d'abord, on
remplace $X$ et $S$ par des voisinages affines de $x$ et $s$. Alors $X$ est
noethérien ; en particulier, $U$ est quasi-compact (voir
Morphismes, lemme \ref{morphisms-lemma-finite-type-noetherian}
et
Topologie, lemmes \ref{topology-lemma-Noetherian} et
\ref{topology-lemma-Noetherian-quasi-compact}).
Il existe donc une spécialisation $x' \leadsto x$ avec $x' \in U$ (voir
Morphismes, lemme \ref{morphisms-lemma-reach-points-scheme-theoretic-image}).
Notons que $f(x') \not = s$. Ainsi, toutes les hypothèses du lemme
sont satisfaites, ce qui conclut.
\end{proof}



\section{Factorisation de Stein}
\label{section-stein-factorization}

\noindent
La factorisation de Stein affirme qu'un morphisme propre $f : X \to S$
tel que $f_*\mathcal{O}_X = \mathcal{O}_S$ a des fibres connexes.

\begin{lemma}
\label{lemma-stein-universally-closed}
Soit $S$ un schéma. Soit $f : X \to S$ un morphisme universellement fermé et
quasi-séparé. Il existe une factorisation
$$
\xymatrix{
X \ar[rr]_{f'} \ar[rd]_f & & S' \ar[dl]^\pi \\
& S &
}
$$
ayant les propriétés suivantes :
\begin{enumerate}
\item le morphisme $f'$ est universellement fermé, quasi-compact, quasi-séparé,
et surjectif,
\item le morphisme $\pi : S' \to S$ est intégral,
\item on a $f'_*\mathcal{O}_X = \mathcal{O}_{S'}$,
\item on a $S' = \underline{\Spec}_S(f_*\mathcal{O}_X)$, et
\item $S'$ est la normalisation de $S$ dans $X$, voir
Morphismes, définition \ref{morphisms-definition-normalization-X-in-Y}.
\end{enumerate}
La formation de la factorisation $f = \pi \circ f'$ commute
au changement de base plat.
\end{lemma}

\begin{proof}
D'après Morphismes, lemme \ref{morphisms-lemma-universally-closed-quasi-compact},
le morphisme $f$ est quasi-compact. Par conséquent, la normalisation $S'$ de $S$ dans
$X$ est définie (Morphismes, définition
\ref{morphisms-definition-normalization-X-in-Y})
et l'on a la factorisation $X \to S' \to S$. D'après
Morphismes, lemme \ref{morphisms-lemma-normalization-in-universally-closed},
on a (2), (4) et (5). Le morphisme $f'$ est universellement fermé d'après
Morphismes, lemme \ref{morphisms-lemma-image-proper-scheme-closed}.
Il est quasi-compact d'après
Schémas, lemme \ref{schemes-lemma-quasi-compact-permanence},
et quasi-séparé d'après
Schémas, lemme \ref{schemes-lemma-compose-after-separated}.

\medskip\noindent
Pour démontrer les assertions restantes, on peut supposer que le schéma de base $S$ est affine,
disons $S = \Spec(R)$. Alors $S' = \Spec(A)$ avec
$A = \Gamma(X, \mathcal{O}_X)$ une $R$-algèbre entière.
Il est donc clair que $f'_*\mathcal{O}_X$
est $\mathcal{O}_{S'}$ (car $f'_*\mathcal{O}_X$ est quasi-cohérent,
d'après
Schémas, lemme
\ref{schemes-lemma-push-forward-quasi-coherent},
et donc égal à $\widetilde{A}$). Cela démontre (3).

\medskip\noindent
Montrons que $f'$ est surjectif. Comme $f'$ est universellement fermé (voir ci-dessus),
l'image de $f'$ est un sous-ensemble fermé
$V(I) \subset S' = \Spec(A)$. Prenons $h \in I$. Alors
$h|_X = f^\sharp(h)$ est une section globale du faisceau structural de
$X$ qui s'annule en tout point. Comme $X$ est quasi-compact, cela signifie
que $h|_X$ est une section nilpotente, c'est-à-dire que $h^n|X = 0$ pour un certain $n > 0$.
Mais $A = \Gamma(X, \mathcal{O}_X)$, donc $h^n = 0$.
Tout élément de $I$ étant nilpotent, on conclut que $V(I) = S'$, comme voulu.

\medskip\noindent
D'après Cohomologie des schémas, lemme \ref{coherent-lemma-flat-base-change-cohomology},
on voit que la formation de $f_*\mathcal{O}_X$ commute au changement de base plat.
La formation du spectre relatif commute à tout changement de base d'après
Constructions, lemme \ref{constructions-lemma-spec-properties}.
Ainsi, la formation de la factorisation commute au changement de base plat.
\end{proof}

\begin{lemma}
\label{lemma-stein-universally-closed-residue-fields}
Dans le lemme \ref{lemma-stein-universally-closed}, supposons en outre que
$f$ soit localement de type fini. Alors, pour $s \in S$, la fibre
$\pi^{-1}(\{s\}) = \{s_1, \ldots, s_n\}$ est finie et les extensions de corps
$\kappa(s_i)/\kappa(s)$ sont finies.
\end{lemma}

\begin{proof}
Rappelons qu'il n'y a pas de spécialisations entre les points de $\pi^{-1}(\{s\})$ ;
voir Algèbre, lemme \ref{algebra-lemma-integral-no-inclusion}.
Comme $f'$ est surjectif, on obtient que $|X_s| \to \pi^{-1}(\{s\})$ est surjectif.
Observons que $X_s$ est un schéma quasi-séparé de type fini
sur un corps (la quasi-compacité a été démontrée dans la preuve du
lemme cité). Ainsi, $X_s$ est noethérien
(Morphismes, lemme \ref{morphisms-lemma-finite-type-noetherian}).
Un argument topologique (omis) montre alors que $\pi^{-1}(\{s\})$ est fini.
Pour chaque $i$, on peut choisir un point de type fini $x_i \in X_s$ s'envoyant sur $s_i$
(Morphismes, lemme \ref{morphisms-lemma-enough-finite-type-points}).
On conclut que $\kappa(s_i)/\kappa(s)$ est finie :
$x_i$ peut être représenté par un morphisme $\Spec(k_i) \to X_s$
de type fini (d'après notre définition des points de type fini)
et, par conséquent, $\Spec(k_i) \to s = \Spec(\kappa(s))$ est de type fini
(comme composé de morphismes de type fini) ;
ainsi, $k_i/\kappa(s)$ est finie (Morphismes, lemme
\ref{morphisms-lemma-point-finite-type}).
\end{proof}

\begin{lemma}
\label{lemma-characterize-geometrically-connected-fibres}
Soit $f : X \to S$ un morphisme de schémas.
Soit $s \in S$. Alors $X_s$ est géométriquement connexe si et
seulement si, pour tout voisinage \'etale $(U, u) \to (S, s)$,
le changement de base $X_U \to U$ a une fibre connexe $X_u$.
\end{lemma}

\begin{proof}
Si $X_s$ est géométriquement connexe, tout changement de base de celui-ci est connexe.
Réciproquement, supposons que $X_s$ ne soit pas géométriquement connexe.
Alors, d'après
Variétés, lemme
\ref{varieties-lemma-characterize-geometrically-disconnected},
on voit que $X_s \times_{\Spec(\kappa(s))} \Spec(k)$ est
non connexe pour une certaine
extension de corps séparable finie $k/\kappa(s)$. D'après
le lemme \ref{lemma-realize-prescribed-residue-field-extension-etale},
il existe un voisinage \'etale affine $(U, u) \to (S, s)$ tel que
$\kappa(u)/\kappa(s)$ soit identifié à $k/\kappa(s)$.
Dans ce cas, $X_u$ est non connexe.
\end{proof}

\begin{theorem}[Factorisation de Stein ; cas noethérien]
\label{theorem-stein-factorization-Noetherian}
Soit $S$ un schéma localement noethérien.
Soit $f : X \to S$ un morphisme propre.
Il existe une factorisation
$$
\xymatrix{
X \ar[rr]_{f'} \ar[rd]_f & & S' \ar[dl]^\pi \\
& S &
}
$$
ayant les propriétés suivantes :
\begin{enumerate}
\item le morphisme $f'$ est propre à fibres géométriquement connexes,
\item le morphisme $\pi : S' \to S$ est fini,
\item on a $f'_*\mathcal{O}_X = \mathcal{O}_{S'}$,
\item on a $S' = \underline{\Spec}_S(f_*\mathcal{O}_X)$, et
\item $S'$ est la normalisation de $S$ dans $X$, voir
Morphismes, définition \ref{morphisms-definition-normalization-X-in-Y}.
\end{enumerate}
\end{theorem}

\begin{proof}
Soit $f = \pi \circ f'$ la factorisation du
lemme \ref{lemma-stein-universally-closed}. Notons qu'outre les
conclusions du lemme \ref{lemma-stein-universally-closed}, on
a aussi que $f'$ est séparé
(Schémas, lemme \ref{schemes-lemma-compose-after-separated})
et de type fini
(Morphismes, lemme \ref{morphisms-lemma-permanence-finite-type}).
Ainsi, $f'$ est propre. D'après
Cohomologie des schémas, proposition
\ref{coherent-proposition-proper-pushforward-coherent},
on voit que $f_*\mathcal{O}_X$ est un $\mathcal{O}_S$-module cohérent.
On voit donc que $\pi$ est fini, c'est-à-dire que (2) est satisfaite.

\medskip\noindent
Cela démontre toutes les assertions sauf la plus intéressante, à savoir que
toutes les fibres de $f'$ sont géométriquement connexes.
Il résulte clairement de la discussion précédente que l'on peut remplacer $S$ par $S'$ ;
on peut donc supposer que $S$ est noethérien, affine,
que $f : X \to S$ est propre et que $f_*\mathcal{O}_X = \mathcal{O}_S$.
Soit $s \in S$ un point de $S$. Il faut montrer que $X_s$ est
géométriquement connexe. D'après le lemme
\ref{lemma-characterize-geometrically-connected-fibres},
on voit qu'il suffit de montrer que $X_u$ est connexe
pour tout voisinage \'etale $(U, u) \to (S, s)$.
On peut supposer que $U$ est affine. Alors $U$ est noethérien
(Morphismes, lemme \ref{morphisms-lemma-finite-type-noetherian}),
le changement de base $f_U : X_U \to U$ est propre
(Morphismes, lemme \ref{morphisms-lemma-base-change-proper}),
et l'on a aussi $(f_U)_*\mathcal{O}_{X_U} = \mathcal{O}_U$
(Cohomologie des schémas, lemme \ref{coherent-lemma-flat-base-change-cohomology}).
Ainsi, après avoir remplacé
$(f : X \to S, s)$ par le changement de base $(f_U : X_U \to U, u)$,
il suffit de démontrer que la fibre $X_s$ est connexe lorsque
$f_*\mathcal{O}_X = \mathcal{O}_S$. On peut le déduire de
Catégories dérivées de schémas, lemme
\ref{perfect-lemma-proper-idempotent-on-fibre}
(en considérant les idempotents du faisceau structural de $X_s$),
mais nous donnerons aussi ci-dessous un argument direct.

\medskip\noindent
Plus précisément, appliquons le théorème des fonctions formelles,
sous la forme de Cohomologie des schémas, lemme
\ref{coherent-lemma-formal-functions-stalk}.
Il nous dit que
$$
\mathcal{O}^\wedge_{S, s} = (f_*\mathcal{O}_X)_s^\wedge =
\lim_n H^0(X_n, \mathcal{O}_{X_n})
$$
où $X_n$ est le $n$-ième voisinage infinitésimal de $X_s$.
Puisque l'espace topologique sous-jacent à $X_n$ est égal à celui
de $X_s$, on voit que, si $X_s = T_1 \amalg T_2$ est une réunion disjointe
de sous-schémas ouverts et fermés non vides, alors, de même,
$X_n = T_{1, n} \amalg T_{2, n}$ pour tout $n$. Cela signifie à son tour que
$H^0(X_n, \mathcal{O}_{X_n})$ contient un idempotent non trivial $e_{1, n}$,
à savoir la fonction identiquement égale à $1$ sur $T_{1, n}$ et
identiquement égale à $0$ sur $T_{2, n}$. Il est clair que $e_{1, n + 1}$
se restreint à $e_{1, n}$ sur $X_n$. Ainsi, $e_1 = \lim e_{1, n}$
est un idempotent non trivial de la limite. Cela contredit le fait
que $\mathcal{O}^\wedge_{S, s}$ soit un anneau local. Ainsi,
l'hypothèse était fausse, c'est-à-dire que $X_s$ est connexe, ce qui conclut.
\end{proof}

\begin{theorem}[Factorisation de Stein ; cas général]
\label{theorem-stein-factorization-general}
Soit $S$ un schéma.
Soit $f : X \to S$ un morphisme propre.
Il existe une factorisation
$$
\xymatrix{
X \ar[rr]_{f'} \ar[rd]_f & & S' \ar[dl]^\pi \\
& S &
}
$$
ayant les propriétés suivantes :
\begin{enumerate}
\item le morphisme $f'$ est propre à fibres géométriquement connexes,
\item le morphisme $\pi : S' \to S$ est entier,
\item on a $f'_*\mathcal{O}_X = \mathcal{O}_{S'}$,
\item on a $S' = \underline{\Spec}_S(f_*\mathcal{O}_X)$, et
\item $S'$ est la normalisation de $S$ dans $X$, voir
Morphismes, définition \ref{morphisms-definition-normalization-X-in-Y}.
\end{enumerate}
\end{theorem}

\begin{proof}
On peut appliquer le lemme \ref{lemma-stein-universally-closed} pour obtenir le
morphisme $f' : X \to S'$.
Notons qu'outre les
conclusions du lemme \ref{lemma-stein-universally-closed}, on
a aussi que $f'$ est séparé
(Schémas, lemme \ref{schemes-lemma-compose-after-separated})
et de type fini
(Morphismes, lemme \ref{morphisms-lemma-permanence-finite-type}).
Ainsi, $f'$ est propre. À ce stade, on a démontré toutes les
assertions, sauf celle selon laquelle
$f'$ a des fibres géométriquement connexes.

\medskip\noindent
On peut supposer que $S = \Spec(R)$ est affine.
Posons $R' = \Gamma(X, \mathcal{O}_X)$. Alors $S' = \Spec(R')$.
On peut donc remplacer $S$ par $S'$ et supposer que
$S = \Spec(R)$ est affine et que $R = \Gamma(X, \mathcal{O}_X)$.
Ensuite, soit $s \in S$ un point. Soit $U \to S$ un morphisme \'etale
de schémas affines et soit $u \in U$ un point s'envoyant sur $s$.
Soit $X_U \to U$ le changement de base de $X$. D'après
le lemme \ref{lemma-characterize-geometrically-connected-fibres},
il suffit de montrer que la fibre de $X_U \to U$ au-dessus de $u$ est
connexe. D'après
Cohomologie des schémas, lemme \ref{coherent-lemma-flat-base-change-cohomology},
on voit que
$\Gamma(X_U, \mathcal{O}_{X_U}) = \Gamma(U, \mathcal{O}_U)$.
Il faut donc montrer ceci : étant donnés
$S = \Spec(R)$ affine, $X \to S$ propre avec $\Gamma(X, \mathcal{O}_X) = R$
et $s \in S$ un point, la fibre $X_s$ est connexe.

\medskip\noindent
Pour cela, il suffit de montrer que les seuls idempotents
$e \in H^0(X_s, \mathcal{O}_{X_s})$ sont $0$ et $1$ (on sait déjà
que $X_s$ est non vide d'après le lemme \ref{lemma-stein-universally-closed}).
D'après Catégories dérivées de schémas, lemme
\ref{perfect-lemma-proper-idempotent-on-fibre},
après avoir remplacé $R$ par une localisation principale,
on peut supposer que $e$ est l'image d'un élément de $R$.
Puisque $R \to H^0(X_s, \mathcal{O}_{X_s})$ se factorise par
$\kappa(s)$, on conclut.
\end{proof}

\noindent
Voici une application.

\begin{lemma}
\label{lemma-geometrically-connected-fibres-towards-normal}
Soit $f : X \to S$ un morphisme de schémas. Supposons que
\begin{enumerate}
\item $f$ soit propre,
\item $S$ soit intègre de point générique $\xi$,
\item $S$ soit normal,
\item $X$ soit réduit,
\item tout point générique d'une composante irréductible de $X$ s'envoie sur $\xi$,
\item on ait $H^0(X_\xi, \mathcal{O}) = \kappa(\xi)$.
\end{enumerate}
Alors $f_*\mathcal{O}_X = \mathcal{O}_S$ et $f$
a des fibres géométriquement connexes.
\end{lemma}

\begin{proof}
Appliquons le théorème \ref{theorem-stein-factorization-general} pour obtenir une
factorisation $X \to S' \to S$. Il suffit de montrer que $S' = S$.
Cela résultera de Morphismes, lemme
\ref{morphisms-lemma-finite-birational-over-normal}.
En effet, $S'$ est réduit parce que $X$ est réduit
(Morphismes, lemme \ref{morphisms-lemma-normalization-in-reduced}).
Le morphisme $S' \to S$ est entier d'après le théorème cité ci-dessus.
Tout point générique de $S'$ est au-dessus de $\xi$ d'après
Morphismes, lemme \ref{morphisms-lemma-normalization-generic},
et l'hypothèse (5). D'autre part, puisque $S'$ est le spectre relatif
de $f_*\mathcal{O}_X$, on voit que la fibre schématique
$S'_\xi$ est le spectre de $H^0(X_\xi, \mathcal{O})$, qui est
égal à $\kappa(\xi)$ par hypothèse. Ainsi, $S'$ est un schéma
intègre dont le corps des fonctions est égal à celui de $S$.
Cela conclut la preuve.
\end{proof}

\noindent
Voici une autre application.

\begin{lemma}
\label{lemma-proper-flat-nr-geom-conn-comps-lower-semicontinuous}
Soit $X \to S$ un morphisme plat propre de présentation finie. Soit
$n_{X/S}$ la fonction sur $S$ qui compte le nombre de composantes
connexes géométriques des fibres de $f$ introduite dans le
lemme \ref{lemma-base-change-fibres-nr-geometrically-connected-components}.
Alors $n_{X/S}$ est semi-continue inférieurement.
\end{lemma}

\begin{proof}
Soit $s \in S$. Posons $n = n_{X/S}(s)$. Notons que $n < \infty$, car la fibre
géométrique de $X \to S$ en $s$ est un schéma propre sur un corps, donc noethérien,
et possède donc un nombre fini de composantes connexes. Il faut trouver un
voisinage ouvert $V$ de $s$ tel que $n_{X/S}|_V \geq n$.
Soit $X \to S' \to S$ la factorisation de Stein du
théorème \ref{theorem-stein-factorization-general}.
D'après le lemme \ref{lemma-stein-universally-closed-residue-fields},
il y a un nombre fini de points $s'_1, \ldots, s'_m \in S'$ au-dessus de $s$,
et les extensions $\kappa(s'_i)/\kappa(s)$ sont finies.
Le lemme \ref{lemma-etale-makes-integral-split}
nous dit alors qu'après avoir remplacé $S$ par un voisinage \'etale
de $s$, on peut supposer que $S' = V_1 \amalg \ldots \amalg V_m$ comme schéma,
avec $s'_i \in V_i$ et $\kappa(s'_i)/\kappa(s)$ purement inséparable.
Alors les schémas $X_{s_i'}$ sont géométriquement connexes sur
$\kappa(s)$, donc $m = n$. Les schémas
$X_i = (f')^{-1}(V_i)$, $i = 1, \ldots, n$,
sont plats et de présentation finie sur $S$. Par conséquent, l'image de $X_i \to S$
est ouverte (Morphismes, lemme \ref{morphisms-lemma-fppf-open}).
Ainsi, dans un voisinage de $s$, on voit que $n_{X/S}$ est
au moins égal à $n$.
\end{proof}

\begin{lemma}
\label{lemma-proper-flat-geom-red}
Soit $f : X \to S$ un morphisme de schémas. Supposons que
\begin{enumerate}
\item $f$ soit propre, plat et de présentation finie, et
\item les fibres géométriques de $f$ soient réduites.
\end{enumerate}
Alors la fonction $n_{X/S} : S \to \mathbf{Z}$,
qui compte le nombre de composantes connexes géométriques
des fibres de $f$, est localement constante.
\end{lemma}

\begin{proof}
D'après le lemme \ref{lemma-proper-flat-nr-geom-conn-comps-lower-semicontinuous},
la fonction $n_{X/S}$ est semi-continue inférieurement.
Pour $s \in S$, considérons la $\kappa(s)$-algèbre
$$
A_s = H^0(X_s, \mathcal{O}_{X_s})
$$
Définissons une fonction $\beta_0 : X \to \mathbf{Z}$ qui envoie
$s$ sur $\dim_{\kappa(s)} A_s$.
D'après Variétés, lemme
\ref{varieties-lemma-proper-geometrically-reduced-global-sections},
et le fait que $X_s$ est géométriquement réduit,
$A_s$ est un produit fini d'extensions séparables finies de $\kappa(s)$.
Ainsi, $A_s \otimes_{\kappa(s)} \kappa(\overline{s})$ est un produit
de $\beta_0(s)$ copies de $\kappa(\overline{s})$. Par conséquent,
$X_{\overline{s}}$ a $\beta_0(s)$ composantes connexes.
Autrement dit, on a $n_{X/S} = \beta_0$
comme fonctions sur $S$. Ainsi, $n_{X/S}$ est semi-continue supérieurement d'après
Catégories dérivées de schémas, lemme \ref{perfect-lemma-jump-loci-geometric}.
Cela conclut la preuve.
\end{proof}

\noindent
Une dernière application.

\begin{lemma}
\label{lemma-split-off-proper-part-henselian}
\begin{reference}
Une référence pour le cas d'une base noethérienne adique est
\cite[III, Proposition 5.5.1]{EGA}
\end{reference}
Soit $(A, I)$ un couple hensélien. Soit $X \to \Spec(A)$
séparé et de type fini. Posons $X_0 = X \times_{\Spec(A)} \Spec(A/I)$.
Soit $Y \subset X_0$ un sous-schéma ouvert et fermé tel que
$Y \to \Spec(A/I)$ soit propre. Alors il existe un sous-schéma ouvert et fermé
$W \subset X$, propre sur $A$, tel que
$W \times_{\Spec(A)} \Spec(A/I) = Y$.
\end{lemma}

\begin{proof}
Nous noterons $T \mapsto T_0$ le changement de base par $\Spec(A/I) \to \Spec(A)$.
D'après le lemme de Chow (sous la forme du
lemme de Limites \ref{limits-lemma-chow-finite-type}),
il existe un morphisme propre surjectif $\varphi : X' \to X$ tel
que $X'$ admette une immersion dans $\mathbf{P}^n_A$.
Posons $Y' = \varphi^{-1}(Y)$. C'est un sous-schéma ouvert et fermé
de $X'_0$. Supposons le lemme vrai pour $(X', Y')$. Soit $W' \subset X'$
le sous-schéma ouvert et fermé propre sur $A$ tel que $Y' = W'_0$.
D'après Morphismes, lemme \ref{morphisms-lemma-image-proper-scheme-closed},
$W = \varphi(W') \subset X$ et
$Q = \varphi(X' \setminus W') \subset X$ sont des sous-ensembles fermés et, d'après
Morphismes, lemme \ref{morphisms-lemma-image-proper-is-proper},
$W$ est propre sur $A$. L'image de $W \cap Q$ dans $\Spec(A)$ est fermée.
Puisque $(A, I)$ est hensélien, si $W \cap Q$ est non vide, alors
il existe dans $W \cap Q$ un point au-dessus de $\Spec(A/I)$.
C'est impossible puisque $W'_0 = Y' = \varphi^{-1}(Y)$.
On conclut que $W$ est un sous-schéma ouvert et fermé
de $X$, propre sur $A$, avec $W_0 = Y$.
On se ramène ainsi au cas décrit au paragraphe suivant.

\medskip\noindent
Supposons qu'il existe une immersion $j : X \to \mathbf{P}^n_A$ sur $A$.
Soit $\overline{X}$ l'image schématique de $j$.
Puisque $j$ est un morphisme quasi-compact
(Schémas, lemme \ref{schemes-lemma-quasi-compact-permanence}),
on voit que $j : X \to \overline{X}$ est une immersion ouverte
(Morphismes, lemme \ref{morphisms-lemma-quasi-compact-immersion}).
Ainsi, le changement de base $j_0 : X_0 \to \overline{X}_0$
est également une immersion ouverte.
Par conséquent, $j_0(Y) \subset \overline{X}_0$ est ouvert.
Il est aussi fermé d'après Morphismes, lemme
\ref{morphisms-lemma-image-proper-scheme-closed}.
Supposons le lemme vrai pour $(\overline{X}, j_0(Y))$.
Soit $\overline{W} \subset \overline{X}$ le
sous-schéma ouvert et fermé correspondant, propre sur $A$,
tel que $j_0(Y) = \overline{W}_0$.
Alors $T = \overline{W} \setminus j(X)$ est fermé dans $\overline{W}$,
et a donc une image fermée dans $\Spec(A)$, par propreté de $\overline{W}$
sur $A$. Puisque $(A, I)$ est hensélien, on voit que si $T$
est non vide, il existe un point de $T$ s'envoyant dans $\Spec(A/I)$.
C'est impossible parce que $j_0(Y) = \overline{W}_0$ est contenu dans $j(X)$.
Ainsi, $\overline{W}$ est contenu dans $j(X)$ et l'on peut
prendre pour $W \subset X$ l'unique sous-schéma ouvert et fermé
s'envoyant isomorphiquement sur $\overline{W}$ par $j$.
On se ramène ainsi au cas décrit au paragraphe suivant.

\medskip\noindent
Supposons que $X \subset \mathbf{P}^n_A$ soit un sous-schéma fermé.
Alors $X \to \Spec(A)$ est un morphisme propre.
Soit $Z = X_0 \setminus Y$. C'est un sous-schéma ouvert et fermé
de $X_0$, et $X_0 = Y \amalg Z$.
Soit $X \to X' \to \Spec(A)$ la factorisation de Stein du
théorème \ref{theorem-stein-factorization-general}.
Soient $Y' \subset X'_0$ et $Z' \subset X'_0$ les images de
$Y$ et de $Z$.
Puisque les fibres de $X \to Z$ sont géométriquement connexes,
on voit que $Y' \cap Z' = \emptyset$.
Ainsi, $X'_0 = Y' \amalg Z'$, car $X \to X'$ est surjectif.
Puisque $X' \to \Spec(A)$ est entier, on voit que
$X'$ est le spectre d'une $A$-algèbre entière sur $A$.
Rappelons que les sous-ensembles ouverts et fermés des spectres correspondent
$1$ à $1$ aux idempotents de l'anneau correspondant, voir
Algèbre, lemme \ref{algebra-lemma-disjoint-decomposition}.
Par conséquent, d'après
Compléments d'algèbre, lemme \ref{more-algebra-lemma-characterize-henselian-pair},
on voit que l'on peut écrire $X' = W' \amalg V'$
avec $W'$ et $V'$ ouverts et fermés, et
avec $Y' = W'_0$ et $Z' = V'_0$.
Soit $W$ l'image réciproque du premier de ces sous-schémas dans $X$,
ce qui conclut la preuve.
\end{proof}









\section{Stratification par platitude générique}
\label{section-generic-flatness-stratification}

\noindent
On peut utiliser la platitude générique pour construire une
stratification de la base telle qu'un module donné
devienne plat sur les strates.

\begin{lemma}[Stratification par platitude générique]
\label{lemma-generic-flatness-stratification}
Soit $f : X \to S$ un morphisme de présentation finie entre des schémas quasi-compacts
et quasi-séparés. Soit $\mathcal{F}$ un $\mathcal{O}_X$-module
de présentation finie. Alors il existe un $t \geq 0$ et des sous-schémas
fermés
$$
S \supset S_0 \supset S_1 \supset \ldots \supset S_t = \emptyset
$$
tels que $S_i \to S$ soit défini par un faisceau d'idéaux de type fini,
que $S_0 \subset S$ soit un épaississement et que le tiré en arrière de $\mathcal{F}$ sur
$X \times_S (S_i \setminus S_{i + 1})$ soit plat sur $S_i \setminus S_{i + 1}$.
\end{lemma}

\begin{proof}
On peut trouver un diagramme cartésien
$$
\xymatrix{
X \ar[d] \ar[r] & X_0 \ar[d] \\
S \ar[r] & S_0
}
$$
et un $\mathcal{O}_{X_0}$-module de présentation finie $\mathcal{F}_0$
dont le tiré en arrière est $\mathcal{F}$, où $X_0$ et $S_0$ sont de
type fini sur $\mathbf{Z}$. Voir
Limites, proposition \ref{limits-proposition-approximate}, et
les lemmes \ref{limits-lemma-descend-finite-presentation} et
\ref{limits-lemma-descend-modules-finite-presentation}.
Ainsi, on peut supposer que $X$ et $S$ sont de type fini sur $\mathbf{Z}$
et que $\mathcal{F}$ est un $\mathcal{O}_X$-module cohérent.

\medskip\noindent
Supposons que $X$ et $S$ soient de type fini sur $\mathbf{Z}$
et que $\mathcal{F}$ soit un $\mathcal{O}_X$-module cohérent.
Dans ce cas, tout idéal quasi-cohérent est de type fini ; ainsi,
on n'a pas à vérifier que $S_i$ est défini
par un idéal de type fini. Posons $S_0 = S_{red}$, le réduit de $S$.
D'après la platitude générique telle qu'énoncée dans Morphismes, proposition
\ref{morphisms-proposition-generic-flatness-reduced},
il existe un ouvert dense $U_0 \subset S_0$ tel que le tiré en arrière de $\mathcal{F}$
sur $X \times_S U_0$ soit plat sur $U_0$.
Soit $S_1 \subset S_0$ le sous-schéma fermé réduit dont
le sous-ensemble fermé sous-jacent est $S \setminus U_0$. On continue de cette
manière, pourvu que $S_1 \not = \emptyset$, afin d'obtenir
$S_0 \supset S_1 \supset \ldots$. Comme $S$
est noethérien, toute chaîne descendante de sous-ensembles fermés se stabilise ;
on voit donc que $S_t = \emptyset$ pour un certain $t \geq 0$.
\end{proof}

\begin{lemma}
\label{lemma-generic-flatness-stratification-scheme}
Soit $f : X \to S$ un morphisme de présentation finie entre des schémas quasi-compacts
et quasi-séparés. Alors il existe un $t \geq 0$ et des sous-schémas
fermés
$$
S \supset S_0 \supset S_1 \supset \ldots \supset S_t = \emptyset
$$
tels que $S_i \to S$ soit défini par un faisceau d'idéaux de type fini,
que $S_0 \subset S$ soit un épaississement et que
$X \times_S (S_i \setminus S_{i + 1})$ soit plat sur $S_i \setminus S_{i + 1}$.
\end{lemma}

\begin{proof}
Appliquons le lemme \ref{lemma-generic-flatness-stratification}
avec $\mathcal{F} = \mathcal{O}_X$.
\end{proof}

\begin{lemma}[Longke Tang]
\label{lemma-flatness-criterion}
Soit $f : X \to S$ un morphisme surjectif de présentation finie.
Soit $M$ un objet de $D_\QCoh(\mathcal{O}_S)$ tel que
$H^i(M) = 0$ pour $i > 0$. Les conditions suivantes sont équivalentes :
\begin{enumerate}
\item $M$ est isomorphe à un $\mathcal{O}_S$-module plat
placé en degré $0$,
\item $Lf^*M$ est isomorphe à un $\mathcal{O}_X$-module plat
placé en degré $0$.
\end{enumerate}
\end{lemma}

\begin{proof}
On omet la preuve que (1) implique (2). Supposons (2). Démontrer (1)
est une question locale ; on peut donc supposer que $S$ est affine.
Choisissons une stratification
$S \supset S_0 \supset S_1 \supset \ldots \supset S_t = \emptyset$
comme dans le lemme \ref{lemma-generic-flatness-stratification-scheme}.
Posons $T_i = S_i \setminus S_{i - 1}$ et $Y_i = X \times_S T_i$.
Notons $g_i : T_i \to S$ et $h_i : Y_i \to X$ les morphismes d'inclusion.
Comme $f$ est surjectif, les morphismes $f_i : X_i \to T_i$
sont surjectifs et plats. Puisque
$Lh_i^* \circ Lf^* = Lf_i^* \circ Lg_i^*$, on déduit de (2) que
$Lf_i^* Lg_i^*M$ est isomorphe à un $\mathcal{O}_{Y_i}$-module plat
placé en degré $0$. Comme $f_i$ est plat et surjectif, on voit que
$Lg_i^*M$ est isomorphe à un $\mathcal{O}_{T_i}$-module plat
placé en degré $0$.

\medskip\noindent
On démontrera que cela implique le résultat par récurrence sur $t$.
On utilisera les notations suivantes : $S = \Spec(A)$ et $M$
est l'objet de $D_\QCoh(\mathcal{O}_S)$ correspondant à
$N$ dans $D(A)$, voir Catégories dérivées de schémas, lemme
\ref{perfect-lemma-affine-compare-bounded}.

\medskip\noindent
Cas de base : $t = 1$. Écrivons $S_0 = \Spec(A/I)$, où
$I$ est un idéal de type fini. Alors $I$ est nilpotent, puisque $S_0 = S$
ensemblistement. L'hypothèse est que $N \otimes_A^\mathbf{L} A/I$
a une amplitude de Tor dans $[0, 0]$. D'après Compléments d'algèbre, lemme
\ref{more-algebra-lemma-check-tor-dimension-modulo-nilpotent},
il en va de même pour $N$.

\medskip\noindent
Étape de récurrence. Supposons que $t > 1$. Écrivons $S_{t - 1} = \Spec(A/I)$,
où $I = (f_1, \ldots, f_r)$ est un idéal de type fini.
On raisonne par récurrence sur $r$. Observons que
$N_{f_r} = N \otimes_A^\mathbf{L} A_{f_r}$ dans $D(A_{f_r})$
a une amplitude de Tor dans $[0, 0]$ par l'hypothèse de récurrence
(car la stratification sur l'ouvert principal $D(f_r)$ est de longueur
au plus $t - 1$). D'autre part, considérons
$N' = N \otimes_A^\mathbf{L} A/f_rA$ dans $D(A/f_rA)$.
En posant $S' = \Spec(A/f_rA)$, on obtient une stratification
$$
S' \supset S' \cap S_0 \supset S' \cap S_1 \supset \ldots \supset
S' \cap S_{t - 1} \supset
S' \cap S_t = \emptyset
$$
où $S' \cap S_{t - 1}$ est découpé par $r - 1$ équations dans $S'$.
Comme précédemment, les tirés en arrière dérivés de $M$ sur les parties
$S' \cap S_i \setminus S' \cap S_{i + 1}$ ont une amplitude de Tor
dans $[0, 0]$. Par récurrence sur $r$, on conclut donc que
$N'$ a une amplitude de Tor dans $[0, 0]$.
On conclut finalement que $N$ a une amplitude de Tor dans $[0, 0]$
d'après Compléments d'algèbre, lemme \ref{more-algebra-lemma-f-flat}.
\end{proof}

\begin{lemma}
\label{lemma-cokernel-flat}
Soit $R$ un anneau intègre noethérien. Soient $R \to A \to B$ des morphismes d'anneaux de type fini.
Soit $M$ un $A$-module de type fini et soit $N$ un $B$-module de type fini.
Soit $M \to N$ un morphisme $A$-linéaire. Il existe un élément non nul $f \in R$
tel que le conoyau de $M_f \to N_f$ soit un $R_f$-module plat.
\end{lemma}

\begin{proof}
En remplaçant $M$ par l'image de $M \to N$, on peut supposer que $M \subset N$.
Choisissons une filtration $0 = N_0 \subset N_1 \subset \ldots \subset N_t = N$
telle que $N_i/N_{i - 1} = B/\mathfrak q_i$ pour un idéal premier
$\mathfrak q_i \subset B$, voir
Algèbre, lemme \ref{algebra-lemma-filter-Noetherian-module}.
Posons $M_i = M \cap N_i$. Alors $Q = N/M$ possède une filtration par les sous-modules
$Q_i = N_i/M_i$. Il suffit de montrer que $Q_i/Q_{i - 1}$ devient plat
après localisation en un élément non nul $f$ (puisque les extensions de
modules plats sont plates d'après Algèbre, lemme \ref{algebra-lemma-flat-ses}).
Puisque $Q_i/Q_{i - 1}$ est isomorphe au conoyau
du morphisme $M_i/M_{i - 1} \to N_i/N_{i - 1}$, on se ramène au
cas traité au paragraphe suivant.

\medskip\noindent
Supposons que $B$ soit un anneau intègre et que $M \subset N = B$. Après avoir remplacé $A$ par
l'image de $A$ dans $B$, on peut supposer que $A \subset B$. Par platitude générique,
on peut supposer que $A$ et $B$ sont plats sur $R$
(Algèbre, lemme \ref{algebra-lemma-generic-flatness-Noetherian}).
Il suffit maintenant de montrer que $M \to B$ devient $R$-universellement
injectif après remplacement de $R$ par une localisation
principale (Algèbre, lemme \ref{algebra-lemma-ui-flat-domain}).
Par liberté générique, on peut trouver un élément non nul $g \in A$ tel que
$B_g$ soit un $A_g$-module libre
(Algèbre, lemme \ref{algebra-lemma-generic-flatness-Noetherian}).
On peut donc choisir un facteur direct $M' \subset B_g$ comme $A_g$-module,
qui soit un $A_g$-module libre de type fini, et
tel que $M \to B \to B_g$ se factorise par $M'$.
Il suffit manifestement de montrer que $M \to M'$
devient $R$-universellement injectif après remplacement de
$R$ par une localisation principale.

\medskip\noindent
Écrivons $M' = A_g^{\oplus n}$. Puisque $M \subset M'$ est un $A$-module de type fini,
on voit que $M$ est contenu dans $(1/g^m)A^{\oplus n}$ pour un certain $m \geq 0$.
Après changement de la base choisie de $M'$, on peut supposer que $M \subset A^{\oplus n}$.
Il suffit alors de montrer que $A^{\oplus n}/M$ et $A_g/A$ deviennent
$R$-plats après remplacement de $R$ par une localisation principale. En effet, alors
$M \to A^{\oplus n}$ et $A^{\oplus n} \to A_g^{\oplus n}$ sont
universellement injectifs d'après Algèbre, lemme \ref{algebra-lemma-flat-tor-zero},
et il en va donc de même du composé $M \to M' = A_g^{\oplus n}$.

\medskip\noindent
Par platitude générique (voir la référence ci-dessus), on peut supposer que le
module $A^{\oplus n}/M$ est $R$-plat. Pour le quotient $A_g/A$,
on utilise le fait que
$$
A_g/A = \colim (1/g^m)A/A \cong \colim A/g^mA
$$
et que le module $A/g^mA$ possède une filtration de longueur $m$ dont les
quotients successifs sont isomorphes à $A/gA$. De nouveau par platitude
générique, on peut supposer que $A/gA$ est $R$-plat, donc chaque $A/g^mA$
est $R$-plat, et par suite $A_g/A$ l'est aussi.
\end{proof}

\noindent
Soit $f : X \to Y$ un morphisme de schémas sur un schéma de base $S$.
Soit $Z \subset Y$ l'image schématique de $f$, voir
Morphismes, section \ref{morphisms-section-scheme-theoretic-image}.
Soit $g : S' \to S$ un morphisme de schémas et soit
$f' : X \times_S S' \to Y \times_S S'$ le changement de base de $f$ par $g$.
Il n'est pas toujours vrai que $Z \times_S S' \subset Y \times_S S'$
soit l'image schématique de $f'$.
Disons que {\it la formation de l'image schématique de $f/S$
commute à tout changement de base} si, pour tout $g$ comme ci-dessus,
l'image schématique de $f'$ est égale à $Z \times_S S'$.

\begin{lemma}
\label{lemma-base-change-scheme-theoretic-image}
Soit $S$ un schéma quasi-compact et quasi-séparé.
Soit $f : X \to Y$ un morphisme de schémas sur $S$, les schémas
$X$ et $Y$ étant tous deux de présentation finie sur $S$.
Il existe alors un entier $t \geq 0$ et des sous-schémas fermés
$$
S \supset S_0 \supset S_1 \supset \ldots \supset S_t = \emptyset
$$
possédant les propriétés suivantes :
\begin{enumerate}
\item $S_i \to S$ est défini par un faisceau d'idéaux de type fini,
\item $S_0 \subset S$ est un épaississement, et
\item si $T_i = S_i \setminus S_{i + 1}$ et si $f_i$ est le changement de
base de $f$ à $T_i$, alors :
la formation de l'image schématique de $f_i/T_i$
commute à tout changement de base (voir la discussion précédant le lemme).
\end{enumerate}
\end{lemma}

\begin{proof}
On peut trouver un diagramme commutatif
$$
\xymatrix{
X \ar[d] \ar[r] & Y \ar[d] \ar[r] & S \ar[d] \\
U \ar[r] & V \ar[r] & W
}
$$
à carrés cartésiens tel que $U$, $V$, $W$
soient de type fini sur $\mathbf{Z}$. En effet, on écrit d'abord $S$
comme limite cofiltrée de schémas de type fini sur $\mathbf{Z}$
dont les morphismes de transition sont affines,
en utilisant Limites, proposition \ref{limits-proposition-approximate},
puis on fait descendre le morphisme $X \to Y$ à l'aide de
Limites, lemme \ref{limits-lemma-descend-finite-presentation}.
On se ramène ainsi au cas traité au paragraphe suivant.

\medskip\noindent
Supposons que $S$ soit noethérien.  Dans ce cas, tout idéal quasi-cohérent
est de type fini ; il n'est donc pas nécessaire de vérifier que
$S_i$ est défini par un idéal de type fini. Posons $S_0 = S_{red}$, le
réduit de $S$. Soit $\eta \in S_0$ un point générique
d'une composante irréductible de $S_0$. Par récurrence noethérienne sur
l'espace topologique sous-jacent de $S_0$, on peut supposer que le résultat
vaut pour tout sous-schéma fermé de $S_0$ ne contenant pas $\eta$.
Il suffit donc de montrer qu'il existe un voisinage ouvert
$U_0 \subset S_0$ tel que le changement de base $f_0$ de $f$ à $U_0$
possède la propriété (3).

\medskip\noindent
Soit $R$ un anneau intègre noethérien. Soit $f : X \to Y$ un morphisme
de schémas de type fini sur $R$. D'après la discussion du paragraphe précédent,
il suffit de montrer qu'après remplacement de $R$ par $R_g$ pour un certain $g \in R$
non nul, et de $X$, $Y$ par leurs changements de base à $R_g$,
la formation de l'image schématique de $f/R$ commute
à tout changement de base.

\medskip\noindent
Soit $Y = V_1 \cup \ldots V_n$ un recouvrement ouvert affine.
Posons $U_i = f^{-1}(V_i)$. Si l'assertion est vraie pour chacun des
morphismes $U_i \to V_i$ sur $R$, alors elle est vraie pour $f$.
En effet, l'image schématique de $U_i \to V_i$
est l'intersection de $V_i$ avec l'image schématique de
$f : X \to Y$, d'après Morphismes, lemme
\ref{morphisms-lemma-quasi-compact-scheme-theoretic-image}.
On peut donc supposer que $Y$ est affine.

\medskip\noindent
Soit $X = U_1 \cup \ldots U_n$ un recouvrement ouvert affine.
Alors l'image schématique de $X \to Y$ est la même que
l'image schématique de $\coprod U_i \to Y$.
On peut donc supposer que $X$ est affine.

\medskip\noindent
Écrivons $X = \Spec(A)$ et $Y = \Spec(B)$, et supposons que $f$ corresponde
au morphisme de $R$-algèbres $\varphi : A \to B$.
Alors l'image schématique de $f$ est $\Spec(A/\Ker(\varphi))$,
et il en va de même après changement de base (par un morphisme affine, mais il
suffit de le vérifier pour ceux-ci).
Ainsi, la formation de l'image schématique commute
au changement de base si $\Ker(\varphi \otimes_R R') = \Ker(\varphi) \otimes_R R'$
pour tout morphisme d'anneaux $R \to R'$.

\medskip\noindent
Après remplacement de $R$, $A$, $B$ par $R_g$, $A_g$, $B_g$ pour un
élément non nul convenable $g$ de $R$, on peut supposer que $A$ et $B$ sont plats sur $R$.
D'après le lemme \ref{lemma-cokernel-flat},
on peut également supposer que $B/A$ est un $R$-module plat.
Alors $0 \to \Ker(\varphi) \to A \to B \to B/A \to 0$
est une suite exacte de $R$-modules plats, ce qui implique
l'assertion voulue sur le changement de base.
\end{proof}






\section{Stratification d'un morphisme}
\label{section-stratifying-morphisms}

\noindent
Soit $f : X \to S$ un morphisme de présentation finie entre schémas quasi-compacts
et quasi-séparés.
Dans la section \ref{section-generic-flatness-stratification},
nous avons vu que l'on peut stratifier $S$ de façon que $X$ soit plat
sur les strates. Dans cette section, nous cherchons des stratifications de $S$
et de $X$ telles que les strates obtenues soient lisses ; cela ne fonctionne pas tout à fait,
et nous aurons également besoin d'un changement de base par des morphismes finis localement libres.

\begin{lemma}
\label{lemma-smoothness-stratification-at-generic-point}
Soit $f : X \to S$ un morphisme de schémas de présentation finie.
Soit $\eta \in S$ un point générique d'une composante irréductible de $S$.
Supposons que $S$ soit réduit. Il existe alors
\begin{enumerate}
\item un sous-schéma ouvert $U \subset S$ contenant $\eta$,
\item un morphisme surjectif, universellement injectif et fini localement libre
$V \to U$,
\item un entier $t \geq 0$ et des sous-schémas fermés
$$
X \times_S V \supset Z_0 \supset Z_1 \supset \ldots \supset Z_t = \emptyset
$$
tels que $Z_i \to X \times_S V$ soit défini par un faisceau d'idéaux de type fini,
que $Z_0 \subset X \times_S V$ soit un épaississement et que le morphisme
$Z_i \setminus Z_{i + 1} \to V$ soit lisse.
\end{enumerate}
\end{lemma}

\begin{proof}
Il est clair que l'on peut remplacer $S$ par un voisinage ouvert de $\eta$
et $X$ par sa restriction à cet ouvert.
On peut donc supposer que $S = \Spec(A)$, où $A$ est un anneau réduit,
et que $\eta$ correspond à un idéal premier minimal $\mathfrak p$.
Rappelons que, dans ce cas, l'anneau local $\mathcal{O}_{S, \eta} = A_\mathfrak p$
est égal à $\kappa(\mathfrak p)$, voir
Algèbre, lemme \ref{algebra-lemma-minimal-prime-reduced-ring}.

\medskip\noindent
Appliquons Variétés, lemme \ref{varieties-lemma-smooth-stratification},
au schéma $X_\eta$ sur $k = \kappa(\eta)$.
Notons $k'/k$ l'extension de corps purement inséparable ainsi obtenue.
Au paragraphe suivant, on se ramène au cas $k' = k$.
(Cette étape correspond à la construction du morphisme $V \to U$ de
l'énoncé du lemme ; en particulier, on peut prendre $V = U$
si le corps $\kappa(\mathfrak p)$ est de caractéristique zéro.)

\medskip\noindent
Si le corps $k = \kappa(\mathfrak p)$ est de caractéristique zéro, alors
$k' = k$. Si le corps $k = \kappa(\mathfrak p)$
est de caractéristique $p > 0$, alors $p$ s'envoie sur zéro dans $A_\mathfrak p = \kappa(\mathfrak p)$.
Ainsi, après remplacement de $A$ par une localisation principale (c'est-à-dire après
restriction de $S$), on peut supposer que $p = 0$ dans $A$. Si $k' \not = k$, alors
il existe un élément $\beta \in k'$, $\beta \not \in k$,
tel que $\beta^p \in k$. Après remplacement de $A$ par une localisation
principale, on peut supposer qu'il existe un élément $a \in A$ tel
que $\beta^p = a$. Posons $A' = A[x]/(x^p - a)$.
Alors $S' = \Spec(A') \to \Spec(A) = S$ est fini localement libre,
surjectif et universellement injectif. De plus, si $\mathfrak p' \subset A'$
désigne l'unique idéal premier au-dessus de $\mathfrak p$,
alors $A'_{\mathfrak p'} = k(\beta)$ et $k'/k(\beta)$
est de degré plus petit. Ainsi, après remplacement de $S$ par $S'$
et de $\eta$ par le point $\eta'$ correspondant à $\mathfrak p'$,
on voit que le degré de $k'$ sur le corps résiduel de $\eta$
a diminué. En poursuivant ainsi, une récurrence nous ramène
au cas $k' = \kappa(\mathfrak p) = \kappa(\eta)$.

\medskip\noindent
On peut donc supposer que $S$ est affine et réduit, et qu'il existe un entier $t \geq 0$
et des sous-schémas fermés
$$
X_\eta \supset Z_{\eta, 0} \supset Z_{\eta, 1}
\supset \ldots \supset Z_{\eta, t} = \emptyset
$$
tels que $Z_{\eta, 0} = (X_\eta)_{red}$ et que
$Z_{\eta, i} \setminus Z_{\eta, i + 1}$
soit lisse sur $\eta$ pour tout $i$.
Rappelons que $\kappa(\eta) = \kappa(\mathfrak p) = A_\mathfrak p$ est
la limite inductive filtrante des $A_a$ pour $a \in A$, $a \not \in \mathfrak p$.
Voir Algèbre, lemme \ref{algebra-lemma-localization-colimit}.
On peut donc faire descendre le diagramme ci-dessus en un diagramme correspondant
sur $\Spec(A_a)$ pour un certain $a \in A$, $a \not \in \mathfrak p$.
Plus précisément, après remplacement de $S$ par $\Spec(A_a)$,
on peut supposer qu'il existe un entier $t \geq 0$ et des sous-schémas fermés
$$
X \supset Z_0 \supset Z_1 \supset \ldots \supset Z_t = \emptyset
$$
tels que $Z_i \to X$ soit une immersion fermée de présentation finie,
que $Z_0 \to X$ soit un épaississement et que $Z_i \setminus Z_{i + 1}$
soit lisse sur $S$. Autrement dit, le lemme est démontré.
Plus précisément, on utilise d'abord
Limites, lemme \ref{limits-lemma-descend-finite-presentation},
pour obtenir des morphismes
$$
Z_t \to Z_{t - 1} \to \ldots \to Z_0 \to X
$$
sur $S$, chacun de présentation finie, et dont le
changement de base à $\eta$ redonne les inclusions entre les
sous-schémas fermés donnés ci-dessus. Quitte à rétrécir encore $S$, on peut supposer
que chacun de ces morphismes est une immersion fermée, voir
Limites, lemme
\ref{limits-lemma-descend-closed-immersion-finite-presentation}.
Quitte à rétrécir $S$, on peut supposer que $Z_0 \to X$ est surjectif,
donc est un épaississement, voir
Limites, lemme \ref{limits-lemma-descend-surjective}.
Quitte à rétrécir encore une fois $S$, on peut supposer que
$Z_i \setminus Z_{i + 1} \to S$ est lisse, voir
Limites, lemme \ref{limits-lemma-descend-smooth}.
Cela achève la démonstration.
\end{proof}

\begin{lemma}
\label{lemma-smoothness-stratification}
Soit $f : X \to S$ un morphisme de présentation finie entre schémas quasi-compacts
et quasi-séparés. Il existe alors un entier $t \geq 0$ et des
sous-schémas fermés
$$
S \supset S_0 \supset S_1 \supset \ldots \supset S_t = \emptyset
$$
tels que
\begin{enumerate}
\item $S_i \to S$ soit défini par un faisceau d'idéaux de type fini,
\item $S_0 \subset S$ soit un épaississement,
\item pour chaque $i$, il existe un morphisme fini localement libre et surjectif
$T_i \to S_i \setminus S_{i + 1}$,
\item pour chaque $i$, il existe un entier $t_i \geq 0$ et des sous-schémas fermés
$$
X_i = X \times_S T_i \supset Z_{i, 0}
\supset Z_{i, 1} \supset \ldots \supset Z_{i, t_i} = \emptyset
$$
tels que $Z_{i, j} \to X_i$ soit défini par un faisceau d'idéaux de type fini,
que $Z_{i, 0} \subset X_i$ soit un épaississement et que le morphisme
$Z_{i, j} \setminus Z_{i, j + 1} \to T_i$ soit lisse.
\end{enumerate}
\end{lemma}

\begin{proof}
On peut trouver un diagramme cartésien
$$
\xymatrix{
X \ar[d] \ar[r] & X_0 \ar[d] \\
S \ar[r] & S_0
}
$$
tel que $X_0$ et $S_0$ soient de type fini sur $\mathbf{Z}$. Voir
Limites, proposition \ref{limits-proposition-approximate} et
lemme \ref{limits-lemma-descend-finite-presentation}.
On peut donc supposer que $X$ et $S$ sont de type fini sur $\mathbf{Z}$.
En effet, une solution du problème posé par le lemme pour $X_0 \to S_0$
donnera par changement de base une solution sur $S$ ; détails omis.

\medskip\noindent
Supposons que $X$ et $S$ soient de type fini sur $\mathbf{Z}$.
Dans ce cas, tout idéal quasi-cohérent est de type fini ; il n'est donc
pas nécessaire de vérifier que $S_i$ est défini
par un idéal de type fini. Posons $S_0 = S_{red}$, le réduit de $S$.
Soit $\eta \in S_0$ un point générique d'une composante irréductible.
D'après le lemme \ref{lemma-smoothness-stratification-at-generic-point},
on peut trouver un sous-schéma ouvert $U \subset S_0$,
un morphisme surjectif, universellement injectif et fini localement libre
$V \to U$, un entier $t_0 \geq 0$ et des sous-schémas fermés
$$
X \times_S V \supset Z_{0, 0} \supset Z_{0, 1} \supset \ldots \supset
Z_{0, t_0} = \emptyset
$$
tels que $Z_{0, i} \to X \times_S V$ soit défini par un faisceau d'idéaux de type fini,
que $Z_{0, 0} \subset X \times_S V$ soit un épaississement et que le morphisme
$Z_{0, i} \setminus Z_{0, i + 1} \to V$ soit lisse.
On munit alors $S_1 \subset S_0$ de la structure de sous-schéma réduit induit
sur $S_0 \setminus U$. Par récurrence noethérienne sur l'espace
topologique sous-jacent de $S$, on peut supposer que le lemme vaut pour
$X \times_S S_1 \to S_1$. Cela fournit un entier $t \geq 1$ et
$$
S_1 = S_1 \supset S_2 \supset \ldots \supset S_t = \emptyset
$$
ainsi que des $t_i$ et des $Z_{i, j}$ comme dans l'énoncé du lemme.
Cela démontre le lemme.
\end{proof}


















\section{Amélioration des morphismes de dimension relative égale à un}
\label{section-make-good-curves}

\noindent
On peut rendre toute courbe lisse et projective après extension du
corps de base, compactification et normalisation. Cela implique aussi des
résultats sur les morphismes de type fini dont les fibres génériques sont de dimension $1$.

\begin{lemma}
\label{lemma-make-good-curves}
Soit $f : X \to S$ un morphisme de schémas. Soit $\eta \in S$ un
point générique d'une composante irréductible de $S$. Supposons que $f$ soit
séparé, de présentation finie, et que $\dim(X_\eta) \leq 1$.
Il existe alors un diagramme commutatif
$$
\xymatrix{
\overline{Y}_1 \amalg \ldots \amalg \overline{Y}_n \ar[rd] &
Y_1 \amalg \ldots \amalg Y_n \ar[r]_-\nu \ar[d] \ar[l]^j &
X_V \ar[r] \ar[d] &
X_U \ar[r] \ar[d] &
X \ar[d]^f \\
& T_1 \amalg \ldots \amalg T_n \ar[r] &
V \ar[r] &
U \ar[r] &
S
}
$$
de schémas possédant les propriétés suivantes :
\begin{enumerate}
\item $U \subset S$ est un voisinage ouvert de $\eta$,
\item $V \to U$ est un morphisme fini, surjectif et universellement injectif,
\item $X_U = U \times_S X$ et $X_V = V \times_S X$ sont les changements de base,
\item $\nu$ est fini, surjectif, et il existe un ouvert $W \subset X_V$
tel que
\begin{enumerate}
\item $W$ soit dense dans toutes les fibres de $X_V \to V$,
\item $\nu^{-1}(W) \cap Y_i$ soit dense dans toutes les fibres de $Y_i \to T_i$, et
\item $\nu^{-1}(W) \to W$ soit un épaississement,
\end{enumerate}
\item $j$ est une immersion ouverte,
\item $T_i \to V$ est fini \'etale,
\item $Y_i \to T_i$ est surjectif et lisse,
\item $\overline{Y}_i \to T_i$ est lisse, propre, à fibres géométriquement
connexes de dimension $\leq 1$.
\end{enumerate}
\end{lemma}

\begin{proof}
Il est clair que l'on peut remplacer $S$ par un voisinage ouvert de $\eta$
et $X$ par sa restriction à cet ouvert.
On peut de plus remplacer $S$ par son réduit et $X$ par
le changement de base à ce réduit.
On peut donc supposer que $S = \Spec(A)$, où $A$ est un anneau réduit,
et que $\eta$ correspond à un idéal premier minimal $\mathfrak p$.
Rappelons que, dans ce cas, l'anneau local $\mathcal{O}_{S, \eta} = A_\mathfrak p$
est égal à $\kappa(\mathfrak p)$, voir
Algèbre, lemme \ref{algebra-lemma-minimal-prime-reduced-ring}.

\medskip\noindent
Appliquons Variétés, lemme
\ref{varieties-lemma-dim-1-projective-completion-after-insep},
au schéma $X_\eta$ sur $k = \kappa(\eta)$.
Notons $k'/k$ l'extension de corps purement inséparable ainsi obtenue.
Au paragraphe suivant, on se ramène au cas $k' = k$.
(Cette étape correspond à la construction du morphisme $V \to U$ de
l'énoncé du lemme ; en particulier, on peut prendre $V = U$
si le corps $\kappa(\mathfrak p)$ est de caractéristique zéro.)

\medskip\noindent
Si le corps $k = \kappa(\mathfrak p)$ est de caractéristique zéro, alors
$k' = k$. Si le corps $k = \kappa(\mathfrak p)$
est de caractéristique $p > 0$, alors $p$ s'envoie sur zéro dans $A_\mathfrak p = \kappa(\mathfrak p)$.
Ainsi, après remplacement de $A$ par une localisation principale (c'est-à-dire après
restriction de $S$), on peut supposer que $p = 0$ dans $A$. Si $k' \not = k$, alors
il existe un élément $\beta \in k'$, $\beta \not \in k$,
tel que $\beta^p \in k$. Après remplacement de $A$ par une localisation
principale, on peut supposer qu'il existe un élément $a \in A$ tel
que $\beta^p = a$. Posons $A' = A[x]/(x^p - a)$.
Alors $S' = \Spec(A') \to \Spec(A) = S$ est fini, surjectif et
universellement injectif. De plus, si $\mathfrak p' \subset A'$
désigne l'unique idéal premier au-dessus de $\mathfrak p$,
alors $A'_{\mathfrak p'} = k(\beta)$ et $k'/k(\beta)$
est de degré plus petit. Ainsi, après remplacement de $S$ par $S'$
et de $\eta$ par le point $\eta'$ correspondant à $\mathfrak p'$,
on voit que le degré de $k'$ sur le corps résiduel de $\eta$
a diminué. En poursuivant ainsi, une récurrence nous ramène
au cas $k' = \kappa(\mathfrak p) = \kappa(\eta)$.

\medskip\noindent
On peut donc supposer que $S$ est affine et réduit, et que l'on dispose d'un diagramme
$$
\xymatrix{
\overline{Y}_{1, \eta} \amalg \ldots \amalg \overline{Y}_{n, \eta} \ar[rd] &
Y_{1, \eta} \amalg \ldots \amalg Y_{n, \eta} \ar[r]_-\nu \ar[d] \ar[l]^j &
X_\eta \ar[d] \\
& \Spec(k_1) \amalg \ldots \amalg \Spec(k_n) \ar[r] &
\eta
}
$$
de schémas possédant les propriétés suivantes :
\begin{enumerate}
\item $\nu$ est la normalisation de $X_\eta$,
\item $j$ est une immersion ouverte d'image dense,
\item $k_i/\kappa(\eta)$
est une extension finie séparable pour $i = 1, \ldots, n$,
\item $\overline{Y}_{i, \eta}$ est lisse, projectif et
géométriquement irréductible de dimension $\leq 1$ sur $k_i$.
\end{enumerate}
Rappelons que $\kappa(\eta) = \kappa(\mathfrak p) = A_\mathfrak p$ est
la limite inductive filtrante des $A_a$ pour $a \in A$, $a \not \in \mathfrak p$.
Voir Algèbre, lemme \ref{algebra-lemma-localization-colimit}.
On peut donc faire descendre le diagramme ci-dessus en un diagramme correspondant
sur $\Spec(A_a)$ pour un certain $a \in A$, $a \not \in \mathfrak p$.
Plus précisément, après remplacement de $S$ par $\Spec(A_a)$,
on peut supposer que l'on dispose d'un diagramme commutatif
$$
\xymatrix{
\overline{Y}_1 \amalg \ldots \amalg \overline{Y}_n \ar[rd] &
Y_1 \amalg \ldots \amalg Y_n \ar[r]_-\nu \ar[d] \ar[l]^j &
X \ar[d] \\
& T_1 \amalg \ldots \amalg T_n \ar[r] & S
}
$$
de schémas dont le changement de base à $\eta$ est le diagramme ci-dessus
et qui possèdent les propriétés suivantes :
\begin{enumerate}
\item $\nu$ est un morphisme fini et surjectif,
\item $j$ est une immersion ouverte,
\item $T_i \to S$ est fini \'etale pour $i = 1, \ldots, n$,
\item $Y_i \to T_i$ est lisse et surjectif,
\item $\overline{Y}_i \to T_i$ est lisse et propre, à fibres
géométriquement connexes de dimension $\leq 1$.
\end{enumerate}
Pour cela, on utilise d'abord
Limites, lemme \ref{limits-lemma-descend-finite-presentation},
afin d'obtenir un diagramme dont le changement de base redonne le diagramme précédent.
On utilise ensuite Limites, lemmes \ref{limits-lemma-descend-etale},
\ref{limits-lemma-descend-smooth},
\ref{limits-lemma-descend-finite-finite-presentation},
\ref{limits-lemma-limit-affine},
\ref{limits-lemma-descend-open-immersion},
\ref{limits-lemma-eventually-proper}, et
\ref{limits-lemma-descend-surjective},
pour obtenir que $\nu$ est fini et surjectif, que $j$ est une immersion ouverte, que $T_i \to S$ est fini
\'etale, que $Y_i \to T$ est lisse, et que $\overline{Y}_i \to T_i$ est propre et
lisse. Puisque $Y_i$ ne peut être vide, puisque les morphismes lisses
sont ouverts, et puisque $T_i \to S$ est fini \'etale, quitte à rétrécir $S$,
on peut supposer que $Y_i \to T_i$ est surjectif. Enfin, la fibre de
$\overline{Y}_i \to T_i$ au-dessus de l'unique point $\eta_i = \Spec(k_i)$ 
de $T_i$ au-dessus de $\eta$ est géométriquement connexe.
Par un nouveau rétrécissement, on peut donc supposer qu'il en va de même
pour toutes les fibres, voir
lemme \ref{lemma-proper-flat-geom-red}.

\medskip\noindent
Il reste à démontrer l'existence d'un ouvert $W \subset X$
satisfaisant (a), (b) et (c). Puisque $\nu_\eta : \coprod Y_{i, \eta} \to X_\eta$
est le morphisme de normalisation, on sait d'après
Variétés, lemme \ref{varieties-lemma-normalization-locally-algebraic},
qu'il existe un ouvert dense $W_\eta \subset X_\eta$
tel que $\nu^{-1}(W_\eta) \to W_\eta$ soit égal à
l'inclusion du réduit de $W_\eta$ dans $W_\eta$.
Soit $W \subset X$ un ouvert quasi-compact dont la fibre
au-dessus de $\eta$ est l'ouvert $W_\eta$ que l'on vient de trouver.
Après remplacement de $A = \Gamma(S, \mathcal{O}_S)$
par une autre localisation, on peut supposer que $\nu^{-1}(W) \to W$
est une immersion fermée, voir Limites, lemme
\ref{limits-lemma-descend-closed-immersion-finite-presentation}.
Comme $\nu$ est aussi surjectif, on en conclut que
$\nu^{-1}(W) \to W$ est un épaississement.
Posons $W_i = \nu^{-1}(W) \cap Y_i$.
Quitte à rétrécir encore $S$, on peut supposer que $W_i \to T_i$ est
surjectif pour tout $i$ (par le même argument que ci-dessus).
On en déduit que $W_i \subset Y_i$
est dense dans toutes les fibres de $Y_i \to T_i$, puisque $Y_i \to T_i$ a des
fibres géométriquement irréductibles.
Comme $\nu$ est fini et surjectif, il en résulte alors
que $W = \nu(\nu^{-1}(W))$ est également dense dans toutes les fibres
de $X \to S$.
\end{proof}











\section{Descente des morphismes séparés localement quasi-finis}
\label{section-separated-locally-quasi-finite}

\noindent
Dans cette section, nous montrons que « les morphismes séparés localement quasi-finis
satisfont à la descente pour les recouvrements fppf ». Voir Descente, définition
\ref{descent-definition-descending-types-morphisms} pour la terminologie.
Ce résultat se trouve dans le remarquable article de Raynaud et Gruson
(à bien des égards), dissimulé dans la démonstration
de \cite[Lemme 5.7.1]{GruRay}.
On le trouve également dans \cite{Murre-representation}, et dans
\cite[Expos\'e X, Lemme 5.4]{SGA3},
sous l'hypothèse supplémentaire
que le morphisme est localement de présentation finie.
Voici l'énoncé formel.

\begin{lemma}
\label{lemma-separated-locally-quasi-finite-morphisms-fppf-descend}
Soit $S$ un schéma.
Soit $\{X_i \to S\}_{i\in I}$ un recouvrement fppf, voir
Topologies, définition \ref{topologies-definition-fppf-covering}.
Soit $(V_i/X_i, \varphi_{ij})$ une donnée de descente
relative à $\{X_i \to S\}$. Si chaque morphisme
$V_i \to X_i$ est séparé et localement quasi-fini,
alors la donnée de descente est effective.
\end{lemma}

\begin{proof}
Être séparé et être localement quasi-fini
sont des propriétés des morphismes de schémas
qui sont préservées par tout changement de base, voir
Schémas, lemme \ref{schemes-lemma-separated-permanence} et
Morphismes, lemme \ref{morphisms-lemma-base-change-quasi-finite}.
Par conséquent, Descente, lemme \ref{descent-lemma-descending-types-morphisms},
s'applique, et il suffit de démontrer l'assertion du lemme
lorsque le recouvrement fppf est donné par un unique
morphisme $\{X \to S\}$ plat, surjectif et de présentation finie entre schémas affines.
Écrivons $X = \Spec(A)$ et $S = \Spec(R)$, de sorte
que $R \to A$ soit un morphisme d'anneaux fidèlement plat.
Soit $(V, \varphi)$ une donnée de descente relativement à $X$ sur $S$
et supposons que $\pi : V \to X$ soit séparé et
localement quasi-fini.

\medskip\noindent
Soit $W^1 \subset V$ un ouvert affine quelconque.
Considérons $W = \text{pr}_1(\varphi(W^1 \times_S X)) \subset V$.
Voici un diagramme
$$
\xymatrix@C=1.2pc{
W^1 \times_S X \ar[rrrrr] \ar[ddd] \ar[rd]
& & & & &
\varphi(W^1 \times_S X) \ar[ddd] \ar[ld] \\
& V \times_S X \ar[rrr]^\varphi \ar[rd] \ar[dd]
& & &
X \times_S V \ar[ld] \ar[dd] & \\
& &
X \times_S X \ar[r]^1 \ar[d]_{\text{pr}_0}
&
X \times_S X \ar[d]^{\text{pr}_1}
& & \\
W^1 \ar[r] &
V \ar[r] &
X &
X &
V \ar[l] &
W \ar[l]
}
$$
Or, puisque $X \to S$ est plat et de présentation finie, il
est universellement ouvert (Morphismes, lemme \ref{morphisms-lemma-fppf-open}).
Nous en déduisons que $W$ est ouvert. De plus, il est
également clair que $W$ est quasi-compact, et que
$W^1 \subset W$. En outre, nous remarquons que
$\varphi(W \times_S X) = X \times_S W$ par la condition de cocycle
pour $\varphi$. Nous obtenons donc une nouvelle donnée de descente
$(W, \varphi')$ en restreignant $\varphi$ à $W \times_S X$.
Remarquons que le morphisme $W \to X$ est quasi-compact, séparé
et localement quasi-fini. Cela implique qu'il est
séparé et quasi-fini par définition. Il est donc quasi-affine d'après le
lemme \ref{lemma-quasi-finite-separated-quasi-affine}.
Ainsi, d'après
Descente, lemme \ref{descent-lemma-quasi-affine},
la donnée de descente
$(W, \varphi')$ est effective.

\medskip\noindent
Autrement dit, on trouve un recouvrement ouvert
$V = \bigcup W_i$ par des ouverts quasi-compacts $W_i$ qui sont
stables par le morphisme de descente $\varphi$.
De plus, pour chacun de ces ouverts quasi-compacts $W \subset V$,
la donnée de descente correspondante $(W, \varphi')$ est effective.
Cela signifie que la donnée de descente initiale est effective, par recollement des
schémas obtenus en descendant les ouverts $W_i$ ; voir
Descente, lemme \ref{descent-lemma-effective-for-fpqc-is-local-upstairs}.
\end{proof}






\section{Présentation finie relative}
\label{section-finite-type-finite-presentation}

\noindent
Soit $R \to A$ un homomorphisme d'anneaux de type fini. Soit $M$ un $A$-module. Dans
Compléments d'algèbre,
section \ref{more-algebra-section-relative-finite-presentation},
nous avons défini ce que signifie pour $M$ d'être de présentation finie relativement à $R$.
Nous avons aussi démontré que cette notion se localise bien et se recolle.
Nous pouvons donc définir la notion globale correspondante comme suit.

\begin{definition}
\label{definition-relatively-finitely-presented-sheaf}
Soit $f : X \to S$ un morphisme de schémas localement de type fini.
Soit $\mathcal{F}$ un $\mathcal{O}_X$-module quasi-cohérent. On dit que
$\mathcal{F}$ est {\it finiment présenté relativement à $S$} ou
{\it de présentation finie relativement à $S$}
s'il existe un recouvrement par des ouverts affines $S = \bigcup V_i$ et,
pour tout $i$, un recouvrement par des ouverts affines
$f^{-1}(V_i) = \bigcup_j U_{ij}$ tels que $\mathcal{F}(U_{ij})$
soit un $\mathcal{O}_X(U_{ij})$-module de présentation finie relativement
à $\mathcal{O}_S(V_i)$.
\end{definition}

\noindent
Remarquons que cela implique que $\mathcal{F}$ est un
$\mathcal{O}_X$-module de type fini. Si $X \to S$ est seulement localement de type
fini, alors $\mathcal{F}$ peut être de présentation finie relativement
à $S$, sans que $X \to S$ soit localement de présentation finie.
Nous verrons que $X \to S$ est localement de présentation finie si
et seulement si $\mathcal{O}_X$ est de présentation finie relativement à $S$.

\begin{lemma}
\label{lemma-relative-finite-presentation-characterize}
Soit $f : X \to S$ un morphisme de schémas localement de type fini.
Soit $\mathcal{F}$ un $\mathcal{O}_X$-module quasi-cohérent. Les conditions suivantes
sont équivalentes :
\begin{enumerate}
\item $\mathcal{F}$ est de présentation finie relativement à $S$,
\item pour tous ouverts affines $U \subset X$, $V \subset S$
tels que $f(U) \subset V$, le $\mathcal{O}_X(U)$-module $\mathcal{F}(U)$
est de présentation finie relativement à $\mathcal{O}_S(V)$.
\end{enumerate}
De plus, si ces conditions sont satisfaites, alors, pour tous sous-schémas ouverts
$U \subset X$ et $V \subset S$ tels que $f(U) \subset V$,
la restriction $\mathcal{F}|_U$ est de présentation finie relativement à $V$.
\end{lemma}

\begin{proof}
La dernière assertion résulte clairement de l'équivalence de (1) et (2).
Il est également clair que (2) implique (1). Supposons (1) satisfaite.
Soient $S = \bigcup V_i$ et $f^{-1}(V_i) = \bigcup U_{ij}$ des
recouvrements par des ouverts affines comme dans la
définition \ref{definition-relatively-finitely-presented-sheaf}.
Soient $U \subset X$ et $V \subset S$ comme dans (2).
D'après Compléments d'algèbre, lemme
\ref{more-algebra-lemma-glue-relative-finite-presentation},
il suffit de trouver un recouvrement $U = \bigcup U_k$ de $U$ par des ouverts principaux
tel que $\mathcal{F}(U_k)$ soit de présentation finie relativement à
$\mathcal{O}_S(V)$. Autrement dit, pour tout $u \in U$, il suffit
de trouver un ouvert affine principal $u \in U' \subset U$ tel que
$\mathcal{F}(U')$ soit de présentation finie relativement à $\mathcal{O}_S(V)$.
Choisissons $i$ tel que $f(u) \in V_i$, puis $j$ tel que
$u \in U_{ij}$. D'après
Schémas, lemme \ref{schemes-lemma-standard-open-two-affines},
on peut trouver $v \in V' \subset V \cap V_i$ qui est un ouvert affine principal
dans $V'$ et dans $V_i$. Alors $f^{-1}V'  \cap U$ et, respectivement,\ $f^{-1}V' \cap U_{ij}$
sont des ouverts affines principaux de $U$ et, respectivement,\ de $U_{ij}$.
En appliquant de nouveau le lemme, on peut trouver
$u \in U' \subset f^{-1}V' \cap U \cap U_{ij}$ qui est un ouvert affine principal
à la fois dans $f^{-1}V'  \cap U$ et dans $f^{-1}V' \cap U_{ij}$.
Ainsi, $U'$ est aussi un ouvert affine principal de $U$ et de $U_{ij}$.
D'après Compléments d'algèbre, lemme
\ref{more-algebra-lemma-localize-relative-finite-presentation},
l'hypothèse selon laquelle $\mathcal{F}(U_{ij})$ est de présentation finie
relativement à $\mathcal{O}_S(V_i)$ implique que
$\mathcal{F}(U')$ est de présentation finie relativement à $\mathcal{O}_S(V_i)$.
Puisque $\mathcal{O}_X(U') =
\mathcal{O}_X(U') \otimes_{\mathcal{O}_S(V_i)} \mathcal{O}_S(V')$,
Compléments d'algèbre, lemme
\ref{more-algebra-lemma-base-change-relative-finite-presentation},
montre que $\mathcal{F}(U')$ est de présentation finie relativement à $\mathcal{O}_S(V')$.
En appliquant encore une fois Compléments d'algèbre, lemme
\ref{more-algebra-lemma-localize-relative-finite-presentation},
nous concluons que
$\mathcal{F}(U')$ est de présentation finie relativement à $\mathcal{O}_S(V)$.
Ceci achève la démonstration.
\end{proof}

\begin{lemma}
\label{lemma-relative-finite-presentation}
Soit $f : X \to S$ un morphisme de schémas localement de type
fini. Soit $\mathcal{F}$ un $\mathcal{O}_X$-module quasi-cohérent.
\begin{enumerate}
\item Si $f$ est localement de présentation finie, alors $\mathcal{F}$
est de présentation finie relativement à $S$ si et seulement si $\mathcal{F}$
est de présentation finie.
\item Le morphisme $f$ est localement de présentation finie si et seulement
si $\mathcal{O}_X$ est de présentation finie relativement à $S$.
\end{enumerate}
\end{lemma}

\begin{proof}
Cela résulte immédiatement des définitions ; voir la
discussion qui suit la définition dans
Compléments d'algèbre, définition
\ref{more-algebra-definition-relatively-finitely-presented}.
\end{proof}

\begin{lemma}
\label{lemma-finite-morphism-relative-finite-presentation}
Soit $\pi : X \to Y$ un morphisme fini entre des schémas localement de type
fini sur un schéma de base $S$. Soit $\mathcal{F}$ un
$\mathcal{O}_X$-module quasi-cohérent. Alors $\mathcal{F}$ est de présentation finie
relativement à $S$ si et seulement si $\pi_*\mathcal{F}$ est de présentation finie
relativement à $S$.
\end{lemma}

\begin{proof}
C'est la traduction du résultat de
Compléments d'algèbre, lemme \ref{more-algebra-lemma-finite-extension},
dans le langage des schémas.
\end{proof}

\begin{lemma}
\label{lemma-base-change-relative-finite-presentation}
Soit $f : X \to S$ un morphisme de schémas localement de type
fini. Soit $\mathcal{F}$ un $\mathcal{O}_X$-module quasi-cohérent.
Soit $S' \to S$ un morphisme de schémas, posons $X' = X \times_S S'$
et notons $\mathcal{F}'$ l'image inverse de $\mathcal{F}$ sur $X'$.
Si $\mathcal{F}$ est de présentation finie relativement à $S$, alors
$\mathcal{F}'$ est de présentation finie relativement à $S'$.
\end{lemma}

\begin{proof}
C'est la traduction du résultat de
Compléments d'algèbre, lemme
\ref{more-algebra-lemma-base-change-relative-finite-presentation},
dans le langage des schémas.
\end{proof}

\begin{lemma}
\label{lemma-pull-relative-finite-presentation}
Soient $X \to Y \to S$ des morphismes de schémas localement de type
fini. Soit $\mathcal{G}$ un $\mathcal{O}_Y$-module quasi-cohérent.
Si $f : X \to Y$ est localement de présentation finie et si
$\mathcal{G}$ est de présentation finie relativement à $S$, alors
$f^*\mathcal{G}$ est de présentation finie relativement à $S$.
\end{lemma}

\begin{proof}
C'est la traduction du résultat de
Compléments d'algèbre, lemme
\ref{more-algebra-lemma-pull-relative-finite-presentation},
dans le langage des schémas.
\end{proof}

\begin{lemma}
\label{lemma-composition-relative-finite-presentation}
Soient $X \to Y \to S$ des morphismes de schémas localement de type
fini. Soit $\mathcal{F}$ un $\mathcal{O}_X$-module quasi-cohérent.
Si $Y \to S$ est localement de présentation finie et si $\mathcal{F}$
est de présentation finie relativement à $Y$, alors $\mathcal{F}$
est de présentation finie relativement à $S$.
\end{lemma}

\begin{proof}
C'est la traduction du résultat de
Compléments d'algèbre, lemme
\ref{more-algebra-lemma-composition-relative-finite-presentation},
dans le langage des schémas.
\end{proof}

\begin{lemma}
\label{lemma-ses-relatively-finite-presentation}
Soit $X \to S$ un morphisme de schémas localement de type fini.
Soit $0 \to \mathcal{F}' \to \mathcal{F} \to \mathcal{F}'' \to 0$
une suite exacte courte de $\mathcal{O}_X$-modules quasi-cohérents.
\begin{enumerate}
\item Si $\mathcal{F}', \mathcal{F}''$ sont de présentation finie relativement à
$S$, alors il en est de même de $\mathcal{F}$.
\item Si $\mathcal{F}'$ est un $\mathcal{O}_X$-module de type fini
et si $\mathcal{F}$ est de présentation finie relativement à $S$, alors
$\mathcal{F}''$ est de présentation finie relativement à $S$.
\end{enumerate}
\end{lemma}

\begin{proof}
C'est la traduction du résultat de
Compléments d'algèbre, lemme
\ref{more-algebra-lemma-ses-relatively-finite-presentation},
dans le langage des schémas.
\end{proof}

\begin{lemma}
\label{lemma-sum-relatively-finite-presentation}
\begin{slogan}
Un facteur direct d'un module hérite de la propriété d'être de présentation
finie relativement à une base.
\end{slogan}
Soit $X \to S$ un morphisme de schémas localement de type fini.
Soient $\mathcal{F}, \mathcal{F}'$ des $\mathcal{O}_X$-modules quasi-cohérents.
Si $\mathcal{F} \oplus \mathcal{F}'$ est de présentation finie relativement à $S$,
alors il en est de même de $\mathcal{F}$ et de $\mathcal{F}'$.
\end{lemma}

\begin{proof}
C'est la traduction du résultat de
Compléments d'algèbre, lemme
\ref{more-algebra-lemma-sum-relatively-finite-presentation},
dans le langage des schémas.
\end{proof}










\section{Pseudo-cohérence relative}
\label{section-relative-pseudo-coherence}

\noindent
Cette section est, pour les schémas, l'analogue de
Compléments d'algèbre, section \ref{more-algebra-section-relative-pseudo-coherent}.
Nous recommandons vivement au lecteur de consulter d'abord cette
section. Bien que nous ayons développé les résultats de cette section
ainsi que ceux sur les complexes pseudo-cohérents dans
Cohomologie, sections \ref{cohomology-section-strictly-perfect},
\ref{cohomology-section-pseudo-coherent},
\ref{cohomology-section-tor} et
\ref{cohomology-section-perfect},
pour des complexes arbitraires de $\mathcal{O}_X$-modules, si
$X$ est un schéma, travailler exclusivement avec des objets de
$D_\QCoh(\mathcal{O}_X)$
simplifie considérablement de nombreux lemmes et arguments, en
ramenant souvent d'emblée le problème à son analogue algébrique.
De plus, l'une des premières choses que nous faisons est de montrer que
la pseudo-cohérence relative implique que les faisceaux de cohomologie
sont quasi-cohérents ; voir le lemme \ref{lemma-relative-pseudo-coherence}.
Pour une première lecture, nous suggérons donc au lecteur de travailler exclusivement
avec des objets de $D_\QCoh(\mathcal{O}_X)$.

\begin{lemma}
\label{lemma-relatively-pseudo-coherent}
Soit $X \to S$ un morphisme de type fini de schémas affines.
Soit $E$ un objet de $D(\mathcal{O}_X)$.
Soit $m \in \mathbf{Z}$.
Les conditions suivantes sont équivalentes :
\begin{enumerate}
\item pour une immersion fermée $i : X \to \mathbf{A}^n_S$,
l'objet $Ri_*E$ de $D(\mathcal{O}_{\mathbf{A}^n_S})$
est $m$-pseudo-cohérent, et
\item pour toute immersion fermée $i : X \to \mathbf{A}^n_S$,
l'objet $Ri_*E$ de $D(\mathcal{O}_{\mathbf{A}^n_S})$
est $m$-pseudo-cohérent.
\end{enumerate}
\end{lemma}

\begin{proof}
Écrivons $S = \Spec(R)$ et $X = \Spec(A)$. Soit $i$ l'immersion correspondant à la surjection
$\alpha : R[x_1, \ldots, x_n] \to A$ et soit $X \to \mathbf{A}^m_S$
correspondant à $\beta : R[y_1, \ldots, y_m] \to A$.
Choisissons $f_j \in R[x_1, \ldots, x_n]$ tels que $\alpha(f_j) = \beta(y_j)$
et $g_i \in R[y_1, \ldots, y_m]$ tels que $\beta(g_i) = \alpha(x_i)$.
On obtient alors un diagramme commutatif
$$
\xymatrix{
R[x_1, \ldots, x_n, y_1, \ldots, y_m]
\ar[d]^{x_i \mapsto g_i} \ar[rr]_-{y_j \mapsto f_j} & &
R[x_1, \ldots, x_n] \ar[d] \\
R[y_1, \ldots, y_m] \ar[rr] & & A
}
$$
qui correspond au diagramme commutatif d'immersions fermées
$$
\xymatrix{
\mathbf{A}^{n + m}_S & \mathbf{A}^n_S \ar[l] \\
\mathbf{A}^m_S \ar[u] & X \ar[u] \ar[l]
}
$$
Il suffit donc de montrer que, pour une immersion fermée
$$
f : \mathbf{A}^m_S \to \mathbf{A}^{n + m}_S
$$
un objet $E$ de $D(\mathcal{O}_{\mathbf{A}^m_S})$ est
$m$-pseudo-cohérent si et seulement si $Rf_*E$ est $m$-pseudo-cohérent.
Cela résulte de
Catégories dérivées de schémas, lemme
\ref{perfect-lemma-closed-push-pseudo-coherent},
et du fait que $f_*\mathcal{O}_{\mathbf{A}^m_S}$ est
un $\mathcal{O}_{\mathbf{A}^{n + m}_S}$-module pseudo-cohérent.
La pseudo-cohérence de $f_*\mathcal{O}_{\mathbf{A}^m_S}$
se démontre directement sans difficulté, mais résulte aussi de
Catégories dérivées de schémas, lemme \ref{perfect-lemma-pseudo-coherent-affine},
et de
Compléments d'algèbre, lemme \ref{more-algebra-lemma-relatively-pseudo-coherent}.
\end{proof}

\noindent
Rappelons que, si $f : X \to S$ est un morphisme de schémas localement de
type fini, alors, pour toute paire d'ouverts affines $U \subset X$ et
$V \subset S$ tels que $f(U) \subset V$, l'homomorphisme d'anneaux
$\mathcal{O}_S(V) \to \mathcal{O}_X(U)$ est de type fini
(Morphismes, lemme \ref{morphisms-lemma-locally-finite-type-characterize}).
Il existe donc toujours des
immersions fermées $U \to \mathbf{A}^n_V$, et la définition suivante a un sens.

\begin{definition}
\label{definition-relative-pseudo-coherence}
Soit $f : X \to S$ un morphisme de schémas localement de type fini.
Soit $E$ un objet de $D(\mathcal{O}_X)$. Soit $\mathcal{F}$ un
$\mathcal{O}_X$-module. Fixons $m \in \mathbf{Z}$.
\begin{enumerate}
\item On dit que $E$ est {\it $m$-pseudo-cohérent relativement à $S$}
s'il existe un recouvrement par des ouverts affines $S = \bigcup V_i$ et,
pour tout $i$, un recouvrement par des ouverts affines $f^{-1}(V_i) = \bigcup U_{ij}$
tel que les conditions équivalentes du
lemme \ref{lemma-relatively-pseudo-coherent}
soient satisfaites pour chacune des paires $(U_{ij} \to V_i, E|_{U_{ij}})$.
\item On dit que $E$ est {\it pseudo-cohérent relativement à $S$}
si $E$ est $m$-pseudo-cohérent relativement à $S$ pour tout $m \in \mathbf{Z}$.
\item On dit que $\mathcal{F}$ est {\it $m$-pseudo-cohérent relativement à $S$} si
$\mathcal{F}$, considéré comme un objet de $D(\mathcal{O}_X)$, est
$m$-pseudo-cohérent relativement à $S$.
\item On dit que $\mathcal{F}$ est {\it pseudo-cohérent relativement à $S$} si
$\mathcal{F}$, considéré comme un objet de $D(\mathcal{O}_X)$, est
pseudo-cohérent relativement à $S$.
\end{enumerate}
\end{definition}

\noindent
Si $X$ est quasi-compact et si $E$ est $m$-pseudo-cohérent relativement à $S$
pour un certain $m$, alors $E$ est borné supérieurement. Si $E$ est
pseudo-cohérent relativement à $S$, alors $E$ a des faisceaux de cohomologie
quasi-cohérents.

\begin{lemma}
\label{lemma-relative-pseudo-coherence}
Soit $f : X \to S$ un morphisme de schémas localement de type fini.
Si $E$ dans $D(\mathcal{O}_X)$ est $m$-pseudo-cohérent relativement à $S$,
alors $H^i(E)$ est un $\mathcal{O}_X$-module quasi-cohérent pour $i > m$.
Si $E$ est pseudo-cohérent relativement à $S$, alors $E$ est un objet de
$D_\QCoh(\mathcal{O}_X)$.
\end{lemma}

\begin{proof}
Choisissons un recouvrement par des ouverts affines $S = \bigcup V_i$ et,
pour tout $i$, un recouvrement par des ouverts affines $f^{-1}(V_i) = \bigcup U_{ij}$
tel que les conditions équivalentes du
lemme \ref{lemma-relatively-pseudo-coherent}
soient satisfaites pour chacune des paires $(U_{ij} \to V_i, E|_{U_{ij}})$.
La quasi-cohérence étant une propriété locale sur $X$, nous pouvons supposer
qu'il existe une immersion fermée $i : X \to \mathbf{A}^n_S$
telle que $Ri_*E$ soit $m$-pseudo-cohérent sur $\mathbf{A}^n_S$.
D'après Catégories dérivées de schémas, lemme \ref{perfect-lemma-pseudo-coherent},
cela signifie que $H^q(Ri_*E)$ est quasi-cohérent pour $q > m$.
Puisque $i_*$ est un foncteur exact, on a $i_*H^q(E) = H^q(Ri_*E)$,
qui est quasi-cohérent sur $\mathbf{A}^n_S$.
D'après Morphismes, lemme \ref{morphisms-lemma-i-star-equivalence},
cela implique que $H^q(E)$ est quasi-cohérent, comme voulu
(strictement parlant, cela implique qu'il existe un
$\mathcal{O}_X$-module quasi-cohérent $\mathcal{F}$ tel que
$i_*\mathcal{F} = i_*H^q(E)$, puis, d'après
Modules, lemme \ref{modules-lemma-i-star-equivalence},
on a $\mathcal{F} \cong H^q(E)$, d'où le résultat).
\end{proof}

\noindent
Montrons maintenant que la condition de pseudo-cohérence relative se localise bien.

\begin{lemma}
\label{lemma-localize-relative-pseudo-coherent}
Soit $S$ un schéma affine. Soit $V \subset S$ un ouvert principal.
Soit $X \to V$ un morphisme de type fini de schémas affines.
Soit $U \subset X$ un ouvert affine. Soit $E$ un objet de
$D(\mathcal{O}_X)$. Si les conditions équivalentes du
lemme \ref{lemma-relatively-pseudo-coherent}
sont satisfaites pour la paire $(X \to V, E)$, alors
les conditions équivalentes du
lemme \ref{lemma-relatively-pseudo-coherent}
sont satisfaites pour la paire $(U \to S, E|_U)$.
\end{lemma}

\begin{proof}
Écrivons $S = \Spec(R)$, $V = D(f)$, $X = \Spec(A)$ et $U = D(g)$.
Supposons que les conditions équivalentes du
lemme \ref{lemma-relatively-pseudo-coherent}
soient satisfaites pour la paire $(X \to V, E)$.

\medskip\noindent
Choisissons une surjection $R_f[x_1, \ldots, x_n] \to A$. Écrivons
$R_f = R[x_0]/(fx_0 - 1)$. Alors $R[x_0, x_1, \ldots, x_n] \to A$
est surjectif, et $R_f[x_1, \ldots, x_n]$ est pseudo-cohérent comme
$R[x_0, \ldots, x_n]$-module. Nous avons donc
$$
X \to \mathbf{A}^n_V \to \mathbf{A}^{n + 1}_S
$$
et nous pouvons appliquer
Catégories dérivées de schémas,
lemme \ref{perfect-lemma-closed-push-pseudo-coherent},
pour conclure que l'image directe $E'$ de $E$ sur $\mathbf{A}^{n + 1}_S$
est $m$-pseudo-cohérente.

\medskip\noindent
Choisissons un élément $g' \in R[x_0, x_1, \ldots, x_n]$ dont l'image est
$g \in A$. Considérons la surjection
$R[x_0, \ldots, x_{n + 1}] \to R[x_0, \ldots, x_n, 1/g']$.
On obtient
$$
\xymatrix{
X \ar[d] & U \ar[d] \ar[l] \ar[dr] \\
\mathbf{A}^{n + 1}_S & D(g')\ar[l] \ar[r] & \mathbf{A}^{n + 2}_S
}
$$
où la flèche inférieure gauche est une immersion ouverte et la flèche inférieure droite
une immersion fermée. On conclut comme précédemment que l'image directe de
$E'|_{D(g')}$ sur $\mathbf{A}^{n + 2}_S$ est $m$-pseudo-cohérente.
Comme c'est également l'image directe de $E|_U$ sur $\mathbf{A}^{n + 2}_S$,
ce qui démontre le lemme.
\end{proof}

\begin{lemma}
\label{lemma-glue-relative-pseudo-coherent}
Soit $X \to S$ un morphisme de type fini de schémas affines. Soit $E$ un
objet de $D(\mathcal{O}_X)$. Soit $m \in \mathbf{Z}$.
Soit $X = \bigcup U_i$ un recouvrement par des ouverts affines principaux.
Les conditions suivantes sont équivalentes :
\begin{enumerate}
\item les conditions équivalentes du
lemme \ref{lemma-relatively-pseudo-coherent}
sont satisfaites pour les paires $(U_i \to S, E|_{U_i})$,
\item les conditions équivalentes du
lemme \ref{lemma-relatively-pseudo-coherent}
sont satisfaites pour la paire $(X \to S, E)$.
\end{enumerate}
\end{lemma}

\begin{proof}
L'implication (2) $\Rightarrow$ (1) est le
lemme \ref{lemma-localize-relative-pseudo-coherent}.
Supposons (1). Écrivons $S = \Spec(R)$ et $X = \Spec(A)$, et
$U_i = D(f_i)$. Écrivons $1 = \sum f_ig_i$ dans $A$.
Considérons les surjections
$$
R[x_i, y_i, z_i] \to R[x_i, y_i, z_i]/(\sum y_iz_i - 1) \to A.
$$
La composée envoie $y_i$ sur $f_i$ et $z_i$ sur $g_i$. Remarquons que
$R[x_i, y_i, z_i]/(\sum y_iz_i - 1)$ est pseudo-cohérent comme
$R[x_i, y_i, z_i]$-module. Il suffit donc de démontrer que
l'image directe de $E$ sur $T = \Spec(R[x_i, y_i, z_i]/(\sum y_iz_i - 1))$
est $m$-pseudo-cohérente ; voir
Catégories dérivées de schémas,
lemme \ref{perfect-lemma-closed-push-pseudo-coherent}.
Pour chaque $i_0$, il suffit de démontrer que la restriction de cette
image directe à
$W_{i_0} = \Spec(R[x_i, y_i, z_i, 1/y_{i_0}]/(\sum y_iz_i - 1))$
est $m$-pseudo-cohérente. Remarquons qu'il existe un diagramme commutatif
$$
\xymatrix{
X \ar[d] & U_{i_0} \ar[l] \ar[d] \\
T & W_{i_0} \ar[l]
}
$$
qui implique que l'image directe de $E$ sur $T$, restreinte à $W_{i_0}$,
est l'image directe de $E|_{U_{i_0}}$ sur $W_{i_0}$. Puisque
$R[x_i, y_i, z_i, 1/y_{i_0}]/(\sum y_iz_i - 1)$ est isomorphe
à un anneau de polynômes sur $R$, cela démontre ce que nous voulons.
\end{proof}

\begin{lemma}
\label{lemma-relative-pseudo-coherence-characterize}
Soit $f : X \to S$ un morphisme de schémas localement de type fini.
Soit $E$ un objet de $D(\mathcal{O}_X)$.
Fixons $m \in \mathbf{Z}$. Les conditions suivantes sont équivalentes :
\begin{enumerate}
\item $E$ est $m$-pseudo-cohérent relativement à $S$,
\item pour tous ouverts affines $U \subset X$ et $V \subset S$
tels que $f(U) \subset V$, les conditions équivalentes du
lemme \ref{lemma-relatively-pseudo-coherent}
sont satisfaites pour la paire $(U \to V, E|_U)$.
\end{enumerate}
De plus, si ces conditions sont satisfaites, alors, pour tous sous-schémas ouverts
$U \subset X$ et $V \subset S$ tels que $f(U) \subset V$,
la restriction $E|_U$ est $m$-pseudo-cohérente relativement à $V$.
\end{lemma}

\begin{proof}
La dernière assertion résulte clairement de l'équivalence de (1) et (2).
Il est également clair que (2) implique (1). Supposons (1) satisfaite.
Soient $S = \bigcup V_i$ et $f^{-1}(V_i) = \bigcup U_{ij}$ des
recouvrements par des ouverts affines comme dans la
définition \ref{definition-relative-pseudo-coherence}.
Soient $U \subset X$ et $V \subset S$ comme dans (2).
D'après le lemme \ref{lemma-glue-relative-pseudo-coherent},
il suffit de trouver un recouvrement $U = \bigcup U_k$ de $U$ par des ouverts principaux
tel que les conditions équivalentes du
lemme \ref{lemma-relatively-pseudo-coherent}
soient satisfaites pour les paires $(U_k \to V, E|_{U_k})$.
Autrement dit, pour tout $u \in U$, il suffit
de trouver un ouvert affine principal $u \in U' \subset U$ tel que
les conditions équivalentes du
lemme \ref{lemma-relatively-pseudo-coherent}
soient satisfaites pour la paire $(U' \to V, E|_{U'})$.
Choisissons $i$ tel que $f(u) \in V_i$, puis $j$ tel que
$u \in U_{ij}$. D'après
Schémas, lemme \ref{schemes-lemma-standard-open-two-affines},
on peut trouver $v \in V' \subset V \cap V_i$ qui est un ouvert affine principal
dans $V'$ et dans $V_i$. Alors $f^{-1}V'  \cap U$ et, respectivement,\ $f^{-1}V' \cap U_{ij}$
sont des ouverts affines principaux de $U$ et, respectivement,\ de $U_{ij}$.
En appliquant de nouveau le lemme, on peut trouver
$u \in U' \subset f^{-1}V' \cap U \cap U_{ij}$ qui est un ouvert affine principal
à la fois dans $f^{-1}V'  \cap U$ et dans $f^{-1}V' \cap U_{ij}$.
Ainsi, $U'$ est aussi un ouvert affine principal de $U$ et de $U_{ij}$.
D'après le lemme \ref{lemma-localize-relative-pseudo-coherent},
l'hypothèse selon laquelle les conditions équivalentes du
lemme \ref{lemma-relatively-pseudo-coherent}
sont satisfaites pour la paire $(U_{ij} \to V_i, E|_{U_{ij}})$
implique que les conditions équivalentes du
lemme \ref{lemma-relatively-pseudo-coherent}
sont satisfaites pour la paire $(U' \to V, E|_{U'})$.
\end{proof}

\noindent
Pour les objets de la catégorie dérivée dont les faisceaux de cohomologie sont
quasi-cohérents, on peut relier la $m$-pseudo-cohérence relative
à la notion définie dans Compléments d'algèbre, définition
\ref{more-algebra-definition-relatively-pseudo-coherent}.
Nous utiliserons le fait que, pour un schéma affine
$U = \Spec(A)$, le foncteur $R\Gamma(U, -)$ induit une équivalence
entre $D_\QCoh(\mathcal{O}_U)$ et $D(A)$ ; voir
Catégories dérivées de schémas, lemme
\ref{perfect-lemma-affine-compare-bounded}.
Ce foncteur est compatible aux images inverses :
si $E$ est un objet de $D_\QCoh(\mathcal{O}_U)$
et $A \to B$ est un homomorphisme d'anneaux correspondant à un morphisme de
schémas affines $g : V = \Spec(B) \to \Spec(A) = U$, alors
$R\Gamma(V, Lg^*E) = R\Gamma(U, E) \otimes_A^\mathbf{L} B$.
Voir Catégories dérivées de schémas, lemme
\ref{perfect-lemma-quasi-coherence-pullback}.

\begin{lemma}
\label{lemma-qcoh-relative-pseudo-coherence-characterize}
Soit $f : X \to S$ un morphisme de schémas localement de type fini.
Soit $E$ un objet de $D_\QCoh(\mathcal{O}_X)$.
Fixons $m \in \mathbf{Z}$. Les assertions suivantes sont équivalentes :
\begin{enumerate}
\item $E$ est $m$-pseudo-cohérent relativement à $S$ ;
\item il existe un recouvrement ouvert affine $S = \bigcup V_i$ et,
pour tout $i$, un recouvrement ouvert affine $f^{-1}(V_i) = \bigcup U_{ij}$
tels que le complexe de $\mathcal{O}_X(U_{ij})$-modules
$R\Gamma(U_{ij}, E)$ soit $m$-pseudo-cohérent relativement à
$\mathcal{O}_S(V_i)$ ;
\item pour tous ouverts affines $U \subset X$ et $V \subset S$
tels que $f(U) \subset V$, le complexe de $\mathcal{O}_X(U)$-modules
$R\Gamma(U, E)$ est $m$-pseudo-cohérent relativement à $\mathcal{O}_S(V)$.
\end{enumerate}
\end{lemma}

\begin{proof}
Soient $U$ et $V$ comme dans (2), et choisissons une immersion fermée
$i : U \to \mathbf{A}^n_V$. Un argument formel, utilisant le
lemme \ref{lemma-relative-pseudo-coherence-characterize}, montre qu'il
suffit de prouver que
$Ri_*(E|_U)$ est $m$-pseudo-cohérent si et seulement si $R\Gamma(U, E)$
est $m$-pseudo-cohérent relativement à $\mathcal{O}_S(V)$.
Écrivons $U = \Spec(A)$, $V = \Spec(R)$ et
$\mathbf{A}^n_V = \Spec(R[x_1, \ldots, x_n]$. D'après les remarques
précédant le lemme, $E|_U$ est quasi-isomorphe au
complexe de faisceaux quasi-cohérents sur $U$ associé à l'objet
$R\Gamma(U, E)$ de $D(A)$. Remarquons que
$R\Gamma(U, E) = R\Gamma(\mathbf{A}^n_V, Ri_*(E|_U))$, car $i$ est une
immersion fermée (et donc $i_*$ est exact). Ainsi, $Ri_*E$
est associé à $R\Gamma(U, E)$ considéré comme un objet de
$D(R[x_1, \ldots, x_n])$. On conclut, puisque la $m$-pseudo-cohérence
de $Ri_*(E|_U)$ équivaut à la $m$-pseudo-cohérence de
$R\Gamma(U, E)$ dans $D(R[x_1, \ldots, x_n])$ d'après
Catégories dérivées de schémas, lemme \ref{perfect-lemma-pseudo-coherent-affine},
ce qui équivaut à dire que $R\Gamma(U, E)$ est $m$-pseudo-cohérent
relativement à $R = \mathcal{O}_S(V)$ par définition.
\end{proof}

\begin{lemma}
\label{lemma-closed-morphism-relative-pseudo-coherence}
Soit $i : X \to Y$ un morphisme de schémas localement de type fini sur un
schéma de base $S$. Supposons que $i$ induise un homéomorphisme de $X$ sur une partie
fermée de $Y$. Soit $E$ un objet de $D(\mathcal{O}_X)$.
Alors $E$ est $m$-pseudo-cohérent relativement à $S$ si et seulement si
$Ri_*E$ est $m$-pseudo-cohérent relativement à $S$.
\end{lemma}

\begin{proof}
D'après Morphismes, lemme \ref{morphisms-lemma-homeomorphism-affine},
le morphisme $i$ est affine. On peut donc supposer $S$, $Y$ et $X$ affines.
Écrivons $S = \Spec(R)$, $Y = \Spec(A)$ et $X = \Spec(B)$.
La condition signifie que $A/\text{rad}(A) \to B/\text{rad}(B)$ est
surjectif ; ici $\text{rad}(A)$ et $\text{rad}(B)$ désignent
les radicaux de Jacobson de $A$ et de $B$.
Comme $B$ est de type fini sur $A$, on peut trouver
$b_1, \ldots, b_m \in \text{rad}(B)$ qui engendrent $B$ comme
$A$-algèbre. Écrivons $b_j^N = 0$ pour tout $j$. Considérons le diagramme
d'anneaux
$$
\xymatrix{
B & R[x_i, y_j]/(y_j^N) \ar[l] & R[x_i, y_j] \ar[l] \\
A \ar[u] & R[x_i] \ar[l] \ar[u] \ar[ru]
}
$$
qui se traduit par un diagramme
$$
\xymatrix{
X \ar[d] \ar[r] & T \ar[d] \ar[r] & \mathbf{A}^{n + m}_S \ar[ld] \\
Y \ar[r] & \mathbf{A}^n_S
}
$$
de schémas affines. D'après le lemme \ref{lemma-relative-pseudo-coherence-characterize},
on voit que $E$ est $m$-pseudo-cohérent relativement à $S$ si et seulement si son
image directe dans $\mathbf{A}^{n + m}_S$ est $m$-pseudo-cohérente.
D'après Catégories dérivées de schémas, lemme
\ref{perfect-lemma-closed-push-pseudo-coherent},
on voit que cela a lieu si et seulement si son image directe dans $T$ est
$m$-pseudo-cohérente. Le même lemme montre que cela a lieu si et seulement si
l'image directe dans $\mathbf{A}^n_S$ est $m$-pseudo-cohérente.
De nouveau, d'après le
lemme \ref{lemma-relative-pseudo-coherence-characterize},
cela a lieu si et seulement si $Ri_*E$ est $m$-pseudo-cohérent relativement à $S$.
\end{proof}

\begin{lemma}
\label{lemma-finite-morphism-relative-pseudo-coherence}
Soit $\pi : X \to Y$ un morphisme fini de schémas localement de type
fini sur un schéma de base $S$. Soit $E$ un objet de
$D_\QCoh(\mathcal{O}_X)$. Alors $E$ est $m$-pseudo-cohérent
relativement à $S$ si et seulement si $R\pi_*E$ est $m$-pseudo-cohérent
relativement à $S$.
\end{lemma}

\begin{proof}
C'est la traduction du résultat de
Plus sur l'algèbre, lemme
\ref{more-algebra-lemma-finite-extension-pseudo-coherent}
dans le langage des schémas. Observons que $R\pi_*$ envoie bien
$D_\QCoh(\mathcal{O}_X)$ dans $D_\QCoh(\mathcal{O}_Y)$
d'après Catégories dérivées de schémas, lemme
\ref{perfect-lemma-quasi-coherence-direct-image}.
Pour effectuer la traduction, utiliser le
lemme \ref{lemma-relative-pseudo-coherence-characterize}.
\end{proof}

\begin{lemma}
\label{lemma-cone-relatively-pseudo-coherent}
Soit $f : X \to S$ un morphisme de schémas localement de type fini.
Soit $(E, E', E'')$ un triangle distingué de
$D(\mathcal{O}_X)$. Soit $m \in \mathbf{Z}$.
\begin{enumerate}
\item Si $E$ est $(m + 1)$-pseudo-cohérent relativement à $S$ et
$E'$ est $m$-pseudo-cohérent relativement à $S$, alors $E''$ est
$m$-pseudo-cohérent relativement à $S$.
\item Si $E, E''$ sont $m$-pseudo-cohérents relativement à $S$,
alors $E'$ est $m$-pseudo-cohérent relativement à $S$.
\item Si $E'$ est $(m + 1)$-pseudo-cohérent relativement à $S$
et $E''$ est $m$-pseudo-cohérent relativement à $S$, alors
$E$ est $(m + 1)$-pseudo-cohérent relativement à $S$.
\end{enumerate}
De plus, si deux des trois objets $E, E', E''$ sont pseudo-cohérents
relativement à $S$, alors le troisième l'est aussi.
\end{lemma}

\begin{proof}
Cela résulte immédiatement du lemme \ref{lemma-relative-pseudo-coherence-characterize} et de
Cohomologie, lemme \ref{cohomology-lemma-cone-pseudo-coherent}.
\end{proof}

\begin{lemma}
\label{lemma-rel-n-pseudo-module}
Soit $X \to S$ un morphisme de schémas localement de type fini.
Soit $\mathcal{F}$ un $\mathcal{O}_X$-module. Alors
\begin{enumerate}
\item $\mathcal{F}$ est $m$-pseudo-cohérent relativement à $S$ pour tout $m > 0$ ;
\item $\mathcal{F}$ est $0$-pseudo-cohérent relativement à $S$ si et seulement si
$\mathcal{F}$ est un $\mathcal{O}_X$-module de type fini ;
\item $\mathcal{F}$ est $(-1)$-pseudo-cohérent relativement à $S$ si et seulement si
$\mathcal{F}$ est quasi-cohérent et de présentation finie relativement à $S$.
\end{enumerate}
\end{lemma}

\begin{proof}
L'assertion (1) résulte immédiatement de la définition. Pour voir l'assertion (3),
on peut travailler localement sur $X$ (les deux propriétés sont locales). Ainsi, on
peut supposer $X$ et $S$ affines. Choisissons une immersion fermée
$i : X \to \mathbf{A}^n_S$. On voit alors que $\mathcal{F}$ est
$(-1)$-pseudo-cohérent relativement à $S$ si et seulement si $i_*\mathcal{F}$
est $(-1)$-pseudo-cohérent, ce qui a lieu si et seulement si $i_*\mathcal{F}$
est un $\mathcal{O}_{\mathbf{A}^n_S}$-module de présentation finie, voir
Cohomologie, lemme \ref{cohomology-lemma-finite-cohomology}.
Un module de présentation finie est quasi-cohérent, voir
Modules, lemme \ref{modules-lemma-finite-presentation-quasi-coherent}.
D'après Morphismes, lemme \ref{morphisms-lemma-i-star-equivalence},
on voit que $\mathcal{F}$ est quasi-cohérent si et seulement si $i_*\mathcal{F}$
est quasi-cohérent. Cela étant dit, l'assertion (3) en découle. La démonstration de (2)
est analogue, mais plus simple.
\end{proof}

\begin{lemma}
\label{lemma-summands-relative-pseudo-coherent}
Soit $X \to S$ un morphisme de schémas localement de type fini.
Soit $m \in \mathbf{Z}$. Soient $E, K$ des objets de $D(\mathcal{O}_X)$.
Si $E \oplus K$ est $m$-pseudo-cohérent relativement à $S$, alors $E$ et $K$ le sont aussi.
\end{lemma}

\begin{proof}
Cela résulte de
Cohomologie, lemme \ref{cohomology-lemma-summands-pseudo-coherent},
et des définitions.
\end{proof}

\begin{lemma}
\label{lemma-complex-relative-pseudo-coherent-modules}
Soit $X \to S$ un morphisme de schémas localement de type fini.
Soit $m \in \mathbf{Z}$. Soit $\mathcal{F}^\bullet$ un complexe (localement) borné
supérieurement de $\mathcal{O}_X$-modules tel que $\mathcal{F}^i$ soit
$(m - i)$-pseudo-cohérent relativement à $S$ pour tout $i$. Alors
$\mathcal{F}^\bullet$ est $m$-pseudo-cohérent relativement à $S$.
\end{lemma}

\begin{proof}
Cela résulte de
Cohomologie, lemme \ref{cohomology-lemma-complex-pseudo-coherent-modules},
et des définitions.
\end{proof}

\begin{lemma}
\label{lemma-cohomology-relative-pseudo-coherent}
Soit $X \to S$ un morphisme de schémas localement de type fini.
Soit $m \in \mathbf{Z}$. Soit $E$ un objet de $D(\mathcal{O}_X)$.
Si $E$ est (localement) borné supérieurement et si $H^i(E)$ est $(m - i)$-pseudo-cohérent
relativement à $S$ pour tout $i$, alors $E$ est $m$-pseudo-cohérent relativement à $S$.
\end{lemma}

\begin{proof}
Cela résulte de
Cohomologie, lemme \ref{cohomology-lemma-cohomology-pseudo-coherent},
et des définitions.
\end{proof}

\begin{lemma}
\label{lemma-base-change-relative-pseudo-coherent}
Soit $X \to S$ un morphisme de schémas localement de type fini.
Soit $m \in \mathbf{Z}$. Soit $E$ un objet de $D(\mathcal{O}_X)$
qui est $m$-pseudo-cohérent relativement à $S$. Soit $S' \to S$ un
morphisme de schémas. Posons $X' = X \times_S S'$ et notons $E'$
l'image inverse dérivée de $E$ sur $X'$. Si $S'$ et $X$ sont
Tor-indépendants sur $S$, alors $E'$
est $m$-pseudo-cohérent relativement à $S'$.
\end{lemma}

\begin{proof}
Le problème est local sur $X$ et $X'$ ; on peut donc supposer $X$, $S$, $S'$ et
$X'$ affines. Choisissons une immersion fermée $i : X \to \mathbf{A}^n_S$
et notons $i' : X' \to \mathbf{A}^n_{S'}$ son changement de base à $S'$.
Notons $g : X' \to X$ et $g' : \mathbf{A}^n_{S'} \to \mathbf{A}^n_S$
les projections, de sorte que $E' = Lg^*E$. Puisque $X$ et $S'$ sont tor-indépendants
sur $S$, le morphisme de changement de base
(Cohomologie, remarque \ref{cohomology-remark-base-change})
induit un isomorphisme
$$
Ri'_*(Lg^*E) = L(g')^*Ri_*E
$$
En effet, pour un point $x' \in X'$ au-dessus de $x \in X$, le morphisme de changement de base
sur les fibres en $x'$ est le morphisme
$$
E_x \otimes_{\mathcal{O}_{\mathbf{A}^n_S, x}}^\mathbf{L}
\mathcal{O}_{\mathbf{A}^n_{S'}, x'}
\longrightarrow
E_x \otimes_{\mathcal{O}_{X, x}}^\mathbf{L}
\mathcal{O}_{X', x'}
$$
provenant des immersions fermées $i$ et $i'$. Remarquons que la source
est quasi-isomorphe à une localisation de
$E_x \otimes_{\mathcal{O}_{S, s}}^\mathbf{L} \mathcal{O}_{S', s'}$,
qui est isomorphe au but, puisque
$\mathcal{O}_{X', x'}$ est isomorphe à la (même) localisation de
$\mathcal{O}_{X, x} \otimes_{\mathcal{O}_{S, s}}^\mathbf{L}
\mathcal{O}_{S', s'}$ par hypothèse. On conclut le lemme
en appliquant
Cohomologie, lemme \ref{cohomology-lemma-pseudo-coherent-pullback}.
\end{proof}

\begin{lemma}
\label{lemma-pull-relative-pseudo-coherent}
Soit $f : X \to Y$ un morphisme de schémas localement de type fini
sur une base $S$. Soit $m \in \mathbf{Z}$. Soit $E$ un objet de
$D(\mathcal{O}_Y)$. Supposons que
\begin{enumerate}
\item $\mathcal{O}_X$ soit pseudo-cohérent relativement à $Y$\footnote{Cela
signifie que $f$ est pseudo-cohérent, voir la
définition \ref{definition-pseudo-coherent}.} ;
\item $E$ soit $m$-pseudo-cohérent relativement à $S$.
\end{enumerate}
Alors $Lf^*E$ est $m$-pseudo-cohérent relativement à $S$.
\end{lemma}

\begin{proof}
Le problème est local sur $X$. On peut donc supposer $X$, $Y$ et $S$ affines.
En raisonnant comme dans la démonstration de Plus sur l'algèbre, lemme
\ref{more-algebra-lemma-pull-relative-pseudo-coherent},
on peut trouver un diagramme commutatif
$$
\xymatrix{
X \ar[r]_i \ar[d]_f &
\mathbf{A}^d_Y \ar[r]_j \ar[ld]^p &
\mathbf{A}^{n + d}_S \ar[ld] \\
Y \ar[r] &
\mathbf{A}^n_S
}
$$
Observons que
$$
Ri_* Lf^*E = Ri_* Li^* Lp^*E =
Lp^*E \otimes_{\mathcal{O}_{\mathbf{A}_Y^n}}^\mathbf{L} Ri_*\mathcal{O}_X
$$
d'après Cohomologie, lemme
\ref{cohomology-lemma-projection-formula-closed-immersion}.
Par hypothèse et puisque $Y$ est affine, on peut représenter
$Ri_*\mathcal{O}_X = i_*\mathcal{O}_X$ par un complexe de modules libres de type fini
sur $\mathcal{O}_{\mathbf{A}_Y^n}$, noté $\mathcal{F}^\bullet$, avec
$\mathcal{F}^q = 0$ pour $q > 0$
(détails omis ; utiliser Catégories dérivées de schémas, lemme
\ref{perfect-lemma-pseudo-coherent-affine}
et
Plus sur l'algèbre, lemme
\ref{more-algebra-lemma-rel-n-pseudo-module}).
Par hypothèse, $E$ est borné supérieurement ; écrivons $H^q(E) = 0$ pour $q > a$.
Représentons $E$ par un complexe $\mathcal{E}^\bullet$ de $\mathcal{O}_Y$-modules
tel que $\mathcal{E}^q = 0$ pour $q > a$. Alors le produit tensoriel dérivé ci-dessus
est représenté par $\text{Tot}(p^*\mathcal{E}^\bullet
\otimes_{\mathcal{O}_{\mathbf{A}_Y^n}} \mathcal{F}^\bullet)$.

\medskip\noindent
Puisque $j$ est une immersion fermée, le foncteur $j_*$ est exact et
$Rj_*$ se calcule en appliquant $j_*$ à n'importe quel complexe représentatif
de faisceaux. Il faut donc montrer que
$j_*\text{Tot}(p^*\mathcal{E}^\bullet \otimes_{\mathcal{O}_{\mathbf{A}_Y^n}}
\mathcal{F}^\bullet)$ est $m$-pseudo-cohérent
comme complexe de $\mathcal{O}_{\mathbf{A}^{n + m}_S}$-modules.
Remarquons que
$\text{Tot}(p^*\mathcal{E}^\bullet \otimes_{\mathcal{O}_{\mathbf{A}_Y^n}}
\mathcal{F}^\bullet)$ possède une filtration par
sous-complexes dont les quotients successifs sont les complexes
$p^*\mathcal{E}^\bullet
\otimes_{\mathcal{O}_{\mathbf{A}_Y^n}} \mathcal{F}^q[-q]$.
Remarquons que, pour $q \ll 0$, les complexes
$p^*\mathcal{E}^\bullet \otimes_{\mathcal{O}_{\mathbf{A}_Y^n}}
\mathcal{F}^q[-q]$
ont une cohomologie nulle en degrés $\leq m$ et sont donc $m$-pseudo-cohérents.
Ainsi, en appliquant le
lemme \ref{lemma-cone-relatively-pseudo-coherent}
et en raisonnant par récurrence, il suffit de montrer que
$p^*\mathcal{E}^\bullet \otimes_{\mathcal{O}_{\mathbf{A}_Y^n}}
\mathcal{F}^q[-q]$
est pseudo-cohérent relativement à $S$ pour tout $q$. Remarquons que $\mathcal{F}^q = 0$
pour $q > 0$. Puisque $\mathcal{F}^q$ est aussi libre de type fini, cela
se ramène à prouver que $p^*\mathcal{E}^\bullet$ est
$m$-pseudo-cohérent relativement à $S$, ce qui résulte du
lemme \ref{lemma-base-change-relative-pseudo-coherent}
par exemple.
\end{proof}

\begin{lemma}
\label{lemma-composition-relative-pseudo-coherent}
Soit $f : X \to Y$ un morphisme de schémas localement de type fini
sur une base $S$. Soit $m \in \mathbf{Z}$. Soit $E$ un objet de
$D(\mathcal{O}_X)$. Supposons que $\mathcal{O}_Y$ soit pseudo-cohérent relativement
à $S$\footnote{Cela signifie que $Y \to S$ est pseudo-cohérent, voir la
définition \ref{definition-pseudo-coherent}.}.
Alors les assertions suivantes sont équivalentes :
\begin{enumerate}
\item $E$ est $m$-pseudo-cohérent relativement à $Y$ ;
\item $E$ est $m$-pseudo-cohérent relativement à $S$.
\end{enumerate}
\end{lemma}

\begin{proof}
La question est locale sur $X$ ; on peut donc supposer $X$, $Y$ et $S$ affines.
En raisonnant comme dans la démonstration de Plus sur l'algèbre, lemme
\ref{more-algebra-lemma-pull-relative-pseudo-coherent},
on peut trouver un diagramme commutatif
$$
\xymatrix{
X \ar[r]_i \ar[d]_f &
\mathbf{A}^m_Y \ar[r]_j \ar[ld]^p &
\mathbf{A}^{n + m}_S \ar[ld] \\
Y \ar[r] &
\mathbf{A}^n_S
}
$$
L'hypothèse que $\mathcal{O}_Y$ est pseudo-cohérent relativement à $S$
entraîne que $\mathcal{O}_{\mathbf{A}^m_Y}$ est pseudo-cohérent relativement
à $\mathbf{A}^m_S$ (par changement de base plat ; on peut le voir en utilisant,
par exemple, le lemme \ref{lemma-base-change-relative-pseudo-coherent}).
Cela entraîne à son tour que $j_*\mathcal{O}_{\mathbf{A}^n_Y}$ est
pseudo-cohérent comme
$\mathcal{O}_{\mathbf{A}^{n + m}_S}$-module. L'équivalence du
lemme résulte alors de
Catégories dérivées de schémas, lemme
\ref{perfect-lemma-closed-push-pseudo-coherent}.
\end{proof}

\begin{lemma}
\label{lemma-check-relative-pseudo-coherence-on-charts}
Soit
$$
\xymatrix{
X \ar[rd] \ar[rr]_i & & P \ar[ld] \\
& S
}
$$
un diagramme commutatif de schémas. Supposons que $i$ soit une immersion fermée
et que $P \to S$ soit plat et localement de présentation finie. Soit $E$
un objet de $D(\mathcal{O}_X)$. Alors les assertions suivantes
sont équivalentes :
\begin{enumerate}
\item $E$ est $m$-pseudo-cohérent relativement à $S$ ;
\item $Ri_*E$ est $m$-pseudo-cohérent relativement à $S$ ;
\item $Ri_*E$ est $m$-pseudo-cohérent sur $P$.
\end{enumerate}
\end{lemma}

\begin{proof}
L'équivalence de (1) et (2) est le
lemme \ref{lemma-finite-morphism-relative-pseudo-coherence}.
L'équivalence de (2) et (3) résulte du
lemme
\ref{lemma-composition-relative-pseudo-coherent}
appliqué à $\text{id} : P \to P$,
pourvu que l'on montre que $\mathcal{O}_P$ est
pseudo-cohérent relativement à $S$. Cela résulte de
Plus sur l'algèbre, lemme
\ref{more-algebra-lemma-flat-finite-presentation-perfect},
et des définitions.
\end{proof}









\section{Morphismes pseudo-cohérents}
\label{section-pseudo-coherent}

\noindent
Évitez à tout prix de lire cette section.
Si vous avez besoin d'une partie de ce matériau, consultez d'abord les
sections algébriques correspondantes, voir
Plus sur l'algèbre, sections \ref{more-algebra-section-pseudo-coherent},
\ref{more-algebra-section-relative-pseudo-coherent} et
\ref{more-algebra-section-pseudo-coherent-perfect-ring-map}.
Pour l'instant, la seule chose à savoir est qu'un homomorphisme d'anneaux
$A \to B$ est pseudo-cohérent si et seulement si $B = A[x_1, \ldots, x_n]/I$
et si $B$, comme $A[x_1, \ldots, x_n]$-module, possède une résolution par des
$A[x_1, \ldots, x_n]$-modules libres de type fini.

\begin{lemma}
\label{lemma-pseudo-coherent}
Soit $f : X \to S$ un morphisme de schémas. Les assertions suivantes sont équivalentes :
\begin{enumerate}
\item il existe un recouvrement ouvert affine $S = \bigcup V_j$ et, pour chaque $j$,
un recouvrement ouvert affine $f^{-1}(V_j) = \bigcup U_{ji}$ tel que
$\mathcal{O}_S(V_j) \to \mathcal{O}_X(U_{ij})$ soit un homomorphisme
d'anneaux pseudo-cohérent ;
\item pour toute paire d'ouverts affines $U \subset X$, $V \subset S$
telle que $f(U) \subset V$, l'homomorphisme d'anneaux
$\mathcal{O}_S(V) \to \mathcal{O}_X(U)$ est pseudo-cohérent ;
\item $f$ est localement de type fini et $\mathcal{O}_X$
est pseudo-cohérent relativement à $S$.
\end{enumerate}
\end{lemma}

\begin{proof}
Pour voir l'équivalence de (1) et (2), il suffit de vérifier les conditions
(1)(a), (b), (c) de
Morphismes, définition \ref{morphisms-definition-property-local},
pour la propriété d'être un homomorphisme d'anneaux pseudo-cohérent.
Ces propriétés résultent (en utilisant que la localisation est plate) de
Plus sur l'algèbre, lemmes
\ref{more-algebra-lemma-base-change-relative-pseudo-coherent},
\ref{more-algebra-lemma-localize-relative-pseudo-coherent} et
\ref{more-algebra-lemma-glue-relative-pseudo-coherent}.

\medskip\noindent
Si (1) est vérifiée, alors $f$ est localement de type fini, car un homomorphisme
d'anneaux pseudo-cohérent est de type fini par définition. De plus, (1) entraîne,
via le lemme \ref{lemma-qcoh-relative-pseudo-coherence-characterize}
et les définitions, que $\mathcal{O}_X$ est pseudo-cohérent
relativement à $S$. Réciproquement, si (3) est vérifiée, on voit
que, pour tous $U$ et $V$ comme dans (2), l'anneau
$\mathcal{O}_X(U)$ est de type fini sur $\mathcal{O}_S(V)$
et que $\mathcal{O}_X(U)$ est, comme module,
pseudo-cohérent relativement à $\mathcal{O}_S(V)$, voir les
lemmes \ref{lemma-relative-pseudo-coherence-characterize} et
\ref{lemma-qcoh-relative-pseudo-coherence-characterize}.
C'est la définition d'un homomorphisme d'anneaux pseudo-cohérent ; ainsi,
(2) et (1) sont vérifiées.
\end{proof}

\begin{definition}
\label{definition-pseudo-coherent}
Un morphisme de schémas $f : X \to S$ est dit {\it pseudo-cohérent}
si les conditions équivalentes du
lemme \ref{lemma-pseudo-coherent}
sont satisfaites. Dans ce cas, on dit aussi que $X$ est pseudo-cohérent
sur $S$.
\end{definition}

\noindent
Attention : un changement de base d'un morphisme pseudo-cohérent n'est pas
pseudo-cohérent en général.

\begin{lemma}
\label{lemma-flat-base-change-pseudo-coherent}
Un changement de base plat d'un morphisme pseudo-cohérent est pseudo-cohérent.
\end{lemma}

\begin{proof}
Cela se traduit par le résultat algébrique suivant :
soit $A \to B$ un homomorphisme d'anneaux pseudo-cohérent.
Soit $A \to A'$ plat. Alors $A' \to B \otimes_A A'$ est
pseudo-cohérent. Cela résulte du résultat plus général
Plus sur l'algèbre,
lemme \ref{more-algebra-lemma-base-change-relative-pseudo-coherent}.
\end{proof}

\begin{lemma}
\label{lemma-composition-pseudo-coherent}
Un composé de morphismes de schémas pseudo-cohérents est
pseudo-cohérent.
\end{lemma}

\begin{proof}
Cela se traduit par le résultat algébrique suivant :
si $A \to B \to C$ sont des homomorphismes d'anneaux pseudo-cohérents composables,
alors $A \to C$ est pseudo-cohérent. Cela résulte soit de
Plus sur l'algèbre,
lemme \ref{more-algebra-lemma-pull-relative-pseudo-coherent},
soit de
Plus sur l'algèbre,
lemme \ref{more-algebra-lemma-composition-relative-pseudo-coherent}.
\end{proof}

\begin{lemma}
\label{lemma-pseudo-coherent-finite-presentation}
Un morphisme pseudo-cohérent est localement de présentation finie.
\end{lemma}

\begin{proof}
Cela résulte immédiatement des définitions.
\end{proof}

\begin{lemma}
\label{lemma-flat-finite-presentation-pseudo-coherent}
Un morphisme plat et localement de présentation finie est pseudo-cohérent.
\end{lemma}

\begin{proof}
Cela résulte du fait qu'un homomorphisme d'anneaux plat de présentation finie est
pseudo-cohérent (et même parfait), voir
Plus sur l'algèbre,
lemme \ref{more-algebra-lemma-flat-finite-presentation-perfect}.
\end{proof}

\begin{lemma}
\label{lemma-permanence-pseudo-coherent}
Soit $f : X \to Y$ un morphisme de schémas pseudo-cohérents
sur un schéma de base $S$. Alors $f$ est pseudo-cohérent.
\end{lemma}

\begin{proof}
Cela se traduit par le résultat algébrique suivant :
si $R \to A \to B$ sont des homomorphismes d'anneaux composables
et si $R \to A$, $R \to B$ sont pseudo-cohérents, alors
$R \to B$ est pseudo-cohérent. Cela résulte de
Plus sur l'algèbre,
lemme \ref{more-algebra-lemma-composition-relative-pseudo-coherent}.
\end{proof}

\begin{lemma}
\label{lemma-finite-pseudo-coherent}
Soit $f : X \to S$ un morphisme fini de schémas.
Alors $f$ est pseudo-cohérent si et seulement si $f_*\mathcal{O}_X$
est pseudo-cohérent comme $\mathcal{O}_S$-module.
\end{lemma}

\begin{proof}
Traduit en algèbre, ce lemme affirme ceci : si $R \to A$
est un homomorphisme d'anneaux fini, alors $R \to A$ est pseudo-cohérent comme homomorphisme d'anneaux
(ce qui signifie par définition que $A$, comme $A$-module, est
pseudo-cohérent relativement à $R$) si et seulement si
$A$ est pseudo-cohérent comme $R$-module. Cela résulte du résultat
plus général de
Plus sur l'algèbre, lemme
\ref{more-algebra-lemma-finite-extension-pseudo-coherent}.
\end{proof}

\begin{lemma}
\label{lemma-Noetherian-pseudo-coherent}
Soit $f : X \to S$ un morphisme de schémas.
Si $S$ est localement noethérien, alors $f$ est pseudo-cohérent si
et seulement si $f$ est localement de type fini.
\end{lemma}

\begin{proof}
Cela se traduit par le résultat algébrique suivant :
si $R \to A$ est un homomorphisme d'anneaux de type fini avec $R$ noethérien, alors
$R \to A$ est pseudo-cohérent si et seulement si $R \to A$ est de type fini.
Pour le voir, remarquons qu'un homomorphisme d'anneaux pseudo-cohérent est de type fini par
définition. Réciproquement, si $R \to A$ est de type fini, alors
on peut écrire $A = R[x_1, \ldots, x_n]/I$ et il résulte de
Plus sur l'algèbre,
lemme \ref{more-algebra-lemma-Noetherian-pseudo-coherent},
que $A$ est pseudo-cohérent comme $R[x_1, \ldots, x_n]$-module, c'est-à-dire que
$R \to A$ est un homomorphisme d'anneaux pseudo-cohérent.
\end{proof}

\begin{lemma}
\label{lemma-descending-property-pseudo-coherent}
La propriété $\mathcal{P}(f) =$``$f$ est pseudo-cohérent''
est fpqc locale sur la base.
\end{lemma}

\begin{proof}
Nous allons utiliser le critère de
Descente, lemme \ref{descent-lemma-descending-properties-morphisms}
pour le démontrer. Par
la définition \ref{definition-pseudo-coherent},
la propriété d'être pseudo-cohérent est locale pour la topologie de Zariski sur la base. Par
le lemme \ref{lemma-flat-base-change-pseudo-coherent},
la propriété d'être pseudo-cohérent est préservée par changement de base plat.
La dernière hypothèse (3) du
lemme \ref{descent-lemma-descending-properties-morphisms} de Descente
se traduit par l'énoncé algébrique suivant :
Soit $A \to B$ un homomorphisme fidèlement plat d'anneaux.
Soit $C = A[x_1, \ldots, x_n]/I$ une $A$-algèbre.
Si $C \otimes_A B$ est pseudo-cohérent comme $B[x_1, \ldots, x_n]$-module,
alors $C$ est pseudo-cohérent comme $A[x_1, \ldots, x_n]$-module.
C'est le
lemme \ref{more-algebra-lemma-flat-descent-pseudo-coherent} de Compléments d'algèbre.
\end{proof}

\begin{lemma}
\label{lemma-quotient-of-flat-finitely-presented}
Soit $A \to B$ un homomorphisme plat d'anneaux de présentation finie.
Soit $I \subset B$ un idéal. Alors $A \to B/I$ est pseudo-cohérent
si et seulement si $I$ est pseudo-cohérent comme $B$-module.
\end{lemma}

\begin{proof}
Choisissons une présentation $B = A[x_1, \ldots, x_n]/J$.
Remarquons que $B$ est pseudo-cohérent comme $A[x_1, \ldots, x_n]$-module
car $A \to B$ est un homomorphisme pseudo-cohérent d'anneaux d'après le
lemme \ref{lemma-flat-finite-presentation-pseudo-coherent}.
Remarquons que $A \to B/I$ est pseudo-cohérent si et seulement si
$B/I$ est pseudo-cohérent comme $A[x_1, \ldots, x_n]$-module. D'après le
lemme \ref{more-algebra-lemma-finite-push-pseudo-coherent} de Compléments d'algèbre,
nous voyons que cela équivaut à demander que $B/I$ soit
pseudo-cohérent comme $B$-module. Cela démontre le lemme, car la
suite exacte courte $0 \to I \to B \to B/I \to 0$
montre que $I$ est pseudo-cohérent si et seulement si $B/I$ l'est (voir
Compléments d'algèbre,
lemme \ref{more-algebra-lemma-two-out-of-three-pseudo-coherent}).
\end{proof}

\noindent
Le lemme suivant sera supplanté par le
lemme \ref{lemma-pseudo-coherent-fppf-local-source}, plus fort.

\begin{lemma}
\label{lemma-pseudo-coherent-syntomic-local-source}
La propriété $\mathcal{P}(f) =$``$f$ est pseudo-cohérent''
est locale pour la topologie syntomique sur la source.
\end{lemma}

\begin{proof}
Nous allons utiliser le critère du
lemme \ref{descent-lemma-properties-morphisms-local-source} de Descente
pour le démontrer. Il résulte des
lemmes \ref{lemma-flat-finite-presentation-pseudo-coherent} et
\ref{lemma-composition-pseudo-coherent}
que la propriété d'être pseudo-cohérent est préservée par précomposition avec des
morphismes plats localement de présentation finie, en particulier par
précomposition avec des morphismes syntomiques (voir
Morphismes, lemmes \ref{morphisms-lemma-syntomic-flat} et
\ref{morphisms-lemma-syntomic-locally-finite-presentation}).
Il ressort clairement de la
définition \ref{definition-pseudo-coherent}
que la propriété d'être pseudo-cohérent est
locale pour la topologie de Zariski sur la source et le but.
Ainsi, d'après le
lemme \ref{descent-lemma-properties-morphisms-local-source} de Descente déjà cité,
il suffit de démontrer ce qui suit : supposons que $X' \to X \to Y$ soient des
morphismes de schémas affines, avec $X' \to X$ syntomique et $X' \to Y$
pseudo-cohérent. Alors $X \to Y$ est pseudo-cohérent.
Pour le voir, remarquons que, dans tous les cas, $X \to Y$ est de présentation finie d'après
Descente, lemme \ref{descent-lemma-flat-finitely-presented-permanence-algebra}.
Choisissons une immersion fermée $X \to \mathbf{A}^n_Y$. D'après
Algèbre, lemme \ref{algebra-lemma-lift-syntomic},
nous pouvons trouver un recouvrement d'ouverts affines $X' = \bigcup_{i = 1, \ldots, n} X'_i$
et des morphismes syntomiques $W_i \to \mathbf{A}^n_Y$ relevant les morphismes
$X'_i \to X$, c'est-à-dire tels qu'il existe des diagrammes cartésiens
$$
\xymatrix{
X'_i \ar[d] \ar[r] & W_i \ar[d] \\
X \ar[r] & \mathbf{A}^n_Y
}
$$
Après avoir remplacé $X'$ par $\coprod X'_i$ et posé $W = \coprod W_i$,
nous obtenons un diagramme cartésien
$$
\xymatrix{
X' \ar[d] \ar[r] & W \ar[d]^h \\
X \ar[r] & \mathbf{A}^n_Y
}
$$
où $W \to \mathbf{A}^n_Y$ est plat et de présentation finie, et
$X' \to Y$ est encore pseudo-cohérent. Puisque
$W \to \mathbf{A}^n_Y$ est ouvert (voir
Morphismes, lemme \ref{morphisms-lemma-fppf-open})
et que $X' \to X$ est surjectif, nous pouvons trouver
$f \in \Gamma(\mathbf{A}^n_Y, \mathcal{O})$ tel que
$X \subset D(f) \subset \Im(h)$. Écrivons
$Y = \Spec(R)$, $X = \Spec(A)$, $X' = \Spec(A')$
et $W = \Spec(B)$, $A = R[x_1, \ldots, x_n]/I$ et
$A' = B/IB$. Alors $R \to A'$ est pseudo-cohérent. Considérons le diagramme
$$
\xymatrix{
A' = B/IB & B \ar[l] \\
A = R[x_1, \ldots, x_n]/I \ar[u] & R[x_1, \ldots, x_n] \ar[l] \ar[u]
}
$$
D'après le
lemme \ref{lemma-quotient-of-flat-finitely-presented},
nous voyons que $IB$ est pseudo-cohérent comme $B$-module.
L'homomorphisme d'anneaux $R[x_1, \ldots, x_n]_f \to B_f$ est fidèlement plat par
notre choix de $f$ ci-dessus. Cela entraîne que
$I_f \subset R[x_1, \ldots, x_n]_f$
est pseudo-cohérent, voir
Compléments d'algèbre, lemme \ref{more-algebra-lemma-flat-descent-pseudo-coherent}.
En appliquant encore une fois le
lemme \ref{lemma-quotient-of-flat-finitely-presented},
nous voyons que $R \to A$ est pseudo-cohérent.
\end{proof}

\begin{lemma}
\label{lemma-pseudo-coherent-fppf-local-source}
La propriété $\mathcal{P}(f) =$``$f$ est pseudo-cohérent''
est locale pour la topologie fppf sur la source.
\end{lemma}

\begin{proof}
Soit $f : X \to S$ un morphisme de schémas.
Soit $\{g_i : X_i \to X\}$ un recouvrement fppf tel que chaque composée
$f \circ g_i$ soit pseudo-cohérente. D'après le
lemme \ref{lemma-dominate-fppf},
il existe
\begin{enumerate}
\item un recouvrement d'ouverts de Zariski $X = \bigcup U_j$,
\item des morphismes surjectifs finis localement libres $W_j \to U_j$,
\item des recouvrements d'ouverts de Zariski $W_j = \bigcup_k W_{j, k}$,
\item des morphismes surjectifs finis localement libres $T_{j, k} \to W_{j, k}$
\end{enumerate}
tels que le recouvrement fppf $\{h_{j, k} : T_{j, k} \to X\}$ raffine le
recouvrement donné $\{X_i \to X\}$. Notons $\psi_{j, k} : T_{j, k} \to X_{\alpha(j, k)}$
les morphismes qui témoignent du fait que $\{T_{j, k} \to X\}$ raffine
le recouvrement donné $\{X_i \to X\}$. Remarquons que $T_{j, k} \to X$ est un morphisme plat
localement de présentation finie ; ainsi, aussi bien $X_i$ que $T_{j, k}$ sont
pseudo-cohérents sur $X$ d'après le
lemme \ref{lemma-flat-finite-presentation-pseudo-coherent}.
Ainsi, $\psi_{j, k} : T_{j, k} \to X_i$ est pseudo-cohérent, voir le
lemme \ref{lemma-permanence-pseudo-coherent}.
Ainsi, $T_{j, k} \to S$ est pseudo-cohérent comme composée
de $\psi_{j, k}$ et de $f \circ g_{\alpha(j, k)}$, voir le
lemme \ref{lemma-composition-pseudo-coherent}.
Nous voyons donc que le lemme est ramené au cas d'un recouvrement
d'ouverts de Zariski (qui convient) et au cas d'un recouvrement
donné par un seul morphisme surjectif fini localement libre, que nous
traitons dans le paragraphe suivant.

\medskip\noindent
Supposons que $X' \to X \to S$ soit une suite de morphismes de schémas,
où $X' \to X$ est surjectif fini localement libre et $X' \to Y$ pseudo-cohérent.
Notre but est de montrer que $X \to S$ est pseudo-cohérent.
Remarquons que, d'après
Descente, lemme \ref{descent-lemma-flat-finitely-presented-permanence},
le morphisme $X \to S$ est localement de présentation finie.
Il est clair que le problème se ramène au cas où $X'$, $X$ et $S$
sont affines et où $X' \to X$ est libre de rang $r > 0$. Le problème
algébrique correspondant est le suivant : supposons que $R \to A \to A'$ soient des homomorphismes d'anneaux
tels que $R \to A'$ soit pseudo-cohérent, que $R \to A$ soit de présentation finie,
et que $A' \cong A^{\oplus r}$ comme $A$-module. But : montrer que $R \to A$ est
pseudo-cohérent. L'hypothèse que $R \to A'$ est pseudo-cohérent
signifie que $A'$, comme $A'$-module, est pseudo-cohérent relativement à $R$. D'après
Compléments d'algèbre,
lemme \ref{more-algebra-lemma-finite-extension-pseudo-coherent},
cela entraîne que $A'$, comme $A$-module, est pseudo-cohérent relativement à $R$.
Puisque $A' \cong A^{\oplus r}$ comme $A$-module, nous voyons que
$A$, comme $A$-module, est pseudo-cohérent relativement à $R$, voir
Compléments d'algèbre,
lemme \ref{more-algebra-lemma-summands-relative-pseudo-coherent}.
Cela signifie, par définition, que $R \to A$ est pseudo-cohérent, ce qui conclut.
\end{proof}








\section{Morphismes parfaits}
\label{section-perfect}

\noindent
Pour comprendre les résultats de cette
section, il suffit de connaître un peu ceux de la section
sur les morphismes pseudo-cohérents.
Pour le moment, la seule chose à savoir est qu'un homomorphisme d'anneaux
$A \to B$ est parfait si et seulement s'il est pseudo-cohérent
et si $B$ est de dimension de Tor finie comme $A$-module.

\begin{lemma}
\label{lemma-perfect}
Soit $f : X \to S$ un morphisme de schémas localement de type fini.
Les assertions suivantes sont équivalentes :
\begin{enumerate}
\item il existe un recouvrement d'ouverts affines $S = \bigcup V_j$ et, pour chaque $j$,
un recouvrement d'ouverts affines $f^{-1}(V_j) = \bigcup U_{ji}$ tel que
$\mathcal{O}_S(V_j) \to \mathcal{O}_X(U_{ij})$ soit un homomorphisme
d'anneaux parfait, et
\item pour toute paire d'ouverts affines $U \subset X$, $V \subset S$
tels que $f(U) \subset V$, l'homomorphisme d'anneaux
$\mathcal{O}_S(V) \to \mathcal{O}_X(U)$ est parfait.
\end{enumerate}
\end{lemma}

\begin{proof}
Supposons (1) et soient $U, V$ comme dans (2).
Il résulte du
lemme \ref{lemma-pseudo-coherent}
que $\mathcal{O}_S(V) \to \mathcal{O}_X(U)$ est pseudo-cohérent.
Il suffit donc de démontrer que la propriété, pour un homomorphisme d'anneaux,
d'être « de dimension de Tor finie » satisfait aux
conditions (1)(a), (b), (c) de
Morphismes, définition \ref{morphisms-definition-property-local}.
Ces propriétés résultent de
Compléments d'algèbre,
lemmes \ref{more-algebra-lemma-flat-push-tor-amplitude},
\ref{more-algebra-lemma-flat-base-change-finite-tor-dimension} et
\ref{more-algebra-lemma-glue-tor-amplitude}.
Certains détails sont omis.
\end{proof}

\begin{definition}
\label{definition-perfect}
Un morphisme de schémas $f : X \to S$ est dit {\it parfait}
si les conditions équivalentes du
lemme \ref{lemma-perfect}
sont satisfaites. Dans ce cas, nous disons aussi que $X$ est parfait
sur $S$.
\end{definition}

\noindent
Remarquons qu'un morphisme parfait est en particulier pseudo-cohérent, donc
localement de présentation finie. Attention : un changement de base d'un
morphisme parfait n'est pas parfait en général.

\begin{lemma}
\label{lemma-flat-base-change-perfect}
Un changement de base plat d'un morphisme parfait est parfait.
\end{lemma}

\begin{proof}
Cela se traduit par le résultat algébrique suivant :
soit $A \to B$ un homomorphisme parfait d'anneaux.
Soit $A \to A'$ plat. Alors $A' \to B \otimes_A A'$ est
parfait. Nous avons vu ce résultat pour les homomorphismes pseudo-cohérents d'anneaux dans le
lemme \ref{lemma-flat-base-change-pseudo-coherent}.
Le fait correspondant pour la dimension de Tor finie résulte de
Compléments d'algèbre,
lemme \ref{more-algebra-lemma-flat-base-change-finite-tor-dimension}.
\end{proof}

\begin{lemma}
\label{lemma-composition-perfect}
Une composée de morphismes parfaits de schémas est parfaite.
\end{lemma}

\begin{proof}
Cela se traduit par le résultat algébrique suivant :
si $A \to B \to C$ sont des homomorphismes parfaits d'anneaux composables,
alors $A \to C$ est parfait. Nous avons vu que c'est le cas pour la
pseudo-cohérence dans le
lemme \ref{lemma-composition-pseudo-coherent}
et sa démonstration. Par hypothèse, il existe des entiers $n$, $m$ tels
que $B$ soit de dimension de Tor $\leq n$ sur $A$ et $C$ de dimension de Tor
$\leq m$ sur $B$. Alors, pour tout $A$-module $M$, nous avons
$$
M \otimes_A^{\mathbf{L}} C =
(M \otimes_A^{\mathbf{L}} B) \otimes_B^{\mathbf{L}} C
$$
et la suite spectrale de
Compléments d'algèbre, exemple \ref{more-algebra-example-tor},
montre que $\text{Tor}^A_p(M, C) = 0$ pour $p > n + m$, comme souhaité.
\end{proof}

\begin{lemma}
\label{lemma-flat-finite-presentation-perfect}
Soit $f : X \to S$ un morphisme de schémas.
Les assertions suivantes sont équivalentes :
\begin{enumerate}
\item $f$ est plat et parfait, et
\item $f$ est plat et localement de présentation finie.
\end{enumerate}
\end{lemma}

\begin{proof}
L'implication (2) $\Rightarrow$ (1) est donnée par
Compléments d'algèbre,
lemme \ref{more-algebra-lemma-flat-finite-presentation-perfect}.
La réciproque résulte du fait qu'un morphisme pseudo-cohérent
est localement de présentation finie, voir le
lemme \ref{lemma-pseudo-coherent-finite-presentation}.
\end{proof}

\begin{lemma}
\label{lemma-regular-target-perfect}
Soit $f : X \to S$ un morphisme de schémas.
Supposons que $S$ soit régulier et que $f$ soit localement de type fini.
Alors $f$ est parfait.
\end{lemma}

\begin{proof}
Voir
Compléments d'algèbre, lemme \ref{more-algebra-lemma-regular-perfect-ring-map}.
\end{proof}

\begin{lemma}
\label{lemma-regular-immersion-perfect}
Une immersion régulière de schémas est parfaite.
Une immersion Koszul-régulière de schémas est parfaite.
\end{lemma}

\begin{proof}
Puisqu'une immersion régulière est une immersion Koszul-régulière, voir
Diviseurs, lemme \ref{divisors-lemma-regular-quasi-regular-immersion},
il suffit de démontrer la seconde assertion. Cela se traduit par l'énoncé
algébrique suivant : supposons que $I \subset A$ soit un
idéal engendré par une suite Koszul-régulière $f_1, \ldots, f_r$ de $A$.
Alors $A \to A/I$ est un homomorphisme parfait d'anneaux. Puisque $A \to A/I$ est surjectif,
ceci est une présentation de $A/I$ par une algèbre polynomiale sur $A$.
Il suffit donc de voir que $A/I$ est pseudo-cohérent comme $A$-module
et est de dimension de Tor finie. Par définition d'une suite de Koszul,
le complexe de Koszul $K(A, f_1, \ldots, f_r)$ est une résolution libre finie
de $A/I$. Ainsi, $A/I$ est un complexe parfait de $A$-modules, ce qui conclut.
\end{proof}

\begin{lemma}
\label{lemma-perfect-permanence}
Soit
$$
\xymatrix{
X \ar[rr]_f \ar[rd] & & Y \ar[ld] \\
& S
}
$$
un diagramme commutatif de morphismes de schémas. Supposons que $Y \to S$ soit
lisse et que $X \to S$ soit parfait. Alors $f : X \to Y$ est parfait.
\end{lemma}

\begin{proof}
Nous pouvons factoriser $f$ comme la composée
$$
X \longrightarrow X \times_S Y \longrightarrow Y
$$
où le premier morphisme est l'application $i = (1, f)$ et le second
morphisme est la projection. Puisque $Y \to S$ est plat, voir
Morphismes, lemme \ref{morphisms-lemma-smooth-flat},
nous voyons que $X \times_S Y \to Y$ est parfait d'après le
lemme \ref{lemma-flat-base-change-perfect}.
Puisque $Y \to S$ est lisse, $X \times_S Y \to X$ est également lisse, voir
Morphismes, lemme \ref{morphisms-lemma-base-change-smooth}.
Ainsi, $i$ est une section d'un morphisme lisse ; par conséquent, $i$ est
une immersion régulière, voir
Diviseurs, lemme \ref{divisors-lemma-section-smooth-regular-immersion}.
Cela entraîne que $i$ est parfait, voir le
lemme \ref{lemma-regular-immersion-perfect}.
Nous concluons que $f$ est parfait, car la composée de morphismes
parfaits est parfaite, voir le
lemme \ref{lemma-composition-perfect}.
\end{proof}

\begin{remark}
\label{remark-perfect-permanence}
Il n'est pas vrai qu'un morphisme entre des schémas $X, Y$ parfaits sur une base $S$
soit parfait. Un exemple est $S = \Spec(k)$, $X = \Spec(k)$,
$Y = \Spec(k[x]/(x^2)$ et $X \to Y$ l'unique $S$-morphisme.
\end{remark}

\begin{lemma}
\label{lemma-descending-property-perfect}
La propriété $\mathcal{P}(f) =$``$f$ est parfait''
est locale pour la topologie fpqc sur la base.
\end{lemma}

\begin{proof}
Nous allons utiliser le critère du
lemme \ref{descent-lemma-descending-properties-morphisms} de Descente
pour le démontrer. D'après la
définition \ref{definition-perfect},
la propriété d'être parfait est locale pour la topologie de Zariski sur la base. D'après le
lemme \ref{lemma-flat-base-change-perfect},
la propriété d'être parfait est préservée par changement de base plat.
La dernière hypothèse (3) du
lemme \ref{descent-lemma-descending-properties-morphisms} de Descente
se traduit par l'énoncé algébrique suivant :
soit $A \to B$ un homomorphisme fidèlement plat d'anneaux.
Soit $C = A[x_1, \ldots, x_n]/I$ une $A$-algèbre.
Si $C \otimes_A B$ est parfait comme $B[x_1, \ldots, x_n]$-module,
alors $C$ est parfait comme $A[x_1, \ldots, x_n]$-module.
C'est le
lemme \ref{more-algebra-lemma-flat-descent-perfect} de Compléments d'algèbre.
\end{proof}

\begin{lemma}
\label{lemma-check-perfect-stalks}
Soit $f : X \to S$ un morphisme pseudo-cohérent de schémas.
Les assertions suivantes sont équivalentes :
\begin{enumerate}
\item $f$ est parfait,
\item $\mathcal{O}_X$ est localement de dimension de Tor finie comme
faisceau de $f^{-1}\mathcal{O}_S$-modules, et
\item pour tout $x \in X$, l'anneau $\mathcal{O}_{X, x}$ est de dimension de Tor
finie comme $\mathcal{O}_{S, f(x)}$-module.
\end{enumerate}
\end{lemma}

\begin{proof}
Le problème est local sur $X$ et $S$. Nous pouvons donc supposer que
$X = \Spec(B)$, $S = \Spec(A)$ et que $f$ correspond à un homomorphisme pseudo-cohérent
d'anneaux $A \to B$.

\medskip\noindent
Si (1) est satisfaite, alors $B$ est de dimension de Tor finie, égale à $d$, comme $A$-module. Alors
$B_\mathfrak q$ est de dimension de Tor $d$ comme $A_\mathfrak p$-module pour tout
idéal premier $\mathfrak q \subset B$ avec $\mathfrak p = A \cap \mathfrak q$, voir
Compléments d'algèbre, lemme \ref{more-algebra-lemma-tor-amplitude-localization}.
Alors $\mathcal{O}_X$ est de dimension de Tor $d$ comme
faisceau de $f^{-1}\mathcal{O}_S$-modules d'après
Cohomologie, lemme \ref{cohomology-lemma-tor-amplitude-stalk}.
Ainsi, (1) implique (2).

\medskip\noindent
D'après Cohomologie, lemme \ref{cohomology-lemma-tor-amplitude-stalk}, (2) implique
(3).

\medskip\noindent
Supposons (3). Nous ne pouvons pas utiliser
Compléments d'algèbre, lemme \ref{more-algebra-lemma-tor-amplitude-localization},
pour conclure, car nous ne savons pas que la dimension de Tor de
$B_\mathfrak q$ sur $A_\mathfrak p$ est bornée indépendamment de $\mathfrak q$.
Choisissons une présentation $A[x_1, \ldots, x_n] \to B$. Alors $B$ est
pseudo-cohérent comme $A[x_1, \ldots, x_n]$-module. Soit
$\mathfrak q \subset A[x_1, \ldots, x_n]$ un idéal premier
au-dessus de $\mathfrak p \subset A$. Alors, soit $B_\mathfrak q$ est nul,
soit, par hypothèse, il est de dimension de Tor finie comme
$A_\mathfrak p$-module. Puisque les fibres de $A \to A[x_1, \ldots, x_n]$
sont de dimension globale finie, nous pouvons appliquer le
lemme \ref{more-algebra-lemma-perfect-over-polynomial-ring} de Compléments d'algèbre
à $A_\mathfrak p \to A[x_1, \ldots, x_n]_\mathfrak q$
pour voir que $B_\mathfrak q$ est un
$A[x_1, \ldots, x_n]_\mathfrak q$-module parfait. Par conséquent,
$B$ est un $A[x_1, \ldots, x_n]$-module parfait d'après
Compléments d'algèbre, lemme \ref{more-algebra-lemma-check-perfect-stalks}.
Ainsi, $A \to B$ est un homomorphisme parfait d'anneaux par définition.
\end{proof}

\begin{lemma}
\label{lemma-perfect-closed-immersion-perfect-direct-image}
Soit $i : Z \to X$ une immersion fermée parfaite de schémas.
Alors $i_*\mathcal{O}_Z$ est un $\mathcal{O}_X$-module parfait, c'est-à-dire
un objet parfait de $D(\mathcal{O}_X)$.
\end{lemma}

\begin{proof}
Cela résulte presque immédiatement de la définition. En effet, soit
$U = \Spec(A)$ un ouvert affine de $X$. Alors $i^{-1}(U) = \Spec(A/I)$
pour un idéal $I \subset A$, et $A/I$ admet une résolution finie par des
$A$-modules projectifs de type fini d'après
Compléments d'algèbre, lemme \ref{more-algebra-lemma-perfect-ring-map}.
Ainsi, $i_*\mathcal{O}_Z|_U$ peut être représenté par un complexe de longueur
finie de $\mathcal{O}_U$-modules localement libres de type fini.
C'est ce qu'il fallait montrer, voir
Cohomologie, section \ref{cohomology-section-perfect}.
\end{proof}

\begin{lemma}
\label{lemma-perfect-proper-perfect-direct-image}
Soit $S$ un schéma noethérien. Soit $f : X \to S$ un morphisme parfait et propre
de schémas. Soit $E \in D(\mathcal{O}_X)$ parfait. Alors
$Rf_*E$ est un objet parfait de $D(\mathcal{O}_S)$.
\end{lemma}

\begin{proof}
Nous affirmons que le lemme
\ref{perfect-lemma-perfect-direct-image} de Catégories dérivées de schémas s'applique.
Les conditions (1) et (2) sont immédiates. La condition (3) est locale
sur $X$. Nous pouvons donc supposer $X$ et $S$ affines et $E$
représenté par un complexe strictement parfait de $\mathcal{O}_X$-modules.
Il suffit donc de montrer que $\mathcal{O}_X$ est de dimension de Tor finie
comme faisceau de $f^{-1}\mathcal{O}_S$-modules.
Cela équivaut à être parfait d'après le
lemme \ref{lemma-check-perfect-stalks}.
\end{proof}

\begin{lemma}
\label{lemma-perfect-fppf-local-source}
La propriété $\mathcal{P}(f) =$``$f$ est parfait''
est locale pour la topologie fppf sur la source.
\end{lemma}

\begin{proof}
Soit $\{g_i : X_i \to X\}_{i \in I}$ un recouvrement fppf de schémas et soit
$f : X \to S$ un morphisme tel que chaque $f \circ g_i$ soit
parfait. D'après le
lemme \ref{lemma-pseudo-coherent-fppf-local-source},
nous concluons que $f$ est pseudo-cohérent.
Ainsi, d'après le
lemme \ref{lemma-check-perfect-stalks},
il suffit de vérifier que $\mathcal{O}_{X, x}$ est un
$\mathcal{O}_{S, f(x)}$-module de dimension de Tor finie pour tout $x \in X$.
Choisissons $i \in I$ et $x_i \in X_i$ s'envoyant sur $x$. Nous voyons alors que
$\mathcal{O}_{X_i, x_i}$ est de dimension de Tor finie sur
$\mathcal{O}_{S, f(x)}$ et que
$\mathcal{O}_{X, x} \to \mathcal{O}_{X_i, x_i}$ est fidèlement plat.
La conclusion voulue résulte du
lemme \ref{more-algebra-lemma-flat-descent-tor-amplitude} de Compléments d'algèbre.
\end{proof}

\begin{lemma}
\label{lemma-factor-regular-immersion}
Soient $i : Z \to Y$ et $j : Y \to X$ des immersions de schémas.
Supposons que
\begin{enumerate}
\item $X$ soit localement noethérien,
\item $j \circ i$ soit une immersion régulière, et
\item $i$ soit parfait.
\end{enumerate}
Alors $i$ et $j$ sont des immersions régulières.
\end{lemma}

\begin{proof}
Puisque $X$ (et donc $Y$) est localement noethérien, les quatre notions d'immersion
régulière coïncident ; de plus, nous pouvons vérifier si un morphisme est une
immersion régulière au niveau des anneaux locaux, voir
Diviseurs, lemme \ref{divisors-lemma-Noetherian-scheme-regular-ideal}.
Le résultat découle donc de
Algèbre à puissances divisées, lemme \ref{dpa-lemma-regular-sequence}.
\end{proof}









\section{Morphismes d'intersection complète locale}
\label{section-lci}

\noindent
Dans
Diviseurs, section \ref{divisors-section-regular-immersions},
nous avons défini quatre types différents d'immersions régulières : régulières,
Koszul-régulières, $H_1$-régulières et quasi-régulières. Dans cette section,
nous considérons les morphismes $f : X \to S$ qui, localement sur $X$, se factorisent comme
$$
\xymatrix{
X \ar[rr]_i \ar[rd] & & \mathbf{A}^n_S \ar[ld] \\
& S
}
$$
où $i$ est une immersion $*$-régulière pour
$* \in \{\emptyset, Koszul, H_1, quasi\}$.
Cependant, nous ne savons pas démontrer que cette condition ne dépend pas
de la factorisation si $* = \emptyset$, c'est-à-dire lorsque nous demandons que $i$
soit une immersion régulière. D'autre part, nous voulons qu'un
morphisme d'intersection complète locale soit parfait, ce qui ne sera vrai
que si $* = Koszul$ ou $* = \emptyset$. Nous définirons donc un
{\it morphisme d'intersection complète locale} ou un
{\it morphisme de Koszul} comme un morphisme de schémas $f : X \to S$
qui, localement sur $X$, admet une factorisation comme ci-dessus où $i$ est une immersion
Koszul-régulière. Pour voir que cette définition convient, nous démontrons d'abord qu'elle ne dépend pas
des factorisations choisies.

\begin{lemma}
\label{lemma-koszul-independence-factorization}
Soit $S$ un schéma. Soient $U$, $P$, $P'$ des schémas sur $S$.
Soit $u \in U$. Soient $i : U \to P$, $i' : U \to P'$ des immersions sur $S$.
Supposons $P$ et $P'$ lisses sur $S$. Les assertions suivantes sont alors équivalentes :
\begin{enumerate}
\item $i$ est une immersion Koszul-régulière dans un voisinage de $u$, et
\item $i'$ est une immersion Koszul-régulière dans un voisinage de $u$.
\end{enumerate}
\end{lemma}

\begin{proof}
Supposons que $i$ soit une immersion Koszul-régulière dans un voisinage de $u$.
Considérons le morphisme $j = (i, i') : U \to P \times_S P' = P''$.
Puisque $P'' = P \times_S P' \to P$ est lisse, il résulte de
Diviseurs, lemme \ref{divisors-lemma-lift-regular-immersion-to-smooth},
que $j$ est une immersion Koszul-régulière ; il résulte alors de
Diviseurs, lemme \ref{divisors-lemma-push-regular-immersion-thru-smooth},
que $i'$ est une immersion Koszul-régulière.
\end{proof}

\noindent
Avant d'énoncer la définition, faisons la remarque élémentaire
suivante. Soit $f : X \to S$ un morphisme de schémas localement
de type fini. Soit $x \in X$. Il existe alors un voisinage ouvert
$U \subset X$ et une factorisation de $f|_U$ comme composée d'une
immersion $i : U \to \mathbf{A}^n_S$ et de la projection
$\mathbf{A}^n_S \to S$, qui est lisse. Considérons le diagramme
$$
\xymatrix{
X \ar[rd] & U \ar[l] \ar[d] \ar[r]_-i & \mathbf{A}^n_S = P \ar[ld]^\pi \\
& S
}
$$
En fait, on peut procéder ainsi avec tout voisinage ouvert affine
$U$ de $x$ dans $X$ ; voir
Morphismes, lemme \ref{morphisms-lemma-quasi-affine-finite-type-over-S}.

\begin{definition}
\label{definition-lci}
Soit $f : X \to S$ un morphisme de schémas.
\begin{enumerate}
\item Soit $x \in X$. On dit que $f$ est {\it de Koszul en $x$} si $f$
est de type fini en $x$ et s'il existe un voisinage ouvert
et une factorisation de $f|_U$ sous la forme $\pi \circ i$, où $i : U \to P$
est une immersion Koszul-régulière et $\pi : P \to S$ est lisse.
\item On dit que $f$ est un {\it morphisme de Koszul}, ou que
$f$ est un {\it morphisme d'intersection complète locale},
si $f$ est de Koszul en tout point.
\end{enumerate}
\end{definition}

\noindent
Nous avons vu plus haut que le choix de la factorisation
$f|_U = \pi \circ i$ est sans importance : étant donnée une factorisation
de $f|_U$ comme une immersion $i$ suivie d'un morphisme lisse $\pi$, le fait que
$i$ soit ou non Koszul-régulière dans un voisinage de $x$ est une propriété
intrinsèque de $f$ en $x$. Consignons-le ici explicitement sous forme de lemme
afin de pouvoir nous y référer.

\begin{lemma}
\label{lemma-lci}
Soit $f : X \to S$ un morphisme d'intersection complète locale.
Soit $P$ un schéma lisse sur $S$. Soit $U \subset X$ un sous-schéma ouvert
et soit $i : U \to P$ une immersion de schémas sur $S$.
Alors $i$ est une immersion Koszul-régulière.
\end{lemma}

\begin{proof}
C'est la propriété définissant un morphisme d'intersection complète
locale. Voir la discussion ci-dessus.
\end{proof}

\noindent
Il paraît opportun de rassembler ici quelques propriétés communes
à tous les morphismes de Koszul.

\begin{lemma}
\label{lemma-lci-properties}
Soit $f : X \to S$ un morphisme d'intersection complète locale.
Alors
\begin{enumerate}
\item $f$ est localement de présentation finie,
\item $f$ est pseudo-cohérent, et
\item $f$ est parfait.
\end{enumerate}
\end{lemma}

\begin{proof}
Puisqu'un morphisme parfait est pseudo-cohérent
(car un homomorphisme d'anneaux parfait est pseudo-cohérent)
et qu'un morphisme pseudo-cohérent est localement de présentation finie
(car un homomorphisme d'anneaux pseudo-cohérent est de présentation finie),
il suffit de démontrer la dernière assertion. La perfection étant une propriété
locale, on peut supposer que $f$ se factorise sous la forme $\pi \circ i$, où
$\pi$ est lisse et $i$ est une immersion Koszul-régulière.
Une immersion Koszul-régulière est parfaite ; voir
le lemme \ref{lemma-regular-immersion-perfect}.
Un morphisme lisse est parfait puisqu'il est plat et localement de
présentation finie ; voir
le lemme \ref{lemma-flat-finite-presentation-perfect}.
Enfin, une composée de morphismes parfaits est parfaite ; voir
le lemme \ref{lemma-composition-perfect}.
\end{proof}

\begin{lemma}
\label{lemma-affine-lci}
Soit $f : X = \Spec(B) \to S = \Spec(A)$ un morphisme de schémas affines.
Alors $f$ est un morphisme d'intersection complète locale si et seulement si
$A \to B$ est un homomorphisme d'intersection complète locale ; voir
Compléments d'algèbre, définition
\ref{more-algebra-definition-local-complete-intersection}.
\end{lemma}

\begin{proof}
Cela résulte immédiatement des définitions.
\end{proof}

\noindent
Attention : un changement de base d'un morphisme de Koszul n'est pas
en général de Koszul.

\begin{lemma}
\label{lemma-flat-base-change-lci}
Un changement de base plat d'un morphisme d'intersection complète locale est
un morphisme d'intersection complète locale.
\end{lemma}

\begin{proof}
Omis. Indication : cela est vrai parce qu'un changement de base d'un morphisme lisse
est lisse et qu'un changement de base plat d'une immersion Koszul-régulière est une
immersion Koszul-régulière ; voir
Diviseurs, lemme \ref{divisors-lemma-regular-immersion-noetherian}.
\end{proof}

\begin{lemma}
\label{lemma-composition-lci}
Une composée de morphismes d'intersection complète locale
est un morphisme d'intersection complète locale.
\end{lemma}

\begin{proof}
Soient $g : Y \to S$ et $f : X \to Y$ des morphismes d'intersection complète
locale. Soit $x \in X$ et posons $y = f(x)$. Choisissons un voisinage ouvert
$V \subset Y$ de $y$ et une factorisation $g|_V = \pi \circ i$, où
$i : V \to P$ est une immersion Koszul-régulière et $\pi : P \to S$ un morphisme lisse.
Choisissons ensuite un voisinage ouvert $U$ de $x \in X$ et une factorisation
$f|_U = \pi' \circ i'$, où $i' : U \to P'$ est une immersion Koszul-régulière
et $\pi' : P' \to Y$ un morphisme lisse. On peut en fait supposer que
$P' = \mathbf{A}^n_V$ ; voir les discussions qui précèdent et qui suivent la
définition \ref{definition-lci}. Diagramme :
$$
\xymatrix{
X \ar[d] & U \ar[l] \ar[r]_-{i'} & P' = \mathbf{A}^n_V \ar[d] \\
Y \ar[d] &  & V \ar[ll] \ar[r]_i & P \ar[d] \\
S & & & S \ar[lll]
}
$$
Posons $P'' = \mathbf{A}^n_P$. Alors $U \to P' \to P''$ est une
immersion Koszul-régulière en tant que composée
d'immersions Koszul-régulières, à savoir $i'$ et le changement de base plat de
$i$ par $P'' \to P$ ; voir
Diviseurs,
lemme \ref{divisors-lemma-regular-immersion-noetherian}
et
Diviseurs, lemme \ref{divisors-lemma-composition-regular-immersion}.
De plus, $P'' \to P \to S$ est lisse en tant que composée de morphismes lisses ;
voir
Morphismes, lemme \ref{morphisms-lemma-composition-smooth}.
On en conclut que $X \to S$ est de Koszul en $x$, comme voulu.
\end{proof}

\begin{lemma}
\label{lemma-flat-lci}
\begin{slogan}
Un morphisme est plat et d'intersection complète locale si et seulement s'il est syntomique.
\end{slogan}
Soit $f : X \to S$ un morphisme de schémas.
Les assertions suivantes sont équivalentes :
\begin{enumerate}
\item $f$ est plat et est un morphisme d'intersection complète locale, et
\item $f$ est syntomique.
\end{enumerate}
\end{lemma}

\begin{proof}
En travaillant localement dans le cas affine, c'est le contenu de
Compléments d'algèbre, lemme \ref{more-algebra-lemma-syntomic-lci}.
Donnons aussi une démonstration plus géométrique.

\medskip\noindent
Supposons (2). D'après
Morphismes, lemme \ref{morphisms-lemma-syntomic-locally-standard-syntomic},
pour tout point $x$ de $X$, il existe des voisinages ouverts affines
$U$ de $x$ et $V$ de $f(x)$ tels que $f|_U : U \to V$ soit syntomique standard.
Cela signifie que $U = \Spec(R[x_1, \ldots, x_n]/(f_1, \ldots, f_c))
\to V = \Spec(R)$, où $R[x_1, \ldots, x_n]/(f_1, \ldots, f_c)$ est une
intersection complète globale relative sur $R$. D'après
Algèbre,
lemme \ref{algebra-lemma-relative-global-complete-intersection-conormal},
la suite $f_1, \ldots, f_c$ est régulière dans tout anneau local
$R[x_1, \ldots, x_n]_{\mathfrak q}$, pour tout idéal premier
$\mathfrak q \supset (f_1, \ldots, f_c)$. Considérons le complexe de Koszul
$K_\bullet = K_\bullet(R[x_1, \ldots, x_n], f_1, \ldots, f_c)$
dont les groupes d'homologie sont $H_i = H_i(K_\bullet)$. D'après
Compléments d'algèbre, lemme \ref{more-algebra-lemma-regular-koszul-regular},
on voit que $(H_i)_{\mathfrak q} = 0$, $i > 0$, pour tout $\mathfrak q$
comme ci-dessus. D'autre part, d'après
Compléments d'algèbre, lemme \ref{more-algebra-lemma-homotopy-koszul},
on voit que $H_i$ est annulé par $(f_1, \ldots, f_c)$. Par conséquent,
$H_i = 0$, $i > 0$, et $f_1, \ldots, f_c$ est une suite Koszul-régulière.
Cela prouve que $U \to V$ se factorise en une immersion Koszul-régulière
$U \to \mathbf{A}^n_V$ suivie d'un morphisme lisse, comme voulu.

\medskip\noindent
Supposons (1). Alors $f$ est plat et localement de présentation finie
(lemme \ref{lemma-lci-properties}).
Ainsi, d'après
Morphismes, lemme \ref{morphisms-lemma-syntomic-locally-standard-syntomic},
il suffit de montrer que les anneaux locaux $\mathcal{O}_{X_s, x}$
sont des anneaux d'intersection complète locale. Choisissons, localement sur $X$, une factorisation
$f = \pi \circ i$, où $i : X \to P$ est une immersion Koszul-régulière
et $\pi : P \to S$ un morphisme lisse. Remarquons que $X \to P$ est
une immersion quasi-régulière relative sur $S$ ; voir
Diviseurs, définition \ref{divisors-definition-relative-H1-regular-immersion}.
Ainsi, d'après
Diviseurs,
lemme \ref{divisors-lemma-relative-regular-immersion-flat-in-neighbourhood},
on voit que $X \to P$ est une immersion régulière et qu'il en reste de même
après tout changement de base. Chaque fibre est donc une immersion régulière, si bien que
tous les anneaux locaux de toutes les fibres de $X$ sont des intersections complètes locales.
\end{proof}

\begin{lemma}
\label{lemma-regular-immersion-lci}
Une immersion régulière de schémas est un morphisme d'intersection complète locale.
Une immersion Koszul-régulière de schémas est un morphisme d'intersection complète
locale.
\end{lemma}

\begin{proof}
Puisqu'une immersion régulière est une immersion Koszul-régulière ; voir
Diviseurs, lemme \ref{divisors-lemma-regular-quasi-regular-immersion},
il suffit de démontrer la seconde assertion. Celle-ci
résulte immédiatement de la définition.
\end{proof}

\begin{lemma}
\label{lemma-lci-permanence}
Soit
$$
\xymatrix{
X \ar[rr]_f \ar[rd] & & Y \ar[ld] \\
& S
}
$$
un diagramme commutatif de morphismes de schémas. Supposons que $Y \to S$
soit lisse et que $X \to S$ soit un morphisme d'intersection complète locale.
Alors $f : X \to Y$ est un morphisme d'intersection complète locale.
\end{lemma}

\begin{proof}
Cela résulte immédiatement des définitions.
\end{proof}

\begin{lemma}
\label{lemma-morphism-regular-schemes-is-lci}
Soit $f : X \to Y$ un morphisme de schémas. Si $f$ est localement
de type fini et si $X$ et $Y$ sont réguliers, alors
$f$ est un morphisme d'intersection complète locale.
\end{lemma}

\begin{proof}
On peut supposer qu'il existe une factorisation $X \to \mathbf{A}^n_Y \to Y$
où la première flèche est une immersion.
Puisque $Y$ est régulier, $\mathbf{A}^n_Y$ l'est aussi d'après
Algèbre, lemme \ref{algebra-lemma-regular-goes-up}.
Ainsi, $X \to \mathbf{A}^n_Y$ est une immersion régulière d'après
Diviseurs, lemme \ref{divisors-lemma-immersion-regular-regular-immersion}.
\end{proof}

\noindent
Le lemme suivant est d'une autre nature.

\begin{lemma}
\label{lemma-lci-avramov}
Soit
$$
\xymatrix{
X \ar[rr]_f \ar[rd] & & Y \ar[ld] \\
& S
}
$$
un diagramme commutatif de morphismes de schémas. Supposons que
\begin{enumerate}
\item $S$ soit localement noethérien,
\item $Y \to S$ soit localement de type fini,
\item $f : X \to Y$ soit parfait,
\item $X \to S$ soit un morphisme d'intersection complète locale.
\end{enumerate}
Alors $X \to Y$ est un morphisme d'intersection complète locale
et $Y \to S$ est de Koszul en $f(x)$ pour tout $x \in X$.
\end{lemma}

\begin{proof}
Au cours de cette démonstration, tous les schémas seront localement noethériens
et tous les anneaux seront noethériens. Nous utiliserons sans autre mention
le fait que les suites régulières et les suites Koszul-régulières coïncident dans ce
cadre ; voir Compléments d'algèbre, lemme
\ref{more-algebra-lemma-noetherian-finite-all-equivalent}.
De plus, la régularité d'un idéal (resp.\ d'un faisceau d'idéaux)
peut se vérifier sur les anneaux locaux (resp.\ sur les germes) ; voir
Algèbre, lemme \ref{algebra-lemma-regular-sequence-in-neighbourhood}
(resp.\ Diviseurs, lemme \ref{divisors-lemma-Noetherian-scheme-regular-ideal}).

\medskip\noindent
La question est locale. On peut donc supposer que $S$, $X$ et $Y$ sont
affines. Dans cette situation, on peut choisir un diagramme commutatif
$$
\xymatrix{
\mathbf{A}^{n + m}_S \ar[d] & X \ar[l] \ar[d] \\
\mathbf{A}^n_S \ar[d] & Y \ar[l] \ar[ld] \\
S
}
$$
dont les flèches horizontales sont des immersions fermées. Soit $x \in X$ un
point, et considérons le diagramme commutatif correspondant d'anneaux
locaux
$$
\xymatrix{
J \ar[r] &
\mathcal{O}_{\mathbf{A}^{n + m}_S, x} \ar[r] &
\mathcal{O}_{X, x} \\
I \ar[r] \ar[u] &
\mathcal{O}_{\mathbf{A}^n_S, f(x)} \ar[r] \ar[u] &
\mathcal{O}_{Y, f(x)} \ar[u]
}
$$
où $J$ et $I$ sont les noyaux des flèches horizontales.
Puisque $X \to S$ est un morphisme d'intersection complète locale, l'idéal
$J$ est engendré par une suite régulière. Puisque $X \to Y$ est
parfait, l'anneau $\mathcal{O}_{X, x}$ est de dimension de Tor finie sur
$\mathcal{O}_{Y, f(x)}$. On peut donc appliquer
Algèbre à puissances divisées, lemme \ref{dpa-lemma-perfect-map-ci},
pour conclure que $I$ et $J/I$ sont engendrés par des suites régulières.
D'après nos remarques initiales, cela achève la démonstration.
\end{proof}

\begin{lemma}
\label{lemma-lci-to-regular}
Soit
$$
\xymatrix{
X \ar[rr]_f \ar[rd] & & Y \ar[ld] \\
& S
}
$$
un diagramme commutatif de morphismes de schémas. Supposons que
$S$ soit localement noethérien, que $Y \to S$ soit localement de type fini,
que $Y$ soit régulier et que $X \to S$ soit un morphisme d'intersection complète locale.
Alors $f : X \to Y$ est un morphisme d'intersection complète locale
et $Y \to S$ est de Koszul en $f(x)$ pour tout $x \in X$.
\end{lemma}

\begin{proof}
C'est un cas particulier du lemme \ref{lemma-lci-avramov}
compte tenu du lemme \ref{lemma-regular-target-perfect}
(et de Morphismes, lemme \ref{morphisms-lemma-permanence-finite-type}).
\end{proof}

\begin{lemma}
\label{lemma-perfect-conormal-free-lci}
Soit $i : X \to Y$ une immersion. Si
\begin{enumerate}
\item $i$ est parfait,
\item $Y$ est localement noethérien, et
\item le faisceau conormal $\mathcal{C}_{X/Y}$ est localement libre de rang fini,
\end{enumerate}
alors $i$ est une immersion régulière.
\end{lemma}

\begin{proof}
Traduit en algèbre, c'est
Algèbre à puissances divisées, proposition \ref{dpa-proposition-regular-ideal}.
\end{proof}

\begin{lemma}
\label{lemma-lci-NL}
Soit $f : X \to Y$ un morphisme d'intersection complète locale.
Alors le complexe cotangent naïf $\NL_{X/Y}$ est un objet parfait
de $D(\mathcal{O}_X)$ d'amplitude de Tor contenue dans $[-1, 0]$.
\end{lemma}

\begin{proof}
Traduit en algèbre, c'est
Compléments d'algèbre, lemme \ref{more-algebra-lemma-lci-NL}.
Pour effectuer la traduction, utiliser
les lemmes \ref{lemma-affine-lci} et
\ref{lemma-NL-affine}, ainsi que
Catégories dérivées des schémas, lemmes
\ref{perfect-lemma-affine-compare-bounded},
\ref{perfect-lemma-tor-dimension-affine} et
\ref{perfect-lemma-perfect-affine}.
\end{proof}

\begin{lemma}
\label{lemma-perfect-NL-lci}
Soit $f : X \to Y$ un morphisme parfait de schémas localement noethériens.
Les assertions suivantes sont équivalentes :
\begin{enumerate}
\item $f$ est un morphisme d'intersection complète locale,
\item $\NL_{X/Y}$ a une amplitude de Tor contenue dans $[-1, 0]$, et
\item $\NL_{X/Y}$ est parfait d'amplitude de Tor contenue dans $[-1, 0]$.
\end{enumerate}
\end{lemma}

\begin{proof}
Traduit en algèbre, c'est
Algèbre à puissances divisées, lemme \ref{dpa-lemma-perfect-NL-lci}.
Pour effectuer la traduction, utiliser
les lemmes \ref{lemma-affine-lci} et
\ref{lemma-NL-affine}, ainsi que
Catégories dérivées des schémas, lemmes
\ref{perfect-lemma-affine-compare-bounded},
\ref{perfect-lemma-tor-dimension-affine} et
\ref{perfect-lemma-perfect-affine}.
\end{proof}

\begin{lemma}
\label{lemma-flat-fp-NL-lci}
Soit $f : X \to Y$ un morphisme plat de présentation finie.
Les assertions suivantes sont équivalentes :
\begin{enumerate}
\item $f$ est un morphisme d'intersection complète locale,
\item $f$ est syntomique,
\item $\NL_{X/Y}$ a une amplitude de Tor contenue dans $[-1, 0]$, et
\item $\NL_{X/Y}$ est parfait d'amplitude de Tor contenue dans $[-1, 0]$.
\end{enumerate}
\end{lemma}

\begin{proof}
Traduit en algèbre, c'est
Algèbre à puissances divisées, lemme \ref{dpa-lemma-flat-fp-NL-lci}.
Pour effectuer la traduction, utiliser
les lemmes \ref{lemma-affine-lci} et
\ref{lemma-NL-affine}, ainsi que
Catégories dérivées des schémas, lemmes
\ref{perfect-lemma-affine-compare-bounded},
\ref{perfect-lemma-tor-dimension-affine} et
\ref{perfect-lemma-perfect-affine}.
\end{proof}

\noindent
Le lemme suivant donne une caractérisation des morphismes lisses comme morphismes
plats dont la diagonale est parfaite.

\begin{lemma}
\label{lemma-smooth-diagonal-perfect}
Soit $f : X \to Y$ un morphisme de type fini entre schémas localement noethériens.
Notons $\Delta : X \to X \times_Y X$ le morphisme diagonal.
Les assertions suivantes sont équivalentes :
\begin{enumerate}
\item $f$ est lisse,
\item $f$ est plat et $\Delta : X \to X \times_Y X$ est une immersion régulière,
\item $f$ est plat et $\Delta : X \to X \times_Y X$ est un
morphisme d'intersection complète locale,
\item $f$ est plat et $\Delta : X \to X \times_Y X$ est parfait.
\end{enumerate}
\end{lemma}

\begin{proof}
Supposons (1). Alors $f$ est plat d'après
Morphismes, lemme \ref{morphisms-lemma-smooth-flat}.
Les projections $X \times_Y X \to X$ sont lisses d'après
Morphismes, lemme \ref{morphisms-lemma-base-change-smooth}. La diagonale
est donc une section d'un morphisme lisse, donc une immersion régulière ; voir
Diviseurs, lemme \ref{divisors-lemma-section-smooth-regular-immersion}.
Ainsi, (1) $\Rightarrow$ (2).
L'implication (2) $\Rightarrow$ (3) est le
lemme \ref{lemma-regular-immersion-lci}.
L'implication (3) $\Rightarrow$ (4) est le
lemme \ref{lemma-lci-properties}.
L'implication intéressante (4) $\Rightarrow$ (1) résulte immédiatement
d'Algèbre à puissances divisées, lemme \ref{dpa-lemma-perfect-diagonal}.
\end{proof}

\begin{lemma}
\label{lemma-descending-property-lci}
La propriété $\mathcal{P}(f) =$``$f$ est un morphisme d'intersection complète
locale'' est locale pour la topologie fpqc sur la base.
\end{lemma}

\begin{proof}
Soit $f : X \to S$ un morphisme de schémas.
Soit $\{S_i \to S\}$ un recouvrement fpqc de $S$.
Supposons que chaque changement de base $f_i : X_i \to S_i$ de $f$ soit
un morphisme d'intersection complète locale.
Remarquons que cela implique notamment que $f$ est localement de type
fini ; voir
le lemme \ref{lemma-lci-properties}
et
Descente, lemme \ref{descent-lemma-descending-property-locally-finite-type}.
Soit $x \in X$. Choisissons un voisinage ouvert $U$ de $x$ et
une immersion $j : U \to \mathbf{A}^n_S$ sur $S$ (voir la
discussion qui précède la
définition \ref{definition-lci}).
Il faut montrer que $j$ est une immersion Koszul-régulière.
Puisque $f_i$ est un morphisme d'intersection complète locale, on voit
que le changement de base $j_i : U \times_S S_i \to \mathbf{A}^n_{S_i}$
est une immersion Koszul-régulière ; voir
le lemme \ref{lemma-lci}.
Comme $\{\mathbf{A}^n_{S_i} \to \mathbf{A}^n_S\}$ est un
recouvrement fpqc, il résulte de
Descente, lemme \ref{descent-lemma-descending-property-regular-immersion},
que $j$ est une immersion Koszul-régulière, comme voulu.
\end{proof}

\begin{lemma}
\label{lemma-lci-syntomic-local-source}
La propriété $\mathcal{P}(f) =$``$f$ est un morphisme d'intersection complète
locale'' est locale pour la topologie syntomique sur la source.
\end{lemma}

\begin{proof}
Nous utiliserons le critère de
Descente, lemme \ref{descent-lemma-properties-morphisms-local-source},
pour le démontrer. Il résulte des
lemmes \ref{lemma-flat-lci} et
\ref{lemma-composition-lci}
que la propriété d'être un morphisme d'intersection complète locale est préservée par
précomposition par des morphismes syntomiques. Il résulte clairement de la
définition \ref{definition-lci}
que la propriété d'être un morphisme d'intersection complète locale est locale pour la topologie de Zariski sur la
source et le but. Ainsi, d'après le lemme susmentionné de
Descente, lemme \ref{descent-lemma-properties-morphisms-local-source},
il suffit de prouver ce qui suit : supposons que $X' \to X \to Y$ soient des
morphismes de schémas affines, avec $X' \to X$ syntomique et $X' \to Y$
un morphisme d'intersection complète locale. Alors $X \to Y$ est un morphisme d'intersection complète
locale. Pour le voir, remarquons que, dans tous les cas, $X \to Y$ est de
présentation finie d'après
Descente, lemme \ref{descent-lemma-flat-finitely-presented-permanence-algebra}.
Choisissons une immersion fermée $X \to \mathbf{A}^n_Y$. D'après
Algèbre, lemme \ref{algebra-lemma-lift-syntomic},
on peut trouver un recouvrement ouvert affine $X' = \bigcup_{i = 1, \ldots, n} X'_i$
et des morphismes syntomiques $W_i \to \mathbf{A}^n_Y$ relevant les morphismes
$X'_i \to X$, c'est-à-dire tels que l'on ait des diagrammes cartésiens
$$
\xymatrix{
X'_i \ar[d] \ar[r] & W_i \ar[d] \\
X \ar[r] & \mathbf{A}^n_Y
}
$$
Après avoir remplacé $X'$ par $\coprod X'_i$ et posé $W = \coprod W_i$,
on obtient un diagramme cartésien de schémas affines
$$
\xymatrix{
X' \ar[d] \ar[r] & W \ar[d]^h \\
X \ar[r] & \mathbf{A}^n_Y
}
$$
où $h : W \to \mathbf{A}^n_Y$ est syntomique et où $X' \to Y$ reste un morphisme d'intersection
complète locale. Puisque $W \to \mathbf{A}^n_Y$ est ouvert (voir
Morphismes, lemme \ref{morphisms-lemma-fppf-open})
et que $X' \to X$ est surjectif, on voit que $X$ est contenu dans l'image
de $W \to \mathbf{A}^n_Y$. Choisissons une immersion fermée
$W \to \mathbf{A}^{n + m}_Y$ au-dessus de $\mathbf{A}^n_Y$. Le diagramme prend alors la forme
$$
\xymatrix{
X' \ar[d] \ar[r] & W \ar[d]^h \ar[r] & \mathbf{A}^{n + m}_Y \ar[ld] \\
X \ar[r] & \mathbf{A}^n_Y
}
$$
Comme $h$ est syntomique, et donc un morphisme d'intersection complète locale (voir
ci-dessus), le morphisme $W \to \mathbf{A}^{n + m}_Y$ est une immersion Koszul-régulière.
Comme $X' \to Y$ est un morphisme d'intersection complète locale, le morphisme
$X' \to \mathbf{A}^{n + m}_Y$ est une immersion Koszul-régulière. On déduit de
Diviseurs, lemme \ref{divisors-lemma-permanence-regular-immersion},
que $X' \to W$ est une immersion Koszul-régulière. Ainsi, puisque la propriété d'être
une immersion Koszul-régulière est locale pour la topologie fpqc sur le but (voir
Descente, lemme \ref{descent-lemma-descending-property-regular-immersion}),
on conclut que $X \to \mathbf{A}^n_Y$ est une immersion Koszul-régulière,
ce qu'il fallait démontrer.
\end{proof}

\begin{lemma}
\label{lemma-base-change-lci-fibres}
Soit $S$ un schéma. Soit $f : X \to Y$ un morphisme de schémas sur $S$.
Supposons que $X$ et $Y$ soient tous deux plats et localement de présentation finie sur $S$.
Alors l'ensemble
$$
\{x \in X \mid f\text{ est de Koszul en }x\}.
$$
est ouvert dans $X$, et sa formation commute à tout changement de base
$S' \to S$.
\end{lemma}

\begin{proof}
L'ensemble est ouvert par définition (voir la
définition \ref{definition-lci}).
Soit $S' \to S$ un morphisme de schémas. Posons $X' = S' \times_S X$,
$Y' = S' \times_S Y$, et notons $f' : X' \to Y'$ le changement de base de $f$.
Soit $x' \in X'$ un point tel que $f'$ soit de Koszul en $x'$. Notons
$s' \in S'$, $x \in X$, $y' \in Y'$, $y \in Y$, $s \in S$ les images
de $x'$. Remarquons que $f$ est localement de présentation finie ; voir
Morphismes, lemme \ref{morphisms-lemma-finite-presentation-permanence}.
On peut donc choisir un voisinage affine $U \subset X$ de
$x$ et une immersion $i : U \to \mathbf{A}^n_Y$. Notons $U' = S' \times_S U$
et $i' : U' \to \mathbf{A}^n_{Y'}$ le changement de base de $i$.
L'hypothèse que $f'$ est de Koszul en $x'$ implique que $i'$ est une
immersion Koszul-régulière dans un voisinage de $x'$ ; voir
le lemme \ref{lemma-lci}.
Le schéma $X'$ est plat et localement de
présentation finie sur $S'$ en tant que changement de base de $X$ (voir
Morphismes, lemmes \ref{morphisms-lemma-base-change-flat} et
\ref{morphisms-lemma-base-change-finite-presentation}).
Ainsi, $i'$ est une immersion $H_1$-régulière relative sur $S'$
dans un voisinage de $x'$ (voir
Diviseurs, définition \ref{divisors-definition-relative-H1-regular-immersion}).
Le changement de base $i'_{s'} : U'_{s'} \to \mathbf{A}^n_{Y'_{s'}}$ est donc
une immersion $H_1$-régulière dans un voisinage ouvert de $x'$ ; voir
Diviseurs, lemme \ref{divisors-lemma-relative-regular-immersion}
et la discussion qui suit
Diviseurs, définition \ref{divisors-definition-relative-H1-regular-immersion}.
Puisque $s' = \Spec(\kappa(s')) \to \Spec(\kappa(s)) = s$
est un morphisme surjectif, plat et universellement ouvert (voir
Morphismes, lemme \ref{morphisms-lemma-scheme-over-field-universally-open}),
on conclut que le changement de base $i_s : U_s \to \mathbf{A}^n_{Y_s}$ est une
immersion $H_1$-régulière dans un voisinage de $x$ ; voir
Descente, lemme \ref{descent-lemma-descending-property-regular-immersion}.
Enfin, remarquons que $\mathbf{A}^n_Y$ est plat et localement de
présentation finie sur $S$ ; par conséquent,
Diviseurs, lemme \ref{divisors-lemma-fibre-quasi-regular}
implique que $i$ est une immersion (Koszul-)régulière dans un voisinage de $x$
comme souhaité.
\end{proof}

\begin{lemma}
\label{lemma-unramified-lci}
Soit $f : X \to Y$ un morphisme d'intersection complète locale de schémas.
Alors $f$ est non ramifié si et seulement si $f$ est formellement non ramifié et,
dans ce cas, le faisceau conormal $\mathcal{C}_{X/Y}$ est localement libre de rang fini
sur $X$.
\end{lemma}

\begin{proof}
La première assertion résulte immédiatement du
lemme \ref{lemma-unramified-formally-unramified}
et du fait qu'un morphisme d'intersection complète locale est localement
de type fini. Pour calculer le faisceau conormal de $f$, choisissons, localement
sur $X$, une factorisation de $f$ sous la forme $f = p \circ i$, où $i : X \to V$
est une immersion Koszul-régulière et $V \to Y$ est lisse. D'après le
lemme \ref{lemma-two-unramified-morphisms-formally-smooth},
on voit que $\mathcal{C}_{X/Y}$ est localement un facteur direct de
$\mathcal{C}_{X/V}$, qui est localement libre de rang fini puisque $i$ est une immersion
Koszul-régulière (donc quasi-régulière) ; voir
Diviseurs, lemme \ref{divisors-lemma-quasi-regular-immersion}.
\end{proof}

\begin{lemma}
\label{lemma-transitivity-conormal-lci}
Soient $Z \to Y \to X$ des morphismes de schémas formellement non ramifiés.
Supposons que $Z \to Y$ soit un morphisme d'intersection complète locale.
La suite exacte
$$
0 \to i^*\mathcal{C}_{Y/X} \to
\mathcal{C}_{Z/X} \to
\mathcal{C}_{Z/Y} \to 0
$$
du
lemme \ref{lemma-transitivity-conormal}
est une suite exacte courte.
\end{lemma}

\begin{proof}
La question est locale sur $Z$ ; on peut donc supposer qu'il existe une factorisation
$Z \to \mathbf{A}^n_Y \to Y$ du morphisme $Z \to Y$. On obtient alors un
diagramme commutatif
$$
\xymatrix{
Z \ar[r]_{i'} \ar@{=}[d] &
\mathbf{A}^n_Y \ar[r] \ar[d] &
\mathbf{A}^n_X \ar[d] \\
Z \ar[r]^i & Y \ar[r] & X
}
$$
Comme $Z \to Y$ est un morphisme d'intersection complète locale, on voit que
$Z \to \mathbf{A}^n_Y$ est une immersion Koszul-régulière. Par conséquent, d'après
Diviseurs, lemme \ref{divisors-lemma-transitivity-conormal-quasi-regular},
la suite
$$
0 \to (i')^*\mathcal{C}_{\mathbf{A}^n_Y/\mathbf{A}^n_X} \to
\mathcal{C}_{Z/\mathbf{A}^n_X} \to
\mathcal{C}_{Z/\mathbf{A}^n_Y} \to 0
$$
est exacte et localement scindée. Remarquons que
$i^*\mathcal{C}_{Y/X} = (i')^*\mathcal{C}_{\mathbf{A}^n_Y/\mathbf{A}^n_X}$
d'après le
lemme \ref{lemma-universal-thickening-fibre-product-flat},
et remarquons que le diagramme
$$
\xymatrix{
(i')^*\mathcal{C}_{\mathbf{A}^n_Y/\mathbf{A}^n_X} \ar[r] &
\mathcal{C}_{Z/\mathbf{A}^n_X} \\
i^*\mathcal{C}_{Y/X} \ar[u]^{\cong} \ar[r] & \mathcal{C}_{Z/X} \ar[u]
}
$$
est commutatif. La flèche horizontale inférieure est donc une injection
localement scindée. Ceci démontre le lemme.
\end{proof}













\section{Suites exactes de différentielles et de faisceaux conormaux}
\label{section-exact}

\noindent
Dans cette section, nous rassemblons quelques résultats sur les suites exactes de faisceaux
conormaux et de faisceaux de différentielles. En un certain sens, ce sont toutes des
réalisations du triangle des complexes cotangents associé à une paire de
morphismes de schémas composables.

\medskip\noindent
Soient $g : Z \to Y$ et $f : Y \to X$ des morphismes de schémas.
\begin{enumerate}
\item Il existe une suite exacte canonique
$$
g^*\Omega_{Y/X} \to \Omega_{Z/X} \to \Omega_{Z/Y} \to 0,
$$
voir
Morphismes, lemme \ref{morphisms-lemma-triangle-differentials}.
Si $g : Z \to Y$ est lisse ou, plus généralement, formellement lisse, alors cette
suite est une suite exacte courte ; voir
Morphismes, lemme \ref{morphisms-lemma-triangle-differentials-smooth},
ou le lemme \ref{lemma-triangle-differentials-formally-smooth}.
\item Si $g$ est une immersion ou, plus généralement, est formellement non ramifié,
alors il existe une suite exacte canonique
$$
\mathcal{C}_{Z/Y} \to g^*\Omega_{Y/X} \to \Omega_{Z/X} \to 0,
$$
voir Morphismes, lemme \ref{morphisms-lemma-differentials-relative-immersion},
ou le lemme \ref{lemma-universally-unramified-differentials-sequence}.
Si $f \circ g : Z \to X$ est lisse ou, plus généralement, formellement lisse, alors cette
suite est une suite exacte courte ; voir
Morphismes, lemme \ref{morphisms-lemma-differentials-relative-immersion-smooth},
ou le lemme \ref{lemma-differentials-formally-unramified-formally-smooth}.
\item Si $g$ et $f \circ g$ sont des immersions ou, plus généralement,
sont formellement non ramifiés, alors il existe une suite exacte canonique
$$
\mathcal{C}_{Z/X} \to \mathcal{C}_{Z/Y} \to g^*\Omega_{Y/X} \to 0,
$$
voir
Morphismes, lemme \ref{morphisms-lemma-two-immersions}, ou le
lemme \ref{lemma-two-unramified-morphisms}. Si $f : Y \to X$ est lisse ou,
plus généralement, formellement lisse, alors cette suite est une suite exacte courte ;
voir Morphismes, lemme \ref{morphisms-lemma-two-immersions-smooth}, ou le
lemme \ref{lemma-two-unramified-morphisms-formally-smooth}.
\item Si $g$ et $f$ sont des immersions ou, plus généralement, sont formellement non ramifiés, alors
il existe une suite exacte canonique
$$
g^*\mathcal{C}_{Y/X} \to \mathcal{C}_{Z/X} \to \mathcal{C}_{Z/Y} \to 0.
$$
Voir
Morphismes, lemme \ref{morphisms-lemma-transitivity-conormal}, ou le
lemme \ref{lemma-transitivity-conormal}.
Si $g : Z \to Y$ est une immersion régulière\footnote{Il suffit que $g$ soit
une immersion $H_1$-régulière. Remarquons qu'une immersion qui est un
morphisme d'intersection complète locale est Koszul-régulière.}, ou
plus généralement un morphisme d'intersection complète locale, alors cette
suite est une suite exacte courte ; voir
Diviseurs, lemme \ref{divisors-lemma-transitivity-conormal-quasi-regular},
ou le lemme \ref{lemma-transitivity-conormal-lci}.
\end{enumerate}










\section{Morphismes faiblement \'etales}
\label{section-weakly-etale}

\noindent
Un homomorphisme d'anneaux $A \to B$ est faiblement \'etale si $A \to B$ et
$B \otimes_A B \to B$ sont tous deux plats ; voir
Compléments d'algèbre, définition \ref{more-algebra-definition-weakly-etale}.
La notion analogue pour les morphismes de schémas est la suivante.

\begin{definition}
\label{definition-weakly-etale}
Un morphisme de schémas $X \to Y$ est dit {\it faiblement \'etale} ou
{\it absolument plat} si $X \to Y$ ainsi que le morphisme diagonal
$X \to X \times_Y X$ sont plats.
\end{definition}

\noindent
Un morphisme \'etale est faiblement \'etale et, réciproquement,
il se trouve qu'un morphisme faiblement \'etale ressemble effectivement, dans une certaine mesure,
à un morphisme \'etale. Par exemple, si $X \to Y$ est faiblement
\'etale, alors $L_{X/Y} = 0$, comme il résulte de
Complexe cotangent, lemme \ref{cotangent-lemma-when-zero}.
Nous démontrerons plus loin un résultat très précis reliant les morphismes faiblement \'etales
aux morphismes \'etales (voir
Cohomologie pro-\'etale, section \ref{proetale-section-weakly-etale}).
Dans cette section, nous nous en tenons aux faits élémentaires.

\begin{lemma}
\label{lemma-check-weakly-etale-stalks}
Soit $f : X \to Y$ un morphisme de schémas. Les assertions suivantes sont équivalentes :
\begin{enumerate}
\item $X \to Y$ est faiblement \'etale, et
\item pour tout $x \in X$, l'homomorphisme d'anneaux
$\mathcal{O}_{Y, f(x)} \to \mathcal{O}_{X, x}$ est faiblement \'etale.
\end{enumerate}
\end{lemma}

\begin{proof}
Remarquons que, sous chacune des hypothèses (1) et (2), le morphisme $f$ est plat.
On peut donc supposer que $f$ est plat. Soit $x \in X$, d'image $y = f(x)$ dans $Y$.
Il existe des homomorphismes canoniques d'anneaux
$$
\mathcal{O}_{X, x} \otimes_{\mathcal{O}_{Y, y}} \mathcal{O}_{X, x}
\longrightarrow
\mathcal{O}_{X \times_Y X, \Delta_{X/Y}(x)}
\longrightarrow
\mathcal{O}_{X, x}
$$
où le premier homomorphisme est une localisation (donc plat) et le second est une
surjection. La condition (1) signifie que la seconde flèche est plate pour tout $x$.
La condition (2) signifie que la composée est plate pour tout $x$. L'équivalence résulte donc
d'Algèbre, lemme \ref{algebra-lemma-flat-localization}, partie (2).
\end{proof}

\begin{lemma}
\label{lemma-key}
Soit $X \to Y$ un morphisme de schémas tel que
$X \to X \times_Y X$ soit plat. Soit $\mathcal{F}$ un $\mathcal{O}_X$-module.
Si $\mathcal{F}$ est plat sur $Y$, alors $\mathcal{F}$ est plat sur $X$.
\end{lemma}

\begin{proof}
Soit $x \in X$, d'image $y = f(x)$ dans $Y$.
Puisque $X \to X \times_Y X$ est plat, on voit que
$\mathcal{O}_{X, x} \otimes_{\mathcal{O}_{Y, y}} \mathcal{O}_{X, x} \to
\mathcal{O}_{X, x}$ est plat. Le résultat découle donc du
lemme \ref{more-algebra-lemma-key} de Compléments d'algèbre
et des définitions.
\end{proof}

\begin{lemma}
\label{lemma-weakly-etale-characterize}
Soit $f : X \to S$ un morphisme de schémas. Les assertions suivantes sont équivalentes :
\begin{enumerate}
\item Le morphisme $f$ est faiblement \'etale.
\item Pour tous ouverts affines $U \subset X$, $V \subset S$
tels que $f(U) \subset V$, l'homomorphisme d'anneaux
$\mathcal{O}_S(V) \to \mathcal{O}_X(U)$ est faiblement \'etale.
\item Il existe un recouvrement ouvert $S = \bigcup_{j \in J} V_j$
et des recouvrements ouverts $f^{-1}(V_j) = \bigcup_{i \in I_j} U_i$ tels
que chacun des morphismes $U_i \to V_j$, $j\in J, i\in I_j$,
soit faiblement \'etale.
\item Il existe un recouvrement ouvert affine $S = \bigcup_{j \in J} V_j$
et des recouvrements ouverts affines $f^{-1}(V_j) = \bigcup_{i \in I_j} U_i$ tels
que l'homomorphisme d'anneaux $\mathcal{O}_S(V_j) \to \mathcal{O}_X(U_i)$ soit
faiblement \'etale, pour tous $j\in J, i\in I_j$.
\end{enumerate}
De plus, si $f$ est faiblement \'etale, alors, pour
tous sous-schémas ouverts $U \subset X$, $V \subset S$ tels que $f(U) \subset V$,
la restriction $f|_U : U \to V$ est faiblement \'etale.
\end{lemma}

\begin{proof}
Supposons donnés des sous-schémas ouverts $U \subset X$, $V \subset S$ tels que
$f(U) \subset V$. Alors $U \times_V U \subset X \times_Y X$ est ouvert
(Schémas, lemme \ref{schemes-lemma-open-fibre-product})
et la diagonale $\Delta_{U/V}$ de $f|_U : U \to V$ est
la restriction $\Delta_{X/Y}|_U : U \to U \times_V U$.
Puisque la platitude est une propriété locale des morphismes de schémas
(Morphismes, lemme \ref{morphisms-lemma-flat-characterize}),
l'assertion finale du lemme en résulte,
ainsi que l'équivalence de (1) et (3).
Si $X$ et $Y$ sont affines, alors $X \to Y$ est faiblement \'etale
si et seulement si $\mathcal{O}_Y(Y) \to \mathcal{O}_X(X)$ est
faiblement \'etale (utiliser de nouveau
Morphismes, lemme \ref{morphisms-lemma-flat-characterize}).
Ainsi, (1) et (3) sont également équivalentes à (2) et (4).
\end{proof}

\begin{lemma}
\label{lemma-composition-weakly-etale}
Soient $X \to Y \to Z$ des morphismes de schémas.
\begin{enumerate}
\item Si $X \to X \times_Y X$ et $Y \to Y \times_Z Y$ sont plats,
alors $X \to X \times_Z X$ est plat.
\item Si $X \to Y$ et $Y \to Z$ sont faiblement \'etales, alors
$X \to Z$ est faiblement \'etale.
\end{enumerate}
\end{lemma}

\begin{proof}
La partie (1) résulte de la factorisation
$$
X \to X \times_Y X \to X \times_Z X
$$
de la diagonale de $X$ sur $Z$, de l'égalité
$$
X \times_Y X = (X \times_Z X) \times_{(Y \times_Z Y)} Y,
$$
du fait qu'un changement de base d'un morphisme plat est plat, et
du fait que la composée de morphismes plats est plate
(Morphismes, lemmes \ref{morphisms-lemma-base-change-flat} et
\ref{morphisms-lemma-composition-flat}).
La partie (2) résulte de la partie (1) et du fait (que l'on vient d'utiliser)
que la composée de morphismes plats est plate.
\end{proof}

\begin{lemma}
\label{lemma-base-change-weakly-etale}
Soient $X \to Y$ et $Y' \to Y$ des morphismes de schémas, et soit
$X' = Y' \times_Y X$ le changement de base de $X$.
\begin{enumerate}
\item Si $X \to X \times_Y X$ est plat, alors $X' \to X' \times_{Y'} X'$
est plat.
\item Si $X \to Y$ est faiblement \'etale, alors $X' \to Y'$ est faiblement \'etale.
\end{enumerate}
\end{lemma}

\begin{proof}
Supposons que $X \to X \times_Y X$ soit plat. Le morphisme $X' \to X' \times_{Y'} X'$
est le changement de base de $X \to X \times_Y X$ par $Y' \to Y$. Il
est donc plat d'après Morphismes, lemme \ref{morphisms-lemma-base-change-flat}.
Ceci démontre (1). La partie (2) résulte de (1) et du fait (que l'on vient d'utiliser)
qu'un changement de base d'un morphisme plat est plat.
\end{proof}

\begin{lemma}
\label{lemma-go-down}
Soient $X \to Y \to Z$ des morphismes de schémas. Supposons que $X \to Y$ soit
plat et surjectif, et que $X \to X \times_Z X$ soit plat.
Alors $Y \to Y \times_Z Y$ est plat.
\end{lemma}

\begin{proof}
Considérons le diagramme commutatif
$$
\xymatrix{
X \ar[r] \ar[d] & X \times_Z X \ar[d] \\
Y \ar[r] & Y \times_Z Y
}
$$
La flèche horizontale supérieure est plate et les flèches verticales sont plates.
Ainsi, $X$ est plat sur $Y \times_Z Y$. D'après
Morphismes, lemme \ref{morphisms-lemma-flat-permanence},
on voit que $Y$ est plat sur $Y \times_Z Y$.
\end{proof}

\begin{lemma}
\label{lemma-weakly-etale-formally-unramified}
Soit $f : X \to Y$ un morphisme de schémas faiblement \'etale.
Alors $f$ est formellement non ramifié, c'est-à-dire que $\Omega_{X/Y} = 0$.
\end{lemma}

\begin{proof}
Rappelons que $f$ est formellement non ramifié si et seulement si $\Omega_{X/Y} = 0$, d'après le
lemme \ref{lemma-formally-unramified-differentials}.
Grâce au lemme \ref{lemma-weakly-etale-characterize} et à
Morphismes, lemme \ref{morphisms-lemma-differentials-affine},
ceci résulte du cas des anneaux, qui est le
lemme \ref{more-algebra-lemma-formally-unramified} de Compléments d'algèbre.
\end{proof}

\begin{lemma}
\label{lemma-when-weakly-etale}
Soit $f : X \to Y$ un morphisme de schémas. Alors $X \to Y$ est faiblement \'etale
dans chacun des cas suivants :
\begin{enumerate}
\item $X \to Y$ est un monomorphisme plat,
\item $X \to Y$ est une immersion ouverte,
\item $X \to Y$ est plat et non ramifié,
\item $X \to Y$ est \'etale.
\end{enumerate}
\end{lemma}

\begin{proof}
Si (1) est satisfaite, alors $\Delta_{X/Y}$ est un isomorphisme
(Schémas, lemme \ref{schemes-lemma-monomorphism}), donc, assurément,
$f$ est faiblement \'etale. Le cas (2) est un cas particulier de (1).
La diagonale d'un morphisme non ramifié est une immersion ouverte
(Morphismes, lemme \ref{morphisms-lemma-diagonal-unramified-morphism}),
donc elle est plate. Ainsi, un morphisme plat et non ramifié est faiblement \'etale.
Un morphisme \'etale est plat et non ramifié
(Morphismes, lemme \ref{morphisms-lemma-etale-smooth-unramified}) ;
ainsi, (4) résulte de (3).
\end{proof}

\begin{lemma}
\label{lemma-reduced-goes-up}
Soit $f : X \to Y$ un morphisme de schémas.
Si $Y$ est réduit et si $f$ est faiblement \'etale, alors $X$ est réduit.
\end{lemma}

\begin{proof}
Grâce au lemme \ref{lemma-weakly-etale-characterize},
ceci résulte du cas des anneaux, qui est le
lemme de Compléments d'algèbre
\ref{more-algebra-lemma-absolutely-flat-over-absolutely-flat}.
\end{proof}

\noindent
Le lemme suivant utilise un résultat non trivial sur les homomorphismes d'anneaux
faiblement \'etales.

\begin{lemma}
\label{lemma-weakly-etale-strictly-henselian-local-rings}
Soit $f : X \to Y$ un morphisme de schémas.
Les assertions suivantes sont équivalentes :
\begin{enumerate}
\item $f$ est faiblement \'etale, et
\item pour $x \in X$, l'homomorphisme d'anneaux locaux
$\mathcal{O}_{Y, f(x)} \to \mathcal{O}_{X, x}$ induit un isomorphisme
sur les hensélisés stricts.
\end{enumerate}
\end{lemma}

\begin{proof}
Soit $x \in X$ un point d'image $y = f(x)$ dans $Y$.
Choisissons une clôture algébrique séparable $\kappa^{sep}$ de $\kappa(x)$.
Soit $\mathcal{O}_{X, x}^{sh}$ le hensélisé strict
correspondant à $\kappa^{sep}$, et $\mathcal{O}_{Y, y}^{sh}$
le hensélisé strict relatif à la clôture algébrique séparable
de $\kappa(y)$ dans $\kappa^{sep}$.
Considérons le diagramme commutatif
$$
\xymatrix{
\mathcal{O}_{X, x} \ar[r] & \mathcal{O}_{X, x}^{sh} \\
\mathcal{O}_{Y, y} \ar[u] \ar[r] & \mathcal{O}_{Y, y}^{sh} \ar[u]
}
$$
d'homomorphismes locaux d'anneaux locaux ; voir
Algèbre, lemme \ref{algebra-lemma-strictly-henselian-functorial}.
Puisque le hensélisé strict est une colimite filtrante d'homomorphismes d'anneaux \'etales,
Compléments d'algèbre, lemme \ref{more-algebra-lemma-when-weakly-etale},
montre que les homomorphismes horizontaux sont faiblement \'etales.
De plus, les homomorphismes horizontaux sont fidèlement plats d'après
Compléments d'algèbre, lemme \ref{more-algebra-lemma-dumb-properties-henselization}.

\medskip\noindent
Supposons $f$ faiblement \'etale. D'après le lemme \ref{lemma-check-weakly-etale-stalks},
la flèche verticale de gauche est faiblement \'etale. D'après
Compléments d'algèbre, lemmes \ref{more-algebra-lemma-composition-weakly-etale} et
\ref{more-algebra-lemma-weakly-etale-permanence},
la flèche verticale de droite est faiblement \'etale. D'après
Compléments d'algèbre, théorème \ref{more-algebra-theorem-olivier},
on conclut que la flèche verticale de droite est un isomorphisme.

\medskip\noindent
Supposons que $\mathcal{O}_{Y, y}^{sh} \to \mathcal{O}_{X, x}^{sh}$ soit un isomorphisme.
Alors $\mathcal{O}_{Y, y} \to \mathcal{O}_{X, x}^{sh}$ est faiblement \'etale.
Comme $\mathcal{O}_{X, x} \to \mathcal{O}_{X, x}^{sh}$ est fidèlement
plat, on conclut que $\mathcal{O}_{Y, y} \to \mathcal{O}_{X, x}$
est faiblement \'etale d'après
Compléments d'algèbre, lemme \ref{more-algebra-lemma-go-down}.
Ainsi, (2) implique (1) d'après le lemme \ref{lemma-check-weakly-etale-stalks}.
\end{proof}

\begin{lemma}
\label{lemma-normal-goes-up}
Soit $f : X \to Y$ un morphisme de schémas. Si $Y$ est un schéma normal
et si $f$ est faiblement \'etale, alors $X$ est un schéma normal.
\end{lemma}

\begin{proof}
D'après Compléments d'algèbre, lemme \ref{more-algebra-lemma-henselization-normal},
un schéma $S$ est normal si et seulement si, pour tout $s \in S$,
le hensélisé strict de $\mathcal{O}_{S, s}$ est un domaine normal.
Le lemme résulte donc du
lemme \ref{lemma-weakly-etale-strictly-henselian-local-rings}.
\end{proof}

\begin{lemma}
\label{lemma-weakly-etale-permanence}
Soit $S$ un schéma. Soit $f : X \to Y$ un morphisme de schémas sur $S$.
Si $X$ et $Y$ sont faiblement \'etales sur $S$, alors $f$ est faiblement \'etale.
\end{lemma}

\begin{proof}
Nous utiliserons sans autre mention Morphismes, lemmes \ref{morphisms-lemma-base-change-flat} et
\ref{morphisms-lemma-composition-flat}.
Écrivons $X \to Y$ comme la composée $X \to X \times_S Y \to Y$.
Le second morphisme est plat, car c'est le changement de base du morphisme plat
$X \to S$. Le premier est le changement de base du morphisme plat
$Y \to Y \times_S Y$ par le morphisme $X \times_S Y \to Y \times_S Y$ ;
il est donc plat. Ainsi, $X \to Y$ est plat. Le morphisme
$X \times_Y X \to X \times_S X$ est une immersion.
Ainsi, le lemme \ref{lemma-key} implique que, puisque
$X$ est plat sur $X \times_S X$, il s'ensuit que $X$ est
plat sur $X \times_Y X$.
\end{proof}

\noindent
Le lemme suivant est une généralisation schématique de l'observation
selon laquelle une extension de corps à la fois séparable et purement inséparable
est nécessairement un isomorphisme.

\begin{lemma}
\label{lemma-weakly-etale-universal-homeomorphism}
Soit $f : X \to Y$ un morphisme de schémas. Si $f$ est faiblement \'etale et
est un homéomorphisme universel, alors c'est un isomorphisme.
\end{lemma}

\begin{proof}
Puisque $f$ est un homéomorphisme universel, la diagonale
$\Delta : X \to X \times_Y X$ est une immersion fermée surjective d'après
Morphismes, lemmes \ref{morphisms-lemma-homeomorphism-affine} et
\ref{morphisms-lemma-universally-injective}. Comme $\Delta$ est aussi
plate, on voit que $\Delta$ est nécessairement un isomorphisme d'après
Morphismes, lemme \ref{morphisms-lemma-characterize-flat-closed-immersions}.
Autrement dit, $f$ est un monomorphisme
(Schémas, lemme \ref{schemes-lemma-monomorphism}).
Puisque $f$ est un homéomorphisme universel, il est certainement
quasi-compact. Ainsi, d'après Descente, lemme
\ref{descent-lemma-flat-surjective-quasi-compact-monomorphism-isomorphism},
on trouve que $f$ est un isomorphisme.
\end{proof}

\noindent
Le résultat suivant est une généralisation faiblement \'etale de
Morphismes \'etales, lemme \ref{etale-lemma-relative-frobenius-etale}.

\begin{lemma}
\label{lemma-relative-frobenius-weakly-etale}
Soit $U \to X$ un morphisme de schémas faiblement \'etale, où $X$ est un schéma
de caractéristique $p$. Alors le Frobenius relatif
$F_{U/X} : U \to U \times_{X, F_X} X$ est un isomorphisme.
\end{lemma}

\begin{proof}
Le morphisme $F_{U/X}$ est un homéomorphisme universel d'après
Variétés, lemme \ref{varieties-lemma-relative-frobenius}.
Le morphisme $F_{U/X}$ est faiblement \'etale en tant que morphisme
entre des schémas faiblement \'etales sur $X$, d'après le
lemme \ref{lemma-weakly-etale-permanence}. Ainsi, $F_{U/X}$
est un isomorphisme d'après le lemme \ref{lemma-weakly-etale-universal-homeomorphism}.
\end{proof}






\section{Théorème de la fibre réduite}
\label{section-reduced-fibre-theorem}

\noindent
Dans cette section, nous examinons le type le plus simple de théorème parmi ceux
annoncés par le titre. Bien que la démonstration du résultat soit quelque peu
laborieuse, elle découle essentiellement de manière directe du
résultat d'Epp sur l'élimination de la ramification ; voir
Compléments d'algèbre, théorème \ref{more-algebra-theorem-epp}.

\medskip\noindent
Soit $A$ un anneau de Dedekind de corps des fractions $K$.
Soit $X$ un schéma plat et de type fini sur $A$.
Soit $L$ une extension finie de $K$. Soit $B$ la clôture intégrale
de $A$ dans $L$. Alors $B$ est un anneau de Dedekind
(Algèbre, lemme \ref{algebra-lemma-integral-closure-Dedekind}).
Soit $X_B = X \times_{\Spec(A)} \Spec(B)$ le changement de base.
Alors $X_B \to \Spec(B)$ est de type fini
(Morphismes, lemme \ref{morphisms-lemma-base-change-finite-type}).
Par conséquent, $X_B$ est noethérien
(Morphismes, lemme \ref{morphisms-lemma-finite-type-noetherian}).
La normalisation $\nu : Y \to X_B$ existe donc (voir
Morphismes, définition \ref{morphisms-definition-normalization},
ainsi que la discussion qui suit). On a le diagramme
\begin{equation}
\label{equation-normalized-base-change}
\xymatrix{
Y \ar[rd] \ar[r]_\nu & X_B \ar[r] \ar[d] & X \ar[d] \\
 & \Spec(B) \ar[r] & \Spec(A)
}
\end{equation}
On appelle parfois $Y$ le {\it changement de base normalisé} de $X$.
En général, le morphisme $\nu$ n'est pas nécessairement fini. Mais si $A$ est
un anneau de Nagata (condition presque toujours satisfaite en
pratique), alors $\nu$ est fini et $Y$ est de type fini sur $B$ ; voir
Morphismes, lemmes \ref{morphisms-lemma-nagata-normalization} et
\ref{morphisms-lemma-finite-type-nagata}.

\medskip\noindent
La formation du changement de base normalisé commute à la composition.
Plus précisément, si $M/L/K$ sont des extensions finies
de corps, de clôtures intégrales $A \subset B \subset C$,
alors le changement de base normalisé $Z$ de $Y \to \Spec(B)$
relativement à $M/L$ est égal au changement de base normalisé
de $X \to \Spec(A)$ relativement à $M/K$.

\begin{theorem}
\label{theorem-normalized-base-change-with-reduced-fibre}
Soit $A$ un anneau de Dedekind de corps des fractions $K$.
Soit $X$ un schéma plat et de type fini sur $A$.
Supposons que $A$ soit un anneau de Nagata.
Il existe une extension finie $L/K$ telle que
le changement de base normalisé $Y$ soit lisse sur $\Spec(B)$
en tous les points génériques de toutes les fibres.
\end{theorem}

\begin{proof}
Au cours de la démonstration, nous utiliserons à plusieurs reprises que la formation de l'ensemble des points
où un morphisme (plat et de présentation finie), tel que $X \to \Spec(A)$, est
lisse commute au changement de base ; voir
Morphismes, lemme \ref{morphisms-lemma-set-points-where-fibres-smooth}.

\medskip\noindent
Choisissons d'abord une extension finie $L/K$ telle que
$(X_L)_{red}$ soit géométriquement réduit sur $L$ ; voir
Variétés, lemme \ref{varieties-lemma-finite-extension-geometrically-reduced}.
Puisque $Y \to (X_B)_{red}$ est birationnel, on voit, en appliquant
Variétés, lemme \ref{varieties-lemma-generic-points-geometrically-reduced},
que $Y_L$ est lui aussi géométriquement réduit sur $L$.
Par conséquent, $Y_L \to \Spec(L)$ est lisse sur un ouvert dense $V \subset Y_L$ d'après
Variétés, lemme \ref{varieties-lemma-geometrically-reduced-dense-smooth-open}.
Ainsi, le lieu de lissité $U \subset Y$ du morphisme $Y \to \Spec(B)$
est ouvert (d'après Morphismes, définition \ref{morphisms-definition-smooth})
et est dense dans la fibre générique. En remplaçant $A$ par $B$ et $X$ par $Y$,
on se ramène au cas traité dans le paragraphe suivant.

\medskip\noindent
Supposons que $X$ soit normal et que le lieu de lissité $U \subset X$ de $X \to \Spec(A)$
soit dense dans la fibre générique. Cela implique que $U$ est dense dans toutes les fibres sauf
un nombre fini ; voir le lemme \ref{lemma-nowhere-dense-generic-fibre}.
Soient $x_1, \ldots, x_r \in X \setminus U$ les points génériques, en nombre fini,
des composantes irréductibles de $X \setminus U$ qui sont en outre
des points génériques de composantes irréductibles de fibres de $X \to \Spec(A)$.
Posons $\mathcal{O}_i = \mathcal{O}_{X, x_i}$. Soit $A_i$ la localisation de
$A$ en l'idéal maximal correspondant à l'image de $x_i$ dans $\Spec(A)$.
D'après
Compléments d'algèbre, proposition
\ref{more-algebra-proposition-epp-essentially-finite-type},
il existe des extensions finies
$K_i/K$ qui sont des solutions pour l'extension d'anneaux de valuation
discrète $A_i \to \mathcal{O}_i$. Soit $L/K$ une extension finie
dominant toutes les extensions $K_i/K$. Alors $L/K$
est encore une solution pour $A_i \to \mathcal{O}_i$ d'après
Compléments d'algèbre, lemme \ref{more-algebra-lemma-solution-goes-up}.

\medskip\noindent
Considérons le diagramme (\ref{equation-normalized-base-change})
avec l'extension $L/K$ que nous venons de construire. Notons que $U_B \subset X_B$
est lisse sur $B$, donc normal (on peut par exemple utiliser
Algèbre, lemme \ref{algebra-lemma-normal-goes-up}).
Ainsi, $Y \to X_B$ est un isomorphisme au-dessus de $U_B$.
Soit $y \in Y$ un point générique d'une composante irréductible
d'une fibre de $Y \to \Spec(B)$ située au-dessus de l'idéal maximal
$\mathfrak m \subset B$. Supposons que $y \not \in U_B$.
Alors $y$ s'envoie sur l'un des points $x_i$. Il s'ensuit que
$\mathcal{O}_{Y, y}$ est un anneau local de la clôture intégrale
de $\mathcal{O}_i$ dans $R(X) \otimes_K L$ (détails omis).
Par conséquent, puisque $L/K$ est une solution pour
$A_i \to \mathcal{O}_i$, on voit que
$B_\mathfrak m \to \mathcal{O}_{Y, y}$ est formellement lisse
pour la topologie $\mathfrak m_y$-adique
(c'est la définition du fait d'être une « solution »).
Autrement dit, $\mathfrak m\mathcal{O}_{Y, y} = \mathfrak m_y$
et l'extension des corps résiduels est séparable ; voir
Compléments d'algèbre, lemme \ref{more-algebra-lemma-extension-dvrs-formally-smooth}.
L'anneau local
de la fibre en $y$ est donc $\kappa(y)$.
Cela implique que la fibre est lisse sur $\kappa(\mathfrak m)$
en $y$, par exemple d'après Algèbre, lemme \ref{algebra-lemma-separable-smooth}.
Ceci achève la démonstration.
\end{proof}

\begin{lemma}[Variante sur les courbes]
\label{lemma-normalized-base-change-with-reduced-fibre-over-curve}
Soit $f : X \to S$ un morphisme de schémas plat et de type fini.
Supposons que $S$ soit un schéma de Nagata, intègre, de corps des fonctions $K$, et
régulier de dimension $1$. Il existe alors une extension finie $L/K$
telle que, dans le diagramme
$$
\xymatrix{
Y \ar[rd]_g \ar[r]_-\nu & X \times_S T \ar[d] \ar[r] & X \ar[d]_f \\
& T \ar[r] & S
}
$$
le morphisme $g$ soit lisse en tous les points génériques des fibres. Ici,
$T$ est la normalisation de $S$ dans $\Spec(L)$ et $\nu : Y \to X \times_S T$
est la normalisation.
\end{lemma}

\begin{proof}
Choisissons un recouvrement ouvert affine fini $S = \bigcup \Spec(A_i)$.
Alors $K$ est égal au corps des fractions de $A_i$ pour tout $i$.
Posons $X_i = X \times_S \Spec(A_i)$.
Choisissons $L_i/K$ comme dans le
théorème \ref{theorem-normalized-base-change-with-reduced-fibre}
pour le morphisme $X_i \to \Spec(A_i)$.
Soit $B_i \subset L_i$ la clôture intégrale de $A_i$ et
soit $Y_i$ le changement de base normalisé de $X$ par $B_i$.
Soit $L/K$ une extension finie dominant chacun des $L_i$.
Soit $T_i \subset T$ l'image réciproque de $\Spec(A_i)$.
Pour tout $i$, on obtient un diagramme commutatif
$$
\xymatrix{
g^{-1}(T_i) \ar[r] \ar[d] & Y_i \ar[r] \ar[d] & X \times_S \Spec(A_i) \ar[d] \\
T_i \ar[r] & \Spec(B_i) \ar[r] & \Spec(A_i)
}
$$
et, en fait, le carré de gauche est un changement de base normalisé,
comme on l'a expliqué au début de la section. Dans la démonstration
du théorème \ref{theorem-normalized-base-change-with-reduced-fibre},
nous avons vu que le lieu de lissité de $Y \to T$ contient
l'image réciproque dans $g^{-1}(T_i)$ de l'ensemble des points
où $Y_i$ est lisse sur $B_i$. Cela démontre le lemme.
\end{proof}

\begin{lemma}[Variante avec une extension séparable]
\label{lemma-normalized-base-change-with-reduced-fibre-separable}
Soit $A$ un anneau de Dedekind de corps des fractions $K$.
Soit $X$ un schéma plat et de type fini sur $A$.
Supposons que $A$ soit un anneau de Nagata et que, pour tout point générique
$\eta$ d'une composante irréductible de $X$, l'extension de corps
$\kappa(\eta)/K$ soit séparable.
Il existe alors une extension finie séparable $L/K$ telle que
le changement de base normalisé $Y$ soit lisse sur $\Spec(B)$
en tous les points génériques de toutes les fibres.
\end{lemma}

\begin{proof}
La démonstration est exactement la même que celle du
théorème \ref{theorem-normalized-base-change-with-reduced-fibre},
à quelques modifications mineures près. Le changement le plus important
consiste à utiliser Compléments d'algèbre, lemme
\ref{more-algebra-lemma-epp-essentially-finite-type-separable},
au lieu de Compléments d'algèbre, proposition
\ref{more-algebra-proposition-epp-essentially-finite-type}.
Au cours de la démonstration, nous utiliserons à plusieurs reprises que la formation de l'ensemble des points
où un morphisme (plat et de présentation finie), tel que $X \to \Spec(A)$, est
lisse commute au changement de base ; voir
Morphismes, lemme \ref{morphisms-lemma-set-points-where-fibres-smooth}.

\medskip\noindent
Puisque $X$ est plat sur $A$, tout point générique $\eta$ de $X$ s'envoie sur le
point générique de $\Spec(A)$.
Après avoir remplacé $X$ par sa réduction, on peut supposer que $X$ est réduit.
Dans ce cas, $X_K$ est géométriquement réduit sur $K$
d'après Variétés, lemme \ref{varieties-lemma-generic-points-geometrically-reduced}.
Par conséquent, $X_K \to \Spec(K)$ est lisse sur un ouvert dense d'après
Variétés, lemme \ref{varieties-lemma-geometrically-reduced-dense-smooth-open}.
Ainsi, le lieu de lissité $U \subset X$ du morphisme $X \to \Spec(A)$
est ouvert (d'après Morphismes, définition \ref{morphisms-definition-smooth})
et est dense dans la fibre générique. On est ainsi ramené à la situation
du paragraphe suivant.

\medskip\noindent
Supposons que $X$ soit normal et que le lieu de lissité $U \subset X$ de $X \to \Spec(A)$
soit dense dans la fibre générique. Cela implique que $U$ est dense dans toutes les fibres sauf
un nombre fini ; voir le lemme \ref{lemma-nowhere-dense-generic-fibre}.
Soient $x_1, \ldots, x_r \in X \setminus U$ les points génériques, en nombre fini,
des composantes irréductibles de $X \setminus U$ qui sont en outre
des points génériques de composantes irréductibles de fibres de $X \to \Spec(A)$.
Posons $\mathcal{O}_i = \mathcal{O}_{X, x_i}$. Observons que le corps des fractions
de $\mathcal{O}_i$ est le corps résiduel d'un point générique de $X$.
Soit $A_i$ la localisation de $A$ en l'idéal maximal correspondant
à l'image de $x_i$ dans $\Spec(A)$. On peut appliquer Compléments d'algèbre, lemme
\ref{more-algebra-lemma-epp-essentially-finite-type-separable},
et l'on trouve des extensions finies séparables $K_i/K$ qui sont des
solutions pour $A_i \to \mathcal{O}_i$. Soit $L/K$ une extension finie
séparable dominant toutes les extensions $K_i/K$.
Alors $L/K$ est encore une solution pour $A_i \to \mathcal{O}_i$ d'après
Compléments d'algèbre, lemme \ref{more-algebra-lemma-solution-goes-up}.

\medskip\noindent
Considérons le diagramme (\ref{equation-normalized-base-change})
avec l'extension $L/K$ que nous venons de construire. Notons que $U_B \subset X_B$
est lisse sur $B$, donc normal (on peut par exemple utiliser
Algèbre, lemme \ref{algebra-lemma-normal-goes-up}).
Ainsi, $Y \to X_B$ est un isomorphisme au-dessus de $U_B$.
Soit $y \in Y$ un point générique d'une composante irréductible
d'une fibre de $Y \to \Spec(B)$ située au-dessus de l'idéal maximal
$\mathfrak m \subset B$. Supposons que $y \not \in U_B$.
Alors $y$ s'envoie sur l'un des points $x_i$. Il s'ensuit que
$\mathcal{O}_{Y, y}$ est un anneau local de la clôture intégrale
de $\mathcal{O}_i$ dans $R(X) \otimes_K L$ (détails omis).
Par conséquent, puisque $L/K$ est une solution pour
$A_i \to \mathcal{O}_i$, on voit que
$B_\mathfrak m \to \mathcal{O}_{Y, y}$ est formellement lisse
(c'est la définition du fait d'être une « solution »).
Autrement dit, $\mathfrak m\mathcal{O}_{Y, y} = \mathfrak m_y$
et l'extension des corps résiduels est séparable. L'anneau local
de la fibre en $y$ est donc $\kappa(y)$.
Cela implique que la fibre est lisse sur $\kappa(\mathfrak m)$
en $y$, par exemple d'après Algèbre, lemme \ref{algebra-lemma-separable-smooth}.
Ceci achève la démonstration.
\end{proof}

\begin{lemma}[Variante avec des extensions séparables sur les courbes]
\label{lemma-normalized-base-change-with-reduced-fibre-over-curve-separable}
Soit $f : X \to S$ un morphisme de schémas plat et de type fini.
Supposons que $S$ soit un schéma de Nagata, intègre, de corps des fonctions $K$, et
régulier de dimension $1$. Supposons que les extensions de corps $\kappa(\eta)/K$
soient séparables pour tout point générique $\eta$ d'une composante
irréductible de $X$. Il existe alors une extension finie séparable $L/K$
telle que, dans le diagramme
$$
\xymatrix{
Y \ar[rd]_g \ar[r]_-\nu & X \times_S T \ar[d] \ar[r] & X \ar[d]_f \\
& T \ar[r] & S
}
$$
le morphisme $g$ soit lisse en tous les points génériques des fibres. Ici,
$T$ est la normalisation de $S$ dans $\Spec(L)$ et $\nu : Y \to X \times_S T$
est la normalisation.
\end{lemma}

\begin{proof}
Cela résulte du
lemme \ref{lemma-normalized-base-change-with-reduced-fibre-separable}
exactement de la même manière que le
lemme \ref{lemma-normalized-base-change-with-reduced-fibre-over-curve}
résulte du
théorème \ref{theorem-normalized-base-change-with-reduced-fibre}.
\end{proof}








\section{Morphismes ind-quasi-affines}
\label{section-ind-quasi-affine}

\noindent
Un peu de théorie qui sera utilisée plus loin.

\begin{definition}
\label{definition-ind-quasi-affine}
Un schéma $X$ est dit {\it ind-quasi-affine} si tout ouvert quasi-compact de
$X$ est quasi-affine. De même, un morphisme de schémas $X \to Y$
est dit {\it ind-quasi-affine} si $f^{-1}(V)$ est ind-quasi-affine
pour tout ouvert affine $V$ de $Y$.
\end{definition}

\noindent
Un ouvert d'un schéma affine est un exemple de schéma ind-quasi-affine.
Si $X = \bigcup_{i \in I} U_i$ est une réunion d'ouverts quasi-affines telle que
deux quelconques des $U_i$ soient contenus dans un troisième, alors $X$ est ind-quasi-affine.
Un schéma ind-quasi-affine $X$ est séparé, car deux ouverts affines
$U, V$ quelconques sont contenus dans un sous-schéma ouvert séparé de $X$, à savoir
$U \cup V$. De même, un morphisme ind-quasi-affine est séparé.

\begin{lemma}
\label{lemma-ind-quasi-affine-alternative-definition}
Pour un morphisme de schémas $f : X \to Y$, les assertions suivantes sont équivalentes :
\begin{enumerate}
\item $f$ est ind-quasi-affine,
\item pour tout sous-schéma ouvert affine $V \subset Y$ et
tout sous-schéma ouvert quasi-compact $U \subset f^{-1}(V)$,
le morphisme induit $U \to V$ est quasi-affine.
\item
pour un recouvrement $\{ V_j \}_{j \in J}$ de $Y$ par des sous-schémas
ouverts quasi-compacts et quasi-séparés
$V_j \subset Y$, pour tout $j \in J$ et tout sous-schéma
ouvert quasi-compact $U \subset f^{-1}(V_j)$, le morphisme induit
$U \to V_j$ est quasi-affine.
\item pour tout sous-schéma ouvert quasi-compact et quasi-séparé
$V \subset Y$ et tout sous-schéma ouvert quasi-compact
$U \subset f^{-1}(V)$, le morphisme induit $U \to V$ est quasi-affine.
\end{enumerate}
En particulier, la propriété d'être un morphisme ind-quasi-affine
est locale pour la topologie de Zariski sur la base.
\end{lemma}

\begin{proof}
L'équivalence (1) $\Leftrightarrow$ (2)
résulte des définitions et de
Morphismes, lemme \ref{morphisms-lemma-characterize-quasi-affine}.
Pour (2) $\Rightarrow$ (4), soient $U$ et $V$ comme dans (4). D'après
Schémas, lemme \ref{schemes-lemma-quasi-compact-permanence}, le
morphisme induit $U \to V$ est quasi-compact. Ainsi, pour tout ouvert
affine $V' \subset V$, le produit fibré $V' \times_V U$ est quasi-compact,
donc, d'après (2), le morphisme induit $V' \times_V U \to V'$ est quasi-affine.
Ainsi, $U \to V$ est lui aussi quasi-affine d'après
Morphismes, lemme \ref{morphisms-lemma-characterize-quasi-affine}.
Cet argument donne également (3) $\Rightarrow$ (4) : en effet, en conservant les
mêmes notations, les ouverts affines $V' \subset V$ qui sont contenus dans l'un
des $V_j$ recouvrent $V$ ; il faut donc montrer que le morphisme
quasi-compact $V' \times_V U \to V'$ est quasi-affine.
Or, d'après (3), le composé $V' \times_V U \to V' \to V_j$
est quasi-affine et, d'après
Schémas, lemme \ref{schemes-lemma-compose-after-separated}, le morphisme
$V' \to V_j$ est quasi-séparé. Ainsi, $V' \times_V U \to V'$
est quasi-affine d'après
Morphismes, lemme \ref{morphisms-lemma-quasi-affine-permanence}.
Les dernières implications (4) $\Rightarrow$ (2) et (4) $\Rightarrow$ (3)
sont évidentes.
\end{proof}

\begin{lemma}
\label{lemma-ind-quasi-affine-composition}
La propriété d'être un morphisme ind-quasi-affine est stable par composition.
\end{lemma}

\begin{proof}
Soient $f : X \to Y$ et $g : Y \to Z$ des morphismes ind-quasi-affines.
Soient $V \subset Z$ et $U \subset f^{-1}(g^{-1}(V))$ des ouverts quasi-compacts,
avec $V$ de plus quasi-séparé. L'image $f(U)$ est une partie
quasi-compacte de $g^{-1}(V)$ ; elle est donc contenue dans un
ouvert quasi-compact $W \subset g^{-1}(V)$ (réunion d'un nombre fini d'ouverts affines).
On obtient une factorisation $U \to W \to V$. Le morphisme $W \to V$ est quasi-affine
d'après le lemme \ref{lemma-ind-quasi-affine-alternative-definition} ; en
particulier, $W$ est quasi-séparé. Alors, encore d'après le
lemme \ref{lemma-ind-quasi-affine-alternative-definition}, $U \to W$
est lui aussi quasi-affine. Par conséquent, d'après Morphismes,
lemme \ref{morphisms-lemma-composition-quasi-affine}, le composé
$U \to V$ est lui aussi quasi-affine, et il reste à appliquer
une fois de plus le lemme \ref{lemma-ind-quasi-affine-alternative-definition}.
\end{proof}

\begin{lemma}
\label{lemma-ind-quasi-affine-examples}
Tout morphisme quasi-affine est ind-quasi-affine.
Toute immersion est ind-quasi-affine.
\end{lemma}

\begin{proof}
La première assertion résulte immédiatement des définitions.
En particulier, les morphismes affines, comme les immersions fermées,
sont ind-quasi-affines. Ainsi, d'après le
lemme \ref{lemma-ind-quasi-affine-composition}, il reste
à montrer qu'une immersion ouverte est ind-quasi-affine.
Mais cela résulte immédiatement des définitions.
\end{proof}

\begin{lemma}
\label{lemma-ind-quasi-affine-permanence}
Si $f : X \to Y$ et $g : Y \to Z$ sont des morphismes de schémas
tels que $g \circ f$ soit ind-quasi-affine, alors $f$ est ind-quasi-affine.
\end{lemma}

\begin{proof}
D'après le lemme \ref{lemma-ind-quasi-affine-alternative-definition}, on peut
travailler localement pour la topologie de Zariski sur $Z$, puis sur $Y$ ; on ne perd donc
rien à supposer que $Z$, puis également $Y$, est affine. Tout ouvert quasi-compact
de $X$ est alors quasi-affine, et le
lemme \ref{lemma-ind-quasi-affine-alternative-definition} donne l'assertion.
\end{proof}

\begin{lemma}
\label{lemma-base-change-ind-quasi-affine}
La propriété d'être ind-quasi-affine est stable par changement de base.
\end{lemma}

\begin{proof}
Soit $f : X \to Y$ un morphisme ind-quasi-affine.
Pour vérifier que tout changement de base de $f$ est ind-quasi-affine, d'après le
lemme \ref{lemma-ind-quasi-affine-alternative-definition}, on peut travailler
localement pour la topologie de Zariski sur $Y$ ; on suppose donc que $Y$ est affine.
En outre, on peut aussi supposer que, dans le morphisme de changement de base
$Z \to Y$, le schéma $Z$ est affine. Le changement de base
$X \times_Y Z \to X$ est un morphisme affine ; d'après les
lemmes \ref{lemma-ind-quasi-affine-composition} et
\ref{lemma-ind-quasi-affine-examples},
le morphisme $X \times_Y Z \to Y$ est donc ind-quasi-affine. Alors, d'après le
lemme \ref{lemma-ind-quasi-affine-permanence}, le changement
de base $X \times_Y Z \to Z$ est ind-quasi-affine, comme souhaité.
\end{proof}

\begin{lemma}
\label{lemma-descending-property-ind-quasi-affine}
La propriété d'être ind-quasi-affine est locale pour la topologie fpqc sur la base.
\end{lemma}

\begin{proof}
La stabilité de l'ind-quasi-affinité par changement de base,
fournie par le lemme \ref{lemma-base-change-ind-quasi-affine},
donne une implication. Pour l'autre, soit $f : X \to Y$
un morphisme de schémas et soit $\{g_i : Y_i \to Y\}$
un recouvrement fpqc tel que le changement de base $f_i : X_i \to Y_i$
soit ind-quasi-affine pour tout $i$. Il faut montrer que $f$ est ind-quasi-affine.

\medskip\noindent
D'après le lemme \ref{lemma-ind-quasi-affine-alternative-definition}, on peut travailler
localement pour la topologie de Zariski sur $Y$ ; on suppose donc que $Y$ est affine.
On utilise alors la stabilité par changement de base assurée par le
lemme \ref{lemma-base-change-ind-quasi-affine} pour raffiner le recouvrement
et supposer qu'il est donné par un unique morphisme affine et fidèlement plat
$g : Y' \to Y$. Pour tout ouvert quasi-compact $U \subset X$, son
changement de base par $Y'$, $U \times_Y Y' \subset X \times_Y Y'$, est lui aussi quasi-compact.
Il reste à observer que, d'après
Descente, lemme \ref{descent-lemma-descending-property-quasi-affine},
le morphisme $U \to Y$ est quasi-affine si et seulement si $U \times_Y Y' \to Y'$ l'est.
\end{proof}

\begin{lemma}
\label{lemma-etale-separated-ind-quasi-affine}
Un morphisme de schémas séparé et localement quasi-fini est ind-quasi-affine.
\end{lemma}

\begin{proof}
Soit $f : X \to Y$ un morphisme de schémas séparé et localement quasi-fini.
Soient $V \subset Y$ un ouvert affine et $U \subset f^{-1}(V)$ un ouvert
quasi-compact. Il faut montrer que $U$ est quasi-affine. Comme $U \to V$ est un
morphisme de schémas séparé et quasi-fini, cela résulte du
théorème principal de Zariski. Voir le
lemme \ref{lemma-quasi-finite-separated-quasi-affine}.
\end{proof}




\section{Sommes amalgamées dans la catégorie des schémas, II}
\label{section-pushouts-II}

\noindent
Cette section prolonge la section \ref{section-pushouts}.
Dans cette section, nous construisons des sommes amalgamées de $Y \leftarrow Z \rightarrow X$
lorsque $Z \to X$ est une immersion fermée, que $Z \to Y$ est entier et
qu'une condition supplémentaire est satisfaite.
On trouvera une discussion détaillée dans \cite{Ferrand-Conducteur}.

\begin{situation}
\label{situation-pushout-along-closed-immersion-and-integral}
Ici, $S$ est un schéma, et $i : Z \to X$ et $j : Z \to Y$
sont des morphismes de schémas sur $S$. On suppose que
\begin{enumerate}
\item $i$ est une immersion fermée,
\item $j$ est un morphisme entier de schémas,
\item pour tout $y \in Y$, il existe un ouvert affine $U \subset X$
tel que $j^{-1}(\{y\}) \subset i^{-1}(U)$.
\end{enumerate}
\end{situation}

\begin{lemma}
\label{lemma-prepare-pushout-along-closed-immersion-and-integral}
Dans la situation \ref{situation-pushout-along-closed-immersion-and-integral},
pour tout $y \in Y$, il existe des ouverts affines $U \subset X$ et
$V \subset Y$ tels que $i^{-1}(U) = j^{-1}(V)$ et $y \in V$.
\end{lemma}

\begin{proof}
Soit $y \in Y$. Choisissons un ouvert affine $U \subset X$
tel que $j^{-1}(\{y\}) \subset i^{-1}(U)$ (ce qui est possible par hypothèse).
Choisissons un voisinage ouvert affine $V \subset Y$ de $y$
tel que $j^{-1}(V) \subset i^{-1}(U)$.
C'est possible parce que $j : Z \to Y$ est un morphisme fermé
(Morphismes, lemme \ref{morphisms-lemma-integral-universally-closed}) et que
$i^{-1}(U)$ contient la fibre au-dessus de $y$.
Comme $j$ est entier, la fibre schématique $Z_y$
est le spectre d'une algèbre entière sur un corps.
D'après Limites, lemme \ref{limits-lemma-ample-profinite-set-in-principal-affine},
on peut trouver un $\overline{f} \in \Gamma(i^{-1}(U), \mathcal{O}_{i^{-1}(U)})$
tel que $Z_y \subset D(\overline{f}) \subset j^{-1}(V)$.
Puisque $i|_{i^{-1}(U)} : i^{-1}(U) \to U$ est une immersion fermée
de schémas affines, on peut choisir un $f \in \Gamma(U, \mathcal{O}_U)$
dont la restriction à $i^{-1}(U)$ est $\overline{f}$.
Après avoir remplacé $U$ par l'ouvert principal $D(f) \subset U$,
on trouve des ouverts affines $y \in V \subset Y$ et $U \subset X$ tels que
$$
j^{-1}(\{y\}) \subset i^{-1}(U) \subset j^{-1}(V)
$$
Nous répétons maintenant (en un certain sens) l'argument. Plus précisément, nous choisissons
$g \in \Gamma(V, \mathcal{O}_V)$ tel que $y \in D(g)$ et
$j^{-1}(D(g)) \subset i^{-1}(U)$ (ce qui est possible par le même argument
que ci-dessus). On peut alors prendre $f \in \Gamma(U, \mathcal{O}_U)$
dont la restriction à $i^{-1}(U)$ est l'image réciproque de $g$
par $i^{-1}(U) \to V$ (à nouveau pour la même raison que ci-dessus).
On dispose finalement d'ouverts affines $y \in V' = D(g) \subset V \subset Y$
et $U' = D(f) \subset U \subset X$ tels que $j^{-1}(V') = i^{-1}(V')$.
\end{proof}

\begin{proposition}
\label{proposition-pushout-along-closed-immersion-and-integral}
\begin{reference}
\cite[Théorème 7.1, partie iii]{Ferrand-Conducteur}
\end{reference}
Dans la situation \ref{situation-pushout-along-closed-immersion-and-integral},
la somme amalgamée $Y \amalg_Z X$ existe dans la catégorie des schémas. On a le diagramme
$$
\xymatrix{
Z \ar[r]_i \ar[d]_j & X \ar[d]^a \\
Y \ar[r]^-b & Y \amalg_Z X
}
$$
Le diagramme est un carré cartésien, le morphisme $a$ est entier,
le morphisme $b$ est une immersion fermée, et
$$
\mathcal{O}_{Y \amalg_Z X} =
b_*\mathcal{O}_Y \times_{c_*\mathcal{O}_Z} a_*\mathcal{O}_X
$$
comme faisceaux d'anneaux, où $c = a \circ i = b \circ j$.
\end{proposition}

\begin{proof}
Comme espace topologique, on définit $Y \amalg_Z X$ comme la somme amalgamée du
diagramme dans la catégorie des espaces topologiques (Topologie, section
\ref{topology-section-colimits}). Il s'agit simplement de la somme amalgamée
des ensembles sous-jacents (Topologie, lemme \ref{topology-lemma-colimits}),
munie de la topologie quotient.
Sur $Y \amalg_Z X$, on dispose des morphismes de faisceaux d'anneaux
$$
b_*\mathcal{O}_Y \longrightarrow c_*\mathcal{O}_Z
\longleftarrow a_*\mathcal{O}_X
$$
et l'on peut définir
$$
\mathcal{O}_{Y \amalg_Z X} =
b_*\mathcal{O}_Y \times_{c_*\mathcal{O}_Z} a_*\mathcal{O}_X
$$
comme le produit fibré dans la catégorie des faisceaux d'anneaux. Pour montrer que l'on
obtient un schéma, il faut montrer que tout point possède un
voisinage ouvert affine. C'est clair pour les points qui ne sont pas dans l'image de $c$,
car l'image de $c$ est une partie fermée dont le complément est
isomorphe, comme espace annelé, à $(Y \setminus j(Z)) \amalg (X \setminus i(Z))$.

\medskip\noindent
Un point de l'image de $c$ correspond à un unique $y \in Y$
dans l'image de $j$. D'après le
lemme \ref{lemma-prepare-pushout-along-closed-immersion-and-integral},
on trouve des ouverts affines $U \subset X$ et $V \subset Y$ tels que
$y \in V$ et $i^{-1}(U) = j^{-1}(V)$.
Comme la construction du premier paragraphe est manifestement compatible
avec la restriction à des sous-schémas ouverts compatibles, pour montrer qu'elle
produit un schéma, on peut supposer que $X$, $Y$ et $Z$ sont affines.

\medskip\noindent
Si $X = \Spec(A)$, $Y = \Spec(B)$ et $Z = \Spec(C)$ sont affines, alors
Compléments d'algèbre, lemme \ref{more-algebra-lemma-points-of-fibre-product},
montre que $Y \amalg_Z X = \Spec(B \times_C A)$ comme espaces topologiques.
Pour achever la démonstration du fait que $Y \times_Z X$ est un schéma, il suffit de montrer
que, sur $\Spec(B \times_C A)$, le faisceau structural est le produit fibré
des images directes. Cela résulte de l'application de
Compléments d'algèbre, lemme \ref{more-algebra-lemma-diagram-localize},
aux ouverts affines principaux de $\Spec(B \times_C A)$.

\medskip\noindent
La discussion ci-dessus montre que le schéma $Y \amalg_Z X$
possède un recouvrement ouvert affine $Y \amalg_Z X = \bigcup W_i$
tel que $U_i = a^{-1}(W_i)$, $V_i = b^{-1}(W_i)$ et
$\Omega_i = c^{-1}(W_i)$ soient des ouverts affines de $X$, $Y$ et $Z$.
Ainsi, $a$ et $b$ sont affines.
De plus, si $A_i$, $B_i$, $C_i$ sont les anneaux correspondant à
$U_i$, $V_i$, $\Omega_i$, alors $A_i \to C_i$ est surjectif et
$W_i$ correspond à $A_i \times_{C_i} B_i$, qui s'envoie surjectivement sur
$B_i$. Ainsi, $b$ est une immersion fermée.
Le morphisme d'anneaux $A_i \times_{C_i} B_i \to A_i$ est entier d'après
Compléments d'algèbre, lemme \ref{more-algebra-lemma-fibre-product-integral} ;
donc $a$ est entier. Le diagramme est cartésien parce que
$$
C_i \cong B_i \otimes_{B_i \times_{C_i} A_i} A_i
$$
Cela résulte du fait que $B_i \times_{C_i} A_i \to B_i$
et $A_i \to C_i$ sont des morphismes surjectifs de même noyau.

\medskip\noindent
Enfin, on peut appliquer les lemmes \ref{lemma-basic-example-pushout} et
\ref{lemma-pushout-fpqc-local} pour conclure que notre construction est une somme amalgamée
dans la catégorie des schémas.
\end{proof}

\begin{lemma}
\label{lemma-pushout-separated}
Dans la situation \ref{situation-pushout-along-closed-immersion-and-integral},
si $X$ et $Y$ sont séparés, alors la somme amalgamée $Y \amalg_Z X$
(proposition \ref{proposition-pushout-along-closed-immersion-and-integral})
est séparée. Il en va de même pour « séparé sur $S$ », « quasi-séparé » et
« quasi-séparé sur $S$ ».
\end{lemma}

\begin{proof}
Le morphisme $Y \amalg X \to Y \amalg_Z X$ est surjectif
et universellement fermé. On peut donc appliquer
Morphismes, lemme \ref{morphisms-lemma-image-universally-closed-separated}.
\end{proof}

\begin{lemma}
\label{lemma-pushout-finite-type}
Dans la situation \ref{situation-pushout-along-closed-immersion-and-integral},
supposons que $S$ soit un schéma localement noethérien et que $X$, $Y$ et $Z$
soient localement de type fini sur $S$. Alors la somme amalgamée $Y \amalg_Z X$
(proposition \ref{proposition-pushout-along-closed-immersion-and-integral})
est localement de type fini sur $S$.
\end{lemma}

\begin{proof}
En se plaçant sur des ouverts affines, on retrouve le résultat de
Compléments d'algèbre, lemme \ref{more-algebra-lemma-fibre-product-finite-type}.
\end{proof}

\begin{lemma}
\label{lemma-pushout-functor}
Dans la situation \ref{situation-pushout-along-closed-immersion-and-integral},
supposons donné un diagramme commutatif
$$
\xymatrix{
Y' \ar[d]^g & Z' \ar[l]^{j'} \ar[r]_{i'} \ar[d]^h & X' \ar[d]^f \\
Y & Z \ar[l] \ar[r] & X
}
$$
dont les carrés sont cartésiens et où $f, g, h$ sont séparés et localement quasi-finis. Alors
\begin{enumerate}
\item les sommes amalgamées $Y \amalg_Z X$ et $Y' \amalg_{Z'} X'$ existent,
\item $Y' \amalg_{Z'} X' \to Y \amalg_Z X$ est
séparé et localement quasi-fini, et
\item les carrés
$$
\xymatrix{
Y' \ar[r] \ar[d] & Y' \amalg_{Z'} X' \ar[d] & X' \ar[l] \ar[d] \\
Y \ar[r] & Y \amalg_Z X & X \ar[l]
}
$$
sont cartésiens.
\end{enumerate}
\end{lemma}

\begin{proof}
La somme amalgamée $Y \amalg_Z X$ existe d'après la
proposition \ref{proposition-pushout-along-closed-immersion-and-integral}.
Pour voir que la somme amalgamée $Y' \amalg_{Z'} X'$ existe, vérifions que la
condition (3) de la
situation \ref{situation-pushout-along-closed-immersion-and-integral}
est satisfaite pour $(X', Y', Z', i', j')$.
En effet, soit $y' \in Y'$ et notons $y \in Y$ son image.
Choisissons un ouvert affine $U \subset X$ tel que $i(j^{-1}(y)) \subset U$.
Choisissons un ouvert quasi-compact $U' \subset X'$ contenu dans
$f^{-1}(U)$ et contenant le sous-ensemble quasi-compact $i'((j')^{-1}(\{y'\}))$.
D'après le lemme \ref{lemma-etale-separated-ind-quasi-affine},
on voit que $U'$ est quasi-affine. Puisque $Z'_{y'}$ est le spectre
d'une algèbre entière sur un corps, on peut appliquer
Limites, lemme \ref{limits-lemma-ample-profinite-set-in-principal-affine},
et l'on trouve qu'il existe un sous-schéma ouvert affine de $U'$ contenant
$i'((j')^{-1}(\{y'\}))$, comme souhaité.

\medskip\noindent
L'existence étant vérifiée, établissons les autres assertions.
Localement pour la topologie affine, nous sommes exactement dans la situation de Compléments d'algèbre, lemme
\ref{more-algebra-lemma-properties-algebras-over-fibre-product},
avec $B \to D$ et $A' \to C'$ localement quasi-finis\footnote{Plus précisément,
$X, Y, Z, Y \amalg_Z X, X', Y', Z', Y' \amalg_{Z'} X'$
correspondent à $A', B, A, B', C', D, C, D'$.}.
En particulier, le morphisme $Y' \amalg_{Z'} X' \to Y \amalg_Z X$ est localement
de type fini. Les carrés du diagramme sont cartésiens d'après
Compléments d'algèbre, lemme \ref{more-algebra-lemma-module-over-fibre-product}.
Puisque la propriété d'être localement quasi-fini se vérifie sur les fibres
(Morphismes, lemme \ref{morphisms-lemma-quasi-finite-at-point-characterize}),
on conclut que $Y' \amalg_{Z'} X' \to Y \amalg_Z X$
est localement quasi-fini.

\medskip\noindent
Il reste à vérifier que $Y' \amalg_{Z'} X' \to Y \amalg_Z X$ est séparé.
Remarquons que $Y' \amalg X' \to Y' \amalg_{Z'} X'$ est universellement fermé
et surjectif d'après la
proposition \ref{proposition-pushout-along-closed-immersion-and-integral}.
Comme le morphisme $Y' \amalg X' \to Y \amalg_Z X$ est aussi séparé
(puisqu'il se factorise sous la forme $Y' \amalg X' \to Y \amalg X \to Y \amalg_Z X$),
on conclut par
Morphismes, lemme \ref{morphisms-lemma-image-universally-closed-separated}.
\end{proof}

\begin{lemma}
\label{lemma-pushout-functor-equivalence-flat}
Dans la situation \ref{situation-pushout-along-closed-immersion-and-integral},
la catégorie des schémas plats, séparés et localement quasi-finis
sur la somme amalgamée $Y \amalg_Z X$ est équivalente à la catégorie des
$(X', Y', Z', i', j', f, g, h)$ comme dans le lemme \ref{lemma-pushout-functor},
où $f, g, h$ sont plats. Même énoncé en remplaçant « plat » par
« \'etale ».
\end{lemma}

\begin{proof}
Si l'on part de $(X', Y', Z', i', j', f, g, h)$
comme dans le lemme \ref{lemma-pushout-functor}, avec $f, g, h$ plats
ou \'etales, alors $Y' \amalg_{Z'} X' \to Y \amalg_Z X$ est plat
ou \'etale d'après Compléments d'algèbre, lemme
\ref{more-algebra-lemma-properties-algebras-over-fibre-product}.

\medskip\noindent
Réciproquement, soit
$W \to Y \amalg_Z X$ un morphisme séparé et localement quasi-fini.
Posons $X' = W \times_{Y \amalg_Z X} X$, $Y' = W \times_{Y \amalg_Z X} Y$ et
$Z' = W \times_{Y \amalg_Z X} Z$, avec les morphismes évidents $i', j', f, g, h$.
Formons la somme amalgamée $Y' \amalg_{Z'} X'$. On obtient un morphisme
$$
Y' \amalg_{Z'} X' \longrightarrow W
$$
de schémas au-dessus de $Y \amalg_X Z$ par la propriété universelle de la
somme amalgamée. Si l'on ne suppose pas que $W \to Y \amalg_Z X$ est plat,
alors ce morphisme n'est en général pas un isomorphisme.
(En fait, Compléments d'algèbre, lemme
\ref{more-algebra-lemma-module-over-fibre-product-bis}
montre que la flèche affichée est une immersion fermée, mais pas un isomorphisme
en général.) Cependant, si $W \to Y \times_Z X$ est plat, alors
c'est un isomorphisme d'après Compléments d'algèbre, lemme
\ref{more-algebra-lemma-properties-algebras-over-fibre-product}.
\end{proof}

\noindent
Nous examinons ensuite l'existence dans le cas où
les deux morphismes sont des immersions fermées.

\begin{lemma}
\label{lemma-pushout-along-closed-immersions}
Soient $i : Z \to X$ et $j : Z \to Y$ des immersions fermées de schémas.
Alors la somme amalgamée $Y \amalg_Z X$ existe dans la catégorie des schémas. Diagramme :
$$
\xymatrix{
Z \ar[r]_i \ar[d]_j & X \ar[d]^a \\
Y \ar[r]^-b & Y \amalg_Z X
}
$$
Le diagramme est un carré cartésien, les morphismes $a$ et $b$
sont des immersions fermées, et l'on a une suite exacte courte
$$
0 \to \mathcal{O}_{Y \amalg_Z X} \to
a_*\mathcal{O}_X \oplus b_*\mathcal{O}_Y \to
c_*\mathcal{O}_Z \to 0
$$
où $c = a \circ i = b \circ j$.
\end{lemma}

\begin{proof}
C'est un cas particulier de la
proposition \ref{proposition-pushout-along-closed-immersion-and-integral}.
Remarquons que l'hypothèse (3) de la
situation \ref{situation-pushout-along-closed-immersion-and-integral}
est immédiate, car
les fibres de $j$ ont au plus un point. Enfin, inversons les rôles des flèches
pour conclure que $a$ et $b$ sont tous deux des immersions fermées.
\end{proof}

\begin{lemma}
\label{lemma-pushout-along-closed-immersions-properties-above}
Soient $i : Z \to X$ et $j : Z \to Y$ des immersions fermées de schémas.
Soient $f : X' \to X$ et $g : Y' \to Y$ des morphismes de schémas, et soit
$\varphi : X' \times_{X, i} Z \to Y' \times_{Y, j} Z$
un isomorphisme de schémas au-dessus de $Z$. Considérons le morphisme
$$
h :
X' \amalg_{X' \times_{X, i} Z, \varphi} Y'
\longrightarrow
X \amalg_Z Y
$$
Alors :
\begin{enumerate}
\item $h$ est localement de type fini si et seulement si $f$ et $g$ sont
localement de type fini ;
\item $h$ est plat si et seulement si $f$ et $g$ sont plats ;
\item $h$ est plat et localement de présentation finie si et seulement si
$f$ et $g$ sont plats et localement de présentation finie ;
\item $h$ est lisse si et seulement si $f$ et $g$ sont lisses ;
\item $h$ est \'etale si et seulement si $f$ et $g$ sont \'etales ; et
\item ajouter d'autres propriétés ici si nécessaire.
\end{enumerate}
\end{lemma}

\begin{proof}
On sait que les sommes amalgamées existent d'après le
lemme \ref{lemma-pushout-along-closed-immersions}.
En particulier, on obtient le morphisme $h$.
On peut donc remplacer tous les schémas considérés par des
schémas affines. Dans ce cas, les assertions du lemme
sont équivalentes aux assertions correspondantes de
Compléments d'algèbre, lemme
\ref{more-algebra-lemma-properties-algebras-over-fibre-product}.
\end{proof}





\section{Morphismes relatifs}
\label{section-relative-morphisms}

\noindent
Dans cette section, nous démontrons un résultat de représentabilité que nous utiliserons dans
Groupes fondamentaux, section \ref{pione-section-finite-etale},
pour démontrer un résultat sur la catégorie des revêtements \'etales finis d'un schéma.
Le contenu de cette section est traité dans la généralité appropriée dans
Critères de représentabilité, section
\ref{criteria-section-relative-morphisms}.

\medskip\noindent
Soit $S$ un schéma. Soient $Z$ et $X$ des schémas sur $S$.
Étant donné un schéma $T$ sur $S$, on peut considérer les morphismes
$b : T \times_S Z \to T \times_S X$ au-dessus de $S$. Diagramme :
\begin{equation}
\label{equation-hom}
\vcenter{
\xymatrix{
T \times_S Z \ar[rd] \ar[rr]_b & &
T \times_S X \ar[ld] & Z \ar[rd] & & X \ar[ld] \\
& T \ar[rrr] & & & S
}
}
\end{equation}
Bien entendu, on peut aussi voir $b$ comme un morphisme
$b : T \times_S Z \to X$ tel que le diagramme
$$
\xymatrix{
T \times_S Z \ar[r] \ar[d] \ar@/^1pc/[rrr]_-b &
Z \ar[rd] & & X \ar[ld] \\
T \ar[rr] & & S
}
$$
soit commutatif. Dans cette situation, on peut définir un foncteur
\begin{equation}
\label{equation-hom-functor}
\mathit{Mor}_S(Z, X) : (\Sch/S)^{opp} \longrightarrow \textit{Sets},
\quad
T \longmapsto \{b\text{ comme ci-dessus}\}
\end{equation}
Voici un résultat élémentaire de représentabilité.

\begin{lemma}
\label{lemma-hom-from-finite-free-into-affine}
Soient $Z \to S$ et $X \to S$ des morphismes de schémas affines.
Supposons que $\Gamma(Z, \mathcal{O}_Z)$ soit un
$\Gamma(S, \mathcal{O}_S)$-module libre de rang fini. Alors $\mathit{Mor}_S(Z, X)$
est représentable par un schéma affine sur $S$.
\end{lemma}

\begin{proof}
Écrivons $S = \Spec(R)$. Choisissons une base $\{e_1, \ldots, e_m\}$
de $\Gamma(Z, \mathcal{O}_Z)$ sur $R$. Choisissons une présentation
$$
\Gamma(X, \mathcal{O}_X) = R[\{x_i\}_{i \in I}]/(\{f_k\}_{k \in K}).
$$
Nous noterons $\overline{x}_i$ l'image de $x_i$ dans ce quotient.
Écrivons
$$
P = R[\{a_{ij}\}_{i \in I, 1 \leq j \leq m}].
$$
Considérons l'homomorphisme de $R$-algèbres
$$
\Psi :
R[\{x_i\}_{i \in I}]
\longrightarrow
P \otimes_R \Gamma(Z, \mathcal{O}_Z), \quad
x_i \longmapsto \sum\nolimits_j a_{ij} \otimes e_j.
$$
Écrivons $\Psi(f_k) = \sum c_{kj} \otimes e_j$ avec $c_{kj} \in P$.
Enfin, notons $J \subset P$ l'idéal engendré par les éléments
$c_{kj}$, $k \in K$, $1 \leq j \leq m$. Nous affirmons que
$W = \Spec(P/J)$ représente le foncteur $\mathit{Mor}_S(Z, X)$.

\medskip\noindent
Remarquons d'abord que, par construction, $P/J$ est une $R$-algèbre, donc définit
un morphisme $W \to S$. Ensuite, par construction, l'homomorphisme
$\Psi$ se factorise par $\Gamma(X, \mathcal{O}_X)$ ; on obtient donc
un homomorphisme de $P/J$-algèbres
$$
P/J \otimes_R \Gamma(X, \mathcal{O}_X)
\longrightarrow
P/J \otimes_R \Gamma(Z, \mathcal{O}_Z)
$$
qui détermine un morphisme
$b_{univ} : W \times_S Z \to W \times_S X$.
Par le lemme de Yoneda, $b_{univ}$ détermine une
transformation de foncteurs $W \to \mathit{Mor}_S(Z, X)$ dont nous
affirmons que c'est un isomorphisme. Pour le montrer, il suffit
de montrer qu'elle induit une bijection d'ensembles
$W(T) \to \mathit{Mor}_S(Z, X)(T)$ pour tout schéma affine
$T$.

\medskip\noindent
Supposons que $T = \Spec(R')$ soit un schéma affine sur $S$
et que $b \in \mathit{Mor}_S(Z, X)(T)$. Le morphisme structural $T \to S$
définit une structure de $R$-algèbre sur $R'$, et $b$ définit un homomorphisme de $R'$-algèbres
$$
b^\sharp :
R' \otimes_R \Gamma(X, \mathcal{O}_X)
\longrightarrow
R' \otimes_R \Gamma(Z, \mathcal{O}_Z).
$$
En particulier, on peut écrire
$b^\sharp(1 \otimes \overline{x}_i) = \sum \alpha_{ij} \otimes e_j$
pour certains $\alpha_{ij} \in R'$. Cela correspond à un homomorphisme de $R$-algèbres
$P \to R'$ déterminé par la règle $a_{ij} \mapsto \alpha_{ij}$. Cet
homomorphisme se factorise par le quotient $P/J$ par construction de l'idéal
$J$, ce qui donne un homomorphisme $P/J \to R'$. Celui-ci correspond à son tour à un morphisme
$T \to W$ tel que $b$ soit le changement de base de $b_{univ}$.
Certains détails sont omis.
\end{proof}

\begin{lemma}
\label{lemma-hom-from-finite-locally-free-into-affine}
Soient $Z \to S$ et $X \to S$ des morphismes de schémas.
Si $Z \to S$ est fini localement libre et $X \to S$ est affine,
alors $\mathit{Mor}_S(Z, X)$ est représentable par un schéma
affine sur $S$.
\end{lemma}

\begin{proof}
Choisissons un recouvrement ouvert affine $S = \bigcup U_i$ tel que
$\Gamma(Z \times_S U_i, \mathcal{O}_{Z \times_S U_i})$ soit
libre de rang fini sur $\mathcal{O}_S(U_i)$. Soit $F_i \subset \mathit{Mor}_S(Z, X)$
le sous-foncteur qui associe à $T/S$ l'ensemble vide si
$T \to S$ ne se factorise pas par $U_i$, et $\mathit{Mor}_S(Z, X)(T)$
sinon. Alors la famille de ces sous-foncteurs satisfait les conditions
(2)(a), (2)(b), (2)(c) de
Schémas, lemme \ref{schemes-lemma-glue-functors}, ce qui démontre le lemme.
La condition (2)(a) résulte du
lemme \ref{lemma-hom-from-finite-free-into-affine},
et les deux autres résultent d'arguments immédiats.
\end{proof}

\noindent
La condition imposée au morphisme $f : X \to S$ dans le lemme ci-dessous est très
utile pour démontrer des énoncés de ce type. Elle est satisfaite si l'une des conditions suivantes
l'est : $X$ est quasi-affine, $f$ est quasi-affine, $f$ est quasi-projectif,
$f$ est localement projectif, il existe un faisceau inversible ample sur $X$,
il existe un faisceau inversible $f$-ample sur $X$, ou
il existe un faisceau inversible $f$-très ample sur $X$.

\begin{lemma}
\label{lemma-hom-from-finite-locally-free-representable}
Soient $Z \to S$ et $X \to S$ des morphismes de schémas.
Supposons que
\begin{enumerate}
\item $Z \to S$ soit fini localement libre ; et
\item pour tout $(s, x_1, \ldots, x_d)$ où $s \in S$ et
$x_1, \ldots, x_d \in X_s$, il existe un ouvert affine $U \subset X$
tel que $x_1, \ldots, x_d \in U$.
\end{enumerate}
Alors $\mathit{Mor}_S(Z, X)$ est représentable par un schéma.
\end{lemma}

\begin{proof}
Considérons l'ensemble $I$ des couples $(U, V)$ où $U \subset X$ et $V \subset S$
sont des ouverts affines et où $U \to S$ se factorise par $V$. Pour $i \in I$, notons
$(U_i, V_i)$ le couple correspondant. Posons
$F_i = \mathit{Mor}_{V_i}(Z_{V_i}, U_i)$.
Il est immédiat que $F_i$ est un sous-foncteur de $\mathit{Mor}_S(Z, X)$.
Nous affirmons alors que les conditions
(2)(a), (2)(b), (2)(c) de
Schémas, lemme \ref{schemes-lemma-glue-functors}, sont satisfaites, ce qui démontre le lemme.

\medskip\noindent
La condition (2)(a) résulte du
lemme \ref{lemma-hom-from-finite-locally-free-into-affine}.

\medskip\noindent
Pour vérifier la condition (2)(b), considérons $T/S$ et $b \in \mathit{Mor}_S(Z, X)(T)$.
En considérant $b$ comme un morphisme $T \times_S Z \to X$, on obtient un ouvert
$b^{-1}(U_i) \subset T \times_S Z$. Il est clair que $b \in F_i(T)$
si et seulement si $b^{-1}(U_i) = T \times_S Z$. Puisque la projection
$p : T \times_S Z \to T$ est finie, donc fermée, l'ensemble
$U_{i, b} \subset T$ des points $t \in T$ tels que
$p^{-1}(\{t\}) \subset b^{-1}(U_i)$ est ouvert.
Alors $f : T' \to T$ se factorise par $U_{i, b}$ si et seulement
si $b \circ f \in F_i(T')$, et la vérification de (2)(b) est achevée.

\medskip\noindent
Enfin, vérifions la condition (2)(c) ; c'est ici qu'intervient notre condition
sur $X \to S$. En effet, considérons
$T/S$ et $b \in \mathit{Mor}_S(Z, X)(T)$.
Il suffit de démontrer que tout $t \in T$
appartient à l'un des ouverts $U_{i, b}$ définis
au paragraphe précédent.
Ceci équivaut à la condition que
$b(p^{-1}(\{t\})) \subset U_i$ pour un certain $i$
où $p : T \times_S Z \to T$ est la projection et
$b : T \times_S Z \to X$ est le morphisme donné.
Puisque $p$ est fini, l'ensemble $b(p^{-1}(\{t\})) \subset X$
est fini et contenu dans la fibre de $X \to S$ au-dessus de
l'image $s$ de $t$ dans $S$.
Ainsi, notre condition sur $X \to S$ montre exactement qu'un
couple convenable existe.
\end{proof}

\begin{lemma}
\label{lemma-hom-from-finite-locally-free-separated-lqf}
Soient $Z \to S$ et $X \to S$ des morphismes de schémas.
Supposons que $Z \to S$ soit fini localement libre et que $X \to S$
soit séparé et localement quasi-fini.
Alors $\mathit{Mor}_S(Z, X)$ est représentable par un schéma.
\end{lemma}

\begin{proof}
Ceci résulte des
lemmes \ref{lemma-hom-from-finite-locally-free-representable} et
\ref{lemma-separated-locally-quasi-finite-over-affine}.
\end{proof}









\section{Caractérisation des complexes pseudo-cohérents, III}
\label{section-characterize-pseudo-coherent}

\noindent
Dans cette section, nous étudions des caractérisations des complexes pseudo-cohérents
en termes de cohomologie. Ceci fait suite à
Catégories dérivées de schémas, section
\ref{perfect-section-pseudo-coherent}.
Un outil fondamental consistera à nous ramener au cas de l'espace projectif
au moyen d'une version dérivée du lemme de Chow ; voir le lemme \ref{lemma-derived-chow}.

\begin{lemma}
\label{lemma-case-of-tor-independence}
Considérons un diagramme commutatif de schémas
$$
\xymatrix{
Z' \ar[d] \ar[r] & Y' \ar[d] \\
X' \ar[r] & S'
}
$$
Soit $S \to S'$ un morphisme. Notons $X$ et $Y$ les changements de
base de $X'$ et $Y'$ à $S$.
Supposons que $Y' \to S'$ et $Z' \to X'$ soient plats.
Alors $X \times_S Y$ et $Z'$ sont Tor-indépendants sur $X' \times_{S'} Y'$.
\end{lemma}

\begin{proof}
La question est locale ; on peut donc supposer tous les schémas affines
(certains détails sont omis). Remarquons que
$$
\xymatrix{
X \times_S Y \ar[r] \ar[d] & X' \times_{S'} Y' \ar[d] \\
X \ar[r] & X'
}
$$
est cartésien, avec des flèches verticales plates.
Écrivons $X = \Spec(A)$, $X' = \Spec(A')$ et
$X' \times_{S'} Y' = \Spec(B')$. Alors
$X \times_S Y = \Spec(A \otimes_{A'} B')$.
Écrivons $Z' = \Spec(C')$. Il faut montrer que
$$
\text{Tor}_p^{B'}(A \otimes_{A'} B', C') = 0,
\quad\text{pour } p > 0
$$
Puisque $A' \to B'$ est plat,
on a $A \otimes_{A'} B' = A \otimes_{A'}^\mathbf{L} B'$.
Par conséquent,
$$
(A \otimes_{A'} B') \otimes_{B'}^\mathbf{L} C' =
(A \otimes_{A'}^\mathbf{L} B') \otimes_{B'}^\mathbf{L} C' =
A \otimes_{A'}^\mathbf{L} C' =
A \otimes_{A'} C'
$$
La deuxième égalité résulte de Compléments d'algèbre, lemme
\ref{more-algebra-lemma-double-base-change}.
La dernière égalité vient de ce que $A' \to C'$ est plat. Ceci démontre le lemme.
\end{proof}

\begin{lemma}[Lemme de Chow dérivé]
\label{lemma-derived-chow}
Soit $A$ un anneau. Soit $X$ un schéma séparé de présentation finie
sur $A$. Soit $x \in X$. Alors il existe
un voisinage ouvert $U \subset X$ de $x$,
un entier $n \geq 0$,
un ouvert $V \subset \mathbf{P}^n_A$,
un sous-schéma fermé $Z \subset X \times_A \mathbf{P}^n_A$,
un point $z \in Z$, et
un objet $E$ de $D(\mathcal{O}_{X \times_A \mathbf{P}^n_A})$ tels que
\begin{enumerate}
\item $Z \to X \times_A \mathbf{P}^n_A$ est de présentation finie ;
\item $b : Z \to X$ est un isomorphisme au-dessus de $U$ et $b(z) = x$ ;
\item $c : Z \to \mathbf{P}^n_A$ est une immersion fermée au-dessus de $V$ ;
\item $b^{-1}(U) = c^{-1}(V)$, en particulier $c(z) \in V$ ;
\item $E|_{X \times_A V} \cong
(b, c)_*\mathcal{O}_Z|_{X \times_A V}$ ;
\item $E$ est pseudo-cohérent et à support dans $Z$.
\end{enumerate}
\end{lemma}

\begin{proof}
On peut trouver une sous-$\mathbf{Z}$-algèbre de type fini $A' \subset A$
et un schéma $X'$ séparé et de présentation finie sur $A'$
dont le changement de base à $A$ est $X$. Voir
Limites, lemmes \ref{limits-lemma-descend-finite-presentation} et
\ref{limits-lemma-descend-separated-finite-presentation}.
Soit $x' \in X'$ l'image de $x$.
Si l'on peut démontrer le lemme pour $x' \in X'/A'$, alors
le lemme en résulte pour $x \in X/A$.
En effet, si $U', n', V', Z', z', E'$ fournissent la solution
pour $x' \in X'/A'$, alors on peut prendre pour
$U \subset X$ l'image réciproque de $U'$,
prendre $n = n'$,
prendre pour $V \subset \mathbf{P}^n_A$ l'image réciproque de $V'$,
prendre pour $Z \subset X \times \mathbf{P}^n$
l'image réciproque schématique de $Z'$,
prendre pour $z \in Z$ l'unique point ayant pour image $x$, et
prendre pour $E$ l'image réciproque dérivée de $E'$.
Remarquons que $E$ est pseudo-cohérent d'après
Cohomologie, lemme \ref{cohomology-lemma-pseudo-coherent-pullback}.
Il reste seulement à vérifier (5). Pour cela,
posons $W = b^{-1}(U) = c^{-1}(V)$ et $W' = (b')^{-1}(U) = (c')^{-1}(V')$ 
et considérons le carré cartésien
$$
\xymatrix{
W \ar[d]_{(b, c)} \ar[r] & W' \ar[d]^{(b', c')} \\
X \times_A V \ar[r] & X' \times_{A'} V'
}
$$
D'après le lemme \ref{lemma-case-of-tor-independence}, les schémas
$X \times_A V$ et $W'$ sont Tor-indépendants sur $X' \times_{A'} V'$.
Par conséquent, l'image réciproque dérivée de
$(b', c')_*\mathcal{O}_{W'}$ sur $X \times_A V$
est $(b, c)_*\mathcal{O}_W$ d'après
Catégories dérivées de schémas,
lemme \ref{perfect-lemma-compare-base-change}.
On utilise aussi le fait que $R(b', c')_*\mathcal{O}_{Z'} = (b', c')_*\mathcal{O}_{Z'}$
car $(b', c')$ est une immersion fermée, et de même pour
$(b, c)_*\mathcal{O}_Z$.
Puisque $E'|_{U' \times_{A'} V'} =
(b', c')_*\mathcal{O}_{W'}$, on obtient
$E|_{U \times_A V} = (b, c)_*\mathcal{O}_W$,
et (5) est vérifiée.
Ceci nous ramène à la situation décrite au
paragraphe suivant.

\medskip\noindent
Supposons que $A$ soit de type fini sur $\mathbf{Z}$.
Choisissons un voisinage ouvert affine $U \subset X$ de $x$.
Alors $U$ est de type fini sur $A$.
Choisissons une immersion fermée $U \to \mathbf{A}^n_A$ et notons
$j : U \to \mathbf{P}^n_A$ l'immersion obtenue en la composant
avec l'immersion ouverte $\mathbf{A}^n_A \to \mathbf{P}^n_A$.
Soit $Z$ l'adhérence schématique de
$$
(\text{id}_U, j) : U \longrightarrow X \times_A \mathbf{P}^n_A
$$
Puisque la projection $X \times \mathbf{P}^n \to X$ est séparée,
on conclut de Morphismes, lemme
\ref{morphisms-lemma-scheme-theoretic-image-of-partial-section},
que $b : Z \to X$ est un isomorphisme au-dessus de $U$.
Soit $z \in Z$ l'unique point situé au-dessus de $x$.

\medskip\noindent
Soit $Y \subset \mathbf{P}^n_A$ l'adhérence schématique
de l'image de $j$. Il est alors clair que $Z \subset X \times_A Y$
est l'adhérence schématique de
$(\text{id}_U, j) : U \to X \times_A Y$.
Comme $X$ est séparé, le morphisme
$X \times_A Y \to Y$ est également séparé.
On voit donc que $Z \to Y$ est un isomorphisme au-dessus du
sous-schéma ouvert $j(U) \subset Y$, d'après le même lemme que ci-dessus.
Choisissons un ouvert $V \subset \mathbf{P}^n_A$ tel que $V \cap Y = j(U)$.
On voit alors que (3) et (4) sont satisfaites.

\medskip\noindent
Puisque $A$ est noethérien, on voit que $X$ et $X \times_A \mathbf{P}^n_A$
sont des schémas noethériens. On peut donc prendre $E = (b, c)_*\mathcal{O}_Z$
dans ce cas ; voir Catégories dérivées de schémas, lemme
\ref{perfect-lemma-identify-pseudo-coherent-noetherian}.
Ceci achève la démonstration.
\end{proof}

\begin{lemma}
\label{lemma-compute-Fourier-Mukai-for-derived-chow}
Soient $A$, $x \in X$, ainsi que
$U, n, V, Z, z, E$ comme dans le lemme \ref{lemma-derived-chow}.
Pour tout $K \in D_\QCoh(\mathcal{O}_X)$, on a
$$
Rq_*(Lp^*K \otimes^\mathbf{L} E)|_V = R(U \to V)_*K|_U
$$
où $p : X \times_A \mathbf{P}^n_A \to X$ et
$q : X \times_A \mathbf{P}^n_A \to \mathbf{P}^n_A$ sont
les projections et où le morphisme $U \to V$ est
l'immersion fermée de présentation finie $c \circ (b|_U)^{-1}$.
\end{lemma}

\begin{proof}
Puisque $b^{-1}(U) = c^{-1}(V)$ et que $c$ est une immersion fermée
au-dessus de $V$, on voit que $c \circ (b|_U)^{-1}$ est une immersion fermée.
Elle est de présentation finie parce que $U$ et $V$ sont de présentation
finie sur $A$ ; voir
Morphismes, lemme \ref{morphisms-lemma-finite-presentation-permanence}.
On a d'abord
$$
Rq_*(Lp^*K \otimes^\mathbf{L} E)|_V =
Rq'_*\left((Lp^*K \otimes^\mathbf{L} E)|_{X \times_A V}\right)
$$
où $q' : X \times_A V \to V$ est la projection, car
la formation de l'image directe totale commute à la localisation.
Posons $W = b^{-1}(U) = c^{-1}(V)$ et notons $i : W \to X \times_A V$
l'immersion fermée $i = (b, c)|_W$. Alors
$$
Rq'_*\left((Lp^*K \otimes^\mathbf{L} E)|_{X \times_A V}\right) =
Rq'_*(Lp^*K|_{X \times_A V} \otimes^\mathbf{L} i_*\mathcal{O}_W)
$$
d'après la propriété (5). Puisque $i$ est une immersion fermée, on a
$i_*\mathcal{O}_W = Ri_*\mathcal{O}_W$.
En utilisant
Catégories dérivées de schémas,
lemme \ref{perfect-lemma-cohomology-base-change},
on peut réécrire ceci sous la forme
$$
Rq'_* Ri_* Li^* Lp^*K|_{X \times_A V} =
R(q' \circ i)_* Lb^*K|_W =
R(U \to V)_* K|_U
$$
ce qui est le résultat voulu.
\end{proof}

\begin{lemma}
\label{lemma-characterize-pseudo-coherent}
Soit $A$ un anneau. Soit $X$ un schéma séparé et
de présentation finie sur $A$. Soit $K \in D_\QCoh(\mathcal{O}_X)$.
Si $R\Gamma(X, E \otimes^\mathbf{L} K)$ est pseudo-cohérent
dans $D(A)$ pour tout objet pseudo-cohérent $E$ de $D(\mathcal{O}_X)$,
alors $K$ est pseudo-cohérent relativement à $A$.
\end{lemma}

\begin{proof}
Supposons que $K \in D_\QCoh(\mathcal{O}_X)$ et que
$R\Gamma(X, E \otimes^\mathbf{L} K)$ soit pseudo-cohérent
dans $D(A)$ pour tout objet pseudo-cohérent $E$ de $D(\mathcal{O}_X)$.
Soit $x \in X$. Nous allons montrer que $K$ est pseudo-cohérent relativement à $A$
dans un voisinage de $x$, ce qui démontrera le lemme.

\medskip\noindent
Choisissons $U, n, V, Z, z, E$ comme dans le lemme \ref{lemma-derived-chow}.
Notons $p : X \times \mathbf{P}^n \to X$ et
$q : X \times \mathbf{P}^n \to \mathbf{P}^n_A$
les projections.
Alors, pour tout $i \in \mathbf{Z}$, on a
\begin{align*}
& R\Gamma(\mathbf{P}^n_A,
Rq_*(Lp^*K \otimes^\mathbf{L} E)
\otimes^\mathbf{L}
\mathcal{O}_{\mathbf{P}^n_A}(i)) \\
& =
R\Gamma(X \times \mathbf{P}^n,
Lp^*K \otimes^\mathbf{L} E \otimes^\mathbf{L}
Lq^*\mathcal{O}_{\mathbf{P}^n_A}(i)) \\
& =
R\Gamma(X,
K \otimes^\mathbf{L} Rp_*(E \otimes^\mathbf{L}
Lq^*\mathcal{O}_{\mathbf{P}^n_A}(i)))
\end{align*}
d'après
Catégories dérivées de schémas,
lemme \ref{perfect-lemma-cohomology-base-change}.
D'après
Catégories dérivées de schémas,
lemme \ref{perfect-lemma-flat-proper-pseudo-coherent-direct-image-general},
le complexe $Rp_*(E \otimes^\mathbf{L} Lq^*\mathcal{O}_{\mathbf{P}^n_A}(i))$
est pseudo-cohérent sur $X$. L'hypothèse nous dit donc que l'expression
de la formule affichée est un objet pseudo-cohérent de $D(A)$.
D'après
Catégories dérivées de schémas,
lemme \ref{perfect-lemma-pseudo-coherent-on-projective-space},
on conclut que $Rq_*(Lp^*K \otimes^\mathbf{L} E)$ est
pseudo-cohérent sur $\mathbf{P}^n_A$.
D'après le lemme \ref{lemma-compute-Fourier-Mukai-for-derived-chow},
on a
$$
Rq_*(Lp^*K \otimes^\mathbf{L} E)|_{X \times_A V} =
R(U \to V)_*K|_U
$$
Comme $U \to V$ est une immersion fermée dans un sous-schéma ouvert de
$\mathbf{P}^n_A$, cela signifie que $K|_U$ est pseudo-cohérent relativement à $A$
d'après le lemme \ref{lemma-check-relative-pseudo-coherence-on-charts}.
\end{proof}

\begin{lemma}
\label{lemma-characterize-pseudo-coh-improved}
Soit $A$ un anneau. Soit $X$ un schéma séparé et
de présentation finie sur $A$. Soit
$K \in D_\QCoh(\mathcal{O}_X).$ Si
$R \Gamma (X, E \otimes ^{\mathbf{L}} K)$ est
pseudo-cohérent dans $D(A)$ pour tout objet parfait
$E \in D(\mathcal{O}_X)$, alors $K$ est pseudo-cohérent
relativement à $A$.
\end{lemma}

\begin{proof}
Compte tenu du lemme \ref{lemma-characterize-pseudo-coherent}, il suffit
de montrer que $R \Gamma (X, E \otimes ^{\mathbf{L}} K)$ est
pseudo-cohérent dans $D(A)$ pour tout objet pseudo-cohérent
$E \in D(\mathcal{O}_X)$. D'après Catégories dérivées de schémas,
proposition \ref{perfect-proposition-detecting-bounded-above},
il s'ensuit que $K \in D^-_\QCoh (\mathcal{O}_X)$. Le
résultat résulte alors de Catégories dérivées de schémas, lemme
\ref{perfect-lemma-perfect-enough}.
\end{proof}

\begin{lemma}
\label{lemma-characterize-relatively-perfect}
Soit $A$ un anneau. Soit $X$ un schéma séparé, de
présentation finie et plat sur $A$. Soit
$K \in D_\QCoh(\mathcal{O}_X).$ Si
$R \Gamma (X, E \otimes^\mathbf{L} K)$ est parfait dans
$D(A)$ pour tout objet parfait $E \in D(\mathcal{O}_X)$, alors $K$ est
parfait relativement à $\Spec(A)$.
\end{lemma}

\begin{proof}
D'après le lemme \ref{lemma-characterize-pseudo-coh-improved},
$K$ est pseudo-cohérent relativement à $A$. D'après le lemme
\ref{lemma-check-relative-pseudo-coherence-on-charts},
$K$ est pseudo-cohérent dans $D( \mathcal{O}_X)$.
D'après Catégories dérivées de schémas, proposition
\ref{perfect-proposition-detecting-bounded-below},
on voit que $K$ appartient à $D^-(\mathcal{O}_X)$.
Soit $\mathfrak{p}$ un idéal premier de $A$ et notons
$i : Y \to X$ l'inclusion de la fibre schématique
au-dessus de $\mathfrak{p}$, c'est-à-dire que $Y$ est un schéma sur $\kappa(\mathfrak p)$.
D'après Catégories dérivées de schémas, lemme \ref{perfect-lemma-bounded-on-fibres},
il suffira de montrer que $Li^*(K)$ est borné inférieurement.
Soit $G \in D_{perf} (\mathcal{O}_X)$ un complexe parfait
qui engendre $D_\QCoh (\mathcal{O}_X)$ ;
voir Catégories dérivées de schémas, théorème
\ref{perfect-theorem-bondal-van-den-Bergh}.
On a
\begin{align*}
R\Hom _{\mathcal{O}_Y}(Li^*(G), Li^*(K))
& =
R\Gamma(Y, Li^*(G ^\vee \otimes ^\mathbf{L} K)) \\
& =
R\Gamma(X, G^\vee \otimes ^{\mathbf{L}} K)
\otimes^\mathbf{L}_A \kappa(\mathfrak{p})
\end{align*}
La première égalité utilise le fait que $Li^*$ préserve les objets parfaits et les duaux,
ainsi que Cohomologie, lemme \ref{cohomology-lemma-dual-perfect-complex} ; nous omettons
quelques détails. La seconde égalité résulte de
Catégories dérivées de schémas, lemme
\ref{perfect-lemma-compare-base-change},
car $X$ est plat sur $A$. Il résulte de notre hypothèse que ceci est un
objet parfait de $D(\kappa(\mathfrak{p}))$. L'objet
$Li^*(G) \in D_{perf}(\mathcal{O}_Y)$ engendre $D_\QCoh(\mathcal{O}_Y)$ d'après
Catégories dérivées de schémas, remarque \ref{perfect-remark-pullback-generator}.
Par conséquent, Catégories dérivées de schémas, proposition
\ref{perfect-proposition-detecting-bounded-below},
implique maintenant que $Li^*(K)$ est borné inférieurement, ce qui conclut.
\end{proof}




\section{Descente des propriétés de finitude des complexes}
\label{section-descent-finiteness}

\noindent
Cette section fait suite à Catégories dérivées de schémas,
section \ref{perfect-section-descent-finiteness}.

\begin{lemma}
\label{lemma-relative-pseudo-coherent-descends-fppf}
Soit $X \to S$ localement de type fini.
Soit $\{f_i : X_i \to X\}$ un recouvrement fppf de schémas.
Soient $E \in D_\QCoh(\mathcal{O}_X)$ et $m \in \mathbf{Z}$.
Alors $E$ est $m$-pseudo-cohérent relativement à $S$
si et seulement si chaque $Lf_i^*E$ est $m$-pseudo-cohérent relativement à $S$.
\end{lemma}

\begin{proof}
Supposons que $E$ soit $m$-pseudo-cohérent relativement à $S$.
Les morphismes $f_i$ sont pseudo-cohérents d'après le
lemme \ref{lemma-flat-finite-presentation-pseudo-coherent}.
Ainsi, $Lf_i^*E$ est $m$-pseudo-cohérent relativement à $S$
d'après le lemme \ref{lemma-pull-relative-pseudo-coherent}.

\medskip\noindent
Réciproquement, supposons que $Lf_i^*E$ soit $m$-pseudo-cohérent relativement à $S$
pour tout $i$. Choisissons $S = \bigcup U_j$, $W_j \to U_j$,
$W_j = \bigcup W_{j, k}$, $T_{j, k} \to W_{j, k}$, et des
morphismes $\alpha_{j, k} : T_{j, k} \to X_{i(j, k)}$ au-dessus de $S$ comme dans le
lemme \ref{lemma-dominate-fppf}.
Comme le morphisme $T_{j, K} \to S$ est plat et de présentation finie,
on voit que $\alpha_{j, k}$ est pseudo-cohérent d'après le
lemme \ref{lemma-permanence-pseudo-coherent}.
Ainsi,
$$
L\alpha_{j, k}^*Lf_{i(j, k)}^*E = L(T_{i, k} \to S)^*E
$$
est $m$-pseudo-cohérent relativement à $S$ d'après le
lemme \ref{lemma-pull-relative-pseudo-coherent}.
Nous voulons maintenant faire descendre cette propriété
le long des recouvrements $\{T_{j, k} \to W_{j, k}\}$,
$W_j = \bigcup W_{j, k}$, $\{W_j \to U_j\}$ et $S = \bigcup U_j$.
Comme le résultat est vrai pour les recouvrements de Zariski
(par définition de la $m$-pseudo-cohérence relativement à $S$),
cela signifie que nous pouvons supposer qu'il existe un morphisme
fini localement libre surjectif $\pi : Y \to X$
tel que $L\pi^*E$ soit pseudo-cohérent relativement à $S$.
Dans ce cas, $R\pi_*L\pi^*E$ est pseudo-cohérent relativement à $S$
d'après le lemme \ref{lemma-finite-morphism-relative-pseudo-coherence}
(c'est la première fois que nous utilisons le fait que les faisceaux de cohomologie de $E$ sont quasi-cohérents).
On a
$R\pi_*L\pi^*E = E \otimes^\mathbf{L}_{\mathcal{O}_X} \pi_*\mathcal{O}_Y$,
par exemple d'après
Catégories dérivées de schémas, lemme
\ref{perfect-lemma-cohomology-base-change},
et, localement sur $X$, l'application $\mathcal{O}_X \to \pi_*\mathcal{O}_Y$
est l'inclusion d'un facteur direct. On conclut donc d'après le
lemme \ref{lemma-summands-relative-pseudo-coherent}.
\end{proof}

\begin{lemma}
\label{lemma-relative-pseudo-coherent-post-compose}
Soient $X \to T \to S$ des morphismes de schémas. Supposons que $T \to S$
soit plat et localement de présentation finie et que $X \to T$ soit
localement de type fini. Soient $E \in D(\mathcal{O}_X)$ et $m \in \mathbf{Z}$.
Alors $E$ est $m$-pseudo-cohérent relativement à $S$
si et seulement si $E$ est $m$-pseudo-cohérent relativement à $T$.
\end{lemma}

\begin{proof}
Localement sur $X$, on peut choisir une immersion fermée $i : X \to \mathbf{A}^n_T$.
Alors $\mathbf{A}^n_T \to S$ est plat et localement de présentation finie.
Nous pouvons donc
appliquer le lemme \ref{lemma-composition-relative-pseudo-coherent}
pour obtenir l'équivalence.
\end{proof}

\begin{lemma}
\label{lemma-relative-pseudo-coherent-descends-fppf-base}
Soit $f : X \to S$ localement de type fini.
Soit $\{S_i \to S\}$ un recouvrement fppf de schémas.
Notons $f_i : X_i \to S_i$ le changement de base de $f$
et $g_i : X_i \to X$ la projection.
Soient $E \in D_\QCoh(\mathcal{O}_X)$ et $m \in \mathbf{Z}$.
Alors $E$ est $m$-pseudo-cohérent relativement à $S$
si et seulement si chaque $Lg_i^*E$ est $m$-pseudo-cohérent relativement à $S_i$.
\end{lemma}

\begin{proof}
Ceci résulte formellement des
lemmes \ref{lemma-relative-pseudo-coherent-descends-fppf} et
\ref{lemma-relative-pseudo-coherent-post-compose}.
En effet, si $E$ est $m$-pseudo-cohérent relativement à $S$,
alors $Lg_i^*E$ est $m$-pseudo-cohérent relativement à $S$ (d'après le premier lemme),
donc $Lg_i^*E$ est $m$-pseudo-cohérent relativement à $S_i$ (d'après le second).
Réciproquement, si
$Lg_i^*E$ est $m$-pseudo-cohérent relativement à $S_i$, alors
$Lg_i^*E$ est $m$-pseudo-cohérent relativement à $S$ (d'après le second lemme),
donc $E$ est $m$-pseudo-cohérent relativement à $S$ (d'après le premier lemme).
\end{proof}






\section{Objets relativement parfaits}
\label{section-relatively-perfect}

\noindent
Cette section fait suite à la discussion de
Catégories dérivées de schémas, section
\ref{perfect-section-relatively-perfect}.

\begin{lemma}
\label{lemma-thickening-pseudo-coherent}
Soit $i : X \to X'$ un épaississement d'ordre fini de schémas. Soit
$K' \in D(\mathcal{O}_{X'})$ un objet tel que
$K = Li^*K'$ soit pseudo-cohérent. Alors $K'$ est pseudo-cohérent.
\end{lemma}

\begin{proof}
Montrons d'abord que les faisceaux de cohomologie de $K'$ sont quasi-cohérents.
Pour ce faire, on peut se ramener au cas d'un épaississement du premier ordre, voir
section \ref{section-thickenings}. Soit $\mathcal{I} \subset \mathcal{O}_{X'}$
le faisceau quasi-cohérent d'idéaux définissant $X$.
En tensorisant la suite exacte courte
$$
0 \to \mathcal{I} \to \mathcal{O}_{X'} \to i_*\mathcal{O}_X \to 0
$$
par $K'$, on obtient un triangle distingué
$$
K' \otimes_{\mathcal{O}_{X'}}^\mathbf{L} \mathcal{I}
\to K' \to
K' \otimes_{\mathcal{O}_{X'}}^\mathbf{L} i_*\mathcal{O}_X
\to
(K' \otimes_{\mathcal{O}_{X'}}^\mathbf{L} \mathcal{I})[1]
$$
Comme $i_* = Ri_*$ et que l'on peut considérer $\mathcal{I}$
comme un $\mathcal{O}_X$-module quasi-cohérent (puisqu'il s'agit d'un épaississement
du premier ordre), on peut réécrire ceci sous la forme
$$
i_*(K \otimes_{\mathcal{O}_X}^\mathbf{L} \mathcal{I})
\to K' \to
i_*K \to
i_*(K \otimes_{\mathcal{O}_X}^\mathbf{L} \mathcal{I})[1]
$$
On utilise Cohomologie, lemme
\ref{cohomology-lemma-projection-formula-closed-immersion},
pour identifier les termes. Comme $K$ appartient à
$D_\QCoh(\mathcal{O}_X)$, on conclut que
$K'$ appartient à $D_\QCoh(\mathcal{O}_{X'})$ ; ceci utilise
Catégories dérivées de schémas, lemmes
\ref{perfect-lemma-pseudo-coherent},
\ref{perfect-lemma-quasi-coherence-tensor-product} et
\ref{perfect-lemma-quasi-coherence-direct-image}.

\medskip\noindent
Supposons que $K'$ appartienne à $D_\QCoh(\mathcal{O}_{X'})$.
La question est locale sur $X'$ ; on peut donc supposer que $X'$ est affine.
Écrivons $X' = \Spec(A')$ et $X = \Spec(A)$ avec $A = A'/I$ et $I$ nilpotent.
Alors $K'$ provient d'un objet $M' \in D(A')$, voir
Catégories dérivées de schémas, lemme
\ref{perfect-lemma-affine-compare-bounded}.
Ainsi, $M = M' \otimes_{A'}^\mathbf{L} A$ est un objet pseudo-cohérent
de $D(A)$ d'après Catégories dérivées de schémas,
lemme \ref{perfect-lemma-pseudo-coherent-affine}, et notre hypothèse sur $K$.
Il s'ensuit que $M'$ est pseudo-cohérent d'après
Compléments d'algèbre, lemme
\ref{more-algebra-lemma-check-pseudo-coherent-modulo-nilpotent}.
Ceci achève la démonstration.
\end{proof}

\begin{lemma}
\label{lemma-thickening-relatively-perfect}
Considérons un diagramme cartésien
$$
\xymatrix{
X \ar[r]_i \ar[d]_f & X' \ar[d]^{f'} \\
Y \ar[r]^j & Y'
}
$$
de schémas. Supposons que $X' \to Y'$ soit plat et localement
de présentation finie et que $Y \to Y'$ soit un épaississement d'ordre fini.
Soit $E' \in D(\mathcal{O}_{X'})$. Si $E = Li^*(E')$ est parfait relativement à $Y$,
alors $E'$ est parfait relativement à $Y'$.
\end{lemma}

\begin{proof}
Rappelons que la perfection relative à $Y$ de $E$ signifie que $E$ est
pseudo-cohérent et possède localement une dimension de Tor finie comme complexe
de $f^{-1}\mathcal{O}_Y$-modules
(Catégories dérivées de schémas,
définition \ref{perfect-definition-relatively-perfect}).
D'après le lemme \ref{lemma-thickening-pseudo-coherent},
on trouve que $E'$ est pseudo-cohérent.
En particulier, $E'$ appartient à $D_\QCoh(\mathcal{O}_{X'})$, voir
Catégories dérivées de schémas, lemme \ref{perfect-lemma-pseudo-coherent}.
Pour montrer que $E'$ possède localement une dimension de Tor finie,
on peut travailler localement sur $X'$. On peut donc supposer que
$X'$, $S'$, $X$, $S$ sont affines, disons associés aux
anneaux $A'$, $R'$, $A$, $R$.
On se ramène alors à la version d'algèbre commutative grâce à
Catégories dérivées de schémas,
lemme \ref{perfect-lemma-affine-locally-rel-perfect}.
La version d'algèbre commutative est donnée dans Compléments d'algèbre, lemme
\ref{more-algebra-lemma-thickening-relatively-perfect}.
\end{proof}

\begin{lemma}
\label{lemma-henselian-relatively-perfect}
Soit $(R, I)$ une paire formée d'un anneau et d'un idéal $I$
contenu dans le radical de Jacobson. Posons $S = \Spec(R)$ et $S_0 = \Spec(R/I)$.
Soit $f : X \to S$ propre, plat et de présentation finie.
Notons $X_0 = S_0 \times_S X$. Soit $E \in D(\mathcal{O}_X)$
pseudo-cohérent. Si la restriction dérivée $E_0$ de $E$
à $X_0$ est parfaite relativement à $S_0$, alors $E$ est parfait relativement à $S$.
\end{lemma}

\begin{proof}
Choisissons un recouvrement fini par des ouverts affines $X = U_1 \cup \ldots \cup U_n$.
Pour tout $i$, on peut choisir une immersion fermée $U_i \to \mathbf{A}^{d_i}_S$.
Posons $U_{i, 0} = S_0 \times_S U_i$.
Pour tout $i$, le complexe $E_0|_{U_{i, 0}}$ a une amplitude de Tor
contenue dans $[a_i, b_i]$ pour certains $a_i, b_i \in \mathbf{Z}$.
Soit $x \in X$ un point.
Nous allons montrer que l'amplitude de Tor de
$E_x$ sur $R$ est contenue dans $[a_i - d_i, b_i]$ pour un certain $i$.
Ceci achèvera la démonstration, car l'amplitude de Tor peut être
lue sur les germes d'après
Cohomologie, lemme \ref{cohomology-lemma-tor-amplitude-stalk}.

\medskip\noindent
Comme $f$ est propre, $f(\overline{\{x\}})$ est un sous-ensemble fermé de $S$.
Comme $I$ est contenu dans le radical de Jacobson, on voit que
$f(\overline{\{x\}})$ rencontre le sous-ensemble fermé $S_0 \subset S$.
Il existe donc une spécialisation $x \leadsto x_0$ avec $x_0 \in X_0$.
Choisissons $i$ tel que $x_0 \in U_i$ ; ainsi $x_0 \in U_{i, 0}$.
Nous fixerons $i$ pour le reste de la démonstration.
Écrivons $U_i = \Spec(A)$. Alors $A$ est une algèbre plate de présentation
finie sur $R$, qui est un quotient d'une $R$-algèbre polynomiale en
$d_i$ variables. La restriction $E|_{U_i}$ correspond
(d'après Catégories dérivées de schémas, lemmes
\ref{perfect-lemma-affine-compare-bounded} et
\ref{perfect-lemma-pseudo-coherent-affine})
à un objet pseudo-cohérent $K$ de $D(A)$.
Remarquons que $E_0$ correspond à $K \otimes_A^\mathbf{L} A/IA$.
Soient $\mathfrak q \subset \mathfrak q_0 \subset A$ les idéaux
premiers correspondant à $x \leadsto x_0$.
Alors $E_x = K_{\mathfrak q}$ et $K_{\mathfrak q}$
est une localisation de $K_{\mathfrak q_0}$. Il suffit donc
de montrer que $K_{\mathfrak q_0}$ a une amplitude de Tor contenue dans
$[a_i - d_i, b_i]$ comme complexe de $R$-modules.
Soit $I \subset \mathfrak p_0 \subset R$ l'idéal
premier correspondant à $f(x_0)$.
Alors on a
\begin{align*}
K \otimes_R^\mathbf{L} \kappa(\mathfrak p_0)
& =
(K \otimes_R^\mathbf{L} R/I) \otimes_{R/I}^\mathbf{L}
\kappa(\mathfrak p_0) \\
& =
(K \otimes_A^\mathbf{L} A/IA) \otimes_{R/I}^\mathbf{L} \kappa(\mathfrak p_0)
\end{align*}
la seconde égalité valant parce que $R \to A$ est plat.
Par notre choix de $a_i, b_i$, ce complexe n'a de cohomologie
qu'aux degrés de l'intervalle $[a_i, b_i]$.
On peut donc finalement appliquer
Compléments d'algèbre, lemme
\ref{more-algebra-lemma-lift-from-fibre-relatively-perfect},
à $R \to A$, $\mathfrak q_0$, $\mathfrak p_0$ et $K$
pour conclure.
\end{proof}










\section{Contraction de courbes rationnelles}
\label{section-contracting}

\noindent
Dans cette section, nous étudions les morphismes propres $f : X \to Y$ dont les fibres
sont de dimension $\leq 1$ et qui vérifient $R^1f_*\mathcal{O}_X = 0$.
Pour comprendre le titre de cette section, voir
Courbes algébriques, sections
\ref{curves-section-contracting-rational-tails},
\ref{curves-section-contracting-rational-bridges} et
\ref{curves-section-contracting-to-stable}.

\begin{lemma}
\label{lemma-check-h1-fibre-zero}
Soit $f : X \to Y$ un morphisme propre de schémas. Soit $y \in Y$
un point tel que $\dim(X_y) \leq 1$. Si
\begin{enumerate}
\item $R^1f_*\mathcal{O}_X = 0$, ou, plus généralement,
\item il existe un morphisme $g : Y' \to Y$ tel que $y$ appartienne à l'image
de $g$ et tel que $R^1f'_*\mathcal{O}_{X'} = 0$, où $f' : X' \to Y'$
est le changement de base de $f$ par $g$,
\end{enumerate}
alors $H^1(X_y, \mathcal{O}_{X_y}) = 0$.
\end{lemma}

\begin{proof}
Pour démontrer le lemme, on peut remplacer $Y$ par un voisinage ouvert de $y$.
On peut donc supposer que $Y$ est affine
et que toutes les fibres de $f$ sont de dimension $\leq 1$, voir
Morphismes, lemme \ref{morphisms-lemma-openness-bounded-dimension-fibres}.
Dans ce cas, $R^1f_*\mathcal{O}_X$ est un $\mathcal{O}_Y$-module quasi-cohérent
de type fini et sa formation commute à tout changement de base, voir
Limites, lemmes \ref{limits-lemma-proper-top-cohomology-finite-type} et
\ref{limits-lemma-higher-direct-images-zero-above-dimension-fibre}.
Le lemme en résulte immédiatement.
\end{proof}

\begin{lemma}
\label{lemma-h1-fibre-zero}
Soit $f : X \to Y$ un morphisme propre de schémas. Soit $y \in Y$
un point tel que $\dim(X_y) \leq 1$ et $H^1(X_y, \mathcal{O}_{X_y}) = 0$.
Alors il existe un voisinage ouvert $V \subset Y$ de $y$ tel que
$R^1f_*\mathcal{O}_X|_V = 0$
et que la même propriété soit vraie après tout changement de base $Y' \to V$.
\end{lemma}

\begin{proof}
Pour démontrer le lemme, on peut remplacer $Y$ par un voisinage ouvert de $y$.
On peut donc supposer que $Y$ est affine et que toutes les fibres de $f$ sont de
dimension $\leq 1$, voir
Morphismes, lemme \ref{morphisms-lemma-openness-bounded-dimension-fibres}.
Dans ce cas, $R^1f_*\mathcal{O}_X$ est un $\mathcal{O}_Y$-module quasi-cohérent
de type fini et sa formation commute à tout changement de base, voir
Limites, lemmes \ref{limits-lemma-proper-top-cohomology-finite-type} et
\ref{limits-lemma-higher-direct-images-zero-above-dimension-fibre}.
Écrivons $Y = \Spec(A)$ ; supposons que $y$ corresponde à l'idéal premier $\mathfrak p \subset A$ et que
$R^1f_*\mathcal{O}_X$ corresponde au $A$-module fini $M$.
Alors $H^1(X_y, \mathcal{O}_{X_y}) = 0$ signifie que
$\mathfrak pM_\mathfrak p = M_\mathfrak p$ d'après l'énoncé
sur le changement de base. Par le lemme de Nakayama,
on conclut que $M_\mathfrak p = 0$. Comme $M$ est fini, on trouve
un $f \in A$, $f \not \in \mathfrak p$, tel que $M_f = 0$.
En prenant pour $V$ l'ouvert principal $D(f)$, on obtient donc le résultat voulu.
\end{proof}

\begin{lemma}
\label{lemma-globally-generated-vanishing}
Soit $f : X \to Y$ un morphisme propre de schémas tel
que $\dim(X_y) \leq 1$ et $H^1(X_y, \mathcal{O}_{X_y}) = 0$
pour tout $y \in Y$. Soit $\mathcal{F}$ quasi-cohérent sur $X$. Alors
\begin{enumerate}
\item $R^pf_*\mathcal{F} = 0$ pour $p > 1$, et
\item $R^1f_*\mathcal{F} = 0$ s'il existe une surjection
$f^*\mathcal{G} \to \mathcal{F}$ avec $\mathcal{G}$ quasi-cohérent
sur $Y$.
\end{enumerate}
Si $Y$ est affine, on a aussi
\begin{enumerate}
\item[(3)] $H^p(X, \mathcal{F}) = 0$ pour $p \not \in \{0, 1\}$, et
\item[(4)] $H^1(X, \mathcal{F}) = 0$ si $\mathcal{F}$ est engendré par ses sections globales.
\end{enumerate}
\end{lemma}

\begin{proof}
L'annulation de (1) est donnée par Limites, lemme
\ref{limits-lemma-higher-direct-images-zero-above-dimension-fibre}.
Pour démontrer (2), on peut travailler localement sur $Y$ et supposer que $Y$ est affine.
Alors $R^1f_*\mathcal{F}$ est le module quasi-cohérent sur $Y$
associé au module $H^1(X, \mathcal{F})$.
On utilise ici le fait que $Y$ est affine, la quasi-cohérence des images directes
supérieures (Cohomologie des schémas, lemme
\ref{coherent-lemma-quasi-coherence-higher-direct-images}) et
Cohomologie des schémas, lemme
\ref{coherent-lemma-quasi-coherence-higher-direct-images-application}.
Comme $Y$ est affine, le module quasi-cohérent $\mathcal{G}$
est engendré par ses sections globales ; il en est donc de même de $f^*\mathcal{G}$
et de $\mathcal{F}$.
On voit ainsi que (4) implique (2).
L'assertion (3) résulte de (1), ainsi que des remarques qui viennent d'être faites
sur la quasi-cohérence des images directes. Ainsi,
il ne reste qu'à démontrer (4).
Si $\mathcal{F}$ est engendré par ses sections globales, il existe une surjection
$\bigoplus_{i \in I} \mathcal{O}_X \to \mathcal{F}$. D'après l'assertion (1)
et la suite exacte longue de cohomologie, celle-ci
induit une surjection sur $H^1$. Comme $H^1(X, \mathcal{O}_X) = 0$
parce que $R^1f_*\mathcal{O}_X = 0$ d'après le
lemme \ref{lemma-h1-fibre-zero}, et
comme $H^1(X, -)$ commute aux sommes directes
(Cohomologie, lemme \ref{cohomology-lemma-quasi-separated-cohomology-colimit}),
on conclut.
\end{proof}

\begin{lemma}
\label{lemma-h1-fibre-zero-check-h0-kappa}
Soit $f : X \to Y$ un morphisme propre de schémas. Supposons que
\begin{enumerate}
\item pour tout $y \in Y$, on ait $\dim(X_y) \leq 1$ et
$H^1(X_y, \mathcal{O}_{X_y}) = 0$, et
\item $\mathcal{O}_Y \to f_*\mathcal{O}_X$ soit surjective.
\end{enumerate}
Alors $\mathcal{O}_{Y'} \to f'_*\mathcal{O}_{X'}$ est surjective
pour tout changement de base $f' : X' \to Y'$ de $f$.
\end{lemma}

\begin{proof}
On peut supposer que $Y$ et $Y'$ sont affines. On peut alors choisir une immersion
fermée $Y' \to Y''$ avec $Y'' \to Y$ un morphisme plat de schémas affines.
Par changement de base plat
(Cohomologie des schémas, lemme \ref{coherent-lemma-flat-base-change-cohomology}),
on voit que le résultat vaut pour $X'' \to Y''$.
On peut donc supposer que $Y'$ est un sous-schéma fermé de $Y$.
Soit $\mathcal{I} \subset \mathcal{O}_Y$ l'idéal définissant $Y'$.
Il existe alors une suite exacte courte
$$
0 \to \mathcal{I}\mathcal{O}_X \to
\mathcal{O}_X \to \mathcal{O}_{X'} \to 0
$$
où l'on considère $\mathcal{O}_{X'}$ comme un module quasi-cohérent sur $X$.
D'après le lemme \ref{lemma-globally-generated-vanishing},
on a $H^1(X, \mathcal{I}\mathcal{O}_X) = 0$.
Il s'ensuit que
$$
H^0(Y, \mathcal{O}_Y) \to
H^0(Y, f_*\mathcal{O}_X) = H^0(X, \mathcal{O}_X) \to
H^0(X, \mathcal{O}_{X'})
$$
est surjective, comme voulu. La première flèche est surjective puisque $Y$
est affine et que nous avons supposé $\mathcal{O}_Y \to f_*\mathcal{O}_X$
surjective, et la seconde l'est par la suite exacte longue de
cohomologie associée à la suite exacte courte ci-dessus et par
l'annulation qui vient d'être démontrée.
\end{proof}

\begin{lemma}
\label{lemma-h1-fibre-zero-h0-kappa}
Considérons un diagramme commutatif
$$
\xymatrix{
X \ar[rr]_f \ar[rd] & & Y \ar[ld] \\
& S
}
$$
de morphismes de schémas. Soit $s \in S$ un point. Supposons que
\begin{enumerate}
\item $X \to S$ soit localement de présentation finie et plat aux
points de $X_s$,
\item $f$ soit propre,
\item les fibres de $f_s : X_s \to Y_s$ soient de dimension $\leq 1$
et que $R^1f_{s, *}\mathcal{O}_{X_s} = 0$,
\item $\mathcal{O}_{Y_s} \to f_{s, *}\mathcal{O}_{X_s}$ soit surjective.
\end{enumerate}
Alors il existe un ouvert $Y_s \subset V \subset Y$ tel que
(a) $f^{-1}(V)$ soit plat sur $S$,
(b) $\dim(X_y) \leq 1$ pour $y \in V$,
(c) $R^1f_*\mathcal{O}_X|_V = 0$,
(d) $\mathcal{O}_V \to f_*\mathcal{O}_X|_V$
soit surjective,
et tel que (b), (c) et (d) restent vrais après tout changement de base $Y' \to V$.
\end{lemma}

\begin{proof}
Soit $y \in Y$ un point au-dessus de $s$.
Il suffit de trouver un voisinage ouvert de $y$
possédant les propriétés voulues. Dans un premier temps, on remplace $Y$
par l'ouvert $V$ trouvé dans le lemme \ref{lemma-h1-fibre-zero}, de sorte que
$R^1f_*\mathcal{O}_X$ soit universellement nul (l'hypothèse du
lemme est satisfaite d'après le lemme \ref{lemma-check-h1-fibre-zero}).
On rétrécit aussi $Y$ afin que toutes les fibres de $f$ soient de dimension $\leq 1$
(on utilise Morphismes, lemme
\ref{morphisms-lemma-openness-bounded-dimension-fibres},
et la propreté de $f$). On peut donc supposer que (b) et (c) sont satisfaites
avec $V = Y$ et après tout changement de base $Y' \to Y$.
Ainsi, d'après le lemme \ref{lemma-h1-fibre-zero-check-h0-kappa},
il suffit maintenant de montrer (d) sur $Y$.
On peut encore rétrécir $Y$ ; par exemple, on peut supposer, et on suppose,
que $Y$ et $S$ sont affines.

\medskip\noindent
D'après le théorème \ref{theorem-openness-flatness},
il existe un ouvert $U \subset X$ sur lequel $X \to S$
est plat et qui contient $X_s$ par hypothèse.
Alors $f(X \setminus U)$ est un fermé
qui ne contient pas $y$. Ainsi, après avoir rétréci $Y$,
on peut supposer que $X$ est plat sur $S$.

\medskip\noindent
Écrivons $S = \Spec(R)$. Choisissons une immersion fermée $Y \to Y'$, où $Y'$
est le spectre d'un anneau de polynômes $R[x_e; e \in E]$ sur un ensemble $E$.
Notons $f' : X \to Y'$ la composée de $f$ et de $Y \to Y'$.
Alors les hypothèses (1) -- (4), ainsi que (b) et (c),
sont satisfaites pour $f'$ et $s$. Si nous montrons que $\mathcal{O}_{Y'} \to f'_*\mathcal{O}_X$
est surjectif dans un voisinage ouvert de $y$, il en va de même pour
$\mathcal{O}_Y \to f_*\mathcal{O}_X$. Ainsi, on peut supposer que $Y$ est
le spectre de $R[x_e; e \in E]$.

\medskip\noindent
À ce stade, $X$ et $Y$ sont plats sur $S$. Alors $Y_s$ et $X$
sont Tor-indépendants au-dessus de $Y$. Nous invitons le lecteur à en trouver sa propre
démonstration, mais cela résulte aussi du
lemme \ref{lemma-case-of-tor-independence} appliqué au
carré de sommets $X, Y, S, S$ et à son changement de base par $s \to S$.
Par conséquent,
$$
Rf_{s, *}\mathcal{O}_{X_s} = L(Y_s \to Y)^*Rf_*\mathcal{O}_X
$$
d'après le lemme
\ref{perfect-lemma-compare-base-change} de Catégories dérivées de schémas.
Compte tenu de l'annulation déjà établie, ceci implique
$f_{s, *}\mathcal{O}_{X_s} = (Y_s \to Y)^*f_*\mathcal{O}_X$.
On en conclut que $\mathcal{O}_Y \to f_*\mathcal{O}_X$ est un morphisme de
$\mathcal{O}_Y$-modules quasi-cohérents dont l'image inverse
sur $Y_s$ est surjective. Nous affirmons que
$f_*\mathcal{O}_X$ est un $\mathcal{O}_Y$-module de type fini.
Si cela est vrai, alors le conoyau $\mathcal{F}$ de
$\mathcal{O}_Y \to f_*\mathcal{O}_X$
est un $\mathcal{O}_Y$-module quasi-cohérent de type fini
tel que $\mathcal{F}_y \otimes \kappa(y) = 0$.
D'après le lemme de Nakayama (Algèbre, lemme \ref{algebra-lemma-NAK}),
on a $\mathcal{F}_y = 0$. Ainsi, $\mathcal{F}$ est nul
dans un voisinage ouvert de $y$
(Modules, lemme \ref{modules-lemma-finite-type-stalk-zero}),
ce qui achève la démonstration.

\medskip\noindent
Démonstration de l'affirmation.
Pour une partie finie $E' \subset E$, posons $Y' = \Spec(R[x_e; e \in E'])$.
Pour $E'$ assez grand, le morphisme $f' : X \to Y \to Y'$ est propre, voir
Limites, lemme \ref{limits-lemma-finite-type-eventually-proper}.
Dans la suite, on fixe $E'$ et $Y'$.
Écrivons $R = \colim R_i$ comme la colimite de ses sous-algèbres de type fini
sur $\mathbf{Z}$. Posons $S_i = \Spec(R_i)$
et $Y'_i = \Spec(R_i[x_e; e \in E'])$.
Pour $i$ assez grand, on peut trouver un diagramme
$$
\xymatrix{
X \ar[d] \ar[r]_{f'} & Y' \ar[d] \ar[r] & S \ar[d] \\
X_i \ar[r]^{f'_i} & Y'_i \ar[r] & S_i
}
$$
à carrés cartésiens, tel que $X_i$ soit plat sur $S_i$ et que
$X_i \to Y'_i$ soit propre. Voir
Limites, lemmes \ref{limits-lemma-descend-finite-presentation},
\ref{limits-lemma-descend-flat-finite-presentation} et
\ref{limits-lemma-eventually-proper}.
Le même argument que ci-dessus montre que $Y'$ et $X_i$
sont Tor-indépendants au-dessus de $Y'_i$ et donc que
$$
R\Gamma(X, \mathcal{O}_X) =
R\Gamma(X_i, \mathcal{O}_{X_i})
\otimes^\mathbf{L}_{R_i[x_e; e \in E']}
R[x_e; e \in E']
$$
d'après la même référence que ci-dessus. D'après Cohomologie des schémas, lemme
\ref{coherent-lemma-proper-over-affine-cohomology-finite},
le complexe $R\Gamma(X_i, \mathcal{O}_{X_i})$ est pseudo-cohérent
dans la catégorie dérivée de l'anneau noethérien $R_i[x_e; e \in E']$ (voir
Compléments d'algèbre, lemme \ref{more-algebra-lemma-Noetherian-pseudo-coherent}).
Ainsi, $R\Gamma(X, \mathcal{O}_X)$ est pseudo-cohérent dans la catégorie dérivée
de $R[x_e; e \in E']$, voir
Compléments d'algèbre, lemme \ref{more-algebra-lemma-pull-pseudo-coherent}.
Puisque son seul module de cohomologie non nul
est $H^0(X, \mathcal{O}_X)$, on en conclut qu'il s'agit d'un
$R[x_e; e \in E']$-module de type fini, voir
Compléments d'algèbre, lemme \ref{more-algebra-lemma-n-pseudo-module}.
Ceci achève la démonstration.
\end{proof}

\begin{lemma}
\label{lemma-h1-fibre-zero-h0-flat}
Considérons un diagramme commutatif
$$
\xymatrix{
X \ar[rr]_f \ar[rd] & & Y \ar[ld] \\
& S
}
$$
de morphismes de schémas. Supposons que $X \to S$
soit plat, que $f$ soit propre, que $\dim(X_y) \leq 1$ pour $y \in Y$, et que
$R^1f_*\mathcal{O}_X = 0$. Alors $f_*\mathcal{O}_X$
est $S$-plat et la formation de $f_*\mathcal{O}_X$ commute
à tout changement de base $S' \to S$.
\end{lemma}

\begin{proof}
On peut supposer que $Y$ et $S$ sont affines, disons $S = \Spec(A)$.
Pour montrer que le $\mathcal{O}_Y$-module quasi-cohérent
$f_*\mathcal{O}_X$ est plat relativement à $S$, il suffit
de montrer que $H^0(X, \mathcal{O}_X)$ est plat sur $A$
(certains détails sont omis).
D'après le lemme \ref{lemma-globally-generated-vanishing}, on a
$H^1(X, \mathcal{O}_X \otimes_A M) = 0$ pour tout $A$-module $M$.
Puisque $\mathcal{O}_X$ est aussi plat sur $A$, on en déduit que le foncteur
$M \mapsto H^0(X, \mathcal{O}_X \otimes_A M)$ est exact.
De plus, ce foncteur commute aux sommes directes d'après
Cohomologie, lemme \ref{cohomology-lemma-quasi-separated-cohomology-colimit}.
Il est alors facile de vérifier que
$H^0(X, \mathcal{O}_X \otimes_A M) = M \otimes_A H^0(X, \mathcal{O}_X)$
fonctoriellement en $M$, ce qui donne la platitude voulue.
Enfin, si $S' \to S$ est un morphisme de schémas affines donné par
l'homomorphisme d'anneaux $A \to A'$, alors, dans le cas affine qui vient d'être traité,
on voit que
$$
H^0(X \times_S S', \mathcal{O}_{X \times_S S'}) =
H^0(X, \mathcal{O}_X \otimes_A A') = H^0(X, \mathcal{O}_X) \otimes_A A'
$$
Cela montre que la formation de $f_*\mathcal{O}_X$ commute
à tout changement de base $S' \to S$. Certains détails sont omis.
\end{proof}

\begin{lemma}
\label{lemma-h1-fibre-zero-isom}
Considérons un diagramme commutatif
$$
\xymatrix{
X \ar[rr]_f \ar[rd] & & Y \ar[ld] \\
& S
}
$$
de morphismes de schémas. Soit $s \in S$ un point. Supposons que
\begin{enumerate}
\item $X \to S$ soit localement de présentation finie et plat aux
points de $X_s$,
\item $Y \to S$ soit localement de présentation finie,
\item $f$ soit propre,
\item les fibres de $f_s : X_s \to Y_s$ soient de dimension $\leq 1$
et que $R^1f_{s, *}\mathcal{O}_{X_s} = 0$,
\item $\mathcal{O}_{Y_s} \to f_{s, *}\mathcal{O}_{X_s}$ soit un isomorphisme.
\end{enumerate}
Alors il existe un ouvert $Y_s \subset V \subset Y$ tel que
(a) $V$ soit plat sur $S$,
(b) $f^{-1}(V)$ soit plat sur $S$,
(c) $\dim(X_y) \leq 1$ pour $y \in V$,
(d) $R^1f_*\mathcal{O}_X|_V = 0$,
(e) $\mathcal{O}_V \to f_*\mathcal{O}_X|_V$
soit un isomorphisme, et tel que (a) -- (e)
restent vraies après le changement de base de $f^{-1}(V) \to V$ par tout $S' \to S$.
\end{lemma}

\begin{proof}
Soit $y \in Y_s$. On peut toujours remplacer $Y$ par un voisinage ouvert de $y$.
On peut donc supposer que $Y$ et $S$ sont affines. On peut également supposer que
$X$ est plat sur $S$, que $\dim(X_y) \leq 1$ pour $y \in Y$,
que $R^1f_*\mathcal{O}_X = 0$ universellement et que
$\mathcal{O}_Y \to f_*\mathcal{O}_X$ est surjectif, voir
le lemme \ref{lemma-h1-fibre-zero-h0-kappa}. (Nous n'utiliserons pas tout ceci.)

\medskip\noindent
Supposons $S$ et $Y$ affines.
Écrivons $S = \lim S_i$ comme limite cofiltrante de schémas affines noethériens $S_i$.
D'après Limites, lemme \ref{limits-lemma-descend-finite-presentation},
il existe un élément $0 \in I$ et un diagramme
$$
\xymatrix{
X_0 \ar[rr]_{f_0} \ar[rd] & & Y_0 \ar[ld] \\
& S_0
}
$$
de morphismes de type fini entre schémas dont le changement de base à $S$
est le diagramme du lemme. Après avoir augmenté $0$, on peut supposer que $Y_0$
est affine et que $X_0 \to S_0$ est propre, voir
Limites, lemmes \ref{limits-lemma-eventually-proper} et
\ref{limits-lemma-limit-affine}.
Soit $s_0 \in S_0$ l'image de $s$.
Puisque $Y_s$ est affine, on voit que
$R^1f_{s, *}\mathcal{O}_{X_s} = 0$ équivaut
à $H^1(X_s, \mathcal{O}_{X_s}) = 0$.
Puisque $X_s$ est le changement de base de $X_{0, s_0}$ par
l'homomorphisme fidèlement plat $\kappa(s_0) \to \kappa(s)$,
on voit que $H^1(X_{0, s_0}, \mathcal{O}_{X_{0, s_0}}) = 0$
et donc que $R^1f_{0, *}\mathcal{O}_{X_{0, s_0}} = 0$.
De même, puisque $\mathcal{O}_{Y_s} \to f_{s, *}\mathcal{O}_{X_s}$
est un isomorphisme, il en va de même de
$\mathcal{O}_{Y_{0, s_0}} \to f_{0, *}\mathcal{O}_{X_{0, s_0}}$.
Puisque les dimensions des fibres de $X_s \to Y_s$ sont au plus $1$,
il en va de même pour le morphisme $X_{0, s_0} \to Y_{0, s_0}$.
Enfin, puisque $X \to S$ est plat, après avoir augmenté $0$,
on peut supposer que $X_0$ est plat sur $S_0$, voir
Limites, lemme \ref{limits-lemma-descend-flat-finite-presentation}.
Il suffit donc de démontrer le lemme pour
$X_0 \to Y_0 \to S_0$ et le point $s_0$.

\medskip\noindent
En combinant les réductions ci-dessus, on se ramène au cas où
$S$ et $Y$ sont affines, $S$ est noethérien, les fibres de $f$ sont de dimension
$\leq 1$, et $R^1f_*\mathcal{O}_X = 0$ universellement.
Soit $y \in Y_s$ un point. Affirmation :
$$
\mathcal{O}_{Y, y} \longrightarrow (f_*\mathcal{O}_X)_y
$$
est un isomorphisme. L'affirmation implique le lemme. En effet, puisque
$f_*\mathcal{O}_X$ est cohérent (Cohomologie des schémas, proposition
\ref{coherent-proposition-proper-pushforward-coherent}),
l'affirmation signifie que l'on peut remplacer $Y$ par un voisinage ouvert de $y$
et obtenir un isomorphisme $\mathcal{O}_Y \to f_*\mathcal{O}_X$.
On en conclut alors que $Y$ est plat sur $S$
d'après le lemme \ref{lemma-h1-fibre-zero-h0-flat}.
Enfin, l'isomorphisme $\mathcal{O}_Y \to f_*\mathcal{O}_X$
reste un isomorphisme après tout changement de base $S' \to S$
d'après la dernière assertion du lemme \ref{lemma-h1-fibre-zero-h0-flat}.

\medskip\noindent
Démonstration de l'affirmation. On sait déjà que
$\mathcal{O}_{Y, y} \longrightarrow (f_*\mathcal{O}_X)_y$
est surjectif (lemme \ref{lemma-h1-fibre-zero-h0-kappa}),
que $(f_*\mathcal{O}_X)_y$ est
plat sur $\mathcal{O}_{S, s}$ (lemme \ref{lemma-h1-fibre-zero-h0-flat})
et que le morphisme induit
$$
\mathcal{O}_{Y_s, y} =  \mathcal{O}_{Y, y}/\mathfrak m_s\mathcal{O}_{Y, y}
\longrightarrow
(f_*\mathcal{O}_X)_y/\mathfrak m_s (f_*\mathcal{O}_X)_y
\to
(f_{s, *}\mathcal{O}_{X_s})_y
$$
est injectif d'après l'hypothèse du lemme. Il résulte alors du
lemme \ref{algebra-lemma-mod-injective} d'Algèbre
que $\mathcal{O}_{Y, y} \longrightarrow (f_*\mathcal{O}_X)_y$
est injectif, comme souhaité.
\end{proof}

\begin{lemma}
\label{lemma-bijection-on-Pic}
Soit $f : X \to Y$ un morphisme propre de schémas noethériens
tel que $f_*\mathcal{O}_X = \mathcal{O}_Y$, dont
les fibres de $f$ sont de dimension $\leq 1$, et tel que
$H^1(X_y, \mathcal{O}_{X_y}) = 0$ pour $y \in Y$.
Alors $f^* : \Pic(Y) \to \Pic(X)$ est une bijection sur
le sous-groupe des $\mathcal{L} \in \Pic(X)$ tels que
$\mathcal{L}|_{X_y} \cong \mathcal{O}_{X_y}$
pour tout $y \in Y$.
\end{lemma}

\begin{proof}
D'après la formule de projection
(Cohomologie, lemme \ref{cohomology-lemma-projection-formula}),
on voit que $f_*f^*\mathcal{N} \cong \mathcal{N}$ pour
$\mathcal{N} \in \Pic(Y)$.
Nous affirmons que, pour $\mathcal{L} \in \Pic(X)$ tel que
$\mathcal{L}|_{X_y} \cong \mathcal{O}_{X_y}$
pour tout $y \in Y$, le module $\mathcal{N} = f_*\mathcal{L}$
est inversible et que $\mathcal{L} \cong f^*\mathcal{N}$.
Cela achèvera la démonstration.

\medskip\noindent
Le $\mathcal{O}_Y$-module $\mathcal{N} = f_*\mathcal{L}$
est cohérent d'après Cohomologie des schémas, proposition
\ref{coherent-proposition-proper-pushforward-coherent}.
Ainsi, pour voir qu'il s'agit d'un $\mathcal{O}_Y$-module inversible,
il suffit de le vérifier sur les tiges (Algèbre, lemme
\ref{algebra-lemma-finite-projective}).
Puisque l'homomorphisme d'un anneau local noethérien vers son complété
est fidèlement plat, il suffit de vérifier que le complété
$(f_*\mathcal{L})_y^\wedge$ est libre (voir
Algèbre, section \ref{algebra-section-completion-noetherian} et
lemme \ref{algebra-lemma-finite-projective-descends}).
Pour cela, nous utiliserons le théorème des fonctions formelles sous la forme donnée dans
Cohomologie des schémas, lemme \ref{coherent-lemma-formal-functions-stalk}.
Puisque $f_*\mathcal{O}_X = \mathcal{O}_Y$ et donc
$(f_*\mathcal{O}_X)_y^\wedge \cong \mathcal{O}_{Y, y}^\wedge$,
il suffit de montrer que $\mathcal{L}|_{X_n} \cong \mathcal{O}_{X_n}$
pour tout $n$ (de manière compatible lorsque $n$ varie.
D'après le lemme \ref{lemma-picard-group-first-order-thickening},
on a une suite exacte
$$
H^1(X_y, \mathfrak m_y^n\mathcal{O}_X/\mathfrak m_y^{n + 1}\mathcal{O}_X)
\to \Pic(X_{n + 1}) \to \Pic(X_n)
$$
avec les notations du théorème des fonctions formelles.
Observons que l'on a une surjection
$$
\mathcal{O}_{X_y}^{\oplus r_n} \cong
\mathfrak m_y^n/\mathfrak m_y^{n + 1} \otimes_{\kappa(y)} \mathcal{O}_{X_y}
\longrightarrow
\mathfrak m_y^n\mathcal{O}_X/\mathfrak m_y^{n + 1}\mathcal{O}_X
$$
pour certains entiers $r_n \geq 0$.
Puisque $\dim(X_y) \leq 1$, cette surjection induit une surjection
sur les groupes de cohomologie de degré un (par l'annulation de la cohomologie en degrés $\geq 2$
qui résulte de Cohomologie, proposition
\ref{cohomology-proposition-vanishing-Noetherian}).
Le $H^1$ de la suite est donc nul et les morphismes de transition
$\Pic(X_{n + 1}) \to \Pic(X_n)$ sont injectifs, comme souhaité.

\medskip\noindent
Il reste à montrer que $f^*\mathcal{N} \cong \mathcal{L}$.
Cela se démontre par la même méthode et nous omettons les détails.
\end{proof}









\section{Stratifications affines}
\label{section-affine-stratifications}

\noindent
Ce contenu est tiré de \cite{RV}. Nous invitons le lecteur à se familiariser quelque peu
avec les stratifications dans
Topologie, section \ref{topology-section-stratifications},
avant de lire cette section.

\medskip\noindent
Si $X$ est un schéma, une stratification de $X$ désigne en général une
stratification de l'espace topologique sous-jacent à $X$.
Les strates sont des parties localement fermées. Nous considérerons ces
strates comme des sous-schémas localement fermés réduits de $X$ au moyen de
Schémas, remarque \ref{schemes-remark-reduced-induced-locally-closed}.

\begin{definition}
\label{definition-affine-stratification}
Soit $X$ un schéma. Une {\it stratification affine} est une
stratification localement finie $X = \coprod_{i \in I} X_i$
dont les strates $X_i$ sont affines et telle que
les morphismes d'inclusion $X_i \to X$ soient affines.
\end{definition}

\noindent
La condition qu'une stratification $X = \coprod X_i$ soit localement finie est,
en présence de la condition que les morphismes d'inclusion $X_i \to X$
soient quasi-compacts, équivalente à la condition que les strates soient des parties localement
constructibles de $X$, voir Propriétés, lemme
\ref{properties-lemma-stratification-locally-finite-constructible}.

\medskip\noindent
La condition que $X_i \to X$ soit un morphisme affine ne dépend pas
de la structure de schéma que l'on met sur la partie localement fermée $X_i$, voir
le lemme \ref{lemma-thicken-property-morphisms}. De plus,
si $X$ est séparé (ou, plus généralement, est à diagonale affine) et si
$X = \coprod X_i$ est une stratification localement finie à strates affines,
alors les morphismes $X_i \to X$ sont affines. Voir
Morphismes, lemme \ref{morphisms-lemma-affine-permanence}.
Cela permet, dans la plupart des cas qui nous intéressent, de faire abstraction de la condition d'affinité des
morphismes d'inclusion $X_i \to X$.

\medskip\noindent
Nous nous intéressons souvent au cas où l'ensemble partiellement ordonné d'indices $I$
de la stratification est fini. Rappelons que la {\it longueur} d'un
ensemble partiellement ordonné $I$ est le supremum des longueurs $p$ des
chaînes $i_0 < i_1 < \ldots < i_p$ d'éléments de $I$.

\begin{lemma}
\label{lemma-improve-stratification}
Soit $X$ un schéma. Soit $X = \coprod_{i \in I} X_i$ une stratification
affine finie. Il existe une stratification affine
dont l'ensemble d'indices est $\{0, \ldots, n\}$, où $n$ est la longueur de $I$.
\end{lemma}

\begin{proof}
Rappelons que $I$ est muni d'un ordre partiel tel que
l'adhérence de $X_i$ soit contenue dans $\bigcup_{j \leq i} X_j$
pour tout $i \in I$.
Soit $I' \subset I$ l'ensemble des indices maximaux de $I$.
Si $i \in I'$, alors $X_i$ est ouvert dans $X$, car la réunion des
adhérences des autres strates est le complémentaire de $X_i$.
Soit $U = \bigcup_{i \in I'} X_i$, considéré comme un sous-schéma ouvert de $X$,
de sorte que $U_{red} = \coprod_{i \in I'} X_i$ en tant que schémas.
Alors $U$ est un schéma affine d'après
Schémas, lemme \ref{schemes-lemma-disjoint-union-affines},
et le lemme \ref{lemma-thickening-affine-scheme}. Le morphisme $U \to X$
est affine, puisque chacun des $X_i \to X$, $i \in I'$, est affine, par le même raisonnement
utilisant le lemme \ref{lemma-thicken-property-morphisms}.
Le complémentaire $Z = X \setminus U$, muni de la
structure de schéma réduite induite, possède la stratification
affine $Z = \bigcup_{i \in I \setminus I'} X_i$.
On utilise ici le fait qu'un morphisme de schémas $T \to Z$ est affine
si et seulement si la composée $T \to X$ est affine ; cela résulte
de Morphismes, lemmes \ref{morphisms-lemma-closed-immersion-affine},
\ref{morphisms-lemma-composition-affine} et
\ref{morphisms-lemma-affine-permanence}.
Observons que l'ensemble partiellement ordonné $I \setminus I'$
a une longueur exactement égale à celle de $I$ moins un.
Par récurrence, on trouve donc que $Z$ possède une stratification
affine $Z = Z_0 \amalg \ldots \amalg Z_{n - 1}$
dont l'ensemble d'indices est $\{1, \ldots, n\}$. En posant $Z_n = U$,
on obtient la stratification voulue de $X$.
\end{proof}

\noindent
Si un schéma $X$ possède une stratification affine finie, alors
$X$ est bien entendu quasi-compact. Il est un peu moins évident que
cela force également $X$ à être quasi-séparé.

\begin{lemma}
\label{lemma-qc-affine-stratification}
Soit $X$ un schéma. Les assertions suivantes sont équivalentes :
\begin{enumerate}
\item $X$ possède une stratification affine finie, et
\item $X$ est quasi-compact et quasi-séparé.
\end{enumerate}
\end{lemma}

\begin{proof}
Soit $X = \bigcup X_i$ une stratification affine finie. Puisque chaque $X_i$
est affine, donc quasi-compact, on en conclut que $X$ est quasi-compact.
Soient $U, V \subset X$ des ouverts affines. Alors $U \cap X_i$ et $V \cap X_i$
sont des ouverts affines de $X_i$, puisque $X_i \to X$ est un morphisme affine. Ainsi,
$U \cap V \cap X_i$ est un ouvert affine du schéma affine $X_i$ (voir, par exemple,
Schémas, lemme \ref{schemes-lemma-characterize-separated}).
Par conséquent, $U \cap V = \coprod U \cap V \cap X_i$ est quasi-compact,
étant une réunion finie de strates affines. On en conclut que $X$ est
quasi-séparé d'après Schémas, lemme
\ref{schemes-lemma-characterize-quasi-separated}.

\medskip\noindent
Supposons que $X$ soit quasi-compact et quasi-séparé. On peut utiliser le
principe de récurrence de
Cohomologie des schémas, lemme \ref{coherent-lemma-induction-principle},
pour démontrer que $X$ possède une stratification affine finie.
Si $X$ est vide, il possède une stratification affine vide.
Si $X$ est affine non vide, il possède une stratification affine
à une strate. Supposons ensuite que $X = U \cup V$, où $U$ est un ouvert quasi-compact,
$V$ est un ouvert affine, et que l'on dispose de stratifications affines finies
$U = \bigcup_{i \in I} U_i$ et $U \cap V = \coprod_{j \in J} W_j$.
Notons $Z = X \setminus V$ et $Z' = X \setminus U$.
Remarquons que $Z$ est fermé dans $U$ et que $Z'$ est fermé dans $V$.
Observons que $U_i \cap Z$ et $U_i \cap W_j = U_i \times_U W_j$ sont des schémas
affines et affines au-dessus de $U$. (Indications : utiliser le fait que $U_i \times_U W_j \to W_j$ est
affine comme changement de base de $U_i \to U$, donc que $U_i \cap W_j$ est affine,
donc que $U_i \cap W_j \to U_i$ est affine, et donc que $U_i \cap W_j \to U$ est affine.)
Il s'ensuit que
$$
U = \coprod\nolimits_{i \in I} (U_i \cap Z) \amalg
\coprod\nolimits_{(i, j) \in I \times J} (U_i \cap W_j)
$$
est une stratification affine finie munie de l'ordre partiel sur
$I \amalg I \times J$ donné par $i' \leq (i, j) \Leftrightarrow i' \leq i$
et $(i', j') \leq (i, j) \Leftrightarrow i' \leq i$ et $j' \leq j$.
Observons que $(U_i \cap Z) \times_X V = \emptyset$ et que
$(U_i \cap W_j) \times_X V = U_i \cap W_j$ sont affines.
Ainsi, les morphismes $U_i \cap Z \to X$ et $U_i \cap W_j \to X$
sont affines, car on peut vérifier l'affinité d'un morphisme
localement sur le but
(Morphismes, lemme \ref{morphisms-lemma-characterize-affine})
et l'on a l'affinité au-dessus de $U$ et de $V$.
Pour achever la démonstration, on prend la stratification ci-dessus
et on ajoute une strate supplémentaire, à savoir $Z'$,
dont on ajoute l'indice comme élément minimal de l'ensemble partiellement ordonné.
\end{proof}

\begin{definition}
\label{definition-affine-stratification-number}
Soit $X$ un schéma non vide, quasi-compact et quasi-séparé. Le
{\it nombre de stratification affine} est le plus petit entier $n \geq 0$
tel que les conditions équivalentes suivantes soient satisfaites :
\begin{enumerate}
\item il existe une stratification affine finie
$X = \coprod_{i \in I} X_i$ où $I$ est de longueur $n$,
\item il existe une stratification affine
$X = X_0 \amalg X_1 \amalg \ldots \amalg X_n$ dont
l'ensemble d'indices est $\{0, \ldots, n\}$.
\end{enumerate}
\end{definition}

\noindent
L'équivalence des conditions résulte du
lemme \ref{lemma-improve-stratification}.
L'existence d'une stratification affine finie
est démontrée dans le lemme \ref{lemma-qc-affine-stratification}.

\begin{lemma}
\label{lemma-affine-stratification-number-bound}
Soit $X$ un schéma séparé qui possède un recouvrement ouvert par
$n + 1$ schémas affines. Alors le nombre de stratification affine de $X$
est au plus $n$.
\end{lemma}

\begin{proof}
Soit $X = U_0 \cup \ldots \cup U_n$ un recouvrement ouvert affine.
Posons
$$
X_i = (U_i \cup \ldots \cup U_n) \setminus
(U_{i + 1} \cup \ldots \cup U_n)
$$
Alors $X_i$ est affine en tant que sous-schéma fermé de $U_i$. Le morphisme $X_i \to X$
est affine d'après Morphismes, lemme \ref{morphisms-lemma-affine-permanence}.
Enfin, on a $\overline{X_i} \subset X_i \cup X_{i - 1} \cup \ldots X_0$.
\end{proof}

\begin{lemma}
\label{lemma-affine-stratification-number-bound-Noetherian}
Soit $X$ un schéma noethérien de dimension $\infty > d \geq 0$.
Alors le nombre de stratification affine de $X$ est au plus $d$.
\end{lemma}

\begin{proof}
On raisonne par récurrence sur $d$. Si $d = 0$, alors $X$ est affine, voir
Propriétés, lemme \ref{properties-lemma-locally-Noetherian-dimension-0}.
Supposons $d > 0$. Soient $\eta_1, \ldots, \eta_n$ les points génériques
des composantes irréductibles de $X$ (Propriétés, lemme
\ref{properties-lemma-Noetherian-irreducible-components}).
On peut recouvrir $X$ par des ouverts affines contenant $\eta_1, \ldots, \eta_n$, voir
Propriétés, lemme \ref{properties-lemma-point-and-maximal-points-affine}.
Puisque $X$ est quasi-compact, on peut trouver un recouvrement ouvert affine fini
$X = \bigcup_{j = 1, \ldots, m} U_j$ tel que
$\eta_1, \ldots, \eta_n \in U_j$ pour tout $j = 1, \ldots, m$.
Choisissons un ouvert affine $U \subset U_1 \cap \ldots \cap U_m$
contenant $\eta_1, \ldots, \eta_n$ (ce qui est possible d'après le lemme
déjà cité). Alors le morphisme $U \to X$ est affine,
car $U \to U_j$ est affine pour tout $j$, voir
Morphismes, lemme \ref{morphisms-lemma-characterize-affine}.
Soit $Z = X \setminus U$.
Par construction, $\dim(Z) < \dim(X)$.
Par hypothèse de récurrence, on peut trouver une
stratification affine $Z = \bigcup_{i \in \{0, \ldots, n\}} Z_i$
de $Z$ avec $n \leq \dim(Z)$.
En posant $U = X_{n + 1}$ et $X_i = Z_i$ pour $i \leq n$,
on conclut.
\end{proof}

\begin{proposition}
\label{proposition-vanishing-affine-stratification-number}
Soit $X$ un schéma non vide, quasi-compact et quasi-séparé, de
nombre de stratification affine $n$. Alors $H^p(X, \mathcal{F}) = 0$, $p > n$,
pour tout $\mathcal{O}_X$-module quasi-cohérent $\mathcal{F}$.
\end{proposition}

\begin{proof}
Nous allons le démontrer par récurrence sur le nombre de stratification affine $n$.
Si $n = 0$, alors $X$ est affine et le résultat est le
lemme
\ref{coherent-lemma-quasi-coherent-affine-cohomology-zero} de Cohomologie des schémas.
Supposons $n > 0$. D'après la définition \ref{definition-affine-stratification-number},
il existe un schéma affine $U$ et une immersion ouverte affine $j : U \to X$
tels que le complémentaire $Z$ ait pour nombre de stratification affine $n - 1$.
Puisque $U$ et $j$ sont affines, on a $H^p(X, j_*(\mathcal{F}|_U)) = 0$
pour $p > 0$, voir Cohomologie des schémas, lemmes
\ref{coherent-lemma-relative-affine-cohomology} et
\ref{coherent-lemma-relative-affine-vanishing}.
Notons $\mathcal{K}$ et $\mathcal{Q}$ le noyau et le conoyau
du morphisme $\mathcal{F} \to j_*(\mathcal{F}|_U)$. On obtient ainsi
une suite exacte
$$
0 \to \mathcal{K} \to \mathcal{F} \to j_*(\mathcal{F}|_U)
\to \mathcal{Q} \to 0
$$
de $\mathcal{O}_X$-modules quasi-cohérents (voir
Schémas, section \ref{schemes-section-quasi-coherent}).
Un argument standard, consistant à décomposer notre suite exacte en suites exactes
courtes et à utiliser la suite exacte longue de cohomologie,
montre qu'il suffit de démontrer $H^p(X, \mathcal{K}) = 0$
et $H^p(X, \mathcal{Q}) = 0$ pour $p \geq n$.
Puisque $\mathcal{F} \to j_*(\mathcal{F}|_U)$ se restreint
à un isomorphisme sur $U$, on voit que $\mathcal{K}$
et $\mathcal{Q}$ sont supportés sur $Z$.
D'après Propriétés, lemme
\ref{properties-lemma-quasi-coherent-colimit-finite-type},
on peut écrire ces modules comme les colimites filtrantes de
leurs sous-modules quasi-cohérents de type fini.
En utilisant le fait que la cohomologie des faisceaux sur $X$ commute
aux colimites filtrantes, voir
Cohomologie, lemme \ref{cohomology-lemma-quasi-separated-cohomology-colimit},
on conclut qu'il suffit de montrer que, si $\mathcal{G}$
est un module quasi-cohérent de type fini dont le support
est contenu dans $Z$, alors $H^p(X, \mathcal{G}) = 0$ pour $p \geq n$.
Soit $Z' \subset X$ le support schématique de
$\mathcal{G} \oplus \mathcal{O}_Z$ ; on peut considérer, et l'on considère,
$\mathcal{G}$ comme un module quasi-cohérent sur $Z'$, voir
Morphismes, section \ref{morphisms-section-support}.
Alors $Z'$ et $Z$ ont le même espace topologique sous-jacent
et donc le même nombre de stratification affine, à savoir $n - 1$.
Ainsi, $H^p(X, \mathcal{G}) = H^p(Z', \mathcal{G})$
(égalité d'après Cohomologie des schémas, lemme
\ref{coherent-lemma-relative-affine-cohomology})
s'annule pour $p \geq n$ par hypothèse de récurrence.
\end{proof}

\begin{example}
\label{example-affine-stratification-number}
Soit $k$ un corps et soit $X = \mathbf{P}^n_k$ l'espace projectif de dimension $n$
sur $k$. Le lemme \ref{lemma-affine-stratification-number-bound}
s'y applique d'après Constructions, lemme
\ref{constructions-lemma-standard-covering-projective-space}.
Par conséquent, le nombre de stratification affine de $\mathbf{P}^n_k$ est
au plus $n$. D'autre part, on a une cohomologie non nulle
en degré $n$ pour certains modules quasi-cohérents sur
$\mathbf{P}^n_k$, voir
Cohomologie des schémas, lemme
\ref{coherent-lemma-cohomology-projective-space-over-ring}.
En appliquant la proposition \ref{proposition-vanishing-affine-stratification-number},
on conclut que le nombre de stratification affine de
$\mathbf{P}^n_k$ est égal à $n$.
\end{example}








\section{Morphismes universellement ouverts}
\label{section-universally-open}

\noindent
Quelques résultats sur les morphismes universellement ouverts.

\begin{lemma}
\label{lemma-test-universally-open}
Soit $f : X \to S$ un morphisme de schémas.
Les conditions suivantes sont équivalentes :
\begin{enumerate}
\item $f$ est universellement ouvert,
\item pour tout morphisme $S' \to S$ localement de présentation finie,
le changement de base $X_{S'} \to S'$ est ouvert, et
\item pour tout $n$, le morphisme
$\mathbf{A}^n \times X \to \mathbf{A}^n \times S$
est ouvert.
\end{enumerate}
\end{lemma}

\begin{proof}
Il est clair que (1) implique (2) et que (2) implique (3).
Montrons que (3) implique (1).
Supposons que le changement de base $X_T \to T$ ne soit pas ouvert
pour un certain morphisme de schémas $g : T \to S$.
On peut alors trouver des ouverts affines
$V \subset S$, $U \subset X$, $W \subset T$
tels que $f(U) \subset V$ et $g(W) \subset V$
et que $U \times_V W \to W$ ne soit pas ouvert.
Si l'on peut montrer que cela implique que
$\mathbf{A}^n \times U \to \mathbf{A}^n \times V$
n'est pas ouvert, alors $\mathbf{A}^n \times X \to \mathbf{A}^n \times S$
n'est pas ouvert et la démonstration est achevée. On est ainsi ramené
au résultat démontré dans le paragraphe suivant.

\medskip\noindent
Soit $A \to B$ un homomorphisme d'anneaux tel que $A' \to B' = A' \otimes_A B$
n'induise pas une application ouverte entre spectres pour une certaine $A$-algèbre $A'$.
Comme les ouverts principaux forment une base de la topologie de $\Spec(B')$,
on conclut que l'image de $D(g)$ dans $\Spec(A')$
n'est pas ouverte pour un certain $g \in B'$. Écrivons
$g = \sum_{i = 1, \ldots, n} a'_i \otimes b_i$
pour un certain $n$, avec $a'_i \in A'$ et $b_i \in B$.
Considérons l'élément $h = \sum_{i = 1, \ldots, n} x_i b_i$
de $B[x_1, \ldots, x_n]$. Supposons, afin d'obtenir une contradiction, que $D(h)$ ait une image
ouverte par le morphisme
$$
\Spec(B[x_1, \ldots, x_n]) \longrightarrow \Spec(A[x_1, \ldots, x_n])
$$
Alors $D(h)$ s'enverrait surjectivement sur un ouvert quasi-compact
$U \subset \Spec(A[x_1, \ldots, x_n])$.
Soit $A[x_1, \ldots, x_n] \to A'$ l'homomorphisme de $A$-algèbres
qui envoie $x_i$ sur $a'_i$. Il induit également un homomorphisme de $B$-algèbres
$B[x_1, \ldots, x_n] \to B'$ qui envoie $h$ sur $g$.
Puisque
$$
\xymatrix{
\Spec(B[x_1, \ldots, x_n]) \ar[d] &
\Spec(B') \ar[l] \ar[d] \\
\Spec(A[x_1, \ldots, x_n]) &
\Spec(A') \ar[l]
}
$$
est cartésien, l'image de $D(g)$ dans $\Spec(A')$ est égale à
l'image réciproque de $U$ dans $\Spec(A')$, et est donc ouverte, ce qui est
la contradiction recherchée.
\end{proof}

\begin{lemma}
\label{lemma-quasi-finite-Noetherian-universally-open}
Soit $f : X \to Y$ un morphisme de schémas. Si
\begin{enumerate}
\item $f$ est localement quasi-fini,
\item $Y$ est géométriquement unibranche et localement noethérien, et
\item toute composante irréductible de $X$ domine
une composante irréductible de $Y$,
\end{enumerate}
alors $f$ est universellement ouvert.
\end{lemma}

\begin{proof}
Pour tout $n$, le schéma $\mathbf{A}^n \times Y$ est géométriquement unibranche
d'après le lemme \ref{lemma-number-of-branches-and-smooth} et
Propriétés, lemme \ref{properties-lemma-number-of-branches-1}. Par conséquent,
les hypothèses du lemme sont satisfaites par les morphismes
$\mathbf{A}^n \times X \to \mathbf{A}^n \times Y$ pour tout $n$.
D'après le lemme \ref{lemma-test-universally-open},
il suffit de montrer que $f$ est ouvert. D'après Morphismes, lemme
\ref{morphisms-lemma-locally-finite-presentation-universally-open},
il suffit de montrer que les générisations se relèvent le long de $f$.
Supposons que $y' \leadsto y$ soit une spécialisation de points de
$Y$ et que $x \in X$ soit un point s'envoyant sur $y$.
Comme dans le lemme \ref{lemma-etale-makes-quasi-finite-finite-at-point},
choisissons un diagramme
$$
\xymatrix{
u \ar[d] & U \ar[d] \ar[r] & X \ar[d] \\
v & V \ar[r] & Y
}
$$
où $(V, v) \to (Y, y)$ est un voisinage \'etale élémentaire,
$U \to V$ est fini, $u$ est l'unique point de $U$ s'envoyant sur $v$,
$U \subset V \times_Y X$ est ouvert, et $v \mapsto y$ et $u \mapsto x$.
Soit $E$ une composante irréductible de $U$ passant par $u$
(il en existe au moins une). Comme $U \to X$ est \'etale, $E$
s'envoie dans une composante irréductible de $X$,
qui à son tour domine une composante irréductible de $Y$ (par hypothèse).
Comme $U \to V$ est fini, donc fermé, on conclut que
l'image $E' \subset V$ de $E$ est une partie fermée irréductible
passant par $v$ et dominant une composante irréductible de $Y$.
Comme $V \to Y$ est \'etale, $E'$ est nécessairement une composante irréductible
de $V$ passant par $v$.
Puisque $Y$ est géométriquement unibranche, on voit que $E'$ est l'unique composante irréductible
de $V$ passant par $v$ (lemme \ref{lemma-nr-branches}).
Comme $V$ est localement noethérien, on peut, après avoir rétréci $V$,
supposer que $E' = V$ (égalité d'ensembles).

\medskip\noindent
Comme $V \to Y$ est \'etale, on peut trouver une spécialisation
$v' \leadsto v$ dont l'image est $y' \leadsto y$.
D'après ce qui précède, on peut trouver $u' \in U$ s'envoyant sur $v'$.
Alors $u' \leadsto u$, car $u$ est l'unique point de
$U$ s'envoyant sur $v$ et $U \to V$ est fermé.
Enfin, l'image $x' \in X$ de $u'$ est un point
qui se spécialise en $x$ et s'envoie sur $y'$, ce qui achève la démonstration.
\end{proof}

\begin{lemma}
\label{lemma-characterize-universally-open-finite}
Soit $A \to B$ un homomorphisme d'anneaux. Supposons que $B$ soit engendré comme $A$-module par
$b_1, \ldots, b_d \in B$. Posons $h = \sum x_ib_i \in B[x_1, \ldots, x_d]$.
Alors $\Spec(B) \to \Spec(A)$ est universellement ouvert si et seulement si l'image de
$D(h)$ dans $\Spec(A[x_1, \ldots, x_d])$ est ouverte.
\end{lemma}

\begin{proof}
Si $\Spec(B) \to \Spec(A)$ est universellement ouvert, alors, bien entendu,
l'image de $D(h)$ est ouverte. Réciproquement, supposons que l'image $U$ de
$D(h)$ soit ouverte. Soit $A \to A'$ un homomorphisme d'anneaux. Il suffit de montrer
que l'image de tout ouvert principal $D(g) \subset \Spec(A' \otimes_A B)$
dans $\Spec(A')$ est ouverte. On peut écrire
$g = \sum_{i = 1, \ldots, d} a'_i \otimes b_i$ avec certains $a'_i \in A'$.
Soit $A[x_1, \ldots, x_n] \to A'$ l'homomorphisme de $A$-algèbres
qui envoie $x_i$ sur $a'_i$. Il induit également un homomorphisme de $B$-algèbres
$B[x_1, \ldots, x_n] \to A' \otimes_A B$ qui envoie $h$ sur $g$.
Puisque
$$
\xymatrix{
\Spec(B[x_1, \ldots, x_n]) \ar[d] &
\Spec(B') \ar[l] \ar[d] \\
\Spec(A[x_1, \ldots, x_n]) &
\Spec(A') \ar[l]
}
$$
est cartésien, l'image de $D(g)$ dans $\Spec(A')$ est égale à
l'image réciproque de $U$ dans $\Spec(A')$, et est donc ouverte.
\end{proof}

\begin{lemma}
\label{lemma-descend-quasi-finite-universally-open}
Soit $S = \lim S_i$ la limite d'un système cofiltrant de schémas
dont les morphismes de transition sont affines.
Soit $0 \in I$ et soit $f_0 : X_0 \to Y_0$ un morphisme de schémas sur $S_0$.
Supposons $S_0$, $X_0$, $Y_0$ quasi-compacts et quasi-séparés.
Soit $f_i : X_i \to Y_i$ le changement de base de $f_0$ à $S_i$ et
soit $f : X \to Y$ le changement de base de $f_0$ à $S$.
Si
\begin{enumerate}
\item $f$ est localement quasi-fini et universellement ouvert, et
\item $f_0$ est localement de présentation finie,
\end{enumerate}
alors il existe $i \geq 0$ tel que $f_i$ soit localement quasi-fini
et universellement ouvert.
\end{lemma}

\begin{proof}
D'après Limites, lemme \ref{limits-lemma-descend-quasi-finite}, après avoir augmenté
$0$, on peut supposer que $f_0$ est localement quasi-fini. Soit $x \in X$.
Par localisation \'etale des morphismes quasi-finis,
on peut trouver un diagramme
$$
\xymatrix{
X \ar[d] & U \ar[l] \ar[d] \\
Y & V \ar[l]
}
$$
où $V \to Y$ est \'etale, $U \subset X_V$ est ouvert, $U \to V$ est fini,
et $x$ appartient à l'image de $U \to X$, voir
le lemme \ref{lemma-etale-makes-quasi-finite-finite-at-point}.
Après avoir rétréci $V$, on peut supposer que $V$ et $U$ sont affines.
Comme $X$ est quasi-compact, il s'ensuit, en prenant une somme disjointe finie
de tels $V$ et $U$, que l'on peut construire un diagramme comme ci-dessus
tel que $U \to X$ soit surjectif. D'après
Limites, lemmes \ref{limits-lemma-descend-finite-presentation},
\ref{limits-lemma-descend-opens},
\ref{limits-lemma-descend-surjective},
\ref{limits-lemma-descend-finite-finite-presentation},
\ref{limits-lemma-descend-etale}, et
\ref{limits-lemma-limit-affine}, après avoir éventuellement augmenté $0$,
on peut supposer que l'on dispose d'un diagramme
$$
\xymatrix{
X_0 \ar[d] & U_0 \ar[l] \ar[d] \\
Y_0 & V_0 \ar[l]
}
$$
où $V_0$ est affine, $V_0 \to Y_0$ est \'etale,
$U_0 \subset (X_0)_{V_0}$ est ouvert, $U_0 \to V_0$ est fini,
et $U_0 \to X_0$ est surjectif. Comme $V_i \to Y_i$ est \'etale
et donc universellement ouvert, il s'ensuit qu'il suffit
de démontrer que $U_i \to V_i$ est universellement ouvert pour
$i$ assez grand. On est ainsi ramené au cas traité dans le paragraphe
suivant.

\medskip\noindent
Soit $A = \colim A_i$ une colimite filtrante d'anneaux.
Soit $A_0 \to B_0$ un homomorphisme d'anneaux. Posons $B = A \otimes_{A_0} B_0$ et
$B_i = A_i \otimes_{A_0} B_0$. Supposons que $A_0 \to B_0$ soit fini,
de présentation finie, et que $A \to B$ soit universellement ouvert.
Il faut montrer que $A_i \to B_i$ est universellement ouvert pour $i$ assez grand.
Choisissons $b_{0, 1}, \ldots, b_{0, d} \in B_0$ qui engendrent $B_0$
comme $A_0$-module. Posons $h_0 = \sum_{j = 1, \ldots, d} x_jb_{0, j}$
dans $B_0[x_1, \ldots, x_d]$. Notons $h$, resp.\ $h_i$, l'image de $h_0$
dans $B[x_1, \ldots, x_d]$, resp.\ $B_i[x_1, \ldots, x_d]$.
L'image $U$ de $D(h)$ dans $\Spec(A[x_1, \ldots, x_d])$ est ouverte,
car $A \to B$ est universellement ouvert. Bien entendu, $U$ est quasi-compact
puisqu'il est l'image d'un schéma affine. Pour $i$ assez grand, il
existe un ouvert quasi-compact $U_i \subset \Spec(A_i[x_1, \ldots, x_d])$
dont l'image réciproque dans $\Spec(A[x_1, \ldots, x_d])$ est $U$, voir
Limites, lemme \ref{limits-lemma-descend-opens}.
Après avoir augmenté $i$, on peut supposer que $D(h_i)$ s'envoie
dans $U_i$ ; cela résulte du même lemme appliqué à
l'image réciproque de $U_i$ dans $D(h_i)$. Enfin, pour $i$ encore plus grand,
le morphisme de schémas $D(h_i) \to U_i$ sera surjectif
par application du lemme déjà utilisé,
Limites, lemme \ref{limits-lemma-descend-surjective}.
On conclut que $A_i \to B_i$ est universellement ouvert d'après le
lemme \ref{lemma-characterize-universally-open-finite}.
\end{proof}

\begin{lemma}
\label{lemma-count-geometric-fibres}
Soit $f : X \to Y$ un morphisme localement quasi-fini. Alors
\begin{enumerate}
\item les fonctions $n_{X/Y}$ des
lemmes \ref{lemma-base-change-fibres-nr-geometrically-irreducible-components}
et \ref{lemma-base-change-fibres-nr-geometrically-connected-components}
coïncident,
\item si $X$ est quasi-compact, alors $n_{X/Y}$ atteint un maximum $d < \infty$.
\end{enumerate}
\end{lemma}

\begin{proof}
L'égalité des fonctions résulte immédiatement du fait que les
fibres (géométriques) d'un morphisme localement quasi-fini sont discrètes, voir
Morphismes, lemme \ref{morphisms-lemma-locally-quasi-finite-fibres}.
Le caractère borné résulte de
Morphismes, lemmes \ref{morphisms-lemma-characterize-universally-bounded} et
\ref{morphisms-lemma-locally-quasi-finite-qc-source-universally-bounded}.
\end{proof}

\begin{lemma}
\label{lemma-count-geometric-fibres-universally-open}
Soit $f : X \to Y$ un morphisme de schémas séparé, localement quasi-fini et universellement ouvert.
Soit $n_{X/Y}$ comme dans le
lemme \ref{lemma-count-geometric-fibres}.
Si $n_{X/Y}(y) \geq d$ pour un certain $y \in Y$ et $d \geq 0$,
alors $n_{X/Y} \geq d$ dans un voisinage ouvert de $y$.
\end{lemma}

\begin{proof}
La question est locale sur $Y$ ; on peut donc supposer $Y$ affine.
Soit $K$ une clôture algébrique du corps résiduel $\kappa(y)$.
Notre hypothèse est que $(X_y)_K$ possède $\geq d$ composantes connexes.
Alors, pour un ouvert quasi-compact convenable $X' \subset X$,
le schéma $(X'_y)_K$ possède $\geq d$ composantes connexes ; détails omis.
Après avoir remplacé $X$ par $X'$, on peut supposer $X$ quasi-compact.
Alors $f$ est quasi-fini. Soient $x_1, \ldots, x_n$ les points de $X$
au-dessus de $y$. Appliquons le
lemme \ref{lemma-etale-splits-off-quasi-finite-part-technical-variant}
pour obtenir un voisinage \'etale $(U, u) \to (Y, y)$ et une décomposition
$$
U \times_Y X =
W \amalg
\ \coprod\nolimits_{i = 1, \ldots, n}
\ \coprod\nolimits_{j = 1, \ldots, m_i}
V_{i, j}
$$
comme au lemme cité. Observons que $n_{X/Y}(y) = \sum_i m_i$
dans cette situation ; certains détails sont omis.
Comme $f$ est universellement ouvert, on voit
que $V_{i, j} \to U$ est ouvert pour tous $i, j$. Après avoir rétréci
$U$, on peut donc supposer $V_{i, j} \to U$ surjectif
pour tous $i, j$. Cela démontre que
$n_{U \times_Y X/U} \geq \sum_i m_i = n_{X/Y}(y) \geq d$.
Comme la construction de $n_{X/Y}$ est compatible avec le
changement de base, la démonstration est achevée.
\end{proof}

\begin{lemma}
\label{lemma-large-open}
Soit $f : X \to Y$ un morphisme de schémas séparé, localement quasi-fini et universellement ouvert.
Soit $n_{X/Y}$ comme dans le
lemme \ref{lemma-count-geometric-fibres}.
Si $n_{X/Y}$ atteint un maximum $d < \infty$, alors l'ensemble
$$
Y_d = \{y \in Y \mid n_{X/Y}(y) = d\}
$$
est ouvert dans $Y$ et le morphisme $f^{-1}(Y_d) \to Y_d$ est fini.
\end{lemma}

\begin{proof}
L'ouverture de $Y_d$ résulte immédiatement du
lemme \ref{lemma-count-geometric-fibres-universally-open}.
Pour démontrer la finitude au-dessus de $Y_d$, reprenons l'argument de la
démonstration de ce lemme. Soit donc $y \in Y_d$. Il y a alors
au plus $d$ points de $X$ au-dessus de $y$. Notons
$x_1, \ldots, x_n$ les points de $X$
au-dessus de $y$. Appliquons le
lemme \ref{lemma-etale-splits-off-quasi-finite-part-technical-variant}
pour obtenir un voisinage \'etale $(U, u) \to (Y, y)$ et une décomposition
$$
U \times_Y X =
W \amalg
\ \coprod\nolimits_{i = 1, \ldots, n}
\ \coprod\nolimits_{j = 1, \ldots, m_i}
V_{i, j}
$$
comme au lemme cité. Observons que $d = n_{X/Y}(y) = \sum_i m_i$
dans cette situation ; certains détails sont omis.
Comme $f$ est universellement ouvert, on voit
que $V_{i, j} \to U$ est ouvert pour tous $i, j$. Après avoir rétréci
$U$, on peut donc supposer $V_{i, j} \to U$ surjectif
pour tous $i, j$, et l'on peut supposer que $U$ s'envoie dans $W$.
Cela démontre que $n_{U \times_Y X/U} \geq \sum_i m_i = d$.
Comme la construction de $n_{X/Y}$ est compatible avec le
changement de base, on sait que $n_{U \times_Y X/U} = d$. Cela signifie que
$W$ doit être vide, et on conclut que $U \times_Y X \to U$ est fini.
D'après Descente, lemme \ref{descent-lemma-descending-property-finite},
cela implique que $X \to Y$ est fini au-dessus de
l'image du morphisme ouvert $U \to Y$. Autrement dit,
on voit que $f$ est fini au-dessus d'un voisinage ouvert de $y$,
comme souhaité.
\end{proof}






\section{Pondérations}
\label{section-weightings}

\noindent
Le contenu de cette section est tiré de \cite[Exposé XVII, 6.2.4]{SGA4}.

\medskip\noindent
Soit $\pi : U \to V$ un morphisme de schémas localement quasi-fini à
fibres finies. Étant donnée une fonction $w : U \to \mathbf{Z}$, on définit une fonction
$$
\textstyle{\int}_\pi w : V \longrightarrow \mathbf{Z},\quad
v \longmapsto
\sum\nolimits_{u \in U,\ \pi(u) = v} w(u) [\kappa(u) : \kappa(v)]_s
$$
Remarquons que les extensions de corps sont finies
(Morphismes, lemme \ref{morphisms-lemma-residue-field-quasi-finite}),
que $[\kappa' : \kappa]_s$ est le degré séparable
(Corps, définition \ref{fields-definition-insep-degree}),
et que la somme est finie puisque les fibres de $\pi$ sont supposées finies.
Une autre manière de calculer la valeur de $\int_\pi w$ en un point $v \in V$
est la suivante.
Choisissons un corps algébriquement clos $k$ et un morphisme
$\overline{v} : \Spec(k) \to V$ dont l'image est $v$. On a alors
$$
(\textstyle{\int}_\pi w)(v) =
\sum\nolimits_{\overline{u} \in U_{\overline{v}}} w(\overline{u})
$$
où, bien entendu, $w(\overline{u})$ désigne la valeur de $w$ en
l'image $u$ du point $\overline{u}$ par le morphisme
$U_{\overline{v}} \to U$.
Remarquons que l'on peut considérer tout $\overline{u} \in U_{\overline{v}}$ comme un morphisme
$\overline{u} : \Spec(k) \to U$ tel que
$\pi \circ \overline{u} = \overline{v}$. En effet, puisque
$U \to V$ est localement quasi-fini à fibres finies,
le schéma $U_{\overline{v}}$ est le spectre
d'une algèbre de dimension finie sur $k$ dont tous les idéaux premiers
sont des idéaux maximaux de corps résiduel $k$. Pour voir que l'égalité
est satisfaite, remarquons que le nombre de morphismes $\overline{u}$ au-dessus
d'un $u$ donné est égal à $[\kappa(u) : \kappa(v)]_s$ d'après
Corps, lemme \ref{fields-lemma-separable-degree}.

\begin{lemma}
\label{lemma-weighting-check-after-etale-base-change}
Étant donnés un carré cartésien
$$
\xymatrix{
U \ar[d]_\pi & U' \ar[l]^h \ar[d]^{\pi'} \\
V & V' \ar[l]_g
}
$$
où $\pi$ est localement quasi-fini à fibres finies,
et une fonction $w : U \to \mathbf{Z}$,
on a $(\int_\pi w) \circ g = \int_{\pi'} (w \circ h)$.
\end{lemma}

\begin{proof}
Cela résulte immédiatement de la seconde description de $\int_\pi w$
donnée ci-dessus. Pour le démontrer à partir de la définition, on utilise le fait que, si
$E/F$ est une extension finie de corps et $F'/F$ une autre extension de corps,
alors, en écrivant $(E \otimes_F F')_{red} = \prod E'_i$ comme un produit de corps
finis sur $F'$, on a
$$
[E : F]_s = \sum [E'_i : F']_s
$$
Pour démontrer cette égalité, choisissons une extension algébriquement close
$\Omega/F'$ et observons que
\begin{align*}
[E : F]_s
& =
|\Mor_F(E, \Omega)| \\
& =
|\Mor_{F'}(E \otimes_F F', \Omega)| \\
& =
|\Mor_{F'}((E \otimes_F F')_{red}, \Omega)| \\
& =
\sum |\Mor_{F'}(E'_i, \Omega)| \\
& =
\sum [E'_i : F']_s
\end{align*}
où l'on a utilisé Corps, lemme \ref{fields-lemma-separable-degree}.
\end{proof}

\begin{definition}
\label{definition-weighting}
Soit $f : X \to Y$ un morphisme localement quasi-fini. Une
{\it fonction de poids} ou une {\it pond\'eration} de $f$ est une application
$w : X \to \mathbf{Z}$ telle que, pour tout diagramme
$$
\xymatrix{
X \ar[d]_f & U \ar[l]^h \ar[d]^\pi \\
Y & V \ar[l]_g
}
$$
où $V \to Y$ est \'etale, $U \subset X_V$ est ouvert et $U \to V$ fini,
la fonction $\int_\pi (w \circ h)$ soit localement constante.
\end{definition}

\noindent
Bien entendu, en prenant $w = 0$, on obtient une pondération de tout morphisme
$f$ localement quasi-fini, bien qu'elle ne soit guère intéressante. Il s'avérera que les
pondérations {\it positives}, c'est-à-dire $w : X \to \mathbf{Z}_{> 0}$, sont
les plus intéressantes à diverses fins.

\begin{lemma}
\label{lemma-weighting-base-change}
Soit $f : X \to Y$ un morphisme localement quasi-fini.
Soit $w : X \to \mathbf{Z}$ une pondération. Soit $f' : X' \to Y'$
le changement de base de $f$ par un morphisme $Y' \to Y$. Alors la
composée $w' : X' \to \mathbf{Z}$ de $w$ et de la projection $X' \to X$
est une pondération de $f'$.
\end{lemma}

\begin{proof}
Considérons un diagramme
$$
\xymatrix{
X' \ar[d]_{f'} & U' \ar[l]^{h'} \ar[d]^{\pi'} \\
Y' & V' \ar[l]_{g'}
}
$$
comme dans la définition \ref{definition-weighting} pour le morphisme $f'$.
Pour tout $v' \in V'$, il faut montrer que $\int_{\pi'} (w' \circ h')$ est
constante dans un voisinage ouvert de $v'$. D'après le
lemme \ref{lemma-weighting-check-after-etale-base-change}
(et le fait que les morphismes \'etales sont ouverts),
on peut remplacer $V'$ par n'importe quel voisinage \'etale de $v'$.
Après avoir remplacé $V'$ par un voisinage \'etale de $v'$,
on peut supposer que $U' = U'_1 \amalg \ldots \amalg U'_n$,
où chaque $U'_i$ possède un unique point $u'_i$ au-dessus de $v'$,
tel que $\kappa(u'_i)/\kappa(v')$ soit purement inséparable,
voir le lemme \ref{lemma-etale-splits-off-quasi-finite-part-technical-variant}.
Il suffit clairement de démontrer que $\int_{U'_i \to V'} w'|_{U'_i}$
est constante dans un voisinage de $v'$.
On est ainsi ramené au cas traité dans le paragraphe suivant.

\medskip\noindent
On a $v' \in V'$ et il existe un unique point $u'$ de $U'$
au-dessus de $v'$ tel que $\kappa(u')/\kappa(v')$ soit purement inséparable.
Notons $x \in X$ et $y \in Y$ les images de $u'$ et de $v'$.
On peut trouver un voisinage \'etale $(V, v) \to (Y, y)$
et un ouvert $U \subset X_V$ tels que $\pi : U \to V$ soit fini
et qu'il existe un unique point $u \in U$ au-dessus de $v$
qui s'envoie sur $x \in X$ par la projection $h : U \to X$,
et tel que, de plus, $\kappa(u)/\kappa(v)$ soit
purement inséparable. Cela est possible d'après le lemme utilisé ci-dessus.
Considérons le morphisme
$$
U'' = U \times_X U' \longrightarrow V \times_Y V' = V''
$$
Comme $u$ et $u'$ s'envoient tous deux sur $x \in X$, il existe un point
$u'' \in U''$ s'envoyant sur $(u, u')$. Notons $v'' \in V''$
l'image de $u''$. Après avoir remplacé $V', v'$ par $V'', v''$,
on peut supposer que la composée $V' \to Y' \to Y$ se factorise
par un morphisme de voisinages \'etales $(V', v') \to (V, v)$
tel que le morphisme induit $X'_{V'} = X_{V'} \to X_V$ envoie
$u'$ sur $u$. Dans le changement de base $X'_{V'} = X_{V'}$, on a
deux sous-schémas ouverts, à savoir $U'$ et l'image réciproque $U_{V'}$ de
$U \subset X_V$. Par construction, tous deux contiennent un unique point
au-dessus de $v'$, à savoir $u'$ dans les deux cas.
Ainsi, après avoir rétréci $V'$, on peut supposer que ces ouverts
sont égaux ; en effet, $U' \setminus (U' \cap U_{V'})$ et
$U_{V'} \setminus (U' \cap U_{V'})$ ont
une image fermée dans $V'$, et ces images ne contiennent pas $v'$.
Ainsi $U' = U_{V'}$, et l'on obtient un diagramme cartésien comme dans le
lemme \ref{lemma-weighting-check-after-etale-base-change}.
Comme $\int_\pi (w \circ h)$ est localement constante
par hypothèse, on conclut.
\end{proof}

\begin{lemma}
\label{lemma-weighting-open}
Soit $f : X \to Y$ un morphisme localement quasi-fini. Soit
$w : X \to \mathbf{Z}$ une pondération de $f$. Si $X' \subset X$ est ouvert,
alors $w|_{X'}$ est une pondération de $f|_{X'} : X' \to Y$.
\end{lemma}

\begin{proof}
Cela résulte immédiatement de la définition.
\end{proof}

\begin{lemma}
\label{lemma-weighting-composition}
Soient $f : X \to Y$ et $g : Y \to Z$ des morphismes localement quasi-finis.
Soit $w_f : X \to \mathbf{Z}$ une pondération de $f$ et soit
$w_g : Y \to \mathbf{Z}$ une pondération de $g$. Alors la fonction
$$
X \longrightarrow \mathbf{Z},\quad
x \longmapsto w_f(x) w_g(f(x))
$$
est une pondération de $g \circ f$.
\end{lemma}

\begin{proof}
Posons $w_{g \circ f}(x) = w_f(x) w_g(f(x))$ pour $x \in X$.
Considérons un diagramme
$$
\xymatrix{
X \ar[d]_{g \circ f} & U \ar[l] \ar[d]^\pi \\
Z & W \ar[l]
}
$$
où $W \to Z$ est \'etale, $U \subset X_W$ est ouvert et $U \to W$ fini.
Il faut montrer que $\int_\pi w_{g \circ f}|_U$ est localement constante.
Choisissons un point $w \in W$. D'après le
lemme \ref{lemma-weighting-check-after-etale-base-change}
(et le fait que les morphismes \'etales sont ouverts), il suffit de montrer que
$\int_\pi w_{g \circ f}|_U$ est constante après avoir remplacé
$(W, w)$ par un voisinage \'etale.
Après avoir remplacé $(W, w)$ par un voisinage \'etale, on peut
supposer $U = U_1 \amalg \ldots \amalg U_n$, où chaque $U_i$
possède un unique point $u_i$ au-dessus de $w$, tel que
$\kappa(u_i)/\kappa(w)$ soit purement inséparable, voir le
lemme \ref{lemma-etale-splits-off-quasi-finite-part-technical-variant}.
Il suffit clairement de montrer que $\int_{U_i \to W} w_{g \circ f}|_{U_i}$
est constante dans un voisinage \'etale de $w$.
On est ainsi ramené au cas traité dans le paragraphe suivant.

\medskip\noindent
On a $w \in W$ et il existe un unique point $u \in U$
au-dessus de $w$ tel que $\kappa(u)/\kappa(w)$ soit purement inséparable.
Considérons le point $v = f(u) \in Y$. Après avoir remplacé
$(W, w)$ par un voisinage \'etale élémentaire, on peut
supposer qu'il existe un voisinage ouvert $V \subset Y_W$
de $v$ tel que $V \to W$ soit fini, voir le
lemme \ref{lemma-etale-makes-quasi-finite-finite-at-point}.
Alors $f_W^{-1}(V) \cap U$ est un voisinage ouvert de $u$,
où $f_W : X_W \to Y_W$ est le changement de base de $f$ à $W$.
Après avoir rétréci $W$ au sens de Zariski, on peut donc supposer $f_W(U) \subset V$.
On obtient ainsi des morphismes
$$
U \xrightarrow{a} V \xrightarrow{b} W
$$
et $U \to V$ est fini puisque $V \to W$ est séparé (car fini).
Puisque $w_f$ et $w_g$ sont des pondérations de $f$ et $g$,
nous voyons que $\int_a w_f|_U$ est localement constante sur $V$ et que
$\int_b w_g|_V$ est localement constante sur $W$. Ainsi, quitte à rétrécir
$W$ une fois encore, nous pouvons supposer que ces fonctions sont constantes,
disons de valeurs $n$ et $m$. Il s'ensuit immédiatement que
$\int_\pi w_{g \circ f}|_U = \int_{b \circ a} w_{g \circ f}|_U$
est constante de valeur $nm$, comme souhaité.
\end{proof}

\begin{lemma}
\label{lemma-weighting-universally-open}
Soit $f : X \to Y$ un morphisme localement quasi-fini.
Soit $w : X \to \mathbf{Z}$ une pondération. Si $w(x) > 0$
pour tout $x \in X$, alors $f$ est universellement ouvert.
\end{lemma}

\begin{proof}
Puisque la propriété est préservée par changement de base, voir le
lemme \ref{lemma-weighting-base-change}, il suffit
de montrer que $f$ est ouvert. Puisque nous pouvons aussi remplacer
$X$ par n'importe quel ouvert de $X$, il suffit de montrer que $f(X)$
est ouvert. Soit $y \in f(X)$. Choisissons $x \in X$ tel que $f(x) = y$.
Il suffit de montrer que $f(X)$ contient un voisinage ouvert
de $y$, et il suffit de le faire après avoir remplacé $Y$ par un
voisinage \'etale de $y$. Par localisation \'etale des
morphismes quasi-finis, voir la section \ref{section-etale-localization},
nous pouvons supposer
qu'il existe un voisinage ouvert $U \subset X$ de $x$
tel que $\pi = f|_U : U \to Y$ soit fini. Alors
$\int_\pi w|_U$ est localement constante et prend une valeur positive en $y$.
Par conséquent, $\pi(U)$ contient un voisinage ouvert de $y$,
ce qui achève la démonstration.
\end{proof}

\begin{lemma}
\label{lemma-weighting-flat-quasi-finite}
Soit $f : X \to Y$ un morphisme de schémas. Supposons $f$
localement quasi-fini, localement de présentation finie et plat.
Alors il existe une pondération positive $w : X \to \mathbf{Z}_{> 0}$ de $f$
donnée par la règle qui associe à $x \in X$, situé au-dessus de $y \in Y$,
$$
w(x) =
\text{longueur}_{\mathcal{O}_{X, x}}
(\mathcal{O}_{X, x}/\mathfrak m_y \mathcal{O}_{X, x})
[\kappa(x) : \kappa(y)]_i
$$
où $[\kappa' : \kappa]_i$ est le degré d'inséparabilité
(Corps, définition \ref{fields-definition-insep-degree}).
\end{lemma}

\begin{proof}
Considérons un diagramme comme dans la définition \ref{definition-weighting}.
Soit $u \in U$, d'images $x, y, v$ dans $X, Y, V$. Nous affirmons alors que
$$
\text{longueur}_{\mathcal{O}_{X, x}}
(\mathcal{O}_{X, x}/\mathfrak m_y \mathcal{O}_{X, x}) =
\text{longueur}_{\mathcal{O}_{U, u}}
(\mathcal{O}_{U, u}/\mathfrak m_v \mathcal{O}_{U, u})
$$
et
$$
[\kappa(x) : \kappa(y)]_i =
[\kappa(u) : \kappa(v)]_i
$$
La première égalité découle du fait que $\mathcal{O}_{X, x} \to \mathcal{O}_{U, u}$
est un homomorphisme local plat tel que
$\mathfrak m_y \mathcal{O}_{U, u} = \mathfrak m_v \mathcal{O}_{U, u}$
et $\mathfrak m_x \mathcal{O}_{U, u} = \mathfrak m_u$
(car $\mathcal{O}_{Y, y} \to \mathcal{O}_{V, v}$ et
$\mathcal{O}_{X, x} \to \mathcal{O}_{U, u}$ sont non ramifiés),
d'où l'égalité par Algèbre, lemme \ref{algebra-lemma-pullback-module}.
La seconde égalité découle du fait que $\kappa(v)/\kappa(y)$ est une extension
séparable finie et que $\kappa(u)$ est un facteur de
$\kappa(x) \otimes_{\kappa(y)} \kappa(v)$, si bien que le degré
d'inséparabilité est inchangé. Cela étant dit, nous voyons que la formation
de la fonction du lemme commute aux changements de base \'etales.
Ceci ramène le problème à la discussion du paragraphe suivant.

\medskip\noindent
Supposons que $f$ soit un morphisme fini, plat et de présentation finie.
Nous devons montrer que $\int_f w$ est localement constante sur $Y$.
En fait, $f$ est fini localement libre
(Morphismes, lemme \ref{morphisms-lemma-finite-flat}),
et nous allons montrer que $\int_f w$ est égale au degré de $f$
(qui est une fonction localement constante sur $Y$). En effet,
pour $y \in Y$, nous voyons que
\begin{align*}
(\textstyle{\int}_f w)(y)
& =
\sum\nolimits_{f(x) = y}
\text{longueur}_{\mathcal{O}_{X, x}}
(\mathcal{O}_{X, x}/\mathfrak m_y \mathcal{O}_{X, x})
[\kappa(x) : \kappa(y)]_i
[\kappa(x) : \kappa(y)]_s \\
& =
\sum\nolimits_{f(x) = y}
\text{longueur}_{\mathcal{O}_{X, x}}
(\mathcal{O}_{X, x}/\mathfrak m_y \mathcal{O}_{X, x})
[\kappa(x) : \kappa(y)] \\
& =
\text{longueur}_{\mathcal{O}_{Y, y}}((f_*\mathcal{O}_X)_y/
\mathfrak m_y (f_*\mathcal{O}_X)_y)
\end{align*}
La dernière égalité résulte d'Algèbre, lemme \ref{algebra-lemma-pushdown-module}.
Le dernier nombre est le rang de $f_*\mathcal{O}_X$ en $y$, comme souhaité.
\end{proof}

\begin{lemma}
\label{lemma-weighting-quasi-finite-Noetherian}
Soit $f : X \to Y$ un morphisme de schémas. Supposons que
\begin{enumerate}
\item $f$ soit localement quasi-fini, et
\item $Y$ soit géométriquement unibranche et localement noethérien.
\end{enumerate}
Alors il existe une pondération $w : X \to \mathbf{Z}_{\geq 0}$ donnée par
la règle qui associe à $x \in X$, situé au-dessus de $y \in Y$, le
``degré séparable générique''
de $\mathcal{O}_{X, x}^{sh}$ sur $\mathcal{O}_{Y, y}^{sh}$.
\end{lemma}

\begin{proof}
Il résulte d'Algèbre, lemme
\ref{algebra-lemma-quasi-finite-strict-henselization}
que $\mathcal{O}_{Y, y}^{sh} \to \mathcal{O}_{X, x}^{sh}$
est fini. Puisque $Y$ est géométriquement unibranche,
il existe un unique idéal premier minimal $\mathfrak p$ de
$\mathcal{O}_{Y, y}^{sh}$, voir
Compléments d'algèbre, lemme \ref{more-algebra-lemma-geometrically-unibranch}.
Écrivons
$$
(\kappa(\mathfrak p) \otimes_{\mathcal{O}_{Y, y}^{sh}}
\mathcal{O}_{X, x}^{sh})_{red} =
\prod K_i
$$
comme un produit fini de corps.
Nous posons $w(x) = \sum [K_i : \kappa(\mathfrak p)]_s$.

\medskip\noindent
Puisque cette définition est manifestement insensible à la localisation \'etale,
pour montrer que $w$ est une pondération, nous nous ramenons à montrer que, si
$f$ est un morphisme fini, alors $\int_f w$ est localement constante.
Observons que la valeur de $\int_f w$ en un point générique $\eta$
de $Y$ est simplement le nombre de points de la fibre géométrique
$X_{\overline{\eta}}$ de $X \to Y$ au-dessus de $\eta$. De plus, puisque
$Y$ est unibranche, un point $y$ de $Y$ est la spécialisation d'un unique
point générique $\eta$. Il suffit donc de montrer que $(\int_f w)(y)$
est égal au nombre de points de $X_{\overline{\eta}}$.
Après passage à un voisinage affine de $y$, nous pouvons supposer que
$X \to Y$ est donné par un homomorphisme fini d'anneaux $A \to B$. Supposons que
$\mathcal{O}_{Y, y}^{sh}$ soit construit au moyen d'un homomorphisme
$\kappa(y) \to k$ vers un corps algébriquement clos $k$.
Alors
$$
\mathcal{O}_{Y, y}^{sh} \otimes_A B =
\prod\nolimits_{f(x) = y}
\prod\nolimits_{\varphi \in \Mor_{\kappa(y)}(\kappa(x), k)}
\mathcal{O}_{X, x}^{sh}
$$
par Algèbre, lemme \ref{algebra-lemma-finite-over-henselian}
et le lemme utilisé ci-dessus.
Observons que l'idéal premier minimal $\mathfrak p$ de $\mathcal{O}_{Y, y}^{sh}$
s'envoie sur l'idéal premier de $A$ correspondant à $\eta$. Nous voyons donc que
l'égalité voulue est satisfaite, car le nombre de points d'une fibre
géométrique est inchangé par extension du corps de base.
\end{proof}

\begin{lemma}
\label{lemma-weighting-specialization}
Soit $f : X \to Y$ un morphisme localement quasi-fini de schémas.
Soit $w : X \to \mathbf{Z}$ une pondération de $f$.
Soit $y \leadsto y'$ une spécialisation de points de $Y$.
Supposons $w(x) = 0$ pour tout $x \in X$ tel que $f(x) = y$. Alors
$w(x') = 0$ pour tout $x' \in X$ tel que $f(x') = y'$.
\end{lemma}

\begin{proof}
Soit $x' \in X$ un point s'envoyant sur $y'$. D'après le
lemme \ref{lemma-etale-makes-quasi-finite-finite-at-point},
nous pouvons choisir un diagramme
$$
\xymatrix{
U \ar[r] \ar[dr]_\pi & X_V \ar[d] \ar[r] & X \ar[d]^f \\
& V \ar[r]^\varphi & Y
}
$$
où $V \to Y$ est \'etale, avec un point $v' \in V$ s'envoyant sur $y'$,
avec $U \subset Y_V$ ouvert, tel que $U \to V$ soit fini, et
tel qu'il existe un unique point $u' \in U$ s'envoyant sur $v'$,
qui s'envoie en outre sur $x'$ dans $X$. Notons que, puisque $V \to Y$ est ouvert,
il existe un point $v \in V$ se spécialisant en $v'$ et s'envoyant sur $y$.
Par la définition \ref{definition-weighting}, la fonction
$\int_\pi (w \circ h)$ est localement constante, où $h : U \to X$ est
la composée des flèches horizontales supérieures.
Puisque $\int_\pi (w \circ h)$ est nulle en $v$ d'après
notre hypothèse sur $w$, nous en concluons qu'elle est nulle
en $v'$. Cependant, sa valeur en $v'$ est $w(x')[\kappa(u'):\kappa(v')]_s$,
et nous concluons.
\end{proof}




\section{Compléments sur les pondérations}
\label{section-more-weightings}

\noindent
Nous démontrons quelques propriétés élémentaires supplémentaires des pondérations. Bien
qu'il semble à première vue que les pondérations puissent être très irrégulières, il
s'avère en fait que la condition imposée dans la
définition \ref{definition-weighting} est assez forte.

\begin{lemma}
\label{lemma-jumps-w}
Soit $f : X \to Y$ un morphisme localement quasi-fini.
Soit $w : X \to \mathbf{Z}$ une pondération de $f$. Alors
les ensembles de niveau de la fonction $w$ sont localement constructibles dans $X$.
\end{lemma}

\begin{proof}
Dans la démonstration ci-dessous, nous utiliserons les lemmes \ref{lemma-weighting-open} et
\ref{lemma-weighting-base-change} sans autre mention.
Nous utiliserons également les propriétés élémentaires des parties constructibles
des schémas et des espaces topologiques, voir
Topologie, section \ref{topology-section-constructible} et
Propriétés, section \ref{properties-section-constructible}.
À l'aide de celles-ci, le lecteur voit que la question est locale sur $X$ et $Y$ ;
les détails sont omis. Nous pouvons donc supposer $X$ et $Y$ affines.
Si nous pouvons trouver un morphisme surjectif $Y' \to Y$
de présentation finie tel que les ensembles de niveau de $w$
aient pour images réciproques des parties localement constructibles de $X' = Y' \times_Y X$,
alors nous concluons par Morphismes, théorème \ref{morphisms-theorem-chevalley}.

\medskip\noindent
Supposons $X$ et $Y$ affines. Nous pouvons choisir une immersion $X \to T$
où $T \to Y$ est fini, voir le
lemme \ref{lemma-quasi-finite-separated-pass-through-finite}.
Par Morphismes, lemme \ref{morphisms-lemma-massage-finite},
après avoir remplacé $Y$ par $Y'$, fini localement libre et surjectif sur $Y$,
remplacé $X$ par $Y' \times_Y X$ et $T$ par un schéma fini localement
libre sur $Y'$ contenant $Y' \times_Y T$ comme sous-schéma fermé,
nous pouvons supposer que $T$ est fini localement libre sur $Y$,
et qu'il contient des sous-schémas fermés $T_i$ s'envoyant isomorphiquement sur $Y$
tels que $T = \bigcup_{i = 1, \ldots, n} T_i$ (ensemblistement).
Puisque $T_i \subset T$ est une partie fermée constructible (étant l'image
d'un morphisme de présentation finie $Y \to T$), nous voyons
que, pour $I \subset \{1, \ldots, n\}$, l'intersection
$\bigcap_{i \in I} T_i$ est une partie fermée constructible de $T$
et son image est donc une partie fermée constructible de $Y$.

\medskip\noindent
Pour une décomposition en union disjointe
$\{1, \ldots, n\} = I_1 \amalg \ldots \amalg I_r$ à parties non vides,
considérons la partie $Y_{I_1, \ldots, I_r} \subset Y$
formée des points $y \in Y$ tels que $T_y = \{x_1, \ldots, x_r\}$
soit formé d'exactement $r$ points, avec $x_j \in T_i \Leftrightarrow i \in I_j$.
D'après les remarques précédentes, ceci donne une partition constructible de $Y$.
Il existe un schéma affine $Y'$ de présentation finie
sur $Y$ tel que l'image de $Y' \to Y$ soit exactement
$Y_{I_1, \ldots, I_r}$, voir Algèbre, lemme
\ref{algebra-lemma-constructible-is-image}.
Nous pouvons donc supposer que $Y = Y_{I_1, \ldots, I_r}$ pour une
décomposition en union disjointe
$\{1, \ldots, n\} = I_1 \amalg \ldots \amalg I_r$.
Dans ce cas, l'égalité $T = T(1) \amalg \ldots \amalg T(r)$, où
$T(j) = \bigcap_{i \in I_j} T_i$,
donne une décomposition de $T$ en parties fermées disjointes (et donc ouvertes).
En intersectant avec le sous-schéma localement fermé $X$, nous obtenons une
décomposition analogue $X = X(1) \amalg \ldots \amalg X(r)$ en parties ouvertes et fermées.
Le morphisme $X(j) \to Y$ est une immersion.
Puisque $w$ est une pondération, il s'ensuit que $w|_{X(j)}$
est localement constante\footnote{En fait, si $f : X \to Y$ est une
immersion et si $w$ est une pondération de $f$, alors la restriction de $f$ est une
application ouverte sur le lieu où $w$ est non nulle.}, et nous concluons.
\end{proof}

\begin{lemma}
\label{lemma-weighting-on-dense-set}
Soit $f : X \to Y$ un morphisme localement quasi-fini de schémas.
Soit $w : X \to \mathbf{Z}$ une pondération de $f$.
Soit $E \subset X$ une partie dense.
Supposons $w(x) = 0$ pour tout $x \in E$. Alors $w = 0$.
\end{lemma}

\begin{proof}
Nous conseillons de lire d'abord le lemme \ref{lemma-weighting-specialization}.
Soit $x \in X$, d'image $y \in Y$. Nous allons montrer que $w(x) = 0$. D'après le
lemme \ref{lemma-etale-makes-quasi-finite-finite-at-point},
nous pouvons choisir un diagramme
$$
\xymatrix{
U \ar[r] \ar[dr]_\pi & X_V \ar[d] \ar[r] & X \ar[d]^f \\
& V \ar[r]^\varphi & Y
}
$$
où $V \to Y$ est \'etale, avec un point $v \in V$ s'envoyant sur $y$,
avec $U \subset Y_V$ ouvert, tel que $U \to V$ soit fini, et
tel qu'il existe un unique point $u \in U$ s'envoyant sur $v$,
qui s'envoie en outre sur $x$ dans $X$. Nous pouvons supposer que $V$
(et donc $U$) est affine. Par la définition \ref{definition-weighting},
la fonction $\int_\pi (w \circ h)$ est localement constante, où
$h : U \to X$ est la composée des flèches horizontales supérieures.
L'ensemble $U_0$ des points où $w \circ h$ vaut $0$ est une
partie constructible de $U$ par le lemme \ref{lemma-jumps-w},
et elle est dense dans $U$ puisqu'elle contient $h^{-1}(E)$ (utiliser le fait que $h$ est ouvert).
Par Topologie, lemme \ref{topology-lemma-dense-in-constructible},
nous voyons que $U_0$ contient un ouvert dense,
c'est-à-dire que l'ensemble des points où $w \circ h$ est non nul est contenu
dans une partie fermée nulle part dense $Z \subset U$. Alors $\pi(Z)$
est une partie fermée nulle part dense de $V$ par
Morphismes, lemme \ref{morphisms-lemma-image-nowhere-dense-finite}.
Ainsi, $\int_\pi (w \circ h)$ est nulle sur un ouvert dense, donc
s'annule partout.
Cependant, sa valeur en $v$ est $w(x)[\kappa(u):\kappa(v)]_s$,
et nous concluons.
\end{proof}

\begin{lemma}
\label{lemma-jumps-int-w}
Soit $f : X \to Y$ un morphisme localement quasi-fini de présentation
finie. Soit $w : X \to \mathbf{Z}$ une pondération de $f$. Alors
les ensembles de niveau de la fonction $\int_f w$ sont localement constructibles dans $Y$.
\end{lemma}

\begin{proof}
Par le lemme \ref{lemma-weighting-check-after-etale-base-change},
la formation de la fonction $\int_f w$ commute à tout
changement de base et, par le lemme \ref{lemma-weighting-base-change},
après changement de base, nous avons encore une pondération.
Cela signifie que, si nous pouvons trouver $Y' \to Y$
surjectif et de présentation finie, alors il
suffit de démontrer le résultat après changement de base à $Y'$, voir
Morphismes, théorème \ref{morphisms-theorem-chevalley}.

\medskip\noindent
La question est locale sur $Y$ ; nous pouvons donc supposer $Y$ affine.
Alors $X$ est quasi-compact et quasi-séparé
(puisque $f$ est de présentation finie). Supposons que $X = U \cup V$
soit une réunion d'ouverts quasi-compacts. Alors nous avons
$$
\textstyle{\int}_f w =
\textstyle{\int}_{f|_U} w|_U +
\textstyle{\int}_{f|_V} w|_V -
\textstyle{\int}_{f|_{U \cap V}} w|_{U \cap V}
$$
Ainsi, si nous connaissons le résultat pour $w|_U$, $w|_V$, $w|_{U \cap V}$,
alors nous le connaissons pour $w$. Par le principe d'induction
(Cohomologie des schémas, lemme \ref{coherent-lemma-induction-principle}),
il suffit de démontrer le lemme lorsque $X$ est affine.

\medskip\noindent
Supposons $X$ et $Y$ affines. Nous pouvons choisir une immersion ouverte
$X \to T$ où $T \to Y$ est fini, voir le
lemme \ref{lemma-quasi-finite-separated-pass-through-finite}.
Puisque nous pouvons encore effectuer un changement de base par un $Y' \to Y$ convenable,
nous pouvons utiliser Morphismes, lemme \ref{morphisms-lemma-massage-finite},
pour nous ramener au cas où
toutes les extensions de corps résiduels induites par le morphisme $T \to Y$
(et a fortiori par $X \to Y$) sont triviales.
Dans cette situation, $\int_f w$ consiste simplement à prendre les sommes
des valeurs de $w$ dans les fibres. Les ensembles de niveau de $w$
sont localement constructibles dans $X$ (lemme \ref{lemma-jumps-w}).
La fonction $w$ ne prend qu'un nombre fini de valeurs d'après
Propriétés, lemme
\ref{properties-lemma-stratification-locally-finite-constructible}.
Nous concluons donc par
Morphismes, théorème \ref{morphisms-theorem-chevalley},
et quelques arguments élémentaires sur les sommes d'entiers.
\end{proof}

\begin{lemma}
\label{lemma-semicontinuous-w}
Soit $f : X \to Y$ un morphisme localement quasi-fini.
Soit $w : X \to \mathbf{Z}_{> 0}$ une pondération positive de $f$.
Alors $w$ est semi-continue supérieurement.
\end{lemma}

\begin{proof}
Soit $x \in X$, d'image $y \in Y$. Choisissons un voisinage \'etale
$(V, v) \to (Y, y)$ et un ouvert $U \subset X_V$ tels que
$\pi : U \to V$ soit fini et qu'il existe un unique point $u \in U$
s'envoyant sur $v$, avec $\kappa(u)/\kappa(v)$ purement inséparable.
Voir le lemme \ref{lemma-etale-makes-quasi-finite-finite-multiple-points-var}.
Alors $(\int_\pi w|_U)(v) = w(u)$.
Il résulte de la définition \ref{definition-weighting}
que, quitte à remplacer $V$ par un voisinage de $v$, nous
avons $w|_U(u') \leq w|_U(u) = w(x)$ pour tout $u' \in U$.
En effet, $w|_U(u')$ intervient comme un terme de la somme dans l'expression
de $(\int_\pi w|_U)(\pi(u'))$.
Ceci démontre le lemme, car le morphisme \'etale
$U \to X$ est ouvert.
\end{proof}

\begin{lemma}
\label{lemma-semicontinuous-int-w}
Soit $f : X \to Y$ un morphisme séparé et localement quasi-fini
à fibres finies.
Soit $w : X \to \mathbf{Z}_{> 0}$ une pondération positive de $f$.
Alors $\int_f w$ est semi-continue inférieurement.
\end{lemma}

\begin{proof}
Soit $y \in Y$. Soient $x_1, \ldots, x_r \in X$ les points situés au-dessus de $y$.
Appliquons le
lemme \ref{lemma-etale-splits-off-quasi-finite-part-technical-variant}
pour obtenir un voisinage \'etale $(U, u) \to (Y, y)$ et une décomposition
$$
U \times_Y X =
W \amalg
\ \coprod\nolimits_{i = 1, \ldots, n}
\ \coprod\nolimits_{j = 1, \ldots, m_i}
V_{i, j}
$$
comme dans le passage cité. Observons que $(\int_f w)(y) = \sum w(v_{i, j})$,
où $w(v_{i, j}) = w(x_i)$. Puisque $\int_{V_{i, j} \to U} w|_{V_{i, j}}$
est localement constante par définition, nous pouvons, quitte à rétrécir $U$,
supposer que ces fonctions sont constantes de valeur $w(v_{i, j})$.
Nous en concluons que
$$
\textstyle{\int}_{U \times_Y X \to U} w|_{U \times_Y X} =
\textstyle{\int}_{W \to U} w|_W +
\sum \textstyle{\int}_{V_{i, j} \to U} w|_{V_{i, j}} =
\textstyle{\int}_{W \to U} w|_W + (\int_f w)(y)
$$
Ceci est $\geq (\int_f w)(y)$, et nous concluons puisque $U \to Y$
est ouvert et que la formation de l'intégrale commute aux changements de base
(lemme \ref{lemma-weighting-check-after-etale-base-change}).
\end{proof}

\begin{lemma}
\label{lemma-max-int-w}
Soit $f : X \to Y$ un morphisme localement quasi-fini,
avec $X$ quasi-compact. Soit $w : X \to \mathbf{Z}$ une pondération de $f$.
Alors $\int_f w$ atteint son maximum.
\end{lemma}

\begin{proof}
Il résulte du lemme \ref{lemma-jumps-w} et de
Propriétés, lemme
\ref{properties-lemma-stratification-locally-finite-constructible}
que $w$ ne prend qu'un nombre fini de valeurs sur $X$.
Il résulte de Morphismes, lemme
\ref{morphisms-lemma-locally-quasi-finite-qc-source-universally-bounded}
que $X \to Y$ a des fibres géométriques de cardinal borné.
Ceci montre que $\int_f w$ est bornée.
\end{proof}

\begin{lemma}
\label{lemma-max-int-finite}
Soit $f : X \to Y$ un morphisme séparé et localement quasi-fini.
Soit $w : X \to \mathbf{Z}_{> 0}$ une pondération positive de $f$.
Supposons que $\int_w f$ atteigne son maximum $d$, et soit $Y_d \subset Y$
l'ouvert des points $y$ tels que $(\int_f w)(y) = d$. Alors
le morphisme $f^{-1}(Y_d) \to Y_d$ est fini.
\end{lemma}

\begin{proof}
Observons que $Y_d$ est ouvert par le lemme \ref{lemma-semicontinuous-int-w}.
Soit $y \in Y_d$. Soient $x_1, \ldots, x_n$ les points de $X$
situés au-dessus de $y$. Appliquons le
lemme \ref{lemma-etale-splits-off-quasi-finite-part-technical-variant}
pour obtenir un voisinage \'etale $(U, u) \to (Y, y)$ et une décomposition
$$
U \times_Y X =
W \amalg
\ \coprod\nolimits_{i = 1, \ldots, n}
\ \coprod\nolimits_{j = 1, \ldots, m_i}
V_{i, j}
$$
comme dans le passage cité. Observons que $d = \sum w(v_{i, j})$, où
$w(v_{i, j}) = w(x_i)$. Puisque $\int_{V_{i, j} \to U} w|_{V_{i, j}}$
est localement constante par définition, nous pouvons, quitte à rétrécir $U$,
supposer que ces fonctions sont constantes de valeur $w(v_{i, j})$.
Nous en concluons que
$$
\textstyle{\int}_{U \times_Y X \to U} w|_{U \times_Y X} =
\textstyle{\int}_{W \to U} w|_W +
\sum \textstyle{\int}_{V_{i, j} \to U} w|_{V_{i, j}} =
\textstyle{\int}_{W \to U} w|_W + (\int_f w)(y)
$$
Ceci est $\geq (\int_f w)(y) = d$, et nous concluons que $W$
doit être vide. Ainsi, $U \times_Y X \to U$ est fini.
Par Descente, lemme \ref{descent-lemma-descending-property-finite},
ceci implique que $X \to Y$ est fini au-dessus de
l'image du morphisme ouvert $U \to Y$. Autrement dit,
nous voyons que $f$ est fini au-dessus d'un voisinage ouvert de $y$,
comme souhaité.
\end{proof}

\begin{lemma}
\label{lemma-open-and-closed-in-finite}
Soit $A \to B$ un homomorphisme d'anneaux fini et de présentation finie.
Il existe un homomorphisme d'anneaux de présentation finie $A \to A_{univ}$
et un idempotent $e_{univ} \in B \otimes_A A_{univ}$
tels que, pour tout homomorphisme d'anneaux $A \to A'$ et tout idempotent $e \in B \otimes_A A'$,
il existe un homomorphisme d'anneaux $A_{univ} \to A'$ envoyant $e_{univ}$ sur $e$.
\end{lemma}

\begin{proof}
Choisissons $b_1, \ldots, b_n \in B$ engendrant $B$ comme $A$-module.
Pour chaque $i$, choisissons un polynôme unitaire $P_i \in A[x]$ tel que $P_i(b_i) = 0$
dans $B$, voir Algèbre, lemme \ref{algebra-lemma-finite-is-integral}.
Ainsi, $B$ est un quotient de la $A$-algèbre libre de rang fini
$B' = A[x_1, \ldots, x_n]/(P_1(x_1), \ldots, P_n(x_n))$.
Soit $J \subset B'$ le noyau de la surjection $B' \to B$.
Alors $J =(f_1, \ldots, f_m)$ est engendré par un nombre fini d'éléments, puisque $B$
est une $A$-algèbre de type fini, voir
Algèbre, lemme \ref{algebra-lemma-compose-finite-type}.
Choisissons, sur $A$, une base $b'_1, \ldots, b'_N$ de $B'$.
Considérons l'algèbre
$$
A_{univ} = A[z_1, \ldots, z_N, y_1, \ldots, y_m]/I
$$
où $I$ est l'idéal engendré par les coefficients dans
$A[z_1, \ldots, z_n, y_1, \ldots, y_m]$
de l'expression suivante relativement à la base $b'_1, \ldots, b'_N$ :
$$
(\sum z_j b'_j)^2 - \sum z_j b'_j + \sum y_k f_k
$$
dans $B'[z_1, \ldots, z_N, y_1, \ldots, y_m]$. Par construction,
l'élément $\sum z_j b'_j$ s'envoie sur un idempotent $e_{univ}$ de
l'algèbre $B \otimes_A A_{univ}$. De plus, si $e \in B \otimes_A A'$
est un idempotent, nous pouvons relever $e$ en un élément de la forme
$\sum b'_j \otimes a'_j$ dans $B' \otimes_A A'$, et nous pouvons trouver
des $a''_k \in A'$ tels que
$$
(\sum b'_j \otimes a'_j)^2 - \sum b'_j \otimes a'_j + \sum f_k \otimes a''_k
$$
soit nul dans $B' \otimes_A A'$. Nous obtenons donc un homomorphisme de $A$-algèbres
$A_{univ} \to A$ envoyant $z_j$ sur $a'_j$ et $y_k$ sur $a''_k$,
et $e_{univ}$ sur $e$. Ceci achève la démonstration.
\end{proof}

\begin{lemma}
\label{lemma-open-and-closed-in-quasi-finite}
Soit $X \to Y$ un morphisme de schémas affines quasi-fini et
de présentation finie. Il existe un morphisme $Y_{univ} \to Y$
de présentation finie et un sous-schéma ouvert
$U_{univ} \subset Y_{univ} \times_Y X$ tel que
$U_{univ} \to Y_{univ}$ soit fini et vérifie la propriété suivante :
pour tout morphisme $Y' \to Y$ de schémas affines
et tout sous-schéma ouvert $U' \subset Y' \times_Y X$
tel que $U' \to Y'$ soit fini, il existe un morphisme
$Y' \to Y_{univ}$ tel que l'image réciproque de $U_{univ}$ soit $U'$.
\end{lemma}

\begin{proof}
Rappelons qu'un morphisme de type fini est quasi-fini si et seulement
s'il est de dimension relative $0$, voir
Morphismes, lemme \ref{morphisms-lemma-locally-quasi-finite-rel-dimension-0}.
Par le lemme \ref{lemma-Noetherian-approximation-dimension-d},
appliqué avec $d = 0$, nous nous ramenons au cas où
$X$ et $Y$ sont noethériens. Nous pouvons choisir une immersion ouverte
$X \to X'$ telle que $X' \to Y$ soit fini, voir
Algèbre, lemme \ref{algebra-lemma-quasi-finite-open-integral-closure}.
Notons que, si nous avons $Y' \to Y$ et $U'$ comme en (2), alors
$$
U' \to Y' \times_Y X \to Y' \times_Y X'
$$
est une immersion ouverte entre des schémas finis sur $Y'$, et est donc
aussi fermée. Nous en concluons que $U'$ correspond à un
idempotent dans
$$
\Gamma(Y', \mathcal{O}_{Y'})
\otimes_{\Gamma(Y, \mathcal{O}_Y)}
\Gamma(X', \mathcal{O}_{X'})
$$
dont la partie ouverte et fermée correspondante est contenue dans
l'ouvert $Y' \times_Y X$. Soient $Y'_{univ} \to Y$ et l'idempotent
$$
e'_{univ} \in
\Gamma(Y_{univ}, \mathcal{O}_{Y_{univ}})
\otimes_{\Gamma(Y, \mathcal{O}_Y)}
\Gamma(X', \mathcal{O}_{X'})
$$
le couple construit dans le lemme \ref{lemma-open-and-closed-in-finite}
pour l'homomorphisme d'anneaux $\Gamma(Y, \mathcal{O}_Y) \to \Gamma(X', \mathcal{O}_{X'})$
(ici, nous utilisons le fait que $Y$ est noethérien pour voir que $X'$ est de présentation finie
sur $Y$). Soit $U'_{univ} \subset Y'_{univ} \times_Y X'$ le sous-schéma
ouvert et fermé correspondant. Nous voyons alors que
$$
U'_{univ} \setminus Y'_{univ} \times_Y X
$$
est une partie fermée de $U'_{univ}$, et a donc une image fermée
$T \subset Y'_{univ}$. Si nous posons $Y_{univ} = Y'_{univ} \setminus T$
et si $U_{univ}$ est la restriction de $U'_{univ}$ à
$Y_{univ} \times_Y X$, alors nous voyons que le lemme est vrai.
\end{proof}

\begin{lemma}
\label{lemma-descend-weighting}
Soit $Y = \lim Y_i$ une limite cofiltrante de schémas affines. Soit $0 \in I$,
et soit $f_0 : X_0 \to Y_0$ un morphisme de schémas affines quasi-fini
et de présentation finie. Soient $f : X  \to Y$
et $f_i : X_i \to Y_i$, pour $i \geq 0$, les changements de base de $f_0$.
Si $w : X \to \mathbf{Z}$ est une pondération de $f$, alors, pour tout indice
$i$ assez grand, il existe une pondération $w_i : X_i \to \mathbf{Z}$
de $f_i$ dont le tiré en arrière sur $X$ est $w$.
\end{lemma}

\begin{proof}
D'après le lemme \ref{lemma-jumps-w}, les ensembles de niveau de $w$ sont les sous-ensembles constructibles
$E_k$ de $X$. Il en résulte que la fonction $w$ ne prend qu'un nombre fini
de valeurs, d'après Propriétés, lemme
\ref{properties-lemma-stratification-locally-finite-constructible}.
Il existe donc un $i$ tel que $E_k$ descende en un sous-ensemble
constructible $E_{i, k}$ de $X_i$ pour tout $k$ ; de plus, nous pouvons supposer que
$X_i = \coprod E_{i, k}$. Cela résulte du fait que l'espace topologique
de $X$ est la limite, dans la catégorie des espaces topologiques,
des espaces spectraux $X_i$ le long d'un système dirigé dont les
applications de transition sont spectrales. Voir
Limites, section \ref{limits-section-descent}
et
Topologie, section \ref{topology-section-limits-spectral}.
Nous définissons $w_i : X_i \to \mathbf{Z}$ en demandant que ses ensembles
de niveau soient les ensembles constructibles $E_{i, k}$.

\medskip\noindent
Choisissons $Y_{i, univ} \to Y_i$ et
$U_{i, univ} \subset Y_{i, univ} \times_{Y_i} X_i$
comme dans le lemme \ref{lemma-open-and-closed-in-quasi-finite}.
D'après la propriété universelle de la construction, pour
montrer que $w_i$ est une pondération, il suffirait de montrer
que
$$
\tau_i = \textstyle{\int}_{U_{i, univ} \to Y_{i, univ}} w_i|_{U_{i, univ}}
$$
est localement constante sur $Y_{i, univ}$. D'après le lemme \ref{lemma-jumps-int-w},
cette fonction a des ensembles de niveau constructibles, mais elle peut
ne pas être (encore) localement constante. Posons
$Y_{univ} = Y_{i, univ} \times_{Y_i} Y$
et soit $U_{univ} \subset Y_{univ} \times_Y X$
l'image réciproque de $U_{i, univ}$.
Alors, puisque le tiré en arrière de $w$ sur $Y_{univ} \times_Y X$
est une pondération de $Y_{univ} \times_Y X \to Y_{univ}$
(lemme \ref{lemma-weighting-base-change}),
nous avons bien que
$$
\tau = \textstyle{\int}_{U_{univ} \to Y_{univ}} w_i|_{U_{univ}}
$$
est localement constante sur $Y_{univ}$. Les ensembles de niveau de
$\tau$ sont donc ouverts et fermés. Enfin, nous avons
$Y_{univ} = \lim_{i' \geq i} Y_{i', univ}$,
et les ensembles de niveau de $\tau$ sont les limites projectives des
ensembles de niveau de $\tau_{i'}$ (définis de manière analogue).
Les références ci-dessus impliquent donc que, pour un indice
$i'$ assez grand, les ensembles de niveau de $\tau_{i'}$ sont eux aussi ouverts.
Pour un tel indice $i'$, nous concluons que $w_{i'}$
est une pondération de $f_{i'}$, comme voulu.
\end{proof}







\section{Pondérations et nombres de stratification affine}
\label{section-bounds-asn}

\noindent
Dans cette section, nous donnons une borne du nombre de stratification affine d'un
schéma qui possède un certain type de revêtement par un schéma affine.

\begin{lemma}
\label{lemma-affineness-of-large-open}
Soit $f : X \to Y$ un morphisme de schémas affines qui est
quasi-fini et de présentation finie.
Soit $w : X \to \mathbf{Z}_{> 0}$ une pondération positive de $f$.
Soit $d < \infty$ la valeur maximale de $\int_f w$. L'ouvert
$$
Y_d = \{y \in Y \mid (\textstyle{\int}_f w)(y) = d \}
$$
de $Y$ est affine.
\end{lemma}

\begin{proof}
Observons que $\int_f w$ atteint son maximum d'après le lemme \ref{lemma-max-int-w}.
L'ensemble $Y_d$ est ouvert d'après le lemme \ref{lemma-semicontinuous-int-w}.
L'énoncé du lemme a donc bien un sens.

\medskip\noindent
Réduction au cas noethérien ; veuillez omettre ce paragraphe.
Rappelons qu'un morphisme de type fini est quasi-fini si et seulement
s'il est de dimension relative $0$, voir
Morphismes, lemme \ref{morphisms-lemma-locally-quasi-finite-rel-dimension-0}.
En appliquant le lemme \ref{lemma-Noetherian-approximation-dimension-d}
avec $d = 0$, nous pouvons trouver un morphisme quasi-fini $f_0 : X_0 \to Y_0$
de schémas affines noethériens et un morphisme $Y \to Y_0$ tels que $f$
soit le changement de base de $f_0$. Nous pouvons alors écrire $Y = \lim Y_i$ comme limite
d'un système dirigé de schémas affines de type fini sur $Y_0$, voir
Algèbre, lemme \ref{algebra-lemma-ring-colimit-fp}.
D'après le lemme \ref{lemma-descend-weighting},
nous pouvons trouver un $i$ tel que notre pondération $w$
descende en une pondération $w_i$ du changement de base $f_i : X_i \to Y_i$
de $f_0$. Si le lemme vaut maintenant pour $f_i, w_i$, il implique
alors le lemme pour $f$, puisque la formation de $\int_f w$ commute au
changement de base, voir le lemme \ref{lemma-weighting-check-after-etale-base-change}.

\medskip\noindent
Supposons $X$ et $Y$ noethériens. Soit $X' \to Y'$ le changement de base de $f$
par un morphisme $g : Y' \to Y$. La formation de $\int_f w$, et par suite
celle de l'ouvert $Y_d$, commute au changement de base. Si $g$ est fini et surjectif, alors
$Y'_d \to Y_d$ est fini et surjectif. Dans ce cas, montrer que
$Y_d$ est affine équivaut à montrer que $Y'_d$ est affine, voir
Cohomologie des schémas, lemme
\ref{coherent-lemma-image-affine-finite-morphism-affine-Noetherian}.

\medskip\noindent
Nous pouvons choisir une immersion $X \to T$ avec $T$ fini sur $Y$, voir le lemme
\ref{lemma-quasi-finite-separated-pass-through-finite}.
Nous allons appliquer Morphismes, lemme \ref{morphisms-lemma-massage-finite}
au morphisme fini $T \to Y$. Ce lemme nous dit qu'il
existe un morphisme fini surjectif $Y' \to Y$ tel que
$Y' \times_Y T$ soit un sous-schéma fermé d'un schéma $T'$ fini sur $Y'$
qui possède une forme particulière.
D'après la discussion du premier paragraphe, nous pouvons remplacer $Y$ par $Y'$,
$T$ par $T'$, et $X$ par $Y' \times_Y X$.
Nous pouvons donc supposer qu'il existe une immersion $X \to T$ (qui n'est pas nécessairement
ouverte ou fermée) et des sous-schémas fermés
$T_i \subset T$, $i = 1, \ldots, n$ tels que
\begin{enumerate}
\item $T \to Y$ est fini (et localement libre),
\item $T_i \to Y$ est un isomorphisme, et
\item $T = \bigcup_{i = 1, \ldots, n} T_i$ ensemblistement.
\end{enumerate}
Soit $Y' = \coprod Y_k$ la réunion disjointe des composantes irréductibles
de $Y$ (considérées comme des sous-schémas fermés intègres de $Y$).
Nous pouvons alors effectuer encore une fois le changement de base par $Y' \to Y$ ; nous
utilisons ici le fait que $Y$ est noethérien. Nous pouvons ainsi supposer en
outre que $Y$ est intègre et noethérien.

\medskip\noindent
Nous pouvons et voulons aussi supposer que $T_i \not = T_j$ si $i \not = j$, en
supprimant les répétitions. Puisque $Y$, et donc tous les $T_i$, sont intègres, cela
signifie que, si $T_i$ et $T_j$ se rencontrent, leur intersection est un
sous-ensemble fermé dont l'image est un sous-ensemble fermé propre de $Y$.

\medskip\noindent
Observons que $V_i = X \cap T_i$ est un sous-ensemble localement fermé
qui est en outre un sous-schéma fermé de $X$, donc affine.
Soient $\eta \in Y$ et $\eta_i \in T_i$ les points génériques.
Si $\eta \not \in Y_d$, alors $Y_d = \emptyset$ et nous avons terminé.
Supposons $\eta \in Y_d$. Notons $I \in \{1, \ldots, n\}$
le sous-ensemble des indices $i$ tels que $\eta_i \in V_i$.
Pour $i \in I$, le sous-ensemble localement fermé $V_i \subset T_i$
contient le point générique de l'espace irréductible $T_i$
et est donc ouvert. D'autre part, puisque $f$ est ouvert
(lemme \ref{lemma-weighting-universally-open}),
pour tout $x \in X$, nous pouvons trouver un $i \in I$ et une spécialisation
$\eta_i \leadsto x$. Il s'ensuit que $x \in T_i$, et donc
$x \in V_i$. Autrement dit, nous voyons que $X = \bigcup_{i \in I} V_i$
ensemblistement.
Nous affirmons que $Y_d = \bigcap_{i \in I} \Im(V_i \to Y)$ ; cela
achèvera la preuve, puisque l'intersection des ouverts affines
$\Im(V_i \to Y)$ de $Y$ est affine.

\medskip\noindent
Pour $y \in Y$, écrivons $f^{-1}(\{y\}) = \{x_1, \ldots, x_r\}$ dans $X$.
Pour chaque $i \in I$, il existe au plus un $j(i) \in \{1, \ldots, x_r\}$
tel que $\eta_i \leadsto x_{j(i)}$. En fait, $j(i)$ existe et est
égal à $j$ si et seulement si $x_j \in V_i$. Si $i \in I$ est tel que
$j = j(i)$ soit défini, alors $V_i \to Y$ est un isomorphisme au voisinage
de $x_j \mapsto y$. Par suite, $\bigcup_{i \in I,\ j(i) = j} V_i \to Y$
est fini après remplacement de la source et du but par des voisinages respectifs de
$x_j \mapsto y$. La définition d'une pondération nous donne donc
$w(x_j) = \sum_{i \in I,\ j(i) = j} w(\eta_i)$.
Nous voyons ainsi que
$$
(\textstyle{\int}_f w)(\eta) =
\sum\nolimits_{i \in I} w(\eta_i) \geq
\sum\nolimits_{j(i)\text{ existe}} w(\eta_i) =
\sum\nolimits_j w(x_j) = (\textstyle{\int}_f w)(y)
$$
Il y a donc égalité si et seulement si $y$ appartient à
$\bigcap_{i \in I} \Im(V_i \to Y)$, ce qu'il fallait montrer.
\end{proof}

\begin{proposition}
\label{proposition-asn-weighting}
Soit $f : X \to Y$ un morphisme surjectif quasi-fini de schémas.
Soit $w : X \to \mathbf{Z}_{> 0}$ une pondération positive de $f$.
Supposons $X$ affine et $Y$ séparé\footnote{Il suffit que la
diagonale de $Y$ soit affine.} et non vide. Alors le nombre de stratification
affine de $Y$ est au plus égal au nombre de valeurs distinctes de $\int_f w$
moins $1$.
\end{proposition}

\begin{proof}
Notons que, puisque $Y$ est séparé, le morphisme $X \to Y$ est affine
(Morphismes, lemme \ref{morphisms-lemma-affine-permanence}).
La fonction $\int_f w$ atteint sa valeur maximale $d$ d'après le
lemme \ref{lemma-max-int-w}. Nous allons raisonner par récurrence sur $d$.
Considérons le sous-schéma ouvert $Y_d = \{y \in Y \mid (\int_f w)(y) = d\}$
de $Y$ et rappelons que $f^{-1}(Y_d) \to Y_d$ est fini, voir le
lemme \ref{lemma-max-int-finite}.
D'après le lemme \ref{lemma-affineness-of-large-open},
pour tout ouvert affine $W \subset Y$, l'intersection $Y_d \cap W$ est affine
(on utilise ici que $W \times_Y X$ est affine, puisqu'il est affine sur $X$).
Ainsi $Y_d \to Y$ est un morphisme affine de schémas. Nous
en concluons que $f^{-1}(Y_d) = Y_d \times_Y X$ est
un schéma affine, puisqu'il est affine sur $X$.
Alors $f^{-1}(Y_d) \to Y_d$ est surjectif, et
$Y_d$ est donc affine d'après Limites, lemme \ref{limits-lemma-affine}.
Si $Y_d = Y$, alors le nombre de stratification affine de $Y$ vaut $0$
et nous concluons que le lemme est vrai. Sinon, posons
$Y' = Y \setminus Y_d$, considéré comme un sous-schéma fermé réduit de $Y$.
Posons $X' = Y' \times_Y X$.
Puisque $X'$ est fermé dans $X$, il est affine. Puisque
$Y'$ est fermé dans $Y$, il est séparé.
Le morphisme $f' : X' \to Y'$ est surjectif, et
$w$ induit une pondération $w'$ de $f'$ (voir
le lemme \ref{lemma-weighting-base-change}) tandis que
$\int_{f'} w'$ est la restriction de $\int_f w$ à $Y'$.
Par récurrence, $Y'$ possède une stratification affine dont la
longueur est au plus le nombre ($\leq$) de valeurs distinctes de
$\int_{f'} w'$ moins $1$, ce qui achève la preuve.
\end{proof}









\section{Morphismes complètement décomposés}
\label{section-completely-decomposed}

\noindent
Nishnevich a étudié la notion de famille complètement décomposée
de morphismes \'etales, afin de définir ce que l'on appelle aujourd'hui la
topologie de Nisnevich ; voir par exemple \cite{Nishnevich}.

\begin{definition}
\label{definition-cd-morphism}
Un morphisme $f : X \to Y$ de schémas est dit
{\it complètement décomposé}\footnote{Cette terminologie n'est peut-être pas standard.}
si, pour tout point $y \in Y$, il
existe un point $x \in X$ tel que $f(x) = y$ et que l'extension de corps
$\kappa(x)/\kappa(y)$ soit triviale.
Une famille de morphismes $\{f_i : X_i \to Y\}_{i \in I}$ de
schémas à but fixé est dite {\it complètement décomposée}
si $\coprod f_i : \coprod Y_i \to X$ est complètement décomposé.
\end{definition}

\noindent
Nous commençons par quelques lemmes élémentaires.

\begin{lemma}
\label{lemma-composition-cd}
Le composé de deux morphismes complètement décomposés de schémas
est complètement décomposé.
Si $\{f_i : X_i \to Y\}_{i \in I}$ est complètement décomposée
et si, pour chaque $i$, nous avons une famille $\{X_{ij} \to X_i\}_{j \in J_i}$
qui est complètement décomposée, alors la famille
$\{X_{ij} \to Y\}_{i \in I, j \in J_i}$ est complètement décomposée.
\end{lemma}

\begin{proof}
Omis.
\end{proof}

\begin{lemma}
\label{lemma-base-change-cd}
Tout changement de base d'un morphisme complètement décomposé de schémas
est complètement décomposé.
Si $\{f_i : X_i \to Y\}_{i \in I}$ est complètement décomposée
et si $Y' \to Y$ est un morphisme de schémas, alors
$\{X_i \times_Y Y' \to Y'\}_{i \in I}$ est complètement
décomposée.
\end{lemma}

\begin{proof}
Soient $f : X \to Y$ et $g : Y' \to Y$ des morphismes de schémas.
Soit $y' \in Y'$ un point d'image $y = g(y')$ dans $Y$.
Si $x \in X$ est un point tel que $f(x) = y$ et $\kappa(x) = \kappa(y)$,
alors il existe un unique point $x' \in X' = X \times_Y Y'$
qui s'envoie sur $y'$ dans $Y'$ et sur $x$ dans $X$, et de plus
$\kappa(x') = \kappa(y')$, voir
Schémas, lemme \ref{schemes-lemma-points-fibre-product}.
Le lemme découle aisément de ce fait ; nous omettons les détails.
\end{proof}

\begin{lemma}
\label{lemma-decompose}
\begin{reference}
\cite[Lemme 2.1.2]{EHIK}
\end{reference}
Soit $f : X \to Y$ un morphisme de schémas. Supposons que
$f$ soit complètement décomposé,
$f$ soit localement de présentation finie, et que
$Y$ soit quasi-compact et quasi-séparé.
Alors il existe $n \geq 0$ et des morphismes
$Z_i \to Y$, $i = 1, \ldots, n$, ayant les propriétés suivantes :
\begin{enumerate}
\item $\coprod Z_i \to Y$ est surjectif,
\item $Z_i \to Y$ est une immersion pour tout $i$,
\item $Z_i \to Y$ est de présentation finie pour tout $i$, et
\item le changement de base $X \times_Y Z_i \to Z_i$ possède une section
pour tout $i$.
\end{enumerate}
\end{lemma}

\begin{proof}
Soit $y \in Y$. Par hypothèse, il existe un morphisme
$\sigma : \Spec(\kappa(y)) \to X$ au-dessus de $Y$. Nous pouvons écrire $\Spec(\kappa(y))$
comme la limite d'un système dirigé de schémas affines $Z$ au-dessus de $Y$ tels que
$Z \to Y$ soit une immersion de présentation finie.
En effet, choisissons un ouvert affine $y \in \Spec(A) \subset Y$
et disons que $y$ correspond à l'idéal premier $\mathfrak p$ de $A$.
Alors $\kappa(\mathfrak p)$ est la colimite filtrante des
anneaux $(A/I)_f$, où $I \subset \mathfrak p$ est un idéal de type fini
et $f \in A$, $f \not \in \mathfrak p$.
Les morphismes $Z = \Spec((A/I)_f) \to Y$ sont des immersions de
présentation finie ; la quasi-compacité de $Z \to Y$
résulte du fait que $Y$ est quasi-séparé, voir
Schémas, lemme \ref{schemes-lemma-quasi-compact-permanence}.
D'après Limites, proposition
\ref{limits-proposition-characterize-locally-finite-presentation},
pour l'un de ces $Z$, il existe un morphisme $\sigma' : Z \to X$ au-dessus de $Y$
qui coïncide avec $\sigma$ sur le spectre de $\kappa(\mathfrak p)$.
Comme $\sigma'$ est un morphisme au-dessus de $Y$, nous obtenons une section
de la projection $X \times_Y Z \to Z$.

\medskip\noindent
Nous concluons que $Y$ est la réunion des images d'immersions
$Z \to Y$ de présentation finie telles que $X \times_Y Z \to Z$
possède une section. Puisque l'image de $Z \to Y$ est constructible
(Morphismes, lemme \ref{morphisms-lemma-chevalley})
et puisque $Y$ est compact pour la topologie constructible
(Propriétés, lemme
\ref{properties-lemma-quasi-compact-quasi-separated-spectral} et
Topologie, lemme \ref{topology-lemma-constructible-hausdorff-quasi-compact}),
nous voyons qu'un nombre fini de ces immersions suffit.
\end{proof}

\begin{lemma}
\label{lemma-descend-cd}
Soit $S = \lim_{\lambda \in \Lambda} S_\lambda$
la limite d'un système dirigé de schémas dont les morphismes de transition sont affines.
Soit $0 \in \Lambda$ et soit $f_0 : X_0 \to Y_0$
un morphisme de schémas au-dessus de $S_0$.
Pour $\lambda \geq 0$, soit $f_\lambda : X_\lambda \to Y_\lambda$
le changement de base de $f_0$ à $S_\lambda$, et
soit $f : X \to Y$ le changement de base de $f_0$ à $S$. Si
\begin{enumerate}
\item $f$ est complètement décomposé,
\item $Y_0$ est quasi-compact et quasi-séparé, et
\item $f_0$ est localement de présentation finie,
\end{enumerate}
alors il existe un $\lambda \geq 0$ tel que $f_\lambda$
soit complètement décomposé.
\end{lemma}

\begin{proof}
Puisque $Y_0$ est quasi-compact et quasi-séparé, le schéma $Y$,
qui est affine sur $Y_0$, est quasi-compact et quasi-séparé.
Choisissons $n \geq 0$ et $Z_i \to Y$, $i = 1, \ldots, n$, comme dans le
lemme \ref{lemma-decompose}. Notons $\sigma_i : Z_i \to X$
des morphismes au-dessus de $Y$ qui existent par notre choix des $Z_i$.
Quitte à remplacer $0 \in \Lambda$ par un indice plus grand, nous pouvons supposer qu'il existe
des morphismes $Z_{i, 0} \to Y_0$ de présentation finie
dont les changements de base à $S$ sont les morphismes $Z_i \to Y$, voir
Limites, lemme \ref{limits-lemma-descend-finite-presentation}.
D'après Limites, lemme \ref{limits-lemma-descend-immersion},
nous pouvons supposer, quitte à remplacer $0$ par un indice plus grand, que $Z_{i, 0} \to Y_0$
est une immersion. Puisque $\coprod Z_i \to Y$ est surjectif, nous pouvons supposer,
quitte à remplacer $0$ par un indice plus grand, que
$\coprod Z_{i, 0} \to Y_0$ est surjectif, voir
Limites, lemme \ref{limits-lemma-descend-surjective}. Observons que
$Z_i = \lim_{\lambda \geq 0} Z_{i, \lambda}$,
où $Z_{i, \lambda} = Y_\lambda \times_{Y_0} Z_{i, 0}$.
Considérons les composés
$$
Z_i \xrightarrow{\sigma_i} X \to X_0
$$
comme des morphismes au-dessus de $Y_0$. Puisque $f_0$ est localement de présentation
finie, Limites, proposition
\ref{limits-proposition-characterize-locally-finite-presentation}
nous permet de trouver un $\lambda \geq 0$ tel qu'il existe des
morphismes $\sigma'_{i, \lambda} : Z_{i, \lambda} \to X_0$
au-dessus de $Y_0$ dont les précomposés avec $Z_i \to Z_{i, \lambda}$
sont les flèches affichées. Bien entendu, $\sigma'_{i, \lambda}$
détermine alors un morphisme $\sigma_{i, \lambda} : Z_{i, \lambda} \to
X_\lambda = X_0 \times_{Y_0} Y_\lambda$ au-dessus de $Y_\lambda$.
Puisque $\coprod Z_{i, \lambda} \to Y_\lambda$ est surjectif,
nous concluons que $X_\lambda \to Y_\lambda$ est complètement décomposé.
\end{proof}




\section{Familles de modules inversibles amples}
\label{section-families-ample-invertible-modules}

\noindent
Nous poursuivons la discussion de
Morphismes, section \ref{morphisms-section-families-ample-invertible-modules}.

\begin{lemma}
\label{lemma-ample-family-ample-relative}
Soit $f : X \to Y$ un morphisme de schémas. Supposons que
\begin{enumerate}
\item $Y$ possède une famille ample de modules inversibles,
\item il existe un module inversible $f$-ample sur $X$.
\end{enumerate}
Alors $X$ possède une famille ample de modules inversibles.
\end{lemma}

\begin{proof}
Soit $\mathcal{L}$ un module inversible $f$-ample sur $X$.
Cela implique en particulier que $f$ est quasi-compact, voir
Morphismes, définition \ref{morphisms-definition-relatively-ample}.
Puisque $Y$ est quasi-compact d'après Morphismes, définition
\ref{morphisms-definition-family-ample-invertible-modules},
nous voyons que $X$ est quasi-compact (et que $X$ lui-même satisfait donc la
première condition de Morphismes, définition
\ref{morphisms-definition-family-ample-invertible-modules}).
Soit $x \in X$ d'image $y \in Y$. Par l'hypothèse (2), nous pouvons
trouver un $\mathcal{O}_Y$-module inversible $\mathcal{N}$ et une section
$t \in \Gamma(Y, \mathcal{N})$ tels que le lieu $Y_t$ où $t$
ne s'annule pas soit affine. Alors $\mathcal{L}$ est ample sur
$f^{-1}(Y_t) = X_{f^*t}$, et nous pouvons donc trouver une section
$s \in \Gamma(X_{f^*t}, \mathcal{L})$ telle que $(X_{f^*t})_s$ soit affine
et contienne $x$. D'après
Propriétés, lemme \ref{properties-lemma-invert-s-sections},
pour un certain $n \geq 0$, le produit $(f^*t)^n s$ se prolonge en une section
$s' \in \Gamma(X, f^*\mathcal{N}^{\otimes n} \otimes \mathcal{L})$.
Enfin, la section $s'' = f^* ts'$ de
$f^*\mathcal{N}^{\otimes n + 1} \otimes \mathcal{L}$
s'annule en tout point de $X \setminus X_{f^*t}$ ; ainsi,
nous voyons que $X_{s''} = (X_{f^*t})_s$ est affine, comme voulu.
\end{proof}

\begin{lemma}
\label{lemma-resolution-property-goes-up-affine}
Soit $f : X \to Y$ un morphisme affine ou quasi-affine de schémas.
Si $Y$ possède une famille ample de modules inversibles, il en va de même pour $X$.
\end{lemma}

\begin{proof}
D'après Morphismes, lemme \ref{morphisms-lemma-quasi-affine-O-ample},
c'est un cas particulier du
lemme \ref{lemma-ample-family-ample-relative}.
\end{proof}






\section{Éclatements et familles amples de modules inversibles}
\label{section-ample-family-by-blowing-up}

\noindent
Nous démontrons un résultat de \cite{Gross-thesis}.

\begin{lemma}
\label{lemma-get-ample-family}
Soit $X$ un schéma. Supposons donnés des diviseurs de Cartier effectifs
$D_1, \ldots, D_m$ sur $X$ et des modules inversibles
$\mathcal{L}_1, \ldots, \mathcal{L}_m$ tels que
$\bigcap D_i = \emptyset$ et que $\mathcal{L}_i|_{X \setminus D_i}$
soit ample. Alors $X$ possède une famille ample de modules inversibles.
\end{lemma}

\begin{proof}
Soit $x \in X$. Choisissons un indice $i \in \{1, \ldots, m\}$
tel que $x \not \in D_i$. Posons $U_i = X \setminus D_i$.
Puisque $\mathcal{L}_i|_{U_i}$ est ample, nous pouvons trouver un $n \geq 1$ et
une section $s \in \Gamma(U_i, \mathcal{L}_i^{\otimes n})$
telle que le lieu $(U_i)_s$ où $s$ ne s'annule pas soit affine
(Propriétés, définition \ref{properties-definition-ample}).
Puisque $U_i$ est le lieu où la section canonique
$1 \in \mathcal{O}_X(D_i)$ ne s'annule pas, nous déduisons de
Propriétés, lemme \ref{properties-lemma-invert-s-sections}
qu'il existe un $N \geq 0$ tel que $s$ se prolonge en une section
$$
s' \in \Gamma(X, \mathcal{L}_i^{\otimes n} \otimes_{\mathcal{O}_X}
\mathcal{O}_X(N D_i))
$$
Après avoir remplacé $N$ par $N + 1$, nous voyons que $s'$ s'annule en tout
point de $D_i$, et donc que $X_{s'} = (U_i)_s$ est affine.
Cela prouve que $X$ possède une famille ample de modules inversibles, voir
Morphismes, définition
\ref{morphisms-definition-family-ample-invertible-modules}.
\end{proof}

\begin{lemma}
\label{lemma-blow-up-ample-family}
\begin{reference}
\cite[Proposition 1.3.1]{Gross-thesis}
\end{reference}
Soit $X$ un schéma quasi-compact et quasi-séparé n'ayant qu'un nombre fini de
composantes irréductibles. Il existe un ouvert dense quasi-compact
$U \subset X$ et un éclatement $U$-admissible $X' \to X$ tels que le
schéma $X'$ possède une famille ample de modules inversibles.
\end{lemma}

\begin{proof}
Soient $\eta_1, \ldots, \eta_n \in X$ les points génériques des
composantes irréductibles de $X$. D'après
Propriétés, lemme \ref{properties-lemma-point-and-maximal-points-affine}
et la quasi-compacité de $X$, nous pouvons trouver un recouvrement ouvert affine
fini $X = U_1 \cup \ldots \cup U_m$ tel que
chaque $U_i$ contienne $\eta_1, \ldots, \eta_n$. En particulier,
le sous-ensemble ouvert quasi-compact $U = U_1 \cap \ldots \cap U_m$ est dense dans $X$.
Soit $\mathcal{I}_i \subset \mathcal{O}_X$ un faisceau quasi-cohérent d'idéaux de type
fini tel que $U_i = X \setminus Z_i$, où $Z_i = V(\mathcal{I}_i)$,
voir Propriétés, lemme \ref{properties-lemma-quasi-coherent-finite-type-ideals}.
Soit
$$
f : X' \longrightarrow X
$$
l'éclatement de $X$ le long du faisceau d'idéaux
$\mathcal{I} = \mathcal{I}_1 \cdots \mathcal{I}_m$.
Notons que $f$ est un éclatement $U$-admissible, puisque $V(\mathcal{I})$
est (ensemblistement) la réunion des $Z_i$, qui sont disjoints de $U$.
De plus, $f$ est un morphisme projectif et
$\mathcal{O}_{X'}(1)$ est relativement ample pour $f$, voir
Diviseurs, lemme \ref{divisors-lemma-blowing-up-projective}.
D'après Diviseurs, lemme \ref{divisors-lemma-blowing-up-two-ideals},
pour chaque $i$, le morphisme $f'$ se factorise en $X' \to X'_i \to X$,
où $X'_i \to X$ est l'éclatement le long de $\mathcal{I}_i$
et $X' \to X'_i$ est un autre éclatement (à savoir, le long du tiré en arrière
du produit des idéaux $\mathcal{I}_j$, en omettant $\mathcal{I}_i$).
Il en résulte que $D_i = f^{-1}(Z_i) \subset X'$ est un diviseur de Cartier
effectif, voir
Diviseurs, lemmes \ref{divisors-lemma-blow-up-pullback-effective-Cartier} et
\ref{divisors-lemma-blowing-up-gives-effective-Cartier-divisor}.
Nous avons $X' \setminus D_i = f^{-1}(U_i)$. Comme $\mathcal{O}_{X'}(1)$ est
$f$-ample, la restriction de
$\mathcal{O}_{X'}(1)$ à $X' \setminus D_i$ est ample. Il résulte du
lemme \ref{lemma-get-ample-family}
que $X'$ possède une famille ample de modules inversibles.
\end{proof}

\begin{proposition}
\label{proposition-envelope-with-resolution-property}
Soit $X$ un schéma quasi-compact et quasi-séparé. Il existe un
morphisme $f : Y \to X$ qui est de présentation finie, propre et
complètement décomposé (définition \ref{definition-cd-morphism}),
tel que le schéma $Y$ possède une famille ample de modules inversibles.
\end{proposition}

\begin{proof}
D'après Limites, proposition \ref{limits-proposition-approximate},
il existe un morphisme affine $X \to X_0$, où $X_0$
est un schéma de type fini sur $\mathbf{Z}$. Nous construisons ci-dessous
un morphisme $Y_0 \to X_0$ ayant toutes les propriétés voulues.
En posant alors $Y = X \times_{X_0} Y_0$ et en prenant pour $f$
la projection $f : Y \to X$, nous concluons.
En effet, $f$ est de présentation finie
(Morphismes, lemme \ref{morphisms-lemma-base-change-finite-presentation}),
$f$ est propre
(Morphismes, lemme \ref{morphisms-lemma-base-change-proper}),
$f$ est complètement décomposé
(lemme \ref{lemma-base-change-cd}). Enfin, puisque $Y \to Y_0$ est
affine (étant le changement de base de $X \to X_0$), nous voyons que $Y$ possède
une famille ample de modules inversibles d'après le
lemme \ref{lemma-resolution-property-goes-up-affine}.
Nous sommes ainsi ramenés au cas traité dans le paragraphe suivant.

\medskip\noindent
Supposons $X$ de type fini sur $\mathbf{Z}$. En particulier,
$\dim(X) < \infty$. Nous allons raisonner par récurrence sur $\dim(X)$.
Si $\dim(X) = 0$, alors $X$ est affine et possède la propriété de résolution.
En général, il existe un ouvert dense $U \subset X$ et un
éclatement $U$-admissible $X' \to X$ tel que $X'$ possède
une famille ample de modules inversibles, voir
le lemme \ref{lemma-blow-up-ample-family}.
Puisque $f : X' \to X$ est un isomorphisme au-dessus de $U$,
nous voyons que tout point de $U$ se relève en un point de $X'$
ayant le même corps résiduel.
Soit $Z = X \setminus U$, muni de la structure de sous-schéma réduit induite.
Alors $\dim(Z) < \dim(X)$, puisque $U$ est dense dans $X$ (voir ci-dessus).
Par récurrence, nous trouvons un morphisme propre, complètement décomposé
$W \to Z$ tel que $W$ possède une famille ample de modules inversibles.
Il s'ensuit alors que $Y = X' \amalg W \to X$ est le morphisme
voulu.
\end{proof}




\section{Le critère extensif pour les immersions fermées}
\label{section-extensive-criterion-closed-immersions}

\noindent
Dans cette section, nous donnons un critère pour qu'un morphisme de schémas 
soit une immersion fermée.

\begin{lemma}
\label{lemma-affine-extensive-criterion}
Un morphisme $f : X \to Y$ de schémas affines est une immersion fermée 
si et seulement si, pour tout homomorphisme injectif d'anneaux $A \to B$ et tout 
carré commutatif
$$
\xymatrix{
\Spec(B) \ar[d] \ar[r] & X \ar[d]^f \\
\Spec(A) \ar[r] \ar@{..>}[ur] & Y
}
$$
il existe un relèvement $\Spec(A) \to X$ rendant les deux triangles commutatifs.
\end{lemma}

\begin{proof}
Soit le morphisme $f$ donné par l'homomorphisme d'anneaux $\phi : R \to S$.
Alors $f$ est une immersion fermée si et seulement si $\phi$ est surjectif.

\medskip\noindent
Supposons tout d'abord que $\phi$ soit surjectif.
Soit $\psi : A \to B$ un homomorphisme injectif d'anneaux, et supposons donné
un diagramme commutatif
$$
\xymatrix{
R \ar[r]^\alpha \ar[d]^\phi & A \ar[d]^\psi \\
S \ar[r]^\beta \ar@{..>}[ur] & B
}
$$
Nous définissons alors un relèvement $S \to A$ par $s \mapsto \alpha(r)$, où
$r \in R$ est tel que $\phi(r) = s$.
Cette définition est bien posée, car $\psi$ est injectif et le carré est commutatif.
Comme le passage au spectre établit une anti-équivalence entre les anneaux
commutatifs et les schémas affines, la propriété de relèvement voulue pour
$f$ est vérifiée.

\medskip\noindent
Supposons ensuite que $\phi$ admette des relèvements par rapport à tous les
homomorphismes injectifs d'anneaux $\psi: A \to B$.
Remarquons que $\phi(R)$ est un sous-anneau de $S$ ; on obtient donc un carré
commutatif
$$
\xymatrix{
R \ar[r] \ar[d]^\phi  & \phi(R) \ar[d] \\
S \ar@{=}[r] \ar@{..>}[ur] & S
}
$$
dans lequel il existe un relèvement $S \to \phi(R)$.
Par suite, l'inclusion $\phi(R) \to S$ est nécessairement un isomorphisme, ce qui
montre que $\phi$ est surjectif, et le résultat s'ensuit.
\end{proof}

\begin{lemma}
\label{lemma-scheme-affine-if-aff-is-split-mono}
Soit $X$ un schéma.
Si le morphisme canonique $X \to \Spec(\Gamma(X, \mathcal{O}_X))$
considéré dans Schémas, lemme \ref{schemes-lemma-morphism-into-affine},
admet une rétraction, alors $X$ est un schéma affine.
\end{lemma}

\begin{proof}
Posons $S = \Spec(\Gamma(X, \mathcal{O}_X))$ et notons $f : X \to S$ le morphisme
donné dans le lemme. Soit $s : S \to X$ une rétraction ; ainsi $\text{id}_X = sf$.
Alors $f s f = \text{id}_S f$. Comme $f$ induit un isomorphisme
$\Gamma(S, \mathcal{O}_S) \to \Gamma(X, \mathcal{O}_X)$,
cela signifie que $fs$ et $\text{id}_S$ induisent la même application
sur $\Gamma(S, \mathcal{O}_S)$. D'où $f s = \text{id}_S$, puisque $S$
est affine. Par suite, $f$ est un isomorphisme et $X$ est un schéma affine,
comme il fallait le démontrer.
\end{proof}

\begin{lemma}
\label{lemma-aff-is-injective-on-sheaves}
Soit $X$ un schéma. Soit $f : X \to S = \Spec(\Gamma(X, \mathcal{O}_X))$
le morphisme canonique considéré dans
Schémas, lemme \ref{schemes-lemma-morphism-into-affine}.
Le plus grand $\mathcal{O}_S$-module quasi-cohérent contenu
dans le noyau de $f^\sharp : \mathcal{O}_S \to f_*\mathcal{O}_X$
est nul. Si $X$ est quasi-compact, alors $f^\sharp$ est injectif.
En particulier, si $X$ est quasi-compact, alors $f$ est un
morphisme dominant.
\end{lemma}

\begin{proof}
Soit $M \subset \Gamma(S, \mathcal{O}_S)$ le sous-module correspondant
au plus grand $\mathcal{O}_S$-module quasi-cohérent contenu
dans le noyau de $f^\sharp$. Alors tout élément $a \in M$ est envoyé
sur zéro par $f^\sharp$. Or $f^\sharp(a)$ est l'élément de
$$
\Gamma(S, f_*\mathcal{O}_X) = \Gamma(X, \mathcal{O}_X) =
\Gamma(S, \mathcal{O}_S)
$$
qui correspond à $a$ lui-même ! Ainsi $a = 0$. Par conséquent $M = 0$, ce qui prouve
la première assertion. Remarquons que cela équivaut à dire que le morphisme
$f : X \to S$ est schématiquement surjectif.

\medskip\noindent
Si $X$ est quasi-compact, alors $\Ker(f^\sharp)$ est quasi-cohérent d'après
Morphismes, lemme \ref{morphisms-lemma-quasi-compact-scheme-theoretic-image}.
Ainsi $\Ker(f^\sharp) = 0$ et $f^\sharp$ est injectif.
Dans ce cas, $f$ est un morphisme dominant en vertu de la partie (4) de
Morphismes, lemme \ref{morphisms-lemma-quasi-compact-scheme-theoretic-image}.
\end{proof}

\begin{lemma}
\label{lemma-extensive-criterion}
Soit $f: X \to Y$ un morphisme quasi-compact de schémas.
Alors $f$ est une immersion fermée si et seulement si, pour tout homomorphisme
injectif d'anneaux $A \to B$ et tout carré commutatif
$$
\xymatrix{
\Spec(B) \ar[d] \ar[r] & X \ar[d]^f \\
\Spec(A) \ar[r] \ar@{..>}[ur] & Y
}
$$
il existe un relèvement $\Spec A \to X$ rendant le diagramme commutatif.
\end{lemma}

\begin{proof}
Supposons que $f$ soit une immersion fermée. Soit $A \to B$ un homomorphisme injectif
d'anneaux et considérons un carré commutatif
$$
\xymatrix{
\Spec(B) \ar[d] \ar[r]  & X \ar[d]^f \\
\Spec(A) \ar[r] \ar@{..>}[ur] & Y
}
$$
Alors $\Spec(A) \times_Y X \to \Spec(A)$ est une immersion fermée et
il existe donc un idéal $I \subset A$ et un diagramme commutatif
$$
\xymatrix{
\Spec(B) \ar[d] \ar[r] & \Spec(A/I) \ar[r] \ar[d] & X \ar[d]^f \\
\Spec(A) \ar[r] \ar@{..>}[ur] & \Spec(A) \ar[r] & Y
}
$$
Le lemme \ref{lemma-affine-extensive-criterion} fournit un relèvement.

\medskip\noindent
Supposons que $f$ satisfasse à la propriété de relèvement énoncée dans le lemme.
La propriété pour $f$ d'être une immersion fermée est locale sur $Y$ ;
nous pouvons donc supposer, ce que nous ferons, que $Y$ est affine.
En particulier, $Y$ est quasi-compact, et $X$ est donc quasi-compact.
Il existe donc un recouvrement fini par des ouverts affines $X = U_1 \cup \ldots \cup U_n$.
La source du morphisme
$$
\pi : U = \coprod U_i \longrightarrow X
$$
est affine et l'homomorphisme d'anneaux induit
$\Gamma(X, \mathcal{O}_X) \to \Gamma(U, \mathcal{O}_U)$
est injectif. Par hypothèse, il existe un relèvement dans le diagramme
$$
\xymatrix{
U \ar[r]^\pi \ar[d] & X \ar[d]^f \\
\Spec(\Gamma(X, \mathcal{O}_X)) \ar[r]^-{f'} \ar@{..>}[ur]^h & Y
}
$$
où $f'$ est le morphisme de schémas affines correspondant à
l'homomorphisme d'anneaux $\Gamma(Y, \mathcal{O}_Y) \to \Gamma(X, \mathcal{O}_X)$.
Il résulte du fait que $\pi$ est un épimorphisme
que le morphisme $h$ est une rétraction du morphisme canonique
$X \to \Spec(\Gamma(X, \mathcal{O}_X))$ ; nous omettons les détails. Ainsi $X$ est affine d'après le
lemme \ref{lemma-scheme-affine-if-aff-is-split-mono}.
Le lemme \ref{lemma-affine-extensive-criterion} permet alors de
conclure que $f$ est une immersion fermée.
\end{proof}





\begin{multicols}{2}[\section{Autres chapitres}]
\noindent
Préliminaires
\begin{enumerate}
\item \hyperref[introduction-section-phantom]{Introduction}
\item \hyperref[conventions-section-phantom]{Conventions}
\item \hyperref[sets-section-phantom]{Théorie des ensembles}
\item \hyperref[categories-section-phantom]{Catégories}
\item \hyperref[topology-section-phantom]{Topologie}
\item \hyperref[sheaves-section-phantom]{Faisceaux sur les espaces}
\item \hyperref[sites-section-phantom]{Sites et faisceaux}
\item \hyperref[stacks-section-phantom]{Champs}
\item \hyperref[fields-section-phantom]{Corps}
\item \hyperref[algebra-section-phantom]{Algèbre commutative}
\item \hyperref[brauer-section-phantom]{Groupes de Brauer}
\item \hyperref[homology-section-phantom]{Algèbre homologique}
\item \hyperref[derived-section-phantom]{Catégories dérivées}
\item \hyperref[simplicial-section-phantom]{Méthodes simpliciales}
\item \hyperref[more-algebra-section-phantom]{Compléments d'algèbre}
\item \hyperref[smoothing-section-phantom]{Lissage des morphismes d'anneaux}
\item \hyperref[modules-section-phantom]{Faisceaux de modules}
\item \hyperref[sites-modules-section-phantom]{Modules sur les sites}
\item \hyperref[injectives-section-phantom]{Objets injectifs}
\item \hyperref[cohomology-section-phantom]{Cohomologie des faisceaux}
\item \hyperref[sites-cohomology-section-phantom]{Cohomologie sur les sites}
\item \hyperref[dga-section-phantom]{Algèbre différentielle graduée}
\item \hyperref[dpa-section-phantom]{Algèbre à puissances divisées}
\item \hyperref[sdga-section-phantom]{Faisceaux différentiels gradués}
\item \hyperref[hypercovering-section-phantom]{Hyperrecouvrements}
\end{enumerate}
Schémas
\begin{enumerate}
\setcounter{enumi}{25}
\item \hyperref[schemes-section-phantom]{Schémas}
\item \hyperref[constructions-section-phantom]{Constructions de schémas}
\item \hyperref[properties-section-phantom]{Propriétés des schémas}
\item \hyperref[morphisms-section-phantom]{Morphismes de schémas}
\item \hyperref[coherent-section-phantom]{Cohomologie des schémas}
\item \hyperref[divisors-section-phantom]{Diviseurs}
\item \hyperref[limits-section-phantom]{Limites de schémas}
\item \hyperref[varieties-section-phantom]{Variétés}
\item \hyperref[topologies-section-phantom]{Topologies sur les schémas}
\item \hyperref[descent-section-phantom]{Descente}
\item \hyperref[perfect-section-phantom]{Catégories dérivées de schémas}
\item \hyperref[more-morphisms-section-phantom]{Compléments sur les morphismes}
\item \hyperref[flat-section-phantom]{Compléments sur la platitude}
\item \hyperref[groupoids-section-phantom]{Schémas en groupoïdes}
\item \hyperref[more-groupoids-section-phantom]{Compléments sur les schémas en groupoïdes}
\item \hyperref[etale-section-phantom]{Morphismes étales de schémas}
\end{enumerate}
Thèmes de théorie des schémas
\begin{enumerate}
\setcounter{enumi}{41}
\item \hyperref[chow-section-phantom]{Homologie de Chow}
\item \hyperref[intersection-section-phantom]{Théorie de l'intersection}
\item \hyperref[pic-section-phantom]{Schémas de Picard des courbes}
\item \hyperref[weil-section-phantom]{Théories cohomologiques de Weil}
\item \hyperref[adequate-section-phantom]{Modules adéquats}
\item \hyperref[dualizing-section-phantom]{Complexes dualisants}
\item \hyperref[duality-section-phantom]{Dualité pour les schémas}
\item \hyperref[discriminant-section-phantom]{Discriminants et différentes}
\item \hyperref[derham-section-phantom]{Cohomologie de de Rham}
\item \hyperref[local-cohomology-section-phantom]{Cohomologie locale}
\item \hyperref[algebraization-section-phantom]{Géométrie algébrique et formelle}
\item \hyperref[curves-section-phantom]{Courbes algébriques}
\item \hyperref[resolve-section-phantom]{Résolution des surfaces}
\item \hyperref[models-section-phantom]{Réduction semi-stable}
\item \hyperref[functors-section-phantom]{Foncteurs et morphismes}
\item \hyperref[equiv-section-phantom]{Catégories dérivées de variétés}
\item \hyperref[pione-section-phantom]{Groupes fondamentaux de schémas}
\item \hyperref[etale-cohomology-section-phantom]{Cohomologie étale}
\item \hyperref[crystalline-section-phantom]{Cohomologie cristalline}
\item \hyperref[proetale-section-phantom]{Cohomologie pro-étale}
\item \hyperref[relative-cycles-section-phantom]{Cycles relatifs}
\item \hyperref[more-etale-section-phantom]{Compléments de cohomologie étale}
\item \hyperref[trace-section-phantom]{La formule des traces}
\end{enumerate}
Espaces algébriques
\begin{enumerate}
\setcounter{enumi}{64}
\item \hyperref[spaces-section-phantom]{Espaces algébriques}
\item \hyperref[spaces-properties-section-phantom]{Propriétés des espaces algébriques}
\item \hyperref[spaces-morphisms-section-phantom]{Morphismes d'espaces algébriques}
\item \hyperref[decent-spaces-section-phantom]{Espaces algébriques décents}
\item \hyperref[spaces-cohomology-section-phantom]{Cohomologie des espaces algébriques}
\item \hyperref[spaces-limits-section-phantom]{Limites d'espaces algébriques}
\item \hyperref[spaces-divisors-section-phantom]{Diviseurs sur les espaces algébriques}
\item \hyperref[spaces-over-fields-section-phantom]{Espaces algébriques sur des corps}
\item \hyperref[spaces-topologies-section-phantom]{Topologies sur les espaces algébriques}
\item \hyperref[spaces-descent-section-phantom]{Descente et espaces algébriques}
\item \hyperref[spaces-perfect-section-phantom]{Catégories dérivées d'espaces}
\item \hyperref[spaces-more-morphisms-section-phantom]{Compléments sur les morphismes d'espaces}
\item \hyperref[spaces-flat-section-phantom]{Platitude sur les espaces algébriques}
\item \hyperref[spaces-groupoids-section-phantom]{Groupoïdes dans les espaces algébriques}
\item \hyperref[spaces-more-groupoids-section-phantom]{Compléments sur les groupoïdes dans les espaces}
\item \hyperref[bootstrap-section-phantom]{Amorçage}
\item \hyperref[spaces-pushouts-section-phantom]{Sommes amalgamées d'espaces algébriques}
\end{enumerate}
Thèmes de géométrie
\begin{enumerate}
\setcounter{enumi}{81}
\item \hyperref[spaces-chow-section-phantom]{Groupes de Chow des espaces}
\item \hyperref[groupoids-quotients-section-phantom]{Quotients de groupoïdes}
\item \hyperref[spaces-more-cohomology-section-phantom]{Compléments sur la cohomologie des espaces}
\item \hyperref[spaces-simplicial-section-phantom]{Espaces simpliciaux}
\item \hyperref[spaces-duality-section-phantom]{Dualité pour les espaces}
\item \hyperref[formal-spaces-section-phantom]{Espaces algébriques formels}
\item \hyperref[restricted-section-phantom]{Algébrisation des espaces formels}
\item \hyperref[spaces-resolve-section-phantom]{Résolution des surfaces revisitée}
\end{enumerate}
Théorie des déformations
\begin{enumerate}
\setcounter{enumi}{89}
\item \hyperref[formal-defos-section-phantom]{Théorie formelle des déformations}
\item \hyperref[defos-section-phantom]{Théorie des déformations}
\item \hyperref[cotangent-section-phantom]{Le complexe cotangent}
\item \hyperref[examples-defos-section-phantom]{Problèmes de déformation}
\end{enumerate}
Champs algébriques
\begin{enumerate}
\setcounter{enumi}{93}
\item \hyperref[algebraic-section-phantom]{Champs algébriques}
\item \hyperref[examples-stacks-section-phantom]{Exemples de champs}
\item \hyperref[stacks-sheaves-section-phantom]{Faisceaux sur les champs algébriques}
\item \hyperref[criteria-section-phantom]{Critères de représentabilité}
\item \hyperref[artin-section-phantom]{Axiomes d'Artin}
\item \hyperref[quot-section-phantom]{Espaces de Quot et de Hilbert}
\item \hyperref[stacks-properties-section-phantom]{Propriétés des champs algébriques}
\item \hyperref[stacks-morphisms-section-phantom]{Morphismes de champs algébriques}
\item \hyperref[stacks-limits-section-phantom]{Limites de champs algébriques}
\item \hyperref[stacks-cohomology-section-phantom]{Cohomologie des champs algébriques}
\item \hyperref[stacks-perfect-section-phantom]{Catégories dérivées de champs}
\item \hyperref[stacks-introduction-section-phantom]{Introduction aux champs algébriques}
\item \hyperref[stacks-more-morphisms-section-phantom]{Compléments sur les morphismes de champs}
\item \hyperref[stacks-geometry-section-phantom]{La géométrie des champs}
\end{enumerate}
Thèmes de théorie des espaces de modules
\begin{enumerate}
\setcounter{enumi}{107}
\item \hyperref[moduli-section-phantom]{Champs de modules}
\item \hyperref[moduli-curves-section-phantom]{Espaces de modules de courbes}
\end{enumerate}
Divers
\begin{enumerate}
\setcounter{enumi}{109}
\item \hyperref[examples-section-phantom]{Exemples}
\item \hyperref[exercises-section-phantom]{Exercices}
\item \hyperref[guide-section-phantom]{Guide bibliographique}
\item \hyperref[desirables-section-phantom]{Desiderata}
\item \hyperref[coding-section-phantom]{Style de programmation}
\item \hyperref[obsolete-section-phantom]{Obsolète}
\item \hyperref[fdl-section-phantom]{Licence de documentation libre GNU}
\item \hyperref[index-section-phantom]{Index généré automatiquement}
\end{enumerate}
\end{multicols}


\bibliography{my}
\bibliographystyle{amsalpha}

\end{document}
